\setcounter{section}{0}
\section{Morphismes affines}
\label{section:II.1-fr}

La plupart des résultats de ce paragraphe sont les contreparties
<<~globales~>> de ceux du chapitre~I\textsuperscript{er}, \S~1~; ils ne sont
donc pas essentiellement nouveaux et fournissent simplement un langage
commode pour la suite.

\setcounter{subsection}{0}
\subsection{$S$-préschémas et $\mathcal{O}_S$-Algèbres}
\label{subsection:II.1.1-fr}

\begin{env}[1.1.1]
\label{II.1.1.1-fr}
Soient $S$ un préschéma, $X$ un $S$-préschéma, $f:X\to S$ son morphisme
structural. On sait (\hyperref[0.4.2.4-fr]{0, 4.2.4}) que l'image directe
$f_*(\mathcal{O}_X)$ est une $\mathcal{O}_S$-Algèbre, que nous
\oldpage[II]{6}
noterons $\mathcal{A}(X)$ lorsqu'il n'en résultera pas de confusion~; si $U$
est un ouvert de $S$, on a
\[
  \mathcal{A}(f^{-1}(U))=\mathcal{A}(X)|U.
\]
De même, pour tout $\mathcal{O}_X$-Module $\mathcal{F}$ (resp. toute
$\mathcal{O}_X$-Algèbre $\mathcal{B}$), nous écrirons
$\mathcal{A}(\mathcal{F})$ (resp. $\mathcal{A}(\mathcal{B})$) l'image directe
$f_*(\mathcal{F})$ (resp. $f_*(\mathcal{B})$) qui est un
$\mathcal{A}(X)$-Module (resp. une $\mathcal{A}(X)$-Algèbre), et non seulement
un $\mathcal{O}_S$-Module (resp. $\mathcal{O}_S$-Algèbre).
\end{env}

\begin{env}[1.1.2]
\label{II.1.1.2-fr}
Soient $Y$ un second $S$-préschéma, $g:Y\to S$ son morphisme structural,
$h:X\to Y$ un $S$-morphisme~; on a donc le diagramme commutatif
\[
  \xymatrix{
    X\ar[rr]^h\ar[rd]_f && Y\ar[ld]^g\\
    & S
  }
  \tag{1.1.2.1}\label{II.1.1.2.1-fr}
\]

On a par définition $h=(\psi,\theta)$, où
$\theta:\mathcal{O}_Y\to h_*(\mathcal{O}_X)=\psi_*(\mathcal{O}_X)$ est un
homomorphisme de faisceaux d'anneaux~; on en déduit
(\hyperref[0.4.2.2-fr]{0, 4.2.2}) un homomorphisme de
$\mathcal{O}_S$-Algèbres
\[
  g_*(\theta):g_*(\mathcal{O}_Y)\to
  g_*(h_*(\mathcal{O}_X))=f_*(\mathcal{O}_X),
\]
autrement dit, un homomorphisme de $\mathcal{O}_S$-Algèbres
$\mathcal{A}(Y)\to\mathcal{A}(X)$, que nous noterons aussi
$\mathcal{A}(h)$. Si $h':Y\to Z$ est un second $S$-morphisme, il est immédiat
que $\mathcal{A}(h'\circ h)=\mathcal{A}(h)\circ\mathcal{A}(h')$. Nous avons
donc défini un \emph{foncteur contravariant} $\mathcal{A}(X)$ sur la catégorie
des $S$-préschémas, à valeurs dans la catégorie des
$\mathcal{O}_S$-Algèbres.

Soient maintenant $\mathcal{F}$ un $\mathcal{O}_X$-Module, $\mathcal{G}$ un
$\mathcal{O}_Y$-Module, et $u:\mathcal{G}\to\mathcal{F}$ un $h$-morphisme,
c'est-à-dire (\hyperref[0.4.4.1-fr]{0, 4.4.1}) un homomorphisme de
$\mathcal{O}_Y$-Modules $\mathcal{G}\to h_*(\mathcal{F})$. Alors
\[
  g_*(u):g_*(\mathcal{G})\to g_*(h_*(\mathcal{F}))=f_*(\mathcal{F})
\]
est un homomorphisme $\mathcal{A}(\mathcal{G})\to
\mathcal{A}(\mathcal{F})$ de $\mathcal{O}_S$-Modules, que nous noterons
$\mathcal{A}(u)$~; en outre, le couple
$(\mathcal{A}(h),\mathcal{A}(u))$ constitue un \emph{di-homomorphisme} du
$\mathcal{A}(Y)$-Module $\mathcal{A}(\mathcal{G})$ dans le
$\mathcal{A}(X)$-Module $\mathcal{A}(\mathcal{F})$.
\end{env}

\begin{env}[1.1.3]
\label{II.1.1.3-fr}
Le préschéma $S$ étant fixé, on peut considérer les couples
$(X,\mathcal{F})$, où $X$ est un $S$-préschéma et $\mathcal{F}$ un
$\mathcal{O}_X$-Module, comme formant une \emph{catégorie}, en définissant un
\emph{morphisme} $(X,\mathcal{F})\to(Y,\mathcal{G})$ comme un couple $(h,u)$,
où $h:X\to Y$ est un $S$-morphisme et $u:\mathcal{G}\to\mathcal{F}$ un
$h$-morphisme. On peut alors dire que
$(\mathcal{A}(X),\mathcal{A}(\mathcal{F}))$ est un \emph{foncteur
contravariant} à valeurs dans la catégorie dont les objets sont les couples
formés d'une $\mathcal{O}_S$-Algèbre et d'un Module sur cette Algèbre, et les
morphismes sont les di-homomorphismes.
\end{env}

\subsection{Préschémas affines sur un préschéma}
\label{subsection:II.1.2-fr}

\begin{definition}[1.2.1]
\label{II.1.2.1-fr}
Soient $X$ un $S$-préschéma, $f:X\to S$ son morphisme structural. On dit que
$X$ est \emph{affine au-dessus de $S$} s'il existe un recouvrement
$(S_\alpha)$ de $S$ par des ouverts affines tels que, pour tout $\alpha$, le
préschéma induit par $X$ sur l'ensemble ouvert $f^{-1}(S_\alpha)$ soit affine.
\end{definition}

\begin{example}[1.2.2]
\label{II.1.2.2-fr}
Tout sous-préschéma fermé de $S$ est un $S$-préschéma affine au-dessus de $S$
(\hyperref[I.4.2.3-fr]{I, 4.2.3} et
\hyperref[I.4.2.4-fr]{4.2.4}).
\end{example}

\begin{remark}[1.2.3]
\label{II.1.2.3-fr}
Un préschéma $X$ affine au-dessus de $S$ n'est pas nécessairement un schéma
affine, comme le montre l'exemple $X=S$ (\hyperref[II.1.2.2-fr]{1.2.2}).
D'autre part, si un schéma affine $X$ est un $S$-préschéma, $X$ n'est pas
nécessairement affine au-dessus
\oldpage[II]{7}
de $S$ (voir (\hyperref[II.1.3.3-fr]{1.3.3})). Toutefois, rappelons que si
$S$ est un \emph{schéma}, tout $S$-préschéma qui est un schéma affine est
affine au-dessus de $S$ (\hyperref[I.5.5.10-fr]{I, 5.5.10}).
\end{remark}

\begin{proposition}[1.2.4]
\label{II.1.2.4-fr}
Tout $S$-préschéma qui est affine au-dessus de $S$ est séparé au-dessus de
$S$ (autrement dit, est un $S$-schéma).
\end{proposition}

Cela résulte aussitôt de (\hyperref[I.5.5.5-fr]{I, 5.5.5}) et
(\hyperref[I.5.5.8-fr]{I, 5.5.8}).

\begin{proposition}[1.2.5]
\label{II.1.2.5-fr}
Soient $X$ un $S$-schéma affine au-dessus de $S$, $f:X\to S$ le morphisme
structural. Pour tout ouvert $U\subset S$, $f^{-1}(U)$ est affine au-dessus
de $U$.
\end{proposition}

En vertu de la déf. (\hyperref[II.1.2.1-fr]{1.2.1}), on est ramené au cas où
$S=\operatorname{Spec}(A)$ et $X=\operatorname{Spec}(B)$ sont affines et
alors $f=({}^a\varphi,\widetilde{\varphi})$, où $\varphi:A\to B$ est un
homomorphisme. Comme les $D(g)$, où $g\in A$, forment une base de $S$, on est
ramené au cas où $U=D(g)$~; mais on sait alors que
$f^{-1}(U)=D(\varphi(g))$ (\hyperref[I.1.2.2.2-fr]{I, 1.2.2.2}), d'où la
proposition.

\begin{proposition}[1.2.6]
\label{II.1.2.6-fr}
Soient $X$ un $S$-schéma affine au-dessus de $S$, $f:X\to S$ le morphisme
structural. Pour tout $\mathcal{O}_X$-Module quasi-cohérent $\mathcal{F}$,
$f_*(\mathcal{F})$ est un $\mathcal{O}_S$-Module quasi-cohérent.
\end{proposition}

Compte tenu de (\hyperref[II.1.2.4-fr]{1.2.4}), cela résulte de
(\hyperref[I.9.2.2-fr]{I, 9.2.2, a)}).

En particulier, la $\mathcal{O}_S$-Algèbre
$\mathcal{A}(X)=f_*(\mathcal{O}_X)$ est \emph{quasi-cohérente}.

\begin{proposition}[1.2.7]
\label{II.1.2.7-fr}
Soient $X$ un $S$-schéma affine au-dessus de $S$. Pour tout $S$-préschéma
$Y$, l'application $h\mapsto\mathcal{A}(h)$ de l'ensemble
$\operatorname{Hom}_S(Y,X)$ dans l'ensemble
$\operatorname{Hom}(\mathcal{A}(X),\mathcal{A}(Y))$
(\hyperref[II.1.1.2-fr]{1.1.2}) est bijective.
\end{proposition}

Soient $f:X\to S$, $g:Y\to S$ les morphismes structuraux. Supposons d'abord
$S=\operatorname{Spec}(A)$ et $X=\operatorname{Spec}(B)$ affines~; il faut
prouver que pour tout homomorphisme
$\omega:f_*(\mathcal{O}_X)\to g_*(\mathcal{O}_Y)$ de
$\mathcal{O}_S$-Algèbres, il existe un $S$-morphisme $h:Y\to X$ et un seul
tel que $\mathcal{A}(h)=\omega$. Par définition, pour tout ouvert $U\subset
S$, $\omega$ définit un homomorphisme
\[
  \omega_U=\Gamma(U,\omega):
  \Gamma(f^{-1}(U),\mathcal{O}_X)\to
  \Gamma(g^{-1}(U),\mathcal{O}_Y)
\]
de $\Gamma(U,\mathcal{O}_S)$-algèbres. En particulier, pour $U=S$, cela
donne un homomorphisme
$\varphi:\Gamma(X,\mathcal{O}_X)\to\Gamma(Y,\mathcal{O}_Y)$ de
$\Gamma(S,\mathcal{O}_S)$-algèbres, auquel correspond un $S$-morphisme
$h:Y\to X$ bien déterminé, puisque $X$ est affine
(\hyperref[I.2.2.4-fr]{I, 2.2.4}). Reste à prouver que
$\mathcal{A}(h)=\omega$, autrement dit, que pour tout ouvert $U$ d'une base
de $S$, $\omega_U$ coïncide avec l'homomorphisme d'algèbres $\varphi_U$
correspondant au $S$-morphisme $g^{-1}(U)\to f^{-1}(U)$ restriction de $h$.
On peut se borner au cas où $U=D(\lambda)$, avec $\lambda\in S$~; alors, si
$f=({}^a\rho,\widetilde{\rho})$, où $\rho:A\to B$ est un homomorphisme
d'anneaux, on a $f^{-1}(U)=D(\mu)$, où $\mu=\rho(\lambda)$, et
$\Gamma(f^{-1}(U),\mathcal{O}_X)$ est l'anneau de fractions $B_\mu$~; or,
le diagramme
\[
  \xymatrix{
    B\ar[r]^\varphi\ar[d] & \Gamma(Y,\mathcal{O}_Y)\ar[d]\\
    B_\mu\ar[r]_{\varphi_U} & \Gamma(g^{-1}(U),\mathcal{O}_Y)
  }
\]
est commutatif, et il en est de même du diagramme analogue où $\varphi_U$
est remplacé par $\omega_U$~; l'égalité $\varphi_U=\omega_U$ résulte alors
de la propriété universelle des anneaux de fractions
(\hyperref[0.1.2.4-fr]{0, 1.2.4}).

Passons au cas général, et soit $(S_\alpha)$ un recouvrement de $S$ par des
ouverts affines
\oldpage[II]{8}
tels que les $f^{-1}(S_\alpha)$ soient affines. Alors tout homomorphisme
$\omega:\mathcal{A}(X)\to\mathcal{A}(Y)$ de $\mathcal{O}_S$-Algèbres donne
par restriction une famille d'homomorphismes
\[
  \omega_\alpha:
  \mathcal{A}(f^{-1}(S_\alpha))\to
  \mathcal{A}(g^{-1}(S_\alpha))
\]
de $\mathcal{O}_{S_\alpha}$-Algèbres, donc une famille de
$S_\alpha$-morphismes
$h_\alpha:g^{-1}(S_\alpha)\to f^{-1}(S_\alpha)$ d'après ce qui précède.
Tout revient à voir que pour tout ouvert affine $U$ d'une base de
$S_\alpha\cap S_\beta$, les restrictions de $h_\alpha$ et de $h_\beta$ à
$g^{-1}(U)$ coïncident, ce qui est évident, puisque, en vertu de ce qui
précède, ces restrictions correspondent toutes deux à l'homomorphisme
$\mathcal{A}(X)|U\to\mathcal{A}(Y)|U$ restriction de $\omega$.

\begin{corollary}[1.2.8]
\label{II.1.2.8-fr}
Soient $X$, $Y$ deux $S$-schémas qui sont affines au-dessus de $S$. Pour
qu'un $S$-morphisme $h:Y\to X$ soit un isomorphisme, il faut et il suffit que
$\mathcal{A}(h):\mathcal{A}(X)\to\mathcal{A}(Y)$ soit un isomorphisme.
\end{corollary}

Cela résulte aussitôt de (\hyperref[II.1.2.7-fr]{1.2.7}) et du caractère
fonctoriel de $\mathcal{A}(X)$.

\subsection{Préschéma affine au-dessus de $S$ associé à une
$\mathcal{O}_S$-Algèbre}
\label{subsection:II.1.3-fr}

\begin{proposition}[1.3.1]
\label{II.1.3.1-fr}
Soit $S$ un préschéma. Pour toute $\mathcal{O}_S$-Algèbre quasi-cohérente
$\mathcal{B}$, il existe un préschéma $X$ affine au-dessus de $S$, défini à
un $S$-isomorphisme unique près, tel que $\mathcal{A}(X)=\mathcal{B}$.
\end{proposition}

L'unicité résulte de (\hyperref[II.1.2.8-fr]{1.2.8})~; démontrons l'existence
de $X$. Pour tout ouvert affine $U\subset S$, soit $X_U$ le préschéma
$\operatorname{Spec}(\Gamma(U,\mathcal{B}))$~; comme
$\Gamma(U,\mathcal{B})$ est une $\Gamma(U,\mathcal{O}_S)$-algèbre, $X_U$
est un $S$-préschéma (\hyperref[I.1.6.1-fr]{I, 1.6.1}). En outre, comme
$\mathcal{B}$ est quasi-cohérente, la $\mathcal{O}_S$-Algèbre
$\mathcal{A}(X_U)$ s'identifie canoniquement à $\mathcal{B}|U$
(\hyperref[I.1.3.7-fr]{I, 1.3.7},
\hyperref[I.1.3.13-fr]{1.3.13} et
\hyperref[I.1.6.3-fr]{1.6.3}). Soit $V$ un second ouvert affine de $S$, et
soit $X_{U,V}$ le préschéma induit par $X_U$ sur $f_U^{-1}(U\cap V)$, en
désignant par $f_U$ le morphisme structural $X_U\to S$~; $X_{U,V}$ et
$X_{V,U}$ sont affines sur $U\cap V$
(\hyperref[II.1.2.5-fr]{1.2.5}), et par définition $\mathcal{A}(X_{U,V})$
et $\mathcal{A}(X_{V,U})$ s'identifient canoniquement à
$\mathcal{B}|(U\cap V)$. Il y a donc
(\hyperref[II.1.2.8-fr]{1.2.8}) un $S$-isomorphisme canonique
$\theta_{U,V}:X_{V,U}\to X_{U,V}$~; en outre, si $W$ est un troisième
ouvert affine de $S$, et si $\theta'_{U,V}$, $\theta'_{V,W}$,
$\theta'_{U,W}$ sont les restrictions de $\theta_{U,V}$, $\theta_{V,W}$,
$\theta_{U,W}$ aux images réciproques de $U\cap V\cap W$ dans $X_V$,
$X_W$ et $X_W$ respectivement par les morphismes structuraux, on a
$\theta'_{U,V}\circ\theta'_{V,W}=\theta'_{U,W}$. Il existe par suite un
préschéma $X$, un recouvrement $(T_U)$ de $X$ par des ouverts affines, et
pour chaque $U$ un isomorphisme $\varphi_U:X_U\to T_U$, de sorte que
$\varphi_U$ applique $f_U^{-1}(U\cap V)$ sur $T_U\cap T_V$ et que l'on ait
$\theta_{U,V}=\varphi_U^{-1}\circ\varphi_V$
(\hyperref[I.2.3.1-fr]{I, 2.3.1}). Le morphisme
$g_U=f_U\circ\varphi_U^{-1}$ fait de $T_U$ un $S$-préschéma, et les
morphismes $g_U$ et $g_V$ coïncident dans $T_U\cap T_V$, donc $X$ est un
$S$-préschéma. En outre, il est clair par définition que $X$ est affine
au-dessus de $S$ et que $\mathcal{A}(T_U)=\mathcal{B}|U$, donc
$\mathcal{A}(X)=\mathcal{B}$.

On dit que le $S$-schéma $X$ ainsi défini est \emph{associé à la
$\mathcal{O}_S$-Algèbre $\mathcal{B}$} ou est le \emph{spectre de
$\mathcal{B}$}, et on le note $\operatorname{Spec}(\mathcal{B})$.

\begin{corollary}[1.3.2]
\label{II.1.3.2-fr}
Soient $X$ un préschéma affine au-dessus de $S$, $f:X\to S$ le morphisme
structural. Pour tout ouvert affine $U\subset S$, le préschéma induit sur
$f^{-1}(U)$ est un schéma affine d'anneau
$\Gamma(U,\mathcal{A}(X))$.
\end{corollary}

\oldpage[II]{9}
Comme on peut supposer $X$ associé à une $\mathcal{O}_S$-Algèbre en vertu de
(\hyperref[II.1.2.6-fr]{1.2.6}) et (\hyperref[II.1.3.1-fr]{1.3.1}), le
corollaire résulte de la construction de $X$ décrite dans
(\hyperref[II.1.3.1-fr]{1.3.1}).

\begin{example}[1.3.3]
\label{II.1.3.3-fr}
Soit $S$ le plan affine sur un corps $K$, où le point $0$ a été dédoublé
(\hyperref[I.5.5.11-fr]{I, 5.5.11})~; avec les notations de
(\hyperref[I.5.5.11-fr]{I, 5.5.11}), $S$ est réunion de deux ouverts affines
$Y_1$, $Y_2$~; si $f$ est l'immersion ouverte $Y_1\to S$,
$f^{-1}(Y_2)=Y_1\cap Y_2$ n'est pas un ouvert affine dans $Y_1$
(\emph{loc. cit.}), donc on a un exemple d'un schéma affine qui n'est pas
affine au-dessus de $S$.
\end{example}

\begin{corollary}[1.3.4]
\label{II.1.3.4-fr}
Soit $S$ un schéma affine~; pour qu'un $S$-préschéma $X$ soit affine au-dessus
de $S$, il faut et il suffit que $X$ soit un schéma affine.
\end{corollary}

\begin{corollary}[1.3.5]
\label{II.1.3.5-fr}
Soient $X$ un préschéma affine au-dessus d'un préschéma $S$, $Y$ un
$X$-préschéma. Pour que $Y$ soit affine au-dessus de $S$, il faut et il
suffit que $Y$ soit affine au-dessus de $X$.
\end{corollary}

On est aussitôt ramené au cas où $S$ est un schéma affine, et il en est alors
de même de $X$ (\hyperref[II.1.3.4-fr]{1.3.4})~; les deux conditions de
l'énoncé expriment alors que $Y$ est un schéma affine
(\hyperref[II.1.3.4-fr]{1.3.4}).

\begin{env}[1.3.6]
\label{II.1.3.6-fr}
Soit $X$ un préschéma affine au-dessus de $S$. Pour définir un préschéma $Y$
affine \emph{au-dessus de $X$}, il revient au même, en vertu de
(\hyperref[II.1.3.5-fr]{1.3.5}), de se donner un préschéma $Y$ affine
\emph{au-dessus de $S$}, et un $S$-morphisme $g:Y\to X$~; en d'autres termes
(\hyperref[II.1.3.1-fr]{1.3.1} et
\hyperref[II.1.2.7-fr]{1.2.7}), cela revient à se donner une
$\mathcal{O}_S$-Algèbre quasi-cohérente $\mathcal{B}$ et un homomorphisme
$\mathcal{A}(X)\to\mathcal{B}$ de $\mathcal{O}_S$-Algèbres (que l'on peut
envisager comme définissant sur $\mathcal{B}$ une structure de
$\mathcal{A}(X)$-Algèbre). Si $f:X\to S$ est le morphisme structural, on a
alors $\mathcal{B}=f_*(g_*(\mathcal{O}_Y))$.
\end{env}

\begin{corollary}[1.3.7]
\label{II.1.3.7-fr}
Soit $X$ un préschéma affine au-dessus de $S$~; pour que $X$ soit de type fini
sur $S$, il faut et il suffit que la $\mathcal{O}_S$-Algèbre quasi-cohérente
$\mathcal{A}(X)$ soit de type fini
(\hyperref[I.9.6.2-fr]{I, 9.6.2}).
\end{corollary}

Par définition (\hyperref[I.9.6.2-fr]{I, 9.6.2}), on est ramené au cas où
$S$ est affine~; alors $X$ est un schéma affine
(\hyperref[II.1.3.4-fr]{1.3.4}), et si $S=\operatorname{Spec}(A)$,
$X=\operatorname{Spec}(B)$, $\mathcal{A}(X)$ est la
$\mathcal{O}_S$-Algèbre $\widetilde{B}$~; comme
$\Gamma(U,\widetilde{B})=B$, le corollaire résulte de
(\hyperref[I.9.6.2-fr]{I, 9.6.2}) et
(\hyperref[I.6.3.3-fr]{I, 6.3.3}).

\begin{corollary}[1.3.8]
\label{II.1.3.8-fr}
Soit $X$ un préschéma affine au-dessus de $S$~; pour que $X$ soit réduit il
faut et il suffit que la $\mathcal{O}_S$-Algèbre quasi-cohérente
$\mathcal{A}(X)$ soit réduite
(\hyperref[0.4.1.4-fr]{0, 4.1.4}).
\end{corollary}

En effet, la question est évidemment locale sur $S$~; en vertu de
(\hyperref[II.1.3.2-fr]{1.3.2}), le corollaire résulte de
(\hyperref[I.5.1.1-fr]{I, 5.1.1}) et
(\hyperref[I.5.1.4-fr]{I, 5.1.4}).

\subsection{Faisceaux quasi-cohérents sur un préschéma affine au-dessus de
$S$}
\label{subsection:II.1.4-fr}

\begin{proposition}[1.4.1]
\label{II.1.4.1-fr}
Soient $X$ un préschéma affine au-dessus de $S$, $Y$ un $S$-préschéma,
$\mathcal{F}$ (resp. $\mathcal{G}$) un $\mathcal{O}_X$-Module (resp. un
$\mathcal{O}_Y$-Module) quasi-cohérent. Alors l'application
$(h,u)\mapsto(\mathcal{A}(h),\mathcal{A}(u))$ de l'ensemble des morphismes
$(Y,\mathcal{G})\to(X,\mathcal{F})$ dans l'ensemble des di-homomorphismes
$(\mathcal{A}(X),\mathcal{A}(\mathcal{F}))\to
(\mathcal{A}(Y),\mathcal{A}(\mathcal{G}))$
(\hyperref[II.1.1.2-fr]{1.1.2} et
\hyperref[II.1.1.3-fr]{1.1.3}) est bijective.
\end{proposition}

La démonstration suit exactement la même marche que celle de
(\hyperref[II.1.2.7-fr]{1.2.7}) en utilisant
(\hyperref[I.2.2.5-fr]{I, 2.2.5}) et
(\hyperref[I.2.2.4-fr]{I, 2.2.4}), et les détails en sont laissés au
lecteur.

\begin{corollary}[1.4.2]
\label{II.1.4.2-fr}
Si, en outre des hypothèses de (\hyperref[II.1.4.1-fr]{1.4.1}), on suppose
$Y$ affine au-dessus de $S$, alors, pour que $(h,u)$ soit un isomorphisme,
il faut et il suffit que $(\mathcal{A}(h),\mathcal{A}(u))$ soit un
di-isomorphisme.
\end{corollary}

\oldpage[II]{10}
\begin{proposition}[1.4.3]
\label{II.1.4.3-fr}
Pour tout couple $(\mathcal{B},\mathcal{M})$ formé d'une
$\mathcal{O}_S$-Algèbre quasi-cohérente $\mathcal{B}$ et d'un
$\mathcal{B}$-Module $\mathcal{M}$ quasi-cohérent (en tant que
$\mathcal{O}_S$-Module ou en tant que $\mathcal{B}$-Module, ce qui revient au
même (\hyperref[I.9.6.1-fr]{I, 9.6.1})), il existe un couple
$(X,\mathcal{F})$ formé d'un préschéma $X$ affine au-dessus de $S$ et d'un
$\mathcal{O}_X$-Module quasi-cohérent $\mathcal{F}$, tels que
$\mathcal{A}(X)=\mathcal{B}$ et $\mathcal{A}(\mathcal{F})=\mathcal{M}$~; en
outre ce couple est déterminé à un isomorphisme unique près.
\end{proposition}

L'unicité résulte de (\hyperref[II.1.4.1-fr]{1.4.1}) et
(\hyperref[II.1.4.2-fr]{1.4.2})~; l'existence se démontre comme dans
(\hyperref[II.1.3.1-fr]{1.3.1}), et nous laissons encore les détails au
lecteur. On note $\widetilde{\mathcal{M}}$ le $\mathcal{O}_X$-Module
$\mathcal{F}$, et on dit qu'il est \emph{associé} au $\mathcal{B}$-Module
quasi-cohérent $\mathcal{M}$~; pour tout ouvert affine $U\subset S$,
$\widetilde{\mathcal{M}}|p^{-1}(U)$ (où $p$ est le morphisme structural
$X\to S$) est canoniquement isomorphe à
$(\Gamma(U,\mathcal{M}))^{\sim}$.

\begin{corollary}[1.4.4]
\label{II.1.4.4-fr}
Dans la catégorie des $\mathcal{B}$-Modules quasi-cohérents,
$\widetilde{\mathcal{M}}$ est un foncteur covariant additif exact en
$\mathcal{M}$, qui commute aux limites inductives et aux sommes directes.
\end{corollary}

On est en effet aussitôt ramené au cas où $S$ est affine, et le corollaire
résulte alors de (\hyperref[I.1.3.5-fr]{I, 1.3.5}),
(\hyperref[I.1.3.9-fr]{I, 1.3.9}) et
(\hyperref[I.1.3.11-fr]{I, 1.3.11}).

\begin{corollary}[1.4.5]
\label{II.1.4.5-fr}
Sous les hypothèses de (\hyperref[II.1.4.3-fr]{1.4.3}), pour que
$\widetilde{\mathcal{M}}$ soit un $\mathcal{O}_X$-Module de type fini, il faut
et il suffit que $\mathcal{M}$ soit un $\mathcal{B}$-Module de type fini.
\end{corollary}

On est aussitôt ramené au cas où $S=\operatorname{Spec}(A)$ est un schéma
affine. Alors $\mathcal{B}=\widetilde{B}$, où $B$ est une $A$-algèbre de type
fini (\hyperref[I.9.6.2-fr]{I, 9.6.2}), et
$\mathcal{M}=\widetilde{M}$, où $M$ est un $B$-module
(\hyperref[I.1.3.13-fr]{I, 1.3.13})~; \emph{sur le préschéma $X$},
$\mathcal{O}_X$ est associé à l'anneau $B$ et
$\widetilde{\mathcal{M}}$ au $B$-module $M$~; pour que
$\widetilde{\mathcal{M}}$ soit de type fini, il faut et il suffit donc que
$M$ soit de type fini (\hyperref[I.1.3.13-fr]{I, 1.3.13}), d'où notre
assertion.

\begin{proposition}[1.4.6]
\label{II.1.4.6-fr}
Soient $Y$ un préschéma affine au-dessus de $S$, $X$, $X'$ deux préschémas
affines au-dessus de $Y$ (donc aussi au-dessus de $S$
(\hyperref[II.1.3.5-fr]{1.3.5})). Soient
$\mathcal{B}=\mathcal{A}(Y)$, $\mathcal{A}=\mathcal{A}(X)$,
$\mathcal{A}'=\mathcal{A}(X')$. Alors $X\times_Y X'$ est affine au-dessus de
$Y$ (donc aussi au-dessus de $S$) et $\mathcal{A}(X\times_Y X')$
s'identifie canoniquement à $\mathcal{A}\otimes_{\mathcal{B}}\mathcal{A}'$.
\end{proposition}

En effet (\hyperref[I.9.6.1-fr]{I, 9.6.1})
$\mathcal{A}\otimes_{\mathcal{B}}\mathcal{A}'$ est une
$\mathcal{B}$-Algèbre quasi-cohérente, donc aussi une
$\mathcal{O}_S$-Algèbre quasi-cohérente
(\hyperref[I.9.6.1-fr]{I, 9.6.1})~; soit $Z$ le spectre de
$\mathcal{A}\otimes_{\mathcal{B}}\mathcal{A}'$
(\hyperref[II.1.3.1-fr]{1.3.1}). Les $\mathcal{B}$-homomorphismes canoniques
$\mathcal{A}\to\mathcal{A}\otimes_{\mathcal{B}}\mathcal{A}'$ et
$\mathcal{A}'\to\mathcal{A}\otimes_{\mathcal{B}}\mathcal{A}'$ correspondent
(\hyperref[II.1.2.7-fr]{1.2.7}) à des $Y$-morphismes $p:Z\to X$ et
$p':Z\to X'$. Pour voir que le triplet $(Z,p,p')$ est un produit
$X\times_Y X'$, on peut se borner au cas où $S$ est un schéma affine
d'anneau $C$ (\hyperref[I.3.2.6.4-fr]{I, 3.2.6.4}). Mais alors $Y$, $X$,
$X'$ sont des schémas affines (\hyperref[II.1.3.4-fr]{1.3.4}) dont les
anneaux $B$, $A$, $A'$ sont des $C$-algèbres telles que
$\mathcal{B}=\widetilde{B}$, $\mathcal{A}=\widetilde{A}$,
$\mathcal{A}'=\widetilde{A'}$. On sait alors
(\hyperref[I.1.3.13-fr]{I, 1.3.13}) que
$\mathcal{A}\otimes_{\mathcal{B}}\mathcal{A}'$ s'identifie canoniquement à
la $\mathcal{O}_S$-Algèbre $\widetilde{A\otimes_B A'}$, donc l'anneau
$A(Z)$ s'identifie à $A\otimes_B A'$ et les morphismes $p$, $p'$ correspondent
aux homomorphismes canoniques $A\to A\otimes_B A'$,
$A'\to A\otimes_B A'$. La proposition résulte alors de
(\hyperref[I.3.2.2-fr]{I, 3.2.2}).

\begin{corollary}[1.4.7]
\label{II.1.4.7-fr}
Soit $\mathcal{F}$ (resp. $\mathcal{F}'$) un $\mathcal{O}_X$-Module (resp. un
$\mathcal{O}_{X'}$-Module) quasi-cohérent~; alors
$\mathcal{A}(\mathcal{F}\otimes_Y\mathcal{F}')$ s'identifie canoniquement à
$\mathcal{A}(\mathcal{F})\otimes_{\mathcal{A}(Y)}
\mathcal{A}(\mathcal{F}')$.
\end{corollary}

On sait que $\mathcal{F}\otimes_Y\mathcal{F}'$ est quasi-cohérent sur
$X\times_Y X'$ (\hyperref[I.9.1.2-fr]{I, 9.1.2}). Soient $g:Y\to S$,
$f:X\to Y$, $f':X'\to Y$ les morphismes structuraux, de sorte que le
morphisme structural
\oldpage[II]{11}
$h:Z\to S$ est égal à $g\circ f\circ p$ et à $g\circ f'\circ p'$. On
définit un homomorphisme canonique
\[
  \mathcal{A}(\mathcal{F})\otimes_{\mathcal{A}(Y)}\mathcal{A}(\mathcal{F}')
  \longrightarrow \mathcal{A}(\mathcal{F}\otimes_Y\mathcal{F}')
\]
de la façon suivante : pour tout ouvert $U\subset S$, on a des homomorphismes
canoniques
\[
  \begin{aligned}
  \Gamma(f^{-1}(g^{-1}(U)),\mathcal{F})
  &\longrightarrow \Gamma(h^{-1}(U),p^*(\mathcal{F}))\quad\text{et}\quad
  \Gamma(f^{\prime-1}(g^{-1}(U)),\mathcal{F}')
  \longrightarrow \Gamma(h^{-1}(U),p^{\prime*}(\mathcal{F}')).
  \end{aligned}
\]
(\hyperref[0.4.4.3-fr]{0, 4.4.3}), d'où on déduit un homomorphisme canonique
\[
  \begin{aligned}
  &\Gamma(f^{-1}(g^{-1}(U)),\mathcal{F})
    \otimes_{\Gamma(g^{-1}(U),\mathcal{O}_Y)}
    \Gamma(f^{\prime-1}(g^{-1}(U)),\mathcal{F}')\longrightarrow\\
  &\hspace{17em}
    \Gamma(h^{-1}(U),p^*(\mathcal{F}))
    \otimes_{\Gamma(h^{-1}(U),\mathcal{O}_Z)}
    \Gamma(h^{-1}(U),p^{\prime*}(\mathcal{F}')).
  \end{aligned}
\]

Pour voir qu'on a bien là un isomorphisme de $\mathcal{A}(Z)$-Modules, on
peut se borner au cas où $S$ (et par suite $X$, $X'$, $Y$,
$X\times_Y X'$) sont des schémas affines, et (avec les notations de
(\hyperref[II.1.4.6-fr]{1.4.6})), $\mathcal{F}=\widetilde{M}$,
$\mathcal{F}'=\widetilde{M'}$, $M$ (resp. $M'$) étant un $A$-module (resp. un
$A'$-module). Alors $\mathcal{F}\otimes_Y\mathcal{F}'$ s'identifie au
faisceau sur $X\times_Y X'$ associé au $(A\otimes_B A')$-module
$M\otimes_B M'$ (\hyperref[I.9.1.3-fr]{I, 9.1.3}) et le corollaire résulte de
l'identification canonique des $\mathcal{O}_S$-Modules
$\widetilde{M\otimes_B M'}$ et
$\widetilde{M}\otimes_{\widetilde{B}}\widetilde{M'}$ (où $M$, $M'$ et $B$
sont considérés comme des $C$-modules)
(\hyperref[I.1.3.13-fr]{I, 1.3.13} et
\hyperref[I.1.6.3-fr]{I, 1.6.3}).

Si on applique en particulier (\hyperref[II.1.4.7-fr]{1.4.7}) au cas où
$X=Y$ et $\mathcal{F}'=\mathcal{O}_{X'}$, on voit que le
$\mathcal{A}'$-Module $\mathcal{A}(f^{\prime*}(\mathcal{F}))$ s'identifie à
$\mathcal{A}(\mathcal{F})\otimes_{\mathcal{B}}\mathcal{A}'$.

\begin{env}[1.4.8]
\label{II.1.4.8-fr}
En particulier, lorsque $X=X'=Y$ ($X$ étant affine au-dessus de $S$), on voit
que si $\mathcal{F}$, $\mathcal{G}$ sont deux $\mathcal{O}_X$-Modules
quasi-cohérents, on a
\[
  \mathcal{A}(\mathcal{F}\otimes_{\mathcal{O}_X}\mathcal{G})
  =\mathcal{A}(\mathcal{F})\otimes_{\mathcal{A}(X)}\mathcal{A}(\mathcal{G})
  \tag{1.4.8.1}\label{II.1.4.8.1-fr}
\]
à un isomorphisme canonique fonctoriel près. Si de plus $\mathcal{F}$ admet
une présentation finie, il résulte de
(\hyperref[I.1.6.3-fr]{I, 1.6.3} et
\hyperref[I.1.3.12-fr]{I, 1.3.12}) que
\[
  \mathcal{A}(\mathcal{H}om_X(\mathcal{F},\mathcal{G}))
  =\mathcal{H}om_{\mathcal{A}(X)}
    (\mathcal{A}(\mathcal{F}),\mathcal{A}(\mathcal{G}))
  \tag{1.4.8.2}\label{II.1.4.8.2-fr}
\]
à un isomorphisme canonique près.
\end{env}

\begin{remark}[1.4.9]
\label{II.1.4.9-fr}
Si $X$, $X'$ sont deux préschémas affines au-dessus de $S$, la somme
$X\sqcup X'$ est aussi affine au-dessus de $S$, la somme de deux schémas
affines étant un schéma affine.
\end{remark}

\begin{proposition}[1.4.10]
\label{II.1.4.10-fr}
Soient $S$ un préschéma, $\mathcal{B}$ une $\mathcal{O}_S$-Algèbre
quasi-cohérente, $X=\operatorname{Spec}(\mathcal{B})$. Pour tout faisceau
quasi-cohérent d'idéaux $\mathcal{J}$ de $\mathcal{B}$,
$\widetilde{\mathcal{J}}$ est un faisceau quasi-cohérent d'idéaux de
$\mathcal{O}_X$, et le sous-préschéma fermé $Y$ de $X$ défini par
$\widetilde{\mathcal{J}}$ est canoniquement isomorphe à
$\operatorname{Spec}(\mathcal{B}/\mathcal{J})$.
\end{proposition}

En effet, il résulte aussitôt de (\hyperref[I.4.2.3-fr]{I, 4.2.3}) que $Y$
est affine au-dessus de $S$~; en vertu de
(\hyperref[II.1.3.1-fr]{1.3.1}), on est donc ramené au cas où $S$ est affine,
et la proposition résulte alors aussitôt de
(\hyperref[I.4.1.2-fr]{I, 4.1.2}).

On peut encore exprimer le résultat de
(\hyperref[II.1.4.10-fr]{1.4.10}) en disant que si
$h:\mathcal{B}\to\mathcal{B}'$ est un homomorphisme \emph{surjectif} de
$\mathcal{O}_S$-Algèbres quasi-cohérentes,
${}^a h:\operatorname{Spec}(\mathcal{B}')\to
\operatorname{Spec}(\mathcal{B})$ est une \emph{immersion fermée}.

\oldpage[II]{12}
\begin{proposition}[1.4.11]
\label{II.1.4.11-fr}
Soient $S$ un préschéma, $\mathcal{B}$ une $\mathcal{O}_S$-Algèbre
quasi-cohérente, $X=\operatorname{Spec}(\mathcal{B})$. Pour tout Idéal
quasi-cohérent $\mathcal{K}$ de $\mathcal{O}_S$, on a (en désignant par $f$
le morphisme structural $X\to S$),
$f^*(\mathcal{K})\mathcal{O}_X=\widetilde{\mathcal{K}\mathcal{B}}$ à un
isomorphisme canonique près.
\end{proposition}

La question étant locale sur $S$, on est ramené au cas où
$S=\operatorname{Spec}(A)$ est affine, et dans ce cas la proposition n'est
autre que (\hyperref[I.1.6.9-fr]{I, 1.6.9}).

\subsection{Changement du préschéma de base}
\label{II.1.5-fr}

\begin{proposition}[1.5.1]
\label{II.1.5.1-fr}
Soit $X$ un préschéma affine au-dessus de $S$. Pour toute extension
$g:S'\to S$ du préschéma de base, $X'=X_{(S')}=X\times_S S'$ est affine
au-dessus de $S'$.
\end{proposition}

Si $f'$ est la projection $X'\to S'$, il suffit de prouver que
$f^{\prime-1}(U')$ est un ouvert affine pour tout ouvert affine $U'$ de $S'$
tel que $g(U')$ soit contenu dans un ouvert affine $U$ de $S$
(\hyperref[II.1.2.1-fr]{1.2.1})~; on peut donc se limiter au cas où $S$ et
$S'$ sont affines, et il suffit de prouver que $X'$ est alors un schéma
affine (\hyperref[II.1.3.4-fr]{1.3.4}). Mais alors
(\hyperref[II.1.3.4-fr]{1.3.4}) $X$ est un schéma affine, et si $A$, $A'$ et
$B$ sont les anneaux de $S$, $S'$ et $X$ respectivement, on sait que $X'$ est
un schéma affine d'anneau $A'\otimes_A B$
(\hyperref[I.3.2.2-fr]{I, 3.2.2}).

\begin{corollary}[1.5.2]
\label{II.1.5.2-fr}
Sous les hypothèses de (\hyperref[II.1.5.1-fr]{1.5.1}), soient $f:X\to S$
le morphisme structural, $f':X'\to S'$, $g':X'\to X$ les projections, de
sorte que le diagramme
\[
  \xymatrix{
    X\ar[d] & X'\ar[l]_{g'}\ar[d]^{f'}\\
    S & S'\ar[l]_g
  }
\]
est commutatif. Pour tout $\mathcal{O}_X$-Module quasi-cohérent $\mathcal{F}$,
il existe un isomorphisme canonique de $\mathcal{O}_{S'}$-Modules
\[
  u:g^*(f_*(\mathcal{F}))\xrightarrow{\sim}
  f'_*(g^{\prime*}(\mathcal{F})).
  \tag{1.5.2.1}\label{II.1.5.2.1-fr}
\]
En particulier, il existe un isomorphisme canonique de $\mathcal{A}(X')$ sur
$g^*(\mathcal{A}(X))$.
\end{corollary}

Pour définir $u$, il suffit de définir un homomorphisme
\[
  v:f_*(\mathcal{F})\longrightarrow
  g_*(f'_*(g^{\prime*}(\mathcal{F})))
  =f_*(g'_*(g^{\prime*}(\mathcal{F})))
\]
et de prendre $u=v^\sharp$ (\hyperref[0.4.4.3-fr]{0, 4.4.3}). Nous prendrons
$v=f_*(\rho)$, où $\rho$ est l'homomorphisme canonique
$\mathcal{F}\to g'_*(g^{\prime*}(\mathcal{F}))$
(\hyperref[0.4.4.3-fr]{0, 4.4.3}). Pour prouver que $u$ est un isomorphisme,
on peut se borner au cas où $S$ et $S'$, donc $X$ et $X'$, sont affines~;
avec les notations de (\hyperref[II.1.5.1-fr]{1.5.1}), on a alors
$\mathcal{F}=\widetilde{M}$, où $M$ est un $B$-module. On constate alors
aussitôt que $g^*(f_*(\mathcal{F}))$ et
$f'_*(g^{\prime*}(\mathcal{F}))$ sont tous deux égaux au
$\mathcal{O}_{S'}$-Module associé au $A'$-module $A'\otimes_A M$ (où $M$ est
considéré comme $A$-module), et que $u$ est l'homomorphisme associé à
l'identité (\hyperref[I.1.6.3-fr]{I, 1.6.3},
\hyperref[I.1.6.5-fr]{1.6.5} et \hyperref[I.1.6.7-fr]{1.6.7}).

\begin{remark}[1.5.3]
\label{II.1.5.3-fr}
On se gardera de croire que (\hyperref[II.1.5.2-fr]{1.5.2}) reste valable
lorsque $X$ n'est plus supposé affine au-dessus de $S$, même lorsque
$S'=\operatorname{Spec}(k(s))$ ($s\in S$) et que $S'\to S$ est le morphisme
canonique (\hyperref[I.2.4.5-fr]{I, 2.4.5}) --- auquel cas $X'$ n'est autre
que la \emph{fibre} $f^{-1}(s)$ (\hyperref[I.3.6.2-fr]{I, 3.6.2}). En
d'autres termes, lorsque $X$ n'est pas affine au-dessus de $S$, l'opération
\oldpage[II]{13}
<<~image directe de faisceaux quasi-cohérents~>> ne permute pas à l'opération
de <<~passage aux fibres~>>. Nous verrons cependant au chapitre~III
(\hyperref[III.4.2.4-fr]{III, 4.2.4}) un résultat dans ce sens, de nature
<<~asymptotique~>>, valable pour les faisceaux \emph{cohérents} sur $X$ lorsque
$f$ est propre (\hyperref[II.5.4-fr]{5.4}) et $S$ noethérien.
\end{remark}

\begin{corollary}[1.5.4]
\label{II.1.5.4-fr}
Pour tout préschéma $X$ affine au-dessus de $S$ et tout $s\in S$, la fibre
$f^{-1}(s)$ (où $f$ désigne le morphisme structural $X\to S$) est un schéma
affine.
\end{corollary}

Il suffit d'appliquer (\hyperref[II.1.5.1-fr]{1.5.1}) à
$S'=\operatorname{Spec}(k(s))$ et d'utiliser
(\hyperref[II.1.3.4-fr]{1.3.4}).

\begin{corollary}[1.5.5]
\label{II.1.5.5-fr}
Soient $X$ un $S$-préschéma, $S'$ un préschéma affine au-dessus de $S$~; alors
$X'=X_{(S')}$ est un préschéma affine au-dessus de $X$. En outre, si
$f:X\to S$ est le morphisme structural, il y a un isomorphisme canonique de
$\mathcal{O}_X$-Algèbres
$\mathcal{A}(X')\xrightarrow{\sim}f^*(\mathcal{A}(S'))$, et pour tout
$\mathcal{A}(S')$-Module quasi-cohérent $\mathcal{M}$, un di-isomorphisme
canonique
$f^*(\mathcal{M})\xrightarrow{\sim}
\mathcal{A}(f^{\prime*}(\widetilde{\mathcal{M}}))$, en désignant par
$f'=f_{(S')}$ le morphisme structural $X'\to S'$.
\end{corollary}

Il suffit d'intervertir les rôles de $X$ et de $S'$ dans
(\hyperref[II.1.5.1-fr]{1.5.1}) et
(\hyperref[II.1.5.2-fr]{1.5.2}).

\begin{env}[1.5.6]
\label{II.1.5.6-fr}
Soient maintenant $S$, $S'$ deux préschémas, $q:S'\to S$ un morphisme,
$\mathcal{B}$ (resp. $\mathcal{B}'$) une $\mathcal{O}_S$-Algèbre (resp. une
$\mathcal{O}_{S'}$-Algèbre) quasi-cohérente, $u:\mathcal{B}\to\mathcal{B}'$
un $q$-morphisme (c'est-à-dire un homomorphisme
$\mathcal{B}\to q_*(\mathcal{B}')$ de $\mathcal{O}_S$-Algèbres). Si
$X=\operatorname{Spec}(\mathcal{B})$,
$X'=\operatorname{Spec}(\mathcal{B}')$, on en déduit canoniquement un
morphisme
\[
  v=\operatorname{Spec}(u):X'\to X
\]
tel que le diagramme
\[
  \xymatrix{
    X'\ar[r]^v\ar[d] & X\ar[d]\\
    S'\ar[r]_q & S
  }
  \tag{1.5.6.1}\label{II.1.5.6.1-fr}
\]
soit commutatif (les flèches verticales étant les morphismes structuraux). En
effet, la donnée de $u$ équivaut à celle d'un homomorphisme de
$\mathcal{O}_{S'}$-Algèbres quasi-cohérentes
$u^\sharp:q^*(\mathcal{B})\to\mathcal{B}'$
(\hyperref[0.4.4.3-fr]{0, 4.4.3})~; elle définit donc canoniquement un
$S'$-morphisme
\[
  w:\operatorname{Spec}(\mathcal{B}')\to
  \operatorname{Spec}(q^*(\mathcal{B}))
\]
tel que $\mathcal{A}(w)=u^\sharp$
(\hyperref[II.1.2.7-fr]{1.2.7}). D'autre part, il résulte de
(\hyperref[II.1.5.2-fr]{1.5.2}) que
$\operatorname{Spec}(q^*(\mathcal{B}))$ s'identifie canoniquement à
$X\times_S S'$~; le morphisme $v$ est le composé
$X'\xrightarrow{w}X\times_S S'\xrightarrow{p_1}X$ de $w$ avec la première
projection, et la commutativité de
(\hyperref[II.1.5.6.1-fr]{1.5.6.1}) résulte des définitions. Soient $U$ (resp.
$U'$) un ouvert affine de $S$ (resp. $S'$) tels que $q(U')\subset U$,
$A=\Gamma(U,\mathcal{O}_S)$, $A'=\Gamma(U',\mathcal{O}_{S'})$ leurs anneaux,
$B=\Gamma(U,\mathcal{B})$, $B'=\Gamma(U',\mathcal{B}')$~; la restriction de
$u$ en un $(q|U')$-morphisme
$\mathcal{B}|U\to\mathcal{B}'|U'$ correspond à un di-homomorphisme d'algèbres
$B\to B'$~; si $V$, $V'$ sont les images réciproques de $U$, $U'$ dans $X$,
$X'$ respectivement, par les morphismes structuraux, le morphisme $V'\to V$,
restriction de $v$, correspond (\hyperref[I.1.7.3-fr]{I, 1.7.3}) au
di-homomorphisme précédent.
\end{env}

\begin{env}[1.5.7]
\label{II.1.5.7-fr}
Sous les mêmes hypothèses que dans (\hyperref[II.1.5.6-fr]{1.5.6}), soit
$\mathcal{M}$ un $\mathcal{B}$-Module quasi-cohérent~; il y a alors un
isomorphisme canonique de $\mathcal{O}_{X'}$-Modules
\[
  v^*(\widetilde{\mathcal{M}})\xrightarrow{\sim}
  \bigl(q^*(\mathcal{M})\otimes_{q^*(\mathcal{B})}\mathcal{B}'\bigr)^{\sim}
  \tag{1.5.7.1}\label{II.1.5.7.1-fr}
\]

\oldpage[II]{14}
En effet, l'isomorphisme canonique
(\hyperref[II.1.5.2.1-fr]{1.5.2.1}) fournit un isomorphisme canonique de
$p_1^*(\widetilde{\mathcal{M}})$ sur le faisceau sur
$\operatorname{Spec}(q^*(\mathcal{B}))$ associé au
$q^*(\mathcal{B})$-Module $q^*(\mathcal{M})$, et il suffit ensuite d'appliquer
(\hyperref[II.1.4.7-fr]{1.4.7}).
\end{env}

\subsection{Morphismes affines}
\label{II.1.6-fr}

\begin{env}[1.6.1]
\label{II.1.6.1-fr}
Nous dirons qu'un morphisme $f:X\to Y$ de préschémas est \emph{affine} s'il
définit $X$ comme un préschéma affine au-dessus de $Y$. Les propriétés des
préschémas affines au-dessus d'un autre se traduisent comme suit dans ce
langage~:
\end{env}

\begin{proposition}[1.6.2]
\label{II.1.6.2-fr}
\begin{enumerate}
  \item[\rm(i)] Une immersion fermée est affine.
  \item[\rm(ii)] Le composé de deux morphismes affines est affine.
  \item[\rm(iii)] Si $f:X\to Y$ est un $S$-morphisme affine,
    $f_{(S')}:X_{(S')}\to Y_{(S')}$ est affine pour toute extension $S'\to S$
    de la base.
  \item[\rm(iv)] Si $f:X\to Y$ et $f':X'\to Y'$ sont deux $S$-morphismes
    affines,
    \[
      f\times_S f':X\times_S X'\to Y\times_S Y'
    \]
    est affine.
  \item[\rm(v)] Si $f:X\to Y$ et $g:Y\to Z$ sont deux morphismes tels que
    $g\circ f$ soit affine et $g$ séparé, alors $f$ est affine.
  \item[\rm(vi)] Si $f$ est affine, il en est de même de
    $f_{\mathrm{red}}$.
\end{enumerate}
\end{proposition}

En vertu de (\hyperref[I.5.5.12-fr]{I, 5.5.12}), il suffit de démontrer (i),
(ii) et (iii). Or (i) n'est autre que l'exemple
(\hyperref[II.1.2.2-fr]{1.2.2}), et (ii) n'est autre que
(\hyperref[II.1.3.5-fr]{1.3.5})~; enfin, (iii) découle de
(\hyperref[II.1.5.1-fr]{1.5.1}) puisque $X_{(S')}$ s'identifie au produit
$X\times_Y Y_{(S')}$ (\hyperref[I.3.3.11-fr]{I, 3.3.11}).

\begin{corollary}[1.6.3]
\label{II.1.6.3-fr}
Si $X$ est un schéma affine et $Y$ un schéma, tout morphisme $X\to Y$ est
affine.
\end{corollary}

\begin{proposition}[1.6.4]
\label{II.1.6.4-fr}
Soient $Y$ un préschéma localement noethérien, $f:X\to Y$ un morphisme de type
fini. Pour que $f$ soit affine, il faut et il suffit que $f_{\mathrm{red}}$ le
soit.
\end{proposition}

Vu (\hyperref[II.1.6.2-fr]{1.6.2, (vi)}), nous n'avons à démontrer que la
suffisance de la condition. Il suffit de prouver que si $Y$ est affine et
noethérien, $X$ est affine~; or $Y_{\mathrm{red}}$ est alors affine, donc il
en est de même de $X_{\mathrm{red}}$ par hypothèse. Or, $X$ est noethérien,
donc la conclusion résulte de (\hyperref[I.6.1.7-fr]{I, 6.1.7}).

\subsection{Fibré vectoriel associé à un Faisceau de modules}
\label{II.1.7-fr}

\begin{env}[1.7.1]
\label{II.1.7.1-fr}
Soient $A$ un anneau, $E$ un $A$-module. Rappelons qu'on appelle \emph{algèbre
symétrique} sur $E$ et qu'on note $\mathbf{S}(E)$ (ou $\mathbf{S}_A(E)$)
l'algèbre quotient de l'algèbre tensorielle $\mathbf{T}(E)$ par l'idéal
bilatère engendré par les éléments $x\otimes y-y\otimes x$, où $x,y$
parcourent $E$. L'algèbre $\mathbf{S}(E)$ est caractérisée par la propriété
universelle suivante~: si $\sigma$ est l'application canonique
$E\to\mathbf{S}(E)$ (obtenue par composition de $E\to\mathbf{T}(E)$ et de
l'application canonique $\mathbf{T}(E)\to\mathbf{S}(E)$), toute application
$A$-linéaire $E\to B$, où $B$ est une $A$-algèbre \emph{commutative}, se
factorise de façon unique en
$E\xrightarrow{\sigma}\mathbf{S}(E)\xrightarrow{g}B$, où $g$ est un
$A$-homomorphisme \emph{d'algèbres}. On déduit aussitôt de cette
caractérisation que pour deux $A$-modules $E$, $F$, on a
\[
  \mathbf{S}(E\oplus F)=\mathbf{S}(E)\otimes\mathbf{S}(F)
\]

\oldpage[II]{15}
à un isomorphisme canonique près~; en outre, $\mathbf{S}(E)$ est un foncteur
covariant en $E$, de la catégorie des $A$-modules dans celle des $A$-algèbres
commutatives~; enfin, la caractérisation précédente montre aussi que si
$E=\varinjlim E_\lambda$, on a
$\mathbf{S}(E)=\varinjlim\mathbf{S}(E_\lambda)$ à un isomorphisme canonique
près. Par abus de langage, un produit
$\sigma(x_1)\sigma(x_2)\cdots\sigma(x_n)$, où les $x_i\in E$, se note souvent
$x_1x_2\cdots x_n$ si aucune confusion n'en résulte. L'algèbre
$\mathbf{S}(E)$ est \emph{graduée}, $\mathbf{S}_n(E)$ étant l'ensemble des
combinaisons linéaires des produits de $n$ éléments de $E$ ($n\geq 0$)~;
l'algèbre $\mathbf{S}(A)$ est canoniquement isomorphe à l'algèbre des
polynômes $A[T]$ à une indéterminée et l'algèbre $\mathbf{S}(A^n)$ à
l'algèbre des polynômes à $n$ indéterminées $A[T_1,\ldots,T_n]$.
\end{env}

\begin{env}[1.7.2]
\label{II.1.7.2-fr}
Soit $\varphi$ un homomorphisme d'anneaux $A\to B$. Si $F$ est un $B$-module,
l'application canonique $F\to\mathbf{S}(F)$ donne une application canonique
$F_{[\varphi]}\to\mathbf{S}(F)_{[\varphi]}$, qui se factorise donc en
$F_{[\varphi]}\to\mathbf{S}(F_{[\varphi]})\to
\mathbf{S}(F)_{[\varphi]}$~; l'homomorphisme canonique
$\mathbf{S}(F_{[\varphi]})\to\mathbf{S}(F)_{[\varphi]}$ est surjectif, mais
non nécessairement bijectif. Si $E$ est un $A$-module, tout di-homomorphisme
$E\to F$ (c'est-à-dire tout $A$-homomorphisme $E\to F_{[\varphi]}$) donne donc
canoniquement un $A$-homomorphisme d'algèbres
$\mathbf{S}(E)\to\mathbf{S}(F_{[\varphi]})\to
\mathbf{S}(F)_{[\varphi]}$, c'est-à-dire un di-homomorphisme d'algèbres
$\mathbf{S}(E)\to\mathbf{S}(F)$.

Avec les mêmes notations, pour tout $A$-module $E$,
$\mathbf{S}(E\otimes_A B)$ s'identifie canoniquement à l'algèbre
$\mathbf{S}(E)\otimes_A B$~; cela résulte aussitôt de la propriété universelle
de $\mathbf{S}(E)$ (\hyperref[II.1.7.1-fr]{1.7.1}).
\end{env}

\begin{env}[1.7.3]
\label{II.1.7.3-fr}
Soit $R$ une partie multiplicative de l'anneau $A$~; appliquant
(\hyperref[II.1.7.2-fr]{1.7.2}) à l'anneau $B=R^{-1}A$, et se rappelant que
$R^{-1}E=E\otimes_A R^{-1}A$, on voit que l'on a
$\mathbf{S}(R^{-1}E)=R^{-1}\mathbf{S}(E)$ à un isomorphisme canonique près.
En outre, si $R'\supset R$ est une seconde partie multiplicative de $A$, le
diagramme
\[
  \xymatrix{
    R^{-1}E\ar[r]\ar[d] & {R'}^{-1}E\ar[d]\\
    \mathbf{S}(R^{-1}E)\ar[r] & \mathbf{S}({R'}^{-1}E)
  }
\]
est commutatif.
\end{env}

\begin{env}[1.7.4]
\label{II.1.7.4-fr}
Soit maintenant $(S,\mathcal{A})$ un espace annelé, et soit $\mathcal{E}$ un
$\mathcal{A}$-Module sur $S$. Si, à tout ouvert $U\subset S$ on associe le
$\Gamma(U,\mathcal{A})$-module
$\mathbf{S}(\Gamma(U,\mathcal{E}))$, on définit (vu le caractère fonctoriel
de $\mathbf{S}(E)$ (\hyperref[II.1.7.2-fr]{1.7.2})) un préfaisceau d'algèbres~;
on dit que le faisceau associé, que l'on note $\mathbf{S}(\mathcal{E})$ ou
$\mathbf{S}_{\mathcal{A}}(\mathcal{E})$ est la $\mathcal{A}$-Algèbre
\emph{symétrique} du $\mathcal{A}$-Module $\mathcal{E}$. Il résulte aussitôt
de (\hyperref[II.1.7.1-fr]{1.7.1}) que $\mathbf{S}(\mathcal{E})$ est solution
d'un problème universel~: tout homomorphisme de $\mathcal{A}$-Modules
$\mathcal{E}\to\mathcal{B}$, où $\mathcal{B}$ est une $\mathcal{A}$-Algèbre,
se factorise d'une seule manière en
$\mathcal{E}\to\mathbf{S}(\mathcal{E})\to\mathcal{B}$, la seconde flèche étant
un homomorphisme de $\mathcal{A}$-Algèbres. Il y a donc correspondance
biunivoque entre homomorphismes $\mathcal{E}\to\mathcal{B}$ de
$\mathcal{A}$-Modules et homomorphismes
$\mathbf{S}(\mathcal{E})\to\mathcal{B}$ de $\mathcal{A}$-Algèbres. En
particulier, tout homomorphisme $u:\mathcal{E}\to\mathcal{F}$ de
$\mathcal{A}$-Modules définit un homomorphisme
$\mathbf{S}(u):\mathbf{S}(\mathcal{E})\to\mathbf{S}(\mathcal{F})$ de
$\mathcal{A}$-Algèbres et $\mathbf{S}(\mathcal{E})$ est donc un foncteur
covariant en $\mathcal{E}$.

\oldpage[II]{16}
En vertu de (\hyperref[II.1.7.2-fr]{1.7.2}) et de la permutabilité de
$\mathbf{S}$ avec les limites inductives, on a
$(\mathbf{S}(\mathcal{E}))_s=\mathbf{S}(\mathcal{E}_s)$ pour tout point
$s\in S$. Si $\mathcal{E}$, $\mathcal{F}$ sont deux $\mathcal{A}$-Modules,
$\mathbf{S}(\mathcal{E}\oplus\mathcal{F})$ s'identifie canoniquement à
$\mathbf{S}(\mathcal{E})\otimes_{\mathcal{A}}\mathbf{S}(\mathcal{F})$,
comme on le voit sur les préfaisceaux correspondants.

Notons aussi que $\mathbf{S}(\mathcal{E})$ est une $\mathcal{A}$-Algèbre
graduée, somme directe infinie des $\mathbf{S}_n(\mathcal{E})$, où le
$\mathcal{A}$-Module $\mathbf{S}_n(\mathcal{E})$ est le faisceau associé au
préfaisceau $U\mapsto\mathbf{S}_n(\Gamma(U,\mathcal{E}))$. Si on prend en
particulier $\mathcal{E}=\mathcal{A}$, on voit que
$\mathbf{S}_{\mathcal{A}}(\mathcal{A})$ s'identifie à
$\mathcal{A}[T]=\mathcal{A}\otimes_{\mathbf{Z}}\mathbf{Z}[T]$
($T$ indéterminée, $\mathbf{Z}$ étant considéré comme faisceau simple).
\end{env}

\begin{env}[1.7.5]
\label{II.1.7.5-fr}
Soient $(T,\mathcal{B})$ un second espace annelé, $f$ un morphisme
$(S,\mathcal{A})\to(T,\mathcal{B})$. Si $\mathcal{F}$ est un
$\mathcal{B}$-Module, $\mathbf{S}(f^*(\mathcal{F}))$ s'identifie
canoniquement à $f^*(\mathbf{S}(\mathcal{F}))$~; en effet, si
$f=(\psi,\theta)$, on a par définition
(\hyperref[0.4.3.1-fr]{0, 4.3.1})
\[
  \mathbf{S}(f^*(\mathcal{F}))
  =\mathbf{S}(\psi^*(\mathcal{F})
    \otimes_{\psi^*(\mathcal{B})}\mathcal{A})
  =\mathbf{S}(\psi^*(\mathcal{F}))
    \otimes_{\psi^*(\mathcal{B})}\mathcal{A}
\]
(\hyperref[II.1.7.2-fr]{1.7.2})~; pour tout ouvert $U$ de $S$ et toute
section $h$ de $\mathbf{S}(\psi^*(\mathcal{F}))$ au-dessus de $U$, $h$
coïncide, dans un voisinage $V$ de chaque point $s\in U$, avec un élément de
$\mathbf{S}(\Gamma(V,\psi^*(\mathcal{F})))$~; si on se reporte à la
définition de $\psi^*(\mathcal{F})$
(\hyperref[0.3.7.1-fr]{0, 3.7.1}) et qu'on tient compte de ce que tout élément
de $\mathbf{S}(E)$ pour un module $E$ est combinaison linéaire d'un nombre
fini de produits d'éléments de $E$, on voit qu'il y a un voisinage $W$ de
$\psi(s)$ dans $T$, une section $h'$ de $\mathbf{S}(\mathcal{F})$ au-dessus
de $W$ et un voisinage $V'\subset V\cap\psi^{-1}(W)$ de $s$ tels que $h$
coïncide avec $t\mapsto h'(\psi(t))$ au-dessus de $V'$~; d'où notre
assertion.
\end{env}

\begin{proposition}[1.7.6]
\label{II.1.7.6-fr}
Soient $A$ un anneau, $S=\operatorname{Spec}(A)$ son spectre premier,
$\mathcal{E}=\widetilde{M}$ le $\mathcal{O}_S$-Module associé à un $A$-module
$M$~; alors la $\mathcal{O}_S$-Algèbre $\mathbf{S}(\mathcal{E})$ est associée
à la $A$-algèbre $\mathbf{S}(M)$.
\end{proposition}

\begin{proof}
En effet, pour tout $f\in A$,
$\mathbf{S}(M_f)=(\mathbf{S}(M))_f$
(\hyperref[II.1.7.3-fr]{1.7.3}) et la proposition résulte donc des définitions
(\hyperref[I.1.3.4-fr]{I, 1.3.4}).
\end{proof}

\begin{corollary}[1.7.7]
\label{II.1.7.7-fr}
Si $S$ est un préschéma, $\mathcal{E}$ un $\mathcal{O}_S$-Module
quasi-cohérent, la $\mathcal{O}_S$-Algèbre $\mathbf{S}(\mathcal{E})$ est
quasi-cohérente. Si en outre $\mathcal{E}$ est de type fini, chacun des
$\mathcal{O}_S$-Modules $\mathbf{S}_n(\mathcal{E})$ est de type fini.
\end{corollary}

\begin{proof}
La première assertion est conséquence immédiate de
(\hyperref[II.1.7.6-fr]{1.7.6}) et de
(\hyperref[I.1.4.1-fr]{I, 1.4.1})~; la seconde résulte de ce que, si $E$ est
un $A$-module de type fini, $\mathbf{S}_n(E)$ est un $A$-module de type
fini~; on applique alors (\hyperref[I.1.3.13-fr]{I, 1.3.13}).
\end{proof}

\begin{definition}[1.7.8]
\label{II.1.7.8-fr}
Soit $\mathcal{E}$ un $\mathcal{O}_S$-Module quasi-cohérent. On appelle
\emph{fibré vectoriel sur $S$ défini par $\mathcal{E}$} et on note
$\mathbf{V}(\mathcal{E})$ le spectre
(\hyperref[II.1.3.1-fr]{1.3.1}) de la $\mathcal{O}_S$-Algèbre quasi-cohérente
$\mathbf{S}(\mathcal{E})$.
\end{definition}

En vertu de (\hyperref[II.1.2.7-fr]{1.2.7}), pour tout $S$-préschéma $X$, il
y a correspondance biunivoque canonique entre les $S$-morphismes
$X\to\mathbf{V}(\mathcal{E})$ et les homomorphismes de
$\mathcal{O}_S$-Algèbres $\mathbf{S}(\mathcal{E})\to\mathcal{A}(X)$, donc
aussi entre ces $S$-morphismes et les homomorphismes de
$\mathcal{O}_S$-Modules
$\mathcal{E}\to\mathcal{A}(X)=f_*(\mathcal{O}_X)$ (où $f$ est le morphisme
structural $X\to S$). En particulier~:

\begin{env}[1.7.9]
\label{II.1.7.9-fr}
Prenons pour $X$ un sous-préschéma induit par $S$ sur un ouvert $U\subset S$.
Alors, les $S$-morphismes $U\to\mathbf{V}(\mathcal{E})$ ne sont autres que
les $U$-sections (\hyperref[I.2.5.5-fr]{I, 2.5.5}) du $U$-préschéma induit
par $\mathbf{V}(\mathcal{E})$ sur l'ouvert $p^{-1}(U)$ (où $p$ est le
morphisme structural $\mathbf{V}(\mathcal{E})\to S$). D'après ce qu'on vient
de voir, ces $U$-sections correspondent biunivoquement aux homomorphismes de
$\mathcal{O}_S$-Modules
$\mathcal{E}\to j_*(\mathcal{O}_S|U)$ (où $j$ est l'injection canonique
$U\to S$), ou,

\oldpage[II]{17}
ce qui revient au même (\hyperref[0.4.4.3-fr]{0, 4.4.3}) aux
$(\mathcal{O}_S|U)$-homomorphismes
$j^*(\mathcal{E})=\mathcal{E}|U\to\mathcal{O}_S|U$. En outre, il est
immédiat qu'à la restriction à un ouvert $U'\subset U$ d'un $S$-morphisme
$U\to\mathbf{V}(\mathcal{E})$ correspond la restriction à $U'$ de
l'homomorphisme correspondant
$\mathcal{E}|U\to\mathcal{O}_S|U$. On en conclut que \emph{le faisceau des
germes de $S$-sections} de $\mathbf{V}(\mathcal{E})$ s'identifie
canoniquement au \emph{dual $\check{\mathcal{E}}$} de $\mathcal{E}$.

En particulier, si on prend $X=U=S$, à l'homomorphisme \emph{nul}
$\mathcal{E}\to\mathcal{O}_S$ correspond une $S$-section canonique de
$\mathbf{V}(\mathcal{E})$, dite \emph{$S$-section nulle}
(cf. (\hyperref[II.8.3.3-fr]{8.3.3})).
\end{env}

\begin{env}[1.7.10]
\label{II.1.7.10-fr}
Prenons maintenant pour $X$ le spectre $\{\xi\}$ d'un corps $K$~; le
morphisme structural $f:X\to S$ correspond donc à un monomorphisme
$k(s)\to K$, où $s=f(\xi)$
(\hyperref[I.2.4.6-fr]{I, 2.4.6})~; les $S$-morphismes
$\{\xi\}\to\mathbf{V}(\mathcal{E})$ ne sont donc autres que les
\emph{points géométriques de $\mathbf{V}(\mathcal{E})$ à valeurs dans
l'extension $K$ de $k(s)$} (\hyperref[I.3.4.5-fr]{I, 3.4.5}), points qui
sont localisés aux points de $p^{-1}(s)$. L'ensemble de ces points, qu'on
peut encore appeler \emph{la fibre géométrique rationnelle sur $K$ de
$\mathbf{V}(\mathcal{E})$ au-dessus du point $s$}, s'identifie d'après
(\hyperref[II.1.7.8-fr]{1.7.8}) à l'ensemble des homomorphismes de
$\mathcal{O}_S$-Modules
$\mathcal{E}\to f_*(\mathcal{O}_X)$, ou, ce qui revient au même
(\hyperref[0.4.4.3-fr]{0, 4.4.3}) à l'ensemble des homomorphismes de
$\mathcal{O}_X$-Modules
$f^*(\mathcal{E})\to\mathcal{O}_X=K$. Mais on a par définition
(\hyperref[0.4.3.1-fr]{0, 4.3.1})
$f^*(\mathcal{E})=\mathcal{E}_s\otimes_{\mathcal{O}_s}K
=\mathcal{E}^s\otimes_{k(s)}K$, en posant
$\mathcal{E}^s=\mathcal{E}_s/\mathfrak{m}_s\mathcal{E}_s$~; la fibre
géométrique de $\mathbf{V}(\mathcal{E})$ rationnelle sur $K$ au-dessus de
$s$ s'identifie donc au \emph{dual} du \emph{$K$-espace vectoriel
$\mathcal{E}^s\otimes_{k(s)}K$}~; si $\mathcal{E}^s$ ou $K$ est de
dimension finie sur $k(s)$, ce dual s'identifie aussi à
$(\mathcal{E}^s)^\vee\otimes_{k(s)}K$, en désignant par
$(\mathcal{E}^s)^\vee$ le dual du $k(s)$-espace vectoriel
$\mathcal{E}^s$.
\end{env}

\begin{proposition}[1.7.11]
\label{II.1.7.11-fr}
\begin{enumerate}
  \item[\rm(i)] $\mathbf{V}(\mathcal{E})$ est un foncteur contravariant en
    $\mathcal{E}$ de la catégorie des $\mathcal{O}_S$-Modules
    quasi-cohérents dans la catégorie des $S$-schémas affines.
  \item[\rm(ii)] Si $\mathcal{E}$ est un $\mathcal{O}_S$-Module de type
    fini, $\mathbf{V}(\mathcal{E})$ est de type fini sur $S$.
  \item[\rm(iii)] Si $\mathcal{E}$ et $\mathcal{F}$ sont deux
    $\mathcal{O}_S$-Modules quasi-cohérents,
    $\mathbf{V}(\mathcal{E}\oplus\mathcal{F})$ s'identifie canoniquement à
    $\mathbf{V}(\mathcal{E})\times_S\mathbf{V}(\mathcal{F})$.
  \item[\rm(iv)] Soit $g:S'\to S$ un morphisme~; pour tout
    $\mathcal{O}_S$-Module quasi-cohérent $\mathcal{E}$,
    $\mathbf{V}(g^*(\mathcal{E}))$ s'identifie canoniquement à
    $\mathbf{V}(\mathcal{E})_{(S')}
    =\mathbf{V}(\mathcal{E})\times_S S'$.
  \item[\rm(v)] À un homomorphisme surjectif
    $\mathcal{E}\to\mathcal{F}$ de $\mathcal{O}_S$-Modules quasi-cohérents
    correspond une immersion fermée
    $\mathbf{V}(\mathcal{F})\to\mathbf{V}(\mathcal{E})$.
\end{enumerate}
\end{proposition}

\begin{proof}
(i) est conséquence immédiate de
(\hyperref[II.1.2.7-fr]{1.2.7}), compte tenu de ce que tout homomorphisme de
$\mathcal{O}_S$-Modules $\mathcal{E}\to\mathcal{F}$ définit canoniquement un
homomorphisme de $\mathcal{O}_S$-Algèbres
$\mathbf{S}(\mathcal{E})\to\mathbf{S}(\mathcal{F})$. (ii) découle
immédiatement des définitions (\hyperref[I.6.3.1-fr]{I, 6.3.1}) et du fait
que si $E$ est un $A$-module de type fini, $\mathbf{S}(E)$ est une
$A$-algèbre de type fini. Pour démontrer (iii), il suffit de partir de
l'isomorphisme canonique
$\mathbf{S}(\mathcal{E}\oplus\mathcal{F})
\simeq\mathbf{S}(\mathcal{E})
\otimes_{\mathcal{O}_S}\mathbf{S}(\mathcal{F})$
(\hyperref[II.1.7.4-fr]{1.7.4}) et d'appliquer
(\hyperref[II.1.4.6-fr]{1.4.6}). De même, pour démontrer (iv), il suffit de
partir de l'isomorphisme canonique
$\mathbf{S}(g^*(\mathcal{E}))
\simeq g^*(\mathbf{S}(\mathcal{E}))$
(\hyperref[II.1.7.5-fr]{1.7.5}) et d'appliquer
(\hyperref[II.1.5.2-fr]{1.5.2}). Enfin, pour établir (v), il suffit de
remarquer que si l'homomorphisme $\mathcal{E}\to\mathcal{F}$ est surjectif,
il en est de même de l'homomorphisme correspondant
$\mathbf{S}(\mathcal{E})\to\mathbf{S}(\mathcal{F})$ de
$\mathcal{O}_S$-Algèbres, et la conclusion résulte de
(\hyperref[II.1.4.10-fr]{1.4.10}).
\end{proof}

\begin{env}[1.7.12]
\label{II.1.7.12-fr}
Prenons en particulier $\mathcal{E}=\mathcal{O}_S$~; le préschéma
$\mathbf{V}(\mathcal{O}_S)$ est le $S$-schéma affine, spectre de la
$\mathcal{O}_S$-Algèbre $\mathbf{S}(\mathcal{O}_S)$ qui s'identifie à la
$\mathcal{O}_S$-Algèbre
$\mathcal{O}_S[T]=\mathcal{O}_S\otimes_{\mathbf{Z}}\mathbf{Z}[T]$
\oldpage[II]{18}
($T$ indéterminée)~; c'est évident lorsque
$S=\operatorname{Spec}(\mathbf{Z})$, en vertu de
(\hyperref[II.1.7.6-fr]{1.7.6}), et on passe de là au cas général en
considérant le morphisme structural
$S\to\operatorname{Spec}(\mathbf{Z})$ et utilisant
(\hyperref[II.1.7.11-fr]{1.7.11, (iv)}). En raison de ce résultat, on pose
encore $\mathbf{V}(\mathcal{O}_S)=S[T]$, et on a donc la formule
\[
  S[T]=S\otimes_{\mathbf{Z}}\mathbf{Z}[T].
  \tag{1.7.12.1}\label{II.1.7.12.1-fr}
\]

L'identification du faisceau des germes de $S$-sections de $S[T]$ avec
$\mathcal{O}_S$, déjà vue dans
(\hyperref[I.3.3.15-fr]{I, 3.3.15}), se retrouve ici dans un contexte plus
général, comme cas particulier de (\hyperref[II.1.7.9-fr]{1.7.9}).
\end{env}

\begin{env}[1.7.13]
\label{II.1.7.13-fr}
Pour tout $S$-préschéma $X$, on a vu
(\hyperref[II.1.7.8-fr]{1.7.8}) que
$\operatorname{Hom}_S(X,S[T])$ s'identifie canoniquement à
$\operatorname{Hom}_{\mathcal{O}_S}(\mathcal{O}_S,\mathcal{A}(X))$, lequel
est canoniquement isomorphe à $\Gamma(S,\mathcal{A}(X))$, et est par suite
muni d'une structure d'anneau~; en outre, à tout $S$-morphisme $h:X\to Y$
correspond un morphisme
$\Gamma(\mathcal{A}(h)):\Gamma(S,\mathcal{A}(Y))
\to\Gamma(S,\mathcal{A}(X))$ pour les structures d'anneau
(\hyperref[II.1.1.2-fr]{1.1.2}). Quand on munit
$\operatorname{Hom}_S(X,S[T])$ de la structure d'anneau ainsi définie, on
voit donc qu'on peut considérer $\operatorname{Hom}_S(X,S[T])$ comme un
\emph{foncteur contravariant en $X$}, de la catégorie des $S$-préschémas
dans celle des anneaux. D'autre part,
$\operatorname{Hom}_S(X,\mathbf{V}(\mathcal{E}))$ s'identifie de même à
$\operatorname{Hom}_{\mathcal{O}_S}(\mathcal{E},\mathcal{A}(X))$ (où
$\mathcal{A}(X)$ est considéré comme un
\emph{$\mathcal{O}_S$-Module})~; on peut par suite le munir canoniquement
d'une structure de \emph{module} sur l'anneau
$\operatorname{Hom}_S(X,S[T])$, et on voit comme ci-dessus que le couple
\[
  \bigl(\operatorname{Hom}_S(X,S[T]),
  \operatorname{Hom}_S(X,\mathbf{V}(\mathcal{E}))\bigr)
\]
est un foncteur contravariant en $X$, à valeurs dans la catégorie dont les
éléments sont les couples $(A,M)$ formés d'un anneau $A$ et d'un $A$-module
$M$, les morphismes étant les di-homomorphismes.

Nous interpréterons ces faits en disant que $S[T]$ est un
\emph{$S$-schéma d'anneaux} et que $\mathbf{V}(\mathcal{E})$ est un
\emph{$S$-schéma de modules} sur le $S$-schéma d'anneaux $S[T]$
(cf. chap.~0, \textsection8).
\end{env}

\begin{env}[1.7.14]
\label{II.1.7.14-fr}
Nous allons voir que la structure de $S$-schéma de modules définie sur le
$S$-schéma $\mathbf{V}(\mathcal{E})$ permet de reconstituer le
$\mathcal{O}_S$-Module $\mathcal{E}$ à un isomorphisme unique près~: pour
cela, nous allons montrer que $\mathcal{E}$ est canoniquement isomorphe à un
sous-$\mathcal{O}_S$-Module de
$\mathbf{S}(\mathcal{E})=\mathcal{A}(\mathbf{V}(\mathcal{E}))$, défini au
moyen de cette structure. En effet
(\hyperref[II.1.7.4-fr]{1.7.4}) l'ensemble
$\operatorname{Hom}_{\mathcal{O}_S}(\mathbf{S}(\mathcal{E}),
\mathcal{A}(X))$ des homomorphismes de
\emph{$\mathcal{O}_S$-Algèbres} s'identifie canoniquement à
$\operatorname{Hom}_{\mathcal{O}_S}(\mathcal{E},\mathcal{A}(X))$, ensemble
d'homomorphismes de \emph{$\mathcal{O}_S$-Modules}~: si $h$, $h'$ sont deux
éléments de ce dernier ensemble, $s_i$ ($1\leq i\leq n$) des sections de
$\mathcal{E}$ au-dessus d'un ouvert $U\subset S$, $t$ une section de
$\mathcal{A}(X)$ au-dessus de $U$, on a par définition
\[
  (h+h')(s_1s_2\cdots s_n)
  =\prod_{i=1}^n\bigl(h(s_i)+h'(s_i)\bigr)
\]
et
\[
  (t.h)(s_1s_2\cdots s_n)=t^n\prod_{i=1}^n h(s_i).
\]

Cela étant, si $z$ est une section de $\mathbf{S}(\mathcal{E})$ au-dessus
de $U$, $h\mapsto h(z)$ est une application de
$\operatorname{Hom}_S(X,\mathbf{V}(\mathcal{E}))
=\operatorname{Hom}_{\mathcal{O}_S}(\mathbf{S}(\mathcal{E}),
\mathcal{A}(X))$ dans $\Gamma(U,\mathcal{A}(X))$. Nous allons
\oldpage[II]{19}
montrer que $\mathcal{E}$ s'identifie au sous-Module de
$\mathbf{S}(\mathcal{E})$ tel que, \emph{pour tout ouvert $U\subset S$,
toute section $z$ de ce sous-$\mathcal{O}_S$-Module au-dessus de $U$ et tout
$S$-préschéma $X$, l'application $h\mapsto h(z)$ de
$\operatorname{Hom}_{\mathcal{O}_S|U}(\mathbf{S}(\mathcal{E})|U,
\mathcal{A}(X)|U)$ dans $\Gamma(U,\mathcal{A}(X))$ soit un homomorphisme de
$\Gamma(U,\mathcal{A}(X))$-modules}.

Il est immédiat en effet que $\mathcal{E}$ possède cette propriété~; pour
démontrer la réciproque, on peut se borner à prouver que lorsque
$S=\operatorname{Spec}(A)$, $\mathcal{E}=\widetilde{M}$, une section $z$ de
$\mathbf{S}(\mathcal{E})$ au-dessus de $S$ qui (pour $U=S$) possède la
propriété énoncée ci-dessus, est nécessairement une section de
$\mathcal{E}$~; on a alors $z=\sum_{n=0}^{\infty}z_n$, où
$z_n\in\mathbf{S}_n(M)$, et il s'agit de prouver que $z_n=0$ pour $n\ne1$.
Posons $B=\mathbf{S}(M)$ et prenons pour $X$ le préschéma
$\operatorname{Spec}(B[T])$, où $T$ est une indéterminée. L'ensemble
$\operatorname{Hom}_{\mathcal{O}_S}(\mathbf{S}(\mathcal{E}),
\mathcal{A}(X))$ s'identifie à l'ensemble des homomorphismes d'anneaux
$h:B\to B[T]$ (\hyperref[I.1.3.13-fr]{I, 1.3.13}), et d'après ce qu'on a vu
plus haut, on a
$(T.h)(z)=\sum_{n=0}^{\infty}T^nh(z_n)$~: l'hypothèse sur $z$ implique que
l'on a
$\sum_{n=0}^{\infty}T^nh(z_n)=T.\sum_{n=0}^{\infty}h(z_n)$ pour tout
homomorphisme $h$. Si on prend en particulier pour $h$ l'injection
canonique, il vient
$\sum_{n=0}^{\infty}T^nz_n=T.\sum_{n=0}^{\infty}z_n$, ce qui entraîne bien
la conclusion $z_n=0$ pour $n\ne1$.
\end{env}

\begin{proposition}[1.7.15]
\label{II.1.7.15-fr}
Soit $Y$ un préschéma dont l'espace sous-jacent est noethérien, ou un schéma
quasi-compact. Tout $Y$-schéma affine $X$ de type fini sur $Y$ est
$Y$-isomorphe à un sous-$Y$-schéma fermé d'un $Y$-schéma de la forme
$\mathbf{V}(\mathcal{E})$, où $\mathcal{E}$ est un $\mathcal{O}_Y$-Module
quasi-cohérent de type fini.
\end{proposition}

\begin{proof}
En effet, la $\mathcal{O}_Y$-Algèbre quasi-cohérente $\mathcal{A}(X)$ est de
type fini (\hyperref[II.1.3.7-fr]{1.3.7}). Les hypothèses faites entraînent
que $\mathcal{A}(X)$ est engendrée par un sous-$\mathcal{O}_Y$-Module
quasi-cohérent de type fini $\mathcal{E}$
(\hyperref[I.9.6.5-fr]{I, 9.6.5})~; par définition, cela entraîne que
l'homomorphisme canonique
$\mathbf{S}(\mathcal{E})\to\mathcal{A}(X)$ prolongeant canoniquement
l'injection $\mathcal{E}\to\mathcal{A}(X)$ est \emph{surjectif}~; la
conclusion résulte alors de (\hyperref[II.1.4.10-fr]{1.4.10}).
\end{proof}

\section{Spectres premiers homogènes}
\label{section:II.2-fr}

\subsection{Généralités sur les anneaux et modules gradués.}
\label{subsection:II.2.1-fr}

\begin{notation}[2.1.1]
\label{II.2.1.1-fr}
Étant donné un anneau $S$ \emph{gradué} en degrés positifs, nous désignerons
par $S_n$ la partie de $S$ formée des éléments homogènes de degré $n$
($n\geq0$), par $S_+$ la somme (directe) des $S_n$ tels que $n>0$~; on a
$1\in S_0$, $S_0$ est un sous-anneau de $S$, $S_+$ un idéal gradué de $S$ et
$S$ est somme directe de $S_0$ et de $S_+$. Si $M$ est un \emph{module
gradué} sur $S$ (à degrés positifs ou négatifs), on désignera de même par
$M_n$ le $S_0$-module formé des éléments homogènes de degré $n$ de $M$ (avec
$n\in\mathbf{Z}$).

Pour tout entier $d>0$, on désigne par $S^{(d)}$ la somme directe des
$S_{nd}$~; en considérant les éléments de $S_{nd}$ comme homogènes de degré
$n$, les $S_{nd}$ définissent sur $S^{(d)}$ une structure d'anneau gradué.

Pour tout entier $k$ tel que $0\leq k\leq d-1$, on désignera par
$M^{(d,k)}$ la somme directe
\oldpage[II]{20}
des $M_{nd+k}$ ($n\in\mathbf{Z}$)~; c'est un $S^{(d)}$-module gradué, quand
on considère les éléments de $M_{nd+k}$ comme homogènes de degré $n$. On
écrira $M^{(d)}$ au lieu de $M^{(d,0)}$.

Avec les notations précédentes, pour tout entier $n$ (positif ou négatif),
nous désignerons par $M(n)$ le $S$-module gradué défini par
$(M(n))_k=M_{n+k}$ pour tout $k\in\mathbf{Z}$. En particulier, $S(n)$ sera
un $S$-module gradué tel que $(S(n))_k=S_{n+k}$, en convenant de poser
$S_n=0$ pour $n<0$. On dit qu'un $S$-module gradué $M$ est \emph{libre}
s'il est isomorphe, en tant que \emph{module gradué}, à une somme directe de
modules de la forme $S(n)$~; comme $S(n)$ est un $S$-module monogène,
engendré par l'élément $1$ de $S$ considéré comme élément de degré $-n$, il
revient au même de dire que $M$ admet une \emph{base} sur $S$ formée
d'éléments \emph{homogènes}.

On dit qu'un $S$-module gradué $M$ \emph{admet une présentation finie} s'il
existe une suite exacte $P\to Q\to M\to0$, où $P$ et $Q$ sont des sommes
directes finies de modules de la forme $S(n)$ et les homomorphismes sont de
degré $0$ (cf. (\hyperref[II.2.1.2-fr]{2.1.2})).
\end{notation}

\begin{env}[2.1.2]
\label{II.2.1.2-fr}
Soient $M$, $N$ deux $S$-modules gradués~; on définit sur $M\otimes_S N$ une
structure de $S$-module \emph{gradué} de la façon suivante. Sur le produit
tensoriel $M\otimes_{\mathbf{Z}}N$, on peut définir une structure de
$\mathbf{Z}$-module gradué (où $\mathbf{Z}$ est gradué par
$\mathbf{Z}_0=\mathbf{Z}$, $\mathbf{Z}_n=0$ pour $n\ne0$) en posant
\[
  (M\otimes_{\mathbf{Z}}N)_q
  =\bigoplus_{m+n=q}M_m\otimes_{\mathbf{Z}}N_n
\]
(comme $M$ et $N$ sont respectivement sommes directes des $M_m$ et des
$N_n$, on sait que l'on peut identifier canoniquement
$M\otimes_{\mathbf{Z}}N$ à la somme directe de tous les
$M_m\otimes_{\mathbf{Z}}N_n$). Cela étant, on a
$M\otimes_S N=(M\otimes_{\mathbf{Z}}N)/P$, où $P$ est le
sous-$\mathbf{Z}$-module de $M\otimes_{\mathbf{Z}}N$ engendré par les
éléments $(xs)\otimes y-x\otimes(sy)$ pour $x\in M$, $y\in N$, $s\in S$~;
il est clair que $P$ est un sous-$\mathbf{Z}$-module \emph{gradué} de
$M\otimes_{\mathbf{Z}}N$, et on voit aussitôt qu'on obtient bien en passant
au quotient une structure de $S$-module gradué sur $M\otimes_S N$.

Pour deux $S$-modules gradués $M$, $N$, rappelons qu'un homomorphisme
$u:M\to N$ de $S$-modules est dit \emph{de degré $k$} si
$u(M_j)\subset N_{j+k}$ pour tout $j\in\mathbf{Z}$. Si $H_n$ désigne
l'ensemble de tous les homomorphismes de degré $n$ de $M$ dans $N$, on
désigne par $\operatorname{Hom}_S(M,N)$ la \emph{somme} (directe) des $H_n$
($n\in\mathbf{Z}$) dans le $S$-module $H$ de tous les homomorphismes (de
$S$-module) de $M$ dans $N$~; en général $\operatorname{Hom}_S(M,N)$ n'est
pas égal à ce dernier. Toutefois, on a $H=\operatorname{Hom}_S(M,N)$ lorsque
$M$ est \emph{de type fini}~; on peut en effet supposer alors $M$ engendré
par un nombre fini d'éléments homogènes $x_i$ ($1\leq i\leq n$), et tout
homomorphisme $u\in H$ s'écrit d'une seule manière
$\sum_{k\in\mathbf{Z}}u_k$, où pour chaque $k$, $u_k(x_i)$ est égal au
composant homogène de degré $k+\deg(x_i)$ de $u(x_i)$ ($1\leq i\leq n$), ce
qui entraîne $u_k=0$ sauf pour un nombre fini d'indices~; on a par définition
$u_k\in H_k$, d'où la conclusion.

On dit que les éléments de degré $0$ de $\operatorname{Hom}_S(M,N)$ sont des
\emph{homomorphismes de $S$-modules gradués}. Il est clair que
$S_mH_n\subset H_{m+n}$, donc les $H_n$ définissent sur
$\operatorname{Hom}_S(M,N)$ une structure de $S$-module gradué.

Il résulte aussitôt de ces définitions que l'on a
\[
  \label{II.2.1.2.1-fr}
  M(m)\otimes_S N(n)=(M\otimes_S N)(m+n),
  \tag{2.1.2.1}
\]
\[
  \label{II.2.1.2.2-fr}
  \operatorname{Hom}_S(M(m),N(n))
  =(\operatorname{Hom}_S(M,N))(n-m)
  \tag{2.1.2.2}
\]
pour deux $S$-modules gradués $M$, $N$.
\oldpage[II]{21}
Soient $S$, $S'$ deux anneaux gradués~; un homomorphisme d'\emph{anneaux
gradués} $\varphi:S\to S'$ est un homomorphisme d'anneaux tel que
$\varphi(S_n)\subset S'_n$, pour tout $n\in\mathbf{Z}$ (autrement dit,
$\varphi$ doit être un homomorphisme \emph{de degré $0$} de
$\mathbf{Z}$-modules gradués). La donnée d'un tel homomorphisme définit sur
$S'$ une structure de $S$-module \emph{gradué}~; muni de cette structure et
de sa structure d'anneau gradué, on dit que $S'$ est une \emph{$S$-algèbre
graduée}.

Si alors $M$ est un $S$-module gradué, le produit tensoriel $M\otimes_S S'$
de $S$-modules gradués est muni de façon naturelle d'une structure de
$S'$-module \emph{gradué}, la graduation étant celle définie plus haut.
\end{env}

\begin{lemma}[2.1.3]
\label{II.2.1.3-fr}
Soit $S$ un anneau gradué à degrés positifs. Pour qu'une partie $E$ de $S_+$
formée d'éléments homogènes engendre $S_+$ en tant que $S$-module, il faut et
il suffit que $E$ engendre $S$ en tant que $S_0$-algèbre.
\end{lemma}

\begin{proof}
La condition est évidemment suffisante~; montrons qu'elle est nécessaire.
Soit $E_n$ (resp. $E^n$) l'ensemble des éléments de $E$ de degré égal à $n$
(resp. $\leq n$)~; il suffira de montrer, par récurrence sur $n>0$, que $S_n$
est le $S_0$-module engendré par les éléments de degré $n$ qui sont des
produits d'éléments de $E^n$. Or, cela est évident pour $n=1$ en vertu de
l'hypothèse~; cette dernière montre en outre que
$S_n=\sum_{p=0}^{n-1}S_pE_{n-p}$, et le raisonnement par récurrence est alors
immédiat.
\end{proof}

\begin{corollary}[2.1.4]
\label{II.2.1.4-fr}
Pour que $S_+$ soit un idéal de type fini, il faut et il suffit que $S$ soit
une $S_0$-algèbre de type fini.
\end{corollary}

\begin{proof}
On peut en effet toujours supposer qu'un système fini de générateurs de la
$S_0$-algèbre $S$ (resp. du $S$-idéal $S_+$) est formé d'éléments homogènes,
en remplaçant chacun des générateurs considérés par ses composants homogènes.
\end{proof}

\begin{corollary}[2.1.5]
\label{II.2.1.5-fr}
Pour que $S$ soit noethérien, il faut et il suffit que $S_0$ soit noethérien
et que $S$ soit une $S_0$-algèbre de type fini.
\end{corollary}

\begin{proof}
La condition est évidemment suffisante~; elle est nécessaire, puisque $S_0$
est isomorphe à $S/S_+$ et que $S_+$ doit être un idéal de type fini
(\hyperref[II.2.1.4-fr]{2.1.4}).
\end{proof}

\begin{lemma}[2.1.6]
\label{II.2.1.6-fr}
Soit $S$ un anneau gradué à degrés positifs, qui soit une $S_0$-algèbre de
type fini. Soit $M$ un $S$-module gradué de type fini. Alors~:
\begin{enumerate}
  \item[\rm(i)] Les $M_n$ sont des $S_0$-modules de type fini, et il existe
    un entier $n_0$ tel que $M_n=0$ pour $n\leq n_0$.
  \item[\rm(ii)] Il existe un entier $n_1$ et un entier $h>0$ tel que, pour
    tout entier $n\geq n_1$, on ait $M_{n+h}=S_hM_n$.
  \item[\rm(iii)] Pour tout couple d'entiers $(d,k)$ tels que $d>0$,
    $0\leq k\leq d-1$, $M^{(d,k)}$ est un $S^{(d)}$-module de type fini.
  \item[\rm(iv)] Pour tout entier $d>0$, $S^{(d)}$ est une $S_0$-algèbre de
    type fini.
  \item[\rm(v)] Il existe un entier $h>0$ tel que $S_{mh}=(S_h)^m$ pour tout
    $m>0$.
  \item[\rm(vi)] Pour tout entier $n>0$, il existe un entier $m_0$ tel que
    $S_m\subset S_+^n$ pour tout $m\geq m_0$.
\end{enumerate}
\end{lemma}

\begin{proof}
On peut supposer $S$ engendré (en tant que $S_0$-algèbre) par des éléments
homogènes $f_i$, de degré $h_i$ ($1\leq i\leq r$), et $M$ engendré (en tant
que $S$-module) par des éléments homogènes $x_j$ de degré $k_j$
($1\leq j\leq s$). Il est clair que $M_n$ est formé des combinaisons
\oldpage[II]{22}
linéaires, à coefficients dans $S_0$, des éléments
$f_1^{\alpha_1}\cdots f_r^{\alpha_r}x_j$ tels que les $\alpha_i$ soient des
entiers $\geq 0$ vérifiant $k_j+\sum_i\alpha_i h_i=n$~; pour chaque $j$, il
n'y a qu'un nombre fini de systèmes $(\alpha_i)$ vérifiant cette équation,
puisque les $h_i$ sont $>0$, d'où la première assertion de (i)~; la seconde
est évidente. Soit d'autre part $h$ le p.p.c.m. des $h_i$ et posons
$g_i=f_i^{h/h_i}$ ($1\leq i\leq r$) de sorte que tous les $g_i$ sont de
degré $h$~; soient $z_\mu$ les éléments de $M$ de la forme
$f_1^{\alpha_1}\cdots f_r^{\alpha_r}x_j$ avec
$0\leq\alpha_i<h/h_i$ pour $1\leq i\leq r$~; ces éléments sont en nombre
fini, soit $n_1$ le plus grand de leurs degrés. Il est clair que pour
$n\geq n_1$, tout élément de $M_{n+h}$ est combinaison linéaire des $z_\mu$
dont les coefficients sont des monômes de degré $>0$ par rapport aux $g_i$,
donc on a $M_{n+h}=S_hM_n$, ce qui établit (ii). De la même manière, on voit
(pour tout $d>0$) qu'un élément de $M^{(d,k)}$ est combinaison linéaire, à
coefficients dans $S_0$, d'éléments de la forme
$g^df_1^{\alpha_1}\cdots f_r^{\alpha_r}x_j$ avec
$0\leq\alpha_i<d$, $g$ étant un élément homogène de $S$~; d'où (iii)~; (iv)
résulte alors de (iii) et du lemme (\hyperref[II.2.1.3-fr]{2.1.3}), en
prenant $M=S_+$, puisque $(S_+)^{(d)}=(S^{(d)})_+$. L'assertion de (v) se
déduit de (ii) en prenant $M=S$. Enfin, pour un $n$ donné, il n'y a qu'un
nombre fini de systèmes $(\alpha_i)$ tels que $\alpha_i\geq 0$ et
$\sum_i\alpha_i<n$, donc si $m_0$ est la plus grande valeur de la somme
$\sum_i\alpha_i h_i$ pour ces systèmes, on a $S_m\subset S_+^n$ pour
$m>m_0$, ce qui prouve (vi).
\end{proof}

\begin{corollary}[2.1.7]
\label{II.2.1.7-fr}
Si $S$ est noethérien, il en est de même de $S^{(d)}$ pour tout entier
$d>0$.
\end{corollary}

\begin{proof}
Cela résulte de (\hyperref[II.2.1.5-fr]{2.1.5}) et de
(\hyperref[II.2.1.6-fr]{2.1.6, (iv)}).
\end{proof}

\begin{env}[2.1.8]
\label{II.2.1.8-fr}
Soit $\mathfrak{p}$ un idéal premier gradué de l'anneau gradué $S$~;
$\mathfrak{p}$ est donc somme directe des sous-groupes
$\mathfrak{p}_n=\mathfrak{p}\cap S_n$. Supposons que $\mathfrak{p}$ ne
contienne pas $S_+$. Alors, si $f\in S_+$, n'appartient pas à
$\mathfrak{p}$, la relation $f^nx\in\mathfrak{p}$ est équivalente à
$x\in\mathfrak{p}$~; en particulier, si $f\in S_d$ ($d>0$), pour tout
$x\in S_{m-nd}$, la relation $f^nx\in\mathfrak{p}_m$ équivaut à
$x\in\mathfrak{p}_{m-nd}$.
\end{env}

\begin{proposition}[2.1.9]
\label{II.2.1.9-fr}
Soit $n_0$ un entier $>0$~; pour tout $n\geq n_0$ soit $\mathfrak{p}_n$ un
sous-groupe de $S_n$. Pour qu'il existe un idéal premier gradué
$\mathfrak{p}$ de $S$ ne contenant pas $S_+$ et tel que
$\mathfrak{p}\cap S_n=\mathfrak{p}_n$ pour tout $n\geq n_0$, il faut et il
suffit que les conditions suivantes soient vérifiées~:
\begin{enumerate}
  \item[$1^\circ$] $S_m\mathfrak{p}_n\subset\mathfrak{p}_{m+n}$ pour tout
    $m\geq 0$ et tout $n\geq n_0$.
  \item[$2^\circ$] Pour $m\geq n_0$, $n\geq n_0$, $f\in S_m$, $g\in S_n$,
    la relation $fg\in\mathfrak{p}_{m+n}$ entraîne
    $f\in\mathfrak{p}_m$ ou $g\in\mathfrak{p}_n$.
  \item[$3^\circ$] $\mathfrak{p}_n\neq S_n$ pour un $n\geq n_0$ au moins.
\end{enumerate}
En outre l'idéal premier gradué $\mathfrak{p}$ est alors unique.
\end{proposition}

\begin{proof}
Il est évident que les conditions $1^\circ$ et $2^\circ$ sont nécessaires.
En outre, si $\mathfrak{p}\not\supset S_+$, il existe au moins un $k>0$ tel que
$\mathfrak{p}\cap S_k\neq S_k$~; si $f\in S_k$ n'appartient pas à
$\mathfrak{p}$, la relation $\mathfrak{p}\cap S_n=S_n$ entraîne
$\mathfrak{p}\cap S_{n-mk}=S_{n-mk}$ d'après
(\hyperref[II.2.1.8-fr]{2.1.8})~; donc, si
$\mathfrak{p}\cap S_n=S_n$ à partir d'une certaine valeur de $n$, on aurait
$\mathfrak{p}\supset S_+$ contrairement à l'hypothèse, ce qui prouve que
$3^\circ$ est nécessaire. Inversement, supposons satisfaites les conditions
$1^\circ$, $2^\circ$ et $3^\circ$. Notons que, si pour un entier
$d\geq n_0$, $f\in S_d$ n'appartient pas à $\mathfrak{p}_d$, alors, si
$\mathfrak{p}$ existe, $\mathfrak{p}_m$, pour $m<n_0$, est nécessairement
égal à l'ensemble des $x\in S_m$ tels que
$f^rx\in\mathfrak{p}_{m+rd}$, sauf pour un nombre fini de valeurs de $r$.
Cela prouve déjà que si $\mathfrak{p}$ existe, il est unique. Reste à montrer
que si on définit les $\mathfrak{p}_m$ pour $m<n_0$ par la condition
précédente, $\mathfrak{p}=\sum_{n=0}^{\infty}\mathfrak{p}_n$ est un idéal
premier. Notons d'abord qu'en vertu de $2^\circ$, pour $m\geq n_0$,
$\mathfrak{p}_m$ est aussi défini comme l'ensemble des $x\in S_m$ tels que
$f^rx\in\mathfrak{p}_{m+rd}$, sauf pour un nombre fini de valeurs de $r$.
Cela
\oldpage[II]{23}
étant, si $g\in S_m$, $x\in\mathfrak{p}_n$, on a
$f^rgx\in\mathfrak{p}_{m+n+rd}$ sauf pour un nombre fini de valeurs de $r$,
donc $gx\in\mathfrak{p}_{m+n}$, ce qui prouve que $\mathfrak{p}$ est un
idéal de $S$. Pour établir que cet idéal est premier, autrement dit que
l'anneau $S/\mathfrak{p}$, gradué par les sous-groupes
$S_n/\mathfrak{p}_n$, est un anneau intègre, il suffit (en considérant les
composants de plus haut degré de deux éléments de $S/\mathfrak{p}$) de
prouver que si $x\in S_m$, $y\in S_n$ sont tels que
$x\notin\mathfrak{p}_m$, $y\notin\mathfrak{p}_n$, alors
$xy\notin\mathfrak{p}_{m+n}$. Sinon, pour $r$ assez grand, on aurait
$f^{2r}xy\in\mathfrak{p}_{m+n+2rd}$~; mais on a
$f^ry\notin\mathfrak{p}_{n+rd}$ pour tout $r>0$~; il résulte alors de
$2^\circ$ que, sauf pour un nombre fini de valeurs de $r$, on a
$f^rx\in\mathfrak{p}_{m+rd}$, et on en conclurait que
$x\in\mathfrak{p}_m$ contrairement à l'hypothèse.
\end{proof}

\begin{env}[2.1.10]
\label{II.2.1.10-fr}
Nous dirons qu'une partie $\mathfrak{J}$ de $S_+$ est un
\emph{idéal de $S_+$} si c'est un idéal de $S$, et que $\mathfrak{J}$ est un
\emph{idéal premier gradué de $S_+$} si elle est l'intersection de $S_+$ et
d'un idéal premier gradué de $S$ \emph{ne contenant pas $S_+$} (cet idéal
premier est d'ailleurs unique d'après
(\hyperref[II.2.1.9-fr]{2.1.9})). Si $\mathfrak{J}$ est un idéal de $S_+$,
la \emph{racine de $\mathfrak{J}$ dans $S_+$} est l'ensemble des éléments de
$S_+$ dont une puissance appartient à $\mathfrak{J}$, autrement dit c'est
l'ensemble $r_+(\mathfrak{J})=r(\mathfrak{J})\cap S_+$~; en particulier, la
racine de $0$ dans $S_+$ est encore appelée le \emph{nilradical} de $S_+$ et
notée $\mathfrak{N}_+$~: c'est donc l'ensemble des éléments nilpotents de
$S_+$. Si $\mathfrak{J}$ est un idéal \emph{gradué} de $S_+$, sa racine
$r_+(\mathfrak{J})$ est un idéal gradué~: en passant à l'anneau quotient
$S/\mathfrak{J}$, on peut se ramener au cas $\mathfrak{J}=0$, et tout
revient à voir que si $x=x_h+x_{h+1}+\cdots+x_k$ est nilpotent, il en est de
même des $x_i\in S_i$ ($1\leq h\leq i\leq k$)~; on peut supposer
$x_k\neq 0$ et le composant de plus grand degré de $x^n$ est alors $x_k^n$,
donc $x_k$ est nilpotent, et on raisonne alors par récurrence sur $k$. On dit
que l'anneau gradué $S$ est \emph{essentiellement réduit} si
$\mathfrak{N}_+=0$, autrement dit si $S_+$ ne contient pas d'éléments
nilpotents $\neq 0$.
\end{env}

\begin{env}[2.1.11]
\label{II.2.1.11-fr}
On notera que si, dans l'anneau gradué $S$, un élément $x$ est diviseur de
$0$, il en est de même de son composant de plus haut degré. On dit qu'un
anneau $S$ est \emph{essentiellement intègre} si l'anneau $S_+$
(\emph{sans élément unité}) ne contient pas de diviseur de $0$ et est
$\neq 0$~; il suffit pour cela qu'un élément homogène $\neq 0$ dans $S_+$ ne
soit pas diviseur de $0$ dans cet anneau. Il est clair que si $\mathfrak{p}$
est un idéal premier gradué de $S_+$, $S/\mathfrak{p}$ est essentiellement
intègre.

Soit $S$ un anneau gradué essentiellement intègre, et soit $x_0\in S_0$~;
s'il existe alors \emph{un} élément homogène $f\neq 0$ de $S_+$ tel que
$x_0f=0$, on a $x_0S_+=0$, car on a $(x_0g)f=(x_0f)g=0$ pour tout
$g\in S_+$, et l'hypothèse entraîne donc $x_0g=0$. Pour que $S$ soit
intègre, il faut et il suffit donc que $S_0$ soit intègre et que l'annulateur
de $S_+$ dans $S_0$ soit réduit à $0$.
\end{env}

\subsection{Anneaux de fractions d'un anneau gradué.}
\label{subsection:II.2.2-fr}

\begin{env}[2.2.1]
\label{II.2.2.1-fr}
Soient $S$ un anneau gradué, à degrés positifs, $f$ un élément
\emph{homogène} de $S$, de degré $d>0$~; alors l'anneau de fractions
$S'=S_f$ est gradué, en prenant pour $S'_n$ l'ensemble des $x/f^k$, où
$x\in S_{n+kd}$ avec $k\geq 0$ (on observera ici que $n$ peut prendre des
valeurs négatives arbitraires)~; nous désignerons le sous-anneau
$S'_0=(S_f)_0$ de $S'$ formé des éléments \emph{de degré $0$} par la
notation $S_{(f)}$.

Si $f\in S_d$, les monômes $(f/1)^h$ dans $S_f$ ($h$ entier positif ou
négatif) forment un \emph{système libre} sur l'anneau $S_{(f)}$, et
l'ensemble de leurs combinaisons linéaires n'est autre que
\oldpage[II]{24}
l'anneau $(S^{(d)})_f$, qui est donc \emph{isomorphe à
$S_{(f)}[T,T^{-1}]=S_{(f)}\otimes_{\mathbf{Z}}\mathbf{Z}[T,T^{-1}]$}
(où $T$ est une indéterminée). En effet, si on a une relation
$\sum_{h=-a}^{b}z_h(f/1)^h=0$ avec $z_h=x_h/f^m$, où tous les $x_h$
appartiennent à $S_{md}$, cette relation équivaut par définition à
l'existence d'un $k\geq a$ tel que
$\sum_{h=-a}^{b}f^{h+k}x_h=0$, et comme les degrés des termes de cette somme
sont distincts, on a $f^{h+k}x_h=0$ pour tout $h$, d'où $z_h=0$ pour tout
$h$.

Si $M$ est un $S$-module gradué, $M'=M_f$ est un $S_f$-module gradué,
$M'_n$ étant l'ensemble des $z/f^k$ avec $z\in M_{n+kd}$ ($k\geq 0$)~; on
désigne par $M_{(f)}$ l'ensemble des éléments homogènes de degré $0$ de
$M'$~; il est immédiat que $M_{(f)}$ est un $S_{(f)}$-module et que l'on a
$(M^{(d)})_f=M_{(f)}\otimes_{S_{(f)}}(S^{(d)})_f$.
\end{env}

\begin{lemma}[2.2.2]
\label{II.2.2.2-fr}
Soient $d$, $e$ deux entiers $>0$, $f\in S_d$, $g\in S_e$. Il existe un
isomorphisme canonique d'anneaux
\[
  S_{(fg)}\xrightarrow{\sim}(S_{(f)})_{g^d/f^e}~;
\]
si on identifie canoniquement ces deux anneaux, il existe un isomorphisme
canonique de modules
\[
  M_{(fg)}\xrightarrow{\sim}(M_{(f)})_{g^d/f^e}.
\]
\end{lemma}

\begin{proof}
En effet, $fg$ divise $f^e g^d$, et ce dernier élément divise $(fg)^{de}$,
donc les anneaux gradués $S_{fg}$ et $S_{f^e g^d}$ s'identifient
canoniquement~; d'autre part, $S_{f^e g^d}$ s'identifie aussi à
$(S_{f^e})_{g^d/1}$ (\hyperref[0.1.4.6-fr]{0, I.4.6}), et comme $f^e/1$ est
inversible dans $S_{f^e}$, $S_{f^e g^d}$ s'identifie aussi à
$(S_{f^e})_{g^d/f^e}$. Or l'élément $g^d/f^e$ est de degré $0$ dans
$S_{f^e}$~; on en conclut aussitôt que le sous-anneau de
$(S_{f^e})_{g^d/f^e}$ formé des éléments de degré $0$ est
$(S_{(f^e)})_{g^d/f^e}$, et comme on a évidemment
$S_{(f^e)}=S_{(f)}$, cela démontre la première partie de la proposition~;
la seconde s'établit de même.
\end{proof}

\begin{env}[2.2.3]
\label{II.2.2.3-fr}
Avec les hypothèses de (\hyperref[II.2.2.2-fr]{2.2.2}), il est clair que
l'homomorphisme canonique $S_f\to S_{fg}$
(\hyperref[0.1.4.1-fr]{0, I.4.1}), qui transforme $x/f^k$ en
$g^k x/(fg)^k$ est de degré $0$, donc donne par restriction un
\emph{homomorphisme canonique $S_{(f)}\to S_{(fg)}$}, tel que le diagramme
\[
  \xymatrix{
    & S_{(f)}\ar[dl]\ar[dr]\\
    S_{(fg)}\ar[rr]^-{\sim} && (S_{(f)})_{g^d/f^e}
  }
\]
soit commutatif. On définit de même un homomorphisme canonique
$M_{(f)}\to M_{(fg)}$.
\end{env}

\begin{lemma}[2.2.4]
\label{II.2.2.4-fr}
Si $f$, $g$ sont deux éléments homogènes de $S_+$, l'anneau $S_{(fg)}$ est
engendré par la réunion des images canoniques de $S_{(f)}$ et $S_{(g)}$.
\end{lemma}

\begin{proof}
En vertu de (\hyperref[II.2.2.2-fr]{2.2.2}), il suffit de voir que
$1/(g^d/f^e)=f^{d+e}/(fg)^d$ appartient à l'image canonique de $S_{(g)}$
dans $S_{(fg)}$, ce qui est évident par définition.
\end{proof}

\begin{proposition}[2.2.5]
\label{II.2.2.5-fr}
Soit $d$ un entier $>0$ et soit $f\in S_d$. Alors il existe un isomorphisme
canonique d'anneaux
$S_{(f)}\xrightarrow{\sim}S^{(d)}/(f-1)S^{(d)}$~; si on identifie ces deux
anneaux par cet isomorphisme, il existe un isomorphisme canonique de modules
$M_{(f)}\xrightarrow{\sim}M^{(d)}/(f-1)M^{(d)}$.
\end{proposition}

\begin{proof}
Le premier de ces isomorphismes se définit en faisant correspondre à
$x/f^n$, où $x\in S_{nd}$, l'élément $\overline{x}$, classe de $x$ mod.
$(f-1)S^{(d)}$~; cette application est bien définie, car on a la congruence
$f^h x\equiv x\pmod{(f-1)S^{(d)}}$ pour tout $x\in S^{(d)}$, donc si
$f^h x=0$ pour un $h>0$,
\oldpage[II]{25}
on a $\overline{x}=0$. D'autre part, si $x\in S_{nd}$ est tel que
$x=(f-1)y$ avec
$y=y_{hd}+y_{(h+1)d}+\cdots+y_{kd}$ avec $y_{jd}\in S_{jd}$ et
$y_{hd}\neq 0$, on a nécessairement $h=n$ et $x=-y_{hd}$, ainsi que les
relations $y_{(j+1)d}=fy_{jd}$ pour $h\leq j\leq k-1$, $fy_{kd}=0$, ce qui
donne finalement $f^{k-n+1}x=0$~; à toute classe $\overline{x}$ mod.
$(f-1)S^{(d)}$ d'un élément $x\in S_{nd}$, on peut donc faire correspondre
l'élément $x/f^n$ de $S_{(f)}$, car la remarque qui précède montre que cette
application est bien définie. Il est immédiat que les deux applications ainsi
définies sont des homomorphismes d'anneaux, réciproques l'un de l'autre. On
procède exactement de même pour $M$.
\end{proof}

\begin{corollary}[2.2.6]
\label{II.2.2.6-fr}
Si $S$ est noethérien, il en est de même de $S_{(f)}$ pour tout $f$ homogène
de degré $>0$.
\end{corollary}

\begin{proof}
Cela résulte aussitôt de (\hyperref[II.2.1.7-fr]{2.1.7}) et
(\hyperref[II.2.2.5-fr]{2.2.5}).
\end{proof}

\begin{env}[2.2.7]
\label{II.2.2.7-fr}
Soit $T$ une partie multiplicative de $S_+$ formée d'éléments
\emph{homogènes}~; $T_0=T\cup\{1\}$ est alors une partie multiplicative de
$S$~; comme les éléments de $T_0$ sont homogènes, l'anneau $T_0^{-1}S$ est
encore gradué de façon évidente~; on désignera par $S_{(T)}$ le sous-anneau
de $T_0^{-1}S$ formé des éléments d'ordre $0$, c'est-à-dire des éléments de
la forme $x/h$, où $h\in T$ et $x$ est homogène de degré égal à celui de
$h$. On sait (\hyperref[0.1.4.5-fr]{0, 1.4.5}) que $T_0^{-1}S$ s'identifie
canoniquement à la limite inductive des anneaux $S_f$, où $f$ parcourt $T$
(pour les homomorphismes canoniques $S_f\to S_{fg}$)~; comme cette
identification respecte les degrés, elle identifie $S_{(T)}$ à la
\emph{limite inductive} des $S_{(f)}$ pour $f\in T$. Pour tout $S$-module
gradué $M$, on définit de même le module $M_{(T)}$ (sur l'anneau $S_{(T)}$)
formé des éléments de degré $0$ de $T_0^{-1}M$, et on voit que ce module est
limite inductive des $M_{(f)}$ pour $f\in T$.

Si $\mathfrak{p}$ est un idéal premier gradué de $S_+$, on désignera par
$S_{(\mathfrak{p})}$ et $M_{(\mathfrak{p})}$ l'anneau $S_{(T)}$ et le
module $M_{(T)}$ respectivement, où $T$ est l'ensemble des éléments
\emph{homogènes} de $S_+$ qui n'appartiennent pas à $\mathfrak{p}$.
\end{env}

\subsection{Spectre premier homogène d'un anneau gradué.}
\label{subsection:II.2.3-fr}

\begin{env}[2.3.1]
\label{II.2.3.1-fr}
Étant donné un anneau gradué $S$, à degrés positifs, on appelle
\emph{spectre premier homogène} de $S$ et on désigne par
$\operatorname{Proj}(S)$ l'ensemble des idéaux premiers gradués de $S_+$
(\hyperref[II.2.1.10-fr]{2.1.10}), ou, ce qui revient au même, l'ensemble
des idéaux premiers gradués de $S$ \emph{ne contenant pas} $S_+$~; nous
allons définir une structure de \emph{schéma} ayant
$\operatorname{Proj}(S)$ comme ensemble de base.
\end{env}

\begin{env}[2.3.2]
\label{II.2.3.2-fr}
Pour toute partie $E$ de $S$, soit $V_+(E)$ l'ensemble des idéaux premiers
gradués de $S$ contenant $E$ et ne contenant pas $S_+$~; c'est donc la
partie $V(E)\cap\operatorname{Proj}(S)$ de
$\operatorname{Spec}(S)$. De (\hyperref[I.1.1.2-fr]{I, 1.1.2}) on déduit
donc~:
\[
  \label{II.2.3.2.1-fr}
  V_+(0)=\operatorname{Proj}(S),\quad
  V_+(S)=V_+(S_+)=\varnothing
  \tag{2.3.2.1}
\]
\[
  \label{II.2.3.2.2-fr}
  V_+\left(\bigcup_\lambda E_\lambda\right)
  =\bigcap_\lambda V_+(E_\lambda)
  \tag{2.3.2.2}
\]
\[
  \label{II.2.3.2.3-fr}
  V_+(EE')=V_+(E)\cup V_+(E')
  \tag{2.3.2.3}
\]
On ne change pas $V_+(E)$ en remplaçant $E$ par l'idéal gradué engendré par
$E$~; en outre, si $\mathfrak{J}$ est un idéal gradué de $S$, on a
\[
  \label{II.2.3.2.4-fr}
  V_+(\mathfrak{J})
  =V_+\left(\bigcup_{q\geq n}(\mathfrak{J}\cap S_q)\right)
  \tag{2.3.2.4}
\]
\oldpage[II]{26}
pour tout $n>0$~: en effet, si
$\mathfrak{p}\in\operatorname{Proj}(S)$ contient les éléments homogènes de
$\mathfrak{J}$ de degré $\geq n$, comme par hypothèse il existe un élément
homogène $f\in S_d$ non contenu dans $\mathfrak{p}$, pour tout $m\geq 0$ et
tout $x\in S_m\cap\mathfrak{J}$, on a
$f^r x\in\mathfrak{J}\cap S_{m+rd}$ sauf pour un nombre fini de valeurs de
$r$, donc $f^r x\in\mathfrak{p}\cap S_{m+rd}$, ce qui entraîne
$x\in\mathfrak{p}\cap S_m$
(\hyperref[II.2.1.9-fr]{2.1.9}).

Enfin, on a, pour tout idéal gradué $\mathfrak{J}$ de $S$
\[
  \label{II.2.3.2.5-fr}
  V_+(\mathfrak{J})=V_+(r_+(\mathfrak{J})).
  \tag{2.3.2.5}
\]
\end{env}

\begin{env}[2.3.3]
\label{II.2.3.3-fr}
Par définition, les $V_+(E)$ sont les parties fermées de
$X=\operatorname{Proj}(S)$ pour la topologie induite par la topologie
spectrale de $\operatorname{Spec}(S)$, et qu'on appellera encore
\emph{topologie spectrale} sur $X$. On pose, pour tout $f\in S$
\[
  \label{II.2.3.3.1-fr}
  D_+(f)=D(f)\cap\operatorname{Proj}(S)
  =\operatorname{Proj}(S)\setminus V_+(f)
  \tag{2.3.3.1}
\]
et on a par suite, pour deux éléments $f$, $g$ de $S$
(\hyperref[I.1.1.9.1-fr]{I, 1.1.9.1})
\[
  \label{II.2.3.3.2-fr}
  D_+(fg)=D_+(f)\cap D_+(g).
  \tag{2.3.3.2}
\]
\end{env}

\begin{proposition}[2.3.4]
\label{II.2.3.4-fr}
Lorsque $f$ parcourt l'ensemble des éléments homogènes de $S_+$, les
$D_+(f)$ forment une base de la topologie de
$X=\operatorname{Proj}(S)$.
\end{proposition}

\begin{proof}
Il résulte en effet de (\hyperref[II.2.3.2.2-fr]{2.3.2.2}) et
(\hyperref[II.2.3.2.4-fr]{2.3.2.4}) que toute partie fermée de $X$ est
intersection d'ensembles de la forme $V_+(f)$, où $f$ est homogène de degré
$>0$.
\end{proof}

\begin{env}[2.3.5]
\label{II.2.3.5-fr}
Soit $f$ un élément \emph{homogène} de $S_+$, de degré $d>0$~; pour tout
idéal premier gradué $\mathfrak{p}$ de $S$, ne contenant pas $f$, on sait
que l'ensemble des $x/f^n$, où $x\in\mathfrak{p}$ et $n\geq 0$, est un
idéal premier de l'anneau de fractions $S_f$
(\hyperref[0.1.2.6-fr]{0, 1.2.6})~; sa trace sur $S_{(f)}$ est donc un idéal
premier de cet anneau, que nous désignerons par
$\psi_f(\mathfrak{p})$~: c'est l'ensemble des $x/f^n$ pour $n\geq 0$,
$x\in\mathfrak{p}\cap S_{nd}$. On a ainsi défini une application
\[
  \psi_f:D_+(f)\longrightarrow\operatorname{Spec}(S_{(f)})~;
\]
en outre, si $g\in S_e$ est un second élément homogène de $S_+$, on a un
diagramme commutatif
\[
  \label{II.2.3.5.1-fr}
  \xymatrix{
    D_+(f)\ar[r]^{\psi_f}
    & \operatorname{Spec}(S_{(f)})\\
    D_+(fg)\ar[u]\ar[r]_{\psi_{fg}}
    & \operatorname{Spec}(S_{(fg)})\ar[u]
  }
  \tag{2.3.5.1}
\]
où la flèche verticale de gauche est l'inclusion, et celle de droite est
l'application ${}^a\omega_{fg,f}$ déduite de l'homomorphisme canonique
$\omega=\omega_{fg,f}:S_{(f)}\to S_{(fg)}$
(\hyperref[I.1.2.1-fr]{I, 1.2.1}). En effet, si
$x/f^n\in\omega^{-1}(\psi_{fg}(\mathfrak{p}))$, où
$fg\notin\mathfrak{p}$, on a par définition
$g^n x/(fg)^n\in\psi_{fg}(\mathfrak{p})$, donc
$g^n x\in\mathfrak{p}$ et par suite $x\in\mathfrak{p}$, et la réciproque
est évidente.
\end{env}

\begin{proposition}[2.3.6]
\label{II.2.3.6-fr}
L'application $\psi_f$ est un homéomorphisme de $D_+(f)$ sur
$\operatorname{Spec}(S_{(f)})$.
\end{proposition}

\begin{proof}
En premier lieu, $\psi_f$ est continue~; car si $h\in S_{nd}$ est tel que
$h/f^n\in\psi_f(\mathfrak{p})$, on a par définition
$h\in\mathfrak{p}$ et réciproquement, donc
\[
  \psi_f^{-1}(D(h/f^n))=D_+(hf)
\]
et notre assertion résulte de la formule
(\hyperref[II.2.3.3.2-fr]{2.3.3.2}). En outre, les $D_+(hf)$, où $h$
parcourt les ensembles $S_{nd}$, forment une base de la topologie de
$D_+(f)$, en vertu de (\hyperref[II.2.3.4-fr]{2.3.4}) et de la formule
(\hyperref[II.2.3.3.2-fr]{2.3.3.2})~; ce
\oldpage[II]{27}
qui précède prouve donc, compte tenu de l'axiome $(T_0)$ valable dans
$D_+(f)$ et dans $\operatorname{Spec}(S_{(f)})$, que $\psi_f$ est injective
et que l'application réciproque
$\psi_f(D_+(f))\longrightarrow D_+(f)$ est continue. Enfin, pour voir que
$\psi_f$ est surjective, on remarque que, si $\mathfrak{q}_0$ est un idéal
premier de $S_{(f)}$, et si, pour tout $n>0$, on désigne par
$\mathfrak{p}_n$ l'ensemble des $x\in S_n$ tels que
$x^d/f^n\in\mathfrak{q}_0$, les $\mathfrak{p}_n$ vérifient les conditions
de (\hyperref[II.2.1.9-fr]{2.1.9})~: en effet, si $x\in S_n$, $y\in S_n$
sont tels que $x^d/f^n\in\mathfrak{q}_0$ et
$y^d/f^n\in\mathfrak{q}_0$, on a
$(x+y)^{2d}/f^{2n}\in\mathfrak{q}_0$, d'où
$(x+y)^d/f^n\in\mathfrak{q}_0$ puisque $\mathfrak{q}_0$ est premier~;
cela prouve que les $\mathfrak{p}_n$ sont des sous-groupes des $S_n$, et la
vérification des autres conditions de
(\hyperref[II.2.1.9-fr]{2.1.9}) est immédiate, compte tenu de ce que
$\mathfrak{q}_0$ est premier. Si $\mathfrak{p}$ est l'idéal premier gradué
de $S$ ainsi défini, on a bien
$\psi_f(\mathfrak{p})=\mathfrak{q}_0$, car si $x\in S_{nd}$, les relations
$x/f^n\in\mathfrak{q}_0$ et $x^d/f^{nd}\in\mathfrak{q}_0$ sont
équivalentes, $\mathfrak{q}_0$ étant premier.
\end{proof}

\begin{corollary}[2.3.7]
\label{II.2.3.7-fr}
Pour que $D_+(f)=\varnothing$, il faut et il suffit que $f$ soit nilpotent.
\end{corollary}

\begin{proof}
En effet, pour que $\operatorname{Spec}(S_{(f)})=\varnothing$, il faut et
il suffit que $S_{(f)}=0$, ou encore que $1=0$ dans $S_f$, ce qui signifie
par définition que $f$ est nilpotent.
\end{proof}

\begin{corollary}[2.3.8]
\label{II.2.3.8-fr}
Soit $E$ une partie de $S_+$. Les conditions suivantes sont équivalentes~:
\begin{enumerate}
  \item[\rm a)] $V_+(E)=X=\operatorname{Proj}(S)$.
  \item[\rm b)] Tout élément de $E$ est nilpotent.
  \item[\rm c)] Les composants homogènes de tout élément de $E$ sont
    nilpotents.
\end{enumerate}
\end{corollary}

\begin{proof}
Il est clair que c) entraîne b) et que b) entraîne a). Si $\mathfrak{J}$ est
l'idéal gradué de $S$ engendré par $E$, la condition a) équivaut à
$V_+(\mathfrak{J})=X$~; \emph{a fortiori}, a) entraîne que tout élément
homogène $f\in\mathfrak{J}$ est tel que $V_+(f)=X$, donc $f$ est nilpotent
en vertu de (\hyperref[II.2.3.7-fr]{2.3.7}).
\end{proof}

\begin{corollary}[2.3.9]
\label{II.2.3.9-fr}
Si $\mathfrak{J}$ est un idéal gradué de $S_+$,
$r_+(\mathfrak{J})$ est l'intersection des idéaux premiers gradués de
$S_+$ qui contiennent $\mathfrak{J}$.
\end{corollary}

\begin{proof}
En considérant l'anneau gradué $S/\mathfrak{J}$, on peut se ramener au cas
où $\mathfrak{J}=0$. Il faut prouver que si $f\in S_+$ n'est pas
nilpotent, il y a un idéal premier gradué de $S$ ne contenant pas $f$~; or,
un au moins des composants homogènes de $f$ n'est pas nilpotent, et on peut
donc supposer $f$ homogène~; l'assertion résulte alors de
(\hyperref[II.2.3.7-fr]{2.3.7}).
\end{proof}

\begin{env}[2.3.10]
\label{II.2.3.10-fr}
Pour toute partie $Y$ de $X=\operatorname{Proj}(S)$, soit
$\mathfrak{j}_+(Y)$ l'ensemble des $f\in S_+$ tels que
$Y\subset V_+(f)$~; il revient au même de dire que
$\mathfrak{j}_+(Y)=\mathfrak{j}(Y)\cap S_+$~; $\mathfrak{j}_+(Y)$ est
donc un idéal de $S_+$, égal à sa racine dans $S_+$.
\end{env}

\begin{proposition}[2.3.11]
\label{II.2.3.11-fr}
\begin{enumerate}
  \item[\rm(i)] Pour toute partie $E$ de $S_+$,
    $\mathfrak{j}_+(V_+(E))$ est la racine dans $S_+$ de l'idéal gradué de
    $S_+$ engendré par $E$.
  \item[\rm(ii)] Pour toute partie $Y$ de $X$,
    $V_+(\mathfrak{j}_+(Y))=\overline{Y}$, adhérence de $Y$ dans $X$.
\end{enumerate}
\end{proposition}

\begin{proof}
\begin{enumerate}
  \item[\rm(i)] Si $\mathfrak{J}$ est l'idéal gradué de $S_+$ engendré
    par $E$, on a $V_+(E)=V_+(\mathfrak{J})$ et l'assertion résulte alors de
    (\hyperref[II.2.3.9-fr]{2.3.9}).
  \item[\rm(ii)] Comme
    $V_+(\mathfrak{J})=\bigcap_{f\in\mathfrak{J}}V_+(f)$, la relation
    $Y\subset V_+(\mathfrak{J})$ entraîne $Y\subset V_+(f)$ pour tout
    $f\in\mathfrak{J}$, et par suite
    $\mathfrak{j}_+(Y)\supset\mathfrak{J}$, d'où
    $V_+(\mathfrak{j}_+(Y))\subset V_+(\mathfrak{J})$, ce qui prouve (ii)
    par définition des ensembles fermés.
\end{enumerate}
\end{proof}

\begin{corollary}[2.3.12]
\label{II.2.3.12-fr}
Les parties fermées $Y$ de $X=\operatorname{Proj}(S)$ et les idéaux gradués
de $S_+$ égaux à leur racine dans $S_+$ se correspondent biunivoquement par
les applications décroissantes
$Y\longmapsto\mathfrak{j}_+(Y)$, $\mathfrak{J}\longmapsto
V_+(\mathfrak{J})$~; à la réunion $Y_1\cup Y_2$ de deux parties fermées de
$X$ correspond
\oldpage[II]{28}
$\mathfrak{j}_+(Y_1)\cap\mathfrak{j}_+(Y_2)$, et à l'intersection d'une
famille quelconque $(Y_\lambda)$ de parties fermées correspond la racine
dans $S_+$ de la somme des $\mathfrak{j}_+(Y_\lambda)$.
\end{corollary}

\begin{corollary}[2.3.13]
\label{II.2.3.13-fr}
Soit $\mathfrak{J}$ un idéal gradué dans $S_+$~; pour que
$V_+(\mathfrak{J})=\varnothing$, il faut et il suffit que tout élément de
$S_+$ ait une puissance dans $\mathfrak{J}$.
\end{corollary}

Ce dernier corollaire s'exprime aussi sous l'une des formes équivalentes~:

\begin{corollary}[2.3.14]
\label{II.2.3.14-fr}
Soit $(f_\alpha)$ une famille d'éléments homogènes de $S_+$. Pour que les
$D_+(f_\alpha)$ forment un recouvrement de
$X=\operatorname{Proj}(S)$, il faut et il suffit que tout élément de $S_+$
ait une puissance dans l'idéal engendré par les $f_\alpha$.
\end{corollary}

\begin{corollary}[2.3.15]
\label{II.2.3.15-fr}
Soient $(f_\alpha)$ une famille d'éléments homogènes de $S_+$, $f$ un
élément de $S_+$. Les relations suivantes sont équivalentes~:
\begin{enumerate}
  \item[\rm a)] $D_+(f)\subset\bigcup_\alpha D_+(f_\alpha)$~;
  \item[\rm b)] $V_+(f)\supset\bigcap_\alpha V_+(f_\alpha)$~;
  \item[\rm c)] une puissance de $f$ appartient à l'idéal engendré par les
    $f_\alpha$.
\end{enumerate}
\end{corollary}

\begin{corollary}[2.3.16]
\label{II.2.3.16-fr}
Pour que $X=\operatorname{Proj}(S)$ soit vide, il faut et il suffit que tout
élément de $S_+$ soit nilpotent.
\end{corollary}

\begin{corollary}[2.3.17]
\label{II.2.3.17-fr}
Dans la correspondance biunivoque décrite dans
(\hyperref[II.2.3.12-fr]{2.3.12}), aux parties fermées irréductibles de $X$
correspondent les idéaux gradués premiers dans $S_+$.
\end{corollary}

\begin{proof}
En effet, si $Y=Y_1\cup Y_2$, où $Y_1$ et $Y_2$ sont fermés et distincts de
$Y$, on a
\[
  \mathfrak{j}_+(Y)
  =\mathfrak{j}_+(Y_1)\cap\mathfrak{j}_+(Y_2)
\]
les idéaux $\mathfrak{j}_+(Y_1)$ et $\mathfrak{j}_+(Y_2)$ étant distincts
de $\mathfrak{j}_+(Y)$, donc $\mathfrak{j}_+(Y)$ n'est pas premier.
Inversement, si $\mathfrak{J}$ est un idéal gradué non premier de $S_+$, il
existe deux éléments $f$, $g$ de $S_+$ tels que
$f\notin\mathfrak{J}$, $g\notin\mathfrak{J}$, $fg\in\mathfrak{J}$~; alors
$V_+(f)\not\supset V_+(\mathfrak{J})$,
$V_+(g)\not\supset V_+(\mathfrak{J})$ mais
$V_+(\mathfrak{J})\subset V_+(f)\cup V_+(g)$ par
(\hyperref[II.2.3.2.3-fr]{2.3.2.3})~; on en conclut que
$V_+(\mathfrak{J})$ est réunion des ensembles fermés
$V_+(f)\cap V_+(\mathfrak{J})$ et $V_+(g)\cap V_+(\mathfrak{J})$, qui sont
distincts de $V_+(\mathfrak{J})$.
\end{proof}

\subsection{La structure de schéma sur $\operatorname{Proj}(S)$.}
\label{subsection:II.2.4-fr}

\begin{env}[2.4.1]
\label{II.2.4.1-fr}
Soient $f$, $g$ deux éléments homogènes de $S_+$~; considérons les schémas
affines $Y_f=\operatorname{Spec}(S_{(f)})$,
$Y_g=\operatorname{Spec}(S_{(g)})$ et
$Y_{fg}=\operatorname{Spec}(S_{(fg)})$. En vertu de
(\hyperref[II.2.2.2-fr]{2.2.2}), le morphisme
$w_{fg,f}=({}^a\omega_{fg,f},\widetilde{\omega}_{fg,f})$ de $Y_{fg}$ dans
$Y_f$, correspondant à l'homomorphisme canonique
$\omega_{fg,f}:S_{(f)}\to S_{(fg)}$, est une \emph{immersion ouverte}
(\hyperref[I.1.3.6-fr]{I, 1.3.6}). Au moyen de l'homéomorphisme réciproque
de $\psi_f:D_+(f)\to Y_f$ (\hyperref[II.2.3.6-fr]{2.3.6}), on peut
transporter à $D_+(f)$ la structure de schéma affine de $Y_f$~; en vertu de
la commutativité du diagramme (\hyperref[II.2.3.5.1-fr]{2.3.5.1}), le
schéma affine $D_+(fg)$ s'identifie ainsi au schéma induit sur l'ensemble
ouvert $D_+(fg)$ de l'espace sous-jacent au schéma affine $D_+(f)$. Il est
clair alors (compte tenu de (\hyperref[II.2.3.4-fr]{2.3.4})) que
$X=\operatorname{Proj}(S)$ est muni d'une unique structure de
\emph{préschéma}, dont la restriction à chaque $D_+(f)$ est le schéma
affine qui vient d'être défini. En outre~:
\end{env}

\begin{proposition}[2.4.2]
\label{II.2.4.2-fr}
Le préschéma $\operatorname{Proj}(S)$ est un schéma.
\end{proposition}

\begin{proof}
Il suffit (\hyperref[I.5.5.6-fr]{I, 5.5.6}) de montrer que, quels que soient
$f$, $g$ homogènes dans $S_+$,
$D_+(f)\cap D_+(g)=D_+(fg)$ est affine et que son anneau est engendré par
les images canoniques des anneaux de $D_+(f)$ et $D_+(g)$~; le premier
point est évident par définition, et le second résulte de
(\hyperref[II.2.2.4-fr]{2.2.4}).
\end{proof}
\oldpage[II]{29}
Lorsqu'on parlera du spectre premier homogène $\operatorname{Proj}(S)$ comme
d'un \emph{schéma}, il s'agira toujours de la structure qui vient d'être
définie.

\begin{example}[2.4.3]
\label{II.2.4.3-fr}
Prenons $S=K[T_1,T_2]$, $K$ étant un corps, $T_1$, $T_2$ deux
indéterminées, et $S$ étant gradué par le degré total. Il résulte de
(\hyperref[II.2.3.14-fr]{2.3.14}) que $\operatorname{Proj}(S)$ est réunion
de $D_+(T_1)$ et $D_+(T_2)$~; on voit aussitôt que ces schémas affines sont
canoniquement isomorphes à $K[T]$, et que $\operatorname{Proj}(S)$
s'obtient par le recollement de ces deux schémas affines décrit dans
(\hyperref[I.2.3.2-fr]{I, 2.3.2}) (cf.
(\hyperref[II.7.4.14-fr]{7.4.14})).
\end{example}

\begin{proposition}[2.4.4]
\label{II.2.4.4-fr}
Soit $S$ un anneau gradué à degrés positifs, et soit $X$ le schéma
$\operatorname{Proj}(S)$.
\begin{enumerate}
  \item[\rm(i)] Si $\mathfrak{N}_+$ est le nilradical de $S_+$
    (\hyperref[II.2.1.10-fr]{2.1.10}), le schéma $X_{\mathrm{red}}$ est
    canoniquement isomorphe à $\operatorname{Proj}(S/\mathfrak{N}_+)$~; en
    particulier, si $S$ est essentiellement réduite,
    $\operatorname{Proj}(S)$ est réduit.
  \item[\rm(ii)] Supposons $S$ essentiellement réduit. Alors, pour que $X$
    soit intègre, il faut et il suffit que $S$ soit essentiellement intègre.
\end{enumerate}
\end{proposition}

\begin{proof}
\begin{enumerate}
  \item[\rm(i)] Soit $\overline{S}$ l'anneau gradué
    $S/\mathfrak{N}_+$, et désignons par $x\mapsto\overline{x}$
    l'homomorphisme canonique $S\to\overline{S}$, de degré $0$. Pour tout
    $f\in S_d$ ($d>0$), l'homomorphisme canonique
    $S_f\to\overline{S}_{\overline{f}}$
    (\hyperref[0.1.5.1-fr]{0, 1.5.1}) est surjectif et de degré $0$, donc
    par restriction donne un homomorphisme surjectif
    $S_{(f)}\to\overline{S}_{(\overline{f})}$~; si on suppose
    $f\notin\mathfrak{N}_+$, on vérifie immédiatement que
    $\overline{S}_{(\overline{f})}$ est réduit, et que le noyau de
    l'homomorphisme précédent est le nilradical de $S_{(f)}$, autrement dit
    $\overline{S}_{(\overline{f})}=(S_{(f)})_{\mathrm{red}}$. À cet
    homomorphisme correspond donc une immersion fermée
    $D_+(\overline{f})\to D_+(f)$ qui identifie $D_+(\overline{f})$ à
    $(D_+(f))_{\mathrm{red}}$ (\hyperref[I.5.1.2-fr]{I, 5.1.2}), et est en
    particulier un homéomorphisme des espaces sous-jacents à ces deux
    schémas affines. En outre, si $g\notin\mathfrak{N}_+$ est un second
    élément homogène de $S_+$, le diagramme
    \[
      \xymatrix{
        S_{(f)}\ar[r]\ar[d]
        & \overline{S}_{(\overline{f})}\ar[d]\\
        S_{(fg)}\ar[r]
        & \overline{S}_{(\overline{fg})}
      }
    \]
    est commutatif~; comme en outre les $D_+(f)$, pour $f$ homogène de
    degré $>0$ et $f\notin\mathfrak{N}_+$, forment un recouvrement de
    $X=\operatorname{Proj}(S)$ (\hyperref[II.2.3.7-fr]{2.3.7}), on voit
    que les morphismes $D_+(\overline{f})\to D_+(f)$ sont les restrictions
    d'une immersion fermée
    $\operatorname{Proj}(\overline{S})\to\operatorname{Proj}(S)$ qui est
    un homéomorphisme des espaces sous-jacents~; d'où la conclusion
    (\hyperref[I.5.1.2-fr]{I, 5.1.2}).
  \item[\rm(ii)] Supposons que $S$ soit essentiellement intègre, autrement
    dit que $(0)$ soit un idéal gradué premier de $S_+$ distinct de $S_+$~;
    alors $X$ est réduit en vertu de (i) et irréductible en vertu de
    (\hyperref[II.2.3.17-fr]{2.3.17}). Inversement, supposons que $S$ soit
    essentiellement réduit et $X$ intègre~; alors, pour $f\neq0$ homogène
    dans $S_+$, on a $D_+(f)\neq\varnothing$
    (\hyperref[II.2.3.7-fr]{2.3.7})~; l'hypothèse que $X$ est irréductible
    entraîne que $D_+(f)\cap D_+(g)\neq\varnothing$ pour $f$, $g$
    homogènes et $\neq0$ dans $S_+$~; donc $fg\neq0$ en vertu de
    (\hyperref[II.2.3.3.2-fr]{2.3.3.2}), et on en conclut que $S_+$ n'a pas
    de diviseurs de zéro, d'où la première assertion.
\end{enumerate}
\end{proof}

\begin{env}[2.4.5]
\label{II.2.4.5-fr}
Étant donné un anneau commutatif $A$, rappelons qu'on dit qu'un anneau
gradué $S$ est une \emph{$A$-algèbre graduée} s'il est muni d'une structure
de $A$-algèbre telle que chacun des sous-groupes $S_n$ soit un $A$-module~;
il suffit d'ailleurs pour cela que $S_0$ soit
\oldpage[II]{30}
une $A$-algèbre, autrement dit on définit sur $S$ une structure de
$A$-algèbre graduée en définissant une structure de $A$-algèbre sur $S_0$,
et en posant, pour $\alpha\in A$, $x\in S_n$,
$\alpha\cdot x=(\alpha\cdot1)x$.
\end{env}

\begin{proposition}[2.4.6]
\label{II.2.4.6-fr}
Supposons que $S$ soit une $A$-algèbre graduée. Alors sur
$X=\operatorname{Proj}(S)$, le faisceau structural $\mathcal{O}_X$ est une
$A$-Algèbre ($A$ étant considéré comme faisceau simple sur $X$)~; en
d'autres termes, $X$ est un schéma au-dessus de $\operatorname{Spec}(A)$.
\end{proposition}

\begin{proof}
Il suffit de noter que pour tout $f$ homogène dans $S_+$, $S_{(f)}$ est une
algèbre sur $A$, et que le diagramme
\[
  \xymatrix{
    S_{(f)}\ar[rr] && S_{(fg)}\\
    & A\ar[ul]\ar[ur] &
  }
\]
est commutatif, pour $f$, $g$ homogènes dans $S_+$.
\end{proof}

\begin{proposition}[2.4.7]
\label{II.2.4.7-fr}
Soit $S$ un anneau gradué à degrés positifs.
\begin{enumerate}
  \item[\rm(i)] Pour tout entier $d>0$, il existe un isomorphisme canonique
    du schéma $\operatorname{Proj}(S)$ sur le schéma
    $\operatorname{Proj}(S^{(d)})$.
  \item[\rm(ii)] Soit $S'$ l'anneau gradué tel que $S'_0=\mathbf{Z}$,
    $S'_n=S_n$ (considéré comme $\mathbf{Z}$-module) pour $n>0$. Il existe un
    isomorphisme canonique du schéma $\operatorname{Proj}(S)$ sur le schéma
    $\operatorname{Proj}(S')$.
\end{enumerate}
\end{proposition}

\begin{proof}
\begin{enumerate}
  \item[\rm(i)] Montrons d'abord que l'application
    $\mathfrak{p}\mapsto\mathfrak{p}\cap S^{(d)}$ est une bijection de
    l'ensemble $\operatorname{Proj}(S)$ sur
    $\operatorname{Proj}(S^{(d)})$. En effet, supposons donné un idéal
    premier gradué $\mathfrak{p}'\in\operatorname{Proj}(S^{(d)})$, et posons
    $\mathfrak{p}_{nd}=\mathfrak{p}'\cap S_{nd}$ ($n\geq0$). Pour tout
    $n>0$, non multiple de $d$, définissons $\mathfrak{p}_n$ comme
    l'ensemble des $x\in S_n$ tels que
    $x^d\in\mathfrak{p}_{nd}$~; si $x\in\mathfrak{p}_n$,
    $y\in\mathfrak{p}_n$, on a
    $(x+y)^{2d}\in\mathfrak{p}_{2nd}$, donc
    $(x+y)^d\in\mathfrak{p}_{nd}$ puisque $\mathfrak{p}'$ est premier~; il
    est immédiat que les $\mathfrak{p}_n$ ainsi définis pour tout $n\geq0$
    vérifient les conditions de
    (\hyperref[II.2.1.9-fr]{2.1.9}), donc il existe un idéal premier unique
    $\mathfrak{p}\in\operatorname{Proj}(S)$ tel que
    $\mathfrak{p}\cap S^{(d)}=\mathfrak{p}'$. Comme pour tout $f$ homogène
    dans $S_+$, on a
    $V_+(f)=V_+(f^d)$ (\hyperref[II.2.3.2.3-fr]{2.3.2.3}), on voit que la
    bijection précédente est un homéomorphisme d'espaces topologiques.
    Enfin, avec les mêmes notations, $S_{(f)}$ et
    $S^{(d)}_{(f^d)}$ s'identifient canoniquement
    (\hyperref[II.2.2.2-fr]{2.2.2}), donc $\operatorname{Proj}(S)$ et
    $\operatorname{Proj}(S^{(d)})$ s'identifient canoniquement en tant que
    schémas.
  \item[\rm(ii)] Si, à tout
    $\mathfrak{p}\in\operatorname{Proj}(S)$ on fait correspondre l'unique
    idéal premier $\mathfrak{p}'\in\operatorname{Proj}(S')$ tel que
    $\mathfrak{p}'\cap S_n=\mathfrak{p}\cap S_n$ pour tout $n>0$
    (\hyperref[II.2.1.9-fr]{2.1.9}), il est clair qu'on définit un
    homéomorphisme canonique
    $\operatorname{Proj}(S)\xrightarrow{\sim}\operatorname{Proj}(S')$ des
    espaces sous-jacents, car $V_+(f)$ est le même ensemble pour $S$ et
    $S'$ lorsque $f$ est un élément homogène de $S_+$. Comme en outre
    $S_{(f)}=S'_{(f)}$, $\operatorname{Proj}(S)$ et
    $\operatorname{Proj}(S')$ s'identifient en tant que schémas.
\end{enumerate}
\end{proof}

\begin{corollary}[2.4.8]
\label{II.2.4.8-fr}
Si $S$ est une $A$-algèbre graduée, $S'_A$ la $A$-algèbre graduée telle que
$(S'_A)_0=A$, $(S'_A)_n=S_n$ pour $n>0$, il existe un isomorphisme canonique
de $\operatorname{Proj}(S)$ sur $\operatorname{Proj}(S'_A)$.
\end{corollary}

\begin{proof}
En effet, ces deux schémas sont isomorphes canoniquement à
$\operatorname{Proj}(S')$, avec les notations de
(\hyperref[II.2.4.7-fr]{2.4.7, (ii)}).
\end{proof}

\subsection{Faisceau associé à un module gradué.}
\label{subsection:II.2.5-fr}

\begin{env}[2.5.1]
\label{II.2.5.1-fr}
Soit $M$ un $S$-module gradué. Pour tout $f$ homogène dans $S_+$,
$M_{(f)}$ est un $S_{(f)}$-module, et il lui correspond donc un faisceau
quasi-cohérent associé $\widetilde{M_{(f)}}$ sur le schéma affine
$\operatorname{Spec}(S_{(f)})$, identifié à $D_+(f)$
(\hyperref[I.1.3.4-fr]{I, 1.3.4}).
\end{env}
\oldpage[II]{31}

\begin{proposition}[2.5.2]
\label{II.2.5.2-fr}
Il existe sur $X=\operatorname{Proj}(S)$ un $\mathcal{O}_X$-Module
quasi-cohérent et un seul $\widetilde{M}$ tel que pour tout $f$ homogène
dans $S_+$, on ait
$\Gamma(D_+(f),\widetilde{M})=M_{(f)}$, l'homomorphisme de restriction
$\Gamma(D_+(f),\widetilde{M})\to
\Gamma(D_+(fg),\widetilde{M})$ pour $f$, $g$ homogènes dans $S_+$, étant
l'homomorphisme canonique $M_{(f)}\to M_{(fg)}$
(\hyperref[II.2.2.3-fr]{2.2.3}).
\end{proposition}

\begin{proof}
Supposons $f\in S_d$, $g\in S_e$. Comme $D_+(fg)$ s'identifie au spectre
premier de $(S_{(f)})_{g^d/f^e}$ en vertu de
(\hyperref[II.2.2.2-fr]{2.2.2}), la restriction à $D_+(fg)$ du faisceau
$\widetilde{M_{(f)}}$ sur $D_+(f)$ s'identifie canoniquement au faisceau
associé au module $(M_{(f)})_{g^d/f^e}$
(\hyperref[I.1.3.6-fr]{I, 1.3.6}), donc aussi à
$\widetilde{M_{(fg)}}$ (\hyperref[II.2.2.2-fr]{2.2.2})~; on en conclut
qu'il existe un isomorphisme canonique
\[
  \theta_{g,f}:\widetilde{M_{(f)}}|D_+(fg)
  \xrightarrow{\sim}\widetilde{M_{(g)}}|D_+(fg)
\]
tel que, si $h$ est un troisième élément homogène de $S_+$, on ait
$\theta_{f,h}=\theta_{f,g}\circ\theta_{g,h}$ dans $D_+(fgh)$. Par suite
(\hyperref[0.3.3.1-fr]{0, 3.3.1}) il existe sur $X$ un
$\mathcal{O}_X$-Module quasi-cohérent $\mathcal{F}$, et pour chaque $f$
homogène dans $S_+$, un isomorphisme $\eta_f$ de
$\mathcal{F}|D_+(f)$ sur $\widetilde{M_{(f)}}$ tels que
$\theta_{g,f}=\eta_g\circ\eta_f^{-1}$. Si alors on considère le faisceau
$\mathcal{G}$ associé au préfaisceau (sur la base de la topologie de $X$
formée des $D_+(f)$) défini par $D_+(f)\to M_{(f)}$, avec les
homomorphismes canoniques $M_{(f)}\to M_{(fg)}$ comme homomorphismes de
restriction, ce qui précède prouve que $\mathcal{F}$ et $\mathcal{G}$ sont
isomorphes (compte tenu de (\hyperref[I.1.3.7-fr]{I, 1.3.7}))~; le
faisceau $\mathcal{G}$ sera désigné par $\widetilde{M}$ et vérifie bien les
conditions de l'énoncé. On a en particulier
$\widetilde{S}=\mathcal{O}_X$.
\end{proof}

\begin{definition}[2.5.3]
\label{II.2.5.3-fr}
On dit que le $\mathcal{O}_X$-Module quasi-cohérent $\widetilde{M}$ défini
dans (\hyperref[II.2.5.2-fr]{2.5.2}) est \emph{associé} au $S$-module
gradué $M$.
\end{definition}

Rappelons que les $S$-modules gradués forment une catégorie lorsqu'on
restreint les homomorphismes de modules gradués aux homomorphismes
\emph{de degré $0$}. Avec cette convention~:

\begin{proposition}[2.5.4]
\label{II.2.5.4-fr}
Le foncteur $M\mapsto\widetilde{M}$ est un foncteur covariant additif exact,
de la catégorie des $S$-modules gradués dans celle des
$\mathcal{O}_X$-Modules quasi-cohérents, qui commute aux limites
inductives et aux sommes directes.
\end{proposition}

\begin{proof}
En effet, ces propriétés étant locales, il suffit de les vérifier sur les
faisceaux
$\widetilde{M}|D_+(f)=\widetilde{M_{(f)}}$~; or, le foncteur
$M\mapsto M_f$, le foncteur $N\mapsto N_0$ (dans la catégorie des
$S_f$-modules gradués) et le foncteur $P\mapsto\widetilde{P}$ (dans la
catégorie des $S_{(f)}$-modules) ont tous trois les propriétés d'exactitude
et de permutabilité aux limites inductives et aux sommes directes
(\hyperref[I.1.3.5-fr]{I, 1.3.5} et
\hyperref[I.1.3.9-fr]{I, 1.3.9})~; d'où la proposition.
\end{proof}

On notera $\widetilde{u}$ l'homomorphisme
$\widetilde{M}\to\widetilde{N}$ correspondant à un homomorphisme de degré
$0$, $u:M\to N$. On déduit aussitôt de
(\hyperref[II.2.5.4-fr]{2.5.4}) que les résultats de
(\hyperref[I.1.3.9-fr]{I, 1.3.9} et
\hyperref[I.1.3.10-fr]{I, 1.3.10}) sont encore valables pour les
$S$-modules gradués et les homomorphismes de degré $0$ (avec le sens donné
ici à $\widetilde{M}$), les démonstrations étant purement formelles.

\begin{proposition}[2.5.5]
\label{II.2.5.5-fr}
Pour tout $\mathfrak{p}\in X=\operatorname{Proj}(S)$, on a
$\widetilde{M}_{\mathfrak{p}}=M_{(\mathfrak{p})}$.
\end{proposition}

\begin{proof}
En effet, par définition
$\widetilde{M}_{\mathfrak{p}}=
\varinjlim\Gamma(D_+(f),\widetilde{M})$, où $f$ parcourt l'ensemble des
éléments homogènes de $S_+$ tels que $f\notin\mathfrak{p}$~; comme
$\Gamma(D_+(f),\widetilde{M})=M_{(f)}$, la proposition résulte de la
définition de $M_{(\mathfrak{p})}$
(\hyperref[II.2.2.7-fr]{2.2.7}).
\end{proof}
\oldpage[II]{32}

En particulier, \emph{l'anneau local $(\mathcal{O}_X)_\mathfrak{p}$} n'est
autre que l'anneau $S_{(\mathfrak{p})}$, ensemble des fractions $x/f$, où
$f$ est homogène dans $S_+$ et n'appartient pas à $\mathfrak{p}$, et $x$
homogène est de même degré que $f$.

Plus particulièrement, si $S$ est \emph{essentiellement intègre}, de sorte
que $\operatorname{Proj}(S)=X$ est \emph{intègre}
(\hyperref[II.2.4.4-fr]{2.4.4}), et si $\xi=(0)$ est le point générique de
$X$, le \emph{corps des fonctions rationnelles}
$R(X)=\mathcal{O}_\xi$, est le corps formé des éléments $fg^{-1}$, où $f$
et $g$ sont homogènes de même degré dans $S_+$ et $g\neq0$.

\begin{proposition}[2.5.6]
\label{II.2.5.6-fr}
Si, pour tout $z\in M$ et tout $f$ homogène dans $S_+$, il existe une
puissance de $f$ annulant $z$, on a $\widetilde{M}=0$. Cette condition
suffisante est aussi nécessaire lorsque la $S_0$-algèbre $S$ est engendrée
par l'ensemble $S_1$ des éléments homogènes de degré $1$.
\end{proposition}

\begin{proof}
En effet, la condition $\widetilde{M}=0$ équivaut à $M_{(f)}=0$ pour tout
$f$ homogène dans $S_+$. D'autre part, si $f\in S_d$, dire que $M_{(f)}=0$
signifie que pour tout $z\in M$ homogène et de degré multiple de $d$, il y
a une puissance $f^n$ telle que $f^nz=0$. Si $M_{(f)}=0$ pour tout
$f\in S_1$, la condition de l'énoncé est donc vérifiée pour tout
$f\in S_1$~; elle l'est \emph{a fortiori} pour tout $f$ homogène dans
$S_+$ lorsque $S_1$ engendre $S$, puisque tout élément homogène de $S_+$
est alors combinaison linéaire de produits d'éléments de $S_1$.
\end{proof}

\begin{proposition}[2.5.7]
\label{II.2.5.7-fr}
Soient $d$ un entier $>0$, $f\in S_d$. Alors, pour tout $n\in\mathbf{Z}$,
le $(\mathcal{O}_X|D_+(f))$-Module
$\widetilde{S(nd)}|D_+(f)$ est canoniquement isomorphe à
$\mathcal{O}_X|D_+(f)$.
\end{proposition}

\begin{proof}
En effet, la multiplication par l'élément inversible $(f/1)^n$ de $S_f$
est une bijection de $S_{(f)}=(S_f)_0$ sur
$(S_f)_{nd}=(S_f(nd))_0=(S(nd)_f)_0=S(nd)_{(f)}$~; autrement dit, les
$S_{(f)}$-modules $S_{(f)}$ et $S(nd)_{(f)}$ sont canoniquement
isomorphes, d'où la proposition.
\end{proof}

\begin{corollary}[2.5.8]
\label{II.2.5.8-fr}
Sur l'ensemble ouvert
\[
  U=\bigcup_{f\in S_d}D_+(f),
\]
la restriction du $\mathcal{O}_X$-Module $\widetilde{S(nd)}$ est un
$(\mathcal{O}_X|U)$-Module inversible
(\hyperref[0.5.4.1-fr]{0, 5.4.1}).
\end{corollary}

\begin{corollary}[2.5.9]
\label{II.2.5.9-fr}
Si l'idéal $S_+$ de $S$ est engendré par l'ensemble $S_1$ des éléments
homogènes de degré $1$, les $\mathcal{O}_X$-Modules
$\widetilde{S(n)}$ sont inversibles pour tout $n\in\mathbf{Z}$.
\end{corollary}

\begin{proof}
Il suffit de remarquer que
$X=\bigcup_{f\in S_1}D_+(f)$ en vertu de l'hypothèse
(\hyperref[II.2.3.14-fr]{2.3.14}) et d'appliquer
(\hyperref[II.2.5.8-fr]{2.5.8}) avec $U=X$.
\end{proof}

\begin{env}[2.5.10]
\label{II.2.5.10-fr}
Nous poserons, dans toute la suite de ce paragraphe
\[
\label{II.2.5.10.1-fr}
  \mathcal{O}_X(n)=\widetilde{S(n)}
\tag{2.5.10.1}
\]
pour tout $n\in\mathbf{Z}$, et pour toute partie ouverte $U$ de $X$ et tout
$(\mathcal{O}_X|U)$-Module $\mathcal{F}$
\[
\label{II.2.5.10.2-fr}
  \mathcal{F}(n)=
  \mathcal{F}\otimes_{\mathcal{O}_X|U}(\mathcal{O}_X(n)|U)
\tag{2.5.10.2}
\]
pour tout $n\in\mathbf{Z}$. Si l'idéal $S_+$ est engendré par $S_1$, le
foncteur $\mathcal{F}(n)$ est \emph{exact} en
$\mathcal{F}$ pour tout $n\in\mathbf{Z}$, car $\mathcal{O}_X(n)$ est alors
un $\mathcal{O}_X$-Module inversible.
\end{env}

\begin{env}[2.5.11]
\label{II.2.5.11-fr}
Soient $M$ et $N$ deux $S$-modules gradués. Pour tout $f\in S_d$ ($d>0$)
on définit un homomorphisme canonique fonctoriel de $S_{(f)}$-modules
\[
\label{II.2.5.11.1-fr}
  \lambda_f:M_{(f)}\otimes_{S_{(f)}}N_{(f)}
  \longrightarrow(M\otimes_S N)_{(f)}
\tag{2.5.11.1}
\]
\oldpage[II]{33}

en composant l'homomorphisme
$M_{(f)}\otimes_{S_{(f)}}N_{(f)}\to M_f\otimes_{S_f}N_f$
(provenant des injections $M_{(f)}\to M_f$, $N_{(f)}\to N_f$ et
$S_{(f)}\to S_f$) et l'isomorphisme canonique
$M_f\otimes_{S_f}N_f\xrightarrow{\sim}(M\otimes_S N)_f$
(\hyperref[0.1.3.4-fr]{0, 1.3.4}) et en notant que, par définition du
produit tensoriel de deux modules gradués, ce dernier isomorphisme conserve
les degrés~; pour $x\in M_{md}$, $y\in N_{nd}$ ($m\geq0$, $n\geq0$), on a
donc
\[
  \lambda_f((x/f^m)\otimes(y/f^n))=(x\otimes y)/f^{m+n}.
\]

Il résulte aussitôt de cette définition que si $g\in S_e$ ($e>0$), le
diagramme
\[
  \xymatrix{
    M_{(f)}\otimes_{S_{(f)}}N_{(f)}\ar[r]^{\lambda_f}\ar[d]
    & (M\otimes_S N)_{(f)}\ar[d]\\
    M_{(fg)}\otimes_{S_{(fg)}}N_{(fg)}\ar[r]_{\lambda_{fg}}
    & (M\otimes_S N)_{(fg)}
  }
\]
(où la flèche verticale de droite est l'homomorphisme canonique et celle de
gauche provient des homomorphismes canoniques) est commutatif. On déduit
donc des $\lambda$ un homomorphisme canonique fonctoriel de
$\mathcal{O}_X$-Modules
\[
\label{II.2.5.11.2-fr}
  \lambda:\widetilde{M}\otimes_{\mathcal{O}_X}\widetilde{N}
  \longrightarrow\widetilde{M\otimes_S N}.
\tag{2.5.11.2}
\]

Considérons en particulier deux idéaux gradués $\mathfrak{J}$,
$\mathfrak{K}$ de $S$~; comme $\widetilde{\mathfrak{J}}$ et
$\widetilde{\mathfrak{K}}$ sont des faisceaux d'idéaux de
$\mathcal{O}_X$, on a un homomorphisme canonique
$\widetilde{\mathfrak{J}}\otimes_{\mathcal{O}_X}
\widetilde{\mathfrak{K}}\to\mathcal{O}_X$~; le diagramme
\[
\label{II.2.5.11.3-fr}
  \xymatrix{
    \widetilde{\mathfrak{J}}\otimes_{\mathcal{O}_X}
    \widetilde{\mathfrak{K}}\ar[rr]^\lambda\ar[dr]
    && \widetilde{\mathfrak{J}\otimes_S\mathfrak{K}}\ar[dl]\\
    & \mathcal{O}_X &
  }
\tag{2.5.11.3}
\]
est alors commutatif. On se ramène en effet à le vérifier dans chaque
ouvert $D_+(f)$ ($f$ homogène dans $S_+$) et cela résulte aussitôt de la
définition (\hyperref[II.2.5.11.1-fr]{2.5.11.1}) de $\lambda_f$ et de
(\hyperref[I.1.3.13-fr]{I, 1.3.13}).

Notons enfin que si $M$, $N$, $P$ sont trois $S$-modules gradués, le
diagramme
\[
\label{II.2.5.11.4-fr}
  \xymatrix{
    \widetilde{M}\otimes_{\mathcal{O}_X}\widetilde{N}
    \otimes_{\mathcal{O}_X}\widetilde{P}
    \ar[r]^{\lambda\otimes1}\ar[d]_{1\otimes\lambda}
    & \widetilde{M\otimes_S N}\otimes_{\mathcal{O}_X}\widetilde{P}
    \ar[d]^\lambda\\
    \widetilde{M}\otimes_{\mathcal{O}_X}\widetilde{N\otimes_S P}
    \ar[r]_\lambda
    & \widetilde{M\otimes_S N\otimes_S P}
  }
\tag{2.5.11.4}
\]
\oldpage[II]{34}

est commutatif. Il suffit encore de le vérifier dans chaque $D_+(f)$ et
cela découle aussitôt des définitions et de
(\hyperref[I.1.3.13-fr]{I, 1.3.13}).
\end{env}

\begin{env}[2.5.12]
\label{II.2.5.12-fr}
Sous les hypothèses de (\hyperref[II.2.5.11-fr]{2.5.11}), on définit un
homomorphisme canonique fonctoriel de $S_{(f)}$-modules
\[
\label{II.2.5.12.1-fr}
  \mu_f:\bigl(\operatorname{Hom}_S(M,N)\bigr)_{(f)}
  \longrightarrow
  \operatorname{Hom}_{S_{(f)}}(M_{(f)},N_{(f)})
\tag{2.5.12.1}
\]
en faisant correspondre à $u/f^n$, où $u$ est un homomorphisme de degré
$nd$, l'homomorphisme $M_{(f)}\to N_{(f)}$ qui transforme $x/f^m$
($x\in M_{md}$) en $u(x)/f^{m+n}$. Pour $g\in S_e$ ($e>0$), on a encore
un diagramme commutatif
\[
  \xymatrix{
    \bigl(\operatorname{Hom}_S(M,N)\bigr)_{(f)}
    \ar[r]^{\mu_f}\ar[d]
    & \operatorname{Hom}_{S_{(f)}}(M_{(f)},N_{(f)})\ar[d]\\
    \bigl(\operatorname{Hom}_S(M,N)\bigr)_{(fg)}
    \ar[r]_{\mu_{fg}}
    & \operatorname{Hom}_{S_{(fg)}}(M_{(fg)},N_{(fg)})
  }
\]
(la flèche de gauche étant l'homomorphisme canonique, celle de droite
provenant des homomorphismes canoniques). On en conclut encore (compte
tenu de (\hyperref[I.1.3.8-fr]{I, 1.3.8})) que les $\mu_f$ définissent un
homomorphisme canonique fonctoriel de $\mathcal{O}_X$-Modules
\[
\label{II.2.5.12.2-fr}
  \mu:\widetilde{\operatorname{Hom}_S(M,N)}
  \longrightarrow
  \mathcal{H}om_{\mathcal{O}_X}(\widetilde{M},\widetilde{N}).
\tag{2.5.12.2}
\]
\end{env}

\begin{proposition}[2.5.13]
\label{II.2.5.13-fr}
Supposons l'idéal $S_+$ engendré par $S_1$. Alors l'homomorphisme
$\lambda$ (\hyperref[II.2.5.11.2-fr]{2.5.11.2}) est un isomorphisme~; il
en est de même de l'homomorphisme $\mu$
(\hyperref[II.2.5.12.2-fr]{2.5.12.2}) lorsque le $S$-module gradué $M$
admet une présentation finie (\hyperref[II.2.1.1-fr]{2.1.1}).
\end{proposition}

\begin{proof}
Comme $X$ est réunion des $D_+(f)$ pour $f\in S_1$
(\hyperref[II.2.3.14-fr]{2.3.14}), on est ramené à prouver que $\lambda_f$
et $\mu_f$ sont des isomorphismes, sous les hypothèses envisagées, lorsque
$f$ est homogène et \emph{de degré $1$}. Or, on définit alors une
application $\mathbf{Z}$-bilinéaire
\[
  M_m\times N_n\longrightarrow
  M_{(f)}\otimes_{S_{(f)}}N_{(f)}
\]
en faisant correspondre à $(x,y)$ l'élément
$(x/f^m)\otimes(y/f^n)$ (lorsque $m<0$, nous écrivons $x/f^m$ pour
$f^{-m}x/1$)~; ces applications définissent une application
$\mathbf{Z}$-linéaire
\[
  M\otimes_{\mathbf{Z}}N\longrightarrow
  M_{(f)}\otimes_{S_{(f)}}N_{(f)},
\]
et si $s\in S_q$, cette application transforme $(sx)\otimes y$ en
\[
  (s/f^q)\bigl((x/f^m)\otimes(y/f^n)\bigr)
\]
(pour $x\in M_m$, $y\in N_n$). On en déduit donc un di-homomorphisme de
modules
\[
  \gamma_f:M\otimes_S N\longrightarrow
  M_{(f)}\otimes_{S_{(f)}}N_{(f)},
\]
relatif à l'homomorphisme canonique $S\to S_{(f)}$ (transformant
$s\in S_q$ en $s/f^q$). Supposons en outre que pour un élément
$\sum_i(x_i\otimes y_i)$ de $M\otimes_S N$ ($x_i$, $y_i$ homogènes de
degrés respectifs $m_i$, $n_i$) on ait
\[
  f^r\sum_i(x_i\otimes y_i)=0,
\]
autrement dit $\sum_i(f^rx_i\otimes y_i)=0$. On en déduit en vertu de
(\hyperref[0.1.3.4-fr]{0, 1.3.4}) que l'on a
\[
  \sum_i(f^rx_i/f^{m_i+r})\otimes(y_i/f^{n_i})=0,
\]
c'est-à-dire $\gamma_f(\sum_i(x_i\otimes y_i))=0$. Par suite $\gamma_f$
se factorise en
\[
  M\otimes_S N\longrightarrow(M\otimes_S N)_f
  \xrightarrow{\gamma'_f}
  M_{(f)}\otimes_{S_{(f)}}N_{(f)}~;
\]
si $\lambda'_f$ est la restriction de $\gamma'_f$
\oldpage[II]{35}
à $(M\otimes_S N)_{(f)}$, on vérifie aussitôt que $\lambda_f$ et
$\lambda'_f$ sont des $S_{(f)}$-homomorphismes réciproques, d'où la
première partie de la proposition.

Pour démontrer la seconde partie, supposons que $M$ soit le conoyau d'un
homomorphisme $P\to Q$, de $S$-modules gradués, $P$ et $Q$ étant sommes
directes d'un nombre fini de modules de la forme $S(n)$~; utilisant
l'exactitude à gauche de $\operatorname{Hom}_S(L,N)$ en $L$ et l'exactitude
de $M_{(f)}$ en $M$, on est aussitôt ramené à prouver que $\mu_f$ est un
isomorphisme lorsque $M=S(n)$. Or, pour tout $z$ homogène dans $N$, soit
$u_z$ l'homomorphisme de $S(n)$ dans $N$ tel que $u_z(1)=z$~; on voit
aussitôt que $\eta:z\mapsto u_z$ est un isomorphisme de degré $0$ de
$N(-n)$ sur $\operatorname{Hom}_S(S(n),N)$. Il lui correspond un
isomorphisme
\[
  \eta_f:\bigl(N(-n)\bigr)_{(f)}
  \longrightarrow
  \bigl(\operatorname{Hom}_S(S(n),N)\bigr)_{(f)}.
\]
D'autre part, soit $\eta'_f$ l'isomorphisme
\[
  N_{(f)}\longrightarrow
  \operatorname{Hom}_{S_{(f)}}\bigl(S(n)_{(f)},N_{(f)}\bigr)
\]
qui, à tout $z'\in N_{(f)}$, fait correspondre l'homomorphisme $v_{z'}$ tel
que $v_{z'}(s/f^k)=sz'/f^{n+k}$ (pour
$s\in S_{n+k}=(S(n))_k$). On constate aisément que l'application composée
\[
  \bigl(N(-n)\bigr)_{(f)}
  \xrightarrow{\eta_f}
  \bigl(\operatorname{Hom}_S(S(n),N)\bigr)_{(f)}
  \xrightarrow{\mu_f}
  \operatorname{Hom}_{S_{(f)}}\bigl(S(n)_{(f)},N_{(f)}\bigr)
  \xrightarrow{{\eta'_f}^{-1}}N_{(f)}
\]
est l'isomorphisme $z/f^h\mapsto z/f^{h-n}$ de
$\bigl(N(-n)\bigr)_{(f)}$ sur $N_{(f)}$, donc $\mu_f$ est un isomorphisme.
\end{proof}

Si l'idéal $S_+$ est engendré par $S_1$, on déduit de
(\hyperref[II.2.5.13-fr]{2.5.13}) que pour tout idéal gradué $\mathfrak{J}$
de $S$ et tout $S$-module gradué $M$, on a
\[
\label{II.2.5.13.1-fr}
  \widetilde{\mathfrak{J}}\cdot\widetilde{M}
  =\widetilde{\mathfrak{J}\cdot M}
\tag{2.5.13.1}
\]
à un isomorphisme canonique près~; cela résulte en effet de la commutativité
du diagramme
\[
  \xymatrix{
    \widetilde{\mathfrak{J}}\otimes_{\mathcal{O}_X}\widetilde{M}
    \ar[rr]^\lambda\ar[dr]
    && \widetilde{\mathfrak{J}\otimes_S M}\ar[dl]\\
    & \widetilde{M} &
  }
\]
que l'on vérifie comme celle de
(\hyperref[II.2.5.11.3-fr]{2.5.11.3}).

\begin{corollary}[2.5.14]
\label{II.2.5.14-fr}
Supposons $S$ engendrée par $S_1$. Quels que soient $m$, $n$ dans
$\mathbf{Z}$, on a alors~:
\[
\label{II.2.5.14.1-fr}
  \mathcal{O}_X(m)\otimes_{\mathcal{O}_X}\mathcal{O}_X(n)
  =\mathcal{O}_X(m+n)
\tag{2.5.14.1}
\]
\[
\label{II.2.5.14.2-fr}
  \mathcal{O}_X(n)=\bigl(\mathcal{O}_X(1)\bigr)^{\otimes n}
\tag{2.5.14.2}
\]
à des isomorphismes canoniques près.
\end{corollary}

\begin{proof}
La première formule résulte de (\hyperref[II.2.5.13-fr]{2.5.13}) et de
l'existence de l'isomorphisme canonique de degré $0$,
$S(m)\otimes_S S(n)\xrightarrow{\sim}S(m+n)$, qui, à l'élément
$1\otimes 1$ (où le premier facteur $1$ est dans $(S(m))_{-m}$ et le second
dans $(S(n))_{-n}$) fait correspondre l'élément
$1\in(S(m+n))_{-m-n}$. Il suffit ensuite de démontrer la seconde formule
pour $n=-1$, et en vertu de (\hyperref[II.2.5.13-fr]{2.5.13}), cela revient
à voir que $\operatorname{Hom}_S(S(1),S)$ est canoniquement isomorphe à
$S(-1)$, vérification immédiate en remontant aux définitions
(\hyperref[II.2.1.2-fr]{2.1.2}) et en se souvenant que $S(1)$ est un
$S$-module monogène.
\end{proof}

\oldpage[II]{36}
\begin{corollary}[2.5.15]
\label{II.2.5.15-fr}
Supposons $S$ engendrée par $S_1$. Pour tout $S$-module gradué $M$ et
tout $n\in\mathbf{Z}$, on a
\[
\label{II.2.5.15.1-fr}
  \widetilde{M(n)}=\widetilde{M}(n)
\tag{2.5.15.1}
\]
à un isomorphisme canonique près.
\end{corollary}

\begin{proof}
Cela résulte des définitions (\hyperref[II.2.5.10.2-fr]{2.5.10.2}) et
(\hyperref[II.2.5.10.1-fr]{2.5.10.1}), de la prop.
(\hyperref[II.2.5.13-fr]{2.5.13}) et de l'existence d'un isomorphisme
canonique $M(n)\xrightarrow{\sim}M\otimes_S S(n)$ de degré $0$ qui, à tout
$z\in(M(n))_h=M_{n+h}$, fait correspondre
$z\otimes1\in M_{n+h}\otimes(S(n))_{-n}\subset(M\otimes_S S(n))_h$.
\end{proof}

\begin{env}[2.5.16]
\label{II.2.5.16-fr}
Désignons par $S'$ l'anneau gradué tel que $S'_0=\mathbf{Z}$,
$S'_n=S_n$ pour $n>0$. Alors, si $f\in S_d$ ($d>0$), on a
$(S(n))_{(f)}=(S'(n))_{(f)}$ pour tout $n\in\mathbf{Z}$, car un élément de
$(S'(n))_{(f)}$ est de la forme $x/f^k$ avec $x\in S'_{n+kd}$ ($k>0$), et
on peut toujours prendre $k$ tel que $n+kd\neq0$. Comme
$X=\operatorname{Proj}(S)$ et $X'=\operatorname{Proj}(S')$ s'identifient
canoniquement (\hyperref[II.2.4.7-fr]{2.4.7, (ii)}), on voit que pour tout
$n\in\mathbf{Z}$, $\mathcal{O}_X(n)$ et $\mathcal{O}_{X'}(n)$ sont images
l'un de l'autre par l'identification précédente.

Notons d'autre part que pour tout $d>0$ et tout $n\in\mathbf{Z}$, on a
\[
  (S^{(d)}(n))_h=S_{(n+h)d}=(S(nd))_{hd}
\]
pour $f\in S_d$, on a donc $(S^{(d)}(n))_{(f)}=(S(nd))_{(f)}$. On sait que
les schémas $X=\operatorname{Proj}(S)$ et
$X^{(d)}=\operatorname{Proj}(S^{(d)})$ s'identifient canoniquement
(\hyperref[II.2.4.7-fr]{2.4.7, (i)})~; ce qui précède montre que si la
$S_0$-algèbre $S^{(d)}$ est engendrée par $S_d$, $\mathcal{O}_X(nd)$ et
$\mathcal{O}_{X^{(d)}}(n)$ sont images l'un de l'autre par cette
identification, pour tout $n\in\mathbf{Z}$.
\end{env}

\begin{proposition}[2.5.17]
\label{II.2.5.17-fr}
Soit $d$ un entier $>0$, et soit
$U=\bigcup_{f\in S_d}D_+(f)$. La restriction à $U$ de l'homomorphisme
canonique
$\mathcal{O}_X(nd)\otimes_{\mathcal{O}_X}\mathcal{O}_X(-nd)
\to\mathcal{O}_X$ est un isomorphisme pour tout entier $n$.
\end{proposition}

\begin{proof}
En vertu de (\hyperref[II.2.5.16-fr]{2.5.16}), on peut se borner au cas où
$d=1$, et la conclusion résulte alors de la démonstration de
(\hyperref[II.2.5.13-fr]{2.5.13}).
\end{proof}

\subsection{$S$-module gradué associé à un faisceau sur
$\operatorname{Proj}(S)$.}
\label{subsection:II.2.6-fr}

\emph{On suppose dans tout ce numéro que l'idéal $S_+$ de $S$ est engendré
par l'ensemble $S_1$ des éléments homogènes de degré $1$.}

\begin{env}[2.6.1]
\label{II.2.6.1-fr}
Le $\mathcal{O}_X$-Module $\mathcal{O}_X(1)$ est alors
\emph{inversible} (\hyperref[II.2.5.9-fr]{2.5.9})~; on pose alors, pour tout
$\mathcal{O}_X$-Module $\mathcal{F}$
(\hyperref[0.5.4.6-fr]{0, 5.4.6}),
\[
\label{II.2.6.1.1-fr}
  \Gamma_*(\mathcal{F})
  =\Gamma_*(\mathcal{O}_X(1),\mathcal{F})
  =\bigoplus_{n\in\mathbf{Z}}\Gamma(X,\mathcal{F}(n))
\tag{2.6.1.1}
\]
compte tenu de (\hyperref[II.2.5.14.2-fr]{2.5.14.2}). Rappelons
(\hyperref[0.5.4.6-fr]{0, 5.4.6}) que $\Gamma_*(\mathcal{O}_X)$ est muni
d'une structure d'\emph{anneau gradué}, et $\Gamma_*(\mathcal{F})$ d'une
structure de
\emph{module gradué} sur $\Gamma_*(\mathcal{O}_X)$.

Puisque $\mathcal{O}_X(n)$ est localement libre, $\Gamma_*(\mathcal{F})$
est un foncteur covariant additif
\emph{exact à gauche} en $\mathcal{F}$~; en particulier, si $\mathcal{J}$
est un faisceau d'idéaux de $\mathcal{O}_X$, $\Gamma_*(\mathcal{J})$
s'identifie canoniquement à un
\emph{idéal gradué} de $\Gamma_*(\mathcal{O}_X)$.
\end{env}

\begin{env}[2.6.2]
\label{II.2.6.2-fr}
Soit $M$ un $S$-module gradué~; pour tout $f\in S_d$ ($d>0$),
$x\mapsto x/1$ est un homomorphisme de groupes abéliens
$M_0\to M_{(f)}$, et comme $M_{(f)}$ s'identifie canoni-

\oldpage[II]{37}

quement à $\Gamma(D_+(f),\widetilde{M})$, on obtient ainsi un
homomorphisme de groupes abéliens
\[
  \alpha_0^f:M_0\longrightarrow\Gamma(D_+(f),\widetilde{M}).
\]
Il est clair que, pour tout $g\in S_e$ ($e>0$), le diagramme
\[
  \xymatrix{
    & \Gamma(D_+(f),\widetilde{M})\ar[dd]\\
    M_0\ar[ur]^{\alpha_0^f}\ar[dr]_{\alpha_0^{fg}}\\
    & \Gamma(D_+(fg),\widetilde{M})
  }
\]
est commutatif~; cela signifie que pour tout $x\in M_0$, les sections
$\alpha_0^f(x)$ et $\alpha_0^g(x)$ de $\widetilde{M}$ coïncident dans
$D_+(f)\cap D_+(g)$, et par suite il existe une section unique
$\alpha_0(x)\in\Gamma(X,\widetilde{M})$ dont la restriction à chaque
$D_+(f)$ est $\alpha_0^f(x)$. On a ainsi défini (sans utiliser l'hypothèse
que $S$ est engendrée par $S_1$) un homomorphisme de groupes abéliens
\[
\label{II.2.6.2.1-fr}
  \alpha_0:M_0\longrightarrow\Gamma(X,\widetilde{M}).
\tag{2.6.2.1}
\]

Appliquant ce résultat au $S$-module gradué $M(n)$ (pour tout
$n\in\mathbf{Z}$), on obtient pour chaque $n\in\mathbf{Z}$ un homomorphisme
de groupes abéliens
\[
\label{II.2.6.2.2-fr}
  \alpha_n:M_n=(M(n))_0\longrightarrow\Gamma(X,\widetilde{M}(n))
\tag{2.6.2.2}
\]
(compte tenu de (\hyperref[II.2.5.15-fr]{2.5.15}))~; d'où un homomorphisme
fonctoriel (de degré $0$) de groupes abéliens gradués
\[
\label{II.2.6.2.3-fr}
  \alpha:M\longrightarrow\Gamma_*(\widetilde{M})
\tag{2.6.2.3}
\]
(aussi noté $\alpha_M$) qui, dans chaque $M_n$, coïncide avec $\alpha_n$.

Si on prend en particulier $M=S$, on vérifie aussitôt (compte tenu de la
définition de la multiplication dans $\Gamma_*(\mathcal{O}_X)$
(\hyperref[0.5.4.6-fr]{0, 5.4.6})) que
$\alpha:S\to\Gamma_*(\mathcal{O}_X)$ est un homomorphisme d'anneaux
gradués, et que pour tout $S$-module gradué $M$,
(\hyperref[II.2.6.2.3-fr]{2.6.2.3}) est un di-homomorphisme de modules
gradués.
\end{env}

\begin{proposition}[2.6.3]
\label{II.2.6.3-fr}
Pour tout $f\in S_d$ ($d>0$), $D_+(f)$ est identique à l'ensemble des
$\mathfrak{p}\in X$ où la section $\alpha_d(f)$ de $\mathcal{O}_X(d)$ ne
s'annule pas (\hyperref[0.5.5.2-fr]{0, 5.5.2}).
\end{proposition}

\begin{proof}
Comme $X=\bigcup_{g\in S_1}D_+(g)$ par hypothèse, il suffit de vérifier que
pour tout $g\in S_1$, l'ensemble des $\mathfrak{p}\in D_+(g)$ où
$\alpha_d(f)$ ne s'annule pas est identique à $D_+(fg)$. Or, la restriction
de $\alpha_d(f)$ à $D_+(g)$ est par définition la section correspondant à
l'élément $f/1$ de $(S(d))_{(g)}$~; par l'isomorphisme canonique
$(S(d))_{(g)}\xrightarrow{\sim}S_{(g)}$
(\hyperref[II.2.5.7-fr]{2.5.7}), cette section de $\mathcal{O}_X(d)$
au-dessus de $D_+(g)$ s'identifie à la section de $\mathcal{O}_X$ au-dessus
de $D_+(g)$ qui correspond à l'élément $f/g^d$ de $S_{(g)}$~; dire que cette
section s'annule en $\mathfrak{p}\in D_+(g)$ signifie que
$f/g^d\in\mathfrak{q}$, où $\mathfrak{q}$ est l'idéal premier de $S_{(g)}$
correspondant à $\mathfrak{p}$ (\hyperref[II.2.3.6-fr]{2.3.6})~; par
définition cela veut dire que $f\in\mathfrak{p}$, d'où la proposition.
\end{proof}

\begin{env}[2.6.4]
\label{II.2.6.4-fr}
Soit maintenant $\mathcal{F}$ un $\mathcal{O}_X$ module, et posons
$M=\Gamma_*(\mathcal{F})$~; en vertu de l'existence de l'homomorphisme
d'anneaux gradués $\alpha:S\to\Gamma_*(\mathcal{O}_X)$, $M$ peut être
considéré comme un $S$-module gradué. Pour tout $f\in S_d$ ($d>0$), il
résulte de (\hyperref[II.2.6.3-fr]{2.6.3}) que la restriction à $D_+(f)$ de
la section $\alpha_d(f)$ de $\mathcal{O}_X(d)$ est inversible~; il en est
donc de même de la restriction à $D_+(f)$ de la section $\alpha_d(f^n)$ de
$\mathcal{O}_X(nd)$, pour tout $n>0$. Soit alors
$z\in M_{nd}=\Gamma(X,\mathcal{F}(nd))$ ($n>0$)~; s'il existe un entier
$k\geq0$ tel que la restriction à $D_+(f)$ de $f^kz$, c'est-à-dire la
\oldpage[II]{38}
section
$(z|D_+(f))(\alpha_d(f^k)|D_+(f))$ de $\mathcal{F}((n+k)d)$, soit nulle,
alors, en raison de la remarque précédente, on a aussi $z|D_+(f)=0$. Cela
montre que l'on définit un $S_{(f)}$-homomorphisme
$\beta_f:M_{(f)}\to\Gamma(D_+(f),\mathcal{F})$ en faisant correspondre à
l'élément $z/f^n$ la section
$(z|D_+(f))(\alpha_d(f^n)|D_+(f))^{-1}$ de $\mathcal{F}$ au-dessus de
$D_+(f)$. On vérifie en outre aussitôt que le diagramme
\[
\label{II.2.6.4.1-fr}
  \xymatrix{
    M_{(f)}\ar[r]^{\beta_f}\ar[d]
    & \Gamma(D_+(f),\mathcal{F})\ar[d]\\
    M_{(fg)}\ar[r]_{\beta_{fg}}
    & \Gamma(D_+(fg),\mathcal{F})
  }
\tag{2.6.4.1}
\]
est commutatif pour $g\in S_e$ ($e>0$). Si on se rappelle que $M_{(f)}$
s'identifie canoniquement à $\Gamma(D_+(f),\widetilde{M})$ et que les
$D_+(f)$ forment une base de la topologie de $X$
(\hyperref[II.2.3.4-fr]{2.3.4}), on voit que les $\beta_f$ proviennent d'un
unique homomorphisme canonique de $\mathcal{O}_X$-Modules
\[
\label{II.2.6.4.2-fr}
  \beta:\widetilde{\Gamma_*(\mathcal{F})}
  \longrightarrow\mathcal{F}
\tag{2.6.4.2}
\]
(aussi noté $\beta_{\mathcal{F}}$) qui est évidemment fonctoriel.
\end{env}

\begin{proposition}[2.6.5]
\label{II.2.6.5-fr}
Soient $M$ un $S$-module gradué, $\mathcal{F}$ un $\mathcal{O}_X$-Module~;
alors les homomorphismes composés
\[
\label{II.2.6.5.1-fr}
  \widetilde{M}\xrightarrow{\widetilde{\alpha}}
  \widetilde{\Gamma_*(\widetilde{M})}
  \xrightarrow{\beta}\widetilde{M}
\tag{2.6.5.1}
\]
\[
\label{II.2.6.5.2-fr}
  \Gamma_*(\mathcal{F})\xrightarrow{\alpha}
  \Gamma_*(\widetilde{\Gamma_*(\mathcal{F})})
  \xrightarrow{\Gamma_*(\beta)}\Gamma_*(\mathcal{F})
\tag{2.6.5.2}
\]
sont les isomorphismes identiques.
\end{proposition}

\begin{proof}
La vérification de (\hyperref[II.2.6.5.1-fr]{2.6.5.1}) est locale~: dans un
ouvert $D_+(f)$, elle résulte aussitôt des définitions et de ce que $\beta$,
appliqué à des faisceaux quasi-cohérents, est déterminé par son action sur les
sections au-dessus de $D_+(f)$ (\hyperref[I.1.3.8-fr]{I, 1.3.8}). La
vérification de (\hyperref[II.2.6.5.2-fr]{2.6.5.2}) se fait pour chaque degré
séparément~: si on pose $M=\Gamma_*(\mathcal{F})$, on a
$M_n=\Gamma(X,\mathcal{F}(n))$ et
$(\Gamma_*(\widetilde{M}))_n=\Gamma(X,\widetilde{M}(n))
=\Gamma(X,\widetilde{M(n)})$. Or, si $f\in S_1$ et $z\in M_n$,
$\alpha_n^f(z)$ est l'élément $z/1$ de $(M(n))_{(f)}$, égal à
$(f/1)^n(z/f^n)$~; il lui correspond par $\beta_f$ la section
\[
  ((\alpha_1(f))^n|D_+(f))
  ((z|D_+(f))((\alpha_1(f))^n|D_+(f))^{-1})
\]
au-dessus de $D_+(f)$, c'est-à-dire la restriction de $z$ à $D_+(f)$, ce qui
vérifie (\hyperref[II.2.6.5.2-fr]{2.6.5.2}).
\end{proof}

\subsection{Conditions de finitude.}
\label{subsection:II.2.7-fr}

\begin{proposition}[2.7.1]
\label{II.2.7.1-fr}
\begin{enumerate}
  \item[\rm(i)] Si $S$ est un anneau gradué noethérien,
    $X=\operatorname{Proj}(S)$ est un schéma noethérien.
  \item[\rm(ii)] Si $S$ est une $A$-algèbre graduée de type fini,
    $X=\operatorname{Proj}(S)$ est un schéma de type fini au-dessus de
    $Y=\operatorname{Spec}(A)$.
\end{enumerate}
\end{proposition}

\oldpage[II]{39}
\begin{proof}
\begin{enumerate}
  \item[\rm(i)] Si $S$ est noethérien, l'idéal $S_+$ admet un système fini de
    générateurs homogènes $f_i$ ($1\leq i\leq p$), donc
    (\hyperref[II.2.3.14-fr]{2.3.14}) l'espace sous-jacent $X$ est réunion des
    $D_+(f_i)=\operatorname{Spec}(S_{(f_i)})$, et tout revient à voir que
    chacun des $S_{(f_i)}$ est noethérien, ce qui résulte de
    (\hyperref[II.2.2.6-fr]{2.2.6}).
  \item[\rm(ii)] L'hypothèse entraîne que $S_0$ est une $A$-algèbre de type
    fini et $S$ une $S_0$-algèbre de type fini, donc $S_+$ est un idéal de
    type fini (\hyperref[II.2.1.4-fr]{2.1.4}). On est alors ramené, comme dans
    (i), à prouver que pour $f\in S_d$, $S_{(f)}$ est une $A$-algèbre de type
    fini. En vertu de (\hyperref[II.2.2.5-fr]{2.2.5}), il suffit de montrer
    que $S^{(d)}$ est une $A$-algèbre de type fini, ce qui résulte de
    (\hyperref[II.2.1.6-fr]{2.1.6}).
\end{enumerate}
\end{proof}

\begin{env}[2.7.2]
\label{II.2.7.2-fr}
Nous considérerons dans ce qui suit les conditions de finitude suivantes pour
un $S$-module gradué $M$~:
\begin{enumerate}
  \item[\rm(TF)] Il existe un entier $n$ tel que le sous-module
    $\bigoplus_{k\geq n}M_k$ soit un $S$-module de type fini.
  \item[\rm(TN)] Il existe un entier $n$ tel que $M_k=0$ pour $k\geq n$.
\end{enumerate}

Si $M$ vérifie (TN), on a $M_{(f)}=0$ pour tout $f$ homogène dans $S_+$,
donc $\widetilde{M}=0$.

Soient $M$, $N$ deux $S$-modules gradués~; on dit qu'un homomorphisme
$u:M\to N$ de degré $0$ est (TN)-\emph{injectif} (resp.
(TN)-\emph{surjectif}, (TN)-\emph{bijectif}) s'il existe un entier $n$ tel
que $u_k:M_k\to N_k$ soit injectif (resp. surjectif, bijectif) pour
$k\geq n$. Dire que $u$ est (TN)-injectif (resp. (TN)-surjectif) revient donc
à dire que $\operatorname{Ker}u$ (resp. $\operatorname{Coker}u$) vérifie
(TN). En vertu de (\hyperref[II.2.5.4-fr]{2.5.4}), si $u$ est (TN)-injectif
(resp. (TN)-surjectif, (TN)-bijectif), $\widetilde{u}$ est injectif (resp.
surjectif, bijectif)~; lorsque $u$ est (TN)-bijectif, on dit encore que $u$
est un (TN)-\emph{isomorphisme}.
\end{env}

\begin{proposition}[2.7.3]
\label{II.2.7.3-fr}
Soient $S$ un anneau gradué tel que l'idéal $S_+$ soit de type fini, $M$ un
$S$-module gradué.
\begin{enumerate}
  \item[\rm(i)] Si $M$ vérifie la condition (TF), le $\mathcal{O}_X$-Module
    $\widetilde{M}$ est de type fini.
  \item[\rm(ii)] Supposons que $M$ vérifie (TF)~; pour que $\widetilde{M}=0$,
    il faut et il suffit que $M$ vérifie (TN).
\end{enumerate}
\end{proposition}

\begin{proof}
On vient de voir que la condition (TN) implique $\widetilde{M}=0$. Si $M$
vérifie (TF), le sous-module gradué $M'=\bigoplus_{k\geq n}M_k$ qui est par
hypothèse de type fini, est tel que $M/M'$ satisfasse à (TN)~; on a par suite
$\widetilde{M/M'}=0$, et l'exactitude du foncteur
$M\mapsto\widetilde{M}$ (\hyperref[II.2.5.4-fr]{2.5.4}) entraîne
$\widetilde{M}=\widetilde{M'}$~; pour prouver que $\widetilde{M}$ est de type
fini, on est donc ramené au cas où $M$ est de type fini. La question étant
locale, il suffit de prouver que $M_{(f)}$ est un $S_{(f)}$-module de type
fini (\hyperref[I.1.3.9-fr]{I, 1.3.9})~; mais $M^{(d)}$ est un
$S^{(d)}$-module de type fini (\hyperref[II.2.1.6-fr]{2.1.6, (iii)}) et notre
assertion résulte alors de (\hyperref[II.2.2.5-fr]{2.2.5}).

Supposons maintenant que $M$ vérifie (TF) et que $\widetilde{M}=0$~; alors,
avec les mêmes notations que ci-dessus, on a $\widetilde{M'}=0$, et la
condition (TN) pour $M'$ est équivalente à la condition (TN) pour $M$, donc
pour prouver que $\widetilde{M}=0$ implique que $M$ vérifie (TN), on peut
encore se borner au cas où $M$ est engendré par un nombre fini d'éléments
homogènes $x_i$ ($1\leq i\leq p$)~; soit d'autre part
$(f_j)_{1\leq j\leq q}$ un système de générateurs homogènes de l'idéal $S_+$.
On a par hypothèse $M_{(f_j)}=0$ pour tout $j$, donc il existe un entier $n$
tel que $f_j^nx_i=0$ quels que soient $i$ et $j$. Soit $n_j$ le degré de
$f_j$, et soit $m$ la plus grande valeur de $\sum_jr_jn_j$ pour les systèmes
d'entiers $(r_j)$ en nombre fini tels que $\sum_jr_j\leq nq$~; il est clair
alors

\oldpage[II]{40}
que si $k>m$, on a $S_kx_i=0$ pour tout $i$~; si $h$ est le plus grand des
degrés des $x_i$, on en conclut que $M_k=0$ pour $k>h+m$, ce qui termine la
démonstration.
\end{proof}

\begin{corollary}[2.7.4]
\label{II.2.7.4-fr}
Soit $S$ un anneau gradué tel que l'idéal $S_+$ soit de type fini~; pour que
$X=\operatorname{Proj}(S)=\varnothing$, il faut et il suffit qu'il existe $n$
tel que $S_k=0$ pour $k\geq n$.
\end{corollary}

\begin{proof}
La condition $X=\varnothing$ est en effet équivalente à
$\mathcal{O}_X=\widetilde{S}=0$, et $S$ est un $S$-module monogène.
\end{proof}

\begin{theorem}[2.7.5]
\label{II.2.7.5-fr}
On suppose que l'idéal $S_+$ est engendré par un nombre fini d'éléments
homogènes de degré $1$~; soit $X=\operatorname{Proj}(S)$. Alors, pour tout
$\mathcal{O}_X$-Module quasi-cohérent $\mathcal{F}$, l'homomorphisme
canonique $\beta:\widetilde{\Gamma_*(\mathcal{F})}
\to\mathcal{F}$ (\hyperref[II.2.6.4-fr]{2.6.4}) est un isomorphisme.
\end{theorem}

\begin{proof}
En effet, si $S_+$ est engendré par un nombre fini d'éléments $f_i\in S_1$,
$X$ est réunion des sous-espaces $\operatorname{Spec}(S_{(f_i)})$
(\hyperref[II.2.3.6-fr]{2.3.6}) qui sont quasi-compacts, donc est
quasi-compact~; en outre $X$ est un schéma
(\hyperref[II.2.4.2-fr]{2.4.2})~; d'après
(\hyperref[I.9.3.2-fr]{I, 9.3.2}),
(\hyperref[II.2.5.14.2-fr]{2.5.14.2}) et
(\hyperref[II.2.6.3-fr]{2.6.3}) on a, pour tout $f\in S_d$ ($d>0$), un
isomorphisme canonique
$(\Gamma_*(\mathcal{F}))_{(\alpha_d(f))}\xrightarrow{\sim}
\Gamma(D_+(f),\mathcal{F})$~; d'ailleurs, par définition,
$(\Gamma_*(\mathcal{F}))_{(\alpha_d(f))}$ (où $\Gamma_*(\mathcal{F})$ est
considéré comme $\Gamma_*(\mathcal{O}_X)$-module) n'est autre que
$(\Gamma_*(\mathcal{F}))_{(f)}$ (où $\Gamma_*(\mathcal{F})$ est considéré
comme $S$-module)~; si on se reporte à la définition
(\hyperref[I.9.3.1-fr]{I, 9.3.1}) de l'isomorphisme canonique précédent, on
voit qu'il coïncide avec l'homomorphisme $\beta_f$, d'où le théorème.
\end{proof}

\begin{remark}[2.7.6]
\label{II.2.7.6-fr}
Si on suppose l'anneau gradué $S$ \emph{noethérien}, la condition de
(\hyperref[II.2.7.5-fr]{2.7.5}) est vérifiée \emph{ipso facto} dès que l'on
suppose l'idéal $S_+$ engendré par l'ensemble $S_1$ des éléments homogènes de
degré $1$.
\end{remark}

\begin{corollary}[2.7.7]
\label{II.2.7.7-fr}
Sous les hypothèses de (\hyperref[II.2.7.5-fr]{2.7.5}), tout
$\mathcal{O}_X$-Module quasi-cohérent $\mathcal{F}$ est isomorphe à un
$\mathcal{O}_X$-Module de la forme $\widetilde{M}$, où $M$ est un $S$-module
gradué.
\end{corollary}

\begin{corollary}[2.7.8]
\label{II.2.7.8-fr}
Sous les hypothèses de (\hyperref[II.2.7.5-fr]{2.7.5}), tout
$\mathcal{O}_X$-Module quasi-cohérent de type fini $\mathcal{F}$ est
isomorphe à un $\mathcal{O}_X$-Module de la forme $\widetilde{N}$, où $N$ est
un $S$-module gradué de type fini.
\end{corollary}

\begin{proof}
On peut supposer $\mathcal{F}=\widetilde{M}$, où $M$ est un $S$-module gradué
(\hyperref[II.2.7.7-fr]{2.7.7}). Soit $(f_\lambda)_{\lambda\in L}$ un
système de générateurs homogènes de $M$~; pour toute partie finie $H$ de $L$,
soit $M_H$ le sous-module gradué de $M$ engendré par les $f_\lambda$ tels que
$\lambda\in H$~; il est clair que $M$ est limite inductive de ses sous-modules
$M_H$, donc $\mathcal{F}$ est limite inductive de ses sous-$\mathcal{O}_X$-Modules
$\widetilde{M_H}$ (\hyperref[II.2.5.4-fr]{2.5.4}). Mais comme $\mathcal{F}$
est de type fini et l'espace sous-jacent $X$ quasi-compact, il résulte de
(\hyperref[0.5.2.3-fr]{0, 5.2.3}) que $\mathcal{F}=\widetilde{M_H}$ pour une
partie finie $H$ de $L$.
\end{proof}

\begin{corollary}[2.7.9]
\label{II.2.7.9-fr}
Sous les hypothèses de (\hyperref[II.2.7.5-fr]{2.7.5}), soit $\mathcal{F}$ un
$\mathcal{O}_X$-Module quasi-cohérent de type fini. Il existe alors un entier
$n_0$ tel que, pour tout $n\geq n_0$, $\mathcal{F}(n)$ soit isomorphe à un
quotient d'un $\mathcal{O}_X$-Module de la forme $\mathcal{O}_X^k$ ($k>0$
dépendant de $n$), et par suite est engendré par un nombre fini de ses sections
au-dessus de $X$ (\hyperref[0.5.1.1-fr]{0, 5.1.1}).
\end{corollary}

\begin{proof}
En vertu de (\hyperref[II.2.7.8-fr]{2.7.8}), on peut supposer que
$\mathcal{F}=\widetilde{M}$, où $M$ est quotient d'une somme directe finie de
$S$-modules de la forme $S(m_i)$~; en vertu de
(\hyperref[II.2.5.4-fr]{2.5.4}) on peut donc se borner au cas où $M=S(m)$,
donc $\mathcal{F}(n)=\widetilde{S(m+n)}
=\mathcal{O}_X(m+n)$ (\hyperref[II.2.5.15-fr]{2.5.15}). Il suffira donc de
démontrer le

\oldpage[II]{41}
\begin{lemma}[2.7.9.1]
\label{II.2.7.9.1-fr}
Sous les hypothèses de (\hyperref[II.2.7.5-fr]{2.7.5}), pour tout $n\geq0$,
il existe un entier $k$ (dépendant de $n$) et un homomorphisme surjectif
$\mathcal{O}_X^k\to\mathcal{O}_X(n)$.
\end{lemma}

Il suffit en effet (\hyperref[II.2.7.2-fr]{2.7.2}) de montrer que, pour un $k$
convenable, il y a un homomorphisme (TN)-surjectif $u$ \emph{de degré $0$},
du $S$-module gradué produit $S^k$ dans $S(n)$. Or, on a $(S(n))_0=S_n$, et
par hypothèse $S_h=S_1^h$ pour tout $h>0$, donc
$SS_n=S_n+S_{n+1}+\cdots$. Comme $S_n$ est un $S_0$-module de type fini
(\hyperref[II.2.1.5-fr]{2.1.5} et
\hyperref[II.2.1.6-fr]{2.1.6, (i)}), considérons un système de générateurs
$(a_i)_{1\leq i\leq k}$ de ce module~; l'homomorphisme $u$ fera correspondre
$a_i$ au $i$-ème élément de la base canonique de $S^k$ ($1\leq i\leq k$)~;
comme $\operatorname{Coker}u$ s'identifie alors à
$(S(n))_{-n}+\cdots+(S(n))_{-1}$, $u$ répond bien à la question.
\end{proof}

\begin{corollary}[2.7.10]
\label{II.2.7.10-fr}
Sous les hypothèses de (\hyperref[II.2.7.5-fr]{2.7.5}), soit $\mathcal{F}$ un
$\mathcal{O}_X$-Module quasi-cohérent de type fini. Il existe un entier $n_0$
tel que, pour tout $n\geq n_0$, $\mathcal{F}$ soit isomorphe à un quotient
d'un $\mathcal{O}_X$-Module de la forme $(\mathcal{O}_X(-n))^k$ ($k$
dépendant de $n$).
\end{corollary}

\begin{proposition}[2.7.11]
\label{II.2.7.11-fr}
Supposons vérifiées les hypothèses de
(\hyperref[II.2.7.5-fr]{2.7.5}) et soit $M$ un $S$-module gradué. Alors~:
\begin{enumerate}
  \item[\rm(i)] L'homomorphisme canonique
    $\widetilde{\alpha}:\widetilde{M}\to
    \widetilde{\Gamma_*(\widetilde{M})}$ est un isomorphisme.
  \item[\rm(ii)] Soit $\mathcal{G}$ un sous-$\mathcal{O}_X$-Module
    quasi-cohérent de $\widetilde{M}$, et soit $N$ le sous-$S$-module gradué
    de $M$, image réciproque de $\Gamma_*(\mathcal{G})$ par $\alpha$. Alors
    on a $\widetilde{N}=\mathcal{G}$ ($\widetilde{N}$ étant identifié, en
    vertu de (\hyperref[II.2.5.4-fr]{2.5.4}), à un
    sous-$\mathcal{O}_X$-Module de $\widetilde{M}$).
\end{enumerate}
\end{proposition}

\begin{proof}
Comme $\beta:\widetilde{\Gamma_*(\widetilde{M})}
\to\widetilde{M}$ est un isomorphisme en vertu de
(\hyperref[II.2.7.5-fr]{2.7.5}), $\widetilde{\alpha}$ est l'isomorphisme
réciproque en vertu de (\hyperref[II.2.6.5.1-fr]{2.6.5.1}), d'où (i). Soit
$P$ le sous-module gradué $\alpha(M)$ de $\Gamma_*(\widetilde{M})$~; comme
$M\mapsto\widetilde{M}$ est un foncteur exact
(\hyperref[II.2.5.4-fr]{2.5.4}), l'image de $\widetilde{M}$ par
$\widetilde{\alpha}$ est égale à $\widetilde{P}$, donc, en vertu de (i),
$\widetilde{P}=\widetilde{\Gamma_*(\widetilde{M})}$. Posons
$Q=\Gamma_*(\mathcal{G})\cap P$~; en vertu de ce qui précède et de
(\hyperref[II.2.5.4-fr]{2.5.4}), on a
$\widetilde{Q}=\widetilde{\Gamma_*(\mathcal{G})}$, donc la
restriction de $\beta$ à $\widetilde{Q}$ est un \emph{isomorphisme} de ce
$\mathcal{O}_X$-Module sur $\mathcal{G}$ par
(\hyperref[II.2.7.5-fr]{2.7.5}). Mais par définition de $N$ et
(\hyperref[II.2.5.4-fr]{2.5.4}), la restriction de l'isomorphisme
$\widetilde{\alpha}$ à $\widetilde{N}$ est un isomorphisme de $\widetilde{N}$
sur $\widetilde{Q}$, d'où la conclusion par
(\hyperref[II.2.6.5.1-fr]{2.6.5.1}).
\end{proof}

\subsection{Comportements fonctoriels.}
\label{subsection:II.2.8-fr}

\begin{env}[2.8.1]
\label{II.2.8.1-fr}
Soient $S$, $S'$ deux anneaux gradués à degrés positifs,
$\varphi:S'\to S$ un homomorphisme d'anneaux gradués. Nous désignerons par
$G(\varphi)$ la partie ouverte de $X=\operatorname{Proj}(S)$ complémentaire
de $V_+(\varphi(S'_+))$, ou, ce qui revient au même, la réunion des
$D_+(\varphi(f'))$, où $f'$ parcourt l'ensemble des éléments homogènes de
$S'_+$. La restriction à $G(\varphi)$ de l'application continue
${}^a\varphi$ de $\operatorname{Spec}(S)$ dans $\operatorname{Spec}(S')$
(\hyperref[I.1.2.1-fr]{I, 1.2.1}) est donc une application continue de
$G(\varphi)$ dans $\operatorname{Proj}(S')$, que nous noterons encore
${}^a\varphi$ par abus de langage. Si $f'\in S'_+$ est homogène, on a
\[
\label{II.2.8.1.1-fr}
  {}^a\varphi^{-1}(D_+(f'))=D_+(\varphi(f'))
\tag{2.8.1.1}
\]
compte tenu du fait que ${}^a\varphi$ applique $G(\varphi)$ dans
$\operatorname{Proj}(S')$ et de
(\hyperref[I.1.2.2.2-fr]{I, 1.2.2.2}). D'autre part, l'homomorphisme
$\varphi$ définit canoniquement (avec les mêmes notations) un homomorphisme
d'anneaux gradués $S'_{f'}\to S_f$, d'où, par restriction aux éléments de
degré $0$,

\oldpage[II]{42}
un homomorphisme $S'_{(f')}\to S_{(f)}$ que nous désignerons par
$\varphi_{(f)}$~; il lui correspond
(\hyperref[I.1.6.1-fr]{I, 1.6.1}) un morphisme
$({}^a\varphi_{(f)},\widetilde{\varphi}_{(f)}):
\operatorname{Spec}(S_{(f)})\to\operatorname{Spec}(S'_{(f')})$ de schémas
affines. Si on identifie canoniquement $\operatorname{Spec}(S_{(f)})$ au
schéma induit par $\operatorname{Proj}(S)$ sur $D_+(f)$
(\hyperref[II.2.3.6-fr]{2.3.6}), on a défini un morphisme
$\Phi_f:D_+(f)\to D_+(f')$ et ${}^a\varphi_{(f)}$ s'identifie à la
restriction de ${}^a\varphi$ à $D_+(f)$. D'autre part, il est immédiat que si
$g'$ est un second élément homogène de $S'_+$ et $g=\varphi(g')$, le
diagramme
\[
  \xymatrix{
    D_+(f)\ar[r]^{\Phi_f} & D_+(f')\\
    D_+(fg)\ar[u]\ar[r]_{\Phi_{fg}} & D_+(f'g')\ar[u]
  }
\]
est commutatif, en raison de la commutativité du diagramme
\[
  \xymatrix{
    S'_{(f')}\ar[r]^{\varphi_{(f)}}\ar[d]_{\omega_{f'g',f'}}
    & S_{(f)}\ar[d]^{\omega_{fg,f}}\\
    S'_{(f'g')}\ar[r]_{\varphi_{(fg)}} & S_{(fg)}
  }
\]
Compte tenu de la définition de $G(\varphi)$ et de
(\hyperref[II.2.3.3.2-fr]{2.3.3.2}), on voit donc que~:
\end{env}

\begin{proposition}[2.8.2]
\label{II.2.8.2-fr}
Étant donné un homomorphisme d'anneaux gradués $\varphi:S'\to S$, il existe
un morphisme et un seul $({}^a\varphi,\widetilde{\varphi})$ du préschéma
induit $G(\varphi)$ dans $\operatorname{Proj}(S')$ (dit \emph{associé à
$\varphi$} et noté $\operatorname{Proj}(\varphi)$), tel que pour tout élément
homogène $f'\in S'_+$, la restriction de ce morphisme à
$D_+(\varphi(f'))$ coïncide avec le morphisme associé à l'homomorphisme
$S'_{(f')}\to S_{(\varphi(f'))}$ correspondant à $\varphi$.
\end{proposition}

\begin{proof}
On notera qu'avec les notations précédentes, si $f'\in S'_d$, le diagramme
\[
\label{II.2.8.2.1-fr}
  \xymatrix{
    S'_{(f')}\ar[r]^{\varphi_{(f)}}\ar[d]_{\sim}
    & S_{(f)}\ar[d]^{\sim}\\
    {S'}^{(d)}/(f'-1){S'}^{(d)}\ar[r]
    & S^{(d)}/(f-1)S^{(d)}
  }
\tag{2.8.2.1}
\]
est commutatif (les flèches verticales étant les isomorphismes
(\hyperref[II.2.2.5-fr]{2.2.5})).
\end{proof}

\begin{corollary}[2.8.3]
\label{II.2.8.3-fr}
\begin{enumerate}
  \item[\rm(i)] Le morphisme $\operatorname{Proj}(\varphi)$ est affine.
  \item[\rm(ii)] Si $\operatorname{Ker}(\varphi)$ est nilpotent (et en
    particulier si $\varphi$ est injectif), le morphisme
    $\operatorname{Proj}(\varphi)$ est dominant.
\end{enumerate}
\end{corollary}

\begin{proof}
(i) est conséquence immédiate de (\hyperref[II.2.8.2-fr]{2.8.2}) et de
(\hyperref[II.2.8.1.1-fr]{2.8.1.1}). D'autre part, si
$\operatorname{Ker}(\varphi)$ est nilpotent, pour tout $f'$ homogène dans
$S'_+$, on vérifie aussitôt que $\operatorname{Ker}(\varphi_f)$ (avec
$f=\varphi(f')$) est nilpotent, donc aussi
$\operatorname{Ker}(\varphi_{(f)})$~; la conclusion résulte donc de
(\hyperref[II.2.8.2-fr]{2.8.2}) et de
(\hyperref[I.1.2.7-fr]{I, 1.2.7}).
\end{proof}

\oldpage[II]{43}
On notera qu'il y a en général des morphismes de $\operatorname{Proj}(S)$
dans $\operatorname{Proj}(S')$ qui ne sont pas affines, et ne proviennent
donc pas d'homomorphismes d'anneaux gradués $S'\to S$~; un exemple est le
morphisme structural $\operatorname{Proj}(S)\to\operatorname{Spec}(A)$ quand
$A$ est un corps ($\operatorname{Spec}(A)$ étant identifié à
$\operatorname{Proj}(A[T])$ (cf. (\hyperref[II.3.1.7-fr]{3.1.7})))~; cela
résulte en effet de (\hyperref[I.2.3.2-fr]{I, 2.3.2}).

\begin{env}[2.8.4]
\label{II.2.8.4-fr}
Soient $S''$ un troisième anneau gradué à degrés positifs,
$\varphi':S''\to S'$ un homomorphisme d'anneaux gradués, et posons
$\varphi''=\varphi\circ\varphi'$. En raison de
(\hyperref[II.2.8.1.1-fr]{2.8.1.1}) et de la formule
${}^a\varphi''=({}^a\varphi')\circ({}^a\varphi)$ on vérifie aussitôt que l'on
a $G(\varphi'')\subset G(\varphi)$ et que, si $\Phi$, $\Phi'$ et $\Phi''$
sont les morphismes associés à $\varphi$, $\varphi'$ et $\varphi''$, on a
$\Phi''=\Phi'\circ(\Phi|G(\varphi''))$.
\end{env}

\begin{env}[2.8.5]
\label{II.2.8.5-fr}
Supposons que $S$ (resp. $S'$) soit une $A$-algèbre graduée (resp. une
$A'$-algèbre graduée), et soit $\psi:A'\to A$ un homomorphisme d'anneaux tel
que le diagramme
\[
\label{II.2.8.5.1-fr}
  \xymatrix{
    A'\ar[r]^{\psi}\ar[d] & A\ar[d]\\
    S'\ar[r]_{\varphi} & S
  }
\tag{2.8.5.1}
\]
soit commutatif. On peut alors considérer $G(\varphi)$ et
$\operatorname{Proj}(S')$ comme des schémas au-dessus de
$\operatorname{Spec}(A)$ et $\operatorname{Spec}(A')$ respectivement~; si
$\Phi$ et $\Psi$ sont les morphismes associés respectivement à $\varphi$ et
$\psi$, le diagramme
\[
  \xymatrix{
    G(\varphi)\ar[r]^{\Phi}\ar[d] & \operatorname{Proj}(S')\ar[d]\\
    \operatorname{Spec}(A)\ar[r]_{\Psi} & \operatorname{Spec}(A')
  }
\]
est alors commutatif~: il suffit en effet de le démontrer pour la restriction
de $\Phi$ à $D_+(f)$, où $f=\varphi(f')$, $f'$ homogène dans $S'_+$~; et
alors cela résulte de la commutativité du diagramme
\[
  \xymatrix{
    A'\ar[r]^{\psi}\ar[d] & A\ar[d]\\
    S'_{(f')}\ar[r]_{\varphi_{(f)}} & S_{(f)}
  }
\]
\end{env}

\begin{env}[2.8.6]
\label{II.2.8.6-fr}
Soit maintenant $M$ un $S$-module gradué, et considérons le $S'$-module
$M_{[\varphi]}$, qui est évidemment gradué. Soit $f'$ homogène dans $S'_+$,
et soit $f=\varphi(f')$~; on sait
(\hyperref[0.1.5.2-fr]{0, 1.5.2}) qu'il y a un isomorphisme canonique
$(M_{[\varphi]})_{f'}\xrightarrow{\sim}(M_f)_{[\varphi_f]}$, et il est
immédiat que cet isomorphisme conserve les degrés, donc il y a un isomorphisme
canonique
$(M_{[\varphi]})_{(f')}\xrightarrow{\sim}
(M_{(f)})_{[\varphi_{(f)}]}$. Il correspond canoniquement à cet isomorphisme
un isomorphisme de faisceaux
$\widetilde{M_{[\varphi]}}|D_+(f')
\xrightarrow{\sim}(\Phi_f)_*(\widetilde{M}|D_+(f))$
(\hyperref[II.2.5.2-fr]{2.5.2} et
\hyperref[I.1.6.3-fr]{I, 1.6.3}). En outre,

\oldpage[II]{44}
si $g'$ est un second élément homogène de $S'_+$ et $g=\varphi(g')$, le
diagramme
\[
  \xymatrix{
    (M_{[\varphi]})_{(f')}\ar[r]^{\sim}\ar[d]
    & (M_{(f)})_{[\varphi_{(f)}]}\ar[d]\\
    (M_{[\varphi]})_{(f'g')}\ar[r]^{\sim}
    & (M_{(fg)})_{[\varphi_{(fg)}]}
  }
\]
est commutatif, d'où on conclut aussitôt que l'isomorphisme
\[
  \widetilde{M_{[\varphi]}}|D_+(f'g')
  \xrightarrow{\sim}(\Phi_{fg})_*(\widetilde{M}|D_+(fg))
\]
est la restriction à $D_+(f'g')$ de l'isomorphisme
$\widetilde{M_{[\varphi]}}|D_+(f')
\xrightarrow{\sim}(\Phi_f)_*(\widetilde{M}|D_+(f))$. Comme $\Phi_f$ est la
restriction à $D_+(f)$ du morphisme $\Phi$, on voit donc, compte tenu de
(\hyperref[II.2.8.1.1-fr]{2.8.1.1}), et en posant
$X'=\operatorname{Proj}(S')$~:
\end{env}

\begin{proposition}[2.8.7]
\label{II.2.8.7-fr}
Il existe un isomorphisme canonique fonctoriel du $\mathcal{O}_{X'}$-Module
$\widetilde{M_{[\varphi]}}$ sur le
$\mathcal{O}_{X'}$-Module $\Phi_*(\widetilde{M}|G(\varphi))$.
\end{proposition}

On en déduit aussitôt une application canonique fonctorielle de l'ensemble
des $\varphi$-morphismes $M'\to M$ d'un $S'$-module gradué dans le $S$-module
gradué $M$, dans l'ensemble des $\Phi$-morphismes
$\widetilde{M'}\to\widetilde{M}|G(\varphi)$. Avec les notations de
(\hyperref[II.2.8.4-fr]{2.8.4}), si $M''$ est un $S''$-module gradué, au
composé d'un $\varphi$-morphisme $M'\to M$ et d'un $\varphi'$-morphisme
$M''\to M'$ correspond canoniquement le composé de
$\widetilde{M'}|G(\varphi')\to\widetilde{M}|G(\varphi'')$ et de
$\widetilde{M''}\to\widetilde{M'}|G(\varphi')$.

\begin{proposition}[2.8.8]
\label{II.2.8.8-fr}
Sous les hypothèses de (\hyperref[II.2.8.1-fr]{2.8.1}), soit $M'$ un
$S'$-module gradué. Il existe un homomorphisme canonique fonctoriel $\nu$ du
$(\mathcal{O}_X|G(\varphi))$-Module $\Phi^*(\widetilde{M'})$ dans le
$(\mathcal{O}_X|G(\varphi))$-Module
$\widetilde{M'\otimes_{S'}S}|G(\varphi)$. Si l'idéal $S'_+$
est engendré par $S'_1$, $\nu$ est un isomorphisme.
\end{proposition}

\begin{proof}
En effet, pour $f'\in S'_d$ ($d>0$), on définit un homomorphisme canonique
fonctoriel de $S_{(f)}$-modules (où $f=\varphi(f')$)
\[
\label{II.2.8.8.1-fr}
  \nu_f:M'_{(f')}\otimes_{S'_{(f')}}S_{(f)}
  \longrightarrow(M'\otimes_{S'}S)_{(f)}
\tag{2.8.8.1}
\]
en composant l'homomorphisme
$M'_{(f')}\otimes_{S'_{(f')}}S_{(f)}\to
M'_{f'}\otimes_{S'_{f'}}S_f$ et l'isomorphisme canonique
$M'_{f'}\otimes_{S'_{f'}}S_f\xrightarrow{\sim}(M'\otimes_{S'}S)_f$
(\hyperref[0.1.5.4-fr]{0, 1.5.4}) et notant que ce dernier conserve les
degrés. On vérifie aussitôt la compatibilité des $\nu_f$ avec les opérateurs
de restriction de $D_+(f)$ à $D_+(fg)$ (pour un second $g'\in S'_+$ et
$g=\varphi(g')$), d'où la définition de l'homomorphisme
\[
  \nu:\Phi^*(\widetilde{M'})
  \longrightarrow\widetilde{M'\otimes_{S'}S}|G(\varphi)
\]
compte tenu de (\hyperref[I.1.6.5-fr]{I, 1.6.5}). Pour prouver la seconde
assertion, il suffit de montrer que $\nu_f$ est un isomorphisme pour tout
$f'\in S'_1$, puisque $G(\varphi)$ est alors réunion des
$D_+(\varphi(f'))$. On définit d'abord une application $\mathbf{Z}$-bilinéaire
$M'_m\times S_n\to M'_{(f')}\otimes_{S'_{(f')}}S_{(f)}$ en faisant
correspondre à $(x',s)$ l'élément
$(x'/{f'}^m)\otimes(s/f^n)$ (avec la convention que $x'/{f'}^m$ est
${f'}^{-m}x'/1$

\oldpage[II]{45}
lorsque $m<0$). On constate comme dans la démonstration de
(\hyperref[II.2.5.13-fr]{2.5.13}) que cette application donne naissance à un
di-homomorphisme de modules
\[
  \eta_f:M'\otimes_{S'}S
  \longrightarrow M'_{(f')}\otimes_{S'_{(f')}}S_{(f)}.
\]
En outre, si, pour $r>0$, on a
$f^r\sum_i(x'_i\otimes s_i)=0$, cela s'écrit aussi
$\sum_i({f'}^rx'_i\otimes s_i)=0$, d'où, par
(\hyperref[0.1.5.4-fr]{0, 1.5.4}),
$\sum_i({f'}^rx'_i/{f'}^{m_i+r})\otimes(s_i/f^{n_i})=0$, c'est-à-dire
$\eta_f(\sum_i x'_i\otimes s_i)=0$, ce qui prouve que $\eta_f$ se factorise
en
\[
  M'\otimes_{S'}S\longrightarrow(M'\otimes_{S'}S)_f
  \xrightarrow{\eta'_f}M'_{(f')}\otimes_{S'_{(f')}}S_{(f)}~;
\]
on vérifie enfin que $\eta'_f$ et $\nu_f$ sont des isomorphismes réciproques
l'un de l'autre.

En particulier, il résulte de (\hyperref[II.2.1.2.1-fr]{2.1.2.1}) que l'on a
un homomorphisme canonique
\[
\label{II.2.8.8.2-fr}
  \Phi^*(\mathcal{O}_{X'}(n))\xrightarrow{\sim}
  \mathcal{O}_X(n)|G(\varphi)
\tag{2.8.8.2}
\]
pour tout $n\in\mathbf{Z}$.
\end{proof}

\begin{env}[2.8.9]
\label{II.2.8.9-fr}
Soient $A$, $A'$ deux anneaux, $\psi:A'\to A$ un homomorphisme d'anneaux,
définissant un morphisme
$\Psi:\operatorname{Spec}(A)\to\operatorname{Spec}(A')$. Soit $S'$ une
$A'$-algèbre graduée à degrés positifs, et posons $S=S'\otimes_{A'}A$, qui
est de façon évidente une $A$-algèbre graduée par les $S'_n\otimes_{A'}A$~;
l'application $\varphi:s'\mapsto s'\otimes1$ est alors un homomorphisme
d'anneaux gradués rendant commutatif le diagramme
(\hyperref[II.2.8.5.1-fr]{2.8.5.1}). Comme ici $S_+$ est le $A$-module
engendré par $\varphi(S'_+)$, on a $G(\varphi)=\operatorname{Proj}(S)=X$~;
d'où, en posant $X'=\operatorname{Proj}(S')$, un diagramme commutatif
\[
\label{II.2.8.9.1-fr}
  \xymatrix{
    X\ar[r]^{\Phi}\ar[d]_p & X'\ar[d]\\
    Y\ar[r]_{\Psi} & Y'
  }
\tag{2.8.9.1}
\]

Soit maintenant $M'$ un $S'$-module gradué, et posons
$M=M'\otimes_{A'}A=M'\otimes_{S'}S$. Dans ces conditions~:
\end{env}

\begin{proposition}[2.8.10]
\label{II.2.8.10-fr}
Le diagramme (\hyperref[II.2.8.9.1-fr]{2.8.9.1}) identifie le schéma $X$ au
produit $X'\times_{Y'}Y$~; en outre, l'homomorphisme canonique
$\nu:\Phi^*(\widetilde{M'})\to\widetilde{M}$
(\hyperref[II.2.8.8-fr]{2.8.8}) est un isomorphisme.
\end{proposition}

\begin{proof}
La première assertion sera démontrée si l'on prouve que pour tout $f'$
homogène dans $S'_+$, en posant $f=\varphi(f')$, les restrictions de $\Phi$
et de $p$ à $D_+(f)$ identifient ce schéma au produit
$D_+(f')\times_{Y'}Y$ (\hyperref[I.3.2.6.2-fr]{I, 3.2.6.2})~; autrement
dit, il suffit (\hyperref[I.3.2.2-fr]{I, 3.2.2}) de prouver que $S_{(f)}$
s'identifie canoniquement à $S'_{(f')}\otimes_{A'}A$, ce qui est immédiat
en vertu de l'existence de l'isomorphisme canonique
$S_f\xrightarrow{\sim}S'_{f'}\otimes_{A'}A$ conservant les degrés
(\hyperref[0.1.5.4-fr]{0, 1.5.4}). La seconde assertion résulte de ce que
$M'_{(f')}\otimes_{S'_{(f')}}S_{(f)}$ s'identifie d'après ce qui précède à
$M'_{(f')}\otimes_{A'}A$, et ce dernier à $M_{(f)}$, puisque $M_f$
s'identifie canoniquement à $M'_{f'}\otimes_{A'}A$ par un isomorphisme
conservant les degrés.
\end{proof}

\begin{corollary}[2.8.11]
\label{II.2.8.11-fr}
Pour tout entier $n\in\mathbf{Z}$, $\widetilde{M(n)}$ s'identifie à
$\Phi^*(\widetilde{M'(n)})=\widetilde{M'(n)}\otimes_{Y'}\mathcal{O}_Y$~;
en particulier
$\mathcal{O}_X(n)=\Phi^*(\mathcal{O}_{X'}(n))
=\mathcal{O}_{X'}(n)\otimes_{Y'}\mathcal{O}_Y$.
\end{corollary}

\begin{proof}
Cela résulte de (\hyperref[II.2.8.10-fr]{2.8.10}) et de
(\hyperref[II.2.5.15-fr]{2.5.15}).
\end{proof}

\oldpage[II]{46}

\begin{env}[2.8.12]
\label{II.2.8.12-fr}
Sous les hypothèses de (\hyperref[II.2.8.9-fr]{2.8.9}), pour
$f'\in S'_d$ ($d>0$) et $f=\varphi(f')$, le diagramme
\[
\label{II.2.8.12.1-fr}
  \xymatrix{
    M'_{(f')}\ar[r]^{\sim}\ar[d]
    & M'^{(d)}/(f'-1)M'^{(d)}\ar[d]\\
    M_{(f)}\ar[r]^{\sim}
    & M^{(d)}/(f-1)M^{(d)}
  }
\tag{2.8.12.1}
\]
(cf. (\hyperref[II.2.2.5-fr]{2.2.5})) est commutatif.
\end{env}

\begin{env}[2.8.13]
\label{II.2.8.13-fr}
Gardons les notations et hypothèses de (\hyperref[II.2.8.9-fr]{2.8.9}), et
soit $\mathcal{F}'$ un $\mathcal{O}_{X'}$-Module~; si on pose
$\mathcal{F}=\Phi^*(\mathcal{F}')$, on a, pour tout $n\in\mathbf{Z}$,
$\mathcal{F}(n)=\Phi^*(\mathcal{F}'(n))$ en vertu de
(\hyperref[II.2.8.11-fr]{2.8.11}) et
(\hyperref[0.4.3.3-fr]{0, 4.3.3}). Par suite
(\hyperref[0.3.7.1-fr]{0, 3.7.1}), on a un homomorphisme canonique
\[
  \Gamma(\rho):\Gamma(X',\mathcal{F}'(n))
  \longrightarrow\Gamma(X,\mathcal{F}(n))
\]
ce qui donne un di-homomorphisme canonique de modules gradués
\[
  \Gamma_\bullet(\mathcal{F}')\longrightarrow
  \Gamma_\bullet(\mathcal{F}).
\]

Supposons l'idéal $S_+$ engendré par $S_1$, et
$\mathcal{F}'=\widetilde{M'}$, donc $\mathcal{F}=\widetilde{M}$ avec
$M=M'\otimes_{A'}A$. Si $f'$ est homogène dans $S'_+$ et
$f=\varphi(f')$, on a vu que
$M_{(f)}=M'_{(f')}\otimes_{A'}A$, et le diagramme
\[
  \xymatrix{
    M'_0\ar[r]\ar[d]
    & M'_{(f')}\ar[d]
    & =\Gamma(D_+(f'),\widetilde{M'})\\
    M_0\ar[r]
    & M_{(f)}
    & =\Gamma(D_+(f),\widetilde{M})
  }
\]
est donc commutatif~; on conclut aussitôt de cette remarque et de la
définition de l'homomorphisme $\alpha$
(\hyperref[II.2.6.2-fr]{2.6.2}) que le diagramme
\[
\label{II.2.8.13.1-fr}
  \xymatrix{
    M'\ar[r]^{\alpha_{M'}}\ar[d]
    & \Gamma_\bullet(\widetilde{M'})\ar[d]\\
    M\ar[r]_{\alpha_M}
    & \Gamma_\bullet(\widetilde{M})
  }
\tag{2.8.13.1}
\]
est commutatif. De même, le diagramme
\[
\label{II.2.8.13.2-fr}
  \xymatrix{
    \widetilde{\Gamma_\bullet(\mathcal{F}')}\ar[r]^{\beta_{\mathcal{F}'}}
      \ar[d]
    & \mathcal{F}'\ar[d]\\
    \widetilde{\Gamma_\bullet(\mathcal{F})}\ar[r]_{\beta_{\mathcal{F}}}
    & \mathcal{F}
  }
\tag{2.8.13.2}
\]
est commutatif (la flèche verticale de droite étant le $\Phi$-morphisme
canonique
$\mathcal{F}'\to\Phi^*(\mathcal{F}')=\mathcal{F}$).
\end{env}

\oldpage[II]{47}

\begin{env}[2.8.14]
\label{II.2.8.14-fr}
Conservant toujours les hypothèses et notations de
(\hyperref[II.2.8.9-fr]{2.8.9}), soit $N'$ un second $S'$-module gradué, et
soit $N=N'\otimes_{A'}A$. Il est immédiat que les di-homomorphismes
canoniques $M'\to M$, $N'\to N$ donnent un di-homomorphisme
$M'\otimes_{S'}N'\to M\otimes_S N$ (relatif à l'homomorphisme canonique
d'anneaux $S'\to S$), et par suite aussi un $S$-homomorphisme de degré~$0$
\[
  (M'\otimes_{S'}N')\otimes_{A'}A\longrightarrow M\otimes_S N,
\]
auquel correspond (compte tenu de
(\hyperref[II.2.8.10-fr]{2.8.10})) un homomorphisme de
$\mathcal{O}_X$-Modules
\[
\label{II.2.8.14.1-fr}
  \Phi^*(\widetilde{M'\otimes_{S'}N'})
  \longrightarrow\widetilde{M\otimes_S N}.
\tag{2.8.14.1}
\]

En outre, on vérifie aussitôt que le diagramme
\[
\label{II.2.8.14.2-fr}
  \xymatrix{
    \Phi^*(\widetilde{M'}\otimes_{\mathcal{O}_{X'}}\widetilde{N'})
      \ar[r]^{\sim}\ar[d]_{\Phi^*(\lambda)}
    & \widetilde{M}\otimes_{\mathcal{O}_X}\widetilde{N}
      \ar[d]^{\lambda}
    & =\Phi^*(\widetilde{M'})\otimes_{\mathcal{O}_X}
      \Phi^*(\widetilde{N'})\\
    \Phi^*(\widetilde{M'\otimes_{S'}N'})\ar[r]
    & \widetilde{M\otimes_S N}
  }
\tag{2.8.14.2}
\]
est commutatif, la première ligne étant l'isomorphisme canonique
(\hyperref[0.4.3.3-fr]{0, 4.3.3}). Si l'idéal $S'_+$ est engendré par
$S'_1$, il est clair que $S_+$ est engendré par $S_1$, et les deux flèches
verticales de (\hyperref[II.2.8.14.2-fr]{2.8.14.2}) sont alors des
isomorphismes (\hyperref[II.2.5.13-fr]{2.5.13})~; il en est donc de même de
(\hyperref[II.2.8.14.1-fr]{2.8.14.1}).

On a de même un di-homomorphisme canonique
$\operatorname{Hom}_{S'}(M',N')\to\operatorname{Hom}_S(M,N)$ en faisant
correspondre à un homomorphisme $u'$ de degré $k$ l'homomorphisme
$u'\otimes1$, qui est aussi de degré $k$~; on en déduit encore un
$S$-homomorphisme de degré~$0$
\[
  (\operatorname{Hom}_{S'}(M',N'))\otimes_{A'}A
  \longrightarrow\operatorname{Hom}_S(M,N)
\]
d'où un homomorphisme de $\mathcal{O}_X$-Modules
\[
\label{II.2.8.14.3-fr}
  \Phi^*(\widetilde{\operatorname{Hom}_{S'}(M',N')})
  \longrightarrow\widetilde{\operatorname{Hom}_S(M,N)}.
\tag{2.8.14.3}
\]

En outre, le diagramme
\[
  \xymatrix{
    \Phi^*(\widetilde{\operatorname{Hom}_{S'}(M',N')})\ar[r]
      \ar[d]_{\Phi^*(\mu)}
    & \widetilde{\operatorname{Hom}_S(M,N)}\ar[d]^{\mu}\\
    \Phi^*(\mathcal{H}om_{\mathcal{O}_{X'}}
      (\widetilde{M'},\widetilde{N'}))\ar[r]
    & \mathcal{H}om_{\mathcal{O}_X}(\widetilde{M},\widetilde{N})
  }
\]
est commutatif (la seconde ligne horizontale étant l'homomorphisme canonique
(\hyperref[0.4.4.6-fr]{0, 4.4.6})).
\end{env}

\oldpage[II]{48}

\begin{env}[2.8.15]
\label{II.2.8.15-fr}
Les notations et hypothèses étant celles de
(\hyperref[II.2.8.1-fr]{2.8.1}), il résulte de
(\hyperref[II.2.4.7-fr]{2.4.7}) que l'on ne change pas le morphisme $\Phi$,
à des isomorphismes près, en remplaçant $S_0$ et $S'_0$ par
$\mathbf{Z}$ et $\varphi_0$ par l'application identique, puis en remplaçant
$S$ et $S'$ par $S^{(d)}$ et ${S'}^{(d)}$ respectivement ($d>0$) et
$\varphi$ par sa restriction $\varphi^{(d)}$ à $S^{(d)}$.
\end{env}

\subsection{Sous-schémas fermés d'un schéma $\operatorname{Proj}(S)$.}
\label{subsection:II.2.9-fr}

\begin{env}[2.9.1]
\label{II.2.9.1-fr}
Si $\varphi:S\to S'$ est un homomorphisme d'anneaux gradués, on dit que
$\varphi$ est (TN)-\emph{surjectif} (resp. (TN)-\emph{injectif},
(TN)-\emph{bijectif}) s'il existe un entier $n$ tel que, pour $k\geq n$,
$\varphi_k:S_k\to S'_k$ soit surjectif (resp. injectif, bijectif). Alors il
résulte de (\hyperref[II.2.8.15-fr]{2.8.15}) que l'étude de $\Phi$ peut se
ramener au cas où $\varphi$ est \emph{surjectif} (resp. \emph{injectif},
\emph{bijectif}). Au lieu de dire que $\varphi$ est (TN)-bijectif, on dit aussi
que c'est alors un (TN)-\emph{isomorphisme}.
\end{env}

\begin{proposition}[2.9.2]
\label{II.2.9.2-fr}
Soit $S$ un anneau gradué à degrés positifs et soit
$X=\operatorname{Proj}(S)$.
\begin{enumerate}
  \item[\rm(i)] Si $\varphi:S\to S'$ est un homomorphisme (TN)-surjectif
    d'anneaux gradués, le morphisme correspondant $\Phi$
    (\hyperref[II.2.8.1-fr]{2.8.1}) est défini dans
    $\operatorname{Proj}(S')$ tout entier et est une immersion fermée de
    $\operatorname{Proj}(S')$ dans $X$. Si $\mathfrak{J}$ est le noyau de
    $\varphi$, le sous-schéma fermé de $X$ associé à $\Phi$ est défini par
    le faisceau d'idéaux quasi-cohérent $\widetilde{\mathfrak{J}}$ de
    $\mathcal{O}_X$.
  \item[\rm(ii)] Supposons en outre l'idéal $S_+$ engendré par un nombre
    fini d'éléments homogènes de degré~$1$. Soit $X'$ un sous-schéma fermé
    de $X$, défini par un faisceau quasi-cohérent d'idéaux $\mathcal{J}$ de
    $\mathcal{O}_X$. Soit $\mathfrak{J}$ l'idéal gradué de $S$, image
    réciproque de $\Gamma_\bullet(\mathcal{J})$ par l'homomorphisme canonique
    $\alpha:S\to\Gamma_\bullet(\mathcal{O}_X)$
    (\hyperref[II.2.6.2-fr]{2.6.2}), et posons
    $S'=S/\mathfrak{J}$. Alors $X'$ est le sous-schéma associé à
    l'immersion fermée $\operatorname{Proj}(S')\to X$ correspondant à
    l'homomorphisme canonique d'anneaux gradués $S\to S'$.
\end{enumerate}
\end{proposition}

\begin{proof}
\begin{enumerate}
  \item[\rm(i)] On peut supposer $\varphi$ surjectif
    (\hyperref[II.2.9.1-fr]{2.9.1}). Comme par hypothèse
    $\varphi(S_+)$ engendre $S'_+$, on a
    $G(\varphi)=\operatorname{Proj}(S')$. D'autre part, la seconde assertion
    se vérifie localement sur $X$~; soit donc $f$ un élément homogène de
    $S_+$ et posons $f'=\varphi(f)$. Comme $\varphi$ est un homomorphisme
    surjectif d'anneaux gradués, on vérifie aussitôt que
    $\varphi_{(f')}:S_{(f)}\to S'_{(f')}$ est surjectif et que son noyau est
    $\mathfrak{J}_{(f)}$, ce qui achève de prouver (i)
    (\hyperref[I.4.2.3-fr]{I, 4.2.3}).
  \item[\rm(ii)] En vertu de (i), on est ramené à vérifier que
    l'homomorphisme $\widetilde{j}:\widetilde{\mathfrak{J}}\to\mathcal{O}_X$
    déduit de l'injection canonique $j:\mathfrak{J}\to S$ est un
    isomorphisme de $\widetilde{\mathfrak{J}}$ sur $\mathcal{J}$, ce qui
    résulte de (\hyperref[II.2.7.11-fr]{2.7.11}).
\end{enumerate}
\end{proof}

On notera que $\mathfrak{J}$ est le \emph{plus grand} des idéaux gradués
$\mathfrak{J}'$ de $S$ tels que
$\widetilde{j}(\widetilde{\mathfrak{J}'})=\mathcal{J}$, car on vérifie
immédiatement en remontant aux définitions
(\hyperref[II.2.6.2-fr]{2.6.2}) que cette relation implique
$\alpha(\mathfrak{J}')\subset\Gamma_\bullet(\mathcal{J})$.

\begin{corollary}[2.9.3]
\label{II.2.9.3-fr}
On suppose vérifiées les hypothèses de
(\hyperref[II.2.9.2-fr]{2.9.2}, (i)) et en outre que l'idéal $S_+$ est
engendré par $S_1$~; alors $\Phi^*(\widetilde{S(n)})$ est canoniquement
isomorphe à $\widetilde{S'(n)}$ pour tout $n\in\mathbf{Z}$, et par suite
$\Phi^*(\mathcal{F}(n))$ à $\Phi^*(\mathcal{F})(n)$ pour tout
$\mathcal{O}_X$-Module $\mathcal{F}$.
\end{corollary}

\begin{proof}
C'est un cas particulier de (\hyperref[II.2.8.8-fr]{2.8.8}), compte tenu de
(\hyperref[II.2.5.10.2-fr]{2.5.10.2}).
\end{proof}

\begin{corollary}[2.9.4]
\label{II.2.9.4-fr}
On suppose vérifiées les hypothèses de
(\hyperref[II.2.9.2-fr]{2.9.2}, (ii)). Pour que le sous-préschéma fermé $X'$
de $X$ soit intègre, il faut et il suffit que l'idéal gradué
$\mathfrak{J}$ soit premier dans $S$.
\end{corollary}

\oldpage[II]{49}

\begin{proof}
Comme $X'$ est isomorphe à $\operatorname{Proj}(S/\mathfrak{J})$, la
condition est suffisante en vertu de (\hyperref[II.2.4.4-fr]{2.4.4}). Pour
voir qu'elle est nécessaire, considérons la suite exacte
$0\to\mathcal{J}\to\mathcal{O}_X\to\mathcal{O}_X/\mathcal{J}$, qui donne une
suite exacte
\[
  0\longrightarrow\Gamma_\bullet(\mathcal{J})
  \longrightarrow\Gamma_\bullet(\mathcal{O}_X)
  \longrightarrow\Gamma_\bullet(\mathcal{O}_X/\mathcal{J}).
\]
Il suffit de prouver que si $f\in S_m$, $g\in S_n$ sont tels que l'image dans
$\Gamma_\bullet(\mathcal{O}_X/\mathcal{J})$ de $\alpha_{n+m}(fg)$ est nulle,
alors l'une des images de $\alpha_m(f)$, $\alpha_n(g)$ est nulle. Or, par
définition, ces images sont des sections des
$(\mathcal{O}_X/\mathcal{J})$-Modules inversibles
$\mathcal{L}=(\mathcal{O}_X/\mathcal{J})(m)$ et
$\mathcal{L}'=(\mathcal{O}_X/\mathcal{J})(n)$ au-dessus du schéma intègre
$X'$~; l'hypothèse entraîne que le produit de ces deux sections est nul dans
$\mathcal{L}\otimes\mathcal{L}'$
((\hyperref[II.2.9.3-fr]{2.9.3}) et
(\hyperref[II.2.5.14.1-fr]{2.5.14.1})), donc l'une d'elles est nulle en vertu
de (\hyperref[I.7.4.4-fr]{I, 7.4.4}).
\end{proof}

\begin{corollary}[2.9.5]
\label{II.2.9.5-fr}
Soient $A$ un anneau, $M$ un $A$-module, $S$ une $A$-algèbre graduée
engendrée par l'ensemble $S_1$ des éléments homogènes de degré~$1$,
$u:M\to S_1$ un homomorphisme surjectif de $A$-modules,
$\overline{u}:\mathbf{S}(M)\to S$ l'homomorphisme (de $A$-algèbres) de
l'algèbre symétrique $\mathbf{S}(M)$ de $M$ dans $S$ qui prolonge $u$. Le
morphisme correspondant à $\overline{u}$ est alors une immersion fermée de
$\operatorname{Proj}(S)$ dans $\operatorname{Proj}(\mathbf{S}(M))$.
\end{corollary}

\begin{proof}
En effet, $\overline{u}$ est surjectif par hypothèse et il suffit d'appliquer
(\hyperref[II.2.9.2-fr]{2.9.2}).
\end{proof}

\section{Spectre homogène d'un faisceau d'algèbres graduées}
\label{section:II.3-fr}

\subsection{Spectre homogène d'une $\mathcal{O}_Y$-Algèbre graduée
quasi-cohérente.}
\label{subsection:II.3.1-fr}
\label{II.3.1-fr}

\begin{env}[3.1.1]
\label{II.3.1.1-fr}
Soient $Y$ un préschéma, $\mathcal{S}$ une $\mathcal{O}_Y$-Algèbre graduée,
$\mathcal{M}$ un $\mathcal{S}$-Module gradué. Si $\mathcal{S}$ est
\emph{quasi-cohérente}, chacun de ses composants homogènes $\mathcal{S}_n$ est
un $\mathcal{O}_Y$-Module \emph{quasi-cohérent}, étant l'image de
$\mathcal{S}$ par un homomorphisme de $\mathcal{S}$ dans lui-même
((\hyperref[I.1.3.8-fr]{I, 1.3.8}) et
(\hyperref[I.1.3.9-fr]{I, 1.3.9}))~; de même, si $\mathcal{M}$ est
quasi-cohérent en tant que $\mathcal{O}_Y$-Module, ses composants homogènes
$\mathcal{M}_n$ le sont aussi, et réciproquement. Si $d$ est un entier $>0$,
on désignera par $\mathcal{S}^{(d)}$ la somme directe des
$\mathcal{O}_Y$-Modules $\mathcal{S}_{nd}$ ($n\in\mathbf{Z}$), qui est
quasi-cohérente si $\mathcal{S}$ l'est
(\hyperref[I.1.3.9-fr]{I, 1.3.9})~; pour tout entier $k$ tel que
$0\leq k\leq d-1$, on désignera par $\mathcal{M}^{(d,k)}$ (ou
$\mathcal{M}^{(d)}$ pour $k=0$) la somme directe des
$\mathcal{M}_{nd+k}$ ($n\in\mathbf{Z}$), qui est un
$\mathcal{S}^{(d)}$-Module gradué, quasi-cohérent lorsque $\mathcal{S}$ et
$\mathcal{M}$ sont quasi-cohérents (\hyperref[I.9.6.1-fr]{I, 9.6.1}). On
notera $\mathcal{M}(n)$ le $\mathcal{S}$-Module gradué tel que
$(\mathcal{M}(n))_k=\mathcal{M}_{n+k}$ pour tout $k\in\mathbf{Z}$~; si
$\mathcal{S}$ et $\mathcal{M}$ sont quasi-cohérents, $\mathcal{M}(n)$ est un
$\mathcal{S}$-Module gradué quasi-cohérent
(\hyperref[I.9.6.1-fr]{I, 9.6.1}).

Nous dirons que $\mathcal{M}$ est un $\mathcal{S}$-Module gradué
\emph{de type fini} (resp. qu'il admet une \emph{présentation finie}) si pour
tout $y\in Y$, il existe un voisinage ouvert $U$ de $y$ et des entiers $n_i$
(resp. des entiers $m_i$ et $n_j$) tels qu'il y ait un homomorphisme surjectif
de degré~$0$
\[
  \bigoplus_{i=1}^r(\mathcal{S}(n_i)|U)\longrightarrow\mathcal{M}|U
\]
(resp. que $\mathcal{M}|U$ soit isomorphe au conoyau d'un homomorphisme de
degré~$0$
\[
  \bigoplus_{i=1}^r(\mathcal{S}(m_i)|U)
  \longrightarrow\bigoplus_{j=1}^s(\mathcal{S}(n_j)|U)).
\]

Soient $U$ un ouvert affine de $Y$, $A=\Gamma(U,\mathcal{O}_Y)$ son anneau~;
par hypothèse, la $(\mathcal{O}_Y|U)$-Algèbre graduée $\mathcal{S}|U$ est
isomorphe à $\widetilde{S}$, où $S=\Gamma(U,\mathcal{S})$ est une
$A$-algèbre

\oldpage[II]{50}

graduée (\hyperref[I.1.4.3-fr]{I, 1.4.3})~; posons
$X_U=\operatorname{Proj}(\Gamma(U,\mathcal{S}))$. Soient $U'\subset U$ un
second ouvert affine de $Y$, $A'$ son anneau, $j:U'\to U$ l'injection
canonique, qui correspond à l'homomorphisme de restriction $A\to A'$~; on a
$\mathcal{S}|U'=j^*(\mathcal{S}|U)$, et par suite
$S'=\Gamma(U',\mathcal{S})$ s'identifie canoniquement à $S\otimes_A A'$
(\hyperref[I.1.6.4-fr]{I, 1.6.4}). On en conclut
(\hyperref[II.2.8.10-fr]{2.8.10}) que $X_{U'}$ s'identifie canoniquement à
$X_U\times_U U'$, et par suite aussi à $f_U^{-1}(U')$, en désignant par
$f_U$ le morphisme structural $X_U\to U$
(\hyperref[I.4.4.1-fr]{I, 4.4.1}). Nous désignerons par $\sigma_{U',U}$
l'isomorphisme canonique $f_U^{-1}(U')\xrightarrow{\sim}X_{U'}$ ainsi défini,
par $\rho_{U',U}$ l'immersion ouverte $X_{U'}\to X_U$ obtenue en composant
$\sigma_{U',U}^{-1}$ et l'injection canonique $f_U^{-1}(U')\to X_U$. Il est
immédiat que si $U''\subset U'$ est un troisième ouvert affine de $Y$, on a
$\rho_{U'',U}=\rho_{U',U}\circ\rho_{U'',U'}$.
\end{env}

\begin{proposition}[3.1.2]
\label{II.3.1.2-fr}
Soit $Y$ un préschéma. Pour toute $\mathcal{O}_Y$-Algèbre graduée
quasi-cohérente $\mathcal{S}$ à degrés positifs, il existe un préschéma $X$
au-dessus de $Y$, et un seul à un $Y$-isomorphisme près, ayant la propriété
suivante~: si $f$ est le morphisme structural $X\to Y$, pour tout ouvert
affine $U$ de $Y$, il existe un \emph{isomorphisme} $\eta_U$ du préschéma
induit $f^{-1}(U)$ sur $X_U=\operatorname{Proj}(\Gamma(U,\mathcal{S}))$ tel
que, si $V$ est un second ouvert affine de $Y$ contenu dans $U$, le diagramme
\[
\label{II.3.1.2.1-fr}
  \xymatrix{
    f^{-1}(V)\ar[r]^{\eta_V}\ar[d]
    & X_V\ar[d]^{\rho_{V,U}}\\
    f^{-1}(U)\ar[r]_{\eta_U}
    & X_U
  }
\tag{3.1.2.1}
\]
soit commutatif.
\end{proposition}

\begin{proof}
Pour deux ouverts affines $U$, $V$ de $Y$, soit $X_{U,V}$ le préschéma
induit sur $f_U^{-1}(U\cap V)$ par $X_U$~; nous allons définir un
$Y$-isomorphisme $\theta_{U,V}:X_{V,U}\xrightarrow{\sim}X_{U,V}$. Pour cela,
considérons un ouvert affine $W\subset U\cap V$~: en composant les
isomorphismes
\[
  f_U^{-1}(W)\xrightarrow{\sigma_{W,U}}X_W
  \xrightarrow{\sigma_{W,V}^{-1}}f_V^{-1}(W),
\]
on obtient un isomorphisme $\tau_W$, et on vérifie aussitôt que si
$W'\subset W$ est un ouvert affine, $\tau_{W'}$ est la restriction à
$f_U^{-1}(W')$ de $\tau_W$~; les $\tau_W$ sont donc bien les restrictions
d'un $Y$-isomorphisme $\theta_{V,U}$. En outre, si $U$, $V$, $W$ sont trois
ouverts affines de $Y$, $\theta'_{U,V}$, $\theta'_{V,W}$ et
$\theta'_{U,W}$ les restrictions de $\theta_{U,V}$, $\theta_{V,W}$,
$\theta_{U,W}$ aux images réciproques de $U\cap V\cap W$ dans $X_V$, $X_W$,
$X_W$ respectivement, il résulte des définitions précédentes que l'on a
$\theta'_{U,V}\circ\theta'_{V,W}=\theta'_{U,W}$. L'existence de $X$
vérifiant les propriétés énoncées résulte donc de
(\hyperref[I.2.3.1-fr]{I, 2.3.1})~; son unicité à un $Y$-isomorphisme près
est triviale, compte tenu de (\hyperref[II.3.1.2.1-fr]{3.1.2.1}).
\end{proof}

\begin{env}[3.1.3]
\label{II.3.1.3-fr}
Nous dirons que le préschéma $X$ défini dans
(\hyperref[II.3.1.2-fr]{3.1.2}) est le \emph{spectre homogène} de la
$\mathcal{O}_Y$-Algèbre graduée quasi-cohérente $\mathcal{S}$ et nous le
noterons $\operatorname{Proj}(\mathcal{S})$. Il est immédiat que
$\operatorname{Proj}(\mathcal{S})$ est \emph{séparé au-dessus de $Y$}
((\hyperref[II.2.4.2-fr]{2.4.2}) et
(\hyperref[I.5.5.5-fr]{I, 5.5.5}))~; si $\mathcal{S}$ est une
$\mathcal{O}_Y$-Algèbre \emph{de type fini}
(\hyperref[I.9.6.2-fr]{I, 9.6.2}), $\operatorname{Proj}(\mathcal{S})$ est
\emph{de type fini} sur $Y$ ((\hyperref[II.2.7.1-fr]{2.7.1, (ii)}) et
(\hyperref[I.6.3.1-fr]{I, 6.3.1})).

Si $f$ est le morphisme structural $X\to Y$, il est immédiat que pour tout
préschéma induit par $Y$ sur un ouvert $U$ de $Y$, $f^{-1}(U)$ s'identifie au
spectre homogène $\operatorname{Proj}(\mathcal{S}|U)$.
\end{env}

\begin{proposition}[3.1.4]
\label{II.3.1.4-fr}
Soit $f\in\Gamma(Y,\mathcal{S}_d)$ pour $d>0$. Il existe une partie ouverte
$X_f$ de l'espace sous-jacent à $X=\operatorname{Proj}(\mathcal{S})$ ayant la
propriété suivante~: pour tout ouvert affine $U$ de $Y$,

\oldpage[II]{51}

$X_f\cap\varphi^{-1}(U)=D_+(f|U)$ dans $\varphi^{-1}(U)$ identifié à
$X_U=\operatorname{Proj}(\Gamma(U,\mathcal{S}))$ ($\varphi$ désignant le
morphisme structural $X\to Y$). En outre, le préschéma induit sur $X_f$ est
affine au-dessus de $Y$ et est canoniquement isomorphe à
$\operatorname{Spec}(\mathcal{S}^{(d)}/(f-1)\mathcal{S}^{(d)})$
(\hyperref[II.1.3.1-fr]{1.3.1}).
\end{proposition}

\begin{proof}
On a $f|U\in\Gamma(U,\mathcal{S}_d)=(\Gamma(U,\mathcal{S}))_d$. Si $U$,
$U'$ sont deux ouverts affines de $Y$ tels que $U'\subset U$, $f|U'$ est
l'image de $f|U$ par l'homomorphisme de restriction
\[
  \Gamma(U,\mathcal{S})\longrightarrow\Gamma(U',\mathcal{S}),
\]
donc $D_+(f|U')$ est égal (avec les notations de
(\hyperref[II.3.1.1-fr]{3.1.1})) au préschéma induit sur l'image réciproque
$\rho_{U',U}^{-1}(D_+(f|U))$ dans $X_{U'}$
(\hyperref[II.2.8.1-fr]{2.8.1})~; d'où la première assertion. En outre, le
préschéma induit sur $D_+(f|U)$ par $X_U$ s'identifie canoniquement à
$\operatorname{Spec}((\Gamma(U,\mathcal{S}))_{(f|U)})$, ces identifications
étant compatibles avec les homomorphismes de restriction
(\hyperref[II.2.8.1-fr]{2.8.1})~; la seconde assertion résulte alors de
(\hyperref[II.2.2.5-fr]{2.2.5}) et de la commutativité du diagramme
(\hyperref[II.2.8.2.1-fr]{2.8.2.1}).
\end{proof}

On dira encore que $X_f$ (en tant qu'ouvert dans l'espace sous-jacent $X$)
est l'ensemble des $x\in X$ où $f$ \emph{ne s'annule pas}.

\begin{corollary}[3.1.5]
\label{II.3.1.5-fr}
Si $f\in\Gamma(Y,\mathcal{S}_d)$, $g\in\Gamma(Y,\mathcal{S}_e)$, on a
\[
\label{II.3.1.5.1-fr}
  X_{fg}=X_f\cap X_g.
\tag{3.1.5.1}
\]
\end{corollary}

\begin{proof}
Il suffit en effet de considérer l'intersection des deux membres avec un
ensemble $\varphi^{-1}(U)$, où $U$ est un ouvert affine dans $Y$, et
d'appliquer la formule (\hyperref[II.2.3.3.2-fr]{2.3.3.2}).
\end{proof}

\begin{corollary}[3.1.6]
\label{II.3.1.6-fr}
Soit $(f_\alpha)$ une famille de sections de $\mathcal{S}$ au-dessus de $Y$
telle que $f_\alpha\in\Gamma(Y,\mathcal{S}_{d_\alpha})$~; si le faisceau
d'idéaux de $\mathcal{S}$ engendré par cette famille
(\hyperref[0.5.1.1-fr]{0, 5.1.1}) contient tous les $\mathcal{S}_n$ à partir
d'un certain rang, l'espace sous-jacent $X$ est réunion des $X_{f_\alpha}$.
\end{corollary}

\begin{proof}
En effet, pour tout ouvert affine $U$ de $Y$, $\varphi^{-1}(U)$ est réunion
des $X_{f_\alpha}\cap\varphi^{-1}(U)$
(\hyperref[II.2.3.14-fr]{2.3.14}).
\end{proof}

\begin{corollary}[3.1.7]
\label{II.3.1.7-fr}
Soit $\mathcal{A}$ une $\mathcal{O}_Y$-Algèbre quasi-cohérente~; posons
\[
  \mathcal{S}=\mathcal{A}[T]
  =\mathcal{A}\otimes_{\mathbf{Z}}\mathbf{Z}[T]
\]
où $T$ est une indéterminée ($\mathbf{Z}$ et $\mathbf{Z}[T]$ étant
considérés comme faisceaux simples sur $Y$). Alors
$X=\operatorname{Proj}(\mathcal{S})$ s'identifie canoniquement à
$\operatorname{Spec}(\mathcal{A})$. En particulier,
$\operatorname{Proj}(\mathcal{O}_Y[T])$ s'identifie à $Y$.
\end{corollary}

\begin{proof}
En appliquant (\hyperref[II.3.1.6-fr]{3.1.6}) à l'unique section
$f\in\Gamma(Y,\mathcal{S})$ égale à $T$ en chaque point de $Y$, on voit que
$X_f=X$. En outre, on a ici $d=1$, et
$\mathcal{S}^{(1)}/(f-1)\mathcal{S}^{(1)}
=\mathcal{S}/(f-1)\mathcal{S}$ est canoniquement isomorphe à $\mathcal{A}$,
d'où le corollaire (\hyperref[II.1.2.2-fr]{1.2.2}).
\end{proof}

Soit $g\in\Gamma(Y,\mathcal{O}_Y)$~; si on prend
$\mathcal{S}=\mathcal{O}_Y[T]$, on a donc $g\in\Gamma(Y,\mathcal{S}_0)$~; soit
\[
  h=gT\in\Gamma(Y,\mathcal{S}_1).
\]
Si $X=\operatorname{Proj}(\mathcal{S})$, l'identification canonique définie
dans (\hyperref[II.3.1.7-fr]{3.1.7}) identifie $X_h$ à l'ensemble ouvert
$Y_g$ de $Y$ (au sens de (\hyperref[0.5.5.2-fr]{0, 5.5.2}))~: en effet, on
peut se borner au cas où $Y=\operatorname{Spec}(A)$ est affine, et tout se
ramène alors (compte tenu de (\hyperref[II.2.2.5-fr]{2.2.5})) au fait que
l'anneau de fractions $A_g$ s'identifie canoniquement à
$A[T]/(gT-1)A[T]$ (\hyperref[0.1.2.3-fr]{0, 1.2.3}).

\begin{proposition}[3.1.8]
\label{II.3.1.8-fr}
Soit $\mathcal{S}$ une $\mathcal{O}_Y$-Algèbre graduée quasi-cohérente à
degrés positifs.
\begin{enumerate}
  \item[\rm(i)] Pour tout $d>0$, il existe un $Y$-isomorphisme canonique de
    $\operatorname{Proj}(\mathcal{S})$ sur
    $\operatorname{Proj}(\mathcal{S}^{(d)})$.

\oldpage[II]{52}

  \item[\rm(ii)] Soit $\mathcal{S}'$ la $\mathcal{O}_Y$-Algèbre graduée
    somme directe de $\mathcal{O}_Y$ et des $\mathcal{S}_n$ ($n\geq0$)~;
    alors $\operatorname{Proj}(\mathcal{S}')$ et
    $\operatorname{Proj}(\mathcal{S})$ sont canoniquement $Y$-isomorphes.
  \item[\rm(iii)] Soit $\mathcal{L}$ un $\mathcal{O}_Y$-Module inversible
    (\hyperref[0.5.4.1-fr]{0, 5.4.1}) et soit $\mathcal{S}_{(\mathcal{L})}$ la
    $\mathcal{O}_Y$-Algèbre graduée somme directe des
    $\mathcal{S}_d\otimes\mathcal{L}^{\otimes d}$ ($d\geq0$)~; alors
    $\operatorname{Proj}(\mathcal{S})$ et
    $\operatorname{Proj}(\mathcal{S}_{(\mathcal{L})})$ sont canoniquement
    $Y$-isomorphes.
\end{enumerate}
\end{proposition}

\begin{proof}
Dans chacun des trois cas, il suffit de définir l'isomorphisme localement sur
$Y$, la vérification des compatibilités avec les opérations de restriction
d'un ouvert à un ouvert plus petit étant triviale. On peut donc supposer $Y$
affine, et alors (i) résulte de (\hyperref[II.2.4.7-fr]{2.4.7, (i)}) et (ii)
de (\hyperref[II.2.4.8-fr]{2.4.8}). En ce qui concerne (iii), lorsqu'on
suppose en outre $\mathcal{L}$ isomorphe à $\mathcal{O}_Y$ (ce qui est permis,
la question étant locale sur $Y$), l'isomorphie de
$\operatorname{Proj}(\mathcal{S})$ et
$\operatorname{Proj}(\mathcal{S}_{(\mathcal{L})})$ est évidente~; pour
définir un isomorphisme \emph{canonique}, soient
$Y=\operatorname{Spec}(A)$, $\mathcal{S}=\widetilde{S}$, où $S$ est une
$A$-algèbre graduée, et soit $c$ un générateur du $A$-module libre $L$ tel
que $\mathcal{L}=\widetilde{L}$~; alors pour tout $n>0$,
$x_n\mapsto x_n\otimes c^{\otimes n}$ est un $A$-isomorphisme de $S_n$ sur
$S_n\otimes L^{\otimes n}$, et ces $A$-isomorphismes définissent un
$A$-isomorphisme d'algèbres graduées
\[
  \varphi_c:S\longrightarrow S_{(L)}
  =\bigoplus_{n\geq0}S_n\otimes L^{\otimes n}.
\]
Soit alors $f\in S_+$, homogène de degré $d$~; pour tout $x\in S_{nd}$, on a
\[
  \frac{x\otimes c^{\otimes nd}}{(f\otimes c^{\otimes d})^n}
  =\frac{x\otimes(\varepsilon c)^{\otimes nd}}
  {(f\otimes(\varepsilon c)^{\otimes d})^n}
\]
pour tout élément inversible $\varepsilon\in A$, ce qui montre que
l'isomorphisme $S_{(f)}\to(S_{(L)})_{(f\otimes c^d)}$ déduit de $\varphi_c$
est \emph{indépendant} du générateur $c$ de $L$ considéré et achève la
démonstration.
\end{proof}

\begin{env}[3.1.9]
\label{II.3.1.9-fr}
Rappelons ((\hyperref[0.4.1.3-fr]{0, 4.1.3}) et
(\hyperref[I.1.3.14-fr]{I, 1.3.14})) que pour que la
$\mathcal{O}_Y$-Algèbre graduée quasi-cohérente $\mathcal{S}$ soit
\emph{engendrée par le $\mathcal{O}_Y$-Module $\mathcal{S}_1$}, il faut et il
suffit qu'il existe un recouvrement $(U_\alpha)$ de $Y$ par des ouverts
affines tels que l'algèbre graduée $\Gamma(U_\alpha,\mathcal{S})$ sur
$\Gamma(U_\alpha,\mathcal{S}_0)$ soit engendrée par l'ensemble
$\Gamma(U_\alpha,\mathcal{S}_1)$ de ses éléments homogènes de degré~$1$. Pour
tout ouvert $V$ de $Y$, $\mathcal{S}|V$ est alors engendrée par le
$(\mathcal{O}_Y|V)$-Module $\mathcal{S}_1|V$.
\end{env}

\begin{proposition}[3.1.10]
\label{II.3.1.10-fr}
Supposons qu'il existe un recouvrement ouvert affine fini $(U_i)$ de $Y$ tel
que chacune des algèbres graduées $\Gamma(U_i,\mathcal{S})$ soit de type fini
sur $\Gamma(U_i,\mathcal{O}_Y)$. Alors il existe $d>0$ tel que
$\mathcal{S}^{(d)}$ soit engendrée par $\mathcal{S}_d$, $\mathcal{S}_d$ étant
un $\mathcal{O}_Y$-Module de type fini.
\end{proposition}

\begin{proof}
En effet, il résulte de (\hyperref[II.2.1.6-fr]{2.1.6, (v)}) que pour chaque
$i$, il existe un entier $m_i$ tel que
$\Gamma(U_i,\mathcal{S}_{nm_i})=(\Gamma(U_i,\mathcal{S}_{m_i}))^n$ pour tout
$n>0$~; il suffit de prendre pour $d$ un multiple commun des $m_i$, compte
tenu de (\hyperref[II.2.1.6-fr]{2.1.6, (i)}).
\end{proof}

\begin{corollary}[3.1.11]
\label{II.3.1.11-fr}
Sous les hypothèses de (\hyperref[II.3.1.10-fr]{3.1.10}),
$\operatorname{Proj}(\mathcal{S})$ est $Y$-isomorphe à un spectre homogène
$\operatorname{Proj}(\mathcal{S}')$, où $\mathcal{S}'$ est une
$\mathcal{O}_Y$-Algèbre graduée engendrée par $\mathcal{S}'_1$,
$\mathcal{S}'_1$ étant supposé être un $\mathcal{O}_Y$-Module de type fini.
\end{corollary}

\begin{proof}
Il suffit en effet de prendre $\mathcal{S}'=\mathcal{S}^{(d)}$, $d$ étant
déterminé par la propriété de (\hyperref[II.3.1.10-fr]{3.1.10}), et
d'appliquer (\hyperref[II.3.1.8-fr]{3.1.8, (i)}).
\end{proof}

\begin{env}[3.1.12]
\label{II.3.1.12-fr}
Si $\mathcal{S}$ est une $\mathcal{O}_Y$-Algèbre graduée quasi-cohérente à
degrés positifs, on sait (\hyperref[I.5.1.1-fr]{I, 5.1.1}) que son
\emph{nilradical} $\mathcal{N}$ est un $\mathcal{O}_Y$-Module
quasi-cohérent~; nous dirons que $\mathcal{N}_+=\mathcal{N}\cap\mathcal{S}_+$
est le \emph{nilradical de $\mathcal{S}_+$}~; c'est un
$\mathcal{S}_0$-Module quasi-cohérent gradué, car on se ramène immédiatement
au cas où $Y$ est affine, et la proposition résulte alors de
(\hyperref[II.2.1.10-fr]{2.1.10}). Pour tout $y\in Y$,
$(\mathcal{N}_+)_y$ est alors le nilradical de
$(\mathcal{S}_+)_y=(\mathcal{S}_y)_+$
(\hyperref[I.5.1.1-fr]{I, 5.1.1}). On dit que la
$\mathcal{O}_Y$-Algèbre graduée $\mathcal{S}$ est \emph{essentiellement
réduite} si $\mathcal{N}_+=0$, ce qui équivaut

\oldpage[II]{53}

à dire que $\mathcal{S}_y$
est une $\mathcal{O}_y$-algèbre graduée essentiellement réduite pour tout
$y\in Y$. Pour toute $\mathcal{O}_Y$-Algèbre graduée $\mathcal{S}$,
$\mathcal{S}/\mathcal{N}_+$ est essentiellement réduite.

Nous dirons que $\mathcal{S}$ est \emph{intègre} si $\mathcal{S}_y$ est un
anneau intègre pour tout $y\in Y$ et si en outre
$(\mathcal{S}_y)_+=(\mathcal{S}_+)_y\neq0$ pour tout $y\in Y$.
\end{env}

\begin{proposition}[3.1.13]
\label{II.3.1.13-fr}
Soit $\mathcal{S}$ une $\mathcal{O}_Y$-Algèbre graduée à degrés positifs.
Si $X=\operatorname{Proj}(\mathcal{S})$, le $Y$-schéma $X_{\mathrm{red}}$ est
canoniquement isomorphe à $\operatorname{Proj}(\mathcal{S}/\mathcal{N}_+)$~;
en particulier, si $\mathcal{S}$ est essentiellement réduite, $X$ est réduit.
\end{proposition}

\begin{proof}
Le fait que $X'=\operatorname{Proj}(\mathcal{S}/\mathcal{N}_+)$ est réduit
découle immédiatement de (\hyperref[II.2.4.4-fr]{2.4.4, (i)}), la propriété
étant locale~; en outre, pour tout ouvert affine $U\subset Y$,
${\varphi'}^{-1}(U)$ est égal à $(\varphi^{-1}(U))_{\mathrm{red}}$ (en
désignant par $\varphi$ et $\varphi'$ les morphismes structuraux $X\to Y$,
$X'\to Y$)~; on vérifie aussitôt que les $U$-morphismes canoniques
${\varphi'}^{-1}(U)\to\varphi^{-1}(U)$ sont compatibles avec les opérations
de restriction et définissent par suite une immersion fermée $X'\to X$ qui
est un homéomorphisme des espaces sous-jacents~; d'où la conclusion
(\hyperref[I.5.1.2-fr]{I, 5.1.2}).
\end{proof}

\begin{proposition}[3.1.14]
\label{II.3.1.14-fr}
Soient $Y$ un préschéma intègre, $\mathcal{S}$ une
$\mathcal{O}_Y$-Algèbre graduée quasi-cohérente telle que
$\mathcal{S}_0=\mathcal{O}_Y$.
\begin{enumerate}
  \item[\rm(i)] Si $\mathcal{S}$ est intègre
    (\hyperref[II.3.1.12-fr]{3.1.12}), alors
    $X=\operatorname{Proj}(\mathcal{S})$ est intègre et le morphisme
    structural $\varphi:X\to Y$ est dominant.
  \item[\rm(ii)] Supposons en outre que $\mathcal{S}$ soit essentiellement
    réduite. Alors, inversement, si $X$ est intègre et si $\varphi$ est
    dominant, $\mathcal{S}$ est intègre.
\end{enumerate}
\end{proposition}

\begin{proof}
\begin{enumerate}
  \item[\rm(i)] Si $(U_\alpha)$ est une base de $Y$ formée d'ouverts affines
    non vides, il suffit de démontrer la proposition lorsque $Y$ est remplacé
    par un des $U_\alpha$ et $\mathcal{S}$ par $\mathcal{S}|U_\alpha$~: en
    effet, il en résultera d'une part que les espaces sous-jacents
    $\varphi^{-1}(U_\alpha)$ sont des ouverts irréductibles (donc non vides)
    de $X$ tels que
    $\varphi^{-1}(U_\alpha)\cap\varphi^{-1}(U_\beta)\neq\varnothing$ pour tout
    couple d'indices (puisque $U_\alpha\cap U_\beta$ contient un $U_\gamma$),
    donc que $X$ est irréductible (\hyperref[0.2.1.4-fr]{0, 2.1.4})~; d'autre
    part $X$ sera réduit, puisque c'est là une propriété locale, donc $X$
    sera bien intègre et $\varphi(X)$ dense dans $Y$.

    Supposons donc que $Y=\operatorname{Spec}(A)$ où $A$ est intègre
    (\hyperref[I.5.1.4-fr]{I, 5.1.4}) et $\mathcal{S}=\widetilde{S}$, où $S$
    est une $A$-algèbre graduée~; l'hypothèse est que pour tout $y\in Y$,
    $\widetilde{S}_y=S_y$ est un anneau gradué intègre tel que
    $(S_y)_+\neq0$. Il suffit de prouver que $S$ est un anneau \emph{intègre},
    car alors on aura $S_+\neq0$ et on pourra appliquer
    (\hyperref[II.2.4.4-fr]{2.4.4, (ii)}). Or, soient $f$, $g$ deux éléments
    $\neq0$ de $S$ et supposons que $fg=0$~; pour tout $y\in Y$ on aura donc
    $(f/1)(g/1)=0$ dans $S_y$, donc $f/1=0$ ou $g/1=0$ par hypothèse.
    Supposons par exemple $f/1=0$ dans $S_y$~; cela signifie qu'il existe
    $a\in A$ tel que $a\notin\mathfrak{j}_y$ et $af=0$~; on a alors pour tout
    $z\in Y$, $(a/1)(f/1)=0$ dans l'anneau \emph{intègre} $S_z$, et comme
    $a/1\neq0$ (puisque $A$ est intègre), $f/1=0$, ce qui implique $f=0$.
  \item[\rm(ii)] La question étant locale sur $Y$, on peut encore supposer
    que $Y=\operatorname{Spec}(A)$, $A$ étant intègre, et
    $\mathcal{S}=\widetilde{S}$. Par hypothèse, pour tout $y\in Y$, $(S_y)_+$
    ne contient pas d'élément nilpotent $\neq0$, et il en est de même de
    $(S_0)_y=A_y$ par hypothèse~; donc $S_y$ est un anneau réduit pour tout
    $y\in Y$, et on en conclut d'abord que $S$ lui-même est réduit
    (\hyperref[I.5.1.1-fr]{I, 5.1.1}). L'hypothèse que $X$ est intègre
    entraîne d'autre part que $S$ est essentiellement intègre
    (\hyperref[II.2.4.4-fr]{2.4.4, (ii)}), et tout revient alors à voir que
    l'annulateur $\mathfrak{J}$ de $S_+$ dans $A=S_0$ est réduit

\oldpage[II]{54}

    à $0$
    (\hyperref[II.2.1.11-fr]{2.1.11}). Dans le cas contraire, on aurait
    $(S_h)_+=0$ pour un $h\neq0$ dans $\mathfrak{J}$, donc
    (\hyperref[II.3.1.1-fr]{3.1.1}) $\varphi^{-1}(D(h))=\varnothing$, et
    $\varphi(X)$ ne serait pas dense dans $Y$ contrairement à l'hypothèse
    (puisque $D(h)\neq\varnothing$, $h$ n'étant pas nilpotent).
\end{enumerate}
\end{proof}

\subsection{Faisceau sur $\operatorname{Proj}(\mathcal{S})$ associé à un
$\mathcal{S}$-Module gradué.}
\label{subsection:II.3.2-fr}

\begin{env}[3.2.1]
\label{II.3.2.1-fr}
Soient $Y$ un préschéma, $\mathcal{S}$ une $\mathcal{O}_Y$-Algèbre graduée
quasi-cohérente à degrés positifs, $\mathcal{M}$ un $\mathcal{S}$-module
gradué quasi-cohérent (sur $(Y,\mathcal{O}_Y)$ ou sur l'espace annelé
$(Y,\mathcal{S})$, ce qui revient au même
(\hyperref[I.9.6.1-fr]{I, 9.6.1})). Avec les notations de
(\hyperref[II.3.1.1-fr]{3.1.1}), désignons par $\widetilde{\mathcal{M}}_U$ le
$\mathcal{O}_{X_U}$-Module quasi-cohérent
$\widetilde{\Gamma(U,\mathcal{M})}$~; pour $U'\subset U$,
$\Gamma(U',\mathcal{M})$ s'identifie canoniquement à
$\Gamma(U,\mathcal{M})\otimes_A A'$ (\hyperref[I.1.6.4-fr]{I, 1.6.4})~; donc
on a $\widetilde{\mathcal{M}}_{U'}
=\rho_{U',U}^*(\widetilde{\mathcal{M}}_U)$
(\hyperref[II.2.8.11-fr]{2.8.11}).
\end{env}

\begin{proposition}[3.2.2]
\label{II.3.2.2-fr}
Il existe sur $\operatorname{Proj}(\mathcal{S})=X$ un
$\mathcal{O}_X$-Module quasi-cohérent et un seul $\widetilde{\mathcal{M}}$ tel
que, pour tout ouvert affine $U$ de $Y$, on ait
\[
  \eta_U^*(\widetilde{\Gamma(U,\mathcal{M})})
  =\widetilde{\mathcal{M}}|f^{-1}(U)
\]
(en désignant par $\eta_U$ l'isomorphisme
$f^{-1}(U)\xrightarrow{\sim}\operatorname{Proj}(\Gamma(U,\mathcal{S}))$, où
$f$ est le morphisme structural $X\to Y$).
\end{proposition}

\begin{proof}
Comme $\rho_{U',U}$ s'identifie au morphisme d'injection
$f^{-1}(U')\to f^{-1}(U)$ (\hyperref[II.3.1.2.1-fr]{3.1.2.1}), la proposition
résulte aussitôt de la relation
$\widetilde{\mathcal{M}}_{U'}=\rho_{U',U}^*(\widetilde{\mathcal{M}}_U)$ et du
principe de recollement des faisceaux (\hyperref[0.3.3.1-fr]{0, 3.3.1}).
\end{proof}

On dit que $\widetilde{\mathcal{M}}$ est le $\mathcal{O}_X$-Module associé au
$\mathcal{S}$-Module gradué quasi-cohérent $\mathcal{M}$.

\begin{proposition}[3.2.3]
\label{II.3.2.3-fr}
Soit $\mathcal{M}$ un $\mathcal{S}$-Module gradué quasi-cohérent, et soit
$f\in\Gamma(Y,\mathcal{S}_d)$ ($d>0$). Si $\xi_f$ est l'isomorphisme canonique
de $X_f$ sur le $Y$-préschéma
$Z_f=\operatorname{Spec}(\mathcal{S}^{(d)}/(f-1)\mathcal{S}^{(d)})$
(\hyperref[II.3.1.4-fr]{3.1.4}), $(\xi_f)_*(\widetilde{\mathcal{M}}|X_f)$ est
le $\mathcal{O}_{Z_f}$-Module
$\widetilde{\mathcal{M}^{(d)}/(f-1)\mathcal{M}^{(d)}}$
(\hyperref[II.1.4.3-fr]{1.4.3}).
\end{proposition}

\begin{proof}
La question étant locale sur $Y$, on est aussitôt ramené à
(\hyperref[II.2.2.5-fr]{2.2.5}), compte tenu de la commutativité du diagramme
(\hyperref[II.2.8.12.1-fr]{2.8.12.1}).
\end{proof}

\begin{proposition}[3.2.4]
\label{II.3.2.4-fr}
Le $\mathcal{O}_X$-Module $\widetilde{\mathcal{M}}$ est un foncteur covariant
additif exact en $\mathcal{M}$, de la catégorie des $\mathcal{S}$-Modules
gradués quasi-cohérents dans celle des $\mathcal{O}_X$-Modules
quasi-cohérents, qui commute aux limites inductives et aux sommes directes.
\end{proposition}

\begin{proof}
La question étant locale sur $Y$, se ramène à
(\hyperref[I.1.3.11-fr]{I, 1.3.11} et
\hyperref[I.1.3.9-fr]{I, 1.3.9}) et à
(\hyperref[II.2.5.4-fr]{2.5.4}).
\end{proof}

En particulier, si $\mathcal{N}$ est un sous-$\mathcal{S}$-Module gradué
quasi-cohérent de $\mathcal{M}$, $\widetilde{\mathcal{N}}$ s'identifie
canoniquement à un sous-$\mathcal{O}_X$-Module quasi-cohérent de
$\widetilde{\mathcal{M}}$~; plus particulièrement, pour tout faisceau gradué
quasi-cohérent $\mathcal{J}$ d'idéaux de $\mathcal{S}$,
$\widetilde{\mathcal{J}}$ est un faisceau quasi-cohérent d'idéaux de
$\mathcal{O}_X$.

Si $\mathcal{M}$ est un $\mathcal{S}$-Module gradué quasi-cohérent et
$\mathcal{I}$ un faisceau quasi-cohérent d'idéaux de $\mathcal{O}_Y$,
$\mathcal{I}\mathcal{M}$ est un sous-$\mathcal{S}$-Module gradué
quasi-cohérent de $\mathcal{M}$ et on a
\[
\label{II.3.2.4.1-fr}
  \widetilde{\mathcal{I}\mathcal{M}}
  =\mathcal{I}\cdot\widetilde{\mathcal{M}}
\tag{3.2.4.1}
\]
(le second membre ayant le sens défini en
(\hyperref[0.4.3.5-fr]{0, 4.3.5})). Il suffit en effet de vérifier cette
formule lorsque $Y=\operatorname{Spec}(A)$ est affine,
$\mathcal{S}=\widetilde{S}$, où $S$ est une $A$-algèbre graduée,
$\mathcal{M}=\widetilde{M}$,

\oldpage[II]{55}

où $M$ est un $S$-module gradué, et
$\mathcal{I}=\widetilde{\mathfrak{I}}$, où $\mathfrak{I}$ est un idéal de
$A$. Pour tout $f$ homogène de $S_+$, la restriction à
$D_+(f)=\operatorname{Spec}(S_{(f)})$ du premier membre de
(\hyperref[II.3.2.4.1-fr]{3.2.4.1}) est associée à
$(\mathfrak{I}M)_{(f)}=\mathfrak{I}\cdot M_{(f)}$, et il en est de même de la
restriction du second membre, vu
(\hyperref[I.1.3.13-fr]{I, 1.3.13} et
\hyperref[I.1.6.9-fr]{I, 1.6.9}).

\begin{proposition}[3.2.5]
\label{II.3.2.5-fr}
Soit $f\in\Gamma(Y,\mathcal{S}_d)$ ($d>0$). Sur l'ensemble ouvert $X_f$, le
$(\mathcal{O}_X|X_f)$-Module
$\widetilde{\mathcal{S}(nd)}|X_f$ est canoniquement isomorphe à
$\mathcal{O}_X|X_f$ pour tout $n\in\mathbf{Z}$. En particulier, si la
$\mathcal{O}_Y$-Algèbre $\mathcal{S}$ est engendrée par $\mathcal{S}_1$
(\hyperref[II.3.1.9-fr]{3.1.9}), les $\mathcal{O}_X$-Modules
$\widetilde{\mathcal{S}(n)}$ sont inversibles pour tout $n\in\mathbf{Z}$.
\end{proposition}

\begin{proof}
En effet, pour tout ouvert affine $U$ de $Y$, on a défini dans
(\hyperref[II.2.5.7-fr]{2.5.7}) un isomorphisme canonique de
$\widetilde{\mathcal{S}(nd)}|(X_f\cap\varphi^{-1}(U))$ sur
$\mathcal{O}_X|(X_f\cap\varphi^{-1}(U))$, compte tenu de
(\hyperref[II.3.1.4-fr]{3.1.4}) (où $\varphi$ est le morphisme structural
$X\to Y$)~; il est immédiat de vérifier que ces isomorphismes sont compatibles
avec la restriction de $U$ à un ouvert affine $U'\subset U$, d'où la première
assertion. Pour établir la seconde, il suffit de remarquer que si
$\mathcal{S}$ est engendrée par $\mathcal{S}_1$, il y a un recouvrement
$(U_\alpha)$ de $Y$ par des ouverts affines tels que
$\Gamma(U_\alpha,\mathcal{S})$ soit engendrée par
$\Gamma(U_\alpha,\mathcal{S})_1=\Gamma(U_\alpha,\mathcal{S}_1)$~; on est alors
ramené à appliquer le résultat de (\hyperref[II.2.5.9-fr]{2.5.9}), la
propriété d'être inversible étant locale.
\end{proof}

On posera encore, pour tout $n\in\mathbf{Z}$,
\[
\label{II.3.2.5.1-fr}
  \mathcal{O}_X(n)=\widetilde{\mathcal{S}(n)}
\tag{3.2.5.1}
\]
et pour tout $\mathcal{O}_X$-Module $\mathcal{F}$
\[
\label{II.3.2.5.2-fr}
  \mathcal{F}(n)=\mathcal{F}\otimes_{\mathcal{O}_X}\mathcal{O}_X(n).
\tag{3.2.5.2}
\]
Il résulte aussitôt de ces définitions que pour tout ouvert $U$ de $Y$, on a
\[
  \widetilde{(\mathcal{S}|U)(n)}
  =\mathcal{O}_X(n)|f^{-1}(U)
\]
$f$ étant le morphisme structural $X\to Y$.

\begin{proposition}[3.2.6]
\label{II.3.2.6-fr}
Soient $\mathcal{M}$ et $\mathcal{N}$ deux $\mathcal{S}$-Modules gradués
quasi-cohérents, il existe un homomorphisme canonique fonctoriel (en
$\mathcal{M}$ et $\mathcal{N}$)
\[
\label{II.3.2.6.1-fr}
  \lambda:\widetilde{\mathcal{M}}
  \otimes_{\mathcal{O}_X}\widetilde{\mathcal{N}}
  \longrightarrow\widetilde{\mathcal{M}\otimes_{\mathcal{S}}\mathcal{N}}
\tag{3.2.6.1}
\]
et un homomorphisme canonique fonctoriel (en $\mathcal{M}$ et
$\mathcal{N}$)
\[
\label{II.3.2.6.2-fr}
  \mu:\widetilde{\mathcal{H}om_{\mathcal{S}}(\mathcal{M},\mathcal{N})}
  \longrightarrow
  \mathcal{H}om_{\mathcal{O}_X}
  (\widetilde{\mathcal{M}},\widetilde{\mathcal{N}}).
\tag{3.2.6.2}
\]
En outre, si $\mathcal{S}$ est engendrée par $\mathcal{S}_1$
(\hyperref[II.3.1.9-fr]{3.1.9}), $\lambda$ est un isomorphisme~; si de plus
$\mathcal{M}$ admet une présentation finie
(\hyperref[II.3.1.1-fr]{3.1.1}), $\mu$ est un isomorphisme.
\end{proposition}

\begin{proof}
Les isomorphismes $\lambda$ et $\mu$ ont été définis dans
(\hyperref[II.2.5.11.2-fr]{2.5.11.2}) et
(\hyperref[II.2.5.12.2-fr]{2.5.12.2}) lorsque $Y$ est affine~; ces définitions
étant locales se transportent aussitôt au cas général considéré ici, compte
tenu de (\hyperref[II.2.8.14-fr]{2.8.14}).
\end{proof}

\begin{corollary}[3.2.7]
\label{II.3.2.7-fr}
Si $\mathcal{S}$ est engendrée par $\mathcal{S}_1$, on a, quels que soient
$m$, $n$ dans $\mathbf{Z}$,
\[
\label{II.3.2.7.1-fr}
  \mathcal{O}_X(m)\otimes_{\mathcal{O}_X}\mathcal{O}_X(n)
  =\mathcal{O}_X(m+n)
\tag{3.2.7.1}
\]
\[
\label{II.3.2.7.2-fr}
  \mathcal{O}_X(n)=(\mathcal{O}_X(1))^{\otimes n}
\tag{3.2.7.2}
\]
à des isomorphismes canoniques près.
\end{corollary}

\oldpage[II]{56}

\begin{corollary}[3.2.8]
\label{II.3.2.8-fr}
Supposons $\mathcal{S}$ engendrée par $\mathcal{S}_1$. Pour tout
$\mathcal{S}$-Module gradué $\mathcal{M}$ et tout $n\in\mathbf{Z}$, on a
\[
\label{II.3.2.8.1-fr}
  (\mathcal{M}(n))^{\sim}=\widetilde{\mathcal{M}}(n)
\tag{3.2.8.1}
\]
à un isomorphisme canonique près.
\end{corollary}

\begin{proof}
Cela résulte des propriétés correspondantes pour $Y$ affine
(\hyperref[II.2.5.14-fr]{2.5.14} et
\hyperref[II.2.5.15-fr]{2.5.15}) ainsi que de
(\hyperref[II.2.8.11-fr]{2.8.11}).
\end{proof}

\par\medskip\noindent\textit{Remarques (3.2.9). ---\ }
\label{II.3.2.9-fr}
\begin{enumerate}
  \item[\rm(i)] Si $\mathcal{S}=\mathcal{A}[T]$, $\mathcal{A}$ étant une
    $\mathcal{O}_Y$-Algèbre quasi-cohérente
    (\hyperref[II.3.1.7-fr]{3.1.7}), on vérifie immédiatement que tous les
    $\mathcal{O}_X$-Modules inversibles $\mathcal{O}_X(n)$ sont canoniquement
    isomorphes à $\mathcal{O}_X$.

    En outre, soit $\mathcal{N}$ un $\mathcal{A}$-Module quasi-cohérent, et
    posons $\mathcal{M}=\mathcal{N}\otimes_{\mathcal{A}}\mathcal{A}[T]$. Il
    résulte alors de (\hyperref[II.3.2.3-fr]{3.2.3}) et de
    (\hyperref[II.3.1.7-fr]{3.1.7}) que dans l'identification canonique de
    $X=\operatorname{Proj}(\mathcal{A}[T])$ et de
    $X'=\operatorname{Spec}(\mathcal{A})$, le $\mathcal{O}_X$-Module
    $\widetilde{\mathcal{M}}$ s'identifie au $\mathcal{O}_{X'}$-Module
    $\widetilde{\mathcal{N}}$ associé à $\mathcal{N}$ (au sens de
    (\hyperref[II.1.4.3-fr]{1.4.3})).
  \item[\rm(ii)] Soit $\mathcal{S}$ une $\mathcal{O}_Y$-Algèbre graduée
    quelconque, $\mathcal{S}'$ la $\mathcal{O}_Y$-Algèbre graduée telle que
    $\mathcal{S}'_0=\mathcal{O}_Y$, $\mathcal{S}'_n=\mathcal{S}_n$ pour tout
    $n>0$~; l'isomorphisme canonique de
    $X=\operatorname{Proj}(\mathcal{S})$ sur
    $X'=\operatorname{Proj}(\mathcal{S}')$
    (\hyperref[II.3.1.8-fr]{3.1.8, (ii)}) identifie $\mathcal{O}_X(n)$ et
    $\mathcal{O}_{X'}(n)$ pour tout $n\in\mathbf{Z}$~: cela résulte de la même
    proposition pour le cas affine (\hyperref[II.2.5.16-fr]{2.5.16}) et du
    fait que ces identifications, pour les ouverts affines de $Y$, commutent
    avec les opérations de restriction. De même, soit
    $X^{(d)}=\operatorname{Proj}(\mathcal{S}^{(d)})$~; l'isomorphisme
    canonique de $X$ sur $X^{(d)}$
    (\hyperref[II.3.1.8-fr]{3.1.8, (i)}) identifie $\mathcal{O}_X(nd)$ et
    $\mathcal{O}_{X^{(d)}}(n)$ pour tout $n\in\mathbf{Z}$.
\end{enumerate}

\begin{proposition}[3.2.10]
\label{II.3.2.10-fr}
Soient $\mathcal{L}$ un $\mathcal{O}_Y$-Module inversible, $g$ l'isomorphisme
canonique
$X_{(\mathcal{L})}=\operatorname{Proj}(\mathcal{S}_{(\mathcal{L})})
\to X=\operatorname{Proj}(\mathcal{S})$
(\hyperref[II.3.1.8-fr]{3.1.8, (iii)}). Pour tout entier $n\in\mathbf{Z}$,
$g_*(\mathcal{O}_{X_{(\mathcal{L})}}(n))$ est canoniquement isomorphe à
$\mathcal{O}_X(n)\otimes_Y\mathcal{L}^{\otimes n}$.
\end{proposition}

\begin{proof}
Supposons d'abord $Y$ affine, d'anneau $A$, et
$\mathcal{L}=\widetilde{L}$, où $L$ est un $A$-module libre monogène. Avec les
notations de la démonstration de
(\hyperref[II.3.1.8-fr]{3.1.8, (iii)}) on définit, pour $f\in S_d$, un
isomorphisme de $(S(n))_{(f)}\otimes_A L^{\otimes n}$ sur
$(S_{(L)}(n))_{(f\otimes c^{\otimes d})}$ en faisant correspondre à
$(x/f^k)\otimes c^{\otimes n}$, où $x\in S_{kd+n}$, l'élément
$(x\otimes c^{\otimes(n+kd)})/(f\otimes c^{\otimes d})^k$~; il est immédiat
que cet isomorphisme est indépendant du générateur $c$ choisi pour $L$~; en
outre, les isomorphismes ainsi définis pour chaque $f\in S_+$ sont compatibles
avec les opérations de restriction $D_+(f)\to D_+(fg)$. Enfin, dans le cas
général, on voit sans peine en revenant aux définitions
(\hyperref[II.3.1.1-fr]{3.1.1}) que les isomorphismes ainsi définis pour
chaque ouvert affine $U$ de $Y$ sont compatibles avec le passage de $U$ à un
ouvert affine $U'\subset U$.
\end{proof}

\subsection{$\mathcal{S}$-Module gradué associé à un faisceau sur
$\operatorname{Proj}(\mathcal{S})$.}
\label{subsection:II.3.3-fr}

\emph{On suppose dans tout ce numéro que la $\mathcal{O}_Y$-Algèbre graduée
$\mathcal{S}$ est engendrée par $\mathcal{S}_1$
(\hyperref[II.3.1.9-fr]{3.1.9}).} Rappelons qu'en raison de
(\hyperref[II.3.1.8-fr]{3.1.8, (i)}), cette restriction n'est pas essentielle,
moyennant les conditions de finitude (\hyperref[II.3.1.10-fr]{3.1.10}).

\begin{env}[3.3.1]
\label{II.3.3.1-fr}
Soit $p$ le morphisme structural $X=\operatorname{Proj}(\mathcal{S})\to Y$.
Pour tout $\mathcal{O}_X$-Module $\mathcal{F}$, nous poserons
\[
\label{II.3.3.1.1-fr}
  \Gamma_*(\mathcal{F})
  =\bigoplus_{n\in\mathbf{Z}}p_*(\mathcal{F}(n))
\tag{3.3.1.1}
\]

\oldpage[II]{57}

et en particulier
\[
\label{II.3.3.1.2-fr}
  \Gamma_*(\mathcal{O}_X)
  =\bigoplus_{n\in\mathbf{Z}}p_*(\mathcal{O}_X(n)).
\tag{3.3.1.2}
\]
On sait (\hyperref[0.4.2.2-fr]{0, 4.2.2}) qu'il existe un homomorphisme
canonique
\[
  p_*(\mathcal{F})\otimes_{\mathcal{O}_Y}p_*(\mathcal{G})
  \longrightarrow
  p_*(\mathcal{F}\otimes_{\mathcal{O}_X}\mathcal{G})
\]
pour deux $\mathcal{O}_X$-Modules $\mathcal{F}$ et $\mathcal{G}$~; on déduit
donc de (\hyperref[II.3.2.7.1-fr]{3.2.7.1}) que
$\Gamma_*(\mathcal{O}_X)$ est muni d'une structure de
$\mathcal{O}_Y$-\emph{Algèbre graduée}, et
(\hyperref[II.3.2.5.2-fr]{3.2.5.2}) définit de même sur
$\Gamma_*(\mathcal{F})$ une structure de \emph{Module gradué} sur
$\Gamma_*(\mathcal{O}_X)$.

En vertu de (\hyperref[II.3.2.5-fr]{3.2.5}), et de l'exactitude à gauche du
foncteur $p_*$ (\hyperref[0.4.2.1-fr]{0, 4.2.1}),
$\Gamma_*(\mathcal{F})$ est un foncteur covariant additif et
\emph{exact à gauche} en $\mathcal{F}$ dans la catégorie des
$\mathcal{O}_X$-Modules, à valeurs dans la catégorie des
$\mathcal{O}_Y$-Modules gradués (où les morphismes sont les homomorphismes de
degré $0$). En particulier, si $\mathcal{J}$ est un faisceau d'idéaux dans
$\mathcal{O}_X$, $\Gamma_*(\mathcal{J})$ s'identifie à un
\emph{faisceau gradué d'idéaux} de $\Gamma_*(\mathcal{O}_X)$.
\end{env}

\begin{env}[3.3.2]
\label{II.3.3.2-fr}
Soit $\mathcal{M}$ un $\mathcal{S}$-Module gradué quasi-cohérent. Pour tout
ouvert affine $U$ de $Y$, on a défini en
(\hyperref[II.2.6.2-fr]{2.6.2}) un homomorphisme de groupes abéliens
\[
  \alpha_{0,U}:\Gamma(U,\mathcal{M}_0)
  \longrightarrow\Gamma(p^{-1}(U),\widetilde{\mathcal{M}}).
\]
Il est immédiat que ces homomorphismes commutent aux opérations de restriction
(\hyperref[II.2.8.13.1-fr]{2.8.13.1}) et définissent donc (sans utiliser
l'hypothèse que $\mathcal{S}$ est engendrée par $\mathcal{S}_1$) un
homomorphisme de faisceaux de groupes abéliens
\[
\label{II.3.3.2.1-fr}
  \alpha_0:\mathcal{M}_0\longrightarrow p_*(\widetilde{\mathcal{M}}).
\tag{3.3.2.1}
\]
Appliquant ce résultat à chacun des
$\mathcal{M}_n=(\mathcal{M}(n))_0$, et tenant compte de
(\hyperref[II.3.2.8.1-fr]{3.2.8.1}), on définit un homomorphisme de faisceaux
de groupes abéliens
\[
\label{II.3.3.2.2-fr}
  \alpha_n:\mathcal{M}_n\longrightarrow
  p_*(\widetilde{\mathcal{M}}(n))
  \qquad\text{pour tout }n\in\mathbf{Z},
\tag{3.3.2.2}
\]
d'où un homomorphisme fonctoriel (de degré $0$) de faisceaux gradués de groupes
abéliens
\[
\label{II.3.3.2.3-fr}
  \alpha:\mathcal{M}\longrightarrow\Gamma_*(\widetilde{\mathcal{M}})
\tag{3.3.2.3}
\]
(aussi noté $\alpha_{\mathcal{M}}$).

En prenant en particulier $\mathcal{M}=\mathcal{S}$, on vérifie que
$\alpha:\mathcal{S}\to\Gamma_*(\mathcal{O}_X)$ est un homomorphisme de
$\mathcal{O}_Y$-Algèbres graduées et que
(\hyperref[II.3.3.2.3-fr]{3.3.2.3}) est un di-homomorphisme de Modules
gradués, relatif à cet homomorphisme d'Algèbres graduées.

Remarquons encore qu'à chacun des $\alpha_n$ il correspond
(\hyperref[0.4.4.3-fr]{0, 4.4.3}) un homomorphisme canonique de
$\mathcal{O}_X$-Modules
\[
\label{II.3.3.2.4-fr}
  \alpha_n^\sharp:p^*(\mathcal{M}_n)
  \longrightarrow\widetilde{\mathcal{M}}(n).
\tag{3.3.2.4}
\]
On vérifie sans peine que cet homomorphisme n'est autre que celui qui
correspond fonctoriellement (\hyperref[II.3.2.4-fr]{3.2.4}) à
l'homomorphisme canonique (de degré $0$) de $\mathcal{O}_Y$-Modules gradués
\[
\label{II.3.3.2.5-fr}
  \mathcal{M}_n\otimes_{\mathcal{O}_Y}\mathcal{S}
  \longrightarrow\mathcal{M}(n)
\tag{3.3.2.5}
\]

\oldpage[II]{58}

où la graduation du second membre provient naturellement de celle de
$\mathcal{S}$. On peut en effet se borner au cas où
$Y=\operatorname{Spec}(A)$ est affine, $\mathcal{M}=\widetilde{M}$ et
$\mathcal{S}=\widetilde{S}$, la $A$-algèbre graduée $S$ étant engendrée par
$S_1$, de sorte que lorsque $f$ parcourt $S_1$, les $D_+(f)$ forment un
recouvrement de $X$. Revenant aux définitions
(\hyperref[II.2.6.2-fr]{2.6.2}) on voit alors, compte tenu de
(\hyperref[I.1.6.7-fr]{I, 1.6.7}), que la restriction à $D_+(f)$ de
l'homomorphisme (\hyperref[II.3.3.2.4-fr]{3.3.2.4}) correspond
(\hyperref[I.1.3.8-fr]{I, 1.3.8}) à l'homomorphisme de $S_{(f)}$-modules
$M_n\otimes_A S_{(f)}\to(S(n))_{(f)}$ qui, à $x\otimes1$ ($x\in M_n$), fait
correspondre $x/1$~; cela démontre notre assertion.
\end{env}

\begin{proposition}[3.3.3]
\label{II.3.3.3-fr}
Pour toute section $f\in\Gamma(Y,\mathcal{S}_d)$ ($d>0$), $X_f$ est identique
à l'ensemble des points de $X$ où $\alpha_d(f)$ (considérée comme section de
$\mathcal{O}_X(d)$) ne s'annule pas
(\hyperref[0.5.5.2-fr]{0, 5.5.2}).
\end{proposition}

\begin{proof}
($\alpha_d(f)$ est une section de $p_*(\mathcal{O}_X(d))$ au-dessus de $Y$,
mais par définition une telle section est aussi une section de
$\mathcal{O}_X(d)$ au-dessus de $X$
(\hyperref[0.4.2.1-fr]{0, 4.2.1})). La définition de $X_f$
(\hyperref[II.3.1.4-fr]{3.1.4}) ramène au cas où $Y$ est affine, qui a été
traité dans (\hyperref[II.2.6.3-fr]{2.6.3}).
\end{proof}

\begin{env}[3.3.4]
\label{II.3.3.4-fr}
Nous supposerons désormais, outre l'hypothèse du début de ce numéro, que pour
tout $\mathcal{O}_X$-Module quasi-cohérent $\mathcal{F}$, les
$p_*(\mathcal{F}(n))$ sont \emph{quasi-cohérents} sur $Y$, et par suite
$\Gamma_*(\mathcal{F})=\bigoplus_{n\in\mathbf{Z}}p_*(\mathcal{F}(n))$ est
aussi un $\mathcal{O}_Y$-Module quasi-cohérent
(\hyperref[I.1.4.1-fr]{I, 1.4.1} et
\hyperref[I.1.3.9-fr]{I, 1.3.9})~; cette circonstance se produira toujours en
particulier si $X$ est \emph{de type fini} sur $Y$
(\hyperref[I.9.2.2-fr]{I, 9.2.2}). On en conclut que
$(\Gamma_*(\mathcal{F}))^{\sim}$ est défini et est un
$\mathcal{O}_X$-Module quasi-cohérent. Pour tout ouvert affine $U$ de $Y$, on
a (\hyperref[I.1.3.9-fr]{I, 1.3.9} et
\hyperref[II.2.5.4-fr]{2.5.4})
\begin{align*}
  \left(\Gamma\left(U,\bigoplus_{n\in\mathbf{Z}}
  p_*(\mathcal{F}(n))\right)\right)^{\sim}
  &=\bigoplus_{n\in\mathbf{Z}}
  \bigl(\Gamma(U,p_*(\mathcal{F}(n)))\bigr)^{\sim}\\
  &=\bigoplus_{n\in\mathbf{Z}}
  \bigl(\Gamma(p^{-1}(U),\mathcal{F}(n))\bigr)^{\sim}\\
  &=\left(\bigoplus_{n\in\mathbf{Z}}
  \Gamma(p^{-1}(U),\mathcal{F}(n))\right)^{\sim}\\
  &=\bigl(\Gamma_*(\mathcal{F}|p^{-1}(U))\bigr)^{\sim}
\end{align*}
et par suite (\hyperref[II.2.6.4-fr]{2.6.4}) on a un homomorphisme canonique
\[
  \beta_U:\left(\Gamma\left(U,\bigoplus_{n\in\mathbf{Z}}
  p_*(\mathcal{F}(n))\right)\right)^{\sim}
  \longrightarrow\mathcal{F}|p^{-1}(U).
\]

En outre, la commutativité de
(\hyperref[II.2.8.13.2-fr]{2.8.13.2}) montre que ces homomorphismes commutent
avec les opérations de restriction sur $Y$~; on en déduit donc un
homomorphisme canonique fonctoriel
\[
\label{II.3.3.4.1-fr}
  \beta:(\Gamma_*(\mathcal{F}))^{\sim}\longrightarrow\mathcal{F}
\tag{3.3.4.1}
\]
(aussi noté $\beta_{\mathcal{F}}$) pour les $\mathcal{O}_X$-Modules
quasi-cohérents.
\end{env}

\begin{proposition}[3.3.5]
\label{II.3.3.5-fr}
Soient $\mathcal{M}$ un $\mathcal{S}$-Module gradué quasi-cohérent,
$\mathcal{F}$ un $\mathcal{O}_X$-Module quasi-cohérent~; les homomorphismes
composés
\[
\label{II.3.3.5.1-fr}
  \widetilde{\mathcal{M}}
  \xrightarrow{\widetilde{\alpha}}
  (\Gamma_*(\widetilde{\mathcal{M}}))^{\sim}
  \xrightarrow{\beta}\widetilde{\mathcal{M}}
\tag{3.3.5.1}
\]
\[
\label{II.3.3.5.2-fr}
  \Gamma_*(\mathcal{F})\xrightarrow{\alpha}
  \Gamma_*((\Gamma_*(\mathcal{F}))^{\sim})
  \xrightarrow{\Gamma_*(\beta)}\Gamma_*(\mathcal{F})
\tag{3.3.5.2}
\]
sont les isomorphismes identiques.
\end{proposition}

\begin{proof}
La question est en effet locale sur $Y$, et on est ramené à
(\hyperref[II.2.6.5-fr]{2.6.5}).
\end{proof}

\subsection{Conditions de finitude.}
\label{subsection:II.3.4-fr}

\oldpage[II]{59}

\begin{proposition}[3.4.1]
\label{II.3.4.1-fr}
Soient $Y$ un préschéma, $\mathcal{S}$ une $\mathcal{O}_Y$-Algèbre
quasi-cohérente engendrée par $\mathcal{S}_1$
(\hyperref[II.3.1.9-fr]{3.1.9})~; on suppose en outre $\mathcal{S}_1$ de type
fini. Alors $X=\operatorname{Proj}(\mathcal{S})$ est de type fini au-dessus de
$Y$.
\end{proposition}

\begin{proof}
En effet, la question étant locale sur $Y$, on peut supposer $Y$ affine
d'anneau $A$~; on a alors $\mathcal{S}=\widetilde{S}$, où
$S=\Gamma(Y,\mathcal{S})$, et par hypothèse $S$ est une $A$-algèbre engendrée
par $S_1=\Gamma(Y,\mathcal{S}_1)$, où on peut en outre supposer que $S_1$ est
un $A$-module de type fini
(\hyperref[I.1.3.9-fr]{I, 1.3.9} et
\hyperref[I.1.3.12-fr]{I, 1.3.12}). Par suite $S$ est une $A$-algèbre graduée
de type fini, et on est ramené à
(\hyperref[II.2.7.1-fr]{2.7.1, (ii)}).
\end{proof}

\begin{env}[3.4.2]
\label{II.3.4.2-fr}
Soit $\mathcal{S}$ une $\mathcal{O}_Y$-Algèbre graduée quasi-cohérente~; pour
un $\mathcal{S}$-Module gradué quasi-cohérent $\mathcal{M}$, nous considérerons
les conditions de finitude suivantes~:
\begin{enumerate}
  \item[(\textbf{TF})] \emph{Il existe un entier $n$ tel que le
    $\mathcal{S}$-Module $\bigoplus_{k\geq n}\mathcal{M}_k$ soit de type fini.}
  \item[(\textbf{TN})] \emph{Il existe un entier $n$ tel que
    $\mathcal{M}_k=0$ pour $k\geq n$.}
\end{enumerate}
Si $\mathcal{M}$ vérifie (\textbf{TN}), on a
$\widetilde{\mathcal{M}}=0$, la question étant locale sur $Y$
(\hyperref[II.2.7.2-fr]{2.7.2}).

Soient $\mathcal{M}$, $\mathcal{N}$ deux $\mathcal{S}$-Modules gradués
quasi-cohérents~; on dit qu'un homomorphisme
$u:\mathcal{M}\to\mathcal{N}$ de degré $0$ est
(\textbf{TN})-\emph{injectif} (resp. (\textbf{TN})-\emph{surjectif},
(\textbf{TN})-\emph{bijectif}) s'il existe un entier $n$ tel que
$u_k:\mathcal{M}_k\to\mathcal{N}_k$ soit injectif (resp. surjectif, bijectif)
pour $k\geq n$~; alors
$\widetilde{u}:\widetilde{\mathcal{M}}\to\widetilde{\mathcal{N}}$ est injectif
(resp. surjectif, bijectif) en vertu de
(\hyperref[II.2.7.2-fr]{2.7.2}), la question étant locale sur $Y$, et compte
tenu de (\hyperref[I.1.3.9-fr]{I, 1.3.9})~; lorsque $u$ est
(\textbf{TN})-bijectif, on dit aussi que $u$ est un
(\textbf{TN})-\emph{isomorphisme}.
\end{env}

\begin{proposition}[3.4.3]
\label{II.3.4.3-fr}
Soient $Y$ un préschéma, $\mathcal{S}$ une $\mathcal{O}_Y$-Algèbre graduée
quasi-cohérente engendrée par $\mathcal{S}_1$, $\mathcal{S}_1$ étant supposé
de type fini. Soit $\mathcal{M}$ un $\mathcal{S}$-Module gradué
quasi-cohérent.
\begin{enumerate}
  \item[\rm(i)] Si $\mathcal{M}$ vérifie la condition (\textbf{TF}),
    $\widetilde{\mathcal{M}}$ est de type fini.
  \item[\rm(ii)] Supposons que $\mathcal{M}$ vérifie (\textbf{TF})~; pour que
    $\widetilde{\mathcal{M}}=0$, il faut et il suffit que $\mathcal{M}$
    vérifie (\textbf{TN}).
\end{enumerate}
\end{proposition}

\begin{proof}
Les questions étant locales sur $Y$, on est ramené au cas où $Y$ est affine
d'anneau $A$, $\mathcal{S}=\widetilde{S}$, où $S$ est une $A$-algèbre graduée
telle que l'idéal $S_+$ soit de type fini,
$\mathcal{M}=\widetilde{M}$ où $M$ est un $S$-module gradué~; la proposition
résulte alors de (\hyperref[II.2.7.3-fr]{2.7.3}).
\end{proof}

\begin{theorem}[3.4.4]
\label{II.3.4.4-fr}
Soient $Y$ un préschéma, $\mathcal{S}$ une $\mathcal{O}_Y$-Algèbre graduée
quasi-cohérente engendrée par $\mathcal{S}_1$, $\mathcal{S}_1$ étant supposé
de type fini~; soit $X=\operatorname{Proj}(\mathcal{S})$. Pour tout
$\mathcal{O}_X$-Module quasi-cohérent $\mathcal{F}$, l'homomorphisme canonique
$\beta$ (\hyperref[II.3.3.4-fr]{3.3.4}) est un isomorphisme.
\end{theorem}

\begin{proof}
Notons d'abord que $\beta$ est défini en vertu de
(\hyperref[II.3.4.1-fr]{3.4.1}). Pour voir que $\beta$ est un isomorphisme,
on est ramené au cas où $Y$ est affine d'anneau $A$,
$\mathcal{S}=\widetilde{S}$, où $S$ est une $A$-algèbre graduée engendrée par
$S_1$ et $S_1$ est un $A$-module de type fini. Il suffit alors d'appliquer
(\hyperref[II.2.7.5-fr]{2.7.5}).
\end{proof}

\begin{corollary}[3.4.5]
\label{II.3.4.5-fr}
Sous les hypothèses de (\hyperref[II.3.4.4-fr]{3.4.4}), tout
$\mathcal{O}_X$-Module quasi-cohérent $\mathcal{F}$ est isomorphe à un
$\mathcal{O}_X$-Module de la forme $\widetilde{\mathcal{M}}$, où
$\mathcal{M}$ est un $\mathcal{S}$-Module gradué quasi-cohérent. Si en outre
$\mathcal{F}$ est de type fini, et si on suppose que $Y$ est un schéma
quasi-compact, ou que l'espace sous-jacent à $Y$ est noethérien, on peut
supposer $\mathcal{M}$ de type fini.
\end{corollary}

\oldpage[II]{60}

\begin{proof}
La première assertion résulte aussitôt de
(\hyperref[II.3.4.4-fr]{3.4.4}) en prenant
$\mathcal{M}=\Gamma_*(\mathcal{F})$. Pour établir la seconde, il suffira
d'établir que $\mathcal{M}$ est limite inductive de ses sous-$\mathcal{S}$-Modules
gradués de type fini $\mathcal{N}_\lambda$~: en effet, il en résultera que
$\widetilde{\mathcal{M}}$ est limite inductive des
$\widetilde{\mathcal{N}}_\lambda$
(\hyperref[II.3.2.4-fr]{3.2.4}), donc $\mathcal{F}$ est limite inductive des
$\beta(\widetilde{\mathcal{N}}_\lambda)$~; comme $X$ est quasi-compact
((\hyperref[II.3.4.1-fr]{3.4.1}) et
\hyperref[I.6.3.1-fr]{I, 6.3.1}) et que $\mathcal{F}$ est de type fini,
$\mathcal{F}$ sera nécessairement égal à un des
$\beta(\widetilde{\mathcal{N}}_\lambda)$
(\hyperref[0.5.2.3-fr]{0, 5.2.3}).

Pour définir les $\mathcal{N}_\lambda$ ayant $\mathcal{M}$ pour limite
inductive, il suffit de considérer pour chaque $n\in\mathbf{Z}$ le
$\mathcal{O}_Y$-Module quasi-cohérent $\mathcal{M}_n$, qui est limite
inductive de ses sous-$\mathcal{O}_Y$-Modules
$\mathcal{M}_n^{(\mu_n)}$ de type fini, en raison des hypothèses sur $Y$
(\hyperref[I.9.4.9-fr]{I, 9.4.9})~; il est immédiat que
$\mathcal{P}_{\mu_n}=\mathcal{S}\cdot\mathcal{M}_n^{(\mu_n)}$ est un
$\mathcal{S}$-Module gradué de type fini, et on vérifie aussitôt qu'en prenant
pour $\mathcal{N}_\lambda$ les sommes finies de $\mathcal{S}$-Modules de la
forme $\mathcal{P}_{\mu_n}$, on répond bien à la question.
\end{proof}

\begin{corollary}[3.4.6]
\label{II.3.4.6-fr}
Supposons vérifiées les hypothèses de (\hyperref[II.3.4.4-fr]{3.4.4}), et en
outre que l'espace sous-jacent $Y$ soit quasi-compact~; soit $\mathcal{F}$ un
$\mathcal{O}_X$-Module quasi-cohérent de type fini. Il existe $n_0$ tel que
pour $n\geq n_0$, l'homomorphisme canonique
$\sigma:p^*(p_*(\mathcal{F}(n)))\to\mathcal{F}(n)$
(\hyperref[0.4.4.3-fr]{0, 4.4.3}) soit surjectif.
\end{corollary}

\begin{proof}
En effet, pour tout $y\in Y$, soit $U$ un voisinage ouvert affine de $y$ dans
$Y$. Il existe un entier $n_0(U)$ tel que, pour $n\geq n_0(U)$,
$\mathcal{F}(n)|p^{-1}(U)$ soit engendré par un nombre fini de ses sections
au-dessus de $p^{-1}(U)$ (\hyperref[II.2.7.9-fr]{2.7.9})~; mais ces dernières
sont images canoniques de sections de $p^*(p_*(\mathcal{F}(n)))$ au-dessus de
$p^{-1}(U)$ (\hyperref[0.3.7.1-fr]{0, 3.7.1} et
\hyperref[0.4.4.3-fr]{0, 4.4.3}), donc
$\mathcal{F}(n)|p^{-1}(U)$ est égal à l'image canonique de
$p^*(p_*(\mathcal{F}(n)))|p^{-1}(U)$. Enfin, comme $Y$ est quasi-compact, il
existe un recouvrement fini de $Y$ par des ouverts affines $U_i$, et en
prenant pour $n_0$ le plus grand des $n_0(U_i)$, on achève la démonstration.
\end{proof}

\par\medskip\noindent\textit{Remarques (3.4.7). ---\ }
\label{II.3.4.7-fr}
Si $p=(\psi,\theta):X\to Y$ est un morphisme d'espaces annelés et
$\mathcal{F}$ un $\mathcal{O}_X$-Module, le fait que l'homomorphisme canonique
$\sigma:p^*(p_*(\mathcal{F}))\to\mathcal{F}$ soit surjectif s'explicite de la
façon suivante (\hyperref[0.4.4.1-fr]{0, 4.4.1})~: pour tout $x\in X$ et toute
section $s$ de $\mathcal{F}$ au-dessus d'un voisinage ouvert $V$ de $x$, il
existe un voisinage ouvert $U$ de $p(x)$ dans $Y$, un nombre fini de sections
$t_i$ ($1\leq i\leq m$) de $\mathcal{F}$ au-dessus de $p^{-1}(U)$, un
voisinage $W\subset V\cap p^{-1}(U)$ de $x$ et des sections $a_i$
($1\leq i\leq m$) de $\mathcal{O}_X$ au-dessus de $W$ telles que
\[
  s|W=\sum_i a_i\cdot(t_i|W).
\]
Lorsque $Y$ est un \emph{schéma affine} et $p_*(\mathcal{F})$
\emph{quasi-cohérent}, cette condition équivaut au fait que $\mathcal{F}$ est
\emph{engendré par ses sections au-dessus de $X$}
(\hyperref[0.5.5.1-fr]{0, 5.5.1})~: en effet, si
$Y=\operatorname{Spec}(A)$, on peut supposer que $U=D(f)$ avec $f\in A$~; il
existe un entier $n>0$ et des sections $s_i$ de $\mathcal{F}$ au-dessus de
$X$ telles que $t_i$ soit la restriction à $p^{-1}(U)$ de $s_ig^n$, avec
$g=\theta(f)$ (en appliquant (\hyperref[I.1.4.1-fr]{I, 1.4.1}) à
$p_*(\mathcal{F})$)~; comme $g$ est inversible au-dessus de $p^{-1}(U)$, on a
donc
\[
  s|W=\sum_i b_i\cdot(s_i|W)
\]
avec $b_i=a_i(g|W)^{-n}$, d'où notre assertion. Lorsque $Y$ est affine, le
corollaire (\hyperref[II.3.4.6-fr]{3.4.6}) redonne donc
(\hyperref[II.2.7.9-fr]{2.7.9}).

\oldpage[II]{61}

On en conclut que lorsque $Y$ est un préschéma quelconque, les trois
conditions suivantes, pour un $\mathcal{O}_X$-Module quasi-cohérent
$\mathcal{F}$ tel que $p_*(\mathcal{F})$ soit \emph{quasi-cohérent}, sont
équivalentes~:
\begin{enumerate}
  \item[a)] \emph{L'homomorphisme canonique
    $\sigma:p^*(p_*(\mathcal{F}))\to\mathcal{F}$ est surjectif.}
  \item[b)] \emph{Il existe un $\mathcal{O}_Y$-Module quasi-cohérent
    $\mathcal{G}$ et un homomorphisme surjectif
    $p^*(\mathcal{G})\to\mathcal{F}$.}
  \item[c)] \emph{Pour tout ouvert affine $U$ de $Y$,
    $\mathcal{F}|p^{-1}(U)$ est engendré par ses sections au-dessus de
    $p^{-1}(U)$.}
\end{enumerate}
On vient en effet de prouver l'équivalence de \emph{a)} et \emph{c)}. D'autre
part, il est clair que \emph{a)} entraîne \emph{b)}, $p_*(\mathcal{F})$ étant
par hypothèse quasi-cohérent. Inversement, tout homomorphisme
$u:p^*(\mathcal{G})\to\mathcal{F}$ se factorise en
$p^*(\mathcal{G})\to p^*(p_*(\mathcal{F}))\xrightarrow{\sigma}\mathcal{F}$
(\hyperref[0.3.5.4.4-fr]{0, 3.5.4.4}), donc si $u$ est surjectif il en est de
même de $\sigma$, ce qui prouve que \emph{b)} entraîne \emph{a)}.

\begin{corollary}[3.4.8]
\label{II.3.4.8-fr}
Supposons vérifiées les hypothèses de (\hyperref[II.3.4.4-fr]{3.4.4}) et
supposons en outre que $Y$ soit un schéma quasi-compact ou que l'espace
sous-jacent à $Y$ soit noethérien. Soit $\mathcal{F}$ un
$\mathcal{O}_X$-Module quasi-cohérent de type fini~; il existe alors un entier
$n_0$ tel que pour $n\geq n_0$, $\mathcal{F}$ soit isomorphe à un quotient
d'un $\mathcal{O}_X$-Module de la forme $(p^*(\mathcal{G}))(-n)$ où
$\mathcal{G}$ est un $\mathcal{O}_Y$-Module quasi-cohérent de type fini
(dépendant de $n$).
\end{corollary}

\begin{proof}
Comme le morphisme structural $X\to Y$ est séparé et de type fini,
$p_*(\mathcal{F}(n))$ est quasi-cohérent
(\hyperref[I.9.2.2-fr]{I, 9.2.2, b)}), donc limite inductive de ses
sous-$\mathcal{O}_Y$-Modules quasi-cohérents de type fini, en vertu de
l'hypothèse sur $Y$ (\hyperref[I.9.4.9-fr]{I, 9.4.9}). On en déduit par
(\hyperref[II.3.4.6-fr]{3.4.6}),
(\hyperref[0.4.3.2-fr]{0, 4.3.2}) et
(\hyperref[0.5.2.3-fr]{0, 5.2.3}) que $\mathcal{F}(n)$ est l'image canonique
d'un $\mathcal{O}_X$-Module de la forme $p^*(\mathcal{G})$, où $\mathcal{G}$
est un sous-$\mathcal{O}_Y$-Module quasi-cohérent de type fini de
$p_*(\mathcal{F}(n))$~; le corollaire résulte alors de
(\hyperref[II.3.2.5.2-fr]{3.2.5.2}) et
(\hyperref[II.3.2.7.1-fr]{3.2.7.1}).
\end{proof}

\subsection{Comportements fonctoriels.}
\label{subsection:II.3.5-fr}

\begin{env}[3.5.1]
\label{II.3.5.1-fr}
Soient $Y$ un préschéma, $\mathcal{S}$, $\mathcal{S}'$ deux
$\mathcal{O}_Y$-Algèbres graduées quasi-cohérentes à degrés positifs~; posons
$X=\operatorname{Proj}(\mathcal{S})$, $X'=\operatorname{Proj}(\mathcal{S}')$,
et soient $p$, $p'$ les morphismes structuraux de $X$ et $X'$ dans $Y$. Soit
$\varphi:\mathcal{S}'\to\mathcal{S}$ un $\mathcal{O}_Y$-homomorphisme
d'Algèbres graduées. Pour tout ouvert affine $U$ de $Y$, posons
$S_U=\Gamma(U,\mathcal{S})$, $S'_U=\Gamma(U,\mathcal{S}')$~;
l'homomorphisme $\varphi$ définit un homomorphisme
$\varphi_U:S'_U\to S_U$ de $A_U$-algèbres graduées, en posant
$A_U=\Gamma(U,\mathcal{O}_Y)$. Il lui correspond dans $p^{-1}(U)$ un ensemble
ouvert $G(\varphi_U)$ et un morphisme
$\Phi_U:G(\varphi_U)\to p'^{-1}(U)$
(\hyperref[II.2.8.1-fr]{2.8.1}). En outre, si $V\subset U$ est un ouvert
affine, le diagramme
\[
\label{II.3.5.1.1-fr}
  \xymatrix{
    S'_U\ar[r]^{\varphi_U}\ar[d]
    & S_U\ar[d]\\
    S'_V\ar[r]_{\varphi_V}
    & S_V
  }
\tag{3.5.1.1}
\]
est commutatif, et on vérifie aussitôt, à partir des définitions
(\hyperref[II.2.8.1-fr]{2.8.1}), que l'on a
\[
  G(\varphi_V)=G(\varphi_U)\cap p^{-1}(V)
\]
et que $\Phi_V$ est la restriction de $\Phi_U$ à $G(\varphi_V)$. On a ainsi
défini une partie ouverte $G(\varphi)$ de $X$ telle que
$G(\varphi)\cap p^{-1}(U)=G(\varphi_U)$ pour tout ouvert affine $U\subset Y$,
et un $Y$-morphisme

\oldpage[II]{62}

affine $\Phi:G(\varphi)\to X'$, qui est dit \emph{associé} à $\varphi$ et que
nous noterons $\operatorname{Proj}(\varphi)$. Lorsque, pour tout $y\in Y$, il
y a un voisinage affine $U$ de $y$ tel que le
$\Gamma(U,\mathcal{O}_Y)$-module $\Gamma(U,\mathcal{S}_+)$ soit engendré par
$\varphi(\Gamma(U,\mathcal{S}'_+))$, on a
$G(\varphi_U)=p^{-1}(U)$, et par suite $G(\varphi)=X$.
\end{env}

\begin{proposition}[3.5.2]
\label{II.3.5.2-fr}
\begin{enumerate}
  \item[\rm(i)] Si $\mathcal{M}$ est un $\mathcal{S}$-Module gradué
    quasi-cohérent, il existe un isomorphisme canonique fonctoriel du
    $\mathcal{O}_{X'}$-Module $(\mathcal{M}_{[\varphi]})^{\sim}$ sur le
    $\mathcal{O}_{X'}$-Module
    $\Phi_*(\widetilde{\mathcal{M}}|G(\varphi))$.
  \item[\rm(ii)] Si $\mathcal{M}'$ est un $\mathcal{S}'$-Module gradué
    quasi-cohérent, il existe un homomorphisme canonique fonctoriel $\nu$ du
    $(\mathcal{O}_X|G(\varphi))$-Module
    $\Phi^*(\widetilde{\mathcal{M}'})$ dans le
    $(\mathcal{O}_X|G(\varphi))$-Module
    $(\mathcal{M}'\otimes_{\mathcal{S}'}\mathcal{S})^{\sim}|G(\varphi)$. Si
    $\mathcal{S}'$ est engendrée par $\mathcal{S}'_1$, $\nu$ est un
    isomorphisme.
\end{enumerate}
\end{proposition}

\begin{proof}
Les homomorphismes considérés ont en effet déjà été définis lorsque $Y$ est
affine (\hyperref[II.2.8.7-fr]{2.8.7} et
\hyperref[II.2.8.8-fr]{2.8.8}), et dans le cas général il suffit de vérifier
qu'ils sont compatibles avec les opérations de restriction d'un ouvert
affine de $Y$ à un ouvert plus petit, ce qui résulte aussitôt de la
commutativité de (\hyperref[II.3.5.1.1-fr]{3.5.1.1}).
\end{proof}

En particulier, pour tout $n\in\mathbf{Z}$, on a un homomorphisme canonique
\[
\label{II.3.5.2.1-fr}
  \Phi^*(\mathcal{O}_{X'}(n))\longrightarrow
  \mathcal{O}_X(n)|G(\varphi).
\tag{3.5.2.1}
\]

\begin{proposition}[3.5.3]
\label{II.3.5.3-fr}
Soient $Y$, $Y'$ deux préschémas, $\psi:Y'\to Y$ un morphisme,
$\mathcal{S}$ une $\mathcal{O}_Y$-Algèbre graduée quasi-cohérente, et posons
$\mathcal{S}'=\psi^*(\mathcal{S})$. Alors le $Y'$-schéma
$X'=\operatorname{Proj}(\mathcal{S}')$ s'identifie canoniquement à
$\operatorname{Proj}(\mathcal{S})\times_Y Y'$. En outre, si $\mathcal{M}$ est
un $\mathcal{S}$-Module gradué quasi-cohérent, le $\mathcal{O}_{X'}$-Module
$(\psi^*(\mathcal{M}))^{\sim}$ s'identifie à
$\widetilde{\mathcal{M}}\otimes_Y\mathcal{O}_{Y'}$.
\end{proposition}

\begin{proof}
Notons en premier lieu que $\psi^*(\mathcal{S})$ et
$\psi^*(\mathcal{M})$ sont des $\mathcal{O}_{Y'}$-Modules quasi-cohérents
ainsi que leurs composants homogènes
(\hyperref[0.5.1.4-fr]{0, 5.1.4}). Soient $U$ un ouvert affine de $Y$,
$U'\subset\psi^{-1}(U)$ un ouvert affine de $Y'$, $A$, $A'$ les anneaux de
$U$ et $U'$ respectivement~; on a alors
$\mathcal{S}|U=\widetilde{S}$, où $S$ est une $A$-algèbre graduée, et
$\mathcal{S}'|U'$ s'identifie à $(S\otimes_A A')^{\sim}$
(\hyperref[I.1.6.5-fr]{I, 1.6.5})~; la première assertion résulte alors de
(\hyperref[II.2.8.10-fr]{2.8.10}) et de
(\hyperref[I.3.2.6.2-fr]{I, 3.2.6.2}), car on vérifie aussitôt que la
projection
$\operatorname{Proj}(\mathcal{S}'|U')\to
\operatorname{Proj}(\mathcal{S}|U)$ définie par l'identification précédente
est compatible avec les opérations de restrictions sur $U$ et $U'$ et
définit donc bien un morphisme
$\operatorname{Proj}(\mathcal{S}')\to\operatorname{Proj}(\mathcal{S})$.
Soient maintenant
\[
  q:\operatorname{Proj}(\mathcal{S})\to Y,
  \qquad q':\operatorname{Proj}(\mathcal{S}')\to Y'
\]
les morphismes structuraux~; $q'^{-1}(U')$ s'identifie alors à
$q^{-1}(U)\times_U U'$, et les deux faisceaux
$(\psi^*(\mathcal{M}))^{\sim}|q'^{-1}(U')$ et
$(\widetilde{\mathcal{M}}\otimes_Y\mathcal{O}_{Y'})|q'^{-1}(U')$
s'identifient alors canoniquement tous deux à $(M\otimes_A A')^{\sim}$, où
on a posé $M=\Gamma(U,\mathcal{M})$, en vertu de
(\hyperref[II.2.8.10-fr]{2.8.10}) et de
(\hyperref[I.1.6.5-fr]{I, 1.6.5})~; d'où la seconde assertion, car on vérifie
encore de façon immédiate la compatibilité des identifications précédentes
avec les opérations de restriction.
\end{proof}

\begin{corollary}[3.5.4]
\label{II.3.5.4-fr}
Avec les notations de (\hyperref[II.3.5.3-fr]{3.5.3}),
$\mathcal{O}_{X'}(n)$ s'identifie canoniquement à
$\mathcal{O}_X(n)\otimes_Y\mathcal{O}_{Y'}$ pour tout $n\in\mathbf{Z}$ (avec
$X=\operatorname{Proj}(\mathcal{S})$).
\end{corollary}

\begin{proof}
En effet, avec les notations de (\hyperref[II.3.5.3-fr]{3.5.3}), il est clair
que $\psi^*(\mathcal{S}(n))=\mathcal{S}'(n)$ pour tout $n\in\mathbf{Z}$.
\end{proof}

\begin{env}[3.5.5]
\label{II.3.5.5-fr}
Gardant les notations précédentes, désignons par $\Psi$ la projection
canonique $X'\to X$, et posons $\mathcal{M}'=\psi^*(\mathcal{M})$~; nous
supposerons en outre que $\mathcal{S}$ est engendrée par $\mathcal{S}_1$

\oldpage[II]{63}

et que $X$ est de type fini sur $Y$~; il en résulte alors que $\mathcal{S}'$
est engendrée par $\mathcal{S}'_1$ (comme on le voit en se ramenant au cas où
$Y$ et $Y'$ sont affines) et que $X'$ est de type fini sur $Y'$
(\hyperref[I.6.3.4-fr]{I, 6.3.4}). Soit $\mathcal{F}$ un
$\mathcal{O}_X$-Module et posons $\mathcal{F}'=\Psi^*(\mathcal{F})$~; il
résulte alors de (\hyperref[II.3.5.4-fr]{3.5.4}) et de
(\hyperref[0.4.3.3-fr]{0, 4.3.3}) que l'on a
$\mathcal{F}'(n)=\Psi^*(\mathcal{F}(n))$ pour tout $n\in\mathbf{Z}$. On
définit en outre un $\psi$-homomorphisme canonique
$\theta_n:q_*(\mathcal{F}(n))\to q'_*(\mathcal{F}'(n))$ de la façon
suivante~: vu la commutativité du diagramme
\[
  \xymatrix{
    X\ar[d]_{q}
    & X'\ar[l]_{\Psi}\ar[d]^{q'}\\
    Y
    & Y'\ar[l]^{\psi}
  }
\]
il s'agit de définir un homomorphisme
$q_*(\mathcal{F}(n))\to
\psi_*(q'_*(\Psi^*(\mathcal{F}(n))))
=q_*(\Psi_*(\Psi^*(\mathcal{F}(n))))$, et il suffit de prendre
l'homomorphisme $\theta_n=q_*(\rho_n)$, où $\rho_n$ est l'homomorphisme
canonique
$\mathcal{F}(n)\to\Psi_*(\Psi^*(\mathcal{F}(n)))$
(\hyperref[0.4.4.3-fr]{0, 4.4.3}). Il est immédiat que pour tout ouvert affine
$U$ de $Y$ et tout ouvert affine $U'$ de $Y'$ tel que
$U'\subset\psi^{-1}(U)$, l'homomorphisme $\theta_n$ donne pour les sections
l'homomorphisme canonique (\hyperref[0.3.7.2-fr]{0, 3.7.2})
$\Gamma(q^{-1}(U),\mathcal{F}(n))\to
\Gamma(q'^{-1}(U'),\mathcal{F}'(n))$.

La commutativité de (\hyperref[II.2.8.13.2-fr]{2.8.13.2}) montre alors que si
$\mathcal{F}$ est quasi-cohérent, le diagramme
\[
  \xymatrix{
    \mathcal{F}\ar[r]^{\rho}
    & \mathcal{F}'\\
    (\Gamma_*(\mathcal{F}))^{\sim}
      \ar[u]^{\beta_{\mathcal{F}}}\ar[r]_{\widetilde{\theta}}
    & (\Gamma_*(\mathcal{F}'))^{\sim}
      \ar[u]_{\beta_{\mathcal{F}'}}
  }
\]
est commutatif (la flèche horizontale du haut étant le $\Psi$-morphisme
canonique $\mathcal{F}\to\Psi^*(\mathcal{F})$).

De même, la commutativité de
(\hyperref[II.2.8.13.1-fr]{2.8.13.1}) montre que le diagramme
\[
  \xymatrix{
    \Gamma_*(\widetilde{\mathcal{M}})\ar[r]^{\theta}
    & \Gamma_*(\widetilde{\mathcal{M}'})\\
    \mathcal{M}\ar[r]_{\rho}\ar[u]^{\alpha_{\mathcal{M}}}
    & \mathcal{M}'\ar[u]_{\alpha_{\mathcal{M}'}}
  }
\]
est commutatif (la flèche horizontale du bas étant le $\psi$-morphisme
canonique $\mathcal{M}\to\psi^*(\mathcal{M})$).
\end{env}

\begin{env}[3.5.6]
\label{II.3.5.6-fr}
Considérons maintenant deux préschémas $Y$, $Y'$, un morphisme
$g:Y'\to Y$, une $\mathcal{O}_Y$-Algèbre (resp.
$\mathcal{O}_{Y'}$-Algèbre) graduée quasi-cohérente $\mathcal{S}$ (resp.
$\mathcal{S}'$) et un $g$-morphisme d'Algèbres graduées
$u:\mathcal{S}\to\mathcal{S}'$, c'est-à-dire un
$\mathcal{O}_Y$-homomorphisme d'Algèbres graduées
$\mathcal{S}\to g_*(\mathcal{S}')$~; on sait d'ailleurs qu'il revient au même
de se donner un $\mathcal{O}_{Y'}$-homomorphisme d'Algèbres graduées
$u^\sharp:g^*(\mathcal{S})\to\mathcal{S}'$. On déduit alors canoniquement de
$u^\sharp$ un $Y'$-morphisme
$W=\operatorname{Proj}(u^\sharp):G(u^\sharp)\to
\operatorname{Proj}(g^*(\mathcal{S}))$ où $G(u^\sharp)$ est un ouvert de
$X'=\operatorname{Proj}(\mathcal{S}')$
(\hyperref[II.3.5.1-fr]{3.5.1}). D'autre part,
$X''=\operatorname{Proj}(g^*(\mathcal{S}))$ s'identifie

\oldpage[II]{64}

canoniquement à $X\times_Y Y'$, en posant
$X=\operatorname{Proj}(\mathcal{S})$
(\hyperref[II.3.5.3-fr]{3.5.3})~; composant avec
$\operatorname{Proj}(u^\sharp)$ la première projection
$p:X\times_Y Y'\to X$, on obtient donc un morphisme
$v:G(u^\sharp)\to X$, que nous désignerons par $\operatorname{Proj}(u)$, et
qui est tel que le diagramme
\[
  \xymatrix{
    G(u^\sharp)\ar[r]^{v}\ar[d]
    & X\ar[d]\\
    Y'\ar[r]_{g}
    & Y
  }
\]
soit commutatif.

En outre, pour tout $\mathcal{S}$-Module gradué quasi-cohérent $\mathcal{M}$,
on a un $v$-morphisme canonique
\[
\label{II.3.5.6.1-fr}
  \nu:\widetilde{\mathcal{M}}\longrightarrow
  (g^*(\mathcal{M})\otimes_{g^*(\mathcal{S})}\mathcal{S}')^{\sim}
  |G(u^\sharp).
\tag{3.5.6.1}
\]

En effet, $\nu^\sharp$ s'obtient en composant les homomorphismes
\[
  v^*(\widetilde{\mathcal{M}})
  =W^*(p^*(\widetilde{\mathcal{M}}))
  \longrightarrow W^*((g^*(\mathcal{M}))^{\sim})
  \longrightarrow
  (g^*(\mathcal{M})\otimes_{g^*(\mathcal{S})}\mathcal{S}')^{\sim}
  |G(u^\sharp)
\]
où la première flèche provient de l'isomorphisme
(\hyperref[II.3.5.3-fr]{3.5.3}) et la seconde est l'homomorphisme
(\hyperref[II.3.5.2-fr]{3.5.2, (ii)})~; lorsque $\mathcal{S}$ est engendrée
par $\mathcal{S}_1$, il résulte de (\hyperref[II.3.5.2-fr]{3.5.2}) que
$\nu^\sharp$ est un isomorphisme.

Comme cas particulier de (\hyperref[II.3.5.6.1-fr]{3.5.6.1}), on a, pour
tout $n\in\mathbf{Z}$, un $v$-morphisme canonique
\[
\label{II.3.5.6.2-fr}
  \nu:\mathcal{O}_X(n)\longrightarrow
  \mathcal{O}_{X'}(n)|G(u^\sharp).
\tag{3.5.6.2}
\]
\end{env}

\subsection{Sous-préschémas fermés d'un préschéma
$\operatorname{Proj}(\mathcal{S})$.}
\label{subsection:II.3.6-fr}

\begin{env}[3.6.1]
\label{II.3.6.1-fr}
Soient $Y$ un préschéma,
$\varphi:\mathcal{S}\to\mathcal{S}'$ un homomorphisme de
$\mathcal{O}_Y$-Algèbres graduées quasi-cohérentes, de degré $0$. On dit que
$\varphi$ est (\textbf{TN})-\emph{surjectif} (resp.
(\textbf{TN})-\emph{injectif}, (\textbf{TN})-\emph{bijectif}) s'il existe $n$
tel que, pour $k\geq n$, $\varphi_k:\mathcal{S}_k\to\mathcal{S}'_k$ soit
surjectif (resp. injectif, bijectif). Lorsqu'il en est ainsi on ramène l'étude
du morphisme
$\Phi:\operatorname{Proj}(\mathcal{S}')\to\operatorname{Proj}(\mathcal{S})$
correspondant au cas où $\varphi$ est surjectif (resp. injectif, bijectif).
Cela se démontre comme dans (\hyperref[II.2.9.1-fr]{2.9.1}) (qui en est le cas
particulier correspondant à $Y$ affine) en utilisant
(\hyperref[II.3.1.8-fr]{3.1.8}). Au lieu de dire que $\varphi$ est
(\textbf{TN})-bijectif, on dit aussi que c'est alors un
(\textbf{TN})-isomorphisme.
\end{env}

\begin{proposition}[3.6.2]
\label{II.3.6.2-fr}
Soient $Y$ un préschéma, $\mathcal{S}$ une $\mathcal{O}_Y$-Algèbre graduée
quasi-cohérente, et soit $X=\operatorname{Proj}(\mathcal{S})$.
\begin{enumerate}
  \item[\rm(i)] Si $\varphi:\mathcal{S}\to\mathcal{S}'$ est un
    homomorphisme (\textbf{TN})-surjectif de $\mathcal{O}_Y$-Algèbres
    graduées, le morphisme correspondant
    $\Phi=\operatorname{Proj}(\varphi)$
    (\hyperref[II.3.5.1-fr]{3.5.1}) est défini dans
    $\operatorname{Proj}(\mathcal{S}')$ tout entier et est une immersion
    fermée de $\operatorname{Proj}(\mathcal{S}')$ dans $X$. Si $\mathcal{J}$
    est le noyau de $\varphi$, le sous-préschéma fermé de $X$ associé à
    $\Phi$ est défini par le faisceau d'idéaux quasi-cohérent
    $\widetilde{\mathcal{J}}$ de $\mathcal{O}_X$.
  \item[\rm(ii)] Supposons en outre $\mathcal{S}_0=\mathcal{O}_Y$,
    $\mathcal{S}$ engendrée par $\mathcal{S}_1$ et $\mathcal{S}_1$ de type
    fini. Soit $X'$ un sous-préschéma fermé de
    $X=\operatorname{Proj}(\mathcal{S})$, défini par un faisceau
    quasi-cohérent $\mathcal{I}$ d'idéaux de $\mathcal{O}_X$. Soit
    $\mathcal{J}$

\oldpage[II]{65}

    le faisceau gradué quasi-cohérent d'idéaux de $\mathcal{S}$, image
    réciproque de $\Gamma_*(\mathcal{I})$ par l'homomorphisme canonique
    $\alpha:\mathcal{S}\to\Gamma_*(\mathcal{O}_X)$
    (\hyperref[II.3.3.2-fr]{3.3.2}), et posons
    $\mathcal{S}'=\mathcal{S}/\mathcal{J}$. Alors $X'$ est le sous-préschéma
    associé (\hyperref[I.4.2.1-fr]{I, 4.2.1}) à l'immersion fermée
    $\operatorname{Proj}(\mathcal{S}')\to X$ correspondant à l'homomorphisme
    canonique $\mathcal{S}\to\mathcal{S}'$ de $\mathcal{O}_Y$-Algèbres
    graduées.
\end{enumerate}
\end{proposition}

\begin{proof}
\begin{enumerate}
  \item[\rm(i)] On peut supposer $\varphi$ surjectif
    (\hyperref[II.3.6.1-fr]{3.6.1}). Alors, pour tout ouvert affine $U$ de
    $Y$, $\Gamma(U,\mathcal{S})\to\Gamma(U,\mathcal{S}')$ est surjectif
    (\hyperref[I.1.3.9-fr]{I, 1.3.9}), donc
    (\hyperref[II.3.5.1-fr]{3.5.1}) on a $G(\varphi)=X$. On est
    immédiatement ramené à démontrer la proposition lorsque $Y$ est affine,
    et alors elle découle de (\hyperref[II.2.9.2-fr]{2.9.2, (i)}).
  \item[\rm(ii)] On est ramené à prouver que l'homomorphisme
    $\widetilde{\mathcal{J}}\to\mathcal{O}_X$ déduit de l'injection canonique
    $\mathcal{J}\to\mathcal{S}$ est un isomorphisme de
    $\widetilde{\mathcal{J}}$ sur $\mathcal{I}$~; comme la question est locale
    sur $Y$, on peut supposer $Y$ affine d'anneau $A$, ce qui entraîne
    $\mathcal{S}=\widetilde{S}$, où $S$ est une $A$-algèbre graduée engendrée
    par $S_1$, $S_1$ étant de type fini sur $A$. Il suffit alors d'appliquer
    (\hyperref[II.2.9.2-fr]{2.9.2, (ii)}).
\end{enumerate}
\end{proof}

\begin{corollary}[3.6.3]
\label{II.3.6.3-fr}
Sous les conditions de (\hyperref[II.3.6.2-fr]{3.6.2, (i)}), et en supposant
de plus $\mathcal{S}$ engendrée par $\mathcal{S}_1$,
$\Phi^*(\mathcal{O}_X(n))$ s'identifie canoniquement à
$\mathcal{O}_{X'}(n)$ pour tout $n\in\mathbf{Z}$.
\end{corollary}

\begin{proof}
On a défini un tel isomorphisme canonique lorsque $Y$ est affine
(\hyperref[II.2.9.3-fr]{2.9.3})~; dans le cas général, il suffit de vérifier
que les isomorphismes ainsi définis pour chaque ouvert affine $U$ de $Y$ sont
compatibles avec le passage de $U$ à un ouvert affine $U'\subset U$, ce qui
est immédiat.
\end{proof}

\begin{corollary}[3.6.4]
\label{II.3.6.4-fr}
Soient $Y$ un préschéma, $\mathcal{S}$ une $\mathcal{O}_Y$-Algèbre graduée
quasi-cohérente engendrée par $\mathcal{S}_1$, $\mathcal{M}$ un
$\mathcal{O}_Y$-Module quasi-cohérent, $u$ un $\mathcal{O}_Y$-homomorphisme
surjectif $\mathcal{M}\to\mathcal{S}_1$,
$\overline{u}:\mathbf{S}_{\mathcal{O}_Y}(\mathcal{M})\to\mathcal{S}$
l'homomorphisme de $\mathcal{O}_Y$-Algèbres graduées prolongeant $u$
(\hyperref[II.1.7.4-fr]{1.7.4}). Le morphisme correspondant à
$\overline{u}$ est alors une immersion fermée de
$\operatorname{Proj}(\mathcal{S})$ dans
$\operatorname{Proj}(\mathbf{S}_{\mathcal{O}_Y}(\mathcal{M}))$.
\end{corollary}

\begin{proof}
En effet, $\overline{u}$ est surjectif par hypothèse et on applique
(\hyperref[II.3.6.1-fr]{3.6.1, (i)}).
\end{proof}

\subsection{Morphismes d'un préschéma dans un spectre homogène.}
\label{subsection:II.3.7-fr}

\begin{env}[3.7.1]
\label{II.3.7.1-fr}
Soient $q:X\to Y$ un morphisme de préschémas, $\mathcal{L}$ un
$\mathcal{O}_X$-Module inversible, $\mathcal{S}$ une $\mathcal{O}_Y$-Algèbre
graduée quasi-cohérente à degrés positifs~; $q^*(\mathcal{S})$ est alors une
$\mathcal{O}_X$-Algèbre graduée quasi-cohérente à degrés positifs.
Considérons la $\mathcal{O}_X$-Algèbre graduée quasi-cohérente
$\mathcal{S}'=\bigoplus_{n\geq0}\mathcal{L}^{\otimes n}$ et supposons donné
un $\mathcal{O}_X$-homomorphisme d'Algèbres graduées
\[
  \psi:q^*(\mathcal{S})\longrightarrow
  \mathcal{S}'=\bigoplus_{n\geq0}\mathcal{L}^{\otimes n}
\]
ce qui revient d'ailleurs à se donner un $q$-morphisme d'Algèbres graduées
\[
  \psi^\flat:\mathcal{S}\longrightarrow q_*(\mathcal{S}').
\]
On sait que $\operatorname{Proj}(\mathcal{S}')$ s'identifie canoniquement à
$X$ (\hyperref[II.3.1.7-fr]{3.1.7} et
\hyperref[II.3.1.8-fr]{3.1.8, (iii)})~; on déduit canoniquement de $\psi$ un
ensemble ouvert $G(\psi)$ de $X$ et un $Y$-morphisme
\[
\label{II.3.7.1.1-fr}
  r_{\mathcal{L},\psi}:G(\psi)\longrightarrow
  \operatorname{Proj}(\mathcal{S})=P
\tag{3.7.1.1}
\]

\oldpage[II]{66}

que nous appellerons le morphisme \emph{associé à $\mathcal{L}$ et à $\psi$}~;
rappelons (\hyperref[II.3.5.6-fr]{3.5.6}) que ce morphisme s'obtient en
composant avec la première projection
$\pi:\operatorname{Proj}(q^*(\mathcal{S}))=P\times_YX\to P$ le
$Y$-morphisme
\[
  \tau=\operatorname{Proj}(\psi):G(\psi)\longrightarrow
  \operatorname{Proj}(q^*(\mathcal{S})).
\]
\end{env}

\begin{env}[3.7.2]
\label{II.3.7.2-fr}
Explicitons $r=r_{\mathcal{L},\psi}$ lorsque $Y=\operatorname{Spec}(A)$ est
affine, et par suite $\mathcal{S}=\widetilde{S}$, où $S$ est une $A$-algèbre
graduée à degrés positifs. Supposons d'abord que $X=\operatorname{Spec}(B)$
soit aussi affine et que l'on ait $\mathcal{L}=\widetilde{L}$, où $L$ est un
$B$-module libre de rang $1$. On a alors
$q^*(\mathcal{S})=(S\otimes_AB)^{\sim}$
(\hyperref[I.1.6.5-fr]{I, 1.6.5})~; si $c$ est un générateur de $L$,
$\psi_n:q^*(\mathcal{S}_n)\to\mathcal{L}^{\otimes n}$ correspond à un
homomorphisme
$w_n:s\otimes b\mapsto bv_n(s)c^{\otimes n}$ de $S_n\otimes_AB$ dans
$L^{\otimes n}$, où $v_n:S_n\to B$ est un homomorphisme de $A$-modules, les
$v_n$ constituant un homomorphisme d'algèbres $S\to B$. Soit $f\in S_d$
($d>0$) et posons $g=v_d(f)$~; on a
$\pi^{-1}(D_+(f))=D_+(f\otimes1)$ en vertu de
(\hyperref[II.2.8.10-fr]{2.8.10}) et de l'identification de $D_+(f)$ à
$\operatorname{Spec}(S_{(f)})$ (\hyperref[II.2.3.6-fr]{2.3.6})~; d'autre
part, la formule (\hyperref[II.2.8.1.1-fr]{2.8.1.1}) montre (compte tenu de
l'identification canonique de $X$ et de $\operatorname{Proj}(\mathcal{S}')$)
que
\[
  \tau^{-1}(D_+(f\otimes1))=D(g)
\]
d'où
\[
\label{II.3.7.2.1-fr}
  r^{-1}(D_+(f))=D(g).
\tag{3.7.2.1}
\]

En outre, le morphisme $\tau=\operatorname{Proj}(\psi)$, restreint à $D(g)$,
correspond à l'homomorphisme qui applique
$(s\otimes1)/(f\otimes1)^n$ (pour $s\in S_{nd}$) sur
$v_{nd}(s)/g^n$ (\hyperref[II.2.8.1-fr]{2.8.1}), et la projection $\pi$,
restreinte à $D_+(f\otimes1)$, correspond à l'homomorphisme
$s/f^n\mapsto(s\otimes1)/(f\otimes1)^n$~; on en conclut que $r$, restreint à
$D(g)$, correspond à l'homomorphisme de $A$-algèbres
$\omega:S_{(f)}\to B_g$ tel que $\omega(s/f^n)=v_{nd}(s)/g^n$ pour
$s\in S_{nd}$ ($n>0$). Passant au cas où $X$ est quelconque ($Y$ étant
toujours affine), on obtient donc, compte tenu de
(\hyperref[II.2.8.1-fr]{2.8.1})~:
\end{env}

\begin{proposition}[3.7.3]
\label{II.3.7.3-fr}
Si $Y=\operatorname{Spec}(A)$ est affine et $\mathcal{S}=\widetilde{S}$, où
$S$ est une $A$-algèbre graduée, pour tout
$f\in S_d=\Gamma(Y,\mathcal{S}_d)$, on a
\[
\label{II.3.7.3.1-fr}
  r_{\mathcal{L},\psi}^{-1}(D_+(f))=X_{\psi^\flat(f)}
  \qquad\bigl(\text{où }\psi^\flat(f)\in
  \Gamma(X,\mathcal{L}^{\otimes d})\bigr)
\tag{3.7.3.1}
\]
et la restriction $X_{\psi^\flat(f)}\to D_+(f)=\operatorname{Spec}(S_{(f)})$
de $r_{\mathcal{L},\psi}$ correspond
(\hyperref[I.2.2.4-fr]{I, 2.2.4}) à l'homomorphisme d'algèbres
\[
\label{II.3.7.3.2-fr}
  \psi_{(f)}^\flat:S_{(f)}\longrightarrow
  \Gamma(X_{\psi^\flat(f)},\mathcal{O}_X)
\tag{3.7.3.2}
\]
tel que, pour $s\in S_{nd}=\Gamma(Y,\mathcal{S}_{nd})$,
\[
\label{II.3.7.3.3-fr}
  \psi_{(f)}^\flat(s/f^n)
  =(\psi^\flat(s)|X_{\psi^\flat(f)})
   (\psi^\flat(f)|X_{\psi^\flat(f)})^{-n}.
\tag{3.7.3.3}
\]
\end{proposition}

Nous dirons que $r_{\mathcal{L},\psi}$ est \emph{partout défini} si l'on a
$G(\psi)=X$. Pour qu'il en soit ainsi, il faut et il suffit évidemment que
$G(\psi)\cap q^{-1}(U)=q^{-1}(U)$ pour tout ouvert affine $U\subset Y$~;
autrement dit, la question est \emph{locale sur $Y$}. Si $Y$ est affine,
$G(\psi)$ est réunion des $r^{-1}(D_+(f))$ pour $f$ homogène dans $S_+$
(\hyperref[II.2.8.1-fr]{2.8.1})~; en vertu de
(\hyperref[II.3.7.3.1-fr]{3.7.3.1}), les $X_{\psi^\flat(f)}$ doivent donc
former un \emph{recouvrement} de $X$, autrement dit~:

\begin{corollary}[3.7.4]
\label{II.3.7.4-fr}
Sous les hypothèses de (\hyperref[II.3.7.3-fr]{3.7.3}), pour que
$r_{\mathcal{L},\psi}$ soit partout défini, il faut et il suffit que pour tout
$x\in X$, il existe un entier $n>0$ et une section $s$ de $\mathcal{S}_n$
au-dessus de $Y$ telle que si on pose
$t=\psi^\flat(s)\in\Gamma(X,\mathcal{L}^{\otimes n})$, on ait $t(x)\neq0$.
\end{corollary}

\oldpage[II]{67}

On notera que cette condition est toujours vérifiée si $\psi$ est
(\textbf{TN})-surjectif.

De même, la question de savoir si $r_{\mathcal{L},\psi}$ est \emph{dominant}
est locale sur $Y$, et l'on a~:

\begin{corollary}[3.7.5]
\label{II.3.7.5-fr}
Sous les hypothèses de (\hyperref[II.3.7.3-fr]{3.7.3}), pour que
$r_{\mathcal{L},\psi}$ soit dominant, il faut et il suffit que, pour tout
entier $n>0$, toute section $s\in S_n$ telle que
$\psi^\flat(s)\in\Gamma(X,\mathcal{L}^{\otimes n})$ soit localement
nilpotente, soit elle-même nilpotente.
\end{corollary}

\begin{proof}
On doit en effet exprimer que
$r_{\mathcal{L},\psi}^{-1}(D_+(s))$ n'est vide que si $D_+(s)$ est vide,
et le corollaire résulte donc de
(\hyperref[II.3.7.3.1-fr]{3.7.3.1}) et de
(\hyperref[II.2.3.7-fr]{2.3.7}).
\end{proof}

\begin{proposition}[3.7.6]
\label{II.3.7.6-fr}
Soient $q:X\to Y$ un morphisme, $\mathcal{L}$ un $\mathcal{O}_X$-Module
inversible, $\mathcal{S}$, $\mathcal{S}'$ deux $\mathcal{O}_Y$-Algèbres
graduées quasi-cohérentes,
$u:\mathcal{S}'\to\mathcal{S}$ un homomorphisme d'Algèbres graduées,
$\psi:q^*(\mathcal{S})\to\bigoplus_{n\geq0}\mathcal{L}^{\otimes n}$ un
homomorphisme d'Algèbres graduées,
$\psi'=\psi\circ q^*(u)$ l'homomorphisme composé. Si
$r_{\mathcal{L},\psi'}$ est partout défini, il en est de même de
$r_{\mathcal{L},\psi}$~; si $u$ est (\textbf{TN})-surjectif et si
$r_{\mathcal{L},\psi'}$ est dominant, il en est de même de
$r_{\mathcal{L},\psi}$~; inversement, si $u$ est (\textbf{TN})-injectif et si
$r_{\mathcal{L},\psi}$ est dominant, $r_{\mathcal{L},\psi'}$ est dominant.
\end{proposition}

\begin{proof}
On a en effet $G(\psi')\subset G(\psi)$
(\hyperref[II.2.8.4-fr]{2.8.4}), d'où la première assertion~; si $u$ est
(\textbf{TN})-surjectif,
$\operatorname{Proj}(u):\operatorname{Proj}(\mathcal{S})\to
\operatorname{Proj}(\mathcal{S}')$ est partout défini et est une immersion
fermée~; comme $r_{\mathcal{L},\psi'}$ est composé de
$\operatorname{Proj}(u)$ et de la restriction de $r_{\mathcal{L},\psi}$ à
$G(\psi')$, on en conclut que si $r_{\mathcal{L},\psi'}$ est dominant, il en
est de même de $r_{\mathcal{L},\psi}$. Enfin, si $u$ est
(\textbf{TN})-injectif, on sait que $\operatorname{Proj}(u)$ est un morphisme
dominant de $G(u)$ dans $\operatorname{Proj}(\mathcal{S}')$
(\hyperref[II.2.8.3-fr]{2.8.3})~; comme $G(\psi')$ est l'image réciproque de
$G(u)$ par $r_{\mathcal{L},\psi}$, on voit que si $r_{\mathcal{L},\psi}$ est
dominant, il en est de même de $r_{\mathcal{L},\psi'}$.
\end{proof}

\begin{proposition}[3.7.7]
\label{II.3.7.7-fr}
Soient $Y$ un préschéma quasi-compact, $q:X\to Y$ un morphisme quasi-compact,
$\mathcal{L}$ un $\mathcal{O}_X$-Module inversible, $\mathcal{S}$ une
$\mathcal{O}_Y$-Algèbre graduée quasi-cohérente, limite inductive filtrante
d'un système inductif $(\mathcal{S}^\lambda)$ de $\mathcal{O}_Y$-Algèbres
quasi-cohérentes. Soient
$\varphi_\lambda:\mathcal{S}^\lambda\to\mathcal{S}$ l'homomorphisme
canonique,
$\psi:q^*(\mathcal{S})\to\bigoplus_{n\geq0}\mathcal{L}^{\otimes n}$ un
homomorphisme d'Algèbres graduées, et posons
$\psi_\lambda=\psi\circ q^*(\varphi_\lambda)$. Pour que
$r_{\mathcal{L},\psi}$ soit partout défini, il faut et il suffit qu'il existe
un $\lambda$ tel que $r_{\mathcal{L},\psi_\lambda}$ soit partout défini~;
$r_{\mathcal{L},\psi_\mu}$ est alors partout défini pour $\mu\geq\lambda$.
\end{proposition}

\begin{proof}
La condition est suffisante en vertu de
(\hyperref[II.3.7.6-fr]{3.7.6}). Inversement, supposons
$r_{\mathcal{L},\psi}$ partout défini~; on peut se ramener au cas où $Y$ est
affine, car si pour tout ouvert affine $U\subset Y$, il existe $\lambda(U)$
tel que la restriction de $r_{\mathcal{L},\psi_{\lambda(U)}}$ à $q^{-1}(U)$
soit partout définie, il suffira ($Y$ étant quasi-compact) de recouvrir $Y$
par un nombre fini d'ouverts affines $U_i$, et de prendre
$\lambda\geq\lambda(U_i)$ pour tous les indices $i$, en vertu de
(\hyperref[II.3.7.6-fr]{3.7.6}). Si $Y$ est affine, l'hypothèse entraîne que
pour tout $x\in X$, il y a une section $s^{(x)}$ d'un $S_n$ telle que, si on
pose $t^{(x)}=\psi^\flat(s^{(x)})$, on ait $t^{(x)}(x)\neq0$
($t^{(x)}$ étant considérée comme section de $\mathcal{L}^{\otimes n}$
au-dessus de $X$), ce qui entraîne $t^{(x)}(z)\neq0$ pour tout $z$ dans un
voisinage $V(x)$ de $x$. Couvrons $X$ par un nombre fini de $V(x_i)$ et
soient $s^{(i)}$ les sections correspondantes de $S$~; il existe alors un
$\lambda$ tel que les $s^{(i)}$ soient toutes de la forme
$\varphi_\lambda(s_\lambda^{(i)})$, avec $s_\lambda^{(i)}\in S^\lambda$ pour
tous les $i$~; il résulte donc de (\hyperref[II.3.7.4-fr]{3.7.4}) que
$r_{\mathcal{L},\psi_\lambda}$ est partout défini. La dernière assertion est
conséquence triviale de (\hyperref[II.3.7.6-fr]{3.7.6}).
\end{proof}

\begin{corollary}[3.7.8]
\label{II.3.7.8-fr}
Sous les hypothèses de (\hyperref[II.3.7.7-fr]{3.7.7}), si les
$r_{\mathcal{L},\psi_\lambda}$ sont dominants, il en est de même de
$r_{\mathcal{L},\psi}$~; la réciproque est vraie lorsque les
$\varphi_\lambda$ sont injectifs.
\end{corollary}

\oldpage[II]{68}

\begin{proof}
La seconde assertion est un cas particulier de
(\hyperref[II.3.7.6-fr]{3.7.6})~; d'autre part, pour montrer que
$r_{\mathcal{L},\psi}$ est dominant, on peut se borner au cas où $Y$ est
affine~; si $s\in S$ est telle que $\psi^\flat(s)$ soit localement
nilpotente, comme on peut écrire $s=\varphi_\lambda(s_\lambda)$ pour un
$\lambda$ au moins, on conclut de l'hypothèse et de
(\hyperref[II.3.7.5-fr]{3.7.5}) que $s_\lambda$ est nilpotente, donc il en est
de même de $s$, et le critère de (\hyperref[II.3.7.5-fr]{3.7.5}) s'applique.
\end{proof}

\begin{namedtheorem}[3.7.9]{Remarques}
\label{II.3.7.9-fr}
\begin{enumerate}
  \item[(i)] Avec les notations de (\hyperref[II.3.7.1-fr]{3.7.1}) et en
    tenant compte de (\hyperref[II.3.2.10-fr]{3.2.10}), on a, pour tout
    $n\in\mathbf{Z}$, un homomorphisme canonique
    \[
    \label{II.3.7.9.1-fr}
      \theta:r_{\mathcal{L},\psi}^*(\mathcal{O}_P(n))\longrightarrow
      \mathcal{L}^{\otimes n}
    \tag{3.7.9.1}
    \]
    défini de façon générale dans
    (\hyperref[II.3.5.6.2-fr]{3.5.6.2}). On voit aussitôt que sous les
    conditions de (\hyperref[II.3.7.3-fr]{3.7.3}), la restriction de cet
    homomorphisme à $X_{\psi^\flat(f)}$ s'explicite comme faisant
    correspondre à $s/f^k$ ($s\in S_{n+kd}$) l'élément
    $(\psi^\flat(s)|X_{\psi^\flat(f)})
    (\psi^\flat(f)|X_{\psi^\flat(f)})^{-k}$.
  \item[(ii)] Soit $\mathcal{F}$ un $\mathcal{O}_X$-Module quasi-cohérent, et
    supposons que $q$ soit quasi-compact et séparé, de sorte que pour tout
    $n\geq0$, $q_*(\mathcal{F}\otimes\mathcal{L}^{\otimes n})$ soit un
    $\mathcal{O}_Y$-Module quasi-cohérent
    (\hyperref[I.9.2.2-fr]{I, 9.2.2}). Soit
    $\mathcal{M}'=\bigoplus_{n\geq0}\mathcal{F}\otimes
    \mathcal{L}^{\otimes n}$, qui est un $\mathcal{S}'$-Module gradué
    quasi-cohérent, et considérons son image
    $\mathcal{M}=q_*(\mathcal{M}')=
    \bigoplus_{n\geq0}q_*(\mathcal{F}\otimes\mathcal{L}^{\otimes n})$
    (qui est un $\mathcal{S}$-Module quasi-cohérent, au moyen de
    l'homomorphisme $\psi^\flat$). Nous allons voir qu'il y a un
    homomorphisme canonique de $\mathcal{O}_X$-Modules
    \[
    \label{II.3.7.9.2-fr}
      \xi:r_{\mathcal{L},\psi}^*(\widetilde{\mathcal{M}})
      \longrightarrow\mathcal{F}|G(\psi).
    \tag{3.7.9.2}
    \]

    En effet, on a défini (\hyperref[II.3.5.6.1-fr]{3.5.6.1}) un
    homomorphisme canonique
    \[
    \label{II.3.7.9.3-fr}
      r_{\mathcal{L},\psi}^*(\widetilde{\mathcal{M}})
      \longrightarrow
      (q^*(\mathcal{M})\otimes_{q^*(\mathcal{S})}\mathcal{S}')^{\sim}
      |G(\psi),
    \tag{3.7.9.3}
    \]
    où le second membre est considéré comme un faisceau quasi-cohérent sur
    $\operatorname{Proj}(\mathcal{S}')$. D'autre part, on a un
    homomorphisme canonique
    \[
    \label{II.3.7.9.4-fr}
      q^*(q_*(\mathcal{M}'))\otimes_{q^*(\mathcal{S})}\mathcal{S}'
      \longrightarrow\mathcal{M}'
    \tag{3.7.9.4}
    \]
    pour tout $\mathcal{S}'$-Module gradué quasi-cohérent $\mathcal{M}'$~:
    pour tout ouvert $U$ de $X$, toute section $t'$ de
    $q^*(q_*(\mathcal{M}'_h))$ au-dessus de $U$ et toute section $b'$ de
    $\mathcal{S}'_k$ au-dessus de $U$, on fait correspondre à $t'\otimes b'$ la
    section $b'\sigma(t')$ de $\mathcal{M}'_{h+k}$, où $\sigma(t')$ est la
    section de $\mathcal{M}'_h$ au-dessus de $U$ qui correspond canoniquement
    (\hyperref[0.4.4.3-fr]{0, 4.4.3}) à $t'$. On en conclut un homomorphisme
    canonique
    \[
    \label{II.3.7.9.5-fr}
      (q^*(q_*(\mathcal{M}'))\otimes_{q^*(\mathcal{S})}\mathcal{S}')^{\sim}
      |G(\psi)\longrightarrow\widetilde{\mathcal{M}'}|G(\psi)
    \tag{3.7.9.5}
    \]
    et comme finalement $\widetilde{\mathcal{M}'}$ s'identifie
    canoniquement à $\mathcal{F}$
    (\hyperref[II.3.2.9-fr]{3.2.9, (i)}), on obtient bien l'homomorphisme
    canonique cherché.

    Sous les conditions de (\hyperref[II.3.7.3-fr]{3.7.3}), la restriction
    de cet homomorphisme à $X_{\psi^\flat(f)}$ s'explicite de la façon
    suivante~: étant donnée une section $t_{nd}$ de
    $\mathcal{F}\otimes\mathcal{L}^{\otimes nd}$ au-dessus de $X$, si
    $t'_{nd}$ est la section $t_{nd}$ considérée comme section de
    $q_*(\mathcal{F}\otimes\mathcal{L}^{\otimes nd})$ au-dessus de $Y$, à
    l'élément $t'_{nd}/f^n$ il correspond la section
    $(t_{nd}|X_{\psi^\flat(f)})
    (\psi^\flat(f)|X_{\psi^\flat(f)})^{-n}$ de $\mathcal{F}$ au-dessus de
    $X_{\psi^\flat(f)}$.
\end{enumerate}
\end{namedtheorem}

\oldpage[II]{69}

\subsection{Critères d'immersion dans un spectre homogène.}
\label{subsection:II.3.8-fr}

\begin{env}[3.8.1]
\label{II.3.8.1-fr}
Avec les notations de (\hyperref[II.3.7.1-fr]{3.7.1}), la question de savoir
si $r_{\mathcal{L},\psi}$ est une immersion (resp. une immersion ouverte, une
immersion fermée) est évidemment \emph{locale sur $Y$}.
\end{env}

\begin{proposition}[3.8.2]
\label{II.3.8.2-fr}
Sous les hypothèses de (\hyperref[II.3.7.3-fr]{3.7.3}), pour que
$r_{\mathcal{L},\psi}$ soit partout définie et soit une immersion, il faut et
il suffit qu'il existe une famille de sections $s_\alpha\in S_{n_\alpha}$
($n_\alpha>0$) telle que, si on pose
$f_\alpha=\psi^\flat(s_\alpha)$, les conditions suivantes soient vérifiées~:
\begin{enumerate}
  \item[(i)] Les $X_{f_\alpha}$ forment un recouvrement de $X$.
  \item[(ii)] Les $X_{f_\alpha}$ sont des ouverts affines.
  \item[(iii)] Pour tout $\alpha$ et tout
    $t\in\Gamma(X_{f_\alpha},\mathcal{O}_X)$, il existe un entier $m>0$ et
    un $s\in S_{mn_\alpha}$ tels que
    $t=(\psi^\flat(s)|X_{f_\alpha})(f_\alpha|X_{f_\alpha})^{-m}$.
\end{enumerate}

Pour que $r_{\mathcal{L},\psi}$ soit partout définie et soit une immersion
ouverte, il faut et il suffit qu'il existe une famille $(s_\alpha)$ vérifiant
les conditions (i), (ii), (iii) et~:
\begin{enumerate}
  \item[(iv)] Pour tout $n>0$ et tout $s\in S_{nn_\alpha}$ tel que
    $\psi^\flat(s)|X_{f_\alpha}=0$, il existe un entier $k>0$ tel que
    $s_\alpha^ks=0$.
\end{enumerate}

Pour que $r_{\mathcal{L},\psi}$ soit partout définie et soit une immersion
fermée, il faut et il suffit qu'il existe une famille $(s_\alpha)$ vérifiant
(i), (ii), (iii) et~:
\begin{enumerate}
  \item[(v)] Les $D_+(s_\alpha)$ forment un recouvrement de
    $P=\operatorname{Proj}(S)$.
\end{enumerate}
\end{proposition}

\begin{proof}
En effet, pour que $r$ soit une immersion (resp. une immersion fermée), il
faut et il suffit qu'il existe un recouvrement de $r(G(\psi))$ (resp. de
$P$) par des $D_+(s_\alpha)$ tel que si on pose
$V_\alpha=r^{-1}(D_+(s_\alpha))$, la restriction de $r$ à $U_\alpha$ soit
une immersion fermée de $V_\alpha$ dans $D_+(s_\alpha)$
(\hyperref[I.4.2.4-fr]{I, 4.2.4}). La condition (i) exprime à la fois que
$r$ est partout définie et que les $D_+(s_\alpha)$ recouvrent $r(X)$,
d'après (\hyperref[II.3.7.3.1-fr]{3.7.3.1})~; comme $D_+(s_\alpha)$ est
affine, les conditions (ii) et (iii) expriment que la restriction de $r$ à
$X_{f_\alpha}$ est une immersion fermée dans $D_+(s_\alpha)$, d'après
(\hyperref[I.4.2.3-fr]{I, 4.2.3})~; enfin, comme (iii) et (iv) expriment que
$\psi_{(s_\alpha)}^\flat$ est bijective (notation de
(\hyperref[II.3.7.3.2-fr]{3.7.3.2})), (ii), (iii) et (iv) expriment que la
restriction de $r$ à $X_{f_\alpha}$ est un isomorphisme sur
$D_+(s_\alpha)$ pour tout $\alpha$, donc (i), (ii), (iii) et (iv) expriment
que $r$ est une immersion ouverte.
\end{proof}

\begin{corollary}[3.8.3]
\label{II.3.8.3-fr}
Sous les hypothèses de (\hyperref[II.3.7.6-fr]{3.7.6}), si
$r_{\mathcal{L},\psi'}$ est partout définie et est une immersion, il en est
de même de $r_{\mathcal{L},\psi}$. Si on suppose en outre que $u$ est
(\textbf{TN})-surjectif et si $r_{\mathcal{L},\psi'}$ est une immersion
ouverte (resp. fermée), il en est de même de $r_{\mathcal{L},\psi}$.
\end{corollary}

\begin{proof}
En effet, d'après (\hyperref[II.3.8.2-fr]{3.8.2}), il y a une famille de
$s'_\alpha\in S'_{n_\alpha}$ telle que si on pose
$f_\alpha=\psi'^\flat(s'_\alpha)$, les conditions (i), (ii), (iii) soient
vérifiées. Or, si on pose $s_\alpha=u(s'_\alpha)$, on a aussi
$f_\alpha=\psi^\flat(s_\alpha)$, et si on a
$t=(\psi'^\flat(s')|X_{f_\alpha})(f_\alpha|X_{f_\alpha})^{-m}$, on a aussi
$t=(\psi^\flat(s)|X_{f_\alpha})(f_\alpha|X_{f_\alpha})^{-m}$ en posant
$s=u(s')$, d'où la première assertion. La seconde résulte aussitôt de ce que
$\operatorname{Proj}(u)$ est une immersion fermée.
\end{proof}

\begin{proposition}[3.8.4]
\label{II.3.8.4-fr}
Supposons vérifiées les hypothèses de
(\hyperref[II.3.7.7-fr]{3.7.7}) et en outre que $q:X\to Y$ soit un morphisme
de type fini. Alors, pour que $r_{\mathcal{L},\psi}$ soit partout définie et
soit une immersion, il faut et il suffit qu'il existe $\lambda$ tel que
$r_{\mathcal{L},\psi_\lambda}$ soit partout définie et soit une immersion et
dans ce cas $r_{\mathcal{L},\psi_\mu}$ est partout définie et est une
immersion pour $\mu\geq\lambda$.
\end{proposition}

\oldpage[II]{70}

\begin{proof}
Compte tenu de (\hyperref[II.3.8.3-fr]{3.8.3}), il suffit de démontrer que si
$r_{\mathcal{L},\psi}$ est partout définie et est une immersion, il en est de
même de $r_{\mathcal{L},\psi_\lambda}$ pour un $\lambda$ au moins. Par le
même raisonnement que dans (\hyperref[II.3.7.7-fr]{3.7.7}), utilisant la
quasi-compacité de $Y$, on se ramène d'abord au cas où $Y$ est affine. Comme
$X$ est quasi-compact, (\hyperref[II.3.8.2-fr]{3.8.2}) montre l'existence
d'une famille \emph{finie} $(s_i)$ d'éléments de $S$ ($s_i\in S_{n_i}$)
ayant les propriétés (i), (ii) et (iii). Le morphisme $X_{f_i}\to Y$ (où
$f_i=\psi^\flat(s_i)$) est de type fini~: en effet, c'est un morphisme de
schémas affines, donc il est quasi-compact
(\hyperref[I.6.6.1-fr]{I, 6.6.1}), et localement de type fini puisque $q$ est
de type fini (\hyperref[I.6.3.2-fr]{I, 6.3.2}), et la conclusion résulte de
(\hyperref[I.6.6.3-fr]{I, 6.6.3}). L'anneau $B_i$ de $X_{f_i}$ est par suite
une $A$-algèbre de type fini (\hyperref[I.6.3.3-fr]{I, 6.3.3})~; soit
$(t_{ij})$ une famille de générateurs de cette algèbre. Il y a par hypothèse
des éléments $s'_{ij}\in S_{m_{ij}n_i}$ tels que
\[
  t_{ij}=(\psi^\flat(s'_{ij})|X_{f_i})
  (\psi^\flat(s_i)|X_{f_i})^{-m_{ij}}.
\]
Il y a par hypothèse un $\lambda$ et des éléments
$s_{i\lambda}\in S_{n_i}^\lambda$,
$s'_{ij\lambda}\in S_{m_{ij}n_i}^\lambda$ dont les images par
$\varphi_\lambda$ sont respectivement $s_i$ et $s'_{ij}$~; il est clair que
$r_{\mathcal{L},\psi_\lambda}$ vérifie les conditions (i), (ii) et (iii) de
(\hyperref[II.3.8.2-fr]{3.8.2}).
\end{proof}

\begin{proposition}[3.8.5]
\label{II.3.8.5-fr}
Soient $Y$ un schéma quasi-compacte, ou un préschéma dont l'espace sous-jacent
est noethérien, $q:X\to Y$ un morphisme de type fini, $\mathcal{L}$ un
$\mathcal{O}_X$-Module inversible, $\mathcal{S}$ une $\mathcal{O}_Y$-Algèbre
graduée quasi-cohérente,
$\psi:q^*(\mathcal{S})\to\bigoplus_{n\geq0}\mathcal{L}^{\otimes n}$ un
homomorphisme d'Algèbres graduées. Pour que $r_{\mathcal{L},\psi}$ soit
partout défini et soit une immersion, il faut et il suffit qu'il existe un
entier $n>0$ et un sous-$\mathcal{O}_Y$-Module quasi-cohérent de type fini
$\mathcal{E}$ de $\mathcal{S}_n$, tels que~:
\begin{enumerate}
  \item[a)] l'homomorphisme
    $\psi_n\circ q^*(j_n):q^*(\mathcal{E})\to\mathcal{L}^{\otimes n}$
    (où $j_n:\mathcal{E}\to\mathcal{S}$ est l'injection canonique) soit
    surjectif~;
  \item[b)] si on désigne par $\mathcal{S}'$ la sous-$\mathcal{O}_Y$-Algèbre
    (graduée) de $\mathcal{S}$ engendrée par $\mathcal{E}$, par $\psi'$
    l'homomorphisme $\psi\circ q^*(j')$ ($j'$ injection de $\mathcal{S}'$
    dans $\mathcal{S}$), $r_{\mathcal{L},\psi'}$ soit partout défini et soit
    une immersion.
\end{enumerate}

Lorsqu'il en est ainsi, tout sous-$\mathcal{O}_Y$-Module quasi-cohérent
$\mathcal{E}'$ de $\mathcal{S}_n$, contenant $\mathcal{E}$, possède la même
propriété, et il en est de même de l'image de $\bigotimes^k\mathcal{E}$ dans
$\mathcal{S}_{kn}$ pour tout $k>0$.
\end{proposition}

\begin{proof}
La suffisance de la condition et les deux dernières assertions sont des cas
particuliers de (\hyperref[II.3.8.3-fr]{3.8.3}), compte tenu de l'isomorphie
canonique entre $\operatorname{Proj}(\mathcal{S})$ et
$\operatorname{Proj}(\mathcal{S}^{(k)})$
(\hyperref[II.3.1.8-fr]{3.1.8}).

Soit $(U_i)$ un recouvrement ouvert fini de $Y$ par des ouverts affines et
posons $A_i=A(U_i)$. Comme $q^{-1}(U_i)$ est compact, l'hypothèse que
$r_{\mathcal{L},\psi}$ est une immersion partout définie et
(\hyperref[II.3.8.2-fr]{3.8.2}) entraînent l'existence d'une famille finie
$(s_{ij})$ d'éléments de
$S^{(i)}=\Gamma(U_i,\mathcal{S})$ ($s_{ij}\in S_{n_{ij}}^{(i)}$) ayant les
propriétés (i), (ii), (iii). On voit comme dans la démonstration de
(\hyperref[II.3.8.4-fr]{3.8.4}) que le morphisme $X_{f_{ij}}\to U_i$ (où
$f_{ij}=\psi^\flat(s_{ij})$) est de type fini, et l'anneau $B_{ij}$ de
$X_{f_{ij}}$ est par suite une $A_i$-algèbre de type fini, ayant un système de
générateurs de la forme
$(\psi^\flat(t_{ijk})|X_{f_{ij}})(f_{ij}|X_{f_{ij}})^{-m_{ijk}}$, où
$t_{ijk}\in S_{m_{ijk}n_{ij}}^{(i)}$. Soit $n$ un multiple commun des
$m_{ijk}n_{ij}$~; remplaçant (pour chaque $(i,j,k)$) $s_{ij}$ par une
puissance $s_{ij}^\rho$ telle que $\rho m_{ijk}n_{ij}=n$, et multipliant
$t_{ijk}$ par $s_{ij}^{\rho-m_{ijk}}$, on peut supposer que pour chaque $i$
tous les $s_{ij}$ et $t_{ijk}$ appartiennent à $S_n^{(i)}$ et que
$m_{ijk}=1$. Soit $E_i$ le sous-$A_i$-Module de $S_n^{(i)}$ engendré par ces
éléments (pour $i$ fixé). Il existe un sous-$\mathcal{O}_Y$-Module cohérent de
type

\oldpage[II]{71}

fini $\mathcal{E}_i$ de $\mathcal{S}_n$ tel que
$\mathcal{E}_i|U_i=(E_i)^{\sim}$
(\hyperref[I.9.4.7-fr]{I, 9.4.7}). Il est clair que le
sous-$\mathcal{O}_Y$-Module $\mathcal{E}$ de $\mathcal{S}_n$, somme des
$\mathcal{E}_i$, répond à la question (chaque section $f_{ij}$ étant telle
que pour tout $x\in X_{f_{ij}}$, il y a un voisinage affine
$V\subset X_{f_{ij}}$ de $x$ tel que $f|V$ soit une base de
$\Gamma(V,\mathcal{L}^{\otimes n})$).
\end{proof}

\section{Fibres projectifs. Faisceaux amples}
\label{section:II.4-fr}

\subsection{Définition des fibres projectifs.}
\label{subsection:II.4.1-fr}

\begin{definition}[4.1.1]
\label{II.4.1.1-fr}
Soient $Y$ un préschéma, $\mathcal{E}$ un $\mathcal{O}_Y$-Module
quasi-cohérent, $\mathbf{S}_{\mathcal{O}_Y}(\mathcal{E})$ la
$\mathcal{O}_Y$-Algèbre symétrique de $\mathcal{E}$
(\hyperref[II.1.7.4-fr]{1.7.4}) qui est quasi-cohérente
(\hyperref[II.1.7.7-fr]{1.7.7}). On appelle \emph{fibré projectif sur $Y$
défini par $\mathcal{E}$} et on note $\mathbf{P}(\mathcal{E})$ le $Y$-schéma
$P=\operatorname{Proj}(\mathbf{S}_{\mathcal{O}_Y}(\mathcal{E}))$. Le
$\mathcal{O}_P$-Module $\mathcal{O}_P(1)$ est appelé le \emph{faisceau
fondamental sur $P$}.
\end{definition}

Lorsque $Y$ est affine d'anneau $A$, on a $\mathcal{E}=\widetilde{E}$, où
$E$ est un $A$-module, et on écrit alors $\mathbf{P}(E)$ au lieu de
$\mathbf{P}(\widetilde{E})$.

Lorsque l'on prend $\mathcal{E}=\mathcal{O}_Y^n$, on écrit
$\mathbf{P}_Y^{n-1}$ au lieu de $\mathbf{P}(\mathcal{E})$~; si de plus $Y$
est affine d'anneau $A$, on écrit aussi $\mathbf{P}_A^{n-1}$ au lieu de
$\mathbf{P}_Y^{n-1}$. Comme
$\mathbf{S}_{\mathcal{O}_Y}(\mathcal{O}_Y)$ s'identifie canoniquement à
$\mathcal{O}_Y[T]$ (\hyperref[II.1.7.4-fr]{1.7.4}),
$\mathbf{P}_Y^0$ s'identifie canoniquement à $Y$
(\hyperref[II.3.1.7-fr]{3.1.7})~; l'exemple
(\hyperref[II.2.4.3-fr]{2.4.3}) n'est autre que $\mathbf{P}_K^1$.

\begin{env}[4.1.2]
\label{II.4.1.2-fr}
Soient $\mathcal{E}$, $\mathcal{F}$ deux $\mathcal{O}_Y$-Modules
quasi-cohérents~; $u:\mathcal{E}\to\mathcal{F}$ un
$\mathcal{O}_Y$-homomorphisme~; il lui correspond canoniquement un
homomorphisme
$\mathbf{S}(u):\mathbf{S}_{\mathcal{O}_Y}(\mathcal{E})\to
\mathbf{S}_{\mathcal{O}_Y}(\mathcal{F})$ de $\mathcal{O}_Y$-Algèbres
graduées (\hyperref[II.1.7.4-fr]{1.7.4}). Si $u$ est \emph{surjectif}, il en
est de même de $\mathbf{S}(u)$, et par suite
(\hyperref[II.3.6.2-fr]{3.6.2, (i)}),
$\operatorname{Proj}(\mathbf{S}(u))$ est une immersion fermée
$\mathbf{P}(\mathcal{F})\to\mathbf{P}(\mathcal{E})$, que nous noterons
$\mathbf{P}(u)$. On peut donc dire que $\mathbf{P}(\mathcal{E})$ est un
\emph{foncteur contravariant en $\mathcal{E}$}, à condition de ne prendre
pour morphismes des $\mathcal{O}_Y$-Modules quasi-cohérents que les
homomorphismes surjectifs.

En supposant toujours $u$ surjectif et posant
$P=\mathbf{P}(\mathcal{E})$, $Q=\mathbf{P}(\mathcal{F})$ et
$j=\mathbf{P}(u)$, on a, à un isomorphisme près,
\[
\label{II.4.1.2.1-fr}
  j^*(\mathcal{O}_P(n))=\mathcal{O}_Q(n)
  \qquad\text{pour tout }n\in\mathbf{Z}
\tag{4.1.2.1}
\]
comme il résulte de (\hyperref[II.3.6.3-fr]{3.6.3}).
\end{env}

\begin{env}[4.1.3]
\label{II.4.1.3-fr}
Soit maintenant $\psi:Y'\to Y$ un morphisme, et soit
$\mathcal{E}'=\psi^*(\mathcal{E})$~; on a alors
$\mathbf{S}_{\mathcal{O}_{Y'}}(\mathcal{E}')=
\psi^*(\mathbf{S}_{\mathcal{O}_Y}(\mathcal{E}))$
(\hyperref[II.1.7.5-fr]{1.7.5})~; par suite
(\hyperref[II.3.5.3-fr]{3.5.3}), on a
\[
\label{II.4.1.3.1-fr}
  \mathbf{P}(\psi^*(\mathcal{E}))=
  \mathbf{P}(\mathcal{E})\times_YY'
\tag{4.1.3.1}
\]
à un isomorphisme canonique près~; en outre, on a évidemment
\[
  \psi^*((\mathbf{S}_{\mathcal{O}_Y}(\mathcal{E}))(n))
  =(\mathbf{S}_{\mathcal{O}_{Y'}}(\mathcal{E}'))(n)
\]
pour tout $n\in\mathbf{Z}$, donc, en posant $P=\mathbf{P}(\mathcal{E})$,
$P'=\mathbf{P}(\mathcal{E}')$, on a
(\hyperref[II.3.5.4-fr]{3.5.4}), à un isomorphisme près,
\[
\label{II.4.1.3.2-fr}
  \mathcal{O}_{P'}(n)=\mathcal{O}_P(n)\otimes_Y\mathcal{O}_{Y'}
  \qquad\text{pour tout }n\in\mathbf{Z}.
\tag{4.1.3.2}
\]
\end{env}

\oldpage[II]{72}

\begin{proposition}[4.1.4]
\label{II.4.1.4-fr}
Soit $\mathcal{L}$ un $\mathcal{O}_Y$-Module inversible. Pour tout
$\mathcal{O}_Y$-Module quasi-cohérent $\mathcal{E}$, il existe un
$Y$-isomorphisme canonique
$i_{\mathcal{L}}:\mathbf{P}(\mathcal{E})\xrightarrow{\sim}
\mathbf{P}(\mathcal{E}\otimes\mathcal{L})$~; en outre, si l'on pose
$P=\mathbf{P}(\mathcal{E})$, $Q=\mathbf{P}(\mathcal{E}\otimes\mathcal{L})$,
$i_{\mathcal{L}}^*(\mathcal{O}_Q(n))$ est canoniquement isomorphe à
$\mathcal{O}_P(n)\otimes_Y\mathcal{L}^{\otimes n}$ pour tout
$n\in\mathbf{Z}$.
\end{proposition}

\begin{proof}
Remarquons d'abord que si $A$ est un anneau, $E$ un $A$-module, $L$ un
$A$-module \emph{monogène libre}, on définit canoniquement un homomorphisme
de $A$-modules
\[
  \mathbf{S}_n(E\otimes L)\longrightarrow
  \mathbf{S}_n(E)\otimes L^{\otimes n}
\]
en faisant correspondre à
$(x_1\otimes y_1)\cdots(x_n\otimes y_n)$ l'élément
\[
  (x_1x_2\cdots x_n)\otimes(y_1\otimes y_2\otimes\cdots\otimes y_n)
  \qquad(x_i\in E,\ y_i\in L\text{ pour }1\leq i\leq n)~;
\]
on vérifie immédiatement (en se ramenant au cas où $L=A$) que cet
homomorphisme est en fait un isomorphisme. On en conclut un isomorphisme
canonique de $A$-algèbres graduées
$\mathbf{S}_A(E\otimes L)\xrightarrow{\sim}
\bigoplus_{n\geq0}\mathbf{S}_n(E)\otimes L^{\otimes n}$. En revenant aux
conditions de (4.1.4), les remarques précédentes permettent de définir un
isomorphisme canonique de $\mathcal{O}_Y$-Algèbres graduées
\[
\label{II.4.1.4.1-fr}
  \mathbf{S}_{\mathcal{O}_Y}(\mathcal{E}\otimes_{\mathcal{O}_Y}\mathcal{L})
  \xrightarrow{\sim}
  \bigoplus_{n\geq0}\mathbf{S}_n(\mathcal{E})
  \otimes_{\mathcal{O}_Y}\mathcal{L}^{\otimes n}
\tag{4.1.4.1}
\]
en définissant cet isomorphisme comme un isomorphisme de préfaisceaux et
tenant compte de (\hyperref[II.1.7.4-fr]{1.7.4}),
(\hyperref[I.1.3.9-fr]{I, 1.3.9}) et
(\hyperref[I.1.3.12-fr]{I, 1.3.12}). La proposition résulte alors de
(\hyperref[II.3.1.8-fr]{3.1.8, (iii)}) et
(\hyperref[II.3.2.10-fr]{3.2.10}).
\end{proof}

\begin{env}[4.1.5]
\label{II.4.1.5-fr}
Avec les hypothèses de (\hyperref[II.4.1.1-fr]{4.1.1}), posons
$P=\mathbf{P}(\mathcal{E})$, et désignons par $p$ le morphisme structural
$P\to Y$. Comme on a par définition
$\mathcal{E}=(\mathbf{S}_{\mathcal{O}_Y}(\mathcal{E}))_1$, on a un
homomorphisme canonique
$\alpha_1:\mathcal{E}\to p_*(\mathcal{O}_P(1))$
(\hyperref[II.3.3.2.2-fr]{3.3.2.2}) et par suite aussi
(\hyperref[0.4.4.3-fr]{0, 4.4.3}) un homomorphisme canonique
\[
\label{II.4.1.5.1-fr}
  \alpha_1^\sharp:p^*(\mathcal{E})\longrightarrow\mathcal{O}_P(1).
\tag{4.1.5.1}
\]
\end{env}

\begin{proposition}[4.1.6]
\label{II.4.1.6-fr}
L'homomorphisme canonique (\hyperref[II.4.1.5.1-fr]{4.1.5.1}) est surjectif.
\end{proposition}

\begin{proof}
En effet, on a vu dans (\hyperref[II.3.3.2-fr]{3.3.2}) que
$\alpha_1^\sharp$ correspond fonctoriellement à l'homomorphisme canonique
\[
  \mathcal{E}\otimes_{\mathcal{O}_Y}
  \mathbf{S}_{\mathcal{O}_Y}(\mathcal{E})
  \longrightarrow
  (\mathbf{S}_{\mathcal{O}_Y}(\mathcal{E}))(1)~;
\]
comme par définition $\mathcal{E}$ engendre
$\mathbf{S}_{\mathcal{O}_Y}(\mathcal{E})$, cet homomorphisme est surjectif,
d'où la conclusion en vertu de (\hyperref[II.3.2.4-fr]{3.2.4}).
\end{proof}

\subsection{Morphismes d'un préschéma dans un fibré projectif.}
\label{subsection:II.4.2-fr}

\begin{env}[4.2.1]
\label{II.4.2.1-fr}
Gardant les notations de (\hyperref[II.4.1.5-fr]{4.1.5}), soient $X$ un
$Y$-préschéma, $q:X\to Y$ le morphisme structural, et soit $r:X\to P$ un
$Y$-\emph{morphisme}, de sorte que l'on a le diagramme commutatif
\[
  \xymatrix{
    P \ar[d]_p & X \ar[l]_r \ar[dl]^q
    \\Y
  }
\]

\oldpage[II]{73}

Comme le foncteur $r^*$ est exact à droite
(\hyperref[0.4.3.1-fr]{0, 4.3.1}), on déduit de l'homomorphisme
(\hyperref[II.4.1.5.1-fr]{4.1.5.1}), qui est surjectif
(\hyperref[II.4.1.6-fr]{4.1.6}), un homomorphisme surjectif
\[
  r^*(\alpha_1^\sharp):r^*(p^*(\mathcal{E}))\longrightarrow
  r^*(\mathcal{O}_P(1)).
\]
Mais $r^*(p^*(\mathcal{E}))=q^*(\mathcal{E})$, et
$r^*(\mathcal{O}_P(1))$ est localement isomorphe à
$r^*(\mathcal{O}_P)=\mathcal{O}_X$, autrement dit c'est un faisceau
inversible $\mathcal{L}_r$ sur $\mathcal{O}_X$ et on a ainsi défini à partir
de $r$ un $\mathcal{O}_X$-homomorphisme canonique surjectif
\[
\label{II.4.2.1.1-fr}
  \varphi_r:q^*(\mathcal{E})\longrightarrow\mathcal{L}_r.
\tag{4.2.1.1}
\]

Lorsque $Y=\operatorname{Spec}(A)$ est affine, et
$\mathcal{E}=\widetilde{E}$ on explicite encore cet homomorphisme de la façon
suivante~: étant donné $f\in E$, il résulte d'abord de
(\hyperref[II.2.6.3-fr]{2.6.3}) que l'on a
\[
\label{II.4.2.1.2-fr}
  r^{-1}(D_+(f))=X_{\varphi_r^\flat(f)}.
\tag{4.2.1.2}
\]
D'autre part, soit $V$ un ouvert affine de $X$ contenu dans
$r^{-1}(D_+(f))$, et soit $B$ son anneau, qui est une $A$-algèbre~; posons
$S=\mathbf{S}_A(E)$~; la restriction de $r$ à $V$ correspond à un
$A$-homomorphisme $\omega:S_{(f)}\to B$, on a
$q^*(\mathcal{E})|V=\widetilde{E\otimes_A B}$ et
$\mathcal{L}_r|V=\widetilde{L_r}$, où
$L_r=(S(1))_{(f)}\otimes_{S_{(f)}}B_{[\omega]}$
(\hyperref[I.1.6.5-fr]{I, 1.6.5}). La restriction de $\varphi_r$ à
$q^*(\mathcal{E})|V$ correspond au $B$-homomorphisme
$u:E\otimes_A B\to L_r$ qui transforme $x\otimes1$ en
$(x/1)\otimes f=(f/1)\otimes\omega(x/f)$. L'extension canonique de
$\varphi_r$ en un homomorphisme de $\mathcal{O}_X$-Algèbres
\[
  \psi_r:q^*(\mathbf{S}(\mathcal{E}))=
  \mathbf{S}(q^*(\mathcal{E}))\longrightarrow
  \mathbf{S}(\mathcal{L}_r)=\bigoplus_{n\geq0}\mathcal{L}_r^{\otimes n}
\]
est donc telle que la restriction de $\psi_r$ à
$q^*(\mathbf{S}_n(\mathcal{E}))|V$ correspond à l'homomorphisme
$\mathbf{S}_n(E\otimes_A B)=\mathbf{S}_n(E)\otimes_A B\to L_r^{\otimes n}$
qui transforme $s\otimes1$ en
$(f/1)^{\otimes n}\otimes\omega(s/f^n)$.
\end{env}

\begin{env}[4.2.2]
\label{II.4.2.2-fr}
Inversement, supposons donnés un morphisme $q:X\to Y$, un
$\mathcal{O}_X$-Module inversible $\mathcal{L}$ et un
$\mathcal{O}_Y$-Module quasi-cohérent $\mathcal{E}$~; à un homomorphisme
$\varphi:q^*(\mathcal{E})\to\mathcal{L}$ correspond canoniquement un
homomorphisme de $\mathcal{O}_X$-Algèbres quasi-cohérentes
\[
  \psi:\mathbf{S}(q^*(\mathcal{E}))=q^*(\mathbf{S}(\mathcal{E}))
  \longrightarrow\bigoplus_{n\geq0}\mathcal{L}^{\otimes n},
\]
et par suite (\hyperref[II.3.7.1-fr]{3.7.1}) un $Y$-morphisme
$r_{\mathcal{L},\psi}:G(\psi)\to
\operatorname{Proj}(\mathbf{S}(\mathcal{E}))=\mathbf{P}(\mathcal{E})$, que
nous noterons $r_{\mathcal{L},\varphi}$. Si $\varphi$ est surjectif, il en est
de même de $\psi$, donc (\hyperref[II.3.7.4-fr]{3.7.4})
$r_{\mathcal{L},\varphi}$ est partout défini. En outre, avec les notations de
(\hyperref[II.4.2.1-fr]{4.2.1}) et de (4.2.2)~:
\end{env}

\begin{proposition}[4.2.3]
\label{II.4.2.3-fr}
Étant donnés un morphisme $q:X\to Y$ et un $\mathcal{O}_Y$-Module
quasi-cohérent $\mathcal{E}$, les applications
$r\mapsto(\mathcal{L}_r,\varphi_r)$ et
$(\mathcal{L},\varphi)\mapsto r_{\mathcal{L},\varphi}$ mettent en
correspondance biunivoque l'ensemble des $Y$-morphismes
$r:X\to\mathbf{P}(\mathcal{E})$ et l'ensemble des classes d'équivalence des
couples $(\mathcal{L},\varphi)$ formés d'un $\mathcal{O}_X$-Module inversible
$\mathcal{L}$ et d'un homomorphisme surjectif
$\varphi:q^*(\mathcal{E})\to\mathcal{L}$, deux couples
$(\mathcal{L},\varphi)$, $(\mathcal{L}',\varphi')$ étant équivalents s'il
existe un $\mathcal{O}_X$-isomorphisme
$\tau:\mathcal{L}\xrightarrow{\sim}\mathcal{L}'$ tel que
$\varphi'=\tau\circ\varphi$.
\end{proposition}

\begin{proof}
Partons d'abord d'un $Y$-morphisme $r:X\to\mathbf{P}(\mathcal{E})$, formons
$\mathcal{L}_r$ et $\varphi_r$ (\hyperref[II.4.2.1-fr]{4.2.1}), puis
$r'=r_{\mathcal{L}_r,\varphi_r}$~; il résulte aussitôt de
(\hyperref[II.4.2.1-fr]{4.2.1}) et de
(\hyperref[II.3.7.2-fr]{3.7.2}) que les morphismes $r$ et $r'$ sont
identiques (en prenant pour générateur de $\mathcal{L}_r$ l'élément
$(f/1)\otimes1$ pour définir les homomorphismes $v_n$ de
(\hyperref[II.3.7.2-fr]{3.7.2})). Inversement, partons d'un couple
$(\mathcal{L},\varphi)$ et formons

\oldpage[II]{74}

$r''=r_{\mathcal{L},\varphi}$, puis $\mathcal{L}_{r''}$ et
$\varphi_{r''}$~; montrons qu'il y a un isomorphisme canonique
$\tau:\mathcal{L}_{r''}\xrightarrow{\sim}\mathcal{L}$ tel que
$\varphi=\tau\circ\varphi_{r''}$~; pour le définir on peut se placer dans le
cas où $Y=\operatorname{Spec}(A)$ et $X=\operatorname{Spec}(B)$ sont affines,
et (avec les notations de (\hyperref[II.4.2.1-fr]{4.2.1}) et de
(\hyperref[II.3.7.2-fr]{3.7.2})) faire correspondre à tout élément
$(x/1)\otimes1$ de $L_{r''}$ (avec $x\in E$) l'élément $v_1(x)c$ de $L$. On
vérifie aussitôt que $\tau$ ne dépend pas du générateur choisi $c$ de $L$~;
comme $v_1$ est par hypothèse surjectif, pour prouver que $\tau$ est un
isomorphisme, il suffit de voir que si $x/1=0$ dans $(S(1))_{(f)}$, on a
$v_1(x)/1=0$ dans $B_g$~; mais la première relation signifie que $f^nx=0$
dans $\mathbf{S}_{n+1}(E)$ pour un certain $n$, et on en déduit que
$v_{n+1}(f^nx)=g^nv_1(x)=0$ dans $B$, d'où la conclusion. Enfin, il est
immédiat que pour deux couples équivalents $(\mathcal{L},\varphi)$,
$(\mathcal{L}',\varphi')$, on a
$r_{\mathcal{L},\varphi}=r_{\mathcal{L}',\varphi'}$.
\end{proof}

En particulier, pour $X=Y$~:

\begin{theorem}[4.2.4]
\label{II.4.2.4-fr}
L'ensemble des $Y$-sections de $\mathbf{P}(\mathcal{E})$ est en
correspondance biunivoque canonique avec l'ensemble des
sous-$\mathcal{O}_Y$-Modules quasi-cohérents $\mathcal{F}$ de $\mathcal{E}$
tels que $\mathcal{E}/\mathcal{F}$ soit inversible.
\end{theorem}

On notera que cette propriété correspond à la définition classique de
l'\og espace projectif\fg{} comme ensemble des hyperplans d'un espace
vectoriel (le cas classique correspondant à $Y=\operatorname{Spec}(K)$, où
$K$ est un corps, et $\mathcal{E}=\widetilde{E}$, $E$ étant un $K$-espace
vectoriel de dimension finie~; les $\mathcal{F}$ ayant la propriété énoncée
dans (\hyperref[II.4.2.4-fr]{4.2.4}) correspondent alors aux hyperplans de
$E$, et on sait d'autre part que les $Y$-sections de
$\mathbf{P}(\mathcal{E})$ sont alors les points de
$\mathbf{P}(\mathcal{E})$ rationnels sur $K$
(\hyperref[I.3.4.5-fr]{I, 3.4.5})).

\begin{remark}[4.2.5]
\label{II.4.2.5-fr}
Comme il y a correspondance biunivoque canonique entre $Y$-morphismes de $X$
dans $P$ et leurs morphismes graphes, $X$-sections de $P\times_YX$
(\hyperref[I.3.3.14-fr]{I, 3.3.14}), on voit qu'inversement
(\hyperref[II.4.2.3-fr]{4.2.3}) peut se déduire de
(\hyperref[II.4.2.4-fr]{4.2.4}). Désignons par
$\operatorname{Hyp}_Y(X,\mathcal{E})$ l'ensemble des
sous-$\mathcal{O}_X$-Modules quasi-cohérents $\mathcal{F}$ de
$\mathcal{E}\otimes_Y\mathcal{O}_X=q^*(\mathcal{E})$ tels que
$q^*(\mathcal{E})/\mathcal{F}$ soit un $\mathcal{O}_X$-Module inversible. Si
$g:X'\to X$ est un $Y$-morphisme, il résulte du fait que $g^*$ est exact à
droite que l'on a
\[
  g^*(q^*(\mathcal{E})/\mathcal{F})=
  g^*(q^*(\mathcal{E}))/g^*(\mathcal{F}),
\]
donc que ce dernier faisceau est inversible, et par suite
$\operatorname{Hyp}_Y(X,\mathcal{E})$ est un foncteur contravariant dans la
catégorie des $Y$-préschémas. On peut alors interpréter le th.
(\hyperref[II.4.2.4-fr]{4.2.4}) comme définissant un isomorphisme canonique
des foncteurs $\operatorname{Hom}_Y(X,\mathbf{P}(\mathcal{E}))$ et
$\operatorname{Hyp}_Y(X,\mathcal{E})$ contravariants en $X$ variable dans la
catégorie des $Y$-préschémas. Ceci fournit aussi une caractérisation du fibré
projectif $P=\mathbf{P}(\mathcal{E})$ par la propriété universelle suivante,
plus proche de l'intuition géométrique que les constructions des \S\S~2 et
3~: pour tout morphisme $q:X\to Y$ et tout $\mathcal{O}_X$-Module inversible
$\mathcal{L}$, quotient de $\mathcal{E}\otimes_{\mathcal{O}_Y}\mathcal{O}_X$,
il existe un $Y$-morphisme unique $r:X\to P$ tel que
$\mathcal{L}=r^*(\mathcal{O}_P(1))$.

On verra plus tard comment on peut de la même manière définir, entre autres,
les schémas \og grassmanniennes\fg{}.
\end{remark}

\begin{corollary}[4.2.6]
\label{II.4.2.6-fr}
Supposons que tout $\mathcal{O}_Y$-Module inversible soit trivial
(\hyperref[I.2.4.8-fr]{I, 2.4.8}). Soit $V$ le groupe
$\operatorname{Hom}_{\mathcal{O}_Y}(\mathcal{E},\mathcal{O}_Y)$, considéré
comme module sur l'anneau $A=\Gamma(Y,\mathcal{O}_Y)$, et soit $V^*$ le
sous-ensemble de $V$ formé des homomorphismes surjectifs. Alors l'ensemble
des $Y$-sections de $\mathbf{P}(\mathcal{E})$ s'identifie canoniquement à
$V^*/A^*$, où $A^*$ est le groupe des unités de $A$.
\end{corollary}

\oldpage[II]{75}

En particulier~:
\begin{enumerate}
  \item[1\textsuperscript{o}] Le corollaire
  (\hyperref[II.4.2.6-fr]{4.2.6}) s'applique lorsque $Y$ est un
  \emph{schéma local} (\hyperref[I.2.4.8-fr]{I, 2.4.8}). Soient $Y$ un
  préschéma quelconque, $y$ un point de $Y$,
  $Y'=\operatorname{Spec}(k(y))$~; la fibre $p^{-1}(y)$ de
  $\mathbf{P}(\mathcal{E})$ est, en vertu de
  (\hyperref[II.4.1.3.1-fr]{4.1.3.1}), identifiée à
  $\mathbf{P}(\mathcal{E}^y)$, où
  $\mathcal{E}^y=\mathcal{E}_y\otimes_{\mathcal{O}_y}k(y)
  =\mathcal{E}_y/\mathfrak{m}_y\mathcal{E}_y$ est considéré comme espace
  vectoriel sur $k(y)$. Plus généralement, si $K$ est une extension de
  $k(y)$, $p^{-1}(y)\otimes_{k(y)}K$ s'identifie à
  $\mathbf{P}(\mathcal{E}^y\otimes_{k(y)}K)$. Le cor.
  (\hyperref[II.4.2.6-fr]{4.2.6}) montre donc que l'ensemble des
  \emph{points géométriques de $\mathbf{P}(\mathcal{E})$ à valeurs dans
  l'extension $K$ de $k(y)$} (\hyperref[I.3.4.5-fr]{I, 3.4.5}), que l'on
  peut encore appeler \emph{fibre géométrique rationnelle sur $K$ de
  $\mathbf{P}(\mathcal{E})$ au-dessus du point $y$}, s'identifie à
  l'\emph{espace projectif} associé au \emph{dual} du $K$-espace vectoriel
  $\mathcal{E}^y\otimes_{k(y)}K$.
  \item[2\textsuperscript{o}] Supposons que $Y$ soit affine d'anneau $A$, et
  en outre que tout $\mathcal{O}_Y$-Module inversible soit trivial~; prenons
  en outre $\mathcal{E}=\mathcal{O}_Y^n$~; alors, dans
  (\hyperref[II.4.2.6-fr]{4.2.6}), $V$ s'identifie à $A^n$
  (\hyperref[I.1.3.8-fr]{I, 1.3.8}), $V^*$ à l'ensemble des systèmes
  $(f_i)_{1\leq i\leq n}$ d'éléments de $A$ engendrant l'idéal $A$~; deux
  tels systèmes définissent la même $Y$-section de
  $\mathbf{P}_Y^{n-1}=\mathbf{P}_A^{n-1}$, autrement dit le même point de
  $\mathbf{P}_A^{n-1}$ à valeurs dans $A$ si et seulement si l'un se déduit
  de l'autre par multiplication par un élément inversible de $A$.
\end{enumerate}

Ces propriétés justifient la terminologie de \og fibré projectif\fg{} pour
$\mathbf{P}(\mathcal{E})$. On notera que la définition que nous obtenons
ainsi de l'\og espace projectif\fg{} est en fait duale de la définition
classique~; cela nous est imposé par la nécessité de pouvoir définir
$\mathbf{P}(\mathcal{E})$ pour un $\mathcal{O}_Y$-Module quasi-cohérent
quelconque $\mathcal{E}$, non nécessairement localement libre.

\begin{remark}[4.2.7]
\label{II.4.2.7-fr}
Nous verrons au chapitre~V que si $Y$ est localement noethérien et connexe,
et si $\mathcal{E}$ est localement libre, alors, en posant
$P=\mathbf{P}(\mathcal{E})$, tout $\mathcal{O}_P$-Module inversible
$\mathcal{L}$ est isomorphe à un $\mathcal{O}_P$-Module de la forme
$\mathcal{L}'\otimes_{\mathcal{O}_Y}\mathcal{O}_P(m)$, où $\mathcal{L}'$ est
un $\mathcal{O}_Y$-Module inversible, bien déterminé à un isomorphisme près,
et $m$ un entier bien déterminé. En d'autres termes,
$\mathrm{H}^1(P,\mathcal{O}_P^*)$ est isomorphe à
$\mathbf{Z}\times\mathrm{H}^1(Y,\mathcal{O}_Y^*)$
(\hyperref[0.5.4.7-fr]{0, 5.4.7}). Nous verrons aussi
(\hyperref[III.2.1.14-fr]{III, 2.1.14}, compte tenu de
(\hyperref[0.5.4.10-fr]{0, 5.4.10})) que
$p_*(\mathcal{L}^{\otimes m})=0$ si $m<0$ et que
$p_*(\mathcal{L}^{\otimes m})$ est isomorphe à
$\mathcal{L}'\otimes_{\mathcal{O}_Y}
(\mathbf{S}_{\mathcal{O}_Y}(\mathcal{E}))_m$ si $m\geq0$. Si $\mathcal{F}$
est un $\mathcal{O}_Y$-Module quasi-cohérent, tout $Y$-morphisme
$\mathbf{P}(\mathcal{E})\to\mathbf{P}(\mathcal{F})$ est donc déterminé par
la donnée d'un $\mathcal{O}_Y$-Module inversible $\mathcal{L}'$, d'un entier
$m\geq0$ et d'un $\mathcal{O}_Y$-homomorphisme
$\psi:\mathcal{F}\to\mathcal{L}'\otimes_{\mathcal{O}_Y}
(\mathbf{S}_{\mathcal{O}_Y}(\mathcal{E}))_m$ tel que l'homomorphisme
correspondant $\psi^\sharp$ de $\mathcal{O}_{\mathbf{P}(\mathcal{F})}$-Modules
soit surjectif. Nous verrons aussi que si le $Y$-morphisme considéré est un
isomorphisme, alors $m=1$ et $\mathcal{F}$ est isomorphe à
$\mathcal{E}\otimes_{\mathcal{O}_Y}\mathcal{L}'$ (réciproque de
(\hyperref[II.4.1.4-fr]{4.1.4})). Ceci permettra de déterminer le faisceau
des germes d'automorphismes de $\mathbf{P}(\mathcal{E})$ comme quotient du
faisceau de groupes $\mathcal{A}ut(\mathcal{E})$ (localement isomorphe à
$\operatorname{GL}(n,\mathcal{O}_Y)$ si $\mathcal{E}$ est de rang $n$) par
$\mathcal{O}_Y^*$.
\end{remark}

\begin{env}[4.2.8]
\label{II.4.2.8-fr}
Conservant les notations de (\hyperref[II.4.2.1-fr]{4.2.1}), soit
$u:X'\to X$ un morphisme~; si le $Y$-morphisme $r:X\to P$ correspond à
l'homomorphisme $\varphi:q^*(\mathcal{E})\to\mathcal{L}$, le $Y$-morphisme
$r\circ u$ correspond à
$u^*(\varphi):u^*(q^*(\mathcal{E}))\to u^*(\mathcal{L})$, comme il résulte
aussitôt des définitions.
\end{env}

\begin{env}[4.2.9]
\label{II.4.2.9-fr}
Soient $\mathcal{E}$, $\mathcal{F}$ deux $\mathcal{O}_Y$-Modules
quasi-cohérents, $v:\mathcal{E}\to\mathcal{F}$ un homomorphisme surjectif,
$j=\mathbf{P}(v)$ l'immersion fermée correspondante
$\mathbf{P}(\mathcal{F})\to\mathbf{P}(\mathcal{E})$
(\hyperref[II.4.1.2-fr]{4.1.2}). Si le $Y$-morphisme
$r:X\to\mathbf{P}(\mathcal{F})$ correspond à l'homomorphisme
$\varphi:q^*(\mathcal{F})\to\mathcal{L}$, le

\oldpage[II]{76}

$Y$-morphisme $j\circ r$ correspond à
\[
  q^*(\mathcal{E})\xrightarrow{q^*(v)}q^*(\mathcal{F})
  \xrightarrow{\varphi}\mathcal{L}~;
\]
cela résulte encore de la définition donnée en
(\hyperref[II.4.2.1-fr]{4.2.1}).
\end{env}

\begin{env}[4.2.10]
\label{II.4.2.10-fr}
Soit $\psi:Y'\to Y$ un morphisme, et posons
$\mathcal{E}'=\psi^*(\mathcal{E})$. Si le $Y$-morphisme $r:X\to P$ correspond
à l'homomorphisme $\varphi:q^*(\mathcal{E})\to\mathcal{L}$, le
$Y'$-morphisme
\[
  r_{(Y')}:X_{(Y')}\longrightarrow P'=\mathbf{P}(\mathcal{E}')
\]
correspond à
$\varphi_{(Y')}:q_{(Y')}^*(\mathcal{E}')=
  q^*(\mathcal{E})\otimes_Y\mathcal{O}_{Y'}\to
\mathcal{L}\otimes_Y\mathcal{O}_{Y'}$. En effet, en vertu de
(\hyperref[II.4.1.3.1-fr]{4.1.3.1}), on a le diagramme commutatif
\[
  \xymatrix{
    Y' \ar[d]
    & P'=P_{(Y')} \ar[l]_{p_{(Y')}} \ar[d]^{u}
    & X_{(Y')} \ar[l]_{r_{(Y')}} \ar[d]^{v}
    \\Y
    & P \ar[l]_{p}
    & X \ar[l]_{r}
  }
\]
D'après (\hyperref[II.4.1.3.2-fr]{4.1.3.2}), on a
\[
  (r_{(Y')})^*(\mathcal{O}_{P'}(1))
  =(r_{(Y')})^*(u^*(\mathcal{O}_P(1)))
  =v^*(r^*(\mathcal{O}_P(1)))
  =v^*(\mathcal{L})=\mathcal{L}\otimes_Y\mathcal{O}_{Y'}~;
\]
d'autre part, $u^*(\alpha_1^\sharp)$ n'est autre que l'homomorphisme
canonique
$\alpha_1^\sharp:(p_{(Y')})^*(\mathcal{E}')\to\mathcal{O}_{P'}(1)$~; cela
se voit en explicitant les homomorphismes canoniques $\alpha_1^\sharp$ dans
$P$ et $P'$ comme dans (\hyperref[II.4.1.6-fr]{4.1.6}). D'où notre assertion.
\end{env}

\subsection{Le morphisme de Segre.}
\label{subsection:II.4.3-fr}

\begin{env}[4.3.1]
\label{II.4.3.1-fr}
Soient $Y$ un préschéma, $\mathcal{E}$, $\mathcal{F}$ deux
$\mathcal{O}_Y$-Modules quasi-cohérents~; posons
$P_1=\mathbf{P}(\mathcal{E})$, $P_2=\mathbf{P}(\mathcal{F})$, et désignons
par $p_1:P_1\to Y$, $p_2:P_2\to Y$ les morphismes structuraux. Soit
$Q=P_1\times_YP_2$, et soient $q_1:Q\to P_1$, $q_2:Q\to P_2$ les
projections canoniques~; le $\mathcal{O}_Q$-Module
\[
  \mathcal{L}=\mathcal{O}_{P_1}(1)\otimes_Y\mathcal{O}_{P_2}(1)
  :=q_1^*(\mathcal{O}_{P_1}(1))\otimes_{\mathcal{O}_Q}
  q_2^*(\mathcal{O}_{P_2}(1))
\]
est inversible comme produit tensoriel de deux $\mathcal{O}_Q$-Modules
inversibles (\hyperref[0.5.4.4-fr]{0, 5.4.4}). D'autre part, si
$r=p_1\circ q_1=p_2\circ q_2$ est le morphisme structural $Q\to Y$, on a
\[
  r^*(\mathcal{E}\otimes_{\mathcal{O}_Y}\mathcal{F})
  =q_1^*(p_1^*(\mathcal{E}))\otimes_{\mathcal{O}_Q}
  q_2^*(p_2^*(\mathcal{F}))
\]
(\hyperref[0.4.3.3-fr]{0, 4.3.3})~; les homomorphismes canoniques surjectifs
(\hyperref[II.4.1.5.1-fr]{4.1.5.1})
$p_1^*(\mathcal{E})\to\mathcal{O}_{P_1}(1)$,
$p_2^*(\mathcal{F})\to\mathcal{O}_{P_2}(1)$ donnent donc, par produit
tensoriel, un homomorphisme canonique
\[
\label{II.4.3.1.1-fr}
  s:r^*(\mathcal{E}\otimes_{\mathcal{O}_Y}\mathcal{F})\longrightarrow
  \mathcal{L}
\tag{4.3.1.1}
\]
évidemment surjectif~; on en déduit donc
(\hyperref[II.4.2.2-fr]{4.2.2}) un morphisme canonique, dit
\emph{morphisme de Segre}~:
\[
\label{II.4.3.1.2-fr}
  \varsigma:\mathbf{P}(\mathcal{E})\times_Y\mathbf{P}(\mathcal{F})
  \longrightarrow
  \mathbf{P}(\mathcal{E}\otimes_{\mathcal{O}_Y}\mathcal{F}).
\tag{4.3.1.2}
\]
Explicitons le morphisme $\varsigma$, lorsque $Y=\operatorname{Spec}(A)$ est
affine, $\mathcal{E}=\widetilde{E}$, $\mathcal{F}=\widetilde{F}$, $E$ et $F$
étant deux $A$-modules, d'où
$\mathcal{E}\otimes_{\mathcal{O}_Y}\mathcal{F}=\widetilde{E\otimes_A F}$
(\hyperref[I.1.3.12-fr]{I, 1.3.12})~; posons
$R=\mathbf{S}_A(E)$, $S=\mathbf{S}_A(F)$,
$T=\mathbf{S}_A(E\otimes_A F)$~; soient $f\in E$, $g\in F$, et considérons
l'ouvert affine
\[
  D_+(f)\times_YD_+(g)=\operatorname{Spec}(B)
\]

\oldpage[II]{77}

de $Q$, où $B=R_{(f)}\otimes_AS_{(g)}$~; la restriction de $\mathcal{L}$ à
cet ouvert affine est $\widetilde{L}$, où
\[
  L=(R(1))_{(f)}\otimes_A(S(1))_{(g)}
\]
et l'élément $c=(f/1)\otimes(g/1)$ est un générateur de $L$ considéré comme
$B$-module libre (\hyperref[II.2.5.7-fr]{2.5.7}). L'homomorphisme
(\hyperref[II.4.3.1.1-fr]{4.3.1.1}) correspond à l'homomorphisme
\[
  (x\otimes y)\otimes b\longmapsto b((x/1)\otimes(y/1))
\]
de $(E\otimes_A F)\otimes_AB$ dans $L$. Avec les notations de
(\hyperref[II.3.7.2-fr]{3.7.2}), on a donc
$v_1(x\otimes y)=(x/f)\otimes(y/g)$~; la restriction de $\varsigma$ à
$D_+(f)\times_YD_+(g)$ est un morphisme de ce schéma affine dans
$D_+(f\otimes g)$, correspondant à l'homomorphisme d'anneaux
$\omega:T_{(f\otimes g)}\to R_{(f)}\otimes_AS_{(g)}$ tel que
\[
\label{II.4.3.1.3-fr}
  \omega((x\otimes y)/(f\otimes g))=(x/f)\otimes(y/g)
\tag{4.3.1.3}
\]
pour $x\in E$ et $y\in F$.
\end{env}

\begin{env}[4.3.2]
\label{II.4.3.2-fr}
Il résulte de (\hyperref[II.4.2.3-fr]{4.2.3}) que l'on a un isomorphisme
canonique
\[
\label{II.4.3.2.1-fr}
  \tau:\varsigma^*(\mathcal{O}_P(1))\xrightarrow{\sim}
  \mathcal{O}_{P_1}(1)\otimes_Y\mathcal{O}_{P_2}(1)
\tag{4.3.2.1}
\]
où on a posé
$P=\mathbf{P}(\mathcal{E}\otimes_{\mathcal{O}_Y}\mathcal{F})$. Montrons
que, pour $x\in\Gamma(Y,\mathcal{E})$ et
$y\in\Gamma(Y,\mathcal{F})$, on a
\[
\label{II.4.3.2.2-fr}
  \tau(\alpha_1(x\otimes y))=\alpha_1(x)\otimes\alpha_1(y).
\tag{4.3.2.2}
\]
En effet, on est ramené au cas où $Y$ est affine, et on a alors, avec les
notations de (\hyperref[II.4.3.1-fr]{4.3.1}) et de
(\hyperref[II.2.6.2-fr]{2.6.2}),
$\alpha_1^{f\otimes g}(x\otimes y)=(x\otimes y)/1$ dans
$(T(1))_{(f\otimes g)}$, $\alpha_1^f(x)=x/1$ dans $(R(1))_{(f)}$ et
$\alpha_1^g(y)=y/1$ dans $(S(1))_{(g)}$. La définition de $\tau$ donnée dans
(\hyperref[II.4.2.3-fr]{4.2.3}) et le calcul de $v_1$ fait dans
(\hyperref[II.4.3.1-fr]{4.3.1}) prouvent alors aussitôt
(\hyperref[II.4.3.2.2-fr]{4.3.2.2}). On en déduit la formule
\[
\label{II.4.3.2.3-fr}
  \varsigma^{-1}(P_{x\otimes y})=(P_1)_x\times_Y(P_2)_y
\tag{4.3.2.3}
\]
avec les notations de (\hyperref[II.3.1.4-fr]{3.1.4}). En effet, compte tenu
de (\hyperref[II.3.3.3-fr]{3.3.3}), la formule
(\hyperref[II.4.3.2.2-fr]{4.3.2.2}) ramène (en revenant au cas affine, à
l'aide de (\hyperref[I.3.2.7-fr]{I, 3.2.7} et
(\hyperref[I.3.2.3-fr]{3.2.3})) à prouver le lemme suivant~:

\begin{lemma}[4.3.2.4]
\label{II.4.3.2.4-fr}
Soient $B$, $B'$ deux $A$-algèbres, et soient
$Y=\operatorname{Spec}(A)$, $Z=\operatorname{Spec}(B)$,
$Z'=\operatorname{Spec}(B')$~; quels que soient $t\in B$, $t'\in B'$, on a
$D(t\otimes t')=D(t)\times_YD(t')$.
\end{lemma}

\begin{proof}
En effet, si $p:Z\times_YZ'\to Z$ et $p':Z\times_YZ'\to Z'$ sont les
projections canoniques, il résulte de
(\hyperref[I.1.2.2.2-fr]{I, 1.2.2.2}) que l'on a
$p^{-1}(D(t))=D(t\otimes1)$ et
$p'^{-1}(D(t'))=D(1\otimes t')$~; la conclusion résulte de
(\hyperref[I.3.2.7-fr]{I, 3.2.7}) et
(\hyperref[I.1.1.9.1-fr]{I, 1.1.9.1}), puisque
$(t\otimes1)(1\otimes t')=t\otimes t'$.
\end{proof}
\end{env}

\begin{proposition}[4.3.3]
\label{II.4.3.3-fr}
Le morphisme de Segre est une immersion fermée.
\end{proposition}

\begin{proof}
En effet, la question étant locale sur $Y$, on est ramené au cas où $Y$ est
affine. Avec les notations de (\hyperref[II.4.3.1-fr]{4.3.1}) et
(\hyperref[II.4.3.2-fr]{4.3.2}), les $D_+(f\otimes g)$ forment alors une base
de la topologie de $P$, puisque les $f\otimes g$ engendrent $T$ lorsque $f$
parcourt $E$ et $g$ parcourt $F$. D'autre part, on a en vertu de
(\hyperref[II.4.3.2.3-fr]{4.3.2.3})
$\varsigma^{-1}(D_+(f\otimes g))=D_+(f)\times_YD_+(g)$. Il suffit donc
(\hyperref[I.4.2.4-fr]{I, 4.2.4}) de prouver que la restriction de
$\varsigma$ à $D_+(f)\times_YD_+(g)$ est une immersion fermée dans
$D_+(f\otimes g)$. Or, avec les mêmes notations, la formule
(\hyperref[II.4.3.1.3-fr]{4.3.1.3}) montre que $\omega$ est surjectif, ce qui
termine la démonstration.
\end{proof}

\begin{env}[4.3.4]
\label{II.4.3.4-fr}
Le morphisme de Segre est fonctoriel en $\mathcal{E}$ et $\mathcal{F}$,
lorsqu'on restreint les

\oldpage[II]{78}

homomorphismes de $\mathcal{O}_Y$-Modules quasi-cohérents aux homomorphismes
surjectifs. Il faut en effet montrer que si
$\mathcal{E}\to\mathcal{E}'$ est un $\mathcal{O}_Y$-homomorphisme surjectif,
le diagramme
\[
  \xymatrix{
    \mathbf{P}(\mathcal{E}')\times_Y\mathbf{P}(\mathcal{F})
      \ar[r]^{j\times1} \ar[d]_{\varsigma}
    & \mathbf{P}(\mathcal{E})\times_Y\mathbf{P}(\mathcal{F})
      \ar[d]^{\varsigma}
    \\\mathbf{P}(\mathcal{E}'\otimes\mathcal{F}) \ar[r]
    & \mathbf{P}(\mathcal{E}\otimes\mathcal{F})
  }
\]
est commutatif, $j$ désignant l'immersion fermée canonique
$\mathbf{P}(\mathcal{E}')\to\mathbf{P}(\mathcal{E})$. Posons
$P_1'=\mathbf{P}(\mathcal{E}')$ et gardons par ailleurs les notations de
(\hyperref[II.4.3.1-fr]{4.3.1})~; $j\times1$ est une immersion fermée
(\hyperref[I.4.3.1-fr]{I, 4.3.1}) et on a, à des isomorphismes près,
\[
  (j\times1)^*(\mathcal{O}_{P_1}(1)\otimes\mathcal{O}_{P_2}(1))
  =j^*(\mathcal{O}_{P_1}(1))\otimes\mathcal{O}_{P_2}(1)
  =\mathcal{O}_{P_1'}(1)\otimes\mathcal{O}_{P_2}(1)
\]
en vertu de (\hyperref[II.4.1.2.1-fr]{4.1.2.1}) et de
(\hyperref[I.9.1.5-fr]{I, 9.1.5})~; notre assertion résulte donc de
(\hyperref[II.4.2.8-fr]{4.2.8}) et
(\hyperref[II.4.2.9-fr]{4.2.9}).
\end{env}

\begin{env}[4.3.5]
\label{II.4.3.5-fr}
Avec les notations de (\hyperref[II.4.3.1-fr]{4.3.1}), soit
$\psi:Y'\to Y$ un morphisme, et posons
$\mathcal{E}'=\psi^*(\mathcal{E})$, $\mathcal{F}'=\psi^*(\mathcal{F})$~;
alors le morphisme de Segre
$\mathbf{P}(\mathcal{E}')\times\mathbf{P}(\mathcal{F}')
\to\mathbf{P}(\mathcal{E}'\otimes\mathcal{F}')$ s'identifie à
$\varsigma_{(Y')}$. En effet, gardant les notations de
(\hyperref[II.4.3.1-fr]{4.3.1}), posons par ailleurs
$P_1'=\mathbf{P}(\mathcal{E}')$, $P_2'=\mathbf{P}(\mathcal{F}')$~; on sait
(\hyperref[II.4.1.3.1-fr]{4.1.3.1}) que $P_i'$ s'identifie à $(P_i)_{(Y')}$
($i=1,2$), donc le morphisme structural $P_1'\times P_2'\to Y'$ s'identifie
à $r_{(Y')}$. D'autre part,
$\mathcal{E}'\otimes\mathcal{F}'$ s'identifie à
$\psi^*(\mathcal{E}\otimes\mathcal{F})$, donc
$\mathbf{P}(\mathcal{E}'\otimes\mathcal{F}')$ s'identifie à
$(\mathbf{P}(\mathcal{E}\otimes\mathcal{F}))_{(Y')}$
(\hyperref[II.4.1.3.1-fr]{4.1.3.1}). Enfin,
$\mathcal{O}_{P_1'}(1)\otimes_{Y'}\mathcal{O}_{P_2'}(1)=\mathcal{L}'$
s'identifie à $\mathcal{L}\otimes_Y\mathcal{O}_{Y'}$, en vertu de
(\hyperref[II.4.1.3.2-fr]{4.1.3.2}) et de
(\hyperref[I.9.1.11-fr]{I, 9.1.11}). L'homomorphisme canonique
$(r_{(Y')})^*(\mathcal{E}'\otimes\mathcal{F}')\to\mathcal{L}'$ s'identifie
alors à $s_{(Y')}$, et notre assertion résulte de
(\hyperref[II.4.2.10-fr]{4.2.10}).
\end{env}

\begin{remark}[4.3.6]
\label{II.4.3.6-fr}
Le préschéma somme de $\mathbf{P}(\mathcal{E})$ et
$\mathbf{P}(\mathcal{F})$ est de même canoniquement isomorphe à un
sous-préschéma fermé de $\mathbf{P}(\mathcal{E}\oplus\mathcal{F})$. En
effet, les homomorphismes surjectifs
$\mathcal{E}\oplus\mathcal{F}\to\mathcal{E}$,
$\mathcal{E}\oplus\mathcal{F}\to\mathcal{F}$ correspondent à des immersions
fermées
$\mathbf{P}(\mathcal{E})\to\mathbf{P}(\mathcal{E}\oplus\mathcal{F})$,
$\mathbf{P}(\mathcal{F})\to\mathbf{P}(\mathcal{E}\oplus\mathcal{F})$~;
tout revient à voir que les espaces sous-jacents des sous-préschémas fermés
de $\mathbf{P}(\mathcal{E}\oplus\mathcal{F})$ ainsi obtenus sont sans point
commun. La question étant locale sur $Y$, on peut adopter les notations de
(\hyperref[II.4.3.1-fr]{4.3.1})~; or $\mathbf{S}_n(E)$ et
$\mathbf{S}_n(F)$ s'identifient à des sous-modules de
$\mathbf{S}_n(E\oplus F)$ d'intersection réduite à $0$~; et si
$\mathfrak{p}$ est un idéal premier gradué de $\mathbf{S}(E)$ tel que
$\mathfrak{p}\cap\mathbf{S}_n(E)\neq\mathbf{S}_n(E)$ pour tout $n\geq0$,
il lui correspond dans $\mathbf{S}(E\oplus F)$ un idéal premier gradué dont
la trace sur $\mathbf{S}_n(E)$ est
$\mathfrak{p}\cap\mathbf{S}_n(E)$, mais qui contient $\mathbf{S}_+(F)$,
comme on le voit aussitôt~; deux points de
$\operatorname{Proj}(\mathbf{S}(E))$ et
$\operatorname{Proj}(\mathbf{S}(F))$ respectivement ne peuvent donc avoir
même image dans $\operatorname{Proj}(\mathbf{S}(E\oplus F))$.
\end{remark}

\subsection{Immersions dans les fibrés projectifs. Faisceaux très amples}
\label{subsection:II.4.4-fr}

\begin{proposition}[4.4.1]
\label{II.4.4.1-fr}
Soient $Y$ un schéma quasi-compact, ou un préschéma dont l'espace sous-jacent
est noethérien, $q:X\to Y$ un morphisme de type fini, $\mathcal{L}$ un
$\mathcal{O}_X$-Module inversible.
\begin{enumerate}
  \item[\rm{(i)}] Soient $\mathcal{S}$ une $\mathcal{O}_Y$-Algèbre graduée
  quasi-cohérente à degrés positifs,
  $\psi:q^*(\mathcal{S})\to\bigoplus_{n\geq0}\mathcal{L}^{\otimes n}$ un
  homomorphisme d'Algèbres graduées. Pour que $r_{\mathcal{L},\psi}$ soit
  partout défini et soit une immersion, il faut et

\oldpage[II]{79}

  il suffit qu'il existe un entier $n\geq0$, un sous-$\mathcal{O}_Y$-Module
  quasi-cohérent de type fini $\mathcal{E}$ de $\mathcal{S}_n$, tels que
  l'homomorphisme
  \[
    \psi'=\psi_n\circ q^*(j):q^*(\mathcal{E})
      \longrightarrow\mathcal{L}^{\otimes n}=\mathcal{L}'
  \]
  ($j$ étant l'injection $\mathcal{E}\to\mathcal{S}_n$) soit surjectif et que
  le morphisme
  $r_{\mathcal{L}',\psi'}:X\to\mathbf{P}(\mathcal{E})$ soit une immersion.
  \item[\rm{(ii)}] Soient $\mathcal{F}$ un $\mathcal{O}_Y$-Module
  quasi-cohérent, $\varphi:q^*(\mathcal{F})\to\mathcal{L}$ un homomorphisme
  surjectif. Pour que le morphisme $r_{\mathcal{L},\varphi}$ soit une
  immersion $X\to\mathbf{P}(\mathcal{F})$, il faut et il suffit qu'il existe
  un sous-$\mathcal{O}_Y$-Module quasi-cohérent de type fini $\mathcal{E}$ de
  $\mathcal{F}$ tel que l'homomorphisme
  $\varphi'=\varphi\circ q^*(j):q^*(\mathcal{E})\to\mathcal{L}$ (où
  $j:\mathcal{E}\to\mathcal{F}$ est l'injection canonique) soit surjectif et
  que $r_{\mathcal{L},\varphi'}:X\to\mathbf{P}(\mathcal{E})$ soit une
  immersion.
\end{enumerate}
\end{proposition}

\begin{proof}
\begin{enumerate}
  \item[\rm{(i)}] Le fait que $r_{\mathcal{L},\psi}$ soit partout défini et
  soit une immersion équivaut d'après
  (\hyperref[II.3.8.5-fr]{3.8.5}) à l'existence de $n\geq0$ et de
  $\mathcal{E}$ tels que, si $\mathcal{S}'$ est la sous-algèbre de
  $\mathcal{S}$ engendrée par $\mathcal{E}$, l'homomorphisme
  $q^*(\mathcal{E})\to\mathcal{L}^{\otimes n}$ soit surjectif et que
  $r_{\mathcal{L}',\psi'}:X\to\operatorname{Proj}(\mathcal{S}')$ soit partout
  défini et soit une immersion. On a d'autre part un homomorphisme canonique
  surjectif $\mathbf{S}(\mathcal{E})\to\mathcal{S}'$ auquel correspond une
  immersion fermée
  $\operatorname{Proj}(\mathcal{S}')\to\mathbf{P}(\mathcal{E})$
  (\hyperref[II.3.6.2-fr]{3.6.2})~; d'où la conclusion.
  \item[\rm{(ii)}] Comme $\mathcal{F}$ est limite inductive de ses
  sous-Modules quasi-cohérents de type fini $\mathcal{E}_\lambda$
  (\hyperref[I.9.4.9-fr]{I, 9.4.9}), $\mathbf{S}(\mathcal{F})$ est limite
  inductive des $\mathbf{S}(\mathcal{E}_\lambda)$~; la conclusion résulte
  alors de (\hyperref[II.3.8.4-fr]{3.8.4}), en observant que, dans la
  démonstration de (\hyperref[II.3.8.4-fr]{3.8.4}), on peut prendre tous les
  $n_i$ égaux à $1$~: en effet ($Y$ étant supposé affine), si
  $r=r_{\mathcal{L},\varphi}$ est une immersion, $r(X)$ est un sous-espace
  quasi-compact de $\mathbf{P}(\mathcal{F})$ que l'on peut recouvrir par un
  nombre fini d'ouverts de $\mathbf{P}(\mathcal{F})$ de la forme $D_+(f)$ avec
  $f\in F$, tels que $D_+(f)\cap r(X)$ soit fermé.
\end{enumerate}
\end{proof}

\begin{definition}[4.4.2]
\label{II.4.4.2-fr}
Soient $Y$ un préschéma, $q:X\to Y$ un morphisme. On dit qu'un
$\mathcal{O}_X$-Module inversible $\mathcal{L}$ est \emph{très ample pour
$q$} (ou \emph{très ample pour $Y$}, ou simplement \emph{très ample} si
aucune confusion n'en résulte) s'il existe un $\mathcal{O}_Y$-Module
quasi-cohérent $\mathcal{E}$ et une $Y$-immersion $i$ de $X$ dans
$P=\mathbf{P}(\mathcal{E})$ tels que $\mathcal{L}$ soit isomorphe à
$i^*(\mathcal{O}_P(1))$.
\end{definition}

Il revient au même (\hyperref[II.4.2.3-fr]{4.2.3}) de dire qu'il existe un
$\mathcal{O}_Y$-Module quasi-cohérent $\mathcal{E}$ et un homomorphisme
surjectif $\varphi:q^*(\mathcal{E})\to\mathcal{L}$ tels que
$r_{\mathcal{L},\varphi}:X\to\mathbf{P}(\mathcal{E})$ soit une immersion.

On notera que l'existence d'un $\mathcal{O}_X$-Module très ample relativement
à $Y$ implique que $q$ est séparé
((\hyperref[II.3.1.3-fr]{3.1.3}) et
(\hyperref[I.5.5.1-fr]{I, 5.5.1, (i) et (ii)})).

\begin{corollary}[4.4.3]
\label{II.4.4.3-fr}
Supposons qu'il existe une $\mathcal{O}_Y$-Algèbre graduée quasi-cohérente
$\mathcal{S}$ engendrée par $\mathcal{S}_1$, et une $Y$-immersion
$i:X\to P=\operatorname{Proj}(\mathcal{S})$ telle que $\mathcal{L}$ soit
isomorphe à $i^*(\mathcal{O}_P(1))$~; alors $\mathcal{L}$ est très ample
relativement à $q$.
\end{corollary}

\begin{proof}
En effet, si $\mathcal{F}=\mathcal{S}_1$, l'homomorphisme canonique
$\mathbf{S}(\mathcal{F})\to\mathcal{S}$ est surjectif, donc, en composant
l'immersion fermée correspondante
$\operatorname{Proj}(\mathcal{S})\to\mathbf{P}(\mathcal{F})$
(\hyperref[II.3.6.2-fr]{3.6.2}) et l'immersion $i$, on obtient une immersion
$j:X\to\mathbf{P}(\mathcal{F})=P'$ telle que $\mathcal{L}$ soit isomorphe à
$j^*(\mathcal{O}_{P'}(1))$ (\hyperref[II.3.6.3-fr]{3.6.3}).
\end{proof}

\begin{proposition}[4.4.4]
\label{II.4.4.4-fr}
Soient $q:X\to Y$ un morphisme quasi-compact, $\mathcal{L}$ un
$\mathcal{O}_X$-Module inversible. Les propriétés suivantes sont
équivalentes~:
\begin{enumerate}
  \item[\rm{a)}] $\mathcal{L}$ est très ample relativement à $q$.
  \item[\rm{b)}] $q_*(\mathcal{L})$ est quasi-cohérent, l'homomorphisme
  canonique $\sigma:q^*(q_*(\mathcal{L}))\to\mathcal{L}$ est surjectif, et le
  morphisme
  $r_{\mathcal{L},\sigma}:X\to\mathbf{P}(q_*(\mathcal{L}))$ est une immersion.
\end{enumerate}
\end{proposition}

\begin{proof}
Comme $q$ est quasi-compact, on sait que $q_*(\mathcal{L})$ est quasi-cohérent
si $q$ est séparé (\hyperref[I.9.2.2-fr]{I, 9.2.2}).

\oldpage[II]{80}

On sait (\hyperref[II.3.4.7-fr]{3.4.7}) que l'existence d'un homomorphisme
surjectif $\varphi:q^*(\mathcal{E})\to\mathcal{L}$ ($\mathcal{E}$ étant un
$\mathcal{O}_Y$-Module quasi-cohérent) entraîne que $\sigma$ est surjectif~;
en outre, à la factorisation
$q^*(\mathcal{E})\to q^*(q_*(\mathcal{L}))\xrightarrow{\sigma}\mathcal{L}$
de $\varphi$ correspond canoniquement une factorisation
\[
  q^*(\mathbf{S}(\mathcal{E}))
    \longrightarrow q^*(\mathbf{S}(q_*(\mathcal{L})))
    \longrightarrow\bigoplus_{n\geq0}\mathcal{L}^{\otimes n}
\]
donc (\hyperref[II.3.8.3-fr]{3.8.3}) l'hypothèse que
$r_{\mathcal{L},\varphi}$ est une immersion entraîne qu'il en est de même de
$j=r_{\mathcal{L},\sigma}$~; en outre
(\hyperref[II.4.2.4-fr]{4.2.4}), $\mathcal{L}$ est isomorphe à
$j^*(\mathcal{O}_{P'}(1))$ en posant
$P'=\mathbf{P}(q_*(\mathcal{L}))$. On voit donc que a) et b) sont
équivalents.
\end{proof}

\begin{corollary}[4.4.5]
\label{II.4.4.5-fr}
Supposons $q$ quasi-compact. Pour que $\mathcal{L}$ soit très ample
relativement à $Y$, il faut et il suffit qu'il existe un recouvrement ouvert
$(U_\alpha)$ de $Y$ tel que
$\mathcal{L}|q^{-1}(U_\alpha)$ soit très ample relativement à $U_\alpha$ pour
tout $\alpha$.
\end{corollary}

\begin{proof}
En effet, la condition b) de (\hyperref[II.4.4.4-fr]{4.4.4}) est locale sur
$Y$.
\end{proof}

\begin{proposition}[4.4.6]
\label{II.4.4.6-fr}
Soient $Y$ un schéma quasi-compact, ou un préschéma dont l'espace sous-jacent
est noethérien, $q:X\to Y$ un morphisme de type fini, $\mathcal{L}$ un
$\mathcal{O}_X$-Module inversible. Les conditions a) et b) de
(\hyperref[II.4.4.4-fr]{4.4.4}) sont alors équivalentes aux suivantes~:
\begin{enumerate}
  \item[\rm{a')}] Il existe un $\mathcal{O}_Y$-Module quasi-cohérent de type
  fini $\mathcal{E}$ et un homomorphisme surjectif
  $\varphi:q^*(\mathcal{E})\to\mathcal{L}$ tel que
  $r_{\mathcal{L},\varphi}$ soit une immersion.
  \item[\rm{b')}] Il existe un sous-$\mathcal{O}_Y$-Module cohérent de type
  fini $\mathcal{E}$ de $q_*(\mathcal{L})$ ayant les propriétés énoncées dans
  a').
\end{enumerate}
\end{proposition}

\begin{proof}
Il est clair que a') ou b') implique a)~; d'autre part a) implique a') en
vertu de (\hyperref[II.4.4.1-fr]{4.4.1}) et de même b) implique b').
\end{proof}

\begin{corollary}[4.4.7]
\label{II.4.4.7-fr}
Supposons que $Y$ soit un schéma quasi-compact ou un préschéma noethérien. Si
$\mathcal{L}$ est très ample pour $q$, il existe une $\mathcal{O}_Y$-Algèbre
graduée quasi-cohérente $\mathcal{S}$, telle que $\mathcal{S}_1$ soit de type
fini et engendre $\mathcal{S}$, et une $Y$-immersion ouverte dominante
$i:X\to P=\operatorname{Proj}(\mathcal{S})$ telle que $\mathcal{L}$ soit
isomorphe à $i^*(\mathcal{O}_P(1))$.
\end{corollary}

\begin{proof}
En effet, la condition b') de (\hyperref[II.4.4.6-fr]{4.4.6}) est vérifiée~;
le morphisme structural $p:\mathbf{P}(\mathcal{E})=P'\to Y$ est alors séparé
et de type fini (\hyperref[II.3.1.3-fr]{3.1.3}), donc $P'$ est un schéma
quasi-compact (resp. un préschéma noethérien) si $Y$ est un schéma
quasi-compact (resp. un préschéma noethérien). Soit $Z$ l'adhérence
(\hyperref[I.9.5.11-fr]{I, 9.5.11}) du sous-préschéma $X'$ de $P'$ associé à
l'immersion $j=r_{\mathcal{L},\varphi}$ de $X$ dans $P'$~; il est clair que
$j$ se factorise en l'injection canonique $Z\to P'$ et en une immersion
ouverte dominante $i:X\to Z$. Mais $Z$ s'identifie à un préschéma
$\operatorname{Proj}(\mathcal{S})$, où $\mathcal{S}$ est une
$\mathcal{O}_Y$-Algèbre graduée quotient de $\mathbf{S}(\mathcal{E})$ par un
faisceau quasi-cohérent gradué d'idéaux
(\hyperref[II.3.6.2-fr]{3.6.2}), et il est clair que $\mathcal{S}_1$ est de
type fini et engendre $\mathcal{S}$~; en outre $\mathcal{O}_Z(1)$ est l'image
réciproque de $\mathcal{O}_{P'}(1)$ par l'injection canonique
(\hyperref[II.3.6.3-fr]{3.6.3}), donc
$\mathcal{L}=i^*(\mathcal{O}_Z(1))$.
\end{proof}

\begin{proposition}[4.4.8]
\label{II.4.4.8-fr}
Soient $q:X\to Y$ un morphisme, $\mathcal{L}$ un $\mathcal{O}_X$-Module très
ample relativement à $q$, $\mathcal{L}'$ un $\mathcal{O}_X$-Module inversible
tel qu'il existe un $\mathcal{O}_Y$-Module quasi-cohérent $\mathcal{E}'$ et un
homomorphisme surjectif $q^*(\mathcal{E}')\to\mathcal{L}'$. Alors
$\mathcal{L}\otimes_{\mathcal{O}_X}\mathcal{L}'$ est très ample relativement
à $q$.
\end{proposition}

\begin{proof}
En effet, l'hypothèse entraîne l'existence d'un $Y$-morphisme
$r':X\to P'=\mathbf{P}(\mathcal{E}')$ tel que
$\mathcal{L}'=r'^*(\mathcal{O}_{P'}(1))$
(\hyperref[II.4.2.1-fr]{4.2.1}). Il y a par hypothèse un
$\mathcal{O}_Y$-Module quasi-cohérent $\mathcal{E}$ et une

\oldpage[II]{81}

$Y$-immersion $r:X\to P=\mathbf{P}(\mathcal{E})$ telle que
$\mathcal{L}=r^*(\mathcal{O}_P(1))$. Posons
$Q=\mathbf{P}(\mathcal{E}\otimes\mathcal{E}')$ et considérons le morphisme de
Segre $\varsigma:P\times_YP'\to Q$ (\hyperref[II.4.3.1-fr]{4.3.1}). Comme
$r$ est une immersion, il en est de même de
$(r,r')_Y:X\to P\times_YP'$ (\hyperref[I.5.3.14-fr]{I, 5.3.14})~; mais
$\varsigma$ étant une immersion (\hyperref[II.4.3.3-fr]{4.3.3}), il en est de
même de
$r'':X\xrightarrow{(r,r')_Y}P\times_YP'\xrightarrow{\varsigma}Q$. D'autre
part (\hyperref[II.4.3.2.1-fr]{4.3.2.1}),
$\varsigma^*(\mathcal{O}_Q(1))$ est isomorphe à
$\mathcal{O}_P(1)\otimes_Y\mathcal{O}_{P'}(1)$, donc
(\hyperref[I.9.1.4-fr]{I, 9.1.4}), $r''^*(\mathcal{O}_Q(1))$ est isomorphe à
$\mathcal{L}\otimes\mathcal{L}'$, ce qui démontre la proposition.
\end{proof}

\begin{corollary}[4.4.9]
\label{II.4.4.9-fr}
Soit $q:X\to Y$ un morphisme.
\begin{enumerate}
  \item[\rm{(i)}] Soient $\mathcal{L}$ un $\mathcal{O}_X$-Module inversible,
  $\mathcal{K}$ un $\mathcal{O}_Y$-Module inversible. Pour que $\mathcal{L}$
  soit très ample relativement à $q$, il faut et il suffit que
  $\mathcal{L}\otimes q^*(\mathcal{K})$ le soit.
  \item[\rm{(ii)}] Si $\mathcal{L}$ et $\mathcal{L}'$ sont deux
  $\mathcal{O}_X$-Modules très amples relativement à $q$, il en est de même
  de $\mathcal{L}\otimes\mathcal{L}'$~; en particulier
  $\mathcal{L}^{\otimes n}$ est très ample relativement à $q$ pour tout
  $n>0$.
\end{enumerate}
\end{corollary}

\begin{proof}
(ii) est conséquence immédiate de (\hyperref[II.4.4.8-fr]{4.4.8}), ainsi que
la nécessité de la condition (i)~; d'autre part, si
$\mathcal{L}\otimes q^*(\mathcal{K})$ est très ample, il en est de même de
$(\mathcal{L}\otimes q^*(\mathcal{K}))\otimes q^*(\mathcal{K}^{-1})$
d'après ce qui précède, et ce dernier $\mathcal{O}_X$-Module est isomorphe à
$\mathcal{L}$ (\hyperref[0.5.4.3-fr]{0, 5.4.3} et
\hyperref[0.5.4.5-fr]{5.4.5}).
\end{proof}

\begin{proposition}[4.4.10]
\label{II.4.4.10-fr}
\begin{enumerate}
  \item[\rm{(i)}] Pour tout préschéma $Y$, tout $\mathcal{O}_Y$-Module
  inversible $\mathcal{L}$ est très ample relativement au morphisme identique
  $1_Y$.
  \item[\rm{(i bis)}] Soient $f:X\to Y$ un morphisme, $j:X'\to X$ une
  immersion. Si $\mathcal{L}$ est un $\mathcal{O}_X$-Module très ample
  relativement à $f$, $j^*(\mathcal{L})$ est très ample relativement à
  $f\circ j$.
  \item[\rm{(ii)}] Soient $Z$ un préschéma quasi-compact, $f:X\to Y$ un
  morphisme de type fini, $g:Y\to Z$ un morphisme quasi-compact,
  $\mathcal{L}$ un $\mathcal{O}_X$-Module très ample relativement à $f$,
  $\mathcal{K}$ un $\mathcal{O}_Y$-Module très ample relativement à $g$.
  Alors il existe un entier $n_0>0$ tel que
  $\mathcal{L}\otimes f^*(\mathcal{K}^{\otimes n})$ soit très ample
  relativement à $g\circ f$ pour tout $n\geq n_0$.
  \item[\rm{(iii)}] Soient $f:X\to Y$, $g:Y'\to Y$ deux morphismes, et posons
  $X'=X_{(Y')}$. Si $\mathcal{L}$ est un $\mathcal{O}_X$-Module très ample
  relativement à $f$, $\mathcal{L}'=\mathcal{L}\otimes_Y\mathcal{O}_{Y'}$
  est un $\mathcal{O}_{X'}$-Module très ample relativement à $f_{(Y')}$.
  \item[\rm{(iv)}] Soient $f_i:X_i\to Y_i$ ($i=1,2$) deux $S$-morphismes.
  Si $\mathcal{L}_i$ est un $\mathcal{O}_{X_i}$-Module très ample relativement
  à $f_i$ ($i=1,2$), $\mathcal{L}_1\otimes_S\mathcal{L}_2$ est très ample
  relativement à $f_1\times_Sf_2$.
  \item[\rm{(v)}] Soient $f:X\to Y$, $g:Y\to Z$ deux morphismes. Si un
  $\mathcal{O}_X$-Module $\mathcal{L}$ est très ample relativement à
  $g\circ f$, alors $\mathcal{L}$ est très ample relativement à $f$.
  \item[\rm{(vi)}] Soient $f:X\to Y$ un morphisme, $j$ l'injection canonique
  $X_{\mathrm{red}}\to X$. Si $\mathcal{L}$ est un $\mathcal{O}_X$-Module
  très ample relativement à $f$, alors $j^*(\mathcal{L})$ est très ample
  relativement à $f_{\mathrm{red}}$.
\end{enumerate}
\end{proposition}

\begin{proof}
La propriété (i bis) découle aussitôt de la définition
(\hyperref[II.4.4.2-fr]{4.4.2}) et il est immédiat que (vi) se déduit
formellement de (i bis) et de (v) par un raisonnement calqué sur celui de
(\hyperref[I.5.5.12-fr]{I, 5.5.12}), que nous laissons au lecteur. Pour
démontrer (v), considérons, comme dans
(\hyperref[I.5.5.12-fr]{I, 5.5.12}), la factorisation
$X\xrightarrow{\Gamma_f}X\times_ZY\xrightarrow{p_2}Y$, en notant que
$p_2=(g\circ f)\times1_Y$. Il résulte de l'hypothèse et de (i) et (iv) que
$\mathcal{L}\otimes_{\mathcal{O}_Z}\mathcal{O}_Y$ est très ample pour $p_2$~;
d'autre part, on a
$\mathcal{L}=\Gamma_f^*(\mathcal{L}\otimes_{\mathcal{O}_Z}\mathcal{O}_Y)$
(\hyperref[I.9.1.4-fr]{I, 9.1.4}), et $\Gamma_f$ est une immersion
(\hyperref[I.5.3.11-fr]{I, 5.3.11})~; on peut donc appliquer (i bis).

\oldpage[II]{82}

Pour prouver (i), on applique la définition
(\hyperref[II.4.4.2-fr]{4.4.2}) avec $\mathcal{E}=\mathcal{L}$, en notant
qu'alors $\mathbf{P}(\mathcal{E})$ s'identifie à $Y$
(\hyperref[II.4.1.4-fr]{4.1.4}).

Démontrons (iii). Il existe un $\mathcal{O}_Y$-Module quasi-cohérent
$\mathcal{E}$ et une $Y$-immersion $i:X\to\mathbf{P}(\mathcal{E})=P$ telle
que $\mathcal{L}=i^*(\mathcal{O}_P(1))$~; alors, si on pose
$\mathcal{E}'=g^*(\mathcal{E})$, $\mathcal{E}'$ est un
$\mathcal{O}_{Y'}$-Module quasi-cohérent et on a
$P'=\mathbf{P}(\mathcal{E}')=P_{(Y')}$
(\hyperref[II.4.1.3.1-fr]{4.1.3.1}), $i_{(Y')}$ est une immersion de
$X_{(Y')}$ dans $P'$ (\hyperref[I.4.3.2-fr]{I, 4.3.2}) et $\mathcal{L}'$ est
isomorphe à $(i_{(Y')})^*(\mathcal{O}_{P'}(1))$
(\hyperref[II.4.2.10-fr]{4.2.10}).

Pour prouver (iv), remarquons qu'il y a par hypothèse une $Y_i$-immersion
$r_i:X_i\to P_i=\mathbf{P}(\mathcal{E}_i)$, où $\mathcal{E}_i$ est un
$\mathcal{O}_{Y_i}$-Module quasi-cohérent, et
$\mathcal{L}_i=r_i^*(\mathcal{O}_{P_i}(1))$ ($i=1,2$)~;
$r_1\times_Sr_2$ est une $S$-immersion de $X_1\times_SX_2$ dans
$P_1\times_SP_2$ (\hyperref[I.4.3.1-fr]{I, 4.3.1}) et l'image réciproque de
$\mathcal{O}_{P_1}(1)\otimes_S\mathcal{O}_{P_2}(1)$ par cette immersion est
$\mathcal{L}_1\otimes_S\mathcal{L}_2$. D'autre part, posons
$T=Y_1\times_SY_2$, et soient $p_1,p_2$ les projections de $T$ dans $Y_1,Y_2$
respectivement. Si on pose
$P_i'=\mathbf{P}(p_i^*(\mathcal{E}_i))$ ($i=1,2$), on a par
(\hyperref[II.4.1.3.1-fr]{4.1.3.1}), $P_i'=P_i\times_{Y_i}T$ d'où
\[
  P_1'\times_TP_2'
  =(P_1\times_{Y_1}T)\times_T(P_2\times_{Y_2}T)
  =P_1\times_{Y_1}(T\times_{Y_2}P_2)
  =P_1\times_{Y_1}(Y_1\times_SP_2)
  =P_1\times_SP_2
\]
à un isomorphisme canonique près. De même, on a
$\mathcal{O}_{P_i'}(1)=\mathcal{O}_{P_i}(1)\otimes_{Y_i}\mathcal{O}_T$
(\hyperref[II.4.1.3.2-fr]{4.1.3.2}), et un calcul analogue (basé notamment
sur (\hyperref[I.9.1.9.1-fr]{I, 9.1.9.1} et
\hyperref[I.9.1.2-fr]{9.1.2})) montre que dans l'identification précédente,
$\mathcal{O}_{P_1'}(1)\otimes_T\mathcal{O}_{P_2'}(1)$ s'identifie à
$\mathcal{O}_{P_1}(1)\otimes_S\mathcal{O}_{P_2}(1)$. On peut donc considérer
$r_1\times_Sr_2$ comme une $T$-immersion de $X_1\times_SX_2$ dans
$P_1'\times_TP_2'$, l'image réciproque de
$\mathcal{O}_{P_1'}(1)\otimes_T\mathcal{O}_{P_2'}(1)$ par cette immersion
étant $\mathcal{L}_1\otimes_S\mathcal{L}_2$. On termine alors le raisonnement
comme dans (\hyperref[II.4.4.8-fr]{4.4.8}) en utilisant le morphisme de Segre.

Il reste à démontrer (ii). On peut d'abord se restreindre au cas où $Z$ est un
schéma affine, car en général il y a un recouvrement fini $(U_i)$ de $Z$ par
des ouverts affines~; si la proposition est démontrée pour
$\mathcal{K}|g^{-1}(U_i)$, $\mathcal{L}|f^{-1}(g^{-1}(U_i))$ et un entier
$n_i$, il suffira de prendre pour $n_0$ le plus grand des $n_i$ pour démontrer
la proposition pour $\mathcal{K}$ et $\mathcal{L}$
(\hyperref[II.4.4.5-fr]{4.4.5}). L'hypothèse implique que $f$ et $g$ sont des
morphismes séparés, donc $X$ et $Y$ sont des schémas quasi-compacts.

Il y a une immersion $r:X\to P=\mathbf{P}(\mathcal{E})$, où $\mathcal{E}$ est
un $\mathcal{O}_Y$-Module quasi-cohérent de type fini et
$\mathcal{L}=r^*(\mathcal{O}_P(1))$, en vertu de
(\hyperref[II.4.4.6-fr]{4.4.6}). Nous allons voir qu'il existe un
$\mathcal{O}_P$-Module $\mathcal{M}$ très ample relativement au morphisme
composé $P\xrightarrow{h}Y\xrightarrow{g}Z$, tel que $\mathcal{O}_P(1)$ soit
isomorphe à $\mathcal{M}\otimes_Y\mathcal{K}^{\otimes(-m)}$ pour un entier
$m$. Pour $n\geq m+1$,
$\mathcal{O}_P(1)\otimes_Y\mathcal{K}^{\otimes n}$ sera donc très ample pour
$Z$ en vertu de l'hypothèse et de (iv) appliqué aux morphismes $h:P\to Y$ et
$1_Y$~; comme $r$ est une immersion et que
$\mathcal{L}\otimes f^*(\mathcal{K}^{\otimes n})
=r^*(\mathcal{O}_P(1)\otimes_Y\mathcal{K}^{\otimes n})$, la conclusion
résultera de (i bis). Pour établir notre assertion sur $\mathcal{O}_P(1)$,
nous utiliserons le lemme suivant~:

\begin{lemma}[4.4.10.1]
\label{II.4.4.10.1-fr}
Soient $Z$ un schéma quasi-compact ou un préschéma dont l'espace sous-jacent
est noethérien, $g:Y\to Z$ un morphisme quasi-compact, $\mathcal{K}$ un
$\mathcal{O}_Y$-Module inversible très ample pour $g$, $\mathcal{E}$ un
$\mathcal{O}_Y$-Module quasi-cohérent de type fini. Il existe alors un entier
$m_0$ tel que, pour tout $m\geq m_0$, $\mathcal{E}$ soit isomorphe à un
quotient d'un $\mathcal{O}_Y$-Module de la forme
$g^*(\mathcal{F})\otimes\mathcal{K}^{\otimes(-m)}$, où $\mathcal{F}$ est un
$\mathcal{O}_Z$-Module quasi-cohérent de type fini (dépendant de $m$).
\end{lemma}

Ce lemme sera démontré en (\hyperref[II.4.5.10.1-fr]{4.5.10.1})~; le lecteur
pourra vérifier que (\hyperref[II.4.4.10-fr]{4.4.10}) n'est utilisé nulle part
dans le n\textsuperscript{o}~4.5.

\oldpage[II]{83}

Ce lemme étant admis, il existe une immersion fermée $j_1$ de $P$ dans
\[
  P_1=\mathbf{P}(g^*(\mathcal{F})\otimes\mathcal{K}^{\otimes(-m)})
\]
telle que $\mathcal{O}_P(1)$ soit isomorphe à
$j_1^*(\mathcal{O}_{P_1}(1))$ (\hyperref[II.4.1.2-fr]{4.1.2}). D'autre part,
il existe un isomorphisme de $P_1$ sur
$P_2=\mathbf{P}(g^*(\mathcal{F}))$, identifiant
$\mathcal{O}_{P_1}(1)$ à
$\mathcal{O}_{P_2}(1)\otimes_Y\mathcal{K}^{\otimes(-m)}$
(\hyperref[II.4.1.4-fr]{4.1.4})~; on a donc une immersion fermée
$j_2:P\to P_2$ telle que $\mathcal{O}_P(1)$ soit isomorphe à
$j_2^*(\mathcal{O}_{P_2}(1))\otimes_Y\mathcal{K}^{\otimes(-m)}$. Enfin,
$P_2$ s'identifie à $P_3\times_ZY$, où $P_3=\mathbf{P}(\mathcal{F})$, et
$\mathcal{O}_{P_2}(1)$ à
$\mathcal{O}_{P_3}(1)\otimes_Z\mathcal{O}_Y$
(\hyperref[II.4.1.3-fr]{4.1.3}). Par définition,
$\mathcal{O}_{P_3}(1)$ est très ample pour $Z$~; comme il en est de même de
$\mathcal{K}$, on conclut de (iv) que
$\mathcal{O}_{P_2}(1)\otimes_Y\mathcal{K}$ est très ample pour $Z$~; il en
est de même de
$\mathcal{M}=j_2^*(\mathcal{O}_{P_2}(1)\otimes_Y\mathcal{K})$ en vertu de
(i bis), et $\mathcal{O}_P(1)$ est isomorphe à
$\mathcal{M}\otimes_Y\mathcal{K}^{\otimes(-m-1)}$, ce qui achève la
démonstration.
\end{proof}

\begin{proposition}[4.4.11]
\label{II.4.4.11-fr}
Soient $f:X\to Y$, $f':X'\to Y$ deux morphismes, $X''$ le préschéma somme
$X\sqcup X'$, $f''$ le morphisme $X''\to Y$ qui coïncide avec $f$ (resp.
$f'$) dans $X$ (resp. $X'$). Soit $\mathcal{L}$ (resp. $\mathcal{L}'$) un
$\mathcal{O}_X$-Module (resp. un $\mathcal{O}_{X'}$-Module) inversible, et
soit $\mathcal{L}''$ le $\mathcal{O}_{X''}$-Module inversible qui coïncide
avec $\mathcal{L}$ dans $X$ et avec $\mathcal{L}'$ dans $X'$. Pour que
$\mathcal{L}''$ soit très ample relativement à $f''$, il faut et il suffit
que $\mathcal{L}$ soit très ample relativement à $f$ et $\mathcal{L}'$ très
ample relativement à $f'$.
\end{proposition}

\begin{proof}
On est aussitôt ramené au cas où $Y$ est affine. Si $\mathcal{L}''$ est très
ample il en est de même de $\mathcal{L}$ et $\mathcal{L}'$ en vertu de
(\hyperref[II.4.4.10-fr]{4.4.10, (i bis)}). Inversement, si $\mathcal{L}$ et
$\mathcal{L}'$ sont très amples, il résulte aussitôt de la définition
(\hyperref[II.4.4.2-fr]{4.4.2}) et de
(\hyperref[II.4.3.6-fr]{4.3.6}) que $\mathcal{L}''$ est très ample.
\end{proof}

\subsection{Faisceaux amples}
\label{subsection:II.4.5-fr}

\begin{env}[4.5.1]
\label{II.4.5.1-fr}
Étant donnés un préschéma $X$ et un $\mathcal{O}_X$-Module inversible
$\mathcal{L}$, nous poserons, pour tout $\mathcal{O}_X$-Module $\mathcal{F}$
(lorsque aucune confusion ne sera possible sur $\mathcal{L}$),
$\mathcal{F}(n)=\mathcal{F}\otimes_{\mathcal{O}_X}\mathcal{L}^{\otimes n}$
($n\in\mathbf{Z}$)~; nous poserons d'autre part
\[
  S=\bigoplus_{n\geq0}\Gamma(X,\mathcal{L}^{\otimes n})
\]
(sous-anneau gradué de l'anneau $\Gamma_\bullet(\mathcal{L})$ défini dans
(\hyperref[0.5.4.6-fr]{0, 5.4.6})). Si on considère $X$ comme un
$\mathbf{Z}$-préschéma et qu'on désigne par $p$ le morphisme structural
$X\to\operatorname{Spec}(\mathbf{Z})$, il y a correspondance biunivoque
entre les homomorphismes
\[
  p^*(\widetilde{S})\longrightarrow
  \bigoplus_{n\geq0}\mathcal{L}^{\otimes n}
\]
de $\mathcal{O}_X$-Algèbres graduées et les endomorphismes de l'anneau gradué
$S$ (\hyperref[I.2.2.5-fr]{I, 2.2.5})~; l'homomorphisme
\[
  \varepsilon:p^*(\widetilde{S})\longrightarrow
  \bigoplus_{n\geq0}\mathcal{L}^{\otimes n}
\]
qui correspond ainsi à l'automorphisme identique de $S$ sera dit canonique.
Il lui correspond (\hyperref[II.3.7.1-fr]{3.7.1}) un morphisme
$G(\varepsilon)\to\operatorname{Proj}(S)$ qui sera aussi dit canonique.
\end{env}

\begin{theorem}[4.5.2]
\label{II.4.5.2-fr}
Soient $X$ un schéma quasi-compact ou un préschéma dont l'espace sous-jacent
est noethérien, $\mathcal{L}$ un $\mathcal{O}_X$-Module inversible, $S$
l'anneau gradué
$\bigoplus_{n\geq0}\Gamma(X,\mathcal{L}^{\otimes n})$. Les conditions
suivantes sont équivalentes~:
\begin{enumerate}
  \item[\rm{a)}] Lorsque $f$ parcourt l'ensemble des éléments homogènes de
  $S_+$, les $X_f$ forment une base de la topologie de $X$.
  \item[\rm{a')}] Lorsque $f$ parcourt l'ensemble des éléments homogènes de
  $S_+$, ceux des $X_f$ qui sont affines forment un recouvrement de $X$.
  \item[\rm{b)}] Le morphisme canonique
  $G(\varepsilon)\to\operatorname{Proj}(S)$
  (\hyperref[II.4.5.1-fr]{4.5.1}) est partout défini et est une immersion
  ouverte dominante.

\oldpage[II]{84}

  \item[\rm{b')}] Le morphisme canonique
  $G(\varepsilon)\to\operatorname{Proj}(S)$ est partout défini et est un
  homéomorphisme de l'espace sous-jacent à $X$ sur un sous-espace de
  $\operatorname{Proj}(S)$.
  \item[\rm{c)}] Pour tout $\mathcal{O}_X$-Module quasi-cohérent
  $\mathcal{F}$, si on désigne par $\mathcal{F}_n$ le
  sous-$\mathcal{O}_X$-Module de $\mathcal{F}(n)$ engendré par les sections de
  $\mathcal{F}(n)$ au-dessus de $X$, alors $\mathcal{F}$ est la somme des
  sous-$\mathcal{O}_X$-Modules $\mathcal{F}_n(-n)$, pour les entiers $n>0$.
  \item[\rm{c')}] La propriété c) a lieu pour tout faisceau quasi-cohérent
  d'idéaux de $\mathcal{O}_X$.
\end{enumerate}
En outre, si $(f_\alpha)$ est une famille d'éléments homogènes de $S_+$,
telle que les $X_{f_\alpha}$ soient affines, alors la restriction à
$\bigcup_\alpha X_{f_\alpha}$ du morphisme canonique
$X\to\operatorname{Proj}(S)$ est un isomorphisme de
$\bigcup_\alpha X_{f_\alpha}$ sur
$\bigcup_\alpha(\operatorname{Proj}(S))_{f_\alpha}$.
\end{theorem}

\begin{proof}
Il est clair que b) implique b'), et b') implique a) en vertu de la formule
(\hyperref[II.3.7.3.1-fr]{3.7.3.1}) (tenant compte de ce que
$\varepsilon^\flat$ est l'identité). La condition a) implique a'), car tout
$x\in X$ a un voisinage affine $U$ tel que
$\mathcal{L}|U$ soit isomorphe à $\mathcal{O}_X|U$~; si
$f\in\Gamma(X,\mathcal{L}^{\otimes n})$ est telle que $x\in X_f\subset U$,
$X_f$ est aussi l'ensemble des $x'\in U$ tels que $(f|U)(x')\neq0$, et c'est
par suite un ouvert affine (\hyperref[I.1.3.6-fr]{I, 1.3.6}). Pour prouver
que a') entraîne b), il suffit de prouver la dernière assertion de l'énoncé,
et de vérifier de plus que si $X=\bigcup_\alpha X_{f_\alpha}$, la condition
(iv) de (\hyperref[II.3.8.2-fr]{3.8.2}) est satisfaite. Ce dernier point
résulte aussitôt de (\hyperref[I.9.3.1-fr]{I, 9.3.1, (i)}). Quant à la
dernière assertion de (\hyperref[II.4.5.2-fr]{4.5.2}), comme
$X_{f_\alpha}$ est l'image réciproque de
$(\operatorname{Proj}(S))_{f_\alpha}$ par
$G(\varepsilon)\to\operatorname{Proj}(S)$, il suffit d'appliquer
(\hyperref[I.9.3.2-fr]{I, 9.3.2}). Donc a), a'), b), b') sont équivalentes.

Pour montrer que a') entraîne c), notons que si $X_f$ est affine (avec
$f\in S_k$), $\mathcal{F}|X_f$ est engendré par ses sections au-dessus de
$X_f$ (\hyperref[I.1.3.9-fr]{I, 1.3.9})~; d'autre part
(\hyperref[I.9.3.1-fr]{I, 9.3.1, (ii)}) une telle section $s$ est de la forme
$(t|X_f)\otimes(f|X_f)^{-m}$ où
$t\in\Gamma(X,\mathcal{F}(km))$~; par définition, $t$ est aussi une section
de $\mathcal{F}_{km}$, donc $s$ est bien une section de
$\mathcal{F}_{km}(-km)$ au-dessus de $X_f$, ce qui prouve c). Il est clair
que c) implique c'), et il reste à montrer que c') entraîne a). Or, soit $U$
un voisinage ouvert de $x\in X$, et soit $\mathcal{J}$ un faisceau d'idéaux
quasi-cohérent de $\mathcal{O}_X$ définissant un sous-préschéma fermé de $X$
ayant $X-U$ comme espace sous-jacent
(\hyperref[I.5.2.1-fr]{I, 5.2.1}). L'hypothèse c') entraîne qu'il existe un
entier $n>0$ et une section $f$ de $\mathcal{J}(n)$ au-dessus de $X$ telle
que $f(x)\neq0$. Mais on a évidemment $f\in S_n$, et $x\in X_f\subset U$, ce
qui prouve a).
\end{proof}

Lorsque $X$ est un préschéma dont l'espace sous-jacent est noethérien, les
conditions équivalentes de (\hyperref[II.4.5.2-fr]{4.5.2}) impliquent que
$X$ est un schéma, puisqu'il est isomorphe à un sous-préschéma du schéma
$S=\operatorname{Proj}(A)$ en vertu de
(\hyperref[II.4.5.2-fr]{4.5.2, b)}).

\begin{definition}[4.5.3]
\label{II.4.5.3-fr}
On dit qu'un $\mathcal{O}_X$-Module inversible $\mathcal{L}$ est ample si
$X$ est un schéma quasi-compact et si les conditions équivalentes de
(\hyperref[II.4.5.2-fr]{4.5.2}) sont vérifiées.
\end{definition}

Il résulte évidemment du critère (\hyperref[II.4.5.2-fr]{4.5.2, a)}) que si
$\mathcal{L}$ est un $\mathcal{O}_X$-Module ample, alors, pour tout ouvert
$U$ de $X$, $\mathcal{L}|U$ est un $(\mathcal{O}_X|U)$-Module ample.

Il résulte de la démonstration de (\hyperref[II.4.5.2-fr]{4.5.2}) que les
$X_f$ affines forment déjà une base de la topologie de $X$. De plus~:

\begin{corollary}[4.5.4]
\label{II.4.5.4-fr}
Soit $\mathcal{L}$ un $\mathcal{O}_X$-Module ample. Pour tout sous-espace fini
$Z$ de $X$ et tout voisinage $U$ de $Z$, il existe un entier $n$ et un
$f\in\Gamma(X,\mathcal{L}^{\otimes n})$ tels que $X_f$ soit un voisinage
affine de $Z$ contenu dans $U$.
\end{corollary}

\begin{proof}
\oldpage[II]{85}

En vertu de (\hyperref[II.4.5.2-fr]{4.5.2, b)}), on peut se borner à prouver
que pour toute partie finie $Z'$ de $\operatorname{Proj}(S)$ et tout
voisinage ouvert $U$ de $Z'$, il existe un élément homogène $f\in S_+$ tel
que $Z'\subset(\operatorname{Proj}(S))_f\subset U$
(\hyperref[II.2.4.1-fr]{2.4.1}). Or, par définition, l'ensemble fermé $Y$,
complémentaire de $U$ dans $\operatorname{Proj}(S)$, est de la forme
$V_+(\mathfrak{I})$, où $\mathfrak{I}$ est un idéal gradué de $S$, ne
contenant pas $S_+$ (\hyperref[II.2.3.2-fr]{2.3.2})~; d'autre part, les
points de $Z'$ sont par définition des idéaux premiers gradués
$\mathfrak{p}_i$ de $S_+$ ne contenant pas $\mathfrak{I}$
(\hyperref[II.2.3.1-fr]{2.3.1}). Il existe donc un élément
$f\in\mathfrak{I}$ n'appartenant à aucun des $\mathfrak{p}_i$ (Bourbaki,
\emph{Alg. comm.}, chap.~II, \S~1, n\textsuperscript{o}~1, prop.~2), et comme
les $\mathfrak{p}_i$ sont gradués, le raisonnement fait \emph{loc. cit.}
montre qu'on peut même supposer $f$ homogène~; cet élément répond alors à la
question.
\end{proof}

\begin{proposition}[4.5.5]
\label{II.4.5.5-fr}
Supposons que $X$ soit un schéma quasi-compact, ou un préschéma dont l'espace
sous-jacent est noethérien. Les conditions a) à c') de
(\hyperref[II.4.5.2-fr]{4.5.2}) sont aussi équivalentes aux suivantes~:
\begin{enumerate}
  \item[\rm{d)}] Pour tout $\mathcal{O}_X$-Module quasi-cohérent
  $\mathcal{F}$ de type fini, il existe un entier $n_0$ tel que, pour tout
  $n\geq n_0$, $\mathcal{F}(n)$ soit engendré par ses sections au-dessus de
  $X$.
  \item[\rm{d')}] Pour tout $\mathcal{O}_X$-Module quasi-cohérent
  $\mathcal{F}$ de type fini, il existe deux entiers $n>0$, $k>0$ tels que
  $\mathcal{F}$ soit isomorphe à un quotient du $\mathcal{O}_X$-Module
  $\mathcal{L}^{\otimes(-n)}\otimes\mathcal{O}_X^k$.
  \item[\rm{d'')}] La propriété d') a lieu pour tout faisceau d'idéaux
  quasi-cohérent de type fini de $\mathcal{O}_X$.
\end{enumerate}
\end{proposition}

\begin{proof}
Comme $X$ est quasi-compact, si un $\mathcal{O}_X$-Module quasi-cohérent de
type fini $\mathcal{F}$ est tel que $\mathcal{F}(n)$ (qui est de type fini)
est engendré par ses sections au-dessus de $X$, $\mathcal{F}(n)$ est alors
engendré par un nombre fini de ces sections
(\hyperref[0.5.2.3-fr]{0, 5.2.3}), donc d) implique d') et il est clair que
d') implique d''). Comme tout $\mathcal{O}_X$-Module quasi-cohérent
$\mathcal{G}$ est limite inductive de ses sous-$\mathcal{O}_X$-Modules de
type fini (\hyperref[I.9.4.9-fr]{I, 9.4.9}), pour vérifier la condition c')
de (\hyperref[II.4.5.2-fr]{4.5.2}), il suffit de le faire pour un faisceau
quasi-cohérent d'idéaux de $\mathcal{O}_X$ qui est de type fini, et d'')
entraîne donc c'). Reste à voir que si $\mathcal{L}$ est ample la propriété
d) est vérifiée. Considérons un recouvrement fini de $X$ par des
$X_{f_i}$ ($f_i\in S_{n_i}$) qu'on peut supposer affines~; en remplaçant les
$f_i$ par des puissances convenables (ce qui ne change pas les $X_{f_i}$),
on peut supposer tous les $n_i$ égaux à un même entier $m$. Le faisceau
$\mathcal{F}|X_{f_i}$, étant de type fini par hypothèse, est engendré par un
nombre fini de ses sections $h_{ij}$ au-dessus de $X_{f_i}$
(\hyperref[I.1.3.13-fr]{I, 1.3.13})~; il existe donc un entier $k_0$ tel que
la section $h_{ij}\otimes f_i^{\otimes k_0}$ se prolonge en une section de
$\mathcal{F}(k_0m)$ au-dessus de $X$ pour tout couple $(i,j)$
(\hyperref[I.9.3.1-fr]{I, 9.3.1}). \emph{A fortiori} les
$h_{ij}\otimes f_i^{\otimes k}$ se prolongent en sections de
$\mathcal{F}(km)$ au-dessus de $X$ pour tout $k\geq k_0$, et pour ces valeurs
de $k$, $\mathcal{F}(km)$ est donc engendré par ses sections au-dessus de
$X$. Pour tout $p$ tel que $0<p<m$, $\mathcal{F}(p)$ est aussi de type fini,
donc il y a un entier $k_p$ tel que
$\mathcal{F}(p)(km)=\mathcal{F}(p+km)$ soit engendré par ses sections
au-dessus de $X$ pour $k\geq k_p$. Prenant $n_0$ plus grand que tous les
$k_pm$, on en conclut que $\mathcal{F}(n)$ est engendré par ses sections
au-dessus de $X$ pour tout $n\geq n_0$, car un tel $n$ s'écrit $n=km+p$ avec
$k\geq k_p$ et $0\leq p<m$.
\end{proof}

\begin{proposition}[4.5.6]
\label{II.4.5.6-fr}
Soient $X$ un schéma quasi-compact, $\mathcal{L}$ un $\mathcal{O}_X$-Module
inversible.
\begin{enumerate}
  \item[\rm{(i)}] Soit $n$ un entier $>0$. Pour que $\mathcal{L}$ soit
  ample, il faut et il suffit que $\mathcal{L}^{\otimes n}$ soit ample.
  \item[\rm{(ii)}] Soit $\mathcal{L}'$ un $\mathcal{O}_X$-Module inversible
  tel que, pour tout $x\in X$, il existe un entier $n>0$

\oldpage[II]{86}

  et une section $s'$ de $\mathcal{L}'^{\otimes n}$ au-dessus de $X$ tels
  que $s'(x)\neq0$. Alors, si $\mathcal{L}$ est ample, il en est de même de
  $\mathcal{L}\otimes\mathcal{L}'$.
\end{enumerate}
\end{proposition}

\begin{proof}
La propriété (i) est conséquence évidente du critère a) de
(\hyperref[II.4.5.2-fr]{4.5.2}) puisque $X_{f^{\otimes n}}=X_f$. D'autre
part, si $\mathcal{L}$ est ample, pour tout $x\in X$ et tout voisinage $U$ de
$x$, il existe $m>0$ et $f\in\Gamma(X,\mathcal{L}^{\otimes m})$ tel que
$x\in X_f\subset U$ (\hyperref[II.4.5.2-fr]{4.5.2, a)})~; si
$f'\in\Gamma(X,\mathcal{L}'^{\otimes n})$ est tel que $f'(x)\neq0$, on aura
$s(x)\neq0$ pour
$s=f^{\otimes n}\otimes f'^{\otimes m}\in
\Gamma(X,(\mathcal{L}\otimes\mathcal{L}')^{\otimes mn})$, donc
$x\in X_s\subset X_f\subset U$, ce qui prouve que
$\mathcal{L}\otimes\mathcal{L}'$ est ample
(\hyperref[II.4.5.2-fr]{4.5.2, a)}).
\end{proof}

\begin{corollary}[4.5.7]
\label{II.4.5.7-fr}
Le produit tensoriel de deux $\mathcal{O}_X$-Modules amples est ample.
\end{corollary}

\begin{corollary}[4.5.8]
\label{II.4.5.8-fr}
Soient $\mathcal{L}$ un $\mathcal{O}_X$-Module ample, $\mathcal{L}'$ un
$\mathcal{O}_X$-Module inversible~; il existe alors un entier $n_0>0$ tel
que $\mathcal{L}^{\otimes n}\otimes\mathcal{L}'$ soit ample et engendré par
ses sections au-dessus de $X$ pour $n\geq n_0$.
\end{corollary}

\begin{proof}
En effet, il résulte de (\hyperref[II.4.5.5-fr]{4.5.5}) qu'il existe un entier
$m_0$ tel que $\mathcal{L}^{\otimes m}\otimes\mathcal{L}'$ soit engendré par
ses sections au-dessus de $X$ pour $m\geq m_0$~; d'après
(\hyperref[II.4.5.6-fr]{4.5.6}) on peut donc prendre $n_0=m_0+1$.
\end{proof}

\begin{remark}[4.5.9]
\label{II.4.5.9-fr}
Soit $P=\mathrm{H}^1(X,\mathcal{O}_X^*)$ le groupe des classes de
$\mathcal{O}_X$-Modules inversibles (\hyperref[0.5.4.7-fr]{0, 5.4.7}), et
soit $P^+$ la partie de $P$ formée des classes de faisceaux amples. Supposons
que $P^+$ soit non vide. Alors il résulte de
(\hyperref[II.4.5.7-fr]{4.5.7}) et
(\hyperref[II.4.5.8-fr]{4.5.8}) que l'on a
\[
  P^++P^+\subset P^+
  \qquad\text{et}\qquad
  P^+-P^+=P
\]
autrement dit $P^+\cup\{0\}$ est l'ensemble des éléments positifs dans $P$
pour une structure de préordre sur $P$ compatible avec sa structure de
groupe, et qui est même archimédienne en vertu de
(\hyperref[II.4.5.8-fr]{4.5.8}). C'est pourquoi on dit parfois «~faisceau
positif~» au lieu de faisceau ample et «~faisceau négatif~» pour l'inverse
d'un tel faisceau (terminologie que nous ne suivrons pas).
\end{remark}

\begin{proposition}[4.5.10]
\label{II.4.5.10-fr}
Soient $Y$ un schéma affine, $q:X\to Y$ un morphisme séparé quasi-compact,
$\mathcal{L}$ un $\mathcal{O}_X$-Module inversible.
\begin{enumerate}
  \item[\rm{(i)}] Si $\mathcal{L}$ est très ample relativement à $q$, alors
  $\mathcal{L}$ est ample.
  \item[\rm{(ii)}] Supposons en outre que le morphisme $q$ soit de type fini.
  Alors, pour que $\mathcal{L}$ soit ample, il faut et il suffit qu'il possède
  l'une des propriétés équivalentes suivantes~:
  \begin{enumerate}
    \item[\rm{e)}] Il existe $n_0>0$ tel que pour tout entier $n\geq n_0$,
    $\mathcal{L}^{\otimes n}$ soit très ample relativement à $q$.
    \item[\rm{e')}] Il existe $n>0$ tel que $\mathcal{L}^{\otimes n}$ soit
    très ample relativement à $q$.
  \end{enumerate}
\end{enumerate}
\end{proposition}

\begin{proof}
La première assertion résulte de la définition
(\hyperref[II.4.4.2-fr]{4.4.2}) d'un $\mathcal{O}_X$-Module très ample~: si
$A$ est l'anneau de $Y$, il y a un $A$-module $E$ et un homomorphisme
surjectif
\[
  \psi:q^*((\mathbf{S}(E))^{\sim})\longrightarrow
  \bigoplus_{n\geq0}\mathcal{L}^{\otimes n}
\]
tel que $i=r_{\mathcal{L},\psi}$ soit une immersion partout définie
$X\to P=\mathbf{P}(\widetilde{E})$ et que l'on ait
$\mathcal{L}=i^*(\mathcal{O}_P(1))$~; comme les $D_+(f)$ pour $f$ homogène
dans $(\mathbf{S}(E))_+$ forment une base de la topologie de $P$, et que
$i^{-1}(D_+(f))=X_{\psi^\flat(f)}$ par
(\hyperref[II.3.7.3.1-fr]{3.7.3.1}), on voit que la condition a) de
(\hyperref[II.4.5.2-fr]{4.5.2}) est vérifiée, donc $\mathcal{L}$ est ample.

Prouvons maintenant que si $q$ est de type fini et si $\mathcal{L}$ est
ample, il vérifie la condition e). En premier lieu, il résulte du critère b)
de (\hyperref[II.4.5.2-fr]{4.5.2}) et de
(\hyperref[II.4.4.1-fr]{4.4.1, (i)}) qu'il existe

\oldpage[II]{87}

un entier $k_0>0$ tel que $\mathcal{L}^{\otimes k_0}$ soit très ample
relativement à $q$. D'autre part, en vertu de
(\hyperref[II.4.5.5-fr]{4.5.5}), il existe un entier $m_0$ tel que, pour
$m\geq m_0$, $\mathcal{L}^{\otimes m}$ soit engendré par ses sections
au-dessus de $X$. Posons $n_0=k_0+m_0$~; si $n\geq n_0$, on peut écrire
$n=k_0+m$ avec $m\geq m_0$, d'où
$\mathcal{L}^{\otimes n}
=\mathcal{L}^{\otimes k_0}\otimes\mathcal{L}^{\otimes m}$. Comme
$\mathcal{L}^{\otimes m}$ est engendré par ses sections au-dessus de $X$, il
résulte du critère (\hyperref[II.4.4.8-fr]{4.4.8}) et de
(\hyperref[II.3.4.7-fr]{3.4.7}) que $\mathcal{L}^{\otimes n}$ est très ample
relativement à $q$. Enfin, il est clair que e) entraîne e'), et e') implique
que $\mathcal{L}$ est ample en vertu de (i) et de
(\hyperref[II.4.5.6-fr]{4.5.6, (i)}).

\begin{env}[4.5.10.1]
\label{II.4.5.10.1-fr}
\emph{Démonstration du lemme (\hyperref[II.4.4.10.1-fr]{4.4.10.1}).} Posons
$\mathcal{E}(n)=\mathcal{E}\otimes\mathcal{K}^{\otimes n}$~; comme $g$ est
séparé (\hyperref[II.4.4.2-fr]{4.4.2}), le raisonnement de
(\hyperref[II.3.4.8-fr]{3.4.8}) s'applique et ramène à voir que
l'homomorphisme canonique
$g^*(g_*(\mathcal{E}(n)))\to\mathcal{E}(n)$ est surjectif pour $n$ assez
grand. En outre, comme $Z$ est quasi-compact, le raisonnement de
(\hyperref[II.3.4.6-fr]{3.4.6}) ramène à prouver la proposition dans le cas
où $Z$ est affine. Or, $\mathcal{K}$ est alors ample en vertu de
(\hyperref[II.4.5.10-fr]{4.5.10, (i)}), et la conclusion résulte de
(\hyperref[II.4.5.5-fr]{4.5.5, d')}).
\end{env}
\end{proof}

\begin{corollary}[4.5.11]
\label{II.4.5.11-fr}
Soient $Y$ un schéma affine, $q:X\to Y$ un morphisme séparé de type fini,
$\mathcal{L}$ un $\mathcal{O}_X$-Module ample, $\mathcal{L}'$ un
$\mathcal{O}_X$-Module inversible. Il existe un entier $n_0$ tel que, pour
$n\geq n_0$, $\mathcal{L}^{\otimes n}\otimes\mathcal{L}'$ soit très ample
relativement à $q$.
\end{corollary}

\begin{proof}
En effet, il existe $m_0$ tel que pour $m\geq m_0$,
$\mathcal{L}^{\otimes m}\otimes\mathcal{L}'$ soit engendré par ses sections
au-dessus de $X$ (\hyperref[II.4.5.8-fr]{4.5.8})~; d'autre part, il existe
$k_0$ tel que $\mathcal{L}^{\otimes k}$ soit très ample relativement à $q$
pour $k\geq k_0$. Par suite
$\mathcal{L}^{\otimes(k+m_0)}\otimes\mathcal{L}'$ est très ample pour
$k\geq k_0$ ((\hyperref[II.4.4.8-fr]{4.4.8}) et
(\hyperref[II.3.4.7-fr]{3.4.7})).
\end{proof}

\begin{remark}[4.5.12]
\label{II.4.5.12-fr}
On ignore si l'hypothèse qu'un $\mathcal{O}_X$-Module $\mathcal{L}$ est tel
que $\mathcal{L}^{\otimes n}$ soit très ample (relativement à $q$) entraîne
la même conclusion pour $\mathcal{L}^{\otimes(n+1)}$.
\end{remark}

\begin{proposition}[4.5.13]
\label{II.4.5.13-fr}
Soient $X$ un préschéma quasi-compact, $Z$ un sous-préschéma fermé de $X$
défini par un faisceau quasi-cohérent nilpotent $\mathcal{J}$ d'idéaux de
$\mathcal{O}_X$, $j$ l'injection canonique $Z\to X$. Pour qu'un
$\mathcal{O}_X$-Module inversible $\mathcal{L}$ soit ample, il faut et il
suffit que $\mathcal{L}'=j^*(\mathcal{L})$ soit un
$\mathcal{O}_Z$-Module ample.
\end{proposition}

\begin{proof}
La condition est nécessaire. En effet, pour toute section $f$ de
$\mathcal{L}^{\otimes n}$ au-dessus de $X$, soit $f'$ son image canonique
$f\otimes1$, section de
$\mathcal{L}'^{\otimes n}=\mathcal{L}^{\otimes n}
\otimes_{\mathcal{O}_X}(\mathcal{O}_X/\mathcal{J})$ au-dessus de l'espace
$Z$ (qui est identique à $X$)~; il est clair que $X_f=Z_{f'}$, donc le
critère a) de (\hyperref[II.4.5.2-fr]{4.5.2}) montre que $\mathcal{L}'$ est
ample.

Pour voir que la condition est suffisante, notons d'abord qu'on peut se
ramener au cas où $\mathcal{J}^2=0$, en considérant la suite (finie) des
préschémas
$X_k=(X,\mathcal{O}_X/\mathcal{J}^{k+1})$ dont chacun est sous-préschéma
fermé du suivant, défini par un faisceau d'idéaux de carré nul. D'ailleurs
$X$ est un schéma puisque $X_{\mathrm{red}}$ l'est par hypothèse
((\hyperref[II.4.5.3-fr]{4.5.3}) et
(\hyperref[I.5.5.1-fr]{I, 5.5.1})). Le critère a) de
(\hyperref[II.4.5.2-fr]{4.5.2}) montre qu'il suffira de prouver le

\begin{lemma}[4.5.13.1]
\label{II.4.5.13.1-fr}
Sous les hypothèses de (\hyperref[II.4.5.13-fr]{4.5.13}), supposons de plus
$\mathcal{J}$ de carré nul~; $\mathcal{L}$ étant un
$\mathcal{O}_X$-Module inversible, soit $g$ une section de
$\mathcal{L}'^{\otimes n}$ au-dessus de $Z$ telle que $Z_g$ soit affine.
Alors il existe un entier $m>0$ tel que $g^{\otimes m}$ soit l'image
canonique d'une section $f$ de $\mathcal{L}^{\otimes nm}$ au-dessus de $X$.
\end{lemma}

On a la suite exacte de $\mathcal{O}_X$-Modules
\[
  0\longrightarrow\mathcal{J}(n)\longrightarrow
  \mathcal{O}_X(n)=\mathcal{L}^{\otimes n}\longrightarrow
  \mathcal{O}_Z(n)=\mathcal{L}'^{\otimes n}\longrightarrow0
\]

\oldpage[II]{88}

puisque $\mathcal{F}(n)$ est un foncteur exact en $\mathcal{F}$~; d'où la
suite exacte de cohomologie
\[
  0\longrightarrow\Gamma(X,\mathcal{J}(n))\longrightarrow
  \Gamma(X,\mathcal{L}^{\otimes n})\longrightarrow
  \Gamma(X,\mathcal{L}'^{\otimes n})\xrightarrow{\partial}
  \mathrm{H}^1(X,\mathcal{J}(n))
\]
qui associe en particulier à $g$ un élément
$\partial g\in\mathrm{H}^1(X,\mathcal{J}(n))$.

Notons que puisque $\mathcal{J}^2=0$, $\mathcal{J}$ peut être considéré
comme un $\mathcal{O}_Z$-Module quasi-cohérent et l'on a, pour tout $k$,
$\mathcal{L}'^{\otimes k}\otimes_{\mathcal{O}_Z}\mathcal{J}(n)
=\mathcal{J}(n+k)$~; pour toute section
$s\in\Gamma(X,\mathcal{L}'^{\otimes k})$ la multiplication tensorielle par
$s$ est donc un homomorphisme $\mathcal{J}(n)\to\mathcal{J}(n+k)$ de
$\mathcal{O}_Z$-Modules, qui donne par suite un homomorphisme
$\mathrm{H}^i(X,\mathcal{J}(n))\to
\mathrm{H}^i(X,\mathcal{J}(n+k))$ de groupes de cohomologie.

Cela étant, nous allons voir que l'on a
\begin{equation}
\label{II.4.5.13.2-fr}
  g^{\otimes m}\otimes\partial g=0
\tag{4.5.13.2}
\end{equation}
pour un $m>0$ assez grand. En effet, $Z_g$ est un ouvert affine de $Z$ et
on a donc $\mathrm{H}^1(Z_g,\mathcal{J}(n))=0$ quand $\mathcal{J}(n)$ est
considéré comme un $\mathcal{O}_Z$-Module
(\hyperref[I.5.1.9.2-fr]{I, 5.1.9.2}). En particulier, si on pose
$g'=g|Z_g$, et si on considère son image par l'application
$\partial:\Gamma(Z_g,\mathcal{L}'^{\otimes n})\to
\mathrm{H}^1(Z_g,\mathcal{J}(n))$, on a $\partial g'=0$. Pour expliciter
cette relation, observons qu'en dimension $1$ la cohomologie d'un faisceau
de groupes abéliens est la même que sa cohomologie de Čech (G, II, 5.9)~;
pour former $\partial g$, il faut donc considérer un recouvrement ouvert
$(U_\alpha)$ assez fin de $X$, que l'on peut supposer fini et formé
d'ouverts affines, prendre pour chaque $\alpha$ une section
$g_\alpha\in\Gamma(U_\alpha,\mathcal{L}^{\otimes n})$ dont l'image
canonique dans $\Gamma(U_\alpha,\mathcal{L}'^{\otimes n})$ est
$g|U_\alpha$, et considérer la classe du cocycle
$(g_{\alpha|\beta}-g_{\beta|\alpha})$, $g_{\alpha|\beta}$ étant la
restriction de $g_\alpha$ à $U_\alpha\cap U_\beta$ (ce cocycle étant à
valeurs dans $\mathcal{J}(n)$). On peut en outre supposer que $\partial g'$
est calculé de la même manière à l'aide du recouvrement formé des
$U_\alpha\cap Z_g$ et des restrictions $g_\alpha|(U_\alpha\cap Z_g)$ (en
remplaçant au besoin $(U_\alpha)$ par un recouvrement plus fin)~; la
relation $\partial g'=0$ signifie alors qu'il existe pour chaque $\alpha$
une section $h_\alpha\in\Gamma(U_\alpha\cap Z_g,\mathcal{J}(n))$ telle que
\[
  (g_{\alpha|\beta}-g_{\beta|\alpha})
  |(U_\alpha\cap U_\beta\cap Z_g)
  =h_{\alpha|\beta}-h_{\beta|\alpha}
\]
en désignant par $h_{\alpha|\beta}$ la restriction de $h_\alpha$ à
$U_\alpha\cap U_\beta\cap Z_g$ (G, II, 5.11). Cela étant, il y a un entier
$m>0$ tel que $g^{\otimes m}\otimes h_\alpha$ soit la restriction à
$U_\alpha\cap Z_g$ d'une section
$t_\alpha\in\Gamma(U_\alpha,\mathcal{J}(n+nm))$ pour tout $\alpha$
(\hyperref[I.9.3.1-fr]{I, 9.3.1})~; on a donc
\[
  g^{\otimes m}\otimes(g_{\alpha|\beta}-g_{\beta|\alpha})
  =t_{\alpha|\beta}-t_{\beta|\alpha}
\]
pour tout couple d'indices, ce qui prouve
(\hyperref[II.4.5.13.2-fr]{4.5.13.2}).

Remarquons d'autre part que si
$s\in\Gamma(X,\mathcal{O}_Z(p))$,
$t\in\Gamma(X,\mathcal{O}_Z(q))$, on a, dans le groupe
$\mathrm{H}^1(X,\mathcal{J}(p+q))$
\begin{equation}
\label{II.4.5.13.3-fr}
  \partial(s\otimes t)=(\partial s)\otimes t+s\otimes(\partial t).
\tag{4.5.13.3}
\end{equation}

En effet, on peut encore, pour calculer les deux membres, considérer un
recouvrement ouvert $(U_\alpha)$ de $X$, pour chaque $\alpha$ une section
$s_\alpha\in\Gamma(U_\alpha,\mathcal{O}_X(p))$ (resp.
$t_\alpha\in\Gamma(U_\alpha,\mathcal{O}_X(q))$) dont l'image canonique dans
$\Gamma(U_\alpha,\mathcal{O}_Z(p))$ (resp.
$\Gamma(U_\alpha,\mathcal{O}_Z(q))$ soit $s|U_\alpha$ (resp.
$t|U_\alpha$)~; la relation
(\hyperref[II.4.5.13.3-fr]{4.5.13.3}) résulte alors des relations
\[
  (s_{\alpha|\beta}\otimes t_{\alpha|\beta})
  -(s_{\beta|\alpha}\otimes t_{\beta|\alpha})
  =(s_{\alpha|\beta}-s_{\beta|\alpha})\otimes t_{\alpha|\beta}
  +s_{\beta|\alpha}\otimes(t_{\alpha|\beta}-t_{\beta|\alpha})
\]
avec les mêmes notations que ci-dessus. Par récurrence sur $k$, on a donc
\begin{equation}
\label{II.4.5.13.4-fr}
  \partial(g^{\otimes k})=(kg^{\otimes(k-1)})\otimes(\partial g)
\tag{4.5.13.4}
\end{equation}

\oldpage[II]{89}

et on conclut de (\hyperref[II.4.5.13.2-fr]{4.5.13.2}) et
(\hyperref[II.4.5.13.4-fr]{4.5.13.4}) que l'on a
$\partial(g^{\otimes(m+1)})=0$~; $g^{\otimes(m+1)}$ est donc l'image
canonique d'une section $f$ de $\mathcal{L}^{\otimes n(m+1)}$ au-dessus de
$X$, ce qui achève de démontrer (\hyperref[II.4.5.13-fr]{4.5.13}).
\end{proof}

\begin{corollary}[4.5.14]
\label{II.4.5.14-fr}
Soient $X$ un schéma noethérien, $j$ l'injection canonique
$X_{\mathrm{red}}\to X$. Pour qu'un $\mathcal{O}_X$-Module inversible
$\mathcal{L}$ soit ample, il faut et il suffit que $j^*(\mathcal{L})$ soit
un $\mathcal{O}_{X_{\mathrm{red}}}$-Module ample.
\end{corollary}

Cela résulte de (\hyperref[I.6.1.6-fr]{I, 6.1.6}).

\subsection{Faisceaux relativement amples}
\label{II.4.6-fr}

\begin{definition}[4.6.1]
\label{II.4.6.1-fr}
Soient $f:X\to Y$ un morphisme quasi-compact, $\mathcal{L}$ un
$\mathcal{O}_X$-Module inversible. On dit que $\mathcal{L}$ est ample
relativement à $f$, ou relativement à $Y$, ou $f$-ample, ou $Y$-ample (ou
même ample si aucune confusion n'en résulte avec la notion définie dans
(\hyperref[II.4.5.3-fr]{4.5.3})) s'il existe un recouvrement $(U_\alpha)$
de $Y$ formé d'ouverts affines tel que si on pose
$X_\alpha=f^{-1}(U_\alpha)$, $\mathcal{L}|X_\alpha$ soit un
$\mathcal{O}_{X_\alpha}$-Module ample pour tout $\alpha$.
\end{definition}

L'existence d'un $\mathcal{O}_X$-Module $f$-ample entraîne que $f$ est
nécessairement séparé ((\hyperref[II.4.5.3-fr]{4.5.3}) et
(\hyperref[I.5.5.5-fr]{I, 5.5.5})).

\begin{proposition}[4.6.2]
\label{II.4.6.2-fr}
Soit $f:X\to Y$ un morphisme quasi-compact, $\mathcal{L}$ un
$\mathcal{O}_X$-Module inversible. Si $\mathcal{L}$ est très ample
relativement à $f$, il est ample relativement à $f$.
\end{proposition}

Cela résulte du caractère local (sur $Y$) de la notion de faisceau très
ample (\hyperref[II.4.4.5-fr]{4.4.5}), de la définition
(\hyperref[II.4.6.1-fr]{4.6.1}) et du critère
(\hyperref[II.4.5.10-fr]{4.5.10, (i)}).

\begin{proposition}[4.6.3]
\label{II.4.6.3-fr}
Soient $f:X\to Y$ un morphisme quasi-compact, $\mathcal{L}$ un
$\mathcal{O}_X$-Module inversible, et soit $\mathcal{S}$ la
$\mathcal{O}_Y$-Algèbre graduée
$\bigoplus_{n\geq0}f_*(\mathcal{L}^{\otimes n})$. Les conditions suivantes
sont équivalentes~:
\begin{enumerate}
  \item[\rm{a)}] $\mathcal{L}$ est $f$-ample.
  \item[\rm{b)}] $\mathcal{S}$ est quasi-cohérente et l'homomorphisme
  canonique $\sigma:f^*(\mathcal{S})\to
  \bigoplus_{n\geq0}\mathcal{L}^{\otimes n}$
  (\hyperref[0.4.4.3-fr]{0, 4.4.3}) est tel que le $Y$-morphisme
  $r_{\mathcal{L},\sigma}:G(\sigma)\to\operatorname{Proj}(\mathcal{S})=P$
  soit partout défini et soit une immersion ouverte dominante.
  \item[\rm{b')}] Le morphisme $f$ est séparé, le $Y$-morphisme
  $r_{\mathcal{L},\sigma}$ est partout défini et est un homéomorphisme de
  l'espace sous-jacent à $X$ sur un sous-espace de
  $\operatorname{Proj}(\mathcal{S})$.
\end{enumerate}
En outre, lorsqu'il en est ainsi, pour tout $n\in\mathbf{Z}$,
l'homomorphisme canonique
\begin{equation}
\label{II.4.6.3.1-fr}
  r_{\mathcal{L},\sigma}^*(\mathcal{O}_P(n))
  \longrightarrow\mathcal{L}^{\otimes n}
\tag{4.6.3.1}
\end{equation}
défini dans (\hyperref[II.3.7.9.1-fr]{3.7.9.1}), est un isomorphisme.

Enfin, pour tout $\mathcal{O}_X$-Module quasi-cohérent $\mathcal{F}$, si
l'on pose
$\mathcal{M}=\bigoplus_{n\geq0}f_*(\mathcal{F}\otimes
\mathcal{L}^{\otimes n})$, l'homomorphisme canonique
\begin{equation}
\label{II.4.6.3.2-fr}
  r_{\mathcal{L},\sigma}^*(\widetilde{\mathcal{M}})
  \longrightarrow\mathcal{F}
\tag{4.6.3.2}
\end{equation}
défini dans (\hyperref[II.3.7.9.2-fr]{3.7.9.2}), est un isomorphisme.
\end{proposition}

On a remarqué que a) implique que $f$ est séparé, donc $\mathcal{S}$ est
quasi-cohérente (\hyperref[I.9.2.2-fr]{I, 9.2.2, a)}). Comme le fait que
$r_{\mathcal{L},\sigma}$ soit une immersion partout définie est de
caractère local sur $Y$, pour prouver que a) entraîne b), on peut supposer
$Y$ affine et $\mathcal{L}$ ample~;

\oldpage[II]{90}

l'assertion résulte alors de (\hyperref[II.4.5.2-fr]{4.5.2, b)}). Il est
clair que b) entraîne b')~; enfin, pour prouver que b') entraîne a), il
suffit de considérer un recouvrement de $Y$ par des ouverts affines
$U_\alpha$ et d'appliquer le critère
(\hyperref[II.4.5.2-fr]{4.5.2, b')}) à chaque faisceau
$\mathcal{L}|f^{-1}(U_\alpha)$.

Pour les deux dernières assertions, on utilise le fait que
$\sigma^\flat$ est ici l'identité, et l'explicitation des homomorphismes
(\hyperref[II.3.7.9.1-fr]{3.7.9.1}) et
(\hyperref[II.3.7.9.2-fr]{3.7.9.2})~; il en résulte aussitôt que
(\hyperref[II.4.6.3.1-fr]{4.6.3.1}) est un isomorphisme. Quant à
(\hyperref[II.4.6.3.2-fr]{4.6.3.2}), on peut se ramener au cas où $Y$ est
affine, donc $\mathcal{L}$ ample~; il est clair que l'homomorphisme
(\hyperref[II.4.6.3.2-fr]{4.6.3.2}) est injectif, et le critère
(\hyperref[II.4.5.2-fr]{4.5.2, c)}) montre qu'il est surjectif, d'où la
conclusion.

\begin{corollary}[4.6.4]
\label{II.4.6.4-fr}
Soit $(U_\alpha)$ un recouvrement ouvert de $Y$. Pour que $\mathcal{L}$
soit ample relativement à $Y$, il faut et il suffit que
$\mathcal{L}|f^{-1}(U_\alpha)$ soit ample relativement à $U_\alpha$ pour
tout $\alpha$.
\end{corollary}

La condition b) est en effet locale sur $Y$.

\begin{corollary}[4.6.5]
\label{II.4.6.5-fr}
Soit $\mathcal{K}$ un $\mathcal{O}_Y$-Module inversible. Pour que
$\mathcal{L}$ soit $Y$-ample, il faut et il suffit que
$\mathcal{L}\otimes f^*(\mathcal{K})$ le soit.
\end{corollary}

C'est une conséquence évidente de (\hyperref[II.4.6.4-fr]{4.6.4}) en
prenant les $U_\alpha$ tels que $\mathcal{K}|U_\alpha$ soit isomorphe à
$\mathcal{O}_Y|U_\alpha$ pour tout $\alpha$.

\begin{corollary}[4.6.6]
\label{II.4.6.6-fr}
Supposons $Y$ affine~; pour que $\mathcal{L}$ soit $Y$-ample, il faut et il
suffit que $\mathcal{L}$ soit ample.
\end{corollary}

C'est une conséquence immédiate de la définition
(\hyperref[II.4.6.1-fr]{4.6.1}), et des critères
(\hyperref[II.4.6.3-fr]{4.6.3, b)}) et
(\hyperref[II.4.5.2-fr]{4.5.2, b)}), car ici
$\operatorname{Proj}(\mathcal{S})
=\operatorname{Proj}(\Gamma(Y,\mathcal{S}))$ par définition.

\begin{corollary}[4.6.7]
\label{II.4.6.7-fr}
Soit $f:X\to Y$ un morphisme quasi-compact. Supposons qu'il existe un
$\mathcal{O}_Y$-Module quasi-cohérent $\mathcal{E}$ et un $Y$-morphisme
$g:X\to P=\mathbf{P}(\mathcal{E})$ qui soit un homéomorphisme de l'espace
sous-jacent à $X$ sur un sous-espace de $P$~; alors
$\mathcal{L}=g^*(\mathcal{O}_P(1))$ est $Y$-ample.
\end{corollary}

On peut en effet supposer $Y$ affine~; le corollaire résulte alors du
critère (\hyperref[II.4.5.2-fr]{4.5.2, a)}), de la formule
(\hyperref[II.3.7.3.1-fr]{3.7.3.1}) et de
(\hyperref[II.4.2.3-fr]{4.2.3}).

\begin{proposition}[4.6.8]
\label{II.4.6.8-fr}
Soient $X$ un schéma quasi-compact ou un préschéma dont l'espace
sous-jacent est noethérien, $f:X\to Y$ un morphisme séparé quasi-compact.
Pour qu'un $\mathcal{O}_X$-Module inversible $\mathcal{L}$ soit $f$-ample,
il faut et il suffit que l'une des conditions équivalentes suivantes soit
vérifiée~:
\begin{enumerate}
  \item[\rm{c)}] Pour tout $\mathcal{O}_X$-Module quasi-cohérent
  $\mathcal{F}$ de type fini, il existe un entier $n_0>0$ tel que, pour
  tout $n\geq n_0$, l'homomorphisme canonique
  $\sigma:f^*(f_*(\mathcal{F}\otimes\mathcal{L}^{\otimes n}))\to
  \mathcal{F}\otimes\mathcal{L}^{\otimes n}$ soit surjectif.
  \item[\rm{c')}] Pour tout faisceau d'idéaux quasi-cohérent de type fini
  $\mathcal{J}$ de $\mathcal{O}_X$, il existe un entier $n>0$ tel que
  l'homomorphisme canonique
  $\sigma:f^*(f_*(\mathcal{J}\otimes\mathcal{L}^{\otimes n}))\to
  \mathcal{J}\otimes\mathcal{L}^{\otimes n}$ soit surjectif.
\end{enumerate}
\end{proposition}

Comme $X$ est quasi-compact, il en est de même de $f(X)$ et il existe donc
un recouvrement fini $(U_i)$ de $f(X)$ par des ouverts affines de $Y$. Pour
démontrer la condition c) lorsque $\mathcal{L}$ est $f$-ample, on peut
remplacer $Y$ par les $U_i$ et $X$ par les $f^{-1}(U_i)$, car si on obtient
pour chaque $i$ un entier $n_i$ tel que c) ait lieu (pour $U_i$,
$f^{-1}(U_i)$ et $\mathcal{L}|f^{-1}(U_i)$) dès que $n\geq n_i$, il
suffira de prendre pour $n_0$ le plus grand des $n_i$ pour obtenir c) pour
$Y$, $X$ et $\mathcal{L}$. Or lorsque $Y$ est affine, la condition c)
découle de (\hyperref[II.4.5.5-fr]{4.5.5, d)}), compte tenu de
(\hyperref[II.4.6.6-fr]{4.6.6}). Il est trivial que c) entraîne c'). Enfin,
pour prouver que c') implique que $\mathcal{L}$ est $f$-ample, on peut
encore se borner au cas où $Y$ est affine~: en effet, tout faisceau
d'idéaux quasi-cohérent de type fini $\mathcal{J}_i$ de
$\mathcal{O}_X|f^{-1}(U_i)$ est la restriction d'un faisceau cohérent
d'idéaux de type fini de $\mathcal{O}_X$
(\hyperref[I.9.4.7-fr]{I, 9.4.7}) et l'hypothèse c') entraîne que

\oldpage[II]{91}

$\mathcal{J}_i\otimes(\mathcal{L}^{\otimes n}|f^{-1}(U_i))$ est engendré
par ses sections (compte tenu de (\hyperref[I.9.2.2-fr]{I, 9.2.2}) et de
(\hyperref[II.3.4.7-fr]{3.4.7}))~; il suffit donc d'appliquer le critère
(\hyperref[II.4.5.5-fr]{4.5.5, d'')}).

\begin{proposition}[4.6.9]
\label{II.4.6.9-fr}
Soient $f:X\to Y$ un morphisme quasi-compact, $\mathcal{L}$ un
$\mathcal{O}_X$-Module inversible.
\begin{enumerate}
  \item[\rm{(i)}] Soit $n$ un entier $>0$. Pour que $\mathcal{L}$ soit
  $f$-ample, il faut et il suffit que $\mathcal{L}^{\otimes n}$ soit
  $f$-ample.
  \item[\rm{(ii)}] Soit $\mathcal{L}'$ un $\mathcal{O}_X$-Module
  inversible, et supposons qu'il existe un entier $n>0$ tel que
  l'homomorphisme canonique
  $\sigma:f^*(f_*(\mathcal{L}'^{\otimes n}))\to
  \mathcal{L}'^{\otimes n}$ soit surjectif. Alors, si $\mathcal{L}$ est
  $f$-ample, il en est de même de $\mathcal{L}\otimes\mathcal{L}'$.
\end{enumerate}
\end{proposition}

On se ramène en effet aussitôt au cas où $Y$ est affine, et la proposition
est alors conséquence immédiate de (\hyperref[II.4.5.6-fr]{4.5.6}).

\begin{corollary}[4.6.10]
\label{II.4.6.10-fr}
Le produit tensoriel de deux $\mathcal{O}_X$-Modules $f$-amples est
$f$-ample.
\end{corollary}

\begin{proposition}[4.6.11]
\label{II.4.6.11-fr}
Soient $Y$ un préschéma quasi-compact, $f:X\to Y$ un morphisme de type
fini, $\mathcal{L}$ un $\mathcal{O}_X$-Module inversible. Pour que
$\mathcal{L}$ soit $f$-ample, il faut et il suffit qu'il possède une des
propriétés équivalentes suivantes~:
\begin{enumerate}
  \item[\rm{d)}] Il existe $n_0>0$ tel que, pour tout entier $n\geq n_0$,
  $\mathcal{L}^{\otimes n}$ soit très ample relativement à $f$.
  \item[\rm{d')}] Il existe $n>0$ tel que $\mathcal{L}^{\otimes n}$ soit
  très ample relativement à $f$.
\end{enumerate}
\end{proposition}

Si $\mathcal{L}$ est ample relativement à $f$, il y a un recouvrement fini
$(U_i)$ de $Y$ par des ouverts affines tel que les
$\mathcal{L}|f^{-1}(U_i)$ soient amples. On en conclut
(\hyperref[II.4.5.10-fr]{4.5.10}) qu'il existe un entier $n_0$ tel que
$\mathcal{L}^{\otimes n}|f^{-1}(U_i)$ soit très ample relativement à
$f^{-1}(U_i)\to U_i$ pour tout $n\geq n_0$ et tout $i$, donc
$\mathcal{L}^{\otimes n}$ est très ample relativement à $f$
(\hyperref[II.4.4.5-fr]{4.4.5}). Inversement, d') entraîne déjà que
$\mathcal{L}^{\otimes n}$ est $f$-ample
(\hyperref[II.4.6.2-fr]{4.6.2}), donc il en est de même de $\mathcal{L}$
(\hyperref[II.4.6.9-fr]{4.6.9, (i)}).

\begin{corollary}[4.6.12]
\label{II.4.6.12-fr}
Soient $Y$ un préschéma quasi-compact, $f:X\to Y$ un morphisme de type
fini, $\mathcal{L}$, $\mathcal{L}'$ deux $\mathcal{O}_X$-Modules
inversibles. Si $\mathcal{L}$ est $f$-ample, il existe $n_0$ tel que
$\mathcal{L}^{\otimes n}\otimes\mathcal{L}'$ soit très ample relativement
à $f$ pour tout $n\geq n_0$.
\end{corollary}

On raisonne comme dans (\hyperref[II.4.6.11-fr]{4.6.11}) en utilisant un
recouvrement ouvert affine fini de $Y$ et
(\hyperref[II.4.5.11-fr]{4.5.11}).

\begin{proposition}[4.6.13]
\label{II.4.6.13-fr}
\begin{enumerate}
  \item[\rm{(i)}] Pour tout préschéma $Y$, tout
  $\mathcal{O}_Y$-Module inversible $\mathcal{L}$ est ample relativement
  au morphisme identique $1_Y$.
  \item[\rm{(i bis)}] Soient $f:X\to Y$ un morphisme quasi-compact,
  $j:X'\to X$ un morphisme quasi-compact qui est un homéomorphisme de
  l'espace sous-jacent à $X'$ sur un sous-espace de $X$. Si $\mathcal{L}$
  est un $\mathcal{O}_X$-Module ample relativement à $f$,
  $j^*(\mathcal{L})$ est ample relativement à $f\circ j$.
  \item[\rm{(ii)}] Soient $Z$ un préschéma quasi-compact, $f:X\to Y$,
  $g:Y\to Z$ deux morphismes quasi-compacts, $\mathcal{L}$ un
  $\mathcal{O}_X$-Module ample relativement à $f$, $\mathcal{K}$ un
  $\mathcal{O}_Y$-Module ample relativement à $g$. Alors il existe un
  entier $n_0>0$ tel que
  $\mathcal{L}\otimes f^*(\mathcal{K}^{\otimes n})$ soit ample relativement
  à $g\circ f$ pour tout $n\geq n_0$.
  \item[\rm{(iii)}] Soient $f:X\to Y$ un morphisme quasi-compact,
  $g:Y'\to Y$ un morphisme, et posons $X'=X_{(Y')}$. Si $\mathcal{L}$ est
  un $\mathcal{O}_X$-Module ample relativement à $f$,
  $\mathcal{L}'=\mathcal{L}\otimes_Y\mathcal{O}_{Y'}$ est un
  $\mathcal{O}_{X'}$-Module ample relativement à $f_{(Y')}$.
  \item[\rm{(iv)}] Soient $f_i:X_i\to Y_i$ ($i=1,2$) deux
  $S$-morphismes quasi-compacts. Si $\mathcal{L}_i$ est un
  $\mathcal{O}_{X_i}$-Module ample relativement à $f_i$ ($i=1,2$),
  $\mathcal{L}_1\otimes_S\mathcal{L}_2$ est ample relativement à
  $f_1\times_S f_2$.
  \item[\rm{(v)}] Soient $f:X\to Y$, $g:Y\to Z$ deux morphismes tels que
  $g\circ f$ soit quasi-compact. Si un $\mathcal{O}_X$-Module
  $\mathcal{L}$ est ample relativement à $g\circ f$, et si $g$ est séparé
  ou l'espace sous-jacent à $X$ localement noethérien, alors $\mathcal{L}$
  est ample relativement à $f$.

  \oldpage[II]{92}

  \item[\rm{(vi)}] Soient $f:X\to Y$ un morphisme quasi-compact, $j$
  l'injection canonique $X_{\mathrm{red}}\to X$. Si $\mathcal{L}$ est un
  $\mathcal{O}_X$-Module ample relativement à $f$, alors
  $j^*(\mathcal{L})$ est ample relativement à $f_{\mathrm{red}}$.
\end{enumerate}
\end{proposition}

Notons d'abord que (v) et (vi) se dérivent de (i), (i bis) et (iv) par le
même raisonnement que dans (\hyperref[II.4.4.10-fr]{4.4.10}), en utilisant
(\hyperref[II.4.6.4-fr]{4.6.4}) au lieu de
(\hyperref[II.4.4.5-fr]{4.4.5})~; nous laissons le détail du raisonnement au
lecteur. L'assertion (i) est trivialement conséquence de
(\hyperref[II.4.4.10-fr]{4.4.10, (i)}) et de
(\hyperref[II.4.6.2-fr]{4.6.2}). Pour démontrer (i bis), (iii) et (iv), nous
utiliserons le lemme suivant~:

\begin{lemma}[4.6.13.1]
\label{II.4.6.13.1-fr}
\begin{enumerate}
  \item[\rm{(i)}] Soient $u:Z\to S$ un morphisme, $\mathcal{L}$ un
  $\mathcal{O}_S$-Module inversible, $s$ une section de $\mathcal{L}$
  au-dessus de $S$, $s'$ la section de $u^*(\mathcal{L})=\mathcal{L}'$
  au-dessus de $Z$ qui lui correspond canoniquement. On a alors
  $Z_{s'}=u^{-1}(S_s)$.
  \item[\rm{(ii)}] Soient $Z$, $Z'$ deux $S$-préschémas, $p$, $p'$ les
  projections de $T=Z\times_S Z'$, $\mathcal{L}$ (resp.
  $\mathcal{L}'$) un $\mathcal{O}_Z$-Module (resp. un
  $\mathcal{O}_{Z'}$-Module) inversible, $t$ (resp. $t'$) une section de
  $\mathcal{L}$ (resp. $\mathcal{L}'$) au-dessus de $Z$ (resp. $Z'$), $s$
  (resp. $s'$) la section de $p^*(\mathcal{L})$ (resp.
  $p'^*(\mathcal{L}')$) au-dessus de $Z\times_S Z'$ qui lui correspond.
  Alors on a $T_{s\otimes s'}=Z_t\times_S Z'_{t'}$.
\end{enumerate}
\end{lemma}

Il résulte des définitions que l'on peut se ramener au cas où tous les
préschémas considérés sont affines. En outre, dans (i), on peut supposer
$\mathcal{L}=\mathcal{O}_S$~; l'assertion (i) résulte alors aussitôt de
(\hyperref[I.1.2.2.2-fr]{I, 1.2.2.2}). De même, dans (ii), on peut se
limiter au cas où $\mathcal{L}=\mathcal{O}_Z$,
$\mathcal{L}'=\mathcal{O}_{Z'}$, et alors l'assertion se réduit au lemme
(\hyperref[II.4.3.2.4-fr]{4.3.2.4}).

Démontrons alors (i bis). On peut supposer $Y$ affine
(\hyperref[II.4.6.4-fr]{4.6.4}), donc $\mathcal{L}$ ample
(\hyperref[II.4.6.6-fr]{4.6.6})~; lorsque $s$ parcourt l'ensemble réunion
des $\Gamma(X,\mathcal{L}^{\otimes n})$ ($n>0$), les $X_s$ forment une
base de la topologie de $X$ (\hyperref[II.4.5.2-fr]{4.5.2, a)}), donc par
hypothèse les $j^{-1}(X_s)$ forment une base de la topologie de $X'$~; on
conclut donc par le lemme
(\hyperref[II.4.6.13.1-fr]{4.6.13.1, (i)}) et par
(\hyperref[II.4.5.2-fr]{4.5.2, a)}) que $j^*(\mathcal{L})$ est ample.

Démontrons ensuite (iii). On peut de nouveau supposer $Y$ et $Y'$ affines
(\hyperref[II.4.6.4-fr]{4.6.4}), d'où il résulte que la projection
$h:X'\to X$ est affine (\hyperref[II.1.5.5-fr]{1.5.5}). Comme
$\mathcal{L}$ est ample (\hyperref[II.4.6.6-fr]{4.6.6}), lorsque $s$
parcourt l'ensemble des sections des $\mathcal{L}^{\otimes n}$ ($n>0$)
au-dessus de $X$ telles que $X_s$ soit affine, les $X_s$ recouvrent $X$
(\hyperref[II.4.5.2-fr]{4.5.2, a')})~; donc les $h^{-1}(X_s)$ sont affines
(\hyperref[II.1.2.5-fr]{1.2.5}) et recouvrent $X'$~; il résulte donc encore
du lemme (\hyperref[II.4.6.13.1-fr]{4.6.13.1, (i)}) et de
(\hyperref[II.4.5.2-fr]{4.5.2, a')}) que $\mathcal{L}'$ est ample, le
morphisme $f_{(Y')}$ étant quasi-compact
(\hyperref[I.6.6.4-fr]{I, 6.6.4, (iii)}).

Pour prouver (iv), notons d'abord que $f_1\times_S f_2$ est quasi-compact
(\hyperref[I.6.6.4-fr]{I, 6.6.4, (iv)}). On peut en outre supposer $S$,
$Y_1$ et $Y_2$ affines ((\hyperref[II.4.6.4-fr]{4.6.4}) et
(\hyperref[I.3.2.7-fr]{I, 3.2.7})), donc $\mathcal{L}_i$ ample ($i=1,2$)
(\hyperref[II.4.6.6-fr]{4.6.6}). Les ouverts
$(X_1)_{s_1}\times_S(X_2)_{s_2}$ forment un recouvrement de
$X_1\times_S X_2$ lorsque $s_i$ parcourt les sections des
$\mathcal{L}_i^{\otimes n_i}$ telles que $(X_i)_{s_i}$ soit affine
(\hyperref[II.4.5.2-fr]{4.5.2, a')}). D'ailleurs, en remplaçant $s_1$ et
$s_2$ par des puissances convenables, ce qui ne change pas les
$(X_i)_{s_i}$, on peut toujours supposer $n_1=n_2$. On déduit alors de
(\hyperref[II.4.6.13.1-fr]{4.6.13.1, (ii)}) et de
(\hyperref[II.4.5.2-fr]{4.5.2, a')}) que
$\mathcal{L}_1\otimes_S\mathcal{L}_2$ est ample, d'où l'assertion, puisque
$Y_1\times_S Y_2$ est affine (\hyperref[II.4.6.6-fr]{4.6.6}).

Reste à prouver (ii). Par le même raisonnement que dans
(\hyperref[II.4.4.10-fr]{4.4.10}), mais utilisant ici
(\hyperref[II.4.6.4-fr]{4.6.4}), on peut se borner au cas où $Z$ est
affine. Comme $\mathcal{K}$ est alors ample, et $Y$ quasi-compact, il existe
un nombre fini de sections $s_i\in\Gamma(Y,\mathcal{K}^{\otimes k_i})$
telles que les $Y_{s_i}$

\oldpage[II]{93}

soient affines et recouvrent $Y$
(\hyperref[II.4.5.2-fr]{4.5.2, a')})~;
remplaçant les $s_i$ par des puissances convenables, on peut en outre
supposer tous les $k_i$ égaux à un même entier $k$. Soient $s'_i$ les
sections de $f^*(\mathcal{K}^{\otimes k})$ au-dessus de $X$ correspondant
canoniquement aux $s_i$, de sorte que les
$X_{s'_i}=f^{-1}(Y_{s_i})$ (4.6.1.13, (i))
recouvrent $X$. Comme $\mathcal{L}|X_{s'_i}$ est ample
((\hyperref[II.4.6.4-fr]{4.6.4}) et
(\hyperref[II.4.6.6-fr]{4.6.6})), il existe pour chaque $i$ des sections en
nombre fini $t_{ij}\in\Gamma(X,\mathcal{L}^{\otimes n_{ij}})$ tels que les
$X_{t_{ij}}$ soient affines, contenus dans $X_{s'_i}$ et recouvrent
$X_{s'_i}$ (\hyperref[II.4.5.2-fr]{4.5.2, a')})~; on peut d'ailleurs
supposer tous les $n_{ij}$ égaux à un même nombre $n$. Cela étant, $X$ est
séparé et quasi-compact, donc il existe un entier $m>0$, et pour tout
$(i,j)$ une section
\[
  u_{ij}\in\Gamma(X,\mathcal{L}^{\otimes n}\otimes
  f^*(\mathcal{K}^{\otimes mk}))
\]
tels que $t_{ij}\otimes s_i'^{\otimes m}$ soit la restriction à $X_{s'_i}$
de $u_{ij}$ (\hyperref[I.9.3.1-fr]{I, 9.3.1})~; en outre on a
$X_{u_{ij}}=X_{t_{ij}}$, donc les $X_{u_{ij}}$ sont affines et recouvrent
$X$. On peut d'ailleurs supposer que $m$ est de la forme $nr$~; si on pose
$n_0=rk$, on voit donc (\hyperref[II.4.5.2-fr]{4.5.2, a')}) que
$\mathcal{L}\otimes_{\mathcal{O}_X}f^*(\mathcal{K}^{\otimes n_0})$ est
ample. En outre,
il existe $h_0>0$ tel que $\mathcal{K}^{\otimes h}$ soit engendré par ses
sections au-dessus de $Y$ dès que $h\geq h_0$
(\hyperref[II.4.5.5-fr]{4.5.5})~; a fortiori,
$f^*(\mathcal{K}^{\otimes h})$ est engendré par ses sections au-dessus de
$X$ pour $h\geq h_0$, par définition des images réciproques
((\hyperref[0.3.7.1-fr]{0, 3.7.1}) et
(\hyperref[II.4.4.1-fr]{4.4.1})). On en conclut que
$\mathcal{L}\otimes f^*(\mathcal{K}^{\otimes n_0+h})$ est ample dès que
$h\geq h_0$ (\hyperref[II.4.5.6-fr]{4.5.6}), ce qui achève la
démonstration.

\begin{remark}[4.6.14]
\label{II.4.6.14-fr}
Sous les conditions de (ii), on se gardera de croire que
$\mathcal{L}\otimes f^*(\mathcal{K})$ soit ample pour $g\circ f$~; en effet,
comme $\mathcal{L}\otimes f^*(\mathcal{K}^{-1})$ est aussi ample pour $f$
(\hyperref[II.4.6.5-fr]{4.6.5}), on en conclurait que $\mathcal{L}$ serait
ample pour $g\circ f$~; prenant en particulier pour $g$ le morphisme
identique, tout $\mathcal{O}_X$-Module inversible serait ample pour $f$, ce
qui n'est pas le cas en général (voir (\hyperref[II.5.1.6-fr]{5.1.6}),
(\hyperref[II.5.3.4-fr]{5.3.4, (i)}) et
(\hyperref[II.5.3.1-fr]{5.3.1})).
\end{remark}

\begin{proposition}[4.6.15]
\label{II.4.6.15-fr}
Soient $f:X\to Y$ un morphisme quasi-compact, $\mathcal{J}$ un faisceau
quasi-cohérent d'idéaux localement nilpotent de $\mathcal{O}_X$, $Z$ le
sous-préschéma fermé de $X$ défini par $\mathcal{J}$, $j:Z\to X$
l'injection canonique. Pour qu'un $\mathcal{O}_X$-Module inversible
$\mathcal{L}$ soit ample pour $f$, il faut et il suffit que
$j^*(\mathcal{L})$ soit ample pour $f\circ j$.
\end{proposition}

En effet, la question étant locale sur $Y$
(\hyperref[II.4.6.4-fr]{4.6.4}), on peut supposer $Y$ affine~; $X$ étant
alors quasi-compact, on peut supposer $\mathcal{J}$ nilpotent. Compte tenu
de (\hyperref[II.4.6.6-fr]{4.6.6}), la proposition n'est autre alors que
(\hyperref[II.4.5.13-fr]{4.5.13}).

\begin{corollary}[4.6.16]
\label{II.4.6.16-fr}
Soient $X$ un préschéma localement noethérien, $f:X\to Y$ un morphisme
quasi-compact. Pour qu'un $\mathcal{O}_X$-Module inversible $\mathcal{L}$
soit ample pour $f$, il faut et il suffit que son image réciproque
$\mathcal{L}'$ par l'injection canonique $X_{\mathrm{red}}\to X$ soit ample
pour $f_{\mathrm{red}}$.
\end{corollary}

On a déjà vu que la condition est nécessaire
(\hyperref[II.4.6.13-fr]{4.6.13, (vi)})~; inversement, si elle est remplie,
on peut se borner, pour prouver que $\mathcal{L}$ est ample pour $f$, au
cas où $Y$ est affine (\hyperref[II.4.6.4-fr]{4.6.4})~; alors
$Y_{\mathrm{red}}$ est aussi affine, donc $\mathcal{L}'$ est ample
(\hyperref[II.4.6.6-fr]{4.6.6}), et il en est de même de $\mathcal{L}$ en
vertu de (\hyperref[II.4.5.13-fr]{4.5.13}), puisque alors $X$ est
noethérien et $X_{\mathrm{red}}$ un sous-préschéma fermé de $X$ défini par
un faisceau quasi-cohérent nilpotent d'idéaux
(\hyperref[I.6.1.6-fr]{I, 6.1.6}).

\begin{proposition}[4.6.17]
\label{II.4.6.17-fr}
Les notations et hypothèses étant celles de
(\hyperref[II.4.4.11-fr]{4.4.11}), pour que $\mathcal{L}''$ soit ample
relativement à $f''$, il faut et il suffit que $\mathcal{L}$ soit ample
relativement à $f$ et $\mathcal{L}'$ ample relativement à $f'$.
\end{proposition}

\oldpage[II]{94}

La nécessité de la condition résulte de
(\hyperref[II.4.6.13-fr]{4.6.13, (i bis)}). Pour voir que la condition est
suffisante, on peut se borner au cas où $Y$ est affine, et alors le fait
que $\mathcal{L}''$ est ample résulte du critère
(\hyperref[II.4.5.2-fr]{4.5.2, a)}) appliqué à $\mathcal{L}$,
$\mathcal{L}'$ et $\mathcal{L}''$, en observant qu'une section de
$\mathcal{L}$ au-dessus de $X$ se prolonge (par $0$) en une section de
$\mathcal{L}''$ au-dessus de $X''$.

\begin{proposition}[4.6.18]
\label{II.4.6.18-fr}
Soient $Y$ un préschéma quasi-compact, $\mathcal{S}$ une
$\mathcal{O}_Y$-Algèbre graduée quasi-cohérente de type fini,
$X=\operatorname{Proj}(\mathcal{S})$, $f:X\to Y$ le morphisme structural.
Alors $f$ est de type fini, et il existe un entier $d>0$ tel que
$\mathcal{O}_X(d)$ soit inversible et $f$-ample.
\end{proposition}

En vertu de (\hyperref[II.3.1.10-fr]{3.1.10}), il existe un entier $d>0$
tel que $\mathcal{S}^{(d)}$ soit engendrée par $\mathcal{S}_d$. On sait que,
dans l'isomorphisme canonique entre $X$ et
$X^{(d)}=\operatorname{Proj}(\mathcal{S}^{(d)})$, $\mathcal{O}_X(d)$
s'identifie à $\mathcal{O}_{X^{(d)}}(1)$
(\hyperref[II.3.2.9-fr]{3.2.9, (ii)}). On voit donc qu'on est ramené au cas
où $\mathcal{S}$ est engendrée par $\mathcal{S}_1$~; la proposition résulte
alors de (\hyperref[II.4.4.3-fr]{4.4.3}) et de
(\hyperref[II.4.6.2-fr]{4.6.2}) (compte tenu de ce que $f$ est un morphisme
de type fini (\hyperref[II.3.4.1-fr]{3.4.1})).

\section{Morphismes quasi-affines~; morphismes quasi-projectifs. Morphismes
propres~; morphismes projectifs}

\subsection{Morphismes quasi-affines.}

\begin{definition}[5.1.1]
\label{II.5.1.1-fr}
On appelle \emph{schéma quasi-affine} un schéma isomorphe à un sous-schéma
induit sur un ouvert quasi-compact d'un schéma affine. On dit qu'un morphisme
$f:X\to Y$ est \emph{quasi-affine}, ou encore que $X$ (considéré comme
$Y$-préschéma au moyen de $f$) est un \emph{$Y$-schéma quasi-affine}, s'il
existe un recouvrement $(U_\alpha)$ de $Y$ par des ouverts affines tels que
les $f^{-1}(U_\alpha)$ soient des schémas quasi-affines.
\end{definition}

Il est clair qu'un morphisme quasi-affine est \emph{séparé}
(\hyperref[I.5.5.5-fr]{I, 5.5.5} et \hyperref[I.5.5.8-fr]{5.5.8}) et
\emph{quasi-compact} (\hyperref[I.6.6.1-fr]{I, 6.6.1})~; tout morphisme
affine est évidemment quasi-affine.

Rappelons que pour tout préschéma $X$, si on pose
$A=\Gamma(X,\mathcal{O}_X)$, l'homomorphisme identique
$A\to A=\Gamma(X,\mathcal{O}_X)$ définit un morphisme
$X\to\operatorname{Spec}(A)$ dit \emph{canonique}
(\hyperref[I.2.2.4-fr]{I, 2.2.4})~; ce n'est autre d'ailleurs que le morphisme
canonique défini dans (\hyperref[II.4.5.1-fr]{4.5.1}) pour le cas particulier
où $\mathcal{L}=\mathcal{O}_X$, si on se souvient que
$\operatorname{Proj}(A[T])$ s'identifie canoniquement à
$\operatorname{Spec}(A)$ (\hyperref[II.3.1.7-fr]{3.1.7}).

\begin{proposition}[5.1.2]
\label{II.5.1.2-fr}
Soient $X$ un schéma quasi-compact ou un préschéma dont l'espace sous-jacent
est noethérien, $A$ l'anneau $\Gamma(X,\mathcal{O}_X)$. Les conditions
suivantes sont équivalentes~:
\begin{enumerate}
  \item[\rm{a)}] $X$ est un schéma quasi-affine.
  \item[\rm{b)}] Le morphisme canonique
  $u:X\to\operatorname{Spec}(A)$ est une immersion ouverte.
  \item[\rm{b')}] Le morphisme canonique
  $u:X\to\operatorname{Spec}(A)$ est un homéomorphisme de $X$ sur un
  sous-espace de l'espace sous-jacent à $\operatorname{Spec}(A)$.
  \item[\rm{c)}] Le $\mathcal{O}_X$-Module $\mathcal{O}_X$ est très ample
  relativement à $u$ (\hyperref[II.4.4.2-fr]{4.4.2}).
  \item[\rm{c')}] Le $\mathcal{O}_X$-Module $\mathcal{O}_X$ est ample
  (\hyperref[II.4.5.1-fr]{4.5.1}).
  \item[\rm{d)}] Lorsque $f$ parcourt $A$, les $X_f$ forment une base de la
  topologie de $X$.
  \item[\rm{d')}] Lorsque $f$ parcourt $A$, ceux des $X_f$ qui sont affines
  forment un recouvrement de $X$.

\oldpage[II]{95}

  \item[\rm{e)}] Tout $\mathcal{O}_X$-Module quasi-cohérent est engendré par
  ses sections au-dessus de $X$.
  \item[\rm{e')}] Tout faisceau quasi-cohérent d'idéaux de type fini de
  $\mathcal{O}_X$ est engendré par ses sections au-dessus de $X$.
\end{enumerate}
\end{proposition}

Il est clair que b) entraîne a), et a) entraîne c) en vertu du critère b) de
(\hyperref[II.4.4.4-fr]{4.4.4}) appliqué au morphisme identique (compte tenu
de la remarque précédant l'énoncé)~; d'autre part c) entraîne c')
(\hyperref[II.4.5.10-fr]{4.5.10, (i)}), et c'), b) et b') sont équivalents par
(\hyperref[II.4.5.2-fr]{4.5.2, critères b) et b')}). Enfin c') équivaut à
chacun des critères d), d'), e), e') en vertu de
(\hyperref[II.4.5.2-fr]{4.5.2, critères a), a'), c)}) et
(\hyperref[II.4.5.5-fr]{4.5.5, critère d'')}).

On observera en outre qu'avec les notations précédentes, ceux des $X_f$ qui
sont affines forment une \emph{base} de la topologie de $X$, et que le
morphisme canonique $u$ est \emph{dominant}
(\hyperref[II.4.5.2-fr]{4.5.2}).

\begin{corollary}[5.1.3]
\label{II.5.1.3-fr}
Soit $X$ un préschéma quasi-compact. S'il existe un morphisme $v:X\to Y$ de
$X$ dans un schéma affine $Y$, qui soit un homéomorphisme de $X$ sur un
sous-espace ouvert de $Y$, alors $X$ est quasi-affine.
\end{corollary}

En effet, il existe une famille $(g_\alpha)$ de sections de $\mathcal{O}_Y$
au-dessus de $Y$ telles que les $D(g_\alpha)$ forment une base de la
topologie de $v(X)$~; si $v=(\psi,\theta)$ et si l'on pose
$f_\alpha=\theta(g_\alpha)$, on a
$X_{f_\alpha}=\psi^{-1}(D(g_\alpha))$
(\hyperref[I.2.2.4.1-fr]{I, 2.2.4.1}), donc les $X_{f_\alpha}$ forment une
base de la topologie de $X$, et le critère d) de
(\hyperref[II.5.1.2-fr]{5.1.2}) est vérifié.

\begin{corollary}[5.1.4]
\label{II.5.1.4-fr}
Si $X$ est un schéma quasi-affine, tout $\mathcal{O}_X$-Module inversible est
très ample (relativement au morphisme canonique) et a fortiori ample.
\end{corollary}

En effet, un tel Module $\mathcal{L}$ est engendré par ses sections au-dessus
de $X$ (\hyperref[II.5.1.2-fr]{5.1.2, e)}), donc
$\mathcal{L}\otimes\mathcal{O}_X=\mathcal{L}$ est très ample
(\hyperref[II.4.4.8-fr]{4.4.8}).

\begin{corollary}[5.1.5]
\label{II.5.1.5-fr}
Soit $X$ un préschéma quasi-compact. S'il existe un $\mathcal{O}_X$-Module
inversible $\mathcal{L}$ tel que $\mathcal{L}$ et $\mathcal{L}^{-1}$ soient
amples, $X$ est un schéma quasi-affine.
\end{corollary}

En effet, $\mathcal{O}_X=\mathcal{L}\otimes\mathcal{L}^{-1}$ est alors ample
(\hyperref[II.4.5.7-fr]{4.5.7}).

\begin{proposition}[5.1.6]
\label{II.5.1.6-fr}
Soit $f:X\to Y$ un morphisme quasi-compact. Les conditions suivantes sont
équivalentes~:
\begin{enumerate}
  \item[\rm{a)}] $f$ est quasi-affine.
  \item[\rm{b)}] La $\mathcal{O}_Y$-Algèbre
  $f_*(\mathcal{O}_X)=\mathcal{A}(X)$ est quasi-cohérente et le morphisme
  canonique $X\to\operatorname{Spec}(\mathcal{A}(X))$ correspondant à
  l'homomorphisme identique $\mathcal{A}(X)\to\mathcal{A}(X)$
  (\hyperref[II.1.2.7-fr]{1.2.7}) est une immersion ouverte.
  \item[\rm{b')}] La $\mathcal{O}_Y$-Algèbre $\mathcal{A}(X)$ est
  quasi-cohérente et le morphisme canonique
  $X\to\operatorname{Spec}(\mathcal{A}(X))$ est un homéomorphisme de $X$
  sur un sous-espace de $\operatorname{Spec}(\mathcal{A}(X))$.
  \item[\rm{c)}] Le $\mathcal{O}_X$-Module $\mathcal{O}_X$ est très ample
  pour $f$.
  \item[\rm{c')}] Le $\mathcal{O}_X$-Module $\mathcal{O}_X$ est ample pour
  $f$.
  \item[\rm{d)}] Le morphisme $f$ est séparé et, pour tout
  $\mathcal{O}_X$-Module quasi-cohérent $\mathcal{F}$, l'homomorphisme
  canonique $\sigma:f^*(f_*(\mathcal{F}))\to\mathcal{F}$
  (\hyperref[0.4.4.3-fr]{0, 4.4.3}) est surjectif.
\end{enumerate}
En outre, lorsque $f$ est quasi-affine, tout $\mathcal{O}_X$-Module
inversible $\mathcal{L}$ est très ample relativement à $f$.
\end{proposition}

L'équivalence de a) et c') résulte du caractère local sur $Y$ de la
$f$-amplitude (\hyperref[II.4.6.4-fr]{4.6.4}), de la définition
(\hyperref[II.5.1.1-fr]{5.1.1}) et du critère
(\hyperref[II.5.1.2-fr]{5.1.2, c')}). Les autres propriétés sont locales sur
$Y$

\oldpage[II]{96}

et résultent donc aussitôt de (\hyperref[II.5.1.2-fr]{5.1.2}) et
(\hyperref[II.5.1.4-fr]{5.1.4}), compte tenu de ce que
$f_*(\mathcal{F})$ est quasi-cohérent lorsque $f$ est séparé
(\hyperref[I.9.2.2-fr]{I, 9.2.2, a)}).

\begin{corollary}[5.1.7]
\label{II.5.1.7-fr}
Soit $f:X\to Y$ un morphisme quasi-affine. Pour tout ouvert $U$ de $Y$, la
restriction $f^{-1}(U)\to U$ de $f$ est quasi-affine.
\end{corollary}

\begin{corollary}[5.1.8]
\label{II.5.1.8-fr}
Soient $Y$ un schéma affine, $f:X\to Y$ un morphisme quasi-compact. Pour que
$f$ soit quasi-affine, il faut et il suffit que $X$ soit un schéma
quasi-affine.
\end{corollary}

C'est une conséquence immédiate de (\hyperref[II.5.1.6-fr]{5.1.6}) et
(\hyperref[II.4.6.6-fr]{4.6.6}).

\begin{corollary}[5.1.9]
\label{II.5.1.9-fr}
Soient $Y$ un schéma quasi-compact ou un préschéma dont l'espace sous-jacent
est noethérien, $f:X\to Y$ un morphisme de type fini. Si $f$ est
quasi-affine, il existe une sous-$\mathcal{O}_Y$-Algèbre quasi-cohérente
$\mathcal{B}$ de $\mathcal{A}(X)=f_*(\mathcal{O}_X)$, de type fini
(\hyperref[I.9.6.2-fr]{I, 9.6.2}), telle que le morphisme
$X\to\operatorname{Spec}(\mathcal{B})$ correspondant à l'injection canonique
$\mathcal{B}\to\mathcal{A}(X)$ (\hyperref[II.1.2.7-fr]{1.2.7}) soit une
immersion. En outre, toute sous-$\mathcal{O}_Y$-Algèbre quasi-cohérente de
type fini $\mathcal{B}'$ de $\mathcal{A}(X)$, contenant $\mathcal{B}$,
possède la même propriété.
\end{corollary}

En effet, $\mathcal{A}(X)$ est limite inductive de ses
sous-$\mathcal{O}_Y$-Algèbres quasi-cohérentes de type fini
(\hyperref[I.9.6.5-fr]{I, 9.6.5})~; le résultat est alors un cas particulier
de (\hyperref[II.3.8.4-fr]{3.8.4}), compte tenu de l'identification de
$\operatorname{Spec}(\mathcal{A}(X))$ et de
$\operatorname{Proj}(\mathcal{A}(X)[T])$
(\hyperref[II.3.1.7-fr]{3.1.7}).

\begin{proposition}[5.1.10]
\label{II.5.1.10-fr}
\begin{enumerate}
  \item[\rm{(i)}] Un morphisme quasi-compact $X\to Y$ qui est un
  homéomorphisme de l'espace sous-jacent à $X$ sur un sous-espace de
  l'espace sous-jacent à $Y$ (en particulier une immersion fermée) est
  quasi-affine.
  \item[\rm{(ii)}] Le composé de deux morphismes quasi-affines est
  quasi-affine.
  \item[\rm{(iii)}] Si $f:X\to Y$ est un $S$-morphisme quasi-affine,
  $f_{(S')}:X_{(S')}\to Y_{(S')}$ est un morphisme quasi-affine pour toute
  extension $S'\to S$ du préschéma de base.
  \item[\rm{(iv)}] Si $f:X\to Y$ et $g:X'\to Y'$ sont deux $S$-morphismes
  quasi-affines, $f\times_S g$ est quasi-affine.
  \item[\rm{(v)}] Si $f:X\to Y$, $g:Y\to Z$ sont deux morphismes tels que
  $g\circ f$ soit quasi-affine, et si $g$ est séparé ou l'espace sous-jacent
  à $X$ localement noethérien, $f$ est quasi-affine.
  \item[\rm{(vi)}] Si $f$ est un morphisme quasi-affine, il en est de même
  de $f_{\mathrm{red}}$.
\end{enumerate}
\end{proposition}

Compte tenu du critère (\hyperref[II.5.1.6-fr]{5.1.6, c')}), (i), (iii),
(iv), (v) et (vi) découlent aussitôt de
(\hyperref[II.4.6.13-fr]{4.6.13, (i bis), (iii), (iv), (v) et (vi)})
respectivement. Pour démontrer (ii), on peut se borner au cas où $Z$ est
affine, et alors l'assertion résulte directement de
(\hyperref[II.4.6.13-fr]{4.6.13, (ii)}), appliqué à
$\mathcal{L}=\mathcal{O}_X$ et $\mathcal{K}=\mathcal{O}_Y$.

\begin{remark}[5.1.11]
\label{II.5.1.11-fr}
Soient $f:X\to Y$, $g:Y\to Z$ deux morphismes tels que $X\times_ZY$ soit
localement noethérien. Alors l'immersion graphe
$\Gamma_f:X\to X\times_ZY$ est quasi-affine, étant quasi-compacte
(\hyperref[I.6.3.5-fr]{I, 6.3.5}), et
(\hyperref[I.5.5.12-fr]{I, 5.5.12}) montre que, dans (v), la conclusion reste
valable si on supprime l'hypothèse que $g$ est séparé.
\end{remark}

\begin{proposition}[5.1.12]
\label{II.5.1.12-fr}
Soient $f:X\to Y$ un morphisme quasi-compact, $g:X'\to X$ un morphisme
quasi-affine. Si $\mathcal{L}$ est un $\mathcal{O}_X$-Module ample pour
$f$, $g^*(\mathcal{L})$ est un $\mathcal{O}_{X'}$-Module ample pour
$f\circ g$.
\end{proposition}

En effet, comme $\mathcal{O}_{X'}$ est très ample pour $g$, et que la
question est locale sur $Y$ (\hyperref[II.4.6.4-fr]{4.6.4}), il résulte de
(\hyperref[II.4.6.13-fr]{4.6.13, (ii)}) qu'il existe (pour $Y$ affine) un
entier $n$ tel que $g^*(\mathcal{L})^{\otimes n}$ soit ample pour
$f\circ g$, donc $g^*(\mathcal{L})$ est ample pour $f\circ g$
(\hyperref[II.4.6.9-fr]{4.6.9}).

\oldpage[II]{97}

\subsection{Le critère de Serre.}

\begin{theorem}[5.2.1]
\label{II.5.2.1-fr}
\emph{(Critère de Serre.)} Soit $X$ un schéma quasi-compact ou un préschéma
dont l'espace sous-jacent est noethérien. Les conditions suivantes sont
équivalentes~:
\begin{enumerate}
  \item[\rm{a)}] $X$ est un schéma affine.
  \item[\rm{b)}] Il existe une famille d'éléments
  $f_\alpha\in A=\Gamma(X,\mathcal{O}_X)$ tels que les $X_{f_\alpha}$ soient
  affines et que l'idéal engendré par les $f_\alpha$ dans $A$ soit égal à
  $A$.
  \item[\rm{c)}] Le foncteur $\Gamma(X,\mathcal{F})$ est exact en
  $\mathcal{F}$ dans la catégorie des $\mathcal{O}_X$-Modules
  quasi-cohérents~; autrement dit, si
  \[
    0\to\mathcal{F}'\to\mathcal{F}\to\mathcal{F}''\to0 \tag{*}
  \]
  est une suite exacte de $\mathcal{O}_X$-Modules quasi-cohérents, la suite
  \[
    0\to\Gamma(X,\mathcal{F}')\to\Gamma(X,\mathcal{F})
      \to\Gamma(X,\mathcal{F}'')\to0
  \]
  est exacte.
  \item[\rm{c')}] La propriété c) a lieu pour toute suite exacte (*) de
  $\mathcal{O}_X$-Modules quasi-cohérents telle que $\mathcal{F}$ soit
  isomorphe à un sous-$\mathcal{O}_X$-Module d'un produit fini
  $\mathcal{O}_X^n$.
  \item[\rm{d)}] $\mathrm{H}^1(X,\mathcal{F})=0$ pour tout
  $\mathcal{O}_X$-Module quasi-cohérent $\mathcal{F}$.
  \item[\rm{d')}] $\mathrm{H}^1(X,\mathcal{J})=0$ pour tout faisceau
  quasi-cohérent d'idéaux $\mathcal{J}$ de $\mathcal{O}_X$.
\end{enumerate}
\end{theorem}

Il est évident que a) implique b)~; b) implique d'autre part que les
$X_{f_\alpha}$ recouvrent $X$, puisque par hypothèse la section $1$ est
combinaison linéaire des $f_\alpha$ et que les $D(f_\alpha)$ recouvrent
alors $\operatorname{Spec}(A)$. La dernière assertion de
(\hyperref[II.4.5.2-fr]{4.5.2}) implique donc que
$X\to\operatorname{Spec}(A)$ est un isomorphisme.

On sait que a) implique c) (\hyperref[I.1.3.11-fr]{I, 1.3.11}) et c)
entraîne trivialement c'). Prouvons que c') implique b). Tout d'abord, c')
entraîne que, pour tout point fermé $x\in X$ et tout voisinage ouvert $U$ de
$x$, il existe $f\in A$ tel que $x\in X_f\subset X-U$. Soit en effet
$\mathcal{J}$ (resp. $\mathcal{J}'$) le faisceau quasi-cohérent d'idéaux de
$\mathcal{O}_X$ définissant le sous-préschéma fermé réduit de $X$ ayant pour
espace sous-jacent $X-U$ (resp. $(X-U)\cup\{x\}$)
(\hyperref[I.5.2.1-fr]{I, 5.2.1})~; il est clair que l'on a
$\mathcal{J}'\subset\mathcal{J}$ et que
$\mathcal{J}''=\mathcal{J}/\mathcal{J}'$ est un $\mathcal{O}_X$-Module
quasi-cohérent ayant pour support $\{x\}$ et tel que
$\mathcal{J}''_x=k(x)$. L'hypothèse c') appliquée à la suite exacte
$0\to\mathcal{J}'\to\mathcal{J}\to\mathcal{J}''\to0$ montre que
$\Gamma(X,\mathcal{J})\to\Gamma(X,\mathcal{J}'')$ est surjectif. La section
de $\mathcal{J}''$ dont le germe en $x$ est $1_x$ est donc l'image d'une
section $f\in\Gamma(X,\mathcal{J})\subset\Gamma(X,\mathcal{O}_X)$, et on a
par définition $f(x)=1_x$ et $f(y)=0$ dans $X-U$, ce qui établit notre
assertion. D'ailleurs, si $U$ est affine, il en est de même de $X_f$
(\hyperref[I.1.3.6-fr]{I, 1.3.6}), donc la réunion des $X_f$ qui sont affines
($f\in A$) est un ensemble ouvert $Z$ contenant tous les points fermés de
$X$~; comme $X$ est un espace de Kolmogoroff quasi-compact, on a
nécessairement $Z=X$ (\hyperref[0.2.1.3-fr]{0, 2.1.3}). Puisque $X$ est
quasi-compact, il y a un nombre fini d'éléments $f_i\in A$
($1\leq i\leq n$) tels que les $X_{f_i}$ soient affines et recouvrent $X$.
Considérons alors l'homomorphisme $\mathcal{O}_X^n\to\mathcal{O}_X$ défini
par les sections $f_i$ (\hyperref[0.5.1.1-fr]{0, 5.1.1})~; comme, pour tout
$x\in X$, un au moins des $(f_i)_x$ est inversible, cet homomorphisme est
surjectif, et on a donc une suite exacte
$0\to\mathcal{R}\to\mathcal{O}_X^n\to\mathcal{O}_X\to0$, où $\mathcal{R}$
est un sous-$\mathcal{O}_X$-Module quasi-cohérent de $\mathcal{O}_X^n$. Il
résulte alors

\oldpage[II]{98}

de c') que l'homomorphisme correspondant
$\Gamma(X,\mathcal{O}_X^n)\to\Gamma(X,\mathcal{O}_X)$ est surjectif, ce qui
démontre b).

Enfin, a) implique d) (\hyperref[I.5.1.9.2-fr]{I, 5.1.9.2}) et d) implique
trivialement d'). Reste à montrer que d') entraîne c'). Or, si
$\mathcal{F}'$ est un sous-$\mathcal{O}_X$-Module quasi-cohérent de
$\mathcal{O}_X^n$, la filtration
$0\subset\mathcal{O}_X\subset\mathcal{O}_X^2\subset\cdots\subset
\mathcal{O}_X^n$ définit sur $\mathcal{F}'$ une filtration formée des
$\mathcal{F}'_k=\mathcal{O}_X^k\cap\mathcal{F}'$ ($0\leq k\leq n$), qui
sont des $\mathcal{O}_X$-Modules quasi-cohérents
(\hyperref[I.4.1.1-fr]{I, 4.1.1}), et
$\mathcal{F}'_{k+1}/\mathcal{F}'_k$ est isomorphe à un
sous-$\mathcal{O}_X$-Module quasi-cohérent de
$\mathcal{O}_X^{k+1}/\mathcal{O}_X^k=\mathcal{O}_X$, c'est-à-dire à un
faisceau quasi-cohérent d'idéaux de $\mathcal{O}_X$. L'hypothèse d')
entraîne donc
$\mathrm{H}^1(X,\mathcal{F}'_{k+1}/\mathcal{F}'_k)=0$~; la suite exacte de
cohomologie
$\mathrm{H}^1(X,\mathcal{F}'_k)\to\mathrm{H}^1(X,\mathcal{F}'_{k+1})
\to\mathrm{H}^1(X,\mathcal{F}'_{k+1}/\mathcal{F}'_k)=0$ permet alors de
prouver par récurrence sur $k$ que
$\mathrm{H}^1(X,\mathcal{F}'_k)=0$ pour tout $k$. C.Q.F.D.

\begin{remark}[5.2.1.1]
\label{II.5.2.1.1-fr}
Lorsque $X$ est un préschéma noethérien, on peut dans l'énoncé de c') et de
d'), remplacer «~quasi-cohérent~» par «~cohérent~». En effet, dans la
démonstration du fait que c') entraîne b), $\mathcal{J}$ et $\mathcal{J}'$
sont alors des faisceaux cohérents d'idéaux, et, d'autre part, tout
sous-Module quasi-cohérent d'un Module cohérent est cohérent
(\hyperref[I.6.1.1-fr]{I, 6.1.1})~; d'où la conclusion.
\end{remark}

\begin{corollary}[5.2.2]
\label{II.5.2.2-fr}
Soit $f:X\to Y$ un morphisme séparé quasi-compact. Les conditions suivantes
sont équivalentes~:
\begin{enumerate}
  \item[\rm{a)}] $f$ est un morphisme affine.
  \item[\rm{b)}] Le foncteur $f_*$ est exact dans la catégorie des
  $\mathcal{O}_X$-Modules quasi-cohérents.
  \item[\rm{c)}] Pour tout $\mathcal{O}_X$-Module quasi-cohérent
  $\mathcal{F}$, on a $\mathrm{R}^1f_*(\mathcal{F})=0$.
  \item[\rm{c')}] Pour tout faisceau d'idéaux quasi-cohérent $\mathcal{J}$
  de $\mathcal{O}_X$, on a $\mathrm{R}^1f_*(\mathcal{J})=0$.
\end{enumerate}
\end{corollary}

Toutes ces conditions étant locales sur $Y$, par définition du foncteur
$\mathrm{R}^1f_*$ (T, 3.7.3), on peut supposer que $Y$ est affine. Si $f$
est affine, $X$ est alors affine et la propriété b) n'est autre que
(\hyperref[I.1.6.4-fr]{I, 1.6.4}). Inversement, montrons que b) implique
a)~: pour tout $\mathcal{O}_X$-Module quasi-cohérent $\mathcal{F}$,
$f_*(\mathcal{F})$ est un $\mathcal{O}_Y$-Module quasi-cohérent
(\hyperref[I.9.2.2-fr]{I, 9.2.2, a)}). Par hypothèse le foncteur
$f_*(\mathcal{F})$ est exact en $\mathcal{F}$, et le foncteur
$\Gamma(Y,\mathcal{G})$ est exact en $\mathcal{G}$ (dans la catégorie des
$\mathcal{O}_Y$-Modules quasi-cohérents) puisque $Y$ est affine
(\hyperref[I.1.3.11-fr]{I, 1.3.11})~; donc
$\Gamma(Y,f_*(\mathcal{F}))=\Gamma(X,\mathcal{F})$ est exact en
$\mathcal{F}$, ce qui prouve notre assertion en vertu de
(\hyperref[II.5.2.1-fr]{5.2.1, c)}).

Si $f$ est affine, $f^{-1}(U)$ est affine pour tout ouvert affine $U$ de
$Y$ (\hyperref[II.1.3.2-fr]{1.3.2}), donc
$\mathrm{H}^1(f^{-1}(U),\mathcal{F})=0$
(\hyperref[II.5.2.1-fr]{5.2.1, d)}), ce qui par définition entraîne
$\mathrm{R}^1f_*(\mathcal{F})=0$. Supposons enfin vérifiée la condition
c')~; la suite exacte des termes de bas degré dans la suite spectrale de
Leray (G, II, 4.17.1 et I, 4.5.1) donne en particulier la suite exacte
\[
  0\to\mathrm{H}^1(Y,f_*(\mathcal{J}))\to
  \mathrm{H}^1(X,\mathcal{J})\to
  \mathrm{H}^0(Y,\mathrm{R}^1f_*(\mathcal{J})).
\]
Comme $Y$ est affine et $f_*(\mathcal{J})$ quasi-cohérent
(\hyperref[I.9.2.2-fr]{I, 9.2.2, a)}), on a
$\mathrm{H}^1(Y,f_*(\mathcal{J}))=0$
(\hyperref[II.5.2.1-fr]{5.2.1})~; l'hypothèse c') entraîne donc
$\mathrm{H}^1(X,\mathcal{J})=0$, et on conclut par
(\hyperref[II.5.2.1-fr]{5.2.1}) que $X$ est un schéma affine.

\begin{corollary}[5.2.3]
\label{II.5.2.3-fr}
Si $f:X\to Y$ est un morphisme affine, alors, pour tout
$\mathcal{O}_X$-Module quasi-cohérent $\mathcal{F}$, l'homomorphisme
canonique
$\mathrm{H}^1(Y,f_*(\mathcal{F}))\to\mathrm{H}^1(X,\mathcal{F})$ est
bijectif.
\end{corollary}

\oldpage[II]{99}

En effet, on a la suite exacte
\[
  0\to\mathrm{H}^1(Y,f_*(\mathcal{F}))\to
  \mathrm{H}^1(X,\mathcal{F})\to
  \mathrm{H}^0(Y,\mathrm{R}^1f_*(\mathcal{F}))
\]
des termes de bas degré de la suite spectrale de Leray, et la conclusion
résulte de (\hyperref[II.5.2.2-fr]{5.2.2}).

\begin{remark}[5.2.4]
\label{II.5.2.4-fr}
Au chapitre III, § 1, nous prouverons que si $X$ est affine, on a
$\mathrm{H}^i(X,\mathcal{F})=0$ pour tout $i>0$ et tout
$\mathcal{O}_X$-Module quasi-cohérent $\mathcal{F}$.
\end{remark}

\subsection{Morphismes quasi-projectifs.}

\begin{definition}[5.3.1]
\label{II.5.3.1-fr}
On dit qu'un morphisme $f:X\to Y$ est \emph{quasi-projectif}, ou que $X$
(considéré comme $Y$-préschéma au moyen de $f$) est \emph{quasi-projectif
sur $Y$}, ou que $X$ est un \emph{$Y$-schéma quasi-projectif}, si $f$ est de
type fini et s'il existe un $\mathcal{O}_X$-Module inversible qui est
$f$-ample.
\end{definition}

On notera que cette notion n'est pas locale sur $Y$~: les contre-exemples de
Nagata [26] et de Hironaka montrent que, même si $X$ et $Y$ sont des schémas
algébriques non singuliers sur un corps algébriquement clos, tout point de
$Y$ peut avoir un voisinage affine $U$ tel que $f^{-1}(U)$ soit
quasi-projectif sur $U$, sans que $f$ soit quasi-projectif.

On notera qu'un morphisme quasi-projectif est nécessairement séparé
(\hyperref[II.4.6.1-fr]{4.6.1}). Lorsque $Y$ est quasi-compact, il revient au
même de dire que $f$ est quasi-projectif ou que $f$ est de type fini et qu'il
existe un $\mathcal{O}_X$-Module très ample relativement à $f$
(\hyperref[II.4.6.2-fr]{4.6.2} et \hyperref[II.4.6.11-fr]{4.6.11}). De plus~:

\begin{proposition}[5.3.2]
\label{II.5.3.2-fr}
Soit $Y$ un schéma quasi-compact ou un préschéma dont l'espace sous-jacent
est noethérien, et soit $X$ un $Y$-préschéma. Les conditions suivantes sont
équivalentes~:
\begin{enumerate}
  \item[\rm{a)}] $X$ est un $Y$-schéma quasi-projectif.
  \item[\rm{b)}] $X$ est de type fini sur $Y$, et il existe un
  $\mathcal{O}_Y$-Module quasi-cohérent de type fini $\mathcal{E}$ tel que
  $X$ soit $Y$-isomorphe à un sous-préschéma de $\mathbf{P}(\mathcal{E})$.
  \item[\rm{c)}] $X$ est de type fini sur $Y$, et il existe une
  $\mathcal{O}_Y$-Algèbre graduée quasi-cohérente $\mathcal{S}$ telle que
  $\mathcal{S}_1$ soit de type fini et engendre $\mathcal{S}$ et que $X$
  soit $Y$-isomorphe à un sous-préschéma induit sur un ouvert partout dense
  de $\operatorname{Proj}(\mathcal{S})$.
\end{enumerate}
\end{proposition}

Cela résulte aussitôt de la remarque précédente et de
(\hyperref[II.4.4.3-fr]{4.4.3}), (\hyperref[II.4.4.6-fr]{4.4.6}) et
(\hyperref[II.4.4.7-fr]{4.4.7}).

On notera que lorsque $Y$ est un préschéma noethérien, on peut, dans les
conditions b) et c) de (\hyperref[II.5.3.2-fr]{5.3.2}), supprimer
l'hypothèse que $X$ est de type fini sur $Y$, qui est alors automatiquement
vérifiée (\hyperref[I.6.3.5-fr]{I, 6.3.5}).

\begin{corollary}[5.3.3]
\label{II.5.3.3-fr}
Soit $Y$ un schéma quasi-compact tel qu'il existe un $\mathcal{O}_Y$-Module
ample $\mathcal{L}$ (\hyperref[II.4.5.3-fr]{4.5.3}). Pour qu'un $Y$-schéma
$X$ soit quasi-projectif, il faut et il suffit qu'il soit de type fini sur
$Y$ et isomorphe à un sous-$Y$-schéma d'un fibré projectif de la forme
$\mathbf{P}_Y^r$.
\end{corollary}

En effet, si $\mathcal{E}$ est un $\mathcal{O}_Y$-Module quasi-cohérent de
type fini, $\mathcal{E}$ est isomorphe à un quotient d'un
$\mathcal{O}_Y$-Module
$\mathcal{L}^{\otimes(-n)}\otimes_{\mathcal{O}_Y}\mathcal{O}_Y^k$
(\hyperref[II.4.5.5-fr]{4.5.5}), donc $\mathbf{P}(\mathcal{E})$ est
isomorphe à un sous-schéma fermé de $\mathbf{P}_Y^{k-1}$
(\hyperref[II.4.1.2-fr]{4.1.2} et \hyperref[II.4.1.4-fr]{4.1.4}).

\begin{proposition}[5.3.4]
\label{II.5.3.4-fr}
\begin{enumerate}
  \item[\rm{(i)}] Un morphisme quasi-affine de type fini (et en particulier
  une immersion quasi-compacte, ou un morphisme affine de type fini) est
  quasi-projectif.
  \item[\rm{(ii)}] Si $f:X\to Y$ et $g:Y\to Z$ sont quasi-projectifs et si
  $Z$ est quasi-compact, $g\circ f$ est quasi-projectif.

\oldpage[II]{100}

  \item[\rm{(iii)}] Si $f:X\to Y$ est un $S$-morphisme quasi-projectif,
  $f_{(S')}:X_{(S')}\to Y_{(S')}$ est quasi-projectif pour toute extension
  $S'\to S$ du préschéma de base.
  \item[\rm{(iv)}] Si $f:X\to Y$ et $g:X'\to Y'$ sont deux $S$-morphismes
  quasi-projectifs, $f\times_Sg$ est quasi-projectif.
  \item[\rm{(v)}] Si $f:X\to Y$, $g:Y\to Z$ sont deux morphismes tels que
  $g\circ f$ soit quasi-projectif, et si $g$ est séparé ou $X$ localement
  noethérien, alors $f$ est quasi-projectif.
  \item[\rm{(vi)}] Si $f$ est un morphisme quasi-projectif, il en est de
  même de $f_{\mathrm{red}}$.
\end{enumerate}
\end{proposition}

(i) résulte de (\hyperref[II.5.1.6-fr]{5.1.6}) et
(\hyperref[II.5.1.10-fr]{5.1.10, (i)}). Les autres assertions sont des
conséquences immédiates de la définition (\hyperref[II.5.3.1-fr]{5.3.1}),
des propriétés des morphismes de type fini (\hyperref[I.6.3.4-fr]{I, 6.3.4})
et de (\hyperref[II.4.6.13-fr]{4.6.13}).

\begin{remark}[5.3.5]
\label{II.5.3.5-fr}
On notera qu'il peut se faire que $f_{\mathrm{red}}$ soit quasi-projectif
sans que $f$ le soit, même en supposant que $Y$ est le spectre d'une algèbre
de rang fini sur $\mathbf{C}$ et que $f$ est propre.
\end{remark}

\begin{corollary}[5.3.6]
\label{II.5.3.6-fr}
Si $X$ et $X'$ sont deux $Y$-schémas quasi-projectifs, $X\amalg X'$ est un
$Y$-schéma quasi-projectif.
\end{corollary}

Cela résulte de (\hyperref[II.4.6.18-fr]{4.6.18}).

\subsection{Morphismes propres et morphismes universellement fermés.}

\begin{definition}[5.4.1]
\label{II.5.4.1-fr}
On dit qu'un morphisme de préschémas $f:X\to Y$ est \emph{propre} s'il
vérifie les deux conditions suivantes~:
\begin{enumerate}
  \item[\rm{a)}] $f$ est séparé et de type fini.
  \item[\rm{b)}] Pour tout préschéma $Y'$ et tout morphisme $Y'\to Y$, la
  projection $f_{(Y')}:X\times_YY'\to Y'$ est un morphisme fermé
  (\hyperref[I.2.2.6-fr]{I, 2.2.6}).
\end{enumerate}
Lorsqu'il en est ainsi, on dit aussi que $X$ (considéré comme $Y$-préschéma
de morphisme structural $f$) est propre au-dessus de $Y$.
\end{definition}

Il est immédiat que les conditions a) et b) sont locales sur $Y$. Pour
vérifier que l'image d'une partie fermée $Z$ de $X\times_YY'$ par la
projection $q:X\times_YY'\to Y'$ est fermée dans $Y'$, il suffit de voir que
$q(Z)\cap U'$ est fermé dans $U'$ pour tout ouvert affine $U'$ de $Y'$~;
comme $q(Z)\cap U'=q(Z\cap q^{-1}(U'))$ et que $q^{-1}(U')$ s'identifie à
$X\times_YU'$ (\hyperref[I.4.4.1-fr]{I, 4.4.1}), on voit que pour vérifier
la condition b) de la déf. (\hyperref[II.5.4.1-fr]{5.4.1}), on peut se
limiter au cas où $Y'$ est un schéma affine. On verra plus loin
(\hyperref[II.5.6.3-fr]{5.6.3}) que si $Y$ est localement noethérien, on
peut même se borner à vérifier b) lorsque $Y'$ est de type fini sur $Y$.

Il est clair que tout morphisme propre est \emph{fermé}.

\begin{proposition}[5.4.2]
\label{II.5.4.2-fr}
\begin{enumerate}
  \item[\rm{(i)}] Une immersion fermée est un morphisme propre.
  \item[\rm{(ii)}] Le composé de deux morphismes propres est propre.
  \item[\rm{(iii)}] Si $X$, $Y$ sont des $S$-préschémas,
  $f:X\to Y$ un $S$-morphisme propre, alors
  $f_{(S')}:X_{(S')}\to Y_{(S')}$ est propre pour toute extension
  $S'\to S$ du préschéma de base.
  \item[\rm{(iv)}] Si $f:X\to Y$ et $g:X'\to Y'$ sont deux $S$-morphismes
  propres, le $S$-morphisme
  $f\times_Sg:X\times_SX'\to Y\times_SY'$ est propre.
\end{enumerate}
\end{proposition}

\oldpage[II]{101}

Il suffit de démontrer (i), (ii) et (iii) (\hyperref[I.3.5.1-fr]{I, 3.5.1}).
Dans chacun de ces trois cas, la vérification de la condition a) de
(\hyperref[II.5.4.1-fr]{5.4.1}) découle de résultats antérieurs
(\hyperref[I.5.5.1-fr]{I, 5.5.1} et \hyperref[I.6.3.4-fr]{6.3.4})~; reste à
vérifier la condition b). C'est immédiat dans le cas (i), car si $X\to Y$
est une immersion fermée, il en est de même de
$X\times_YY'\to Y\times_YY'=Y'$ (\hyperref[I.4.3.2-fr]{I, 4.3.2} et
\hyperref[II.3.3.3-fr]{3.3.3}). Pour démontrer (ii), considérons deux
morphismes propres $X\to Y$, $Y\to Z$, et un morphisme $Z'\to Z$. On peut
écrire $X\times_ZZ'=X\times_Y(Y\times_ZZ')$
(\hyperref[I.3.3.9.1-fr]{I, 3.3.9.1}), et par suite la projection
$X\times_ZZ'\to Z'$ se factorise en
$X\times_Y(Y\times_ZZ')\to Y\times_ZZ'\to Z'$. Compte tenu de la remarque
du début, (ii) résulte donc de ce que le composé de deux morphismes fermés
est fermé. Enfin, pour tout morphisme $S'\to S$, $X_{(S')}$ s'identifie à
$X\times_YY_{(S')}$ (\hyperref[I.3.3.11-fr]{I, 3.3.11})~; pour tout
morphisme $Z\to Y_{(S')}$, on peut écrire
\[
  X_{(S')}\times_{Y_{(S')}}Z
  =(X\times_YY_{(S')})\times_{Y_{(S')}}Z=X\times_YZ~;
\]
$X\times_YZ\to Z$ est fermé par hypothèse, donc cela prouve (iii).

\begin{corollary}[5.4.3]
\label{II.5.4.3-fr}
Soient $f:X\to Y$, $g:Y\to Z$ deux morphismes tels que $g\circ f$ soit
propre.
\begin{enumerate}
  \item[\rm{(i)}] Si $g$ est séparé, $f$ est propre.
  \item[\rm{(ii)}] Si $g$ est séparé et de type fini, et si $f$ est
  surjectif, $g$ est propre.
\end{enumerate}
\end{corollary}

(i) découle de (\hyperref[II.5.4.2-fr]{5.4.2}) par le procédé général
(\hyperref[I.5.5.12-fr]{I, 5.5.12}). Pour démontrer (ii), il n'y a à
vérifier que la condition b) de la déf. (\hyperref[II.5.4.1-fr]{5.4.1}).
Pour tout morphisme $Z'\to Z$, le diagramme
\[
  \xymatrix{
    X\times_ZZ'\ar[r]^{f\times1_{Z'}}\ar[dr]_p &
    Y\times_ZZ'\ar[d]^{p'}\\
    & Z'
  }
\]
(où $p$ et $p'$ sont les projections) est commutatif
(\hyperref[I.3.2.1-fr]{I, 3.2.1})~; en outre, $f\times1_{Z'}$ est surjectif
puisqu'il en est ainsi de $f$ (\hyperref[I.3.5.2-fr]{I, 3.5.2}), et $p$ est
un morphisme fermé par hypothèse. Toute partie fermée $F$ de
$Y\times_ZZ'$ est alors l'image par $f\times1_{Z'}$ d'une partie fermée $E$
de $X\times_ZZ'$, donc $p'(F)=p(E)$ est fermé dans $Z'$ par hypothèse, d'où
le corollaire.

\begin{corollary}[5.4.4]
\label{II.5.4.4-fr}
Si $X$ est un préschéma propre au-dessus de $Y$, et $\mathcal{S}$ une
$\mathcal{O}_Y$-Algèbre quasi-cohérente, tout $Y$-morphisme
$f:X\to\operatorname{Proj}(\mathcal{S})$ est propre (et a fortiori fermé).
\end{corollary}

En effet, le morphisme structural
$p:\operatorname{Proj}(\mathcal{S})\to Y$ est séparé, et $p\circ f$ est
propre par hypothèse.

\begin{corollary}[5.4.5]
\label{II.5.4.5-fr}
Soit $f:X\to Y$ un morphisme séparé de type fini. Soient
$(X_i)_{1\leq i\leq n}$ (resp. $(Y_i)_{1\leq i\leq n}$) une famille finie
de sous-préschémas fermés de $X$ (resp. $Y$), $j_i$ (resp. $h_i$)
l'injection canonique $X_i\to X$ (resp. $Y_i\to Y$). Supposons que l'espace
sous-jacent $X$ soit réunion des $X_i$, et que pour tout $i$, il y ait un
morphisme $f_i:X_i\to Y_i$ tel que le diagramme
\[
  \xymatrix{
    X_i\ar[r]^{f_i}\ar[d]_{j_i} & Y_i\ar[d]^{h_i}\\
    X\ar[r]_f & Y
  }
\]
soit commutatif. Alors, pour que $f$ soit propre, il faut et il suffit que
chacun des $f_i$ le soit.
\end{corollary}

\oldpage[II]{102}

Si $f$ est propre, il en est de même de $f\circ j_i$, puisque $j_i$ est une
immersion fermée (\hyperref[II.5.4.2-fr]{5.4.2})~; comme $h_i$ est une
immersion fermée, donc un morphisme séparé, $f_i$ est propre par
(\hyperref[II.5.4.3-fr]{5.4.3}). Supposons inversement que chacun des $f_i$
soit propre, et considérons le préschéma $Z$ somme des $X_i$, soit $u$ le
morphisme $Z\to X$ qui se réduit à $j_i$ dans chacun des $X_i$. La
restriction de $f\circ u$ à chaque $X_i$ est égale à
$f\circ j_i=h_i\circ f_i$, donc est propre puisqu'il en est ainsi de $h_i$
et de $f_i$ (\hyperref[II.5.4.2-fr]{5.4.2})~; il résulte alors aussitôt de
la déf. (\hyperref[II.5.4.1-fr]{5.4.1}) que $f\circ u$ est propre. Mais
comme $u$ est surjectif par hypothèse, on conclut que $f$ est propre par
(\hyperref[II.5.4.3-fr]{5.4.3}).

\begin{corollary}[5.4.6]
\label{II.5.4.6-fr}
Soit $f:X\to Y$ un morphisme séparé et de type fini~; pour que $f$ soit
propre, il faut et il suffit que
$f_{\mathrm{red}}:X_{\mathrm{red}}\to Y_{\mathrm{red}}$ le soit.
\end{corollary}

C'est un cas particulier de (\hyperref[II.5.4.5-fr]{5.4.5}) avec $n=1$,
$X_1=X_{\mathrm{red}}$ et $Y_1=Y_{\mathrm{red}}$
(\hyperref[I.5.1.5-fr]{I, 5.1.5}).

\begin{env}[5.4.7]
\label{II.5.4.7-fr}
Si $X$ et $Y$ sont des préschémas noethériens et $f:X\to Y$ un morphisme
séparé de type fini, on peut, pour vérifier que $f$ est propre, se ramener à
le faire pour des morphismes dominants de préschémas intègres. En effet,
soient $X_i$ ($1\leq i\leq n$) les composantes irréductibles en nombre fini
de $X$ et considérons pour chaque $i$ l'unique sous-préschéma fermé réduit de
$X$ ayant $X_i$ pour espace sous-jacent, que nous désignerons encore par
$X_i$ (\hyperref[I.5.2.1-fr]{I, 5.2.1}). Soit d'autre part $Y_i$ l'unique
sous-préschéma fermé réduit de $Y$ ayant $\overline{f(X_i)}$ pour espace
sous-jacent. Si $g_i$ (resp. $h_i$) est le morphisme d'injection $X_i\to X$
(resp. $Y_i\to Y$), on en conclut que $f\circ g_i=h_i\circ f_i$, où $f_i$
est un morphisme dominant $X_i\to Y_i$
(\hyperref[I.5.2.2-fr]{I, 5.2.2})~; on est alors dans les conditions
d'application de (\hyperref[II.5.4.5-fr]{5.4.5}), et pour que $f$ soit
propre, il faut et il suffit que les $f_i$ le soient.
\end{env}

\begin{corollary}[5.4.8]
\label{II.5.4.8-fr}
Soient $X$, $Y$ des $S$-préschémas séparés de type fini sur $S$, et soit
$f:X\to Y$ un $S$-morphisme. Pour que $f$ soit propre, il faut et il suffit
que, pour tout $S$-préschéma $S'$, le morphisme
$f\times_S1_{S'}:X\times_SS'\to Y\times_SS'$ soit fermé.
\end{corollary}

Notons d'abord que si $g:X\to S$ et $h:Y\to S$ sont les morphismes
structuraux, on a par définition $g=h\circ f$, donc $f$ est séparé et de
type fini (\hyperref[I.5.5.1-fr]{I, 5.5.1} et
\hyperref[I.6.3.4-fr]{6.3.4}). Si $f$ est propre, il en est de même de
$f\times_S1_{S'}$ (\hyperref[II.5.4.2-fr]{5.4.2})~; a fortiori,
$f\times_S1_{S'}$ est fermé. Inversement, supposons vérifiée la condition de
l'énoncé et soit $Y'$ un $Y$-préschéma~; $Y'$ peut aussi être considéré
comme un $S$-préschéma, et comme $Y\to S$ est séparé, $X\times_YY'$
s'identifie à un sous-préschéma fermé de $X\times_SY'$
(\hyperref[I.5.4.2-fr]{I, 5.4.2}). Dans le diagramme commutatif
\[
  \xymatrix{
    X\times_YY'\ar[r]^{f\times1_{Y'}}\ar[d] &
    Y\times_YY'=Y'\ar[d]\\
    X\times_SY'\ar[r]_{f\times_S1_{Y'}} & Y\times_SY'
  }
\]
les flèches verticales sont des immersions fermées~; il en résulte aussitôt
que si $f\times_S1_{Y'}$ est un morphisme fermé, il en est de même de
$f\times_Y1_{Y'}$.

\begin{remark}[5.4.9]
\label{II.5.4.9-fr}
Nous dirons qu'un morphisme $f:X\to Y$ est \emph{universellement fermé} s'il
satisfait à la condition b) de la définition
(\hyperref[II.5.4.1-fr]{5.4.1}). Le lecteur observera que

\oldpage[II]{103}

dans les propositions (\hyperref[II.5.4.2-fr]{5.4.2}) à
(\hyperref[II.5.4.8-fr]{5.4.8}), on peut partout remplacer «~propre~» par
«~universellement fermé~» sans modifier la validité des résultats (et dans
les hypothèses de (\hyperref[II.5.4.3-fr]{5.4.3}),
(\hyperref[II.5.4.5-fr]{5.4.5}), (\hyperref[II.5.4.6-fr]{5.4.6}) et
(\hyperref[II.5.4.8-fr]{5.4.8}), on peut omettre les conditions de
finitude).
\end{remark}

\begin{env}[5.4.10]
\label{II.5.4.10-fr}
Soit $f:X\to Y$ un morphisme de type fini. Nous dirons qu'une partie fermée
$Z$ de $X$ est \emph{propre sur $Y$} (ou \emph{$Y$-propre}, ou \emph{propre
pour $f$}) si la restriction de $f$ à un sous-préschéma fermé de $X$,
d'espace sous-jacent $Z$ (\hyperref[I.5.2.1-fr]{I, 5.2.1}), est propre.
Comme cette restriction est alors séparée, il résulte de
(\hyperref[II.5.4.6-fr]{5.4.6}) et (\hyperref[I.5.5.1-fr]{I, 5.5.1, (vi)})
que la propriété précédente ne dépend pas du sous-préschéma fermé de $X$
ayant $Z$ pour espace sous-jacent. Si $g:X'\to X$ est un morphisme propre,
$g^{-1}(Z)$ est une partie propre de $X'$~: si $T$ est un sous-préschéma de
$X$ ayant pour espace sous-jacent $Z$, il suffit en effet de remarquer que
la restriction de $g$ au sous-préschéma fermé $g^{-1}(T)$ de $X'$ est un
morphisme propre $g^{-1}(T)\to T$ par
(\hyperref[II.5.4.2-fr]{5.4.2, (iii)}), et d'appliquer ensuite
(\hyperref[II.5.4.2-fr]{5.4.2, (ii)}). D'autre part, si $X''$ est un
$Y$-schéma de type fini et $u:X\to X''$ un $Y$-morphisme, $u(Z)$ est une
partie propre de $X''$~; en effet, prenons pour $T$ le sous-préschéma fermé
réduit de $X$ ayant $Z$ pour espace sous-jacent~; la restriction de $f$ à
$T$ étant propre, il en est de même de la restriction de $u$ à $T$
(\hyperref[II.5.4.3-fr]{5.4.3, (i)}), donc $u(Z)$ est fermé dans $X''$~;
soit $T''$ un sous-préschéma fermé de $X''$ ayant $u(Z)$ pour espace
sous-jacent (\hyperref[I.5.2.1-fr]{I, 5.2.1}), de sorte que $u|T$ se
factorise en
\[
  T\xrightarrow{v}T''\xrightarrow{j}X'',
\]
où $j$ est l'injection canonique (\hyperref[I.5.2.2-fr]{I, 5.2.2}) et $v$
est donc propre et surjectif (\hyperref[II.5.4.5-fr]{5.4.5})~; si $g$ est la
restriction à $T''$ du morphisme structural $X''\to Y$, $g$ est séparé et
de type fini, et on a $f|T=g\circ v$~; il résulte donc de
(\hyperref[II.5.4.3-fr]{5.4.3, (ii)}) que $g$ est propre, d'où notre
assertion.
\end{env}

Il résulte en particulier de ces remarques que si $Z$ est une partie
$Y$-propre de $X$, alors~:
\begin{enumerate}
  \item[\rm{1\textsuperscript{o}}] Pour tout sous-préschéma fermé $X'$ de
  $X$, $Z\cap X'$ est une partie $Y$-propre de $X'$.
  \item[\rm{2\textsuperscript{o}}] Si $X$ est un sous-préschéma d'un
  $Y$-schéma de type fini $X''$, $Z$ est aussi une partie $Y$-propre de
  $X''$ (et en particulier est fermée dans $X''$).
\end{enumerate}

\subsection{Morphismes projectifs.}

\begin{proposition}[5.5.1]
\label{II.5.5.1-fr}
Soit $X$ un $Y$-préschéma. Les conditions suivantes sont équivalentes~:
\begin{enumerate}
  \item[\rm{a)}] $X$ est $Y$-isomorphe à un sous-préschéma fermé d'un fibré
  projectif $\mathbf{P}(\mathcal{E})$, où $\mathcal{E}$ est un
  $\mathcal{O}_Y$-Module quasi-cohérent de type fini.
  \item[\rm{b)}] Il existe une $\mathcal{O}_Y$-Algèbre graduée
  quasi-cohérente $\mathcal{S}$, telle que $\mathcal{S}_1$ soit de type fini
  et engendre $\mathcal{S}$, et que $X$ soit $Y$-isomorphe à
  $\operatorname{Proj}(\mathcal{S})$.
\end{enumerate}
\end{proposition}

La condition a) entraîne b) en vertu de
(\hyperref[II.3.6.2-fr]{3.6.2, (ii)})~: si $\mathcal{J}$ est un faisceau
gradué quasi-cohérent d'idéaux de $\mathbf{S}(\mathcal{E})$, la
$\mathcal{O}_Y$-Algèbre graduée quasi-cohérente
$\mathcal{S}=\mathbf{S}(\mathcal{E})/\mathcal{J}$ est engendrée par
$\mathcal{S}_1$ et ce dernier, image canonique de $\mathcal{E}$, est un
$\mathcal{O}_Y$-Module de type fini. La condition b) implique a) en vertu de
(\hyperref[II.3.6.2-fr]{3.6.2}) appliqué au cas où
$\mathcal{M}\to\mathcal{S}_1$ est l'application identique.

\oldpage[II]{104}

\begin{definition}[5.5.2]
\label{II.5.5.2-fr}
On dit qu'un $Y$-préschéma $X$ est \emph{projectif sur $Y$} ou est un
\emph{$Y$-schéma projectif} s'il vérifie les conditions équivalentes a) et b)
de (\hyperref[II.5.5.1-fr]{5.5.1}). On dit qu'un morphisme $f:X\to Y$ est
\emph{projectif} s'il fait de $X$ un $Y$-schéma projectif.
\end{definition}

Il est clair que si $f:X\to Y$ est projectif, il existe un
$\mathcal{O}_X$-Module \emph{très ample relativement à $f$}
(\hyperref[II.4.4.2-fr]{4.4.2}).

\begin{theorem}[5.5.3]
\label{II.5.5.3-fr}
\begin{enumerate}
  \item[\rm{(i)}] Tout morphisme projectif est quasi-projectif et propre.
  \item[\rm{(ii)}] Inversement, soit $Y$ un schéma quasi-compact ou un
  préschéma dont l'espace sous-jacent est noethérien~; alors tout morphisme
  $f:X\to Y$ qui est quasi-projectif et propre est projectif.
\end{enumerate}
\end{theorem}

(i) Il est clair que si $f:X\to Y$ est projectif, il est de type fini et
quasi-projectif (donc en particulier séparé)~; d'autre part, il résulte
aussitôt de (\hyperref[II.5.5.1-fr]{5.5.1, b)}) et de
(\hyperref[II.3.5.3-fr]{3.5.3}) que si $f$ est projectif, il en est de même de
$f\times_Y1_{Y'}:X\times_YY'\to Y'$ pour tout morphisme $Y'\to Y$. Pour
montrer que $f$ est universellement fermé, tout revient donc à prouver qu'un
morphisme projectif $f$ est \emph{fermé}. La question étant locale sur $Y$, on
peut supposer que $Y=\operatorname{Spec}(A)$, donc
(\hyperref[II.5.5.1-fr]{5.5.1}) $X=\operatorname{Proj}(S)$, où $S$ est une
$A$-algèbre graduée engendrée par un nombre fini d'éléments de $S_1$. Pour
tout $y\in Y$, la fibre $f^{-1}(y)$ s'identifie à
$\operatorname{Proj}(S)\times_Y\operatorname{Spec}(k(y))$
(\hyperref[I.3.6.1-fr]{I, 3.6.1}), donc à
$\operatorname{Proj}(S\otimes_Ak(y))$
(\hyperref[II.2.8.10-fr]{2.8.10})~; par suite, $f^{-1}(y)$ est vide si et
seulement si $S\otimes_Ak(y)$ vérifie la condition (TN)
(\hyperref[II.2.7.4-fr]{2.7.4}), autrement dit $S_n\otimes_Ak(y)=0$ pour $n$
assez grand. Mais comme $(S_n)_y$ est un $\mathcal{O}_y$-module de type fini,
la condition précédente signifie que $(S_n)_y=0$ pour $n$ assez grand, en
vertu du lemme de Nakayama. Si $\mathfrak a_n$ est l'annulateur dans $A$ du
$A$-module $S_n$, la condition précédente signifie encore que
$\mathfrak a_n\not\subset\mathfrak j_y$ pour $n$ assez grand
(\hyperref[0.1.7.4-fr]{0, 1.7.4}). Or, comme $S_nS_1=S_{n+1}$ par hypothèse,
on a $\mathfrak a_n\subset\mathfrak a_{n+1}$, et si $\mathfrak a$ est la
réunion des $\mathfrak a_n$, on voit donc que $f(X)=V(\mathfrak a)$, ce qui
prouve que $f(X)$ est fermé dans $Y$. Si maintenant $X'$ est une partie
fermée quelconque de $X$, il existe un sous-préschéma fermé de $X$ ayant $X'$
pour espace sous-jacent (\hyperref[I.5.2.1-fr]{I, 5.2.1}) et il est clair
(\hyperref[II.5.5.1-fr]{5.5.1, a)}) que le morphisme
$X'\to X\xrightarrow{f}Y$ est projectif, donc $f(X')$ est fermé dans $Y$.

(ii) L'hypothèse sur $Y$ et le fait que $f$ est quasi-projectif entraînent
l'existence d'un $\mathcal{O}_Y$-Module quasi-cohérent de type fini
$\mathcal{E}$ et d'une $Y$-immersion $j:X\to\mathbf{P}(\mathcal{E})$
(\hyperref[II.5.3.2-fr]{5.3.2}). Mais comme $f$ est propre, $j$ est fermée par
(\hyperref[II.5.4.4-fr]{5.4.4}), donc $f$ est projectif.

\par\medskip\noindent\textit{Remarques (5.5.4). ---\ }
\label{II.5.5.4-fr}
\begin{enumerate}
  \item[\rm{(i)}] Soit $f:X\to Y$ un morphisme tel que~:
  1\textsuperscript{o} $f$ soit propre~; 2\textsuperscript{o} il existe un
  $\mathcal{O}_X$-Module $\mathcal{L}$ très ample relativement à $f$~;
  3\textsuperscript{o} le $\mathcal{O}_Y$-Module quasi-cohérent
  $\mathcal{E}=f_*(\mathcal{L})$ soit de type fini. Alors $f$ est un morphisme
  projectif~: en effet (\hyperref[II.4.4.4-fr]{4.4.4}), il y a alors une
  $Y$-immersion $r:X\to\mathbf{P}(\mathcal{E})$, et comme $f$ est propre, $r$
  est une immersion fermée (\hyperref[II.5.4.4-fr]{5.4.4}). On verra au
  chapitre III, \S~3, que lorsque $Y$ est localement noethérien, la condition
  3\textsuperscript{o} ci-dessus est une conséquence des deux autres, donc les
  conditions 1\textsuperscript{o} et 2\textsuperscript{o} caractérisent dans
  ce cas les morphismes projectifs et si $Y$ est quasi-compact, on peut
  remplacer la condition 2\textsuperscript{o} par l'hypothèse de l'existence
  d'un $\mathcal{O}_X$-Module ample pour $f$
  (\hyperref[II.4.6.11-fr]{4.6.11}).
  \item[\rm{(ii)}] Soit $Y$ un schéma quasi-compact tel qu'il existe un
  $\mathcal{O}_Y$-Module ample. Pour qu'un $Y$-schéma $X$ soit projectif, il
  faut et il suffit qu'il soit $Y$-isomorphe à un sous-$Y$-schéma fermé d'un
  fibré projectif de la forme $\mathbf{P}_Y^r$. La condition est évidemment
  suffisante.

\oldpage[II]{105}

  Inversement, si $X$ est projectif sur $Y$, il est quasi-projectif, donc il
  existe une $Y$-immersion $j$ de $X$ dans un $\mathbf{P}_Y^r$
  (\hyperref[II.5.3.3-fr]{5.3.3}) qui est fermée par
  (\hyperref[II.5.4.4-fr]{5.4.4}) et
  (\hyperref[II.5.5.3-fr]{5.5.3}).
  \item[\rm{(iii)}] Le raisonnement de
  (\hyperref[II.5.5.3-fr]{5.5.3}) montre que, pour tout préschéma $Y$, et tout
  entier $r\geq0$, le morphisme structural $\mathbf{P}_Y^r\to Y$ est
  surjectif, car si on pose
  $\mathcal{S}=\mathbf{S}_{\mathcal{O}_Y}(\mathcal{O}_Y^{r+1})$, on a
  évidemment $\mathcal{S}_y=\mathbf{S}_{k(y)}((k(y))^{r+1})$
  (\hyperref[II.1.7.3-fr]{1.7.3}), donc $(\mathcal{S}_n)_y\neq0$ pour tout
  $y\in Y$ et tout $n\geq0$.
  \item[\rm{(iv)}] Il résulte des exemples de Nagata \cite{II-26} qu'il y a
  des morphismes propres qui ne sont pas quasi-projectifs.
\end{enumerate}

\begin{proposition}[5.5.5]
\label{II.5.5.5-fr}
\begin{enumerate}
  \item[\rm{(i)}] Une immersion fermée est un morphisme projectif.
  \item[\rm{(ii)}] Si $f:X\to Y$ et $g:Y\to Z$ sont des morphismes
  projectifs, et si $Z$ est un schéma quasi-compact ou un préschéma dont
  l'espace sous-jacent est noethérien, $g\circ f$ est projectif.
  \item[\rm{(iii)}] Si $f:X\to Y$ est un $S$-morphisme projectif,
  $f_{(S')}:X_{(S')}\to Y_{(S')}$ est projectif pour toute extension
  $S'\to S$ du préschéma de base.
  \item[\rm{(iv)}] Si $f:X\to Y$ et $g:X'\to Y'$ sont des $S$-morphismes
  projectifs, il en est de même de $f\times_Sg$.
  \item[\rm{(v)}] Si $g\circ f$ est un morphisme projectif et si $g$ est
  séparé, $f$ est projectif.
  \item[\rm{(vi)}] Si $f$ est projectif, il en est de même de
  $f_{\mathrm{red}}$.
\end{enumerate}
\end{proposition}

(i) résulte aussitôt de (\hyperref[II.3.1.7-fr]{3.1.7}). Il faut ici
démontrer (iii) et (iv) séparément, à cause de la restriction introduite sur
$Z$ dans (ii) (cf. (\hyperref[I.3.5.1-fr]{I, 3.5.1})). Pour prouver (iii), on
se ramène au cas où $S=Y$ (\hyperref[I.3.3.11-fr]{I, 3.3.11}) et l'assertion
résulte alors aussitôt de (\hyperref[II.5.5.1-fr]{5.5.1, b)}) et de
(\hyperref[II.3.5.3-fr]{3.5.3}). Pour démontrer (iv), on est aussitôt ramené
au cas où $X=\mathbf{P}(\mathcal{E})$, $X'=\mathbf{P}(\mathcal{E}')$,
$\mathcal{E}$ (resp. $\mathcal{E}'$) étant un $\mathcal{O}_Y$-Module (resp.
un $\mathcal{O}_{Y'}$-Module) quasi-cohérent de type fini. Soient $p$, $p'$ les
projections canoniques de $T=Y\times_SY'$ dans $Y$ et $Y'$ respectivement~;
d'après (\hyperref[II.4.1.3.1-fr]{4.1.3.1}), on a
$\mathbf{P}(p^*(\mathcal{E}))=\mathbf{P}(\mathcal{E})\times_YT$,
$\mathbf{P}(p'^*(\mathcal{E}'))=\mathbf{P}(\mathcal{E}')\times_{Y'}T$, d'où
\begin{align*}
\mathbf{P}(p^*(\mathcal{E}))\times_T\mathbf{P}(p'^*(\mathcal{E}'))
&=(\mathbf{P}(\mathcal{E})\times_YT)\times_T
(T\times_{Y'}\mathbf{P}(\mathcal{E}'))\\
&=\mathbf{P}(\mathcal{E})\times_Y
(T\times_{Y'}\mathbf{P}(\mathcal{E}'))
=\mathbf{P}(\mathcal{E})\times_S\mathbf{P}(\mathcal{E}')
\end{align*}
en remplaçant $T$ par $Y\times_SY'$ et utilisant
(\hyperref[I.3.3.9.1-fr]{I, 3.3.9.1}). Or, $p^*(\mathcal{E})$ et
$p'^*(\mathcal{E}')$ sont de type fini sur $T$
(\hyperref[0.5.2.4-fr]{0, 5.2.4}), donc il en est de même de
$p^*(\mathcal{E})\otimes_{\mathcal{O}_T}p'^*(\mathcal{E}')$~; comme
$\mathbf{P}(p^*(\mathcal{E}))\times_T\mathbf{P}(p'^*(\mathcal{E}'))$
s'identifie à un sous-préschéma fermé de
$\mathbf{P}(p^*(\mathcal{E})\otimes_{\mathcal{O}_T}p'^*(\mathcal{E}'))$
(\hyperref[II.4.3.3-fr]{4.3.3}), cela achève de prouver (iv). Pour obtenir (v)
et (vi), on peut appliquer (\hyperref[I.5.5.13-fr]{I, 5.5.13}), car tout
sous-préschéma fermé d'un $Y$-schéma projectif est un $Y$-schéma projectif par
(\hyperref[II.5.5.1-fr]{5.5.1, a)}).

Il reste à prouver (ii)~; en vertu de l'hypothèse sur $Z$, cela résulte de
(\hyperref[II.5.5.3-fr]{5.5.3}), de
(\hyperref[II.5.3.4-fr]{5.3.4, (ii)}) et
(\hyperref[II.5.4.2-fr]{5.4.2, (ii)}).

\begin{proposition}[5.5.6]
\label{II.5.5.6-fr}
Si $X$ et $X'$ sont deux $Y$-schémas projectifs, $X\amalg X'$ est un
$Y$-schéma projectif.
\end{proposition}

C'est une conséquence évidente de (\hyperref[II.5.5.2-fr]{5.5.2}) et
(\hyperref[II.4.3.6-fr]{4.3.6}).

\begin{proposition}[5.5.7]
\label{II.5.5.7-fr}
Soit $X$ un $Y$-schéma projectif, $\mathcal{L}$ un $\mathcal{O}_X$-Module
$Y$-ample~; pour toute section $f$ de $\mathcal{L}$ au-dessus de $X$, $X_f$
est affine sur $Y$.
\end{proposition}

\oldpage[II]{106}

La question étant locale sur $Y$, on peut supposer
$Y=\operatorname{Spec}(A)$~; en outre
$X_{f^{\otimes n}}=X_f$, donc en remplaçant $\mathcal{L}$ par un
$\mathcal{L}^{\otimes n}$ convenable, on peut supposer $\mathcal{L}$ très
ample pour le morphisme structural $q:X\to Y$
(\hyperref[II.4.6.11-fr]{4.6.11}). L'homomorphisme canonique
$\sigma:q^*(q_*(\mathcal{L}))\to\mathcal{L}$ est donc surjectif et le
morphisme correspondant
\[
r=r_{\mathcal{L},\sigma}:X\to P=\mathbf{P}(q_*(\mathcal{L}))
\]
une immersion telle que $\mathcal{L}=r^*(\mathcal{O}_P(1))$
(\hyperref[II.4.4.4-fr]{4.4.4})~; en outre, comme $X$ est propre sur $Y$,
l'immersion $r$ est fermée (\hyperref[II.5.4.4-fr]{5.4.4}). Par ailleurs, on
a par définition $f\in\Gamma(Y,q_*(\mathcal{L}))$ et $\sigma^\flat$ est
l'identité de $q_*(\mathcal{L})$~; il résulte alors de la formule
(\hyperref[II.3.7.3.1-fr]{3.7.3.1}) que l'on a $X_f=r^{-1}(D_+(f))$~; $X_f$
est par suite un sous-préschéma fermé du schéma affine $D_+(f)$ et est donc
bien un schéma affine.

Lorsqu'on prend en particulier $Y=X$, on obtient (compte tenu de
(\hyperref[II.4.6.13-fr]{4.6.13, (i)})) le corollaire suivant, dont la
démonstration directe est d'ailleurs immédiate~:

\begin{corollary}[5.5.8]
\label{II.5.5.8-fr}
Soit $X$ un préschéma, $\mathcal{L}$ un $\mathcal{O}_X$-Module inversible.
Pour toute section $f$ de $\mathcal{L}$ au-dessus de $X$, $X_f$ est affine
sur $X$ (et par suite un schéma affine lorsque $X$ est un schéma affine).
\end{corollary}

\subsection{Le lemme de Chow.}

\begin{theorem}[5.6.1]
\label{II.5.6.1-fr}
(Lemme de Chow). Soit $S$ un préschéma, $X$ un $S$-schéma de type fini. On
suppose vérifiée l'une des hypothèses suivantes~:
\begin{enumerate}
  \item[\rm{a)}] $S$ est noethérien.
  \item[\rm{b)}] $S$ est un schéma quasi-compact, et $X$ a un nombre fini de
  composantes irréductibles.
\end{enumerate}
Dans ces conditions~:
\begin{enumerate}
  \item[\rm{(i)}] Il existe un $S$-schéma $X'$ quasi-projectif sur $S$ et un
  $S$-morphisme $f:X'\to X$, projectif et surjectif.
  \item[\rm{(ii)}] On peut prendre $X'$ et $f$ tels qu'il existe un ouvert
  $U\subset X$ pour lequel $U'=f^{-1}(U)$ soit dense dans $X'$ et la
  restriction de $f$ à $U'$ un isomorphisme $U'\simeq U$.
  \item[\rm{(iii)}] Si $X$ est réduit (resp. irréductible, intègre), on peut
  supposer $X'$ réduit (resp. irréductible, intègre).
\end{enumerate}
\end{theorem}

La démonstration se fait en plusieurs étapes.

A) On peut d'abord se ramener au cas où $X$ est irréductible. En effet, dans
l'hypothèse a), $X$ est noethérien, donc, dans les deux hypothèses, les
composantes irréductibles $X_i$ de $X$ sont en nombre fini. Si le théorème est
démontré pour chacun des sous-préschémas fermés réduits de $X$ ayant les $X_i$
pour espaces sous-jacents, et si $X'_i$ et $f_i:X'_i\to X_i$ sont le préschéma
et le morphisme correspondant à $X_i$, le préschéma $X'$ somme des $X'_i$, et
le morphisme $f:X'\to X$ dont la restriction à chaque $X'_i$ est $j_i\circ f_i$
($j_i$ injection canonique $X_i\to X$), répondent à la question. Il est
immédiat en effet que $X'$ est réduit si les $X'_i$ le sont~; d'autre part on
satisfait à (ii) en prenant pour $U$ la réunion des ensembles
$U_i\cap(\bigcup_{j\neq i}X_j)^{\complement}$. Enfin, comme les $X'_i$ sont
quasi-projectifs sur $S$, il en est de même de $X'$

\oldpage[II]{107}

(\hyperref[II.5.3.6-fr]{5.3.6})~; de même, les morphismes $X'_i\to X$ sont
projectifs par (\hyperref[II.5.5.5-fr]{5.5.5, (i) et (ii)}), donc $f$ est
projectif (\hyperref[II.5.5.6-fr]{5.5.6}), et est évidemment surjectif par
définition.

B) Supposons maintenant $X$ irréductible. Comme le morphisme structural
$r:X\to S$ est de type fini, il y a un recouvrement fini $(S_i)$ de $S$ par
des ouverts affines, et pour chaque $i$ un recouvrement fini $(T_{ij})$ de
$r^{-1}(S_i)$ par des ouverts affines, le morphisme $T_{ij}\to S_i$ étant
affine et de type fini, donc quasi-projectif
(\hyperref[II.5.3.4-fr]{5.3.4, (i)})~; comme dans chacune des hypothèses a),
b), l'immersion $S_i\to S$ est quasi-compacte, c'est un morphisme
quasi-projectif (\hyperref[II.5.3.4-fr]{5.3.4, (ii)}), donc la restriction de
$r$ à $T_{ij}$ est un morphisme quasi-projectif
(\hyperref[II.5.3.4-fr]{5.3.4, (iii)}). Désignons les $T_{ij}$ par $U_k$
$(1\leq k\leq n)$. Il existe, pour chaque indice $k$ une immersion ouverte
$\varphi_k:U_k\to P_k$, où $P_k$ est projectif sur $S$
(\hyperref[II.5.3.2-fr]{5.3.2} et \hyperref[II.5.5.2-fr]{5.5.2}). Soit
$U=\bigcap_kU_k$~; comme $X$ est irréductible et les $U_k$ non vides, $U$ est
non vide, et par suite partout dense dans $X$~; les restrictions à $U$ des
$\varphi_k$ définissent un morphisme
\[
\varphi:U\to P=P_1\times_SP_2\times_S\cdots\times_SP_n
\]
de sorte que les diagrammes
\[
\tag{5.6.1.1}
\label{II.5.6.1.1-fr}
\xymatrix{
U\ar[r]^\varphi\ar[d]_{j_k} & P\ar[d]^{p_k}\\
U_k\ar[r]_{\varphi_k} & P_k
}
\]
soient commutatifs, $j_k$ étant l'injection canonique $U\to U_k$ et $p_k$ la
projection canonique $P\to P_k$. Si $j$ est l'injection canonique $U\to X$,
le morphisme $\psi=(j,\varphi)_S:U\to X\times_SP$ est une immersion
(\hyperref[I.5.3.14-fr]{I, 5.3.14}). Dans l'hypothèse a), $X\times_SP$ est
localement noethérien (\hyperref[II.3.4.1-fr]{3.4.1} et
\hyperref[I.6.3.7-fr]{I, 6.3.7 et 6.3.8})~; dans l'hypothèse b),
$X\times_SP$ est un schéma quasi-compact
(\hyperref[I.5.5.1-fr]{I, 5.5.1 et 6.6.4})~; dans les deux cas l'adhérence
$X'$ dans $X\times_SP$ du sous-préschéma $Z$ associé à $\psi$ (et dont
$\psi(U)$ est l'espace sous-jacent) existe, et $\psi$ se factorise en
\[
\tag{5.6.1.2}
\label{II.5.6.1.2-fr}
\psi:U\xrightarrow{\psi'}X'\xrightarrow{h}X\times_SP
\]
où $\psi'$ est une immersion ouverte et $h$ une immersion fermée
(\hyperref[I.9.5.10-fr]{I, 9.5.10}). Soient
$q_1:X\times_SP\to X$ et $q_2:X\times_SP\to P$ les projections canoniques~;
nous poserons
\[
\tag{5.6.1.3}
\label{II.5.6.1.3-fr}
f:X'\xrightarrow{h}X\times_SP\xrightarrow{q_1}X
\]
et
\[
\tag{5.6.1.4}
\label{II.5.6.1.4-fr}
g:X'\xrightarrow{h}X\times_SP\xrightarrow{q_2}P.
\]
Nous allons voir que $X'$ et $f$ répondent à la question.

C) Montrons d'abord que $f$ est projectif et surjectif, et que la restriction
de $f$ à $U'=f^{-1}(U)$ est un isomorphisme de $U'$ sur $U$. Comme les $P_k$
sont projectifs sur $S$, il en est de même de $P$
(\hyperref[II.5.5.5-fr]{5.5.5, (iv)}), donc $X\times_SP$ est projectif sur
$X$ (\hyperref[II.5.5.5-fr]{5.5.5, (iii)}), et il en est de même de $X'$ qui
est un sous-préschéma fermé de $X\times_SP$. D'autre part, on a
$f\circ\psi'=q_1\circ(h\circ\psi')=q_1\circ\psi=j$, donc $f(X')$ contient
l'ensemble ouvert partout dense $U$ de $X$~; mais $f$ est un morphisme fermé
(\hyperref[II.5.5.3-fr]{5.5.3}), donc $f(X')=X$. Notons maintenant que
$q_1^{-1}(U)=U\times_SP$ est induit sur un ouvert de $X\times_SP$, et par
définition le préschéma $U'=h^{-1}(U\times_SP)$ est induit par $X'$ sur
l'ouvert $U'$~; c'est par suite l'adhérence par

\oldpage[II]{108}

rapport à $U\times_SP$ du préschéma $Z$
(\hyperref[I.9.5.8-fr]{I, 9.5.8}). Or l'immersion $\psi$ se factorise en
\[
U\xrightarrow{\Gamma_\varphi}U\times_SP\xrightarrow{j\times1}X\times_SP,
\]
et comme $P$ est séparé sur $S$, le morphisme graphe $\Gamma_\varphi$ est une
immersion fermée (\hyperref[I.5.4.3-fr]{I, 5.4.3}), donc $Z$ est un
sous-préschéma fermé de $U\times_SP$, d'où on conclut que $U'=Z$. Comme
d'ailleurs $\psi$ est une immersion, la restriction de $f$ à $U'$ est un
isomorphisme sur $U$, réciproque de $\psi'$~; enfin, par définition de $X'$,
$U'$ est dense dans $X'$.

D) Prouvons maintenant que $g$ est une immersion, ce qui établira que $X'$ est
quasi-projectif sur $S$, puisque $P$ est projectif sur $S$. Posons
\begin{align*}
V_k&=\varphi_k(U_k) &&\text{(ouvert de $P_k$)},\\
W_k&=p_k^{-1}(V_k) &&\text{(ouvert de $P$)},\\
U'_k&=f^{-1}(U_k),\qquad U''_k=g^{-1}(W_k)
&&\text{(ouverts de $X'$)}.
\end{align*}
Il est clair que les $U'_k$ forment un recouvrement ouvert de $X'$~; nous
allons d'abord voir qu'il en est de même des $U''_k$, en prouvant que
$U'_k\subset U''_k$. Il suffira pour cela de montrer que le diagramme
\[
\tag{5.6.1.5}
\label{II.5.6.1.5-fr}
\xymatrix{
U'_k\ar[r]^{g|U'_k}\ar[d]_{f|U'_k} & P\ar[d]^{p_k}\\
U_k\ar[r]_{\varphi_k} & P_k
}
\]
est commutatif. Or le préschéma $U'_k=h^{-1}(U_k\times_SP)$ est induit par
$X'$ sur l'ouvert $U'_k$, et est donc l'adhérence de $Z=U'\subset U'_k$ par
rapport à $U'_k$ (\hyperref[I.9.5.8-fr]{I, 9.5.8}). Pour montrer la
commutativité du diagramme (\hyperref[II.5.6.1.5-fr]{5.6.1.5}), il suffit
donc, puisque $P_k$ est un $S$-schéma, de montrer qu'en composant ce diagramme
avec l'injection canonique $U'\to U'_k$ (ou, ce qui revient au même en vertu de
l'isomorphisme de $U'$ et $U$, avec $\psi$), on obtient un diagramme commutatif
(\hyperref[I.9.5.6-fr]{I, 9.5.6}). Mais par définition le diagramme que l'on
obtient alors n'est autre que (\hyperref[II.5.6.1.1-fr]{5.6.1.1}), d'où notre
assertion.

Les $W_k$ forment donc un recouvrement ouvert de $g(X')$~; pour prouver que
$g$ est une immersion, il suffit de voir que chacune des restrictions
$g|U''_k$ est une immersion dans $W_k$
(\hyperref[I.4.2.4-fr]{I, 4.2.4}). Pour cela, considérons le morphisme
$u_k:W_k\xrightarrow{p_k}V_k\xrightarrow{\varphi_k^{-1}}U_k\to X$~; comme
$X$ est séparé sur $S$, le morphisme graphe
$\Gamma_{u_k}:W_k\to X\times_SW_k$ est une immersion fermée
(\hyperref[I.5.4.3-fr]{I, 5.4.3}), donc le graphe
$T_k=\Gamma_{u_k}(W_k)$ est un sous-préschéma fermé de $X\times_SW_k$~; si
nous montrons que ce sous-préschéma majore $U'$, il majorera aussi le
sous-préschéma induit par $X'$ sur l'ouvert $X''_k$ de $X'$, par
(\hyperref[I.9.5.8-fr]{I, 9.5.8}). Comme la restriction de $q_2$ à $T_k$ est
un isomorphisme sur $W_k$, la restriction de $g$ à $X''_k$ sera une immersion
dans $W_k$, et notre assertion sera démontrée. Soit $v_k$ l'injection
canonique $U'\to X\times_SW_k$~; nous devons montrer qu'il existe un morphisme
$w_k:U'\to W_k$ tel que $v_k=\Gamma_{u_k}\circ w_k$. Par définition du
produit, il suffit de prouver que
$q_1\circ v_k=u_k\circ q_2\circ v_k$
(\hyperref[I.3.2.1-fr]{I, 3.2.1}), ou, en composant à droite

\oldpage[II]{109}

avec l'isomorphisme $\psi':U\to U'$, que
$q_1\circ\psi=u_k\circ q_2\circ\psi$. Mais comme $q_1\circ\psi=j$,
$q_2\circ\psi=\varphi$, notre assertion résulte de la commutativité du
diagramme (\hyperref[II.5.6.1.1-fr]{5.6.1.1}), compte tenu de la définition de
$u_k$.

E) Il est clair que comme $U$, et par suite $U'$, est irréductible, il en est
de même de $X'$ dans la construction précédente, et le morphisme $f$ est donc
birationnel (\hyperref[I.2.2.9-fr]{I, 2.2.9}). Si en outre $X$ est réduit, il
en est de même de $U'$, donc $X'$ est réduit
(\hyperref[I.9.5.9-fr]{I, 9.5.9}). Cela achève la démonstration.

\begin{corollary}[5.6.2]
\label{II.5.6.2-fr}
On suppose vérifiée l'une des hypothèses a), b) de
(\hyperref[II.5.6.1-fr]{5.6.1}). Pour que $X$ soit propre sur $S$, il faut et
il suffit qu'il existe un schéma $X'$ projectif sur $S$ et un $S$-morphisme
surjectif $f:X'\to X$ (qui est alors projectif par
(\hyperref[II.5.5.5-fr]{5.5.5, (v)})). Lorsqu'il en est ainsi, on peut en
outre choisir $f$ de sorte qu'il existe un ouvert dense $U$ de $X$ tel que la
restriction de $f$ à $f^{-1}(U)$ soit un isomorphisme
$f^{-1}(U)\simeq U$ et que $f^{-1}(U)$ soit dense dans $X'$. Si en outre $X$
est irréductible (resp. réduit) on peut supposer qu'il en est de même de
$X'$~; lorsque $X$ et $X'$ sont irréductibles, $f$ est un morphisme
birationnel.
\end{corollary}

La condition est suffisante en vertu de (\hyperref[II.5.5.3-fr]{5.5.3}) et
(\hyperref[II.5.4.3-fr]{5.4.3, (ii)}). Elle est nécessaire, car avec les
notations de (\hyperref[II.5.6.1-fr]{5.6.1}), si $X$ est propre sur $S$, $X'$
est propre sur $S$, car il est projectif sur $X$, donc propre sur $X$
(\hyperref[II.5.5.3-fr]{5.5.3}), et notre assertion résulte de
(\hyperref[II.5.4.2-fr]{5.4.2, (ii)})~; en outre, comme $X'$ est
quasi-projectif sur $S$, il est projectif sur $S$ en vertu de
(\hyperref[II.5.5.3-fr]{5.5.3}).

\begin{corollary}[5.6.3]
\label{II.5.6.3-fr}
Soit $S$ un préschéma localement noethérien, $X$ un $S$-schéma de type fini
sur $S$, $f_0:X\to S$ son morphisme structural. Pour que $X$ soit propre sur
$S$, il faut et il suffit que pour tout morphisme de type fini $S'\to S$,
$(f_0)_{(S')}:X_{(S')}\to S'$ soit un morphisme fermé. Il suffit même que
cette condition soit vérifiée pour tout $S$-préschéma de la forme
$S'=S\otimes_{\mathbf Z}\mathbf Z[T_1,\ldots,T_n]$ ($T_i$ indéterminées).
\end{corollary}

La condition étant évidemment nécessaire, montrons qu'elle est suffisante.
La question étant locale sur $S$ et $S'$ (\hyperref[II.5.4.1-fr]{5.4.1}), on
peut supposer $S$ et $S'$ affines et noethériens. En vertu du lemme de Chow,
il existe un $S$-schéma projectif $P$, une immersion $j:X'\to P$ et un
morphisme projectif et surjectif $f:X'\to X$, tels que le diagramme
\[
\xymatrix{
X\ar[d]_{f_0} & X'\ar[l]_f\ar[d]^j\\
S & P\ar[l]_r
}
\]
soit commutatif. Comme $P$ est de type fini sur $S$, la première hypothèse
entraîne que la projection $q_2:X\times_SP\to P$ est un morphisme fermé. Mais
l'immersion $j$ est composée de $q_2$ et du morphisme $f\times1$ de
$X'\times_SP$ dans $X\times_SP$~; or $f$, étant projectif, est propre
(\hyperref[II.5.5.3-fr]{5.5.3}), donc $f\times1$ est fermé. On en conclut que
$j$ est une immersion fermée, et par suite est propre
(\hyperref[II.5.4.2-fr]{5.4.2, (ii)}). En outre le morphisme structural
$r:P\to S$ est projectif, donc propre (\hyperref[II.5.5.3-fr]{5.5.3}), donc
$f_0\circ f=r\circ j$ est propre
(\hyperref[II.5.4.2-fr]{5.4.2, (ii)})~; finalement, comme $f$ est surjectif,
$f_0$ est propre en vertu de (\hyperref[II.5.4.3-fr]{5.4.3}).

Pour prouver la proposition en utilisant seulement la seconde hypothèse, il
suffit de montrer qu'elle entraîne la première. Or, si $S'$ est affine et de
type fini sur $S=\operatorname{Spec}(A)$,

\oldpage[II]{110}

on a $S'=\operatorname{Spec}(A[c_1,\ldots,c_n])$
(\hyperref[I.6.3.3-fr]{I, 6.3.3}), et $S'$ est donc isomorphe à un
sous-préschéma fermé de $S''=\operatorname{Spec}(A[T_1,\ldots,T_n])$ ($T_i$
indéterminées). Dans le diagramme commutatif
\[
\xymatrix{
X\times_SS'\ar[r]^{1_X\times j}\ar[d]_{(f_0)_{(S')}} &
X\times_SS''\ar[d]^{(f_0)_{(S'')}}\\
S'\ar[r]_j & S''
}
\]
$j$ et $1_X\times j$ sont des immersions fermées
(\hyperref[I.4.3.1-fr]{I, 4.3.1}) et $(f_0)_{(S'')}$ est fermé par hypothèse~;
il en est donc de même de $(f_0)_{(S')}$. 

\section{Morphismes entiers et morphismes finis}

\subsection{Préschémas entiers sur un autre.}

\begin{definition}[6.1.1]
\label{II.6.1.1-fr}
Soit $X$ un $S$-préschéma, $f:X\to S$ son morphisme structural. On dit que
$X$ est \emph{entier au-dessus de $S$}, ou que $f$ est un \emph{morphisme
entier}, s'il existe un recouvrement $(S_\alpha)$ de $S$ par des ouverts
affines tels que, pour tout $\alpha$, le préschéma induit $f^{-1}(S_\alpha)$
soit un schéma affine, dont l'anneau $B_\alpha$ est une algèbre entière sur
l'anneau $A_\alpha$ de $S_\alpha$. On dit que $X$ est \emph{fini au-dessus de
$S$}, ou que $f$ est un \emph{morphisme fini} si $X$ est entier sur $S$ et de
type fini sur $S$.
\end{definition}

Lorsque $S$ est affine d'anneau $A$, on dira aussi «~entier (resp. fini) sur
$A$~» au lieu de «~entier (resp. fini) sur $S$~».

\begin{env}[6.1.2]
\label{II.6.1.2-fr}
Il est clair que si $X$ est entier au-dessus de $S$, il est affine au-dessus
de $S$. Pour qu'un préschéma $X$ affine au-dessus de $S$ soit entier (resp.
fini) au-dessus de $S$, il faut et il suffit que la $\mathcal{O}_S$-Algèbre
quasi-cohérente associée $\mathcal{A}(X)$ soit telle qu'il existe un
recouvrement $(S_\alpha)$ de $S$ par des ouverts affines ayant la propriété
que pour tout $\alpha$, $\Gamma(S_\alpha,\mathcal{A}(X))$ est une algèbre
entière (resp. entière et de type fini) sur
$\Gamma(S_\alpha,\mathcal{O}_S)$. Une $\mathcal{O}_S$-Algèbre quasi-cohérente
ayant cette propriété est dite \emph{entière} (resp. \emph{finie}) sur
$\mathcal{O}_S$. Se donner un préschéma entier (resp. fini) au-dessus de $S$
revient donc (\hyperref[II.1.3.1-fr]{1.3.1}) à se donner une
$\mathcal{O}_S$-Algèbre quasi-cohérente entière (resp. finie) sur
$\mathcal{O}_S$. On notera qu'une $\mathcal{O}_S$-Algèbre quasi-cohérente
$\mathcal{B}$ est finie si et seulement si c'est un $\mathcal{O}_S$-Module de
type fini (\hyperref[I.1.3.9-fr]{I, 1.3.9})~; il revient au même de dire que
$\mathcal{B}$ est une $\mathcal{O}_S$-Algèbre entière et de type fini, car une
algèbre entière et de type fini sur un anneau $A$ est un $A$-module de type
fini.
\end{env}

\begin{proposition}[6.1.3]
\label{II.6.1.3-fr}
Soit $S$ un préschéma localement noethérien. Pour qu'un préschéma $X$ affine
au-dessus de $S$ soit fini au-dessus de $S$, il faut et il suffit que la
$\mathcal{O}_S$-Algèbre $\mathcal{A}(X)$ soit cohérente.
\end{proposition}

Compte tenu de la remarque précédente, cela revient à remarquer que si $S$ est
localement noethérien, les $\mathcal{O}_S$-Modules quasi-cohérents de type fini
ne sont autres que les $\mathcal{O}_S$-Modules cohérents
(\hyperref[I.1.5.1-fr]{I, 1.5.1}).

\begin{proposition}[6.1.4]
\label{II.6.1.4-fr}
Soit $X$ un préschéma entier (resp. fini) au-dessus de $S$, $f:X\to S$ le
morphisme structural. Alors, pour tout ouvert affine $U\subset S$, d'anneau
$A$, $f^{-1}(U)$ est un schéma affine dont l'anneau $B$ est une algèbre
entière (resp. finie) sur $A$.
\end{proposition}

\oldpage[II]{111}

Nous démontrerons d'abord le lemme suivant~:

\begin{lemma}[6.1.4.1]
\label{II.6.1.4.1-fr}
Soient $A$ un anneau, $M$ un $A$-module, $(g_i)_{1\leq i\leq m}$ un système
fini d'éléments de $A$ tels que les $D(g_i)$ ($1\leq i\leq m$) recouvrent
$\operatorname{Spec}(A)$. Si, pour tout $i$, $M_{g_i}$ est un
$A_{g_i}$-module de type fini, alors $M$ est un $A$-module de type fini.
\end{lemma}

On peut en effet supposer que $M_{g_i}$ admet un système de générateurs fini
$(m_{ij}/g_i^n)$ avec $m_{ij}\in M$, $n$ étant le même pour tous les indices
$i$. Montrons que les $m_{ij}$ forment un système de générateurs de $M$. Soit
$M'$ le sous-$A$-module de $M$ engendré par ces éléments, et soit $m$ un
élément de $M$. Par hypothèse, pour chaque $i$, il existe des $a_{ij}\in A$ et
un entier $p$ (indépendant de $i$) tels que, dans $M_{g_i}$,
$m/1=(\sum_j a_{ij}m_{ij})/g_i^p$~; cela entraîne qu'il existe un entier
$r\geq p$ tel que, pour tout $i$, on ait $g_i^r m\in M'$. Or, comme les
$D(g_i^r)=D(g_i)$ recouvrent $\operatorname{Spec}(A)$, l'idéal de $A$
engendré par les $g_i^r$ est égal à $A$, autrement dit il existe des éléments
$a_i\in A$ tels que $\sum_i a_i g_i^r=1$~; par suite
$m=(\sum_i a_i g_i^r)m\in M'$, d'où le lemme.

Cela étant, on sait déjà (\hyperref[II.1.3.2-fr]{1.3.2}) que $f^{-1}(U)$ est
affine. Si $\varphi$ est l'homomorphisme $A\to B$ correspondant à $f$, il
existe un recouvrement fini de $U$ par des ouverts $D(g_i)$ ($g_i\in A$) tels
que, si $h_i=\varphi(g_i)$, $B_{h_i}$ soit une algèbre entière (resp. entière
finie) sur $A_{g_i}$. En effet, il y a un recouvrement de $U$ par des ouverts
affines $V_\alpha\subset U$ tels que si
$A_\alpha=A(V_\alpha)$, $B_\alpha=A(f^{-1}(V_\alpha))$ soit une algèbre
entière (resp. finie) sur $A_\alpha$. Tout $x\in U$ appartient à un
$V_\alpha$, donc il existe $g\in A$ tel que $x\in D(g)\subset V_\alpha$~; si
$g_\alpha$ est l'image de $g$ dans $A_\alpha$, on a
$A(D(g))=A_g=(A_\alpha)_{g_\alpha}$~; soit $h=\varphi(g)$, et soit $h_\alpha$
l'image de $g_\alpha$ dans $B_\alpha$~; on a
\[
A(D(h))=B_h=(B_\alpha)_{h_\alpha}
\]
et comme $B_\alpha$ est entière (resp. finie) sur $A_\alpha$,
$(B_\alpha)_{h_\alpha}$ est entière (resp. finie) sur
$(A_\alpha)_{g_\alpha}$. Il suffit maintenant d'utiliser le fait que $U$ est
quasi-compact pour obtenir le recouvrement cherché.

Si on suppose d'abord que les $B_{h_i}$ sont entières finies sur les
$A_{g_i}$, comme, en tant que $A_{g_i}$-module, $B_{h_i}$ s'écrit aussi
$B_{g_i}$, le lemme (\hyperref[II.6.1.4.1-fr]{6.1.4.1}) montre que dans ce cas
$B$ est un $A$-module de type fini.

Supposons seulement maintenant chaque $B_{h_i}$ entière sur $A_{g_i}$~; soit
$b\in B$, et soit $C$ la sous-$A$-algèbre de $B$ engendrée par $b$. Pour tout
$i$, $C_{h_i}$ est l'algèbre sur $A_{g_i}$ engendrée par $b/1$ dans
$B_{h_i}$~; il résulte de l'hypothèse que chaque $C_{h_i}$ est un
$A_{g_i}$-module de type fini, donc
(\hyperref[II.6.1.4.1-fr]{6.1.4.1}) $C$ est un $A$-module de type fini, ce qui
prouve que $B$ est entière sur $A$.

\begin{proposition}[6.1.5]
\label{II.6.1.5-fr}
\begin{enumerate}
\item[(i)] Une immersion fermée est finie (et a fortiori entière).
\item[(ii)] Le composé de deux morphismes finis (resp. entiers) est fini
(resp. entier).
\item[(iii)] Si $f:X\to Y$ est un $S$-morphisme fini (resp. entier),
$f_{(S')}:X_{(S')}\to Y_{(S')}$ est fini (resp. entier) pour toute extension
$S'\to S$ de la base.
\item[(iv)] Si $f:X\to Y$ et $g:X'\to Y'$ sont deux $S$-morphismes finis
(resp. entiers), il en est de même de
$f\times_S g:X\times_S X'\to Y\times_S Y'$.
\item[(v)] Si $f:X\to Y$ et $g:Y\to Z$ sont deux morphismes tels que $g\circ
f$ soit fini (resp. entier) et $g$ séparé, alors $f$ est fini (resp. entier).
\item[(vi)] Si $f:X\to Y$ est un morphisme fini (resp. entier), il en est de
même de $f_{\mathrm{red}}$.
\end{enumerate}
\end{proposition}

\oldpage[II]{112}

En vertu de (\hyperref[I.5.5.12-fr]{I, 5.5.12}), il suffit de démontrer (i),
(ii) et (iii). Pour prouver qu'une immersion fermée $X\to S$ est finie, on
peut se borner au cas où $S=\operatorname{Spec}(A)$, et tout revient alors à
remarquer qu'un anneau quotient $A/\mathfrak{J}$ est un $A$-module monogène.
Pour prouver que le composé de deux morphismes finis (resp. entiers) $X\to Y$,
$Y\to Z$ est fini (resp. entier), on peut encore supposer $Z$ (et par suite
$X$, $Y$ (\hyperref[II.1.3.4-fr]{1.3.4})) affines, et alors l'assertion
équivaut à dire que si $B$ est une $A$-algèbre finie (resp. entière) et $C$
une $B$-algèbre finie (resp. entière), alors $C$ est une $A$-algèbre finie
(resp. entière), ce qui est immédiat. Enfin, pour démontrer (iii), on peut se
borner au cas où $S=Y$, puisque $X_{(S')}$ s'identifie à
$X\times_Y Y_{(S')}$ (\hyperref[I.3.3.11-fr]{I, 3.3.11})~; on peut en outre
supposer que $S=\operatorname{Spec}(A)$, $S'=\operatorname{Spec}(A')$~; alors
$X$ est affine d'anneau $B$ (\hyperref[II.1.3.4-fr]{1.3.4}), $X_{(S')}$
affine d'anneau $A'\otimes_A B$, et il suffit de remarquer que si $B$ est une
$A$-algèbre finie (resp. entière), $A'\otimes_A B$ est une $A'$-algèbre finie
(resp. entière).

Notons aussi que si $X$ et $Y$ sont deux $S$-préschémas qui sont finis (resp.
entiers) au-dessus de $S$, leur \emph{somme} $X\amalg Y$ est un préschéma fini
(resp. entier) au-dessus de $S$, car cela revient à voir que si $B$ et $C$
sont deux $A$-algèbres finies (resp. entières) sur $A$, il en est de même de
$B\times C$.

\begin{corollary}[6.1.6]
\label{II.6.1.6-fr}
Si $X$ est un préschéma entier (resp. fini) au-dessus de $S$, alors, pour tout
ouvert $U\subset S$, $f^{-1}(U)$ est entier (resp. fini) au-dessus de $U$.
\end{corollary}

C'est un cas particulier de (\hyperref[II.6.1.5-fr]{6.1.5, (iii)}).

\begin{corollary}[6.1.7]
\label{II.6.1.7-fr}
Soit $f:X\to Y$ un morphisme fini. Alors, pour tout $y\in Y$, la fibre
$f^{-1}(y)$ est un schéma algébrique fini sur $k(y)$, et a fortiori son
espace sous-jacent est discret et fini.
\end{corollary}

En effet, en tant que $k(y)$-préschéma, $f^{-1}(y)$ est identifié à
$X\times_Y\operatorname{Spec}(k(y))$ (\hyperref[I.3.6.1-fr]{I, 3.6.1}), qui
est fini au-dessus de $\operatorname{Spec}(k(y))$
(\hyperref[II.6.1.5-fr]{6.1.5, (iii)})~; c'est donc un schéma affine dont
l'anneau est une algèbre de rang fini sur $k(y)$
(\hyperref[II.6.1.4-fr]{6.1.4}). La proposition résulte alors de
(\hyperref[I.6.4.4-fr]{I, 6.4.4}).

\begin{corollary}[6.1.8]
\label{II.6.1.8-fr}
Soient $X$, $S$ deux préschémas intègres, $f:X\to S$ un morphisme dominant. Si
$f$ est entier (resp. fini), le corps $R(X)$ des fonctions rationnelles sur
$X$ est algébrique (resp. algébrique de degré fini) sur le corps $R(S)$ des
fonctions rationnelles sur $S$.
\end{corollary}

En effet, soit $s$ le point générique de $S$~; le $k(s)$-préschéma
$f^{-1}(s)$ est entier (resp. fini) au-dessus de
$\operatorname{Spec}(k(s))$ (\hyperref[II.6.1.5-fr]{6.1.5, (iii)}), et
contient par hypothèse le point générique $x$ de $X$~; l'anneau local de $x$
dans $f^{-1}(s)$ étant égal à $k(x)$ (\hyperref[I.3.6.5-fr]{I, 3.6.5}) est un
anneau local d'une algèbre entière (resp. finie) sur $k(s)$
(\hyperref[II.6.1.4-fr]{6.1.4}), d'où le corollaire.

\begin{remark}[6.1.9]
\label{II.6.1.9-fr}
L'hypothèse que $g$ est \emph{séparé} est essentielle pour la validité de
(\hyperref[II.6.1.5-fr]{6.1.5, (v)})~: en effet, si $Y$ n'est pas séparé sur
$Z$, l'identité $1_Y$ est le morphisme composé
$Y\xrightarrow{\Delta_Y}Y\times_ZY\xrightarrow{p_1}Y$, mais $\Delta_Y$ n'est
pas un morphisme entier, comme il résulte de
(\hyperref[II.6.1.10-fr]{6.1.10})~:
\end{remark}

\begin{proposition}[6.1.10]
\label{II.6.1.10-fr}
Tout morphisme entier est universellement fermé.
\end{proposition}

Soit $f:X\to Y$ un morphisme entier~; en vertu de
(\hyperref[II.6.1.5-fr]{6.1.5, (iii)}), il suffit de montrer que $f$ est
fermé. Soit $Z$ une partie fermée de $X$~; il existe un sous-préschéma de $X$
ayant

\oldpage[II]{113}

pour espace sous-jacent $Z$ (\hyperref[I.5.2.1-fr]{I, 5.2.1}), et il résulte
donc de (\hyperref[II.6.1.5-fr]{6.1.5, (i) et (ii)}) qu'on peut se borner à
prouver que $f(X)$ est fermé dans $Y$. En vertu de
(\hyperref[II.6.1.5-fr]{6.1.5, (vii)}), on peut supposer $X$ et $Y$ réduits~;
en outre, si $T$ est le sous-préschéma fermé réduit de $Y$ ayant pour espace
sous-jacent $\overline{f(X)}$ (\hyperref[I.5.2.1-fr]{I, 5.2.1}), on sait que
$f$ se factorise en $X\xrightarrow{g}T\xrightarrow{j}Y$, où $j$ est le
morphisme d'injection (\hyperref[I.5.2.2-fr]{I, 5.2.2}), et comme $j$ est
séparé (\hyperref[I.5.5.1-fr]{I, 5.5.1, (i)}), il résulte de
(\hyperref[II.6.1.5-fr]{6.1.5, (v)}) que $g$ est un morphisme entier. On peut
donc supposer que $f(X)$ est dense dans $Y$. Enfin, la question étant locale
sur $Y$, on peut se borner au cas où $Y=\operatorname{Spec}(A)$. Alors
$X=\operatorname{Spec}(B)$, où $B$ est une $A$-algèbre entière sur $A$
(\hyperref[II.6.1.4-fr]{6.1.4})~; en outre $A$ est réduit
(\hyperref[I.5.1.4-fr]{I, 5.1.4}), et l'hypothèse que $f(X)$ est dense dans
$Y$ entraîne que l'homomorphisme $\varphi:A\to B$ correspondant à $f$ est
injectif (\hyperref[I.1.2.7-fr]{I, 1.2.7}). Dire que dans ces conditions
$f(X)=Y$ signifie que tout idéal premier de $A$ est la trace sur $A$ d'un
idéal premier de $B$, ce qui n'est autre que le premier théorème de
Cohen-Seidenberg (\cite[t.~I, p.~257, th.~3]{I-13}).

\begin{corollary}[6.1.11]
\label{II.6.1.11-fr}
Tout morphisme fini $f:X\to Y$ est projectif.
\end{corollary}

En effet, comme $f$ est affine, $\mathcal{O}_X$ est un
$\mathcal{O}_X$-Module très ample relativement à $f$
(\hyperref[II.5.1.2-fr]{5.1.2})~; en outre $f_*(\mathcal{O}_X)$ est un
$\mathcal{O}_Y$-Module quasi-cohérent de type fini
(\hyperref[II.6.1.2-fr]{6.1.2})~; enfin $f$ est séparé, de type fini et
universellement fermé (\hyperref[II.6.1.10-fr]{6.1.10}), et on est donc dans
les conditions d'applications du critère
(\hyperref[II.5.5.4-fr]{5.5.4, (ii)}).

\begin{proposition}[6.1.12]
\label{II.6.1.12-fr}
Soient $f:X'\to X$ un morphisme fini, et soit
$\mathcal{B}=f_*(\mathcal{O}_{X'})$ (qui est une $\mathcal{O}_X$-Algèbre
quasi-cohérente et un $\mathcal{O}_X$-Module de type fini). Soit
$\mathcal{F}'$ un $\mathcal{O}_{X'}$-Module quasi-cohérent~; pour que
$\mathcal{F}'$ soit localement libre, de rang $r$, il faut et il suffit que
$f_*(\mathcal{F}')$ soit un $\mathcal{B}$-Module localement libre de rang
$r$.
\end{proposition}

Il est clair que si $f_*(\mathcal{F}')|U$ est isomorphe à
$\mathcal{B}^r|U$ ($U$ ouvert de $X$), $\mathcal{F}'|f^{-1}(U)$ est isomorphe
à $\mathcal{O}_{X'}^r|f^{-1}(U)$
(\hyperref[II.1.4.2-fr]{1.4.2}). Inversement, supposons que $\mathcal{F}'$
soit localement libre de rang $r$ et montrons que $f_*(\mathcal{F}')$ est
localement isomorphe à $\mathcal{B}^r$ en tant que $\mathcal{B}$-Module. Soit
$x$ un point de $X$~; lorsque $U$ parcourt un système fondamental de
voisinages affines de $x$, $f^{-1}(U)$ parcourt un système fondamental de
voisinages affines (\hyperref[II.1.2.5-fr]{1.2.5}) de l'ensemble fini
$f^{-1}(x)$, puisque $f$ est fermé
(\hyperref[II.6.1.10-fr]{6.1.10}). La proposition résulte alors du lemme
suivant~:

\begin{lemma}[6.1.12.1]
\label{II.6.1.12.1-fr}
Soient $Y$ un préschéma, $\mathcal{E}$ un $\mathcal{O}_Y$-Module localement
libre de rang $r$, $Z$ une partie finie de $Y$ contenue dans un ouvert affine
$V$. Il existe alors un voisinage $U\subset V$ de $Z$ tel que
$\mathcal{E}|U$ soit isomorphe à $\mathcal{O}_Y^r|U$.
\end{lemma}

On peut évidemment supposer $Y$ affine~; pour tout $z_i\in Z$, il existe dans
l'adhérence $\overline{\{z_i\}}$ au moins un point fermé $z'_i$
(\hyperref[0.2.1.3-fr]{0, 2.1.3})~; si $Z'$ est l'ensemble des $z'_i$, tout
voisinage de $Z'$ est un voisinage de $Z$, et on peut donc supposer $Z$ fermé
dans $Y$ et discret. Considérons le sous-préschéma fermé réduit de $Y$ ayant
$Z$ pour espace sous-jacent (\hyperref[I.5.2.1-fr]{I, 5.2.1}) et soit
$j:Z\to Y$ l'injection canonique~; $j^*(\mathcal{E})=\mathcal{E}\otimes_Y
\mathcal{O}_Z$ est localement libre de rang $r$ sur le schéma discret $Z$,
et par suite est isomorphe à $\mathcal{O}_Z^r$~; autrement dit il existe $r$
sections $s_i$ ($1\leq i\leq r$) de $\mathcal{E}\otimes_Y\mathcal{O}_Z$
au-dessus de $Z$ telles que l'homomorphisme
$\mathcal{O}_Z^r\to\mathcal{E}\otimes_Y\mathcal{O}_Z$ défini par ces sections
soit bijectif. Mais $Y=\operatorname{Spec}(A)$ est affine, $Z$ est défini par
un idéal $\mathfrak{J}$ de $A$ et on a $\mathcal{E}=\widetilde{M}$, où $M$
est un $A$-module~; les $s_i$ sont des éléments

\oldpage[II]{114}

de $M\otimes_A(A/\mathfrak{J})$ et sont donc images de $r$ éléments
$t_i\in M=\Gamma(Y,\mathcal{E})$. Pour tout $z_j\in Z$ il y a alors un
voisinage $V_j$ de $z_j$ tel que les restrictions des $t_i$ à $V_j$
définissent un isomorphisme
$\mathcal{O}_Y^r|V_j\to\mathcal{E}|V_j$
(\hyperref[0.5.5.4-fr]{0, 5.5.4})~; le voisinage $U$ réunion des $V_j$ répond
donc à la question.

\begin{proposition}[6.1.13]
\label{II.6.1.13-fr}
Soient $g:X'\to X$ un morphisme entier de préschémas, $Y$ un préschéma normal
localement intègre, $f$ une application rationnelle de $Y$ dans $X'$ telle
que $g\circ f$ soit une application rationnelle partout définie
(\hyperref[I.7.2.1-fr]{I, 7.2.1})~; alors $f$ est partout définie.
\end{proposition}

Si $f_1$ et $f_2$ sont deux morphismes (d'ouverts denses de $Y$ dans $X'$) de
la classe $f$, il est clair que $g\circ f_1$ et $g\circ f_2$ sont des
morphismes équivalents, ce qui justifie la notation $g\circ f$ pour leur
classe d'équivalence. On rappelle aussi que lorsqu'on suppose en outre que
$Y$ est localement noethérien, l'hypothèse que $Y$ est normal entraîne déjà
que $Y$ est localement intègre (\hyperref[I.6.1.13-fr]{I, 6.1.13}).

Pour démontrer (\hyperref[II.6.1.13-fr]{6.1.13}), notons d'abord que la
question est locale sur $Y$ et on peut par suite supposer qu'il existe dans
la classe de $g\circ f$ un morphisme $h:Y\to X$. Considérons l'image
réciproque $Y'=X'_{(h)}=X'_{(Y)}$, et notons que le morphisme
$g'=g_{(Y)}:Y'\to Y$ est entier
(\hyperref[II.6.1.5-fr]{6.1.5, (iii)}). Vu la correspondance entre
applications rationnelles de $Y$ dans $X'$ et $Y$-sections rationnelles de
$Y'$ (\hyperref[I.7.1.2-fr]{I, 7.1.2}), on voit qu'on est ramené à prouver le
cas particulier de (\hyperref[II.6.1.13-fr]{6.1.13}) où $X=Y$, autrement dit
le

\begin{corollary}[6.1.14]
\label{II.6.1.14-fr}
Soient $X$ un préschéma normal localement intègre, $g:X'\to X$ un morphisme
entier, $f$ une $X$-section rationnelle de $X'$. Alors $f$ est partout
définie.
\end{corollary}

La question étant locale sur $X$, on peut supposer $X$ intègre, et $f$
s'identifie alors à un morphisme d'un ouvert $U$ de $X$ dans $X'$
(\hyperref[I.7.2.2-fr]{I, 7.2.2}) qui est une $U$-section de $g^{-1}(U)$.
Comme $g$ est séparé, $f$ est une immersion fermée de $U$ dans $g^{-1}(U)$
(\hyperref[I.5.4.6-fr]{I, 5.4.6})~; soit $Z$ le sous-préschéma fermé de
$g^{-1}(U)$ associé à $f$ (\hyperref[I.4.2.1-fr]{I, 4.2.1}), qui est
isomorphe à $U$, donc intègre~; soit $X_1$ le sous-préschéma réduit de $X'$
ayant pour espace sous-jacent l'adhérence $\overline{Z}$ de $Z$ dans $X'$
(\hyperref[I.5.2.1-fr]{I, 5.2.1})~; $Z$ est alors un sous-préschéma induit
sur un ouvert de $X_1$ (\hyperref[I.5.2.3-fr]{I, 5.2.3}) et comme il est
irréductible, il en est de même de $X_1$, qui est par suite intègre. Le
morphisme $f$ peut alors être considéré comme une $X$-section rationnelle de
$X_1$~; comme la restriction de $g$ à $X_1$ est un morphisme entier
(\hyperref[II.6.1.5-fr]{6.1.5, (i) et (ii)}), on est finalement ramené à
prouver (\hyperref[II.6.1.14-fr]{6.1.14}) dans le cas particulier où
$X'=X_1$, autrement dit, le

\begin{corollary}[6.1.15]
\label{II.6.1.15-fr}
Soient $X$ un préschéma intègre normal, $X'$ un préschéma intègre,
$g:X'\to X$ un morphisme entier. S'il existe une $X$-section rationnelle $f$
de $X'$, $g$ est un isomorphisme.
\end{corollary}

La question étant locale sur $X$, on peut supposer $X$ affine d'anneau
intègre $A$, et alors $X'$ est affine d'anneau $A'$ entier sur $A$
(\hyperref[II.6.1.4-fr]{6.1.4}) et intègre~; en outre, le raisonnement de
(\hyperref[II.6.1.14-fr]{6.1.14}) montre qu'il y a un ouvert dense dans $X$
isomorphe à un ouvert dense de $X'$, donc $A$ et $A'$ ont même corps de
fractions. D'autre part, en vertu de
(\hyperref[I.8.2.1.1-fr]{I, 8.2.1.1}) et de l'hypothèse que les
$\mathcal{O}_x$ sont intégralement clos, l'anneau $A$ est intégralement clos,
donc $A'=A$, ce qui achève la démonstration de
(\hyperref[II.6.1.13-fr]{6.1.13}).

\subsection{Morphismes quasi-finis.}

\begin{proposition}[6.2.1]
\label{II.6.2.1-fr}
Soient $f:X\to Y$ un morphisme localement de type fini, $x$ un point de $X$.
Les conditions suivantes sont équivalentes~:

\oldpage[II]{115}

\begin{enumerate}
\item[a)] Le point $x$ est isolé dans sa fibre $f^{-1}(f(x))$.
\item[b)] L'anneau $\mathcal{O}_x$ est un $\mathcal{O}_{f(x)}$-module
quasi-fini (\hyperref[0.7.4.1-fr]{0, 7.4.1}).
\end{enumerate}
\end{proposition}

La question étant évidemment locale sur $X$ et sur $Y$, on peut supposer
$X=\operatorname{Spec}(A)$ et $Y=\operatorname{Spec}(B)$ affines, $A$ étant
une $B$-algèbre de type fini (\hyperref[I.6.3.3-fr]{I, 6.3.3}). En outre, on
peut remplacer $X$ par $X\times_Y\operatorname{Spec}(\mathcal{O}_{f(x)})$
sans changer la fibre $f^{-1}(f(x))$ ni l'anneau local $\mathcal{O}_x$
(\hyperref[I.3.6.5-fr]{I, 3.6.5})~; on peut donc supposer que $B$ est un
anneau local (égal à $\mathcal{O}_{f(x)}$)~; si $\mathfrak{n}$ est l'idéal
maximal de $B$, $f^{-1}(f(x))$ est alors un schéma affine d'anneau
$A/\mathfrak{n}A$, de type fini sur $k(f(x))=B/\mathfrak{n}$
(\hyperref[I.6.4.11-fr]{I, 6.4.11}). Cela étant, si a) est vérifiée, on peut
en outre supposer que $f^{-1}(f(x))$ est réduite au point $x$~; donc
$A/\mathfrak{n}A$ est de rang fini sur $B/\mathfrak{n}$
(\hyperref[I.6.4.4-fr]{I, 6.4.4}), autrement dit $A$ est un $B$-module
quasi-fini. Inversement, s'il en est ainsi, $f^{-1}(f(x))$ est un schéma
affine artinien, donc discret (\hyperref[I.6.4.4-fr]{I, 6.4.4})~; par suite
$x$ est isolé dans sa fibre, et cela montre que b) entraîne a).

\begin{corollary}[6.2.2]
\label{II.6.2.2-fr}
Soit $f:X\to Y$ un morphisme de type fini. Les conditions suivantes sont
équivalentes~:
\begin{enumerate}
\item[a)] Tout point $x\in X$ est isolé dans sa fibre $f^{-1}(f(x))$
(autrement dit le sous-espace $f^{-1}(f(x))$ est discret).
\item[b)] Pour tout $x\in X$, le préschéma $f^{-1}(f(x))$ est un
$k(f(x))$-préschéma fini.
\item[c)] Pour tout $x\in X$, l'anneau $\mathcal{O}_x$ est un
$\mathcal{O}_{f(x)}$-module quasi-fini.
\end{enumerate}
\end{corollary}

L'équivalence de a) et c) résulte de
(\hyperref[II.6.2.1-fr]{6.2.1}). D'autre part, comme $f^{-1}(f(x))$ est un
$k(f(x))$-préschéma algébrique (\hyperref[I.6.4.11-fr]{I, 6.4.11}),
l'équivalence de a) et b) résulte de (\hyperref[I.6.4.4-fr]{I, 6.4.4}).

\begin{definition}[6.2.3]
\label{II.6.2.3-fr}
Lorsque $f:X\to Y$ est un morphisme de type fini vérifiant les conditions
équivalentes de (\hyperref[II.6.2.2-fr]{6.2.2}), on dit que $f$ est
\emph{quasi-fini}, ou que $X$ est quasi-fini sur $Y$.
\end{definition}

Il est clair que tout morphisme fini est quasi-fini
(\hyperref[II.6.1.8-fr]{6.1.8}).

\begin{proposition}[6.2.4]
\label{II.6.2.4-fr}
\begin{enumerate}
\item[(i)] Une immersion $X\to Y$, qui est fermée, ou telle que $X$ soit
noethérien, est un morphisme quasi-fini.
\item[(ii)] Si $f:X\to Y$ et $g:Y\to Z$ sont des morphismes quasi-finis,
$g\circ f$ est quasi-fini.
\item[(iii)] Si $X$ et $Y$ sont des $S$-préschémas, $f:X\to Y$ un
$S$-morphisme quasi-fini, alors $f_{(S')}:X_{(S')}\to Y_{(S')}$ est quasi-fini
pour toute extension $g:S'\to S$ de la base.
\item[(iv)] Si $f:X\to Y$ et $g:X'\to Y'$ sont deux $S$-morphismes
quasi-finis,
\[
f\times_S g:X\times_S Y\to X'\times_S Y'
\]
est quasi-fini.
\item[(v)] Soient $f:X\to Y$ et $g:Y\to Z$ deux morphismes tels que $g\circ
f$ soit quasi-fini~; alors, si en outre $g$ est séparé, ou $X$ noethérien, ou
$X\times_ZY$ localement noethérien, $f$ est quasi-fini.
\item[(vi)] Si $f$ est quasi-fini, il en est de même de $f_{\mathrm{red}}$.
\end{enumerate}
\end{proposition}

Si $f:X\to Y$ est une immersion, toute fibre est réduite à un point, et
l'assertion (i) résulte donc de (\hyperref[I.6.3.4-fr]{I, 6.3.4 (i) et
6.3.5}). Pour prouver (ii), on remarque d'abord que $h=g\circ f$ est de type
fini (\hyperref[I.6.3.4-fr]{I, 6.3.4 (ii)})~; en outre, si $z=h(x)$ et
$y=f(x)$, $y$ est isolé dans $g^{-1}(z)$, donc il existe un voisinage ouvert
$V$ de $y$ dans $Y$ ne rencontrant aucun point de $g^{-1}(z)$ distinct de
$y$~; $f^{-1}(V)$ est donc un voisinage ouvert de $x$ ne rencontrant aucun

\oldpage[II]{116}

des $f^{-1}(y')$, où $y'\neq y$ est dans $g^{-1}(z)$~; comme $x$ est isolé
dans $f^{-1}(y)$, il est isolé dans
$h^{-1}(z)=f^{-1}(g^{-1}(z))$. Pour prouver (iii), on peut se limiter au cas
où $Y=S$ (\hyperref[I.3.3.11-fr]{I, 3.3.11})~; on remarque encore que tout
d'abord $f'=f_{(S')}$ est de type fini
(\hyperref[I.6.3.4-fr]{I, 6.3.4, (iii)})~; d'autre part, si
$x'\in X'=X_{(S')}$ et si on pose $y'=f'(x')$ et $y=g(y')$,
$f'^{-1}(y')$ s'identifie à
$f^{-1}(y)\otimes_{k(y)}k(y')$ (\hyperref[I.3.6.5-fr]{I, 3.6.5})~; comme
$f^{-1}(y)$ est de rang fini sur $k(y)$ par hypothèse, $f'^{-1}(y')$ est de
rang fini sur $k(y')$, donc discret. Les assertions (iv), (v), (vi) se
déduisent de (i), (ii), (iii) par la méthode générale
(\hyperref[I.5.5.12-fr]{I, 5.5.12}), sauf lorsque les hypothèses dans (v) sont
autres que l'hypothèse «~$g$ séparé~»~; pour traiter ce cas, on remarque
d'abord que si $x$ est isolé dans
$f^{-1}(g^{-1}(g(f(x))))$, il l'est a fortiori dans $f^{-1}(f(x))$~; le fait
que $f$ est de type fini résulte alors de
(\hyperref[I.6.3.6-fr]{I, 6.3.6}).

\begin{proposition}[6.2.5]
\label{II.6.2.5-fr}
Soient $A$ un anneau local noethérien complet, $X$ un $A$-schéma localement
de type fini, $x$ un point de $X$ au-dessus du point fermé $y$ de
$Y=\operatorname{Spec}(A)$, et supposons $x$ isolé dans sa fibre
$f^{-1}(y)$ ($f$ étant le morphisme structural $X\to Y$). Alors
$\mathcal{O}_x$ est un $A$-module de type fini et $X$ est $Y$-isomorphe à la
somme (I, 3.1) de
$X'=\operatorname{Spec}(\mathcal{O}_x)$ (qui est un $Y$-schéma fini) et d'un
$A$-schéma $X''$.
\end{proposition}

Il résulte de (\hyperref[II.6.2.1-fr]{6.2.1}) que $\mathcal{O}_x$ est un
$A$-module quasi-fini. Comme $\mathcal{O}_x$ est noethérien
(\hyperref[I.6.3.7-fr]{I, 6.3.7}) et l'homomorphisme
$A\to\mathcal{O}_x$ local, l'hypothèse que $A$ est complet entraîne que
$\mathcal{O}_x$ est un $A$-module de type fini
(\hyperref[0.7.4.3-fr]{0, 7.4.3}). Soit
$X'=\operatorname{Spec}(\mathcal{O}_x)$ le schéma local de $X$ au point $x$
(\hyperref[I.2.4.1-fr]{I, 2.4.1}), et soit $g:X'\to X$ le morphisme
canonique. Comme le composé $X'\xrightarrow{g}X\xrightarrow{f}Y$ est fini
(\hyperref[II.6.1.1-fr]{6.1.1}) et que $f$ est séparé, $g$ est fini
(\hyperref[II.6.1.5-fr]{6.1.5, (v)}), donc $g(X')$ est fermé dans $X$
(\hyperref[II.6.1.10-fr]{6.1.10})~; d'autre part, comme $g$ est de type fini,
$g$ est un isomorphisme local au point fermé $x'$ de $X'$ en vertu de la
définition de $g$ et de (\hyperref[I.6.5.4-fr]{I, 6.5.4})~; mais comme $X'$
est le seul voisinage ouvert de $x'$, cela signifie que $g$ est une immersion
ouverte, donc $g(X')$ est aussi ouvert dans $X$, ce qui achève la
démonstration.

\begin{corollary}[6.2.6]
\label{II.6.2.6-fr}
Soient $A$ un anneau local noethérien complet, $Y=\operatorname{Spec}(A)$,
$f:X\to Y$ un morphisme séparé et quasi-fini. Alors $X$ est $Y$-isomorphe à
une somme $X'\amalg X''$, où $X'$ est un $Y$-schéma fini et $X''$ un
$Y$-schéma quasi-fini tel que, si $y$ est le point fermé de $Y$,
$X''\cap f^{-1}(y)=\varnothing$.
\end{corollary}

En effet, la fibre $f^{-1}(y)$ est finie et discrète par hypothèse, et le
corollaire résulte donc, par récurrence sur le nombre de points de cette
fibre, de (\hyperref[II.6.2.5-fr]{6.2.5}).

\begin{remark}[6.2.7]
\label{II.6.2.7-fr}
Au chapitre V, nous verrons que si $Y$ est localement noethérien, un morphisme
quasi-fini séparé $X\to Y$ est nécessairement quasi-affine.
\end{remark}

\subsection{Fermeture intégrale d'un préschéma.}

\begin{proposition}[6.3.1]
\label{II.6.3.1-fr}
Soient $(X,\mathcal{A})$ un espace annelé, $\mathcal{B}$ une
$\mathcal{A}$-Algèbre (commutative), $f$ une section de $\mathcal{B}$
au-dessus de $X$. Les propriétés suivantes sont équivalentes~:
\begin{enumerate}
\item[a)] Le sous-$\mathcal{A}$-Module de $\mathcal{B}$ engendré par les
$f^n$ pour $n\geq 0$ (\hyperref[0.5.1.1-fr]{0, 5.1.1}) est de type fini.
\item[b)] Il existe une sous-$\mathcal{A}$-Algèbre $\mathcal{C}$ de
$\mathcal{B}$, qui est un $\mathcal{A}$-Module de type fini, telle que
$f\in\Gamma(X,\mathcal{C})$.
\item[c)] Pour tout $x\in X$, $f_x$ est entier sur la fibre
$\mathcal{A}_x$.
\end{enumerate}
\end{proposition}

\oldpage[II]{117}

Comme le sous-$\mathcal{A}$-Module de $\mathcal{B}$ engendré par les $f^n$ est
une $\mathcal{A}$-Algèbre, il est clair que a) entraîne b). D'autre part, b)
implique que pour tout $x\in X$, le $\mathcal{A}_x$-module $\mathcal{C}_x$ est
de type fini, ce qui entraîne que tout élément de l'algèbre $\mathcal{C}_x$ et
en particulier $f_x$, est entier sur $\mathcal{A}_x$. Enfin, si pour un
$x\in X$, on a une relation de la forme
\[
f_x^n+(a_1)_xf_x^{n-1}+\ldots+(a_n)_x=0
\]
les $a_i$ ($1\leq i\leq n$) étant des sections de $\mathcal{A}$ au-dessus
d'un voisinage ouvert $U$ de $x$, la section
$f^n|U+a_1\cdot f^{n-1}|U+\ldots+a_n$ est nulle au-dessus d'un voisinage
$V\subset U$ de $x$, d'où résulte aussitôt que tous les $f^k|V$ ($k\geq0$)
sont combinaisons linéaires à coefficients dans $\Gamma(V,\mathcal{A})$ des
$f^j|V$ pour $0\leq j\leq n-1$~; on en conclut que c) entraîne a).

Lorsque les conditions équivalentes de (\hyperref[II.6.3.1-fr]{6.3.1}) sont
remplies on dit que la section $f$ est \emph{entière sur $\mathcal{A}$}.

\begin{corollary}[6.3.2]
\label{II.6.3.2-fr}
Sous les hypothèses de (\hyperref[II.6.3.1-fr]{6.3.1}), il existe un
sous-$\mathcal{A}$-Module (unique) $\mathcal{A}'$ de $\mathcal{B}$ tel que pour
tout $x\in X$, $\mathcal{A}'_x$ soit l'ensemble des germes
$f_x\in\mathcal{B}_x$ qui sont entiers sur $\mathcal{A}_x$. Pour tout ouvert
$U\subset X$, les sections de $\mathcal{A}'$ au-dessus de $U$ sont les
sections $f\in\Gamma(U,\mathcal{B})$ qui sont entières sur
$\mathcal{A}|U$.
\end{corollary}

L'existence de $\mathcal{A}'$ est immédiate, en prenant pour
$\Gamma(U,\mathcal{A}')$ l'ensemble des $f\in\Gamma(U,\mathcal{B})$ tels que
$f_x$ soit entier sur $\mathcal{A}_x$ pour tout $x\in U$. La seconde
assertion résulte aussitôt de (\hyperref[II.6.3.1-fr]{6.3.1}).

Il est clair que $\mathcal{A}'$ est une sous-$\mathcal{A}$-Algèbre de
$\mathcal{B}$~; on dit que c'est la \emph{fermeture intégrale de
$\mathcal{A}$ dans $\mathcal{B}$}.

\begin{env}[6.3.3]
\label{II.6.3.3-fr}
Soient $(X,\mathcal{A})$, $(Y,\mathcal{B})$ deux espaces annelés,
\[
g=(\psi,\theta):X\to Y
\]
un morphisme. Soit $\mathcal{C}$ (resp. $\mathcal{D}$) une
$\mathcal{A}$-Algèbre (resp. une $\mathcal{B}$-Algèbre) et soit
\[
u:\mathcal{D}\to\mathcal{C}
\]
un $g$-morphisme (\hyperref[0.4.4.1-fr]{0, 4.4.1}). Alors, si
$\mathcal{A}'$ (resp. $\mathcal{B}'$) est la fermeture intégrale de
$\mathcal{A}$ (resp. $\mathcal{B}$) dans $\mathcal{C}$ (resp.
$\mathcal{D}$), la \emph{restriction} de $u$ à $\mathcal{B}'$ est un
$g$-morphisme
\[
u':\mathcal{B}'\to\mathcal{A}'.
\]

En effet, si $j$ est l'injection canonique $\mathcal{B}'\to\mathcal{D}$, il
suffit de montrer que
\[
v=u^\sharp\circ g^*(j):g^*(\mathcal{B}')\to\mathcal{C}
\]
applique $g^*(\mathcal{B}')$ dans $\mathcal{A}'$. Or, un élément de
$(g^*(\mathcal{B}'))_x=\mathcal{B}'_{\psi(x)}\otimes_{\mathcal{B}_{\psi(x)}}
\mathcal{A}_x$ est entier sur $\mathcal{A}_x$ par définition de
$\mathcal{B}'$, donc il en est de même de son image par $v_x$, ce qui établit
notre assertion.
\end{env}

\begin{proposition}[6.3.4]
\label{II.6.3.4-fr}
Soient $X$ un préschéma, $\mathcal{A}$ une $\mathcal{O}_X$-Algèbre
quasi-cohérente. La fermeture intégrale $\mathcal{O}'_X$ de
$\mathcal{O}_X$ dans $\mathcal{A}$ est alors une $\mathcal{O}_X$-Algèbre
quasi-cohérente, et pour tout ouvert affine $U$ de $X$,
$\Gamma(U,\mathcal{O}'_X)$ est la fermeture intégrale de
$\Gamma(U,\mathcal{O}_X)$ dans $\Gamma(U,\mathcal{A})$.
\end{proposition}

On peut se borner au cas où $X=\operatorname{Spec}(B)$ est affine et
$\mathcal{A}=\widetilde{A}$, $A$ étant une

\oldpage[II]{118}

$B$-algèbre~; soit $B'$ la fermeture intégrale de $B$ dans $A$. Tout revient
à voir que pour tout $x\in X$, un élément de $A_x$, entier sur $B_x$,
appartient nécessairement à $B'_x$, ce qui résulte du fait que, pour un anneau
commutatif $C$, les opérations de fermeture intégrale dans une $C$-algèbre et
de passage à un anneau de fractions (pour une partie multiplicative de $C$)
commutent (\cite[t.~I, p.~261 et 257]{I-13}).

Le $X$-schéma $X'=\operatorname{Spec}(\mathcal{O}'_X)$ est alors appelé la
\emph{fermeture intégrale de $X$ relativement à $\mathcal{A}$} (ou
\emph{relativement à $\operatorname{Spec}(\mathcal{A})$})~; il est clair que
$X'$ est \emph{entier} sur $X$ (\hyperref[II.6.1.2-fr]{6.1.2}).

On déduit aussitôt de (\hyperref[II.6.3.4-fr]{6.3.4}) que si $f:X'\to X$ est
le morphisme structural, alors, pour tout ouvert $U$ de $X$, $f^{-1}(U)$ est
la \emph{fermeture intégrale du préschéma induit par $X$ sur $U$, relativement
à $\mathcal{A}|U$}.

\begin{env}[6.3.5]
\label{II.6.3.5-fr}
Soient $X$, $Y$ deux préschémas, $f:X\to Y$ un morphisme, $\mathcal{A}$
(resp. $\mathcal{B}$) une $\mathcal{O}_X$-Algèbre (resp. une
$\mathcal{O}_Y$-Algèbre) quasi-cohérente, et soit
$u:\mathcal{B}\to\mathcal{A}$ un $f$-morphisme. On a vu
(\hyperref[II.6.3.3-fr]{6.3.3}) qu'on en déduit un $f$-morphisme
$u':\mathcal{O}'_Y\to\mathcal{O}'_X$, où $\mathcal{O}'_X$ (resp.
$\mathcal{O}'_Y$) est la fermeture intégrale de $\mathcal{O}_X$ (resp.
$\mathcal{O}_Y$) dans $\mathcal{A}$ (resp. $\mathcal{B}$). Par suite, si
$X'$ (resp. $Y'$) est la fermeture intégrale de $X$ (resp. $Y$) relativement
à $\mathcal{A}$ (resp. $\mathcal{B}$), on déduit canoniquement de $u$ un
morphisme $f'=\operatorname{Spec}(u'):X'\to Y'$
(\hyperref[II.1.5.6-fr]{1.5.6}) rendant commutatif le diagramme
\[
\label{II.6.3.5.1-fr}
\xymatrix{
X'\ar[r]^{f'}\ar[d]&Y'\ar[d]\\
X\ar[r]_{f}&Y
}
\tag{6.3.5.1}
\]
\end{env}

\begin{env}[6.3.6]
\label{II.6.3.6-fr}
Supposons que $X$ n'ait qu'un nombre \emph{fini} de composantes irréductibles
$X_i$ ($1\leq i\leq r$), de points génériques $\xi_i$, et considérons en
particulier la fermeture intégrale de $X$ relativement à une
$\mathcal{R}(X)$-Algèbre $\mathcal{A}$ quasi-cohérente (en tant que
$\mathcal{O}_X$-Algèbre ou en tant que $\mathcal{R}(X)$-Algèbre, ce qui
revient au même). On sait (\hyperref[I.7.3.5-fr]{I, 7.3.5}) que
$\mathcal{A}$ est composée directe de $r$ $\mathcal{O}_X$-Algèbres
quasi-cohérentes $\mathcal{A}_i$, le support de $\mathcal{A}_i$ étant contenu
dans $X_i$, et le faisceau induit par $\mathcal{A}_i$ sur $X_i$ étant un
faisceau simple dont la fibre $A_i$ est une algèbre sur
$\mathcal{O}_{\xi_i}$. Il est clair alors (\hyperref[II.6.3.4-fr]{6.3.4}) que
la fermeture intégrale $\mathcal{O}'_X$ de $\mathcal{O}_X$ dans
$\mathcal{A}$ est \emph{composée directe} des fermetures intégrales
$\mathcal{O}^{(i)}_X$ de $\mathcal{O}_X$ dans chacune des $\mathcal{A}_i$, et
par suite la fermeture intégrale $X'=\operatorname{Spec}(\mathcal{O}'_X)$ de
$X$ relativement à $\mathcal{A}$ est un $X$-schéma \emph{somme} des
$\operatorname{Spec}(\mathcal{O}^{(i)}_X)=X'_i$ ($1\leq i\leq r$).

Supposons en outre que la $\mathcal{O}_X$-Algèbre $\mathcal{A}$ soit
\emph{réduite}, ou, ce qui revient au même, que chacune des algèbres $A_i$
soit réduite, et puisse par suite être considérée comme une algèbre sur le
\emph{corps} $k(\xi_i)$ (égal au corps des fonctions rationnelles du
sous-préschéma réduit de $X$ ayant $X_i$ pour espace sous-jacent)~; alors
(\hyperref[II.1.3.8-fr]{1.3.8}) chacun des $X'_i$ est un $X$-schéma
\emph{réduit} et $X'$ est aussi la fermeture intégrale de $X_{\mathrm{red}}$.
Supposons de plus que chacune des algèbres $A_i$ soit \emph{composée directe
d'un nombre fini de corps} $K_{ij}$ ($1\leq j\leq s_i$)~; si
$\mathcal{K}_{ij}$ est la sous-Algèbre de $\mathcal{A}_i$ correspondant à
$K_{ij}$, il est clair que $\mathcal{O}^{(i)}_X$ est \emph{composée directe}
des fermetures intégrales $\mathcal{O}^{(ij)}_X$ de $\mathcal{O}_X$ dans
chacune des $\mathcal{K}_{ij}$. Par suite, $X'_i$ est alors \emph{somme} des
$X$-schémas $X'_{ij}=\operatorname{Spec}(\mathcal{O}^{(ij)}_X)$
($1\leq j\leq s_i$). En outre, sous ces hypothèses et avec ces notations~:
\end{env}

\oldpage[II]{119}

\begin{proposition}[6.3.7]
\label{II.6.3.7-fr}
Chacun des $X'_{ij}$ est un $X$-préschéma intègre et normal, et son corps des
fonctions rationnelles $R(X'_{ij})$ s'identifie canoniquement à la fermeture
algébrique $K'_{ij}$ de $k(\xi_i)$ dans $K_{ij}$.
\end{proposition}

En vertu de ce qui précède, on peut supposer $X$ intègre, donc $r=1$, et
$s_1=1$, de sorte que l'unique algèbre $A_1$ est un corps $K$~; soit $\xi$ le
point générique de $X$, et soit $f:X'\to X$ le morphisme structural. Pour tout
ouvert affine non vide $U$ de $X$, $f^{-1}(U)$ s'identifie au spectre de la
fermeture intégrale $B'_U$ dans le corps $K$ de l'anneau intègre
$B_U=\Gamma(U,\mathcal{O}_X)$ (\hyperref[II.6.3.4-fr]{6.3.4})~; comme l'anneau
$B'_U$ est intègre et intégralement clos, il en est de même des anneaux locaux
des points de son spectre, donc $f^{-1}(U)$ est par définition un schéma
\emph{intègre et normal} ((\hyperref[0.4.1.4-fr]{0, 4.1.4}) et
(\hyperref[I.5.1.4-fr]{I, 5.1.4})). En outre, comme $(0)$ est le seul idéal
premier de $B'_U$ au-dessus de l'idéal premier $(0)$ de $B_U$
(\cite[t.~I, p.~259]{I-13}), $f^{-1}(\xi)$ se réduit à un seul point $\xi'$,
et $k(\xi')$ est le corps des fractions $K'$ de $B'_U$, qui n'est autre que la
fermeture algébrique de $k(\xi)$ dans $K$. Enfin, $X'$ est irréductible, car
lorsque $U$ parcourt l'ensemble des ouverts affines non vides de $X$, les
$f^{-1}(U)$ constituent un recouvrement ouvert de $X'$ formé d'ouverts
irréductibles~; en outre l'intersection $f^{-1}(U\cap V)$ de deux de ces
ouverts contient $\xi'$, donc n'est pas vide, et on conclut par
(\hyperref[0.2.1.4-fr]{0, 2.1.4}).

\begin{corollary}[6.3.8]
\label{II.6.3.8-fr}
Soit $X$ un préschéma réduit n'ayant qu'un nombre fini de composantes
irréductibles $X_i$ ($1\leq i\leq r$), et soit $\xi_i$ le point générique de
$X_i$. La fermeture intégrale $X'$ de $X$ relativement à $\mathcal{R}(X)$ est
somme de $r$ $X$-schémas $X'_i$ qui sont intègres et normaux. Si $f:X'\to X$
est le morphisme structural, $f^{-1}(\xi_i)$ se réduit au point générique
$\xi'_i$ de $X'_i$ et on a $k(\xi'_i)=k(\xi_i)$, autrement dit $f$ est
birationnel.
\end{corollary}

On dit dans ce cas que $X'$ est le \emph{normalisé} du préschéma réduit $X$~;
on notera que $f$, étant birationnel et entier, est \emph{surjectif}
(\hyperref[II.6.1.10-fr]{6.1.10}). Pour que l'on ait alors $X'=X$, il faut et
il suffit que $X$ soit normal. Lorsque $X$ est un préschéma intègre, il résulte
de (\hyperref[II.6.3.8-fr]{6.3.8}) que son normalisé $X'$ est intègre.

\begin{env}[6.3.9]
\label{II.6.3.9-fr}
Soient $X$, $Y$ deux préschémas intègres, $f:X\to Y$ un morphisme dominant,
$L=R(X)$, $K=R(Y)$ les corps des fonctions rationnelles de $X$ et $Y$~; il
correspond canoniquement à $f$ une injection $K\to L$, et lorsqu'on identifie
$K$ (resp. $L$) au faisceau simple $\mathcal{R}(Y)$ (resp.
$\mathcal{R}(X)$), cette injection est un $f$-morphisme. Soit $K_1$ (resp.
$L_1$) une extension de $K$ (resp. $L$) et supposons donné un monomorphisme
$K_1\to L_1$ tel que le diagramme
\[
\xymatrix{
K_1\ar[r]&L_1\\
K\ar[u]\ar[r]&L\ar[u]
}
\]
soit commutatif~; si $K_1$ (resp. $L_1$) est considéré comme faisceau simple
sur $Y$ (resp. $X$), donc comme une $\mathcal{R}(Y)$-Algèbre (resp. une
$\mathcal{R}(X)$-Algèbre), cela signifie que $K_1\to L_1$ est un
$f$-morphisme. Cela étant, si $X'$ (resp. $Y'$) est la fermeture intégrale de
$X$ (resp. $Y$) relativement à $L_1$ (resp. $K_1$), $X'$ (resp. $Y'$) est un
préschéma intègre et normal (\hyperref[II.6.3.6-fr]{6.3.6}) dont le corps des
fonctions rationnelles s'identifie canoniquement à la fermeture algébrique

\oldpage[II]{120}

$L'$ (resp. $K'$) de $L$ (resp. $K$) dans $L_1$ (resp. $K_1$), et il existe
un morphisme canonique (nécessairement dominant) $f':X'\to Y'$ rendant
commutatif le diagramme (\hyperref[II.6.3.5.1-fr]{6.3.5.1}).

Le cas le plus important est celui où on prend $L_1=L$, $K_1$ étant donc une
extension de $K$ contenue dans $L$, et où on suppose $X$ intègre et
\emph{normal}, donc $X'=X$. Ce qui précède montre alors que lorsque $X$ est
normal, et que $Y'$ est la fermeture intégrale de $Y$ relative à un corps
$K_1\subset L=R(X)$, tout morphisme dominant $f:X\to Y$ se factorise en
\[
f:X\xrightarrow{f'}Y'\to Y
\]
où $f'$ est dominant~; d'ailleurs, lorsque le monomorphisme $K_1\to L$ est
donné, $f'$ est nécessairement unique, comme on le voit en se ramenant au cas
où $X$ et $Y$ sont affines. On voit donc que pour la donnée de $Y$, $L$ et
d'un $K$-monomorphisme $K_1\to L$, la fermeture intégrale $Y'$ de $Y$
relativement à $K_1$ est solution d'un \emph{problème universel}.
\end{env}

\begin{remark}[6.3.10]
\label{II.6.3.10-fr}
Reprenons les hypothèses de (\hyperref[II.6.3.6-fr]{6.3.6}), en supposant de
plus que chacune des algèbres $A_i$ soit de \emph{rang fini} sur $k(\xi_i)$
(ce qui entraîne que $A_i$ est composée directe d'un nombre fini de corps)~;
on peut affirmer dans certains cas que le morphisme structural $X'\to X$ est
non seulement entier, mais même \emph{fini}. Bornons-nous au cas où $X$ est
\emph{réduit}~; la question étant locale sur $X$, on peut en outre supposer
que $X$ est affine d'anneau $C$, et que $C$ n'a qu'un nombre fini d'idéaux
minimaux $\mathfrak{p}_i$ ($1\leq i\leq r$) tels que les
$C_i=C/\mathfrak{p}_i$ soient intègres~; $X'$ sera alors fini sur $X$ si la
fermeture intégrale de chaque $C_i$ dans toute extension de degré fini de son
corps des fractions est un $C$-module de type fini
(\hyperref[II.6.3.4-fr]{6.3.4}). On sait que cette condition est toujours
vérifiée lorsque $C$ est une \emph{algèbre de type fini sur un corps}
(\cite[t.~I, p.~267, th.~9]{I-13}), ou \emph{sur $\mathbf{Z}$}
(\cite[t.~I, p.~93, th.~3]{I-9}), ou \emph{sur un anneau local noethérien
complet} ([25], p.~298). On en conclut que $X'\to X$ sera un morphisme
fini lorsque $X$ est un schéma de \emph{type fini sur un corps}, ou sur
$\mathbf{Z}$, ou sur un anneau local noethérien complet.
\end{remark}

\subsection{Déterminant d'un endomorphisme de $\mathcal{O}_X$-Module.}

\begin{env}[6.4.1]
\label{II.6.4.1-fr}
Soient $A$ un anneau (commutatif), $E$ un $A$-module libre de rang $n$, $u$
un endomorphisme de $E$~; rappelons que pour définir le \emph{polynôme
caractéristique} de $u$, on considère l'endomorphisme $u\otimes1$ du
$A[T]$-module libre de rang $n$, $E\otimes_AA[T]$ ($T$ indéterminée), et on
pose
\[
\label{II.6.4.1.1-fr}
P(u,T)=\det(T\cdot I-(u\otimes1))
\tag{6.4.1.1}
\]
($I$ automorphisme identique de $E\otimes_AA[T]$). On a
\[
\label{II.6.4.1.2-fr}
P(u,T)=T^n-\sigma_1(u)T^{n-1}+\ldots+(-1)^n\sigma_n(u)
\tag{6.4.1.2}
\]
où $\sigma_i(u)$ est un élément de $A$, égal à un polynôme homogène de degré
$i$ (à coefficients entiers) par rapport aux éléments de la matrice de $u$
relative à une base quelconque de $E$~; on dit que les $\sigma_i(u)$ sont les
\emph{fonctions symétriques élémentaires de $u$}, et on a en particulier
$\sigma_1(u)=\operatorname{Tr}u$ et $\sigma_n(u)=\det u$. Rappelons qu'en
vertu du th. de Hamilton-Cayley, on a
\[
\label{II.6.4.1.3-fr}
P(u,u)=u^n-\sigma_1(u)u^{n-1}+\ldots+(-1)^n\sigma_n(u)=0
\tag{6.4.1.3}
\]

\oldpage[II]{121}

ce qui s'écrit aussi
\[
\label{II.6.4.1.4-fr}
(\det u)\cdot I_E=uQ(u)
\tag{6.4.1.4}
\]
($I_E$ automorphisme identique de $E$), avec
\[
\label{II.6.4.1.5-fr}
Q(u)=(-1)^{n+1}(u^{n-1}-\sigma_1(u)u^{n-2}+\ldots+(-1)^{n-1}\sigma_{n-1}(u)).
\tag{6.4.1.5}
\]

Soit $\varphi:A\to B$ un homomorphisme d'anneaux, faisant de $B$ une
$A$-algèbre~; considérons le $B$-module
$E_{(B)}=E\otimes_AB$ qui est libre de rang $n$, et l'extension $u\otimes1$
de $u$ à un endomorphisme de $E_{(B)}$~; il est immédiat que l'on a
$\sigma_i(u\otimes1)=\varphi(\sigma_i(u))$ pour tout indice $i$.
\end{env}

\begin{env}[6.4.2]
\label{II.6.4.2-fr}
Supposons maintenant que $A$ soit un anneau \emph{intègre}, de corps des
fractions $K$, et $E$ un $A$-module de \emph{type fini} (mais non
nécessairement libre). Soit $n$ le \emph{rang} de $E$, c'est-à-dire le rang du
$K$-module libre $E\otimes_AK$~; à tout endomorphisme $u$ de $E$ correspond
canoniquement l'endomorphisme $u\otimes1$ de $E\otimes_AK$~; par abus de
langage, on appelle encore \emph{polynôme caractéristique} de $u$ et on note
$P(u,T)$ le polynôme $P(u\otimes1,T)$, dont les coefficients
$\sigma_i(u\otimes1)$ (qui appartiennent à $K$) seront encore notés
$\sigma_i(u)$ et appelés les \emph{fonctions symétriques élémentaires de
$u$}~; en particulier $\det u=\det(u\otimes1)$ par définition. Avec ces
notations, les formules (\hyperref[II.6.4.1.3-fr]{6.4.1.3}) à
(\hyperref[II.6.4.1.5-fr]{6.4.1.5}) ont un sens et sont encore valables,
lorsque l'on interprète $u^j$ comme l'homomorphisme $E\to E\otimes_AK$
composé de l'endomorphisme $u^j\otimes1=(u\otimes1)^j$ de $E\otimes_AK$ et
de l'homomorphisme canonique $x\mapsto x\otimes1$.

Si $F$ est le sous-module de torsion de $E$, et si $E_0=E/F$, on a
$u(F)\subset F$, donc, par passage aux quotients, $u$ donne un endomorphisme
$u_0$ de $E_0$~; en outre $E\otimes_AK$ s'identifie à $E_0\otimes_AK$ et
$u\otimes1$ à $u_0\otimes1$, donc $\sigma_i(u)=\sigma_i(u_0)$ pour
$1\leq i\leq n$.

Lorsque $E$ est \emph{sans torsion}, $E$ s'identifie canoniquement à un
sous-$A$-module de $E\otimes_AK$, et la relation $u\otimes1=0$ est
équivalente à $u=0$. Lorsque $E$ est un $A$-module libre, les deux définitions
des $\sigma_i(u)$ données dans (\hyperref[II.6.4.1-fr]{6.4.1}) et
(\hyperref[II.6.4.2-fr]{6.4.2}) coïncident d'après ce qui précède, ce qui
justifie les notations adoptées.

On notera aussi que lorsque $E$ est un \emph{module de torsion}, autrement dit
$E_0=\{0\}$, l'algèbre extérieure de $E_0$ est réduite à $K$ et le déterminant
de l'unique endomorphisme $u_0$ de $E_0$ est \emph{égal à $1$}.
\end{env}

\begin{proposition}[6.4.3]
\label{II.6.4.3-fr}
Soient $A$ un anneau intègre, $E$ un $A$-module de type fini, $u$ un
endomorphisme de $E$~; alors les fonctions symétriques élémentaires
$\sigma_i(u)$ de $u$ (et en particulier $\det u$) sont des éléments de $K$
entiers sur $A$.
\end{proposition}

Soit $A'$ la clôture intégrale de $A$~; comme $A'[T]$ est un anneau
intégralement clos ([24], p.~99), c'est la clôture intégrale de
$A[T]$ dans son corps des fractions $K(T)$. Remplaçant $u$ par
$T\cdot I-u\otimes1$ et $A$ par $A[T]$, on voit qu'on est ramené à prouver que
$\det u$ est entier sur $A$. Si $n$ est le rang de $E$, on a
$\det u=\det(\bigwedge^n u)$ et
$(\bigwedge^n u)\otimes1=\bigwedge^n(u\otimes1)$, donc on peut supposer
$n=1$. Mais alors l'application $u\mapsto\det u$ est un homomorphisme du
$A$-module $\operatorname{Hom}_A(E,E)$ dans $K$~; comme $E$ est de type fini,
$\operatorname{Hom}_A(E,E)$ est isomorphe à un sous-module du $A$-module de
type fini $E^n$ (si $E$ admet un système de $n$ générateurs)~; donc les
éléments $\det u$ appartiennent à un sous-$A$-module de $K$ de type fini, et
sont par suite entiers sur $A$.

\oldpage[II]{122}

\begin{corollary}[6.4.4]
\label{II.6.4.4-fr}
Sous les hypothèses de (\hyperref[II.6.4.3-fr]{6.4.3}), si on suppose en outre
$A$ normal, les $\sigma_i(u)$ (en particulier $\operatorname{Tr}u$ et
$\det u$) appartiennent à $A$.
\end{corollary}

\begin{proposition}[6.4.5]
\label{II.6.4.5-fr}
Soient $A$ un anneau intègre, $E$ un $A$-module de type fini, de rang $n$, $u$
un endomorphisme de $E$ tel que les $\sigma_i(u)$ appartiennent à $A$. Pour
que $u$ soit un automorphisme de $E$, il est nécessaire que $\det u$ soit
inversible dans $A$~; cette condition est suffisante lorsque $E$ est sans
torsion.
\end{proposition}

La condition est \emph{suffisante}, car l'hypothèse et les formules
(\hyperref[II.6.4.1.4-fr]{6.4.1.4}) et
(\hyperref[II.6.4.1.5-fr]{6.4.1.5}) (valables dans $E$, et non seulement dans
$E\otimes_AK$, puisque $E$ est sans torsion) prouvent que
$(\det u)^{-1}Q(u)$ est l'inverse de $u$.

La condition est \emph{nécessaire}, car si $u$ est inversible, il résulte de
(\hyperref[II.6.4.3-fr]{6.4.3}) que $\det(u^{-1})$ appartient à la clôture
intégrale $A'$ de $A$ dans son corps des fractions $K$ et est évidemment
inverse de $\det u$ dans $A'$. Notre assertion résulte alors du~:

\begin{lemma}[6.4.5.1]
\label{II.6.4.5.1-fr}
Soit $A$ un sous-anneau d'un anneau $A'$ tel que $A'$ soit entier sur $A$. Si
un élément $x\in A$ est inversible dans $A'$, il est aussi inversible dans
$A$.
\end{lemma}

Dans le cas contraire, $x$ appartiendrait à un idéal maximal $\mathfrak{m}$
de $A$ et il résulte du premier th. de Cohen-Seidenberg
(\cite[t.~I, p.~257, th.~3]{I-13}) qu'il y aurait un idéal maximal
$\mathfrak{m}'$ de $A'$ tel que $\mathfrak{m}=\mathfrak{m}'\cap A$~; on aurait
donc $x\in\mathfrak{m}'$, ce qui est absurde.

\begin{corollary}[6.4.6]
\label{II.6.4.6-fr}
Soient $A$ un anneau intègre et intégralement clos, $E$ un $A$-module de type
fini sans torsion, $u$ un endomorphisme de $E$. Pour que $u$ soit un
automorphisme de $E$, il faut et il suffit que $\det u$ soit inversible dans
$A$.
\end{corollary}

Cela résulte de (\hyperref[II.6.4.4-fr]{6.4.4}) et
(\hyperref[II.6.4.5-fr]{6.4.5}).

\begin{remark}[6.4.7]
\label{II.6.4.7-fr}
Nous aurons besoin plus loin d'une généralisation des résultats précédents.
Considérons un anneau noethérien \emph{réduit} $A$~; soient
$\mathfrak{p}_\alpha$ ($1\leq\alpha\leq r$) ses idéaux minimaux, $K_\alpha$
le corps des fractions de l'anneau intègre $A/\mathfrak{p}_\alpha$, $K$
l'anneau total des fractions de $A$, \emph{composé direct} des corps
$K_\alpha$. Soit $E$ un $A$-module de type fini, et \emph{supposons que}
$E\otimes_AK$ soit un $K$-module libre de rang $n$ (ce qui ici n'est plus
conséquence des autres hypothèses)~; il revient au même de dire que tous les
$K_\alpha$-espaces vectoriels $E\otimes_AK_\alpha=E_\alpha$ ont \emph{même
dimension $n$}~; si alors $u$ est un endomorphisme de $E$, on pose encore
$P(u,T)=P(u\otimes1,T)$ et $\sigma_j(u)=\sigma_j(u\otimes1)$, et en
particulier $\det u=\det(u\otimes1)$~; les $\sigma_j(u)$ sont donc des
éléments de $K$. Il est immédiat que $E\otimes_AK$ est somme directe des
$E_\alpha$ et que ces derniers sont stables par $u\otimes1$~; la restriction
de $u\otimes1$ à $E_\alpha$ n'est autre d'ailleurs que l'extension $u_\alpha$
de $u$ à ce $K_\alpha$-espace vectoriel~; on en conclut que $\sigma_j(u)$ est
l'élément de $K$ dont les composantes dans les $K_\alpha$ sont les
$\sigma_j(u_\alpha)$. Comme la fermeture intégrale de $A$ dans $K$ est
\emph{composée directe} des fermetures intégrales de $A$ dans les
$K_\alpha$, les $\sigma_j(u)$ sont \emph{entiers} sur $A$.

\begin{lemma}[6.4.7.1]
\label{II.6.4.7.1-fr}
La sous-$A$-algèbre de $K$ engendrée par tous les éléments $\sigma_j(u)$
($1\leq j\leq n$) lorsque $u$ parcourt $\operatorname{Hom}_A(E,E)$, est un
$A$-module de type fini.
\end{lemma}

Il suffit de démontrer que la sous-$A[T]$-algèbre de $K[T]$ engendrée par les
$P(u,T)$ est un $A[T]$-module de type fini, car si les $F_i(T)$
($1\leq i\leq m$) forment un système de générateurs de ce $A[T]$-module, les
coefficients des $F_i(T)$ sont entiers sur $A$, donc engendrent une
$A$-algèbre qui est un $A$-module de type fini
(\cite[t.~I, p.~255, th.~1]{I-13}). On peut donc

\oldpage[II]{123}

remplacer $A$ par $A[T]$ (qui est noethérien) et $E$ par
$E\otimes_AA[T]=E'$ qui est tel que
$E'\otimes_{A[T]}K[T]=E\otimes_AK[T]$ soit un $K[T]$-module libre de rang $n$.
Revenant aux notations initiales, on voit donc qu'il suffit de prouver que le
$A$-module engendré par les éléments $\det u$, où $u$ parcourt
$\operatorname{Hom}_A(E,E)$, est de type fini~; \emph{a fortiori} (puisque
tout sous-module d'un $A$-module de type fini est de type fini) il suffit de
prouver que lorsque $v$ parcourt l'ensemble des endomorphismes de
$\bigwedge^nE$, le $A$-module engendré par les $\det v$ est de type fini~;
autrement dit, on est encore ramené au cas où $n=1$. Mais alors la proposition
résulte de ce que $\operatorname{Hom}_A(E,E)$ est un $A$-module de type fini
et de ce que $v\mapsto\det v$ est un homomorphisme de ce $A$-module dans $K$.

Soit $F$ le noyau de l'homomorphisme canonique $E\to E\otimes_AK$ et soit
$E_0=E/F$~; on voit comme ci-dessus que $E\otimes_AK$ s'identifie à
$E_0\otimes_AK$, que $u(F)\subset F$, et que si $u_0$ est l'endomorphisme de
$E_0$ déduit de $u$ par passage aux quotients, $u\otimes1$ s'identifie à
$u_0\otimes1$, et $\sigma_j(u)=\sigma_j(u_0)$ pour tout $j$. \emph{Si l'on a
$F=0$}, les formules (\hyperref[II.6.4.1.3-fr]{6.4.1.3}) et
(\hyperref[II.6.4.1.5-fr]{6.4.1.5}) ont un sens et sont valables lorsque $E$
est identifié à un sous-module de $E\otimes_AK$ et les $u^j$ à des
homomorphismes $E\to E\otimes_AK$~; par suite, la prop.
(\hyperref[II.6.4.5-fr]{6.4.5}) s'étend aussi à ce cas, avec sa démonstration.
\end{remark}

\begin{env}[6.4.8]
\label{II.6.4.8-fr}
Soient $(X,\mathcal{A})$ un espace annelé, $\mathcal{E}$ un
$\mathcal{A}$-Module \emph{localement libre} (de rang fini), $u$ un
endomorphisme de $\mathcal{E}$. Il y a par hypothèse une base $\mathfrak{B}$
de la topologie de $X$ telle que pour tout $V\in\mathfrak{B}$,
$\mathcal{E}|V$ soit isomorphe à $\mathcal{A}^n|V$ (pour un $n$ pouvant varier
avec $V$). Soit $u$ un endomorphisme de $\mathcal{E}$~; pour tout
$V\in\mathfrak{B}$, $u_V$ est donc un endomorphisme du
$\Gamma(V,\mathcal{A})$-module $\Gamma(V,\mathcal{E})$, qui est libre par
hypothèse~; le déterminant $\det u_V$ est donc défini et appartient à
$\Gamma(V,\mathcal{A})$. En outre, si $e_1,\ldots,e_n$ forment une base de
$\Gamma(V,\mathcal{E})$, leurs restrictions à tout ouvert $W\subset V$
forment une base de $\Gamma(W,\mathcal{E})$ sur $\Gamma(W,\mathcal{A})$,
donc $\det u_W$ est la restriction de $\det u_V$ à $W$. Il existe donc une
section et une seule de $\mathcal{A}$ au-dessus de $X$, que nous noterons
$\det u$ et appellerons le \emph{déterminant de $u$}, telle que la restriction
de $\det u$ à tout $V\in\mathfrak{B}$ soit $\det u_V$. Il est clair que pour
tout $x\in X$, on a $(\det u)_x=\det u_x$~; pour deux endomorphismes $u$, $v$
de $\mathcal{E}$, on a
\[
\label{II.6.4.8.1-fr}
\det(u\circ v)=(\det u)(\det v)
\tag{6.4.8.1}
\]
ainsi que
\[
\label{II.6.4.8.2-fr}
\det(1_{\mathcal{E}})=1_{\mathcal{A}}
\tag{6.4.8.2}
\]
et, si le rang de $\mathcal{E}$ est constant (ce qui sera le cas
(\hyperref[0.5.4.1-fr]{0, 5.4.1}) pour $X$ connexe) et égal à $n$
\[
\label{II.6.4.8.3-fr}
\det(s\cdot u)=s^n\det u
\tag{6.4.8.3}
\]
pour tout $s\in\Gamma(X,\mathcal{A})$ (on notera que
$\det(0)=0_{\mathcal{A}}$ si $n\geq1$, mais $\det(0)=1_{\mathcal{A}}$ pour
$n=0$). En outre, pour que $u$ soit un \emph{automorphisme} de
$\mathcal{E}$, il faut et il suffit que $\det u$ soit \emph{inversible} dans
$\Gamma(X,\mathcal{A})$.

Si le rang de $\mathcal{E}$ est constant, on peut définir de la même manière
les fonctions symétriques élémentaires $\sigma_i(u)$, qui sont des éléments
de $\Gamma(X,\mathcal{A})$~; on a encore les relations
(\hyperref[II.6.4.1.3-fr]{6.4.1.3}) à
(\hyperref[II.6.4.1.5-fr]{6.4.1.5}).

\oldpage[II]{124}

On a ainsi défini un homomorphisme
$\det:\operatorname{Hom}_{\mathcal{A}}(\mathcal{E},\mathcal{E})
\to\Gamma(X,\mathcal{A})$ de monoïdes multiplicatifs~; si on note que
$\operatorname{Hom}_{\mathcal{A}}(\mathcal{E},\mathcal{E})=
\Gamma(X,\mathcal{H}om_{\mathcal{A}}(\mathcal{E},\mathcal{E}))$ par
définition, on voit qu'on peut dans cette définition remplacer $X$ par un
ouvert quelconque $U$ de $X$, ce qui définit immédiatement un homomorphisme
$\det:\mathcal{H}om_{\mathcal{A}}(\mathcal{E},\mathcal{E})\to\mathcal{A}$ de
faisceaux de monoïdes multiplicatifs. Lorsque $\mathcal{E}$ a un rang
constant, on définit de même des homomorphismes
$\sigma_i:\mathcal{H}om_{\mathcal{A}}(\mathcal{E},\mathcal{E})
\to\mathcal{A}$ de faisceaux d'ensembles~; pour $i=1$, l'homomorphisme
$\sigma_1=\operatorname{Tr}$ est un homomorphisme de $\mathcal{A}$-Modules.

Soit $(Y,\mathcal{B})$ un second espace annelé, et soit
$f:(X,\mathcal{A})\to(Y,\mathcal{B})$ un morphisme d'espaces annelés~; si
$\mathcal{F}$ est un $\mathcal{B}$-Module localement libre,
$f^*(\mathcal{F})$ est un $\mathcal{A}$-Module localement libre (qui a même
rang que $\mathcal{F}$ si ce dernier est de rang constant)
(\hyperref[0.5.4.5-fr]{0, 5.4.5}). Pour tout endomorphisme $v$ de
$\mathcal{F}$, $f^*(v)$ est un endomorphisme de $f^*(\mathcal{F})$, et il
résulte aussitôt des définitions que $\det f^*(v)$ est la section de
$\mathcal{A}=f^*(\mathcal{B})$ au-dessus de $X$ qui correspond canoniquement
à $\det v\in\Gamma(Y,\mathcal{B})$. On peut encore dire que l'homomorphisme
$f^*(\det):f^*(\mathcal{H}om_{\mathcal{B}}(\mathcal{F},\mathcal{F}))
\to f^*(\mathcal{B})=\mathcal{A}$ est le composé
\[
f^*(\mathcal{H}om_{\mathcal{B}}(\mathcal{F},\mathcal{F}))
\xrightarrow{\gamma^\sharp}
\mathcal{H}om_{\mathcal{A}}(f^*(\mathcal{F}),f^*(\mathcal{F}))
\xrightarrow{\det}\mathcal{A}
\]
(\hyperref[0.4.4.6-fr]{0, 4.4.6}). On a des résultats analogues pour les
$\sigma_i$.
\end{env}

\begin{env}[6.4.9]
\label{II.6.4.9-fr}
Supposons maintenant que $X$ soit un \emph{préschéma localement intègre}, de
sorte que le faisceau $\mathcal{R}(X)$ des fonctions rationnelles sur $X$ est
un faisceau de corps localement simple (\hyperref[I.7.3.4-fr]{I, 7.3.4}),
quasi-cohérent en tant que $\mathcal{O}_X$-Module. Si $\mathcal{E}$ est un
$\mathcal{O}_X$-Module quasi-cohérent \emph{de type fini},
$\mathcal{E}'=\mathcal{E}\otimes_{\mathcal{O}_X}\mathcal{R}(X)$ est alors un
$\mathcal{R}(X)$-Module localement libre
(\hyperref[I.7.3.6-fr]{I, 7.3.6})~; pour tout endomorphisme $u$ de
$\mathcal{E}$, $u\otimes1_{\mathcal{R}(X)}$ est alors un endomorphisme de
$\mathcal{E}'$, et $\det(u\otimes1)$ une section de $\mathcal{R}(X)$ au-dessus
de $X$, que l'on appelle encore \emph{déterminant de $u$} et que l'on note
encore $\det u$. Il résulte de (\hyperref[II.6.4.3-fr]{6.4.3}) que $\det u$ est
une section de la \emph{fermeture intégrale} de $\mathcal{O}_X$ dans
$\mathcal{R}(X)$ (\hyperref[II.6.3.2-fr]{6.3.2})~; si en outre $X$ est
\emph{normal}, $\det u$ est donc une \emph{section de $\mathcal{O}_X$}
au-dessus de $X$, et si on suppose de plus que $\mathcal{E}$ est \emph{sans
torsion}, pour que $u$ soit un automorphisme de $\mathcal{E}$, il faut et il
suffit que $\det u$ soit \emph{inversible} en vertu de
(\hyperref[II.6.4.6-fr]{6.4.6}). Les formules
(\hyperref[II.6.4.8.1-fr]{6.4.8.1}) à
(\hyperref[II.6.4.8.3-fr]{6.4.8.3}) sont encore valables~; de l'homomorphisme
$u\mapsto\det u$, appliqué aux modules de sections de
$\mathcal{H}om_{\mathcal{O}_X}(\mathcal{E},\mathcal{E})$, on déduit un
homomorphisme de faisceaux
$\det:\mathcal{H}om_{\mathcal{O}_X}(\mathcal{E},\mathcal{E})
\to\mathcal{R}(X)$, qui prend ses valeurs dans $\mathcal{O}_X$ quand $X$ est
normal. On a des définitions et résultats analogues pour les autres fonctions
symétriques élémentaires $\sigma_j(u)$, lorsque $\mathcal{E}'$ est de
\emph{rang constant}~; si de plus $X$ est \emph{normal}, les $\sigma_j(u)$
sont des \emph{sections de $\mathcal{O}_X$} au-dessus de $X$.

Enfin, soient $X$, $Y$ deux préschémas intègres, et soit $f:X\to Y$ un
morphisme \emph{dominant}. On sait qu'il existe alors un homomorphisme
canonique $f^*(\mathcal{R}(Y))\to\mathcal{R}(X)$
(\hyperref[I.7.3.8-fr]{I, 7.3.8}\footnote{Voir \emph{Errata et Addenda},
liste 1.}), d'où résulte, pour tout
$\mathcal{O}_Y$-Module quasi-cohérent de type fini $\mathcal{F}$, un
homomorphisme canonique
$\theta:f^*(\mathcal{F}\otimes_{\mathcal{O}_Y}\mathcal{R}(Y))
\to f^*(\mathcal{F})\otimes_{\mathcal{O}_X}\mathcal{R}(X)$. Si $v$ est un
endo-

\oldpage[II]{125}

morphisme de $\mathcal{F}$,
$f^*(v\otimes1_{\mathcal{R}(Y)})$ est un endomorphisme de
$f^*(\mathcal{F}\otimes_{\mathcal{O}_Y}\mathcal{R}(Y))$, et on a un diagramme
commutatif
\[
\xymatrix{
f^*(\mathcal{F}\otimes_{\mathcal{O}_Y}\mathcal{R}(Y))
  \ar[r]^{f^*(v\otimes1)}\ar[d]_{\theta}
& f^*(\mathcal{F}\otimes_{\mathcal{O}_Y}\mathcal{R}(Y))\ar[d]^{\theta}\\
f^*(\mathcal{F})\otimes_{\mathcal{O}_X}\mathcal{R}(X)
  \ar[r]_{f^*(v)\otimes1}
& f^*(\mathcal{F})\otimes_{\mathcal{O}_X}\mathcal{R}(X).
}
\]
On en conclut aisément que $\det f^*(v)$ est l'image canonique, par
l'homomorphisme $f^*(\mathcal{R}(Y))\to\mathcal{R}(X)$, de la section
$\det v$ de $\mathcal{R}(Y)$~; on est en effet aussitôt ramené au cas où
$X=\operatorname{Spec}(A)$ et $Y=\operatorname{Spec}(B)$ sont affines, $A$ et
$B$ étant des anneaux intègres de corps de fractions respectifs $K$ et $L$,
l'homomorphisme $B\to A$ étant injectif et s'étendant donc en un monomorphisme
$L\to K$~; si $\mathcal{F}=\widetilde{M}$, où $M$ est un $B$-module de type
fini, le rang sur $L$ de $M\otimes_B L$ est \emph{égal} à celui de
$(M\otimes_B A)\otimes_A K$ sur $K$, et $\det((u\otimes1)\otimes1)$ est
l'image dans $K$ de $\det(u\otimes1)$ pour tout endomorphisme $u$ de $M$,
d'où la conclusion.
\end{env}

\begin{env}[6.4.10]
\label{II.6.4.10-fr}
Supposons enfin que $X$ soit un \emph{préschéma localement noethérien réduit},
de sorte que le faisceau $\mathcal{R}(X)$ des fonctions rationnelles sur $X$
est encore un $\mathcal{O}_X$-Module quasi-cohérent
(\hyperref[I.7.3.4-fr]{I, 7.3.4})~; soit d'autre part $\mathcal{E}$ un
$\mathcal{O}_X$-Module \emph{cohérent} tel que
$\mathcal{E}'=\mathcal{E}\otimes_{\mathcal{O}_X}\mathcal{R}(X)$ soit
\emph{localement libre} (de rang fini). En vertu de
(\hyperref[II.6.4.7-fr]{6.4.7}), si $\mathcal{E}'$ est de rang constant, on
peut, pour tout endomorphisme $u$ de $\mathcal{E}$, définir les $\sigma_j(u)$,
qui sont des sections de $\mathcal{R}(X)$ au-dessus de $X$. Lorsqu'on ne
suppose plus $\mathcal{E}'$ de rang constant, on peut encore définir
l'homomorphisme
$\det:\mathcal{H}om_{\mathcal{O}_X}(\mathcal{E},\mathcal{E})
\to\mathcal{R}(X)$.
\end{env}

\subsection{Norme d'un faisceau inversible.}
\label{subsection:II.6.5-fr}

\begin{env}[6.5.1]
\label{II.6.5.1-fr}
Soit $(X,\mathcal{A})$ un espace annelé et soit $\mathcal{B}$ une
$\mathcal{A}$-Algèbre (commutative). Le $\mathcal{A}$-Module $\mathcal{B}$
s'identifie canoniquement à un sous-$\mathcal{A}$-Module de
$\mathcal{H}om_{\mathcal{A}}(\mathcal{B},\mathcal{B})$, une section $f$ de
$\mathcal{B}$ au-dessus d'un ouvert $U$ de $X$ s'identifiant à la
multiplication par cette section. Si $(X,\mathcal{A})$ et $\mathcal{B}$
vérifient l'une des conditions énumérées dans
(\hyperref[II.6.4.8-fr]{6.4.8}), (\hyperref[II.6.4.9-fr]{6.4.9}) ou
(\hyperref[II.6.4.10-fr]{6.4.10}), on peut donc définir $\det(f)$ (et dans
certains cas les $\sigma_j(f)$) qui sont des sections de $\mathcal{A}$ ou de
$\mathcal{R}(X)$ au-dessus de $U$, que nous appellerons la \emph{norme de
$f$} (resp. les \emph{fonctions symétriques élémentaires de $f$}) et noterons
$N_{\mathcal{B}/\mathcal{A}}(f)$. Nous allons supposer vérifiée \emph{l'une}
des conditions suivantes~:
\begin{enumerate}
\item[(I)] \emph{$\mathcal{B}$ est un $\mathcal{A}$-Module localement libre
(de rang fini).}
\item[(II)] \emph{$(X,\mathcal{A})$ est un préschéma localement noethérien
réduit, $\mathcal{B}$ est un $\mathcal{A}$-Module cohérent tel que
$\mathcal{B}\otimes_{\mathcal{A}}\mathcal{R}(X)$ soit un
$\mathcal{R}(X)$-Module localement libre, et pour toute section
$f\in\Gamma(U,\mathcal{B})$ au-dessus d'un ouvert $U\subset X$,
$N_{\mathcal{B}/\mathcal{A}}(f)$ est une section de $\mathcal{A}$ au-dessus
de $U$.}
\end{enumerate}

\oldpage[II]{126}

On notera que cette dernière condition est automatiquement vérifiée lorsque
le préschéma localement noethérien $X$ est \emph{normal}
(\hyperref[II.6.4.9-fr]{6.4.9}).

D'autre part, l'hypothèse que
$\mathcal{B}\otimes_{\mathcal{A}}\mathcal{R}(X)$ soit localement libre peut
encore s'exprimer de la façon suivante~: désignons par $X_\alpha$ les
sous-préschémas fermés réduits de $X$ ayant pour espaces sous-jacents les
composantes irréductibles de $X$ (\hyperref[I.5.2.1-fr]{I, 5.2.1}), qui sont
donc des préschémas intègres localement noethériens. Tout $x\in X$ appartient
à un nombre fini des sous-espaces $X_\alpha$~; d'autre part,
$\mathcal{B}\otimes_{\mathcal{A}}\mathcal{R}(X_\alpha)$ est un
$\mathcal{R}(X_\alpha)$-Module localement libre de rang constant $k_\alpha$
(\hyperref[I.7.3.6-fr]{I, 7.3.6})~; dire que
$\mathcal{B}\otimes_{\mathcal{A}}\mathcal{R}(X)$ est un
$\mathcal{R}(X)$-Module localement libre signifie que, pour tout $x\in X$,
\emph{les rangs $k_\alpha$ correspondant aux indices tels que $x\in X_\alpha$
sont égaux}. La question est en effet locale, et on est ramené au cas
$X=\operatorname{Spec}(C)$, où $C$ est un anneau noethérien réduit, et
$\mathcal{B}=\widetilde{D}$, où $D$ est une $C$-algèbre qui est un $C$-module
de type fini~; si $\mathfrak{p}_i$ ($1\leq i\leq m$) sont les idéaux premiers
minimaux de $C$, l'anneau total des fractions $L$ de $C$ est composé direct
des corps des fractions $K_i$ des anneaux intègres $C/\mathfrak{p}_i$, et
$D\otimes_C L$ est somme directe des $D\otimes_C K_i$, d'où la conclusion.

Il est clair que sous les hypothèses (I) ou (II), on a défini ainsi un
\emph{homomorphisme de faisceaux de monoïdes multiplicatifs}
$N_{\mathcal{B}/\mathcal{A}}:\mathcal{B}\to\mathcal{A}$, aussi noté $N$ si
aucune confusion n'en résulte, et appelé l'homomorphisme \emph{norme}. Pour
deux sections $f$, $g$ de $\mathcal{B}$ au-dessus d'un même ouvert $U$, on a
donc
\[
\label{II.6.5.1.1-fr}
N_{\mathcal{B}/\mathcal{A}}(fg)
=N_{\mathcal{B}/\mathcal{A}}(f)N_{\mathcal{B}/\mathcal{A}}(g)
\tag{6.5.1.1}
\]
pour les sections correspondantes de $\mathcal{A}$ au-dessus de $U$~;
\[
\label{II.6.5.1.2-fr}
N_{\mathcal{B}/\mathcal{A}}(1_{\mathcal{B}})=1_{\mathcal{A}}~;
\tag{6.5.1.2}
\]
enfin, pour toute section $s$ de $\mathcal{A}$ au-dessus de $U$
\[
\label{II.6.5.1.3-fr}
N_{\mathcal{B}/\mathcal{A}}(s\cdot1_{\mathcal{B}})=s^n
\tag{6.5.1.3}
\]
si le rang de $\mathcal{B}$ est constant et égal à $n$ (pour
$s=0_{\mathcal{A}}$, cette formule donne
$N(0_{\mathcal{B}})=0_{\mathcal{A}}$ si $n\geq1$,
$N(0_{\mathcal{B}})=N(1_{\mathcal{B}})=1_{\mathcal{A}}$ si $n=0$).

Dans l'hypothèse (I), pour que $f\in\Gamma(U,\mathcal{B})$ soit inversible,
il faut et il suffit que $N(f)\in\Gamma(U,\mathcal{A})$ le soit. Dans
l'hypothèse (II), cette condition est nécessaire~; elle est suffisante (par
(\hyperref[II.6.4.7-fr]{6.4.7})) lorsqu'on suppose que
$\mathcal{B}\to\mathcal{B}\otimes_{\mathcal{A}}\mathcal{R}(X)$ est
\emph{injectif} et que l'hypothèse plus restrictive suivante est vérifiée~:
\begin{enumerate}
\item[(II \emph{bis})] \emph{$(X,\mathcal{A})$ est un préschéma localement
noethérien réduit, $\mathcal{B}$ est un $\mathcal{A}$-Module cohérent tel que
$\mathcal{B}\otimes_{\mathcal{A}}\mathcal{R}(X)$ soit un
$\mathcal{R}(X)$-Module localement libre, et pour toute section
$f\in\Gamma(U,\mathcal{B})$ au-dessus d'un ouvert tel que
$\mathcal{B}\otimes_{\mathcal{A}}\mathcal{R}(X)|U$ soit de rang constant $n$
sur $\mathcal{R}(X)|U$, les $\sigma_j(f)$ ($1\leq j\leq n$) sont des sections
de $\mathcal{A}$ au-dessus de $U$.}
\end{enumerate}
(On notera encore que cette condition est vérifiée lorsque $X$ est
\emph{normal}.)
\end{env}

\begin{env}[6.5.2]
\label{II.6.5.2-fr}
Supposons vérifiée une des hypothèses (I), (II) de
(\hyperref[II.6.5.1-fr]{6.5.1}), et soit $\mathcal{L}'$ un
$\mathcal{B}$-Module \emph{inversible}. Nous allons lui associer canoniquement
(à un isomorphisme unique près) un $\mathcal{A}$-Module \emph{inversible}
$\mathcal{L}$ de la façon suivante. Désignons par $\mathcal{A}^*$ (resp.
$\mathcal{B}^*$) le sous-faisceau de $\mathcal{A}$ (resp. $\mathcal{B}$) tel
que $\Gamma(U,\mathcal{A}^*)$ (resp. $\Gamma(U,\mathcal{B}^*)$) soit
l'ensemble des

\oldpage[II]{127}

éléments inversibles de $\Gamma(U,\mathcal{A})$ (resp.
$\Gamma(U,\mathcal{B})$) pour tout ouvert $U\subset X$~; ce sont des faisceaux
de groupes multiplicatifs, et $N_{\mathcal{B}/\mathcal{A}}$, restreint à
$\mathcal{B}^*$, est un \emph{homomorphisme}
$\mathcal{B}^*\to\mathcal{A}^*$ de faisceaux de groupes
(\hyperref[II.6.5.1-fr]{6.5.1}). Soit $\mathfrak{L}$ l'ensemble des couples
$(U_\lambda,\eta_\lambda)$, ayant la propriété suivante~: $U_\lambda$ est un
ouvert de $X$, et $\eta_\lambda$ est un isomorphisme
$\mathcal{L}'|U_\lambda\simeq\mathcal{B}|U_\lambda$ de
$(\mathcal{B}|U_\lambda)$-Modules. Par hypothèse, les $U_\lambda$ forment un
recouvrement de $X$~; pour deux indices quelconques $\lambda$, $\mu$, on pose
$\omega_{\lambda\mu}=(\eta_\lambda|U_\lambda\cap U_\mu)
\circ(\eta_\mu|U_\lambda\cap U_\mu)^{-1}$, automorphisme de
$\mathcal{B}|U_\lambda\cap U_\mu$, qui s'identifie canoniquement à une section
de $\mathcal{B}^*$ au-dessus de $U_\lambda\cap U_\mu$, et
$(\omega_{\lambda\mu})$ est un $1$-cocycle du recouvrement
$\mathfrak{U}=(U_\lambda)$ à valeurs dans $\mathcal{B}^*$
(\hyperref[0.5.4.7-fr]{0, 5.4.7}). Le fait que
$N_{\mathcal{B}/\mathcal{A}}:\mathcal{B}^*\to\mathcal{A}^*$ est un
homomorphisme entraîne alors que
$(N_{\mathcal{B}/\mathcal{A}}\omega_{\lambda\mu})$ est un $1$-cocycle de
$\mathfrak{U}$ à valeurs dans $\mathcal{A}^*$, et il lui correspond donc (à un
isomorphisme unique près) un $\mathcal{A}$-Module inversible $\mathcal{L}$
(\hyperref[0.5.4.7-fr]{0, 5.4.7}) que nous désignerons par
$N_{\mathcal{B}/\mathcal{A}}(\mathcal{L}')$ et appellerons \emph{la norme du
$\mathcal{B}$-Module inversible $\mathcal{L}'$}.

Soit $\mathfrak{M}$ une partie de $\mathfrak{L}$ telle que les $U_\lambda$
correspondants forment encore un recouvrement de $X$, et soit $\mathfrak{B}$
ce recouvrement~; la restriction du cocycle $(\omega_{\lambda\mu})$ à
$\mathfrak{B}$ définit un $1$-cocycle
$(N_{\mathcal{B}/\mathcal{A}}\omega_{\lambda\mu})$ de $\mathfrak{B}$ à
valeurs dans $\mathcal{A}^*$, restriction du $1$-cocycle
$(N_{\mathcal{B}/\mathcal{A}}\omega_{\lambda\mu})$ de $\mathfrak{U}$~; il est
clair qu'il y a un isomorphisme canonique du $\mathcal{A}$-Module inversible
défini par ce $1$-cocycle de $\mathfrak{B}$ sur
$N_{\mathcal{B}/\mathcal{A}}(\mathcal{L}')$, ce qui permet de définir ce
$\mathcal{A}$-Module inversible par un sous-recouvrement quelconque de
$\mathfrak{U}$. Cette possibilité montre aussitôt que, si
$\mathcal{L}'_1$, $\mathcal{L}'_2$ sont deux $\mathcal{B}$-Modules inversibles,
on a, en vertu de (\hyperref[II.6.5.1.1-fr]{6.5.1.1}) et
(\hyperref[II.6.5.1.2-fr]{6.5.1.2}),
\[
\label{II.6.5.2.1-fr}
N(\mathcal{L}'_1\otimes_{\mathcal{B}}\mathcal{L}'_2)
=N(\mathcal{L}'_1)\otimes_{\mathcal{A}}N(\mathcal{L}'_2)
\tag{6.5.2.1}
\]
et
\[
\label{II.6.5.2.2-fr}
N_{\mathcal{B}/\mathcal{A}}(\mathcal{B})=\mathcal{A}
\tag{6.5.2.2}
\]
ainsi que
\[
\label{II.6.5.2.3-fr}
N((\mathcal{L}')^{-1})=(N(\mathcal{L}'))^{-1}
\tag{6.5.2.3}
\]
à des isomorphismes canoniques près. En outre il résulte de
(\hyperref[II.6.5.1.3-fr]{6.5.1.3}) que si $\mathcal{L}$ est un
$\mathcal{A}$-Module inversible et si $\mathcal{B}$ est de rang constant $n$
sur $\mathcal{A}$ dans le cas (I) (resp.
$\mathcal{B}\otimes_{\mathcal{A}}\mathcal{R}(X)$ de rang constant $n$ sur
$\mathcal{R}(X)$ dans le cas (II)), on a, à un isomorphisme canonique près
\[
\label{II.6.5.2.4-fr}
N_{\mathcal{B}/\mathcal{A}}(\mathcal{L}\otimes_{\mathcal{A}}\mathcal{B})
=\mathcal{L}^{\otimes n}.
\tag{6.5.2.4}
\]
\end{env}

\begin{env}[6.5.3]
\label{II.6.5.3-fr}
Montrons maintenant que $N_{\mathcal{B}/\mathcal{A}}(\mathcal{L}')$ est un
\emph{foncteur covariant} dans la catégorie des $\mathcal{B}$-Modules
inversibles. Soit en effet $h':\mathcal{L}'_1\to\mathcal{L}'_2$ un
homomorphisme de $\mathcal{B}$-Modules inversibles, et soit
$\mathfrak{B}=(U_\lambda)$ un recouvrement ouvert de $X$ tel que pour tout
$\lambda$, on ait des $(\mathcal{B}|U_\lambda)$-isomorphismes
$\eta^{(1)}_\lambda:\mathcal{L}'_1|U_\lambda\simeq\mathcal{B}|U_\lambda$ et
$\eta^{(2)}_\lambda:\mathcal{L}'_2|U_\lambda\simeq\mathcal{B}|U_\lambda$~;
il y a donc pour chaque $\lambda$ un endomorphisme $h'_\lambda$ de
$\mathcal{B}|U_\lambda$ tel que
$h'_\lambda\circ\eta^{(1)}_\lambda
=\eta^{(2)}_\lambda\circ(h'|U_\lambda)$, et on peut évidemment identifier
$h'_\lambda$ à une section de $\mathcal{B}$ au-dessus de $U_\lambda$
(\hyperref[0.5.1.1-fr]{0, 5.1.1}). Donc, pour tout couple $(\lambda,\mu)$
d'indices, les restrictions à $U_\lambda\cap U_\mu$ de
$(\eta^{(2)}_\lambda)^{-1}\circ h'_\lambda\circ\eta^{(1)}_\lambda$ et de
$(\eta^{(2)}_\mu)^{-1}\circ h'_\mu\circ\eta^{(1)}_\mu$ coïncident. On en
déduit pour les $1$-cocycles $(\omega^{(1)}_{\lambda\mu})$ et
$(\omega^{(2)}_{\lambda\mu})$ à valeurs dans $\mathcal{B}^*$ correspondant à
$\mathcal{L}'_1$ et $\mathcal{L}'_2$, la relation
\[
\omega^{(2)}_{\lambda\mu}h'_\mu
=h'_\lambda\omega^{(1)}_{\lambda\mu}.
\]

\oldpage[II]{128}

Si on pose $h_\lambda=N(h'_\lambda)$, on aura donc les relations analogues
\[
N(\omega^{(2)}_{\lambda\mu})h_\mu
=h_\lambda N(\omega^{(1)}_{\lambda\mu})
\]
et par suite les $h_\lambda$ définissent un homomorphisme
$N(\mathcal{L}'_1)\to N(\mathcal{L}'_2)$ que nous désignerons par
$N_{\mathcal{B}/\mathcal{A}}(h')$ ou $N(h')$. Dans l'hypothèse (I), pour que
$h'$ soit un \emph{isomorphisme}, il faut et il suffit (puisqu'il s'agit d'une
question locale) que $N_{\mathcal{B}/\mathcal{A}}(h')$ le soit. Dans
l'hypothèse (II), cette condition est encore nécessaire~; elle est suffisante
si l'hypothèse (II \emph{bis}) est vérifiée et si
$\mathcal{B}\to\mathcal{B}\otimes_{\mathcal{A}}\mathcal{R}(X)$ est injectif.

Prenons en particulier $\mathcal{L}'_1=\mathcal{B}$~; les homomorphismes
$\mathcal{B}\to\mathcal{L}'$ s'identifient alors
(\hyperref[0.5.1.1-fr]{0, 5.1.1}) aux sections de $\mathcal{L}'$ au-dessus de
$X$, d'où une application canonique
\[
N_{\mathcal{B}/\mathcal{A}}:\Gamma(X,\mathcal{L}')
\to\Gamma(X,N_{\mathcal{B}/\mathcal{A}}(\mathcal{L}')).
\]
Il résulte encore de (\hyperref[II.6.5.1.1-fr]{6.5.1.1}) que si
$f'_1\in\Gamma(X,\mathcal{L}'_1)$, $f'_2\in\Gamma(X,\mathcal{L}'_2)$, on a
\[
\label{II.6.5.3.1-fr}
N(f'_1\otimes f'_2)=N(f'_1)\otimes N(f'_2).
\tag{6.5.3.1}
\]
Pour tout $\mathcal{A}$-Module inversible $\mathcal{L}$ et toute section
$f\in\Gamma(X,\mathcal{L})$, on a, compte tenu de
(\hyperref[II.6.5.2.4-fr]{6.5.2.4})
\[
\label{II.6.5.3.2-fr}
N_{\mathcal{B}/\mathcal{A}}(f\otimes1_{\mathcal{B}})=f^{\otimes n}
\tag{6.5.3.2}
\]
lorsque $\mathcal{B}$ est de rang constant $n$ dans l'hypothèse (I) (resp.
lorsque $\mathcal{B}\otimes_{\mathcal{A}}\mathcal{R}(X)$ est de rang constant
$n$ dans l'hypothèse (II)). Enfin, pour que l'homomorphisme
$\mathcal{B}\to\mathcal{L}'$ correspondant à une section $f'$ de
$\mathcal{L}'$ au-dessus de $X$ soit un isomorphisme, il faut et il suffit que
$f'_x$ soit une base de $\mathcal{L}'_x$ pour tout $x\in X$~; dans l'hypothèse
(I), cette condition équivaut donc à dire que $(N(f'))_x$ est une base de
$(N(\mathcal{L}'))_x$ pour tout $x$~; dans l'hypothèse (II), cette condition
est encore nécessaire, et elle est suffisante lorsque $\mathcal{B}$ vérifie
(II \emph{bis}) et que
$\mathcal{B}\to\mathcal{B}\otimes_{\mathcal{A}}\mathcal{R}(X)$ est injectif.
\end{env}

\begin{env}[6.5.4]
\label{II.6.5.4-fr}
Soient $(X,\mathcal{A})$, $(X',\mathcal{A}')$ deux espaces annelés,
$f:X'\to X$ un morphisme, $\mathcal{B}$ une $\mathcal{A}$-Algèbre,
$\mathcal{B}'=f^*(\mathcal{B})$ la $\mathcal{A}'$-Algèbre image réciproque. On
suppose vérifiée l'une des hypothèses suivantes~:
\begin{enumerate}
\item[$1^\circ$] $\mathcal{B}$ vérifie l'hypothèse (I) de
(\hyperref[II.6.5.1-fr]{6.5.1}).
\item[$2^\circ$] $(X,\mathcal{A})$ et $\mathcal{B}$ vérifient l'hypothèse
(II) de (\hyperref[II.6.5.1-fr]{6.5.1}), $(X',\mathcal{A}')$ est un préschéma
localement noethérien réduit, et si l'on désigne par $X_\alpha$ et
$X'_\beta$ les sous-préschémas fermés réduits respectifs de $X$ et $X'$ ayant
pour espaces sous-jacents les composantes irréductibles de ces espaces, la
restriction de $f$ à chaque $X'_\beta$ est un morphisme \emph{dominant} de
$X'_\beta$ dans un $X_\alpha$.
\end{enumerate}
Sous ces conditions, $\mathcal{B}'$ vérifie respectivement l'hypothèse (I) ou
l'hypothèse (II) de (\hyperref[II.6.5.1-fr]{6.5.1})~; la première assertion
est immédiate. Pour établir la seconde, il suffit de voir que pour tout
$x'\in X'$, les rangs des
$\mathcal{B}'\otimes_{\mathcal{O}_{X'}}\mathcal{R}(X'_\beta)$ pour les $\beta$
tels que $x'\in X'_\beta$ sont \emph{les mêmes}. Or, si la restriction de $f$
à $X'_\beta$ est un morphisme dominant dans $X_\alpha$, le rang de
$\mathcal{B}'\otimes_{\mathcal{O}_{X'}}\mathcal{R}(X'_\beta)$ est égal à celui
de $\mathcal{B}\otimes_{\mathcal{O}_X}\mathcal{R}(X_\alpha)$ (on le voit
aussitôt en se ramenant au cas affine comme dans
(\hyperref[II.6.4.9-fr]{6.4.9})), d'où notre assertion en vertu de l'hypothèse
(II) et de (\hyperref[II.6.5.1-fr]{6.5.1}).

\oldpage[II]{129}

Cela étant, il résulte de (\hyperref[II.6.4.8-fr]{6.4.8}),
(\hyperref[II.6.4.9-fr]{6.4.9}) et (\hyperref[II.6.4.10-fr]{6.4.10}) que si
$s$ est une section de $\mathcal{B}$ au-dessus d'un ouvert $U\subset X$,
$s'$ la section correspondante de $\mathcal{B}'$ au-dessus de $f^{-1}(U)$,
$N_{\mathcal{B}'/\mathcal{A}'}(s')$ est la section de $\mathcal{A}'$ au-dessus
de $f^{-1}(U)$ qui correspond à la section
$N_{\mathcal{B}/\mathcal{A}}(s)$ de $\mathcal{A}$ au-dessus de $U$.

Si $\mathcal{M}$ est un $\mathcal{B}$-Module inversible, on déduit de ce qui
précède que si $\mathcal{M}'=f^*(\mathcal{M})$ (qui est un
$\mathcal{B}'$-Module inversible), on a
$N_{\mathcal{B}'/\mathcal{A}'}(\mathcal{M}')
=f^*(N_{\mathcal{B}/\mathcal{A}}(\mathcal{M}))$ à un isomorphisme canonique
près.
\end{env}

\begin{env}[6.5.5]
\label{II.6.5.5-fr}
Supposons désormais que $(X,\mathcal{A})$ soit un \emph{préschéma}. La donnée
d'une $\mathcal{A}$-Algèbre quasi-cohérente $\mathcal{B}$, qui est un
$\mathcal{A}$-Module \emph{de type fini} équivaut alors, comme on sait, à
celle d'un morphisme \emph{fini} $g:X'\to X$ tel que
$g_*(\mathcal{O}_{X'})=\mathcal{B}$, défini à un $X$-isomorphisme près
(\hyperref[II.6.1.2-fr]{6.1.2} et \hyperref[II.1.3.1-fr]{1.3.1}). En outre,
se donner un $\mathcal{O}_{X'}$-Module quasi-cohérent $\mathcal{F}'$ équivaut
à se donner un $\mathcal{B}$-Module quasi-cohérent $\mathcal{F}$ tel que
$g_*(\mathcal{F}')=\mathcal{F}$ (\hyperref[II.1.4.3-fr]{1.4.3}), et pour que
$\mathcal{F}'$ soit inversible, il faut et il suffit que $\mathcal{F}$ le soit
(\hyperref[II.6.1.12-fr]{6.1.12}). Pour traduire les résultats précédents en
termes du morphisme fini $g$, il faudra donc supposer, soit que
$g_*(\mathcal{O}_{X'})$ est un $\mathcal{O}_X$-Module \emph{localement libre}
(de type fini), soit que $(X,\mathcal{O}_X)$ et $g_*(\mathcal{O}_{X'})$
vérifient l'hypothèse (II). Pour tout $\mathcal{O}_{X'}$-Module inversible
$\mathcal{L}'$, on posera alors
\[
\label{II.6.5.5.1-fr}
N_{X'/X}(\mathcal{L}')
=N_{g_*(\mathcal{O}_{X'})/\mathcal{O}_X}(g_*(\mathcal{L}'))
\tag{6.5.5.1}
\]
et on dira que c'est la \emph{norme} (relative à $g$) de $\mathcal{L}'$. De
même, si $h':\mathcal{L}'_1\to\mathcal{L}'_2$ est un homomorphisme de
$\mathcal{O}_{X'}$-Modules inversibles, on posera
\[
\label{II.6.5.5.2-fr}
N_{X'/X}(h')
=N_{g_*(\mathcal{O}_{X'})/\mathcal{O}_X}(g_*(h')):
N_{X'/X}(\mathcal{L}'_1)\to N_{X'/X}(\mathcal{L}'_2).
\tag{6.5.5.2}
\]
En particulier, pour $\mathcal{L}'_1=\mathcal{O}_{X'}$, on obtient ainsi une
application canonique
\[
\label{II.6.5.5.3-fr}
N_{X'/X}:\Gamma(X',\mathcal{L}')
\to\Gamma(X,N_{X'/X}(\mathcal{L}')).
\tag{6.5.5.3}
\]
Nous laisserons au lecteur la plupart de ces traductions, et nous bornerons à
expliciter les suivantes~:
\end{env}

\begin{proposition}[6.5.6]
\label{II.6.5.6-fr}
Soit $g:X'\to X$ un morphisme fini, et supposons, soit que
$g_*(\mathcal{O}_{X'})$ est un $\mathcal{O}_X$-Module localement libre, soit
que $(X,\mathcal{O}_X)$ et $g_*(\mathcal{O}_{X'})$ vérifient (II \emph{bis})
(ce qui sera le cas en particulier lorsque $X$ est localement noethérien et
normal). Pour qu'un homomorphisme $h':\mathcal{L}'_1\to\mathcal{L}'_2$ de
$\mathcal{O}_{X'}$-Modules inversibles soit un isomorphisme, il faut et il
suffit, dans la première hypothèse, que $N_{X'/X}(h')$ soit un isomorphisme~;
dans la seconde hypothèse, cette condition est encore nécessaire, et elle est
suffisante lorsque l'homomorphisme
$g_*(\mathcal{O}_{X'})\to
g_*(\mathcal{O}_{X'})\otimes_{\mathcal{O}_X}\mathcal{R}(X)$ est injectif.
\end{proposition}

On notera qu'on utilise ici le fait que, pour que $g_*(h')$ soit un
isomorphisme, il faut et il suffit que $h'$ en soit un
(\hyperref[II.1.4.2-fr]{1.4.2}).

\begin{corollary}[6.5.7]
\label{II.6.5.7-fr}
Soit $g:X'\to X$ un morphisme fini, et supposons, soit que
$g_*(\mathcal{O}_{X'})$ est un $\mathcal{O}_X$-Module localement libre, soit
que $(X,\mathcal{O}_X)$ et $g_*(\mathcal{O}_{X'})$ vérifient (II \emph{bis})
et que
$g_*(\mathcal{O}_{X'})\to
g_*(\mathcal{O}_{X'})\otimes_{\mathcal{O}_X}\mathcal{R}(X)$ est injectif.
Soient $\mathcal{L}'$ un $\mathcal{O}_{X'}$-Module inversible, $f'$ une
section de $\mathcal{L}'$ au-dessus de $X'$, $f=N_{X'/X}(f')$ la section de
$\mathcal{L}=N_{X'/X}(\mathcal{L}')$ au-dessus de $X$ qui lui correspond
(\hyperref[II.6.5.5.3-fr]{6.5.5.3}).

\oldpage[II]{130}

Alors on a $g(X'-X'_{f'})=X-X_f$ et $X_f$ est le plus grand ouvert $U$ de
$X$ tel que $g^{-1}(U)\subset X'_{f'}$.
\end{corollary}

En effet, $g(X'-X'_{f'})$ est fermé dans $X$
(\hyperref[II.6.1.10-fr]{6.1.10})~; il suffit donc de prouver la dernière
assertion. Or la relation $U\subset X_f$ équivaut au fait que l'homomorphisme
$\mathcal{O}_X|U\to\mathcal{L}|U$ défini par $f|U$ est un isomorphisme. En
vertu de (\hyperref[II.6.5.6-fr]{6.5.6}), cela équivaut à dire que
l'homomorphisme
$\mathcal{O}_{X'}|g^{-1}(U)\to\mathcal{L}'|g^{-1}(U)$ défini par
$f'|g^{-1}(U)$ est un isomorphisme, c'est-à-dire à la relation
$g^{-1}(U)\subset X'_{f'}$.

\begin{proposition}[6.5.8]
\label{II.6.5.8-fr}
Soient $g:X'\to X$ un morphisme fini, $f:Y\to X$ un morphisme~; soient
$Y'=X'_{(Y)}$, $g'=g_{(Y)}$, $f'=f_{(X')}$ de sorte que l'on a le diagramme
commutatif
\[
\xymatrix{
X'\ar[d]_g & Y'\ar[l]_{f'}\ar[d]^{g'}\\
X & Y\ar[l]^f
}
\]
On suppose, soit que $g_*(\mathcal{O}_{X'})$ est localement libre, soit que
$(X,\mathcal{O}_X)$ et $g_*(\mathcal{O}_{X'})$ vérifient (II), que $Y$ est un
préschéma localement noethérien réduit, et que la restriction de $f$ à toute
composante irréductible de $Y$ est un morphisme dominant dans une composante
irréductible de $X$. Alors, pour tout $\mathcal{O}_{X'}$-Module inversible
$\mathcal{L}'$, on a
\[
N_{Y'/Y}({f'}^*(\mathcal{L}'))=f^*(N_{X'/X}(\mathcal{L}'))
\]
à un isomorphisme canonique près.
\end{proposition}

Notons que l'on a $f^*(g_*(\mathcal{L}'))=g'_*({f'}^*(\mathcal{L}'))$ en
vertu de (\hyperref[II.1.5.2-fr]{1.5.2}), et en particulier
$g'_*(\mathcal{O}_{Y'})=f^*(g_*(\mathcal{O}_{X'}))$~; si
$g_*(\mathcal{O}_{X'})$ est localement libre, il en est donc de même de
$g'_*(\mathcal{O}_{Y'})$. La conclusion résulte alors des définitions et de
(\hyperref[II.6.5.4-fr]{6.5.4}).

\begin{remark}[6.5.9]
\label{II.6.5.9-fr}
Nous généraliserons ultérieurement la notion de norme développée ci-dessus,
et la mettrons en relation avec la notion d'image directe d'un diviseur.
\end{remark}

\subsection{Application : critères d'amplitude.}
\label{subsection:II.6.6-fr}

\begin{proposition}[6.6.1]
\label{II.6.6.1-fr}
Soient $Y$ un préschéma, $f:X\to Y$ un morphisme quasi-compact,
$g:X'\to X$ un morphisme fini et surjectif. On suppose, soit que
$g_*(\mathcal{O}_{X'})$ est un $\mathcal{O}_X$-Module localement libre, soit
que $(X,\mathcal{O}_X)$ et $g_*(\mathcal{O}_{X'})$ vérifient (II
\emph{bis}). Alors, pour tout $\mathcal{O}_{X'}$-Module inversible
$\mathcal{L}'$, ample pour $f\circ g$, $N_{X'/X}(\mathcal{L}')=\mathcal{L}$
est ample pour $f$.
\end{proposition}

On peut supposer $Y$ affine (\hyperref[II.4.6.4-fr]{4.6.4}), et alors, en
vertu de (\hyperref[II.4.6.6-fr]{4.6.6}), l'énoncé est équivalent au

\begin{corollary}[6.6.2]
\label{II.6.6.2-fr}
Soient $X$ un préschéma quasi-compact, $g:X'\to X$ un morphisme fini
surjectif, tel que $g_*(\mathcal{O}_{X'})$ soit un $\mathcal{O}_X$-Module
localement libre, ou que $(X,\mathcal{O}_X)$ et $g_*(\mathcal{O}_{X'})$
vérifient (II \emph{bis}). Alors, pour tout $\mathcal{O}_{X'}$-Module ample
$\mathcal{L}'$, $\mathcal{L}=N_{X'/X}(\mathcal{L}')$ est ample.
\end{corollary}

Dans la seconde hypothèse, on peut supposer en outre que l'homomorphisme
canonique
$g_*(\mathcal{O}_{X'})\to
g_*(\mathcal{O}_{X'})\otimes_{\mathcal{O}_X}\mathcal{R}(X)$ est
\emph{injectif}. En effet, dans le cas contraire, soit $\mathcal{T}$

\oldpage[II]{131}

le noyau de cet homomorphisme, qui est un Idéal cohérent de
$g_*(\mathcal{O}_{X'})=\mathcal{B}$
(\hyperref[I.6.1.1-fr]{I, 6.1.1}), et posons
$X''=\operatorname{Spec}(\mathcal{B}/\mathcal{T})$~; on a donc un diagramme
commutatif
\[
\xymatrix{
X''\ar[rr]^h\ar[dr]_{g'} && X'\ar[dl]^g\\
& X
}
\]
où $h$ est une immersion fermée (\hyperref[II.1.4.10-fr]{1.4.10}). En outre,
on sait que le support de $\mathcal{T}$ est un ensemble fermé
(\hyperref[0.5.2.2-fr]{0, 5.2.2}) rare dans $X$
(\hyperref[I.7.4.6-fr]{I, 7.4.6}), d'où on conclut que pour le point générique
$x$ d'une composante irréductible de $X$, il y a un voisinage ouvert affine
$U$ de $x$ tel que
$\mathcal{B}|U=(\mathcal{B}/\mathcal{T})|U$. Comme $g$ est par hypothèse
surjectif, on en conclut que $x\in g'(X'')$~; $g'$ est donc dominant, et étant
un morphisme fini, il est \emph{surjectif}
(\hyperref[II.6.1.10-fr]{6.1.10})~; on a par définition
\[
g'_*(\mathcal{O}_{X''})\otimes_{\mathcal{O}_X}\mathcal{R}(X)
=(\mathcal{B}/\mathcal{T})\otimes_{\mathcal{O}_X}\mathcal{R}(X)
=g_*(\mathcal{O}_{X'})\otimes_{\mathcal{O}_X}\mathcal{R}(X),
\]
donc $(X,\mathcal{O}_X)$ et $g'_*(\mathcal{O}_{X''})$ vérifient
(II \emph{bis}), et en outre
\[
g'_*(\mathcal{O}_{X''})\longrightarrow
g'_*(\mathcal{O}_{X''})\otimes_{\mathcal{O}_X}\mathcal{R}(X)
\]
est injectif. Enfin, $h^*(\mathcal{L}')=\mathcal{L}''$ est un
$\mathcal{O}_{X''}$-Module ample
(\hyperref[II.4.6.13-fr]{4.6.13, (i \emph{bis})}), et on a
$N_{X''/X}(\mathcal{L}'')=N_{X'/X}(\mathcal{L}')$. En effet, pour définir
ces deux $\mathcal{O}_X$-Modules inversibles on peut utiliser un même
recouvrement ouvert affine $(U_\lambda)$ de $X$ tel que les restrictions de
$g_*(\mathcal{L}')$ et $g'_*(\mathcal{L}'')$ à $U_\lambda$ soient
respectivement isomorphes à $\mathcal{B}|U_\lambda$ et
$(\mathcal{B}/\mathcal{T})|U_\lambda$. On voit aussitôt qu'à tout
isomorphisme
\[
\eta_\lambda:g_*(\mathcal{L}')|U_\lambda\longrightarrow
\mathcal{B}|U_\lambda
\]
correspond canoniquement un isomorphisme
\[
\eta'_\lambda:g'_*(\mathcal{L}'')|U_\lambda\longrightarrow
(\mathcal{B}/\mathcal{T})|U_\lambda,
\]
de sorte que, si $(\omega_{\lambda\mu})$ et
$(\omega'_{\lambda\mu})$ sont les $1$-cocycles correspondant aux systèmes
d'isomorphismes $(\eta_\lambda)$ et $(\eta'_\lambda)$
(\hyperref[II.6.5.2-fr]{6.5.2}), $\omega'_{\lambda\mu}$ soit l'image
canonique dans
$\Gamma(U_\lambda\cap U_\mu,\mathcal{B}/\mathcal{T})$ de
$\omega_{\lambda\mu}\in\Gamma(U_\lambda\cap U_\mu,\mathcal{B})$. En vertu
de la définition de $\mathcal{T}$, on en conclut que
\[
N_{\mathcal{B}/\mathcal{A}}\omega_{\lambda\mu}
=N_{(\mathcal{B}/\mathcal{T})/\mathcal{A}}\omega'_{\lambda\mu}
\]
(avec $\mathcal{A}=\mathcal{O}_X$), d'où l'égalité énoncée.

Supposons donc l'homomorphisme
$g_*(\mathcal{O}_{X'})\to
g_*(\mathcal{O}_{X'})\otimes_{\mathcal{O}_X}\mathcal{R}(X)$ injectif
lorsqu'on est dans l'hypothèse (II \emph{bis}). Il suffit de prouver que
lorsque $f$ parcourt les sections de $\mathcal{L}^{\otimes n}$ ($n>0$)
au-dessus de $X$, les $X_f$ forment une base de la topologie de $X$
(\hyperref[II.4.5.2-fr]{4.5.2}). Or, soit $x\in X$, et soit $U$ un voisinage
quelconque de $x$~; comme $g^{-1}(x)$ est fini
(\hyperref[II.6.1.7-fr]{6.1.7}) et que $\mathcal{L}'$ est ample, il existe un
entier $n>0$ et une section $f'$ de ${\mathcal{L}'}^{\otimes n}$ au-dessus de
$X'$, tels que $X'_{f'}$ soit un voisinage de $g^{-1}(x)$ contenu dans
$g^{-1}(U)$ (\hyperref[II.4.5.4-fr]{4.5.4}). Comme
\[
\mathcal{L}^{\otimes n}=N_{X'/X}({\mathcal{L}'}^{\otimes n}),
\]
il suffit de prendre $f=N_{X'/X}(f')$~; en effet, alors
$X-X_f=g(X'-X'_{f'})$ (\hyperref[II.6.5.7-fr]{6.5.7}), donc
$x\in X_f\subset U$.

\begin{corollary}[6.6.3]
\label{II.6.6.3-fr}
Sous les hypothèses de (\hyperref[II.6.6.1-fr]{6.6.1}), pour qu'un
$\mathcal{O}_X$-Module inversible $\mathcal{L}$ soit ample pour $f$, il faut
et il suffit que $\mathcal{L}'=g^*(\mathcal{L})$ soit ample pour $f\circ g$.
\end{corollary}

La condition est nécessaire, puisque $g$ est affine
(\hyperref[II.5.1.12-fr]{5.1.12}). Pour démontrer que la condition est
suffisante, on peut supposer $Y$ affine (\hyperref[II.4.6.4-fr]{4.6.4}), donc
$X$ et $X'$ quasi-compacts et $\mathcal{L}'$ ample
(\hyperref[II.4.6.6-fr]{4.6.6}), et il faut montrer que $\mathcal{L}$ est
ample. Or, l'ensemble des points

\oldpage[II]{132}

$x\in X$ au voisinage desquels $g_*(\mathcal{O}_{X'})$ (resp.
$g_*(\mathcal{O}_{X'})\otimes_{\mathcal{O}_X}\mathcal{R}(X)$) a un rang
donné $n$ dans la première (resp. la seconde) hypothèse, est à la fois ouvert
et fermé dans $X$, donc $X$ est le préschéma somme d'un nombre fini de ces
ouverts, et on peut par suite supposer qu'il est égal à l'un d'eux
(\hyperref[II.4.6.17-fr]{4.6.17}). Mais on a alors
$N_{X'/X}(\mathcal{L}')=\mathcal{L}^{\otimes n}$, donc
$\mathcal{L}^{\otimes n}$ est ample en vertu de
(\hyperref[II.6.6.2-fr]{6.6.2}), et il en est donc de même de $\mathcal{L}$
(\hyperref[II.4.5.6-fr]{4.5.6}).

\begin{corollary}[6.6.4]
\label{II.6.6.4-fr}
Supposons vérifiées les hypothèses de
(\hyperref[II.6.6.1-fr]{6.6.1}) et supposons en outre que
$f:X\to Y$ soit de type fini. Alors, pour que $f$ soit quasi-projectif, il
faut et il suffit que $f\circ g$ le soit. Si on suppose de plus que $Y$ est
un schéma quasi-compact ou un préschéma dont l'espace sous-jacent est
noethérien, alors, pour que $f$ soit projectif, il faut et il suffit que
$f\circ g$ le soit.
\end{corollary}

L'hypothèse entraîne que $f\circ g$ est de type fini. Compte tenu de la
définition des morphismes quasi-projectifs
(\hyperref[II.5.3.1-fr]{5.3.1}), la première assertion résulte de
(\hyperref[II.6.6.1-fr]{6.6.1}) et
(\hyperref[II.6.6.3-fr]{6.6.3}). Compte tenu de ce résultat et de
(\hyperref[II.5.5.3-fr]{5.5.3, (ii)}), il reste à prouver que lorsque $f$ est
supposé quasi-projectif, pour que $f$ soit propre, il faut et il suffit que
$f\circ g$ le soit. Or $f$ est alors séparé
(\hyperref[II.5.3.1-fr]{5.3.1}) et de type fini~; comme $g$ est surjectif,
notre assertion résulte de (\hyperref[II.5.4.2-fr]{5.4.2, (ii)}) et
(\hyperref[II.5.4.3-fr]{5.4.3, (ii)}).

En particulier~:

\begin{corollary}[6.6.5]
\label{II.6.6.5-fr}
Soient $X$ un préschéma de type fini sur un corps $K$, $K'$ une extension de
degré fini de $K$. Pour que $X$ soit projectif (resp. quasi-projectif) sur
$K$, il faut et il suffit que $X'=X\otimes_K K'$ soit projectif (resp.
quasi-projectif) sur $K'$.
\end{corollary}

La condition est en effet nécessaire
(\hyperref[II.5.3.4-fr]{5.3.4, (iii)} et
\hyperref[II.5.5.5-fr]{5.5.5, (iii)}). Inversement, supposons-la vérifiée,
et soit $g:X'\to X$ la projection canonique. Il est clair que $g$ est un
morphisme fini (\hyperref[II.6.1.5-fr]{6.1.5, (iii)}) et surjectif
(\hyperref[I.3.5.2-fr]{I, 3.5.2, (ii)}). En outre,
$g_*(\mathcal{O}_{X'})$ est un $\mathcal{O}_X$-Module localement libre, étant
isomorphe à $\mathcal{O}_X\otimes_K K'$
(\hyperref[II.1.5.2-fr]{1.5.2}). Il résulte alors de l'hypothèse et de
(\hyperref[II.6.1.11-fr]{6.1.11}) et
(\hyperref[II.5.5.5-fr]{5.5.5, (ii)}) que $X'$ est projectif (resp.
quasi-projectif) \emph{sur $K$}~; on déduit alors de
(\hyperref[II.6.6.4-fr]{6.6.4}) que $X$ est projectif (resp.
quasi-projectif) sur $K$.

Au chapitre V, nous montrerons que l'énoncé
(\hyperref[II.6.6.5-fr]{6.6.5}) reste valable lorsque $K'$ est une extension
\emph{quelconque} de $K$.

La fin de ce numéro est consacrée à la démonstration du critère
(\hyperref[II.6.6.11-fr]{6.6.11}), qui est un raffinement assez technique de
(\hyperref[II.6.6.1-fr]{6.6.1})~; elle peut être omise en première lecture.

\begin{lemma}[6.6.6]
\label{II.6.6.6-fr}
Soient $X$ un préschéma noethérien réduit, $\mathcal{E}$ un
$\mathcal{O}_X$-Module cohérent tel que
$\mathcal{E}\otimes_{\mathcal{O}_X}\mathcal{R}(X)$ soit un
$\mathcal{R}(X)$-Module localement libre de rang constant $n$. Il existe
alors un préschéma noethérien réduit $Z$ et un morphisme fini birationnel
$h:Z\to X$ ayant la propriété suivante~: les morphismes de faisceaux
d'ensembles
$\sigma_i:\mathcal{H}om_{\mathcal{O}_X}(\mathcal{E},\mathcal{E})
\to\mathcal{R}(X)$ ($1\leq i\leq n$)
(cf. \hyperref[II.6.4.10-fr]{6.4.10}) appliquent
$\mathcal{H}om_{\mathcal{O}_X}(\mathcal{E},\mathcal{E})$ dans la
$\mathcal{O}_X$-Algèbre cohérente $h_*(\mathcal{O}_Z)$.
\end{lemma}

Considérons en effet un ouvert affine $U$ de $X$, d'anneau $A(U)=A$~; soit
$E=\Gamma(U,\mathcal{E})$, et soit $C_U$ la sous-algèbre de $R(U)$ engendrée
par les $\sigma_i(u)$ lorsque $u$ parcourt $\operatorname{Hom}_A(E,E)$~; on
a vu (\hyperref[II.6.4.7.1-fr]{6.4.7.1}) que cette $A$-algèbre est de
\emph{rang fini}. En outre, il est clair que la formation des algèbres $C_U$
commute avec les opérations de restriction d'un ouvert affine $U$ à un ouvert
affine $U'\subset U$. On a donc défini ainsi une sous-$\mathcal{O}_X$-Algèbre
\emph{finie} $\mathcal{C}$ de $\mathcal{R}(X)$ telle que
$\Gamma(U,\mathcal{C})=C_U$ pour tout ouvert affine $U$ de $X$. On

\oldpage[II]{133}

prendra $Z=\operatorname{Spec}(\mathcal{C})$, et pour $h$ le morphisme
structural, qui est donc fini (\hyperref[II.6.1.2-fr]{6.1.2})~; comme
$\mathcal{C}$ est réduite, $Z$ est un préschéma noethérien réduit
(\hyperref[II.1.3.8-fr]{1.3.8}). Enfin, l'anneau total des fractions de $C_U$
est $R(U)$ par définition, et comme $C_U$ est contenu dans la fermeture
intégrale de $A(U)$ dans $R(U)$, il y a correspondance biunivoque entre idéaux
premiers minimaux de $A(U)$ et idéaux premiers minimaux de $C_U$
(\cite[t.~I, p.~259]{I-13}), ce qui prouve que $h$ est birationnel et achève
la démonstration.

\begin{corollary}[6.6.7]
\label{II.6.6.7-fr}
Sous les hypothèses de (\hyperref[II.6.6.6-fr]{6.6.6}), soit $W$ un ouvert
de $X$ tel que, pour tout $x\in W$, ou bien $X$ est normal au point $x$, ou
bien $\mathcal{E}_x$ est un $\mathcal{O}_x$-Module libre. Alors on peut
supposer $h$ défini de sorte que la restriction de $h$ à $h^{-1}(W)$ soit un
isomorphisme de $h^{-1}(W)$ sur $W$.
\end{corollary}

En effet, l'une ou l'autre hypothèse entraîne que si $U\subset W$ est un
ouvert affine, on a, avec les notations de
(\hyperref[II.6.6.6-fr]{6.6.6}), $(\sigma_i(u))_x\in A_x$ pour tout
$x\in U$ (\hyperref[II.6.4.3-fr]{6.4.3}), donc $\sigma_i(u)\in A$, et la
conclusion résulte de la définition de $h$ donnée dans
(\hyperref[II.6.6.6-fr]{6.6.6}).

\begin{env}[6.6.8]
\label{II.6.6.8-fr}
Soient $X$ un préschéma noethérien réduit, $g:X'\to X$ un morphisme fini
surjectif, de sorte que $\mathcal{B}=g_*(\mathcal{O}_{X'})$ est une
$\mathcal{O}_X$-Algèbre \emph{cohérente}~; supposons en outre que
$\mathcal{B}\otimes_{\mathcal{O}_X}\mathcal{R}(X)$ soit un
$\mathcal{R}(X)$-Module \emph{localement libre de rang constant $n$}. On peut
alors appliquer le lemme (\hyperref[II.6.6.6-fr]{6.6.6}) en prenant
$\mathcal{E}=\mathcal{B}$, d'où, avec les notations de
(\hyperref[II.6.6.6-fr]{6.6.6}), un homomorphisme de faisceaux de monoïdes
multiplicatifs
$\sigma_n:\mathcal{H}om_{\mathcal{O}_X}(\mathcal{B},\mathcal{B})
\to h_*(\mathcal{O}_Z)$, et en composant cet homomorphisme avec
l'homomorphisme canonique
$\mathcal{B}\to
\mathcal{H}om_{\mathcal{O}_X}(\mathcal{B},\mathcal{B})$
(\hyperref[II.6.5.1-fr]{6.5.1}), on obtient donc un homomorphisme de
faisceaux de monoïdes multiplicatifs~:
\[
\label{II.6.6.8.1-fr}
N':\mathcal{B}=g_*(\mathcal{O}_{X'})\longrightarrow
h_*(\mathcal{O}_Z)=\mathcal{C}.
\tag{6.6.8.1}
\]

Cela étant, pour tout $\mathcal{O}_{X'}$-Module inversible $\mathcal{L}'$,
$g_*(\mathcal{L}')$ est un $\mathcal{B}$-Module inversible
(\hyperref[II.6.1.12-fr]{6.1.12}), et la méthode de
(\hyperref[II.6.5.2-fr]{6.5.2}) permet d'associer fonctoriellement à
$\mathcal{L}'$ un $\mathcal{C}$-Module inversible que nous noterons
$N'(g_*(\mathcal{L}'))$.
\end{env}

\begin{lemma}[6.6.9]
\label{II.6.6.9-fr}
Soient $X$ un préschéma noethérien réduit, $g:X'\to X$ un morphisme fini
surjectif tel que
$g_*(\mathcal{O}_{X'})\otimes_{\mathcal{O}_X}\mathcal{R}(X)$ soit un
$\mathcal{R}(X)$-Module localement libre de rang constant. Il existe alors
un préschéma noethérien réduit $Z$ et un morphisme fini birationnel
$h:Z\to X$ ayant la propriété suivante~: pour tout
$\mathcal{O}_{X'}$-Module ample $\mathcal{L}'$, le $\mathcal{O}_Z$-Module
inversible $\mathcal{M}$ tel que
$h_*(\mathcal{M})=N'(g_*(\mathcal{L}'))$ (notations de
\hyperref[II.6.6.8-fr]{6.6.8}) est ample.
\end{lemma}

Supposons d'abord que l'homomorphisme
$\mathcal{B}\to
\mathcal{B}\otimes_{\mathcal{O}_X}\mathcal{R}(X)$ soit \emph{injectif}.
Définissons $Z$ et $h$ comme dans (\hyperref[II.6.6.6-fr]{6.6.6}) (avec
$\mathcal{E}=g_*(\mathcal{O}_{X'})$). Soit $z\in Z$~; il faut montrer qu'il
existe un entier $m>0$ et une section $t$ de $\mathcal{M}^{\otimes m}$
au-dessus de $Z$ telle que $Z_t$ soit un ouvert affine contenant $z$
(\hyperref[II.4.5.2-fr]{4.5.2}). Soient $x=h(z)$, $U$ un ouvert affine de
$X$ contenant $x$~; $h^{-1}(U)$ est alors un voisinage ouvert affine de $z$,
et il suffira de trouver $t$ tel que $z\in Z_t\subset h^{-1}(U)$, car $Z_t$
sera alors nécessairement affine (\hyperref[II.5.5.8-fr]{5.5.8}). Il existe
par hypothèse un entier $n>0$ et une section $s'$ de
${\mathcal{L}'}^{\otimes n}$ au-dessus de $X'$ telle que l'on ait
\[
\label{II.6.6.9.1-fr}
g^{-1}(x)\subset X'_{s'}\subset g^{-1}(U)
\tag{6.6.9.1}
\]
en vertu de (\hyperref[II.4.5.4-fr]{4.5.4}). Par définition, $s'$ est aussi
une section de $g_*({\mathcal{L}'}^{\otimes n})$ au-dessus de $X$, et il lui
correspond comme dans (\hyperref[II.6.5.2-fr]{6.5.2}) une section
$s=N'(s')$ de $N'(g_*({\mathcal{L}'}^{\otimes n}))$ au-dessus

\oldpage[II]{134}

de $X$. Nous allons voir que si $t$ est la section $s$ considérée comme
section de $\mathcal{M}^{\otimes n}$ au-dessus de $Z$, $t$ répond à la
question. Posons en effet
\[
\label{II.6.6.9.2-fr}
V=X-g(X'-X'_{s'})
\tag{6.6.9.2}
\]
qui est un ouvert de $X$ contenant $x$ et contenu dans $U$, en vertu de
(\hyperref[II.6.6.9.1-fr]{6.6.9.1}) et
(\hyperref[II.6.1.10-fr]{6.1.10}). Nous allons montrer que l'on a
\[
\label{II.6.6.9.3-fr}
h^{-1}(V)\subset Z_t\subset h^{-1}(U),
\tag{6.6.9.3}
\]
ce qui achèvera la démonstration. Il revient au même de dire que l'ensemble
$T$ des $y\in X$ tels que $s_y$ soit inversible, contient $V$ et est contenu
dans $U$. Pour cela, considérons d'abord un ouvert affine $W$ contenu dans
$V$~; $g^{-1}(W)$ est ouvert affine dans $X'$ et en vertu de
(\hyperref[II.6.6.9.2-fr]{6.6.9.2}), $s'_{y'}$ est inversible pour tout
$y'\in g^{-1}(W)$~; en vertu des hypothèses sur $X$ et $\mathcal{B}$, on
peut appliquer les résultats de (\hyperref[II.6.4.7-fr]{6.4.7}) et on voit
que si $y=g(y')$, $s_y$ est inversible, autrement dit $V\subset T$. D'autre
part, il résulte aussi de (\hyperref[II.6.4.7-fr]{6.4.7}) que, inversement,
si $s_y$ est inversible, il en est de même de $s'_{y'}$, ce qui implique
$y'\in g^{-1}(U)$ par (\hyperref[II.6.6.9.1-fr]{6.6.9.1}) et par suite
$y\in U$, d'où $T\subset U$ dans ce cas.

On passe de là au cas général par le même raisonnement que dans
(\hyperref[II.6.6.2-fr]{6.6.2}), remplaçant $X'$ par $X''$ tel que
\[
g'_*(\mathcal{O}_{X''})\longrightarrow
g'_*(\mathcal{O}_{X''})\otimes_{\mathcal{O}_X}\mathcal{R}(X)
\]
soit injectif, et $\mathcal{L}'$ par un $\mathcal{O}_{X''}$-Module ample
$\mathcal{L}''$ tel que
$N'(g_*(\mathcal{L}'))=N'(g'_*(\mathcal{L}''))$. Le lemme
(\hyperref[II.6.6.9-fr]{6.6.9}) est donc démontré dans tous les cas (avec un
choix convenable de $h$).

\begin{corollary}[6.6.10]
\label{II.6.6.10-fr}
Supposons vérifiées les hypothèses de
(\hyperref[II.6.6.9-fr]{6.6.9})~; pour tout $\mathcal{O}_X$-Module inversible
$\mathcal{L}$ tel que $g^*(\mathcal{L})$ soit ample, $h^*(\mathcal{L})$ est
ample.
\end{corollary}

En effet, si on pose $\mathcal{L}'=g^*(\mathcal{L})$, on a alors
$g_*(\mathcal{L}')=\mathcal{L}\otimes_{\mathcal{O}_X}\mathcal{B}$
(\hyperref[0.5.4.10-fr]{0, 5.4.10}), donc
\[
N'(g_*(\mathcal{L}'))=
(\mathcal{L}\otimes_{\mathcal{O}_X}\mathcal{C})^{\otimes n}
\]
(par le même raisonnement que pour
\hyperref[II.6.5.2.4-fr]{6.5.2.4}). On en conclut que l'on a
$\mathcal{M}=(h^*(\mathcal{L}))^{\otimes n}$, et comme $\mathcal{M}$ est
ample, il en est de même de $h^*(\mathcal{L})$
(\hyperref[II.4.5.6-fr]{4.5.6}).

\begin{proposition}[6.6.11]
\label{II.6.6.11-fr}
Soient $Y$ un schéma affine, $X$ un préschéma noethérien réduit,
$f:X\to Y$ un morphisme quasi-compact, $g:X'\to X$ un morphisme fini
surjectif. Soit $W$ une partie ouverte de $X$ telle que, pour tout $x\in W$,
ou bien $X$ est normal au point $x$, ou bien il existe un voisinage ouvert
$T\subset W$ de $x$ tel que $(g_*(\mathcal{O}_{X'}))|T$ soit un
$(\mathcal{O}_X|T)$-Module localement libre. Il existe alors un
$Y$-préschéma réduit $Z$ et un $Y$-morphisme fini et birationnel
$h:Z\to X$ tel que la restriction de $h$ à $h^{-1}(W)$ soit un isomorphisme
$h^{-1}(W)\xrightarrow{\sim}W$ et qui possède la propriété suivante~: pour
tout $\mathcal{O}_X$-Module inversible $\mathcal{L}$ tel que
$g^*(\mathcal{L})$ soit ample relativement à $f\circ g$, $h^*(\mathcal{L})$
est ample relativement à $f\circ h$.
\end{proposition}

Comme $Y$ est affine, $g^*(\mathcal{L})$ est ample, et il s'agit alors de
prouver que pour un choix convenable de $h$, $h^*(\mathcal{L})$ est ample
(\hyperref[II.4.6.6-fr]{4.6.6}). Nous allons voir qu'on peut remplacer $g$
par un morphisme fini surjectif $g':X''\to X$ tel que
${g'}^*(\mathcal{L})$ soit ample et que
$g'_*(\mathcal{O}_{X''})\otimes_{\mathcal{O}_X}\mathcal{R}(X)$ soit un
$\mathcal{R}(X)$-Module \emph{localement libre de rang constant}~; on sera
alors ramené aux conditions de (\hyperref[II.6.6.10-fr]{6.6.10}) et la
proposition sera démontrée.

Pour cela, soit $\mathcal{B}=g_*(\mathcal{O}_{X'})$~; désignons par $X_i$
($1\leq i\leq n$) les sous-préschémas fermés réduits de $X$ qui ont pour
espaces sous-jacents les composantes irréductibles de $X$
(\hyperref[I.5.2.1-fr]{I, 5.2.1})~; ils sont intègres par hypothèse. Soient
$X'_i$ le sous-préschéma fermé

\oldpage[II]{135}

$g^{-1}(X_i)$ de $X'$, $g_i:X'_i\to X_i$ le morphisme $g$ restreint à
$X'_i$, qui est fini (\hyperref[II.6.1.5-fr]{6.1.5, (iii)}) et surjectif~;
soit $k_i$ le \emph{rang} de la $\mathcal{O}_{X_i}$-Algèbre
$\mathcal{B}_i=(g_i)_*(\mathcal{O}_{X'_i})$. Comme
$\mathcal{B}_i\otimes_{\mathcal{O}_{X_i}}\mathcal{R}(X_i)$ est un
préfaisceau constant (\hyperref[I.7.3.5-fr]{I, 7.3.5}), le rang $k_i$ est
aussi le rang de la $(\mathcal{O}_X|U)$-Algèbre $\mathcal{B}|U$ pour tout
ouvert $U$ de $X$ ne rencontrant que la \emph{seule} composante irréductible
$X_i$. Si $T$ est un ensemble ouvert dans $X$ tel que $\mathcal{B}|T$ soit
isomorphe à $\mathcal{O}_X^m|T$, il résulte de la remarque précédente que les
nombres $k_i$ sont égaux à $m$ pour tous les indices $i$ tels que
$T\cap X_i\neq\varnothing$. Soit alors $U$ l'ouvert de $X$ formé des points
au voisinage desquels $\mathcal{B}$ est un $\mathcal{O}_X$-Module localement
libre, et soient $U_j$ ($1\leq j\leq s$) ses composantes connexes, qui sont
ouvertes dans $X$ et en nombre fini (puisque $U$ est noethérien)~; désignons
par $V_j$ le sous-préschéma fermé de $X'$, \emph{adhérence} du sous-préschéma
induit sur l'ouvert $g^{-1}(U_j)$
(\hyperref[I.9.5.11-fr]{I, 9.5.11}). D'après ce qui précède, pour tous les
indices $i$ tels que $X_i\cap U_j\neq\varnothing$, les rangs $k_i$ sont tous
égaux à un même entier $m_j$~; on notera d'ailleurs qu'un même $X_i$ ne peut
rencontrer deux $U_j$ d'indices distincts. Soient $i_\lambda$ les indices
$i$ tels que $X_i\cap U=\varnothing$. Considérons le produit $k$ de tous les
$k_i$, posons $n_i=k/k_i$, et soit $X''$ le préschéma défini comme suit.
Pour chaque $j$ ($1\leq j\leq s$), on considère $k/m_j$ préschémas
isomorphes à $V_j$, et pour chaque $\lambda$, $k/k_{i_\lambda}$ préschémas
isomorphes à $X'_{i_\lambda}$~; $X''$ est la \emph{somme} de tous ces
préschémas. On définit un morphisme $g'':X''\to X'$, se réduisant à
l'injection canonique dans chacun des summandes de $X''$~; il est clair que
$g''$ est un morphisme fini dominant, donc surjectif (puisqu'un morphisme
fini est fermé (\hyperref[II.6.1.10-fr]{6.1.10}))~; posons
$g'=g\circ g''$, qui est un morphisme fini surjectif $X''\to X$~; on a
${g'}^*(\mathcal{L})={g''}^*(g^*(\mathcal{L}))$, donc
${g'}^*(\mathcal{L})$ est un $\mathcal{O}_{X''}$-Module ample
(\hyperref[II.5.1.12-fr]{5.1.12}). Il est clair alors que pour ce nouveau
préschéma $X''$, les rangs définis comme les $k_i$ pour $X'$ sont tous
\emph{égaux à $k$}~; compte tenu de
(\hyperref[I.7.3.3-fr]{I, 7.3.3}), on en conclut aussitôt que pour tout ouvert
affine $T$ de $X$,
\[
\bigl(g'_*(\mathcal{O}_{X''})
\otimes_{\mathcal{O}_X}\mathcal{R}(X)\bigr)|T
\]
est un $(\mathcal{R}(X)|T)$-Module isomorphe à
$(\mathcal{R}(X)|T)^k$.

\begin{corollary}[6.6.12]
\label{II.6.6.12-fr}
Si, dans l'énoncé de (\hyperref[II.6.6.11-fr]{6.6.11}), on a $W=X$, alors
pour qu'un $\mathcal{O}_X$-Module inversible $\mathcal{L}$ soit ample
relativement à $f$, il faut et il suffit que $g^*(\mathcal{L})$ soit ample
relativement à $f\circ g$.
\end{corollary}

\begin{remark}[6.6.13]
\label{II.6.6.13-fr}
Au chapitre III, nous verrons que si $Y$ est noethérien, $f$ de type fini, et
si la restriction de $f$ au sous-préschéma fermé réduit de $X$ ayant $X-W$
pour espace sous-jacent est \emph{propre}, alors la conclusion de
(\hyperref[II.6.6.12-fr]{6.6.12}) est encore valable. Mais nous donnerons au
chapitre V des exemples de schémas algébriques $X$ sur un corps $K$ (le
morphisme structural $X\to\operatorname{Spec}(K)$ n'étant pas propre) dont
le normalisé $X'$ est quasi-affine, mais qui n'est pas quasi-affine (de sorte
que $\mathcal{O}_X$ n'est pas ample, bien que $\mathcal{O}_{X'}$ le soit
(\hyperref[II.5.1.12-fr]{5.1.12}) et que le morphisme $X'\to X$ soit fini
surjectif (\hyperref[II.6.3.10-fr]{6.3.10})). Nous allons voir au numéro
suivant que cette circonstance ne peut se produire lorsqu'on remplace
«~quasi-affine~» par «~affine~».
\end{remark}

\subsection{Le théorème de Chevalley.}
\label{subsection:II.6.7-fr}

Nous allons établir (à l'aide du critère de Serre
(\hyperref[II.5.2.1-fr]{5.2.1})) le théorème suivant, qui a été démontré par
C.~Chevalley par d'autres méthodes, dans le cas des schémas
\emph{algébriques}.

\oldpage[II]{136}

\begin{theorem}[6.7.1]
\label{II.6.7.1-fr}
Soient $X$ un schéma affine, $Y$ un préschéma noethérien, $f:X\to Y$ un
morphisme fini surjectif. Alors $Y$ est un schéma affine.
\end{theorem}

Il est clair que
$f_{\mathrm{red}}:X_{\mathrm{red}}\to Y_{\mathrm{red}}$ est fini
(\hyperref[II.6.1.5-fr]{6.1.5, (vi)})~; comme $X_{\mathrm{red}}$ est un
schéma affine, et qu'il revient au même de dire que $Y$ est affine ou que
$Y_{\mathrm{red}}$ est affine, puisque $Y$ est noethérien
(\hyperref[I.6.1.7-fr]{I, 6.1.7}), on voit qu'on peut supposer $Y$
\emph{réduit}. Pour toute partie fermée $Y'$ de $Y$, il y a alors un seul
sous-préschéma réduit de $Y$ ayant $Y'$ comme espace sous-jacent
(\hyperref[I.5.1.2-fr]{I, 5.1.2})~; son image réciproque $f^{-1}(Y')$,
canoniquement isomorphe à $X\times_Y Y'$
(\hyperref[I.4.4.1-fr]{I, 4.4.1}), est affine comme sous-préschéma fermé de
$X$, et la restriction de $f$ à $f^{-1}(Y')$, qui s'identifie à
$f\times_Y 1_{Y'}$, est un morphisme fini surjectif
(\hyperref[II.6.1.5-fr]{6.1.5, (iii)}). En vertu du principe de récurrence
noethérienne (\hyperref[0.2.2.2-fr]{0, 2.2.2}), on peut donc (compte tenu de
\hyperref[I.6.1.7-fr]{I, 6.1.7}) se ramener à démontrer le théorème sous
l'hypothèse que pour \emph{toute partie fermée $Y'\neq Y$}, tout
sous-préschéma fermé de $Y$ ayant $Y'$ pour espace sous-jacent est affine.
On en conclut que, \emph{pour tout $\mathcal{O}_Y$-Module cohérent
$\mathcal{F}$ dont le support (fermé) $Z$ est distinct de $Y$, on a
$\mathrm{H}^1(Y,\mathcal{F})=0$}. En effet, il existe un sous-préschéma fermé
$Y'$ de $Y$ ayant $Z$ pour espace sous-jacent et tel que, si $j:Y'\to Y$ est
l'injection canonique, on ait
$\mathcal{F}=j_*(j^*(\mathcal{F}))$
(\hyperref[I.9.3.5-fr]{I, 9.3.5})~; par suite
(\hyperref[II.5.2.3-fr]{5.2.3}),
$\mathrm{H}^1(Y,\mathcal{F})
=\mathrm{H}^1(Y',j^*(\mathcal{F}))=0$ d'après
(\hyperref[I.5.1.9.2-fr]{I, 5.1.9.2}).

Supposons d'abord que $Y$ ne soit pas irréductible, et soit $Y'$ une
composante irréductible de $Y$, et $Y''=Y-Y'$~; on désignera encore par
$Y'$ le sous-préschéma fermé réduit de $Y$ ayant $Y'$ pour espace
sous-jacent, et par $j$ l'injection canonique $Y'\to Y$. Soit $\mathcal{F}$
un $\mathcal{O}_Y$-Module cohérent, et considérons l'homomorphisme canonique
\[
\rho:\mathcal{F}\longrightarrow
\mathcal{F}'=j_*(j^*(\mathcal{F}))
\]
(\hyperref[0.4.4.3-fr]{0, 4.4.3})~; $\mathcal{F}'$ est un
$\mathcal{O}_Y$-Module cohérent en vertu de
(\hyperref[0.5.3.10-fr]{0, 5.3.10}) et
(\hyperref[0.5.3.12-fr]{0, 5.3.12}), car on a
$j_*(\mathcal{O}_{Y'})=\mathcal{O}_Y/\mathcal{J}$, en désignant par
$\mathcal{J}$ le faisceau d'idéaux de $\mathcal{O}_Y$ définissant le
sous-préschéma $Y'$. Par suite $\mathcal{G}=\operatorname{Ker}\rho$ et
$\mathcal{K}=\operatorname{Im}\rho$ sont aussi des $\mathcal{O}_Y$-Modules
cohérents (\hyperref[0.5.3.4-fr]{0, 5.3.4})~; or par définition la fibre
$\mathcal{F}'_y$ de $\mathcal{F}'$ au point générique $y$ de $Y'$ est égale
à $\mathcal{F}_y$, car $y$ est intérieur à $Y'$ et on a donc
$\mathcal{J}_y=0$, puisque $Y$ est réduit. On en conclut que $y$ n'est pas
contenu dans le support (fermé) de $\mathcal{G}$~; par ailleurs, le support de
$\mathcal{F}'$ (et \emph{a fortiori} celui de $\mathcal{K}$) est contenu
dans $Y'$~; autrement dit les supports de $\mathcal{G}$ et de $\mathcal{K}$
sont \emph{distincts de $Y$}. On en déduit que
$\mathrm{H}^1(Y,\mathcal{G})=\mathrm{H}^1(Y,\mathcal{K})=0$, et la suite
exacte de cohomologie appliquée à la suite exacte
$0\to\mathcal{G}\to\mathcal{F}\to\mathcal{K}\to0$ montre que
$\mathrm{H}^1(Y,\mathcal{F})=0$. On conclut alors par le critère de Serre
(\hyperref[II.5.2.1-fr]{5.2.1}).

Supposons donc $Y$ irréductible, et par suite \emph{intègre}. On peut aussi
supposer que $X$ est \emph{intègre}~: en effet, si on désigne par $X_i$ les
sous-préschémas fermés réduits de $X$ ayant pour espaces sous-jacents les
composantes irréductibles de $X$ (\hyperref[I.5.2.1-fr]{I, 5.2.1}) et par
$f_i$ la restriction de $f$ à $X_i$, un des $f_i$ au moins est dominant, et
comme c'est un morphisme fini (\hyperref[II.6.1.5-fr]{6.1.5}), il est
surjectif (\hyperref[II.6.1.10-fr]{6.1.10})~; comme d'autre part $X_i$ est un
schéma affine, on voit qu'on peut remplacer $X$ par $X_i$ dans l'énoncé.

\begin{lemma}[6.7.1.1]
\label{II.6.7.1.1-fr}
Soient $X$, $Y$ deux préschémas intègres noethériens, $x$ (resp. $y$) le
point générique de $X$ (resp. $Y$), $f:X\to Y$ un morphisme fini surjectif.
Soit $\mathcal{L}$ un $\mathcal{O}_X$-Module inversible tel qu'il existe un
voisinage ouvert affine $U$ de $y$ et une section
$g\in\Gamma(X,\mathcal{L})$ pour lesquels

\oldpage[II]{137}

$x\in X_g\subset f^{-1}(U)$. Alors il existe deux entiers $m>0$, $n>0$, un
homomorphisme
$u:\mathcal{O}_Y^m\to f_*(\mathcal{L}^{\otimes n})$ et un voisinage ouvert
$V$ de $y$ tels que la restriction $u|V$ soit un isomorphisme de
$\mathcal{O}_Y^m|V$ sur $f_*(\mathcal{L}^{\otimes n})|V$.
\end{lemma}

Soient $C$ l'anneau (intègre) de $U$, $k=\mathcal{O}_y$ son corps des
fractions~: comme $f$ est fini, $U'=f^{-1}(U)$ est affine
(\hyperref[II.1.3.2-fr]{1.3.2})~; soit $D$ son anneau (intègre) de corps des
fractions $K=\mathcal{O}_x$~; par hypothèse $D$ est un $C$-module de type
fini (\hyperref[II.6.1.4-fr]{6.1.4}), donc $K$ est une extension de rang fini
de $k$. La fibre
\[
f^{-1}(y)=X\times_Y\operatorname{Spec}(k(y))
=X\times_Y\operatorname{Spec}(\mathcal{O}_y)
\]
s'identifie à $\operatorname{Spec}(K)$
(\hyperref[I.3.6.6-fr]{I, 3.6.6})~; soient $s_i$ ($1\leq i\leq m$) des
éléments de $D$ formant une base de $K$ sur $k$. Il existe $n>0$ tel que les
sections $(s_i|X_g)g^{\otimes n}$ de $\mathcal{L}^{\otimes n}$ au-dessus de
$X_g$ se prolongent en sections $b_i$ ($1\leq i\leq m$) de
$\mathcal{L}^{\otimes n}$ au-dessus de $X$
(\hyperref[I.9.3.1-fr]{I, 9.3.1}). Les $b_i$ sont aussi par définition des
sections de $f_*(\mathcal{L}^{\otimes n})$ au-dessus de $Y$ et définissent
donc un homomorphisme
$u:\mathcal{O}_Y^m\to f_*(\mathcal{L}^{\otimes n})$
(\hyperref[0.5.1.1-fr]{0, 5.1.1})~; nous allons voir que $u$ répond à la
question. On a $\mathcal{L}^{\otimes n}|U'=\widetilde{M}$, où $M$ est un
$D$-module de type fini, donc si $\varphi$ est l'injection $C\to D$
correspondant au morphisme $f^{-1}(U)\to U$ restriction de $f$,
$M_{[\varphi]}$ est un $C$-module de type fini~; par suite
\[
f_*(\mathcal{L}^{\otimes n})|U=(M_{[\varphi]})^\sim
\]
(\hyperref[I.1.6.3-fr]{I, 1.6.3}) est cohérent et comme $U$ est un ouvert
affine quelconque de $Y$, $f_*(\mathcal{L}^{\otimes n})$ est cohérent~; en
outre $u|U=\widetilde{\theta}$, où $\theta$ est un $C$-homomorphisme
$C^m\to M_{[\varphi]}$, et $u_y=\theta_y$ est l'homomorphisme
\[
\theta\otimes1:K^m=C^m\otimes K\longrightarrow
M_{[\varphi]}\otimes K,
\]
or ce dernier est par définition un \emph{isomorphisme}, car les $(b_i)_x$
forment une base de $M_{[\varphi]}\otimes K$ sur $k$, $g_x$ étant par
hypothèse un générateur de $(\mathcal{L}^{\otimes n})_x$. On en conclut que
les supports des $\mathcal{O}_Y$-Modules $\operatorname{Ker}u$ et
$\operatorname{Coker}u$ ne contiennent pas $y$~; comme ces
$\mathcal{O}_Y$-Modules sont cohérents
(\hyperref[0.5.3.3-fr]{0, 5.3.3}), leurs supports sont fermés
(\hyperref[0.5.2.2-fr]{0, 5.2.2}), d'où le lemme.

Cela étant, les hypothèses du lemme
(\hyperref[II.6.7.1.1-fr]{6.7.1.1}) sont remplies dans le cas que nous
considérons en prenant $\mathcal{L}=\mathcal{O}_X$, puisque $X$ est affine
(\hyperref[I.1.1.10-fr]{I, 1.1.10})~; nous poserons
$\mathcal{A}=\mathcal{O}_Y$, $\mathcal{B}=f_*(\mathcal{O}_X)$. En vertu du
critère de Serre (\hyperref[II.5.2.1.1-fr]{5.2.1.1}), il suffit de prouver
que pour tout $\mathcal{O}_Y$-Module cohérent $\mathcal{F}$, on a
$\mathrm{H}^1(Y,\mathcal{F})=0$~; il suffit même de le prouver lorsque
$\mathcal{F}\subset\mathcal{O}_Y$, ce qui entraîne que $\mathcal{F}$ est
sans torsion puisque $Y$ est intègre~; en fait nous allons montrer que
$\mathrm{H}^1(Y,\mathcal{F})=0$ pour \emph{tout} $\mathcal{O}_Y$-Module
cohérent \emph{sans torsion} $\mathcal{F}$. Or, l'homomorphisme
$u:\mathcal{A}^m\to\mathcal{B}$ définit un homomorphisme
\[
v:\mathcal{G}
=\mathcal{H}om_{\mathcal{A}}(\mathcal{B},\mathcal{F})
\longrightarrow
\mathcal{H}om_{\mathcal{A}}(\mathcal{A}^m,\mathcal{F})
=\mathcal{F}^m.
\]
Montrons d'abord que $v$ est \emph{injectif}~: par hypothèse,
$\mathcal{T}=\operatorname{Coker}u$ a un support ne rencontrant pas $V$, donc
est un $\mathcal{O}_Y$-Module de torsion
(\hyperref[I.7.4.6-fr]{I, 7.4.6})~; la suite exacte
\[
\mathcal{A}^m\longrightarrow\mathcal{B}\longrightarrow\mathcal{T}
\longrightarrow0
\]
donne, par exactitude à gauche du foncteur
$\mathcal{H}om_{\mathcal{A}}$, la suite exacte
\[
0\longrightarrow
\mathcal{H}om_{\mathcal{A}}(\mathcal{T},\mathcal{F})
\longrightarrow\mathcal{G}\xrightarrow{v}\mathcal{F}^m.
\]
Mais comme $\mathcal{F}$ est sans torsion, on a
$\mathcal{H}om_{\mathcal{A}}(\mathcal{T},\mathcal{F})=0$
(\hyperref[0.5.2.6-fr]{0, 5.2.6}), d'où notre assertion. On a donc la suite
exacte
\[
0\longrightarrow\mathcal{G}\longrightarrow\mathcal{F}^m
\longrightarrow\operatorname{Coker}v\longrightarrow0
\]

\oldpage[II]{138}

où $\mathcal{G}$ et $\operatorname{Coker}v$ sont des
$\mathcal{O}_Y$-Modules cohérents
(\hyperref[0.5.3.4-fr]{0, 5.3.4} et
\hyperref[0.5.3.5-fr]{5.3.5})~; en vertu de la suite exacte de cohomologie,
il suffira de montrer que
$\mathrm{H}^1(Y,\mathcal{G})
=\mathrm{H}^1(Y,\operatorname{Coker}v)=0$ pour en déduire
$\mathrm{H}^1(Y,\mathcal{F}^m)
=(\mathrm{H}^1(Y,\mathcal{F}))^m=0$, et par suite
$\mathrm{H}^1(Y,\mathcal{F})=0$. Or la restriction $v|V$ est un isomorphisme,
donc le support de $\operatorname{Coker}v$ est distinct de $Y$, d'où
$\mathrm{H}^1(Y,\operatorname{Coker}v)=0$ par l'hypothèse du début. D'autre
part, $\mathcal{G}$ est un $\mathcal{B}$-Module cohérent
(\hyperref[I.9.6.4-fr]{I, 9.6.4})~; comme $X$ est affine au-dessus de $Y$, il
existe un $\mathcal{O}_X$-Module quasi-cohérent $\mathcal{K}$ tel que
$\mathcal{G}$ soit isomorphe à $f_*(\mathcal{K})$
(\hyperref[II.1.4.3-fr]{1.4.3})~; comme $X$ est affine, on a
$\mathrm{H}^1(X,\mathcal{K})=0$
(\hyperref[I.5.1.9.2-fr]{I, 5.1.9.2}), donc aussi
$\mathrm{H}^1(Y,\mathcal{G})=0$ par
(\hyperref[II.5.2.3-fr]{5.2.3}), ce qui achève la démonstration du théorème
(\hyperref[II.6.7.1-fr]{6.7.1}).

\begin{corollary}[6.7.2]
\label{II.6.7.2-fr}
Soient $X$ un préschéma noethérien, $(X_i)_{1\leq i\leq n}$ un recouvrement
fini de l'espace $X$ formé d'ensembles fermés. Pour que $X$ soit affine, il
faut et il suffit que pour chaque $i$, il existe un sous-préschéma fermé de
$X$, ayant $X_i$ pour espace sous-jacent et qui soit affine.
\end{corollary}

En effet, s'il en est ainsi, soit $X'$ le schéma \emph{somme} des $X_i$~; il
est clair que $X'$ est affine, et on définit un morphisme surjectif
$f:X'\to X$ en prenant pour restriction de $f$ à $X_i$ l'injection
canonique. Tout revient à vérifier que $f$ est \emph{fini}, en vertu de
(\hyperref[II.6.7.1-fr]{6.7.1}), et cela a été vu en
(\hyperref[II.6.1.5-fr]{6.1.5}).

\begin{corollary}[6.7.3]
\label{II.6.7.3-fr}
Pour qu'un préschéma noethérien $X$ soit affine, il faut et il suffit que les
sous-préschémas fermés réduits ayant pour espaces sous-jacents les composantes
irréductibles de $X$ soient affines.
\end{corollary}

\section{Critères valuatifs}
\label{section:II.7-fr}

Nous donnons dans ce paragraphe des critères valuatifs de séparation et de
propreté pour un morphisme, c'est-à-dire des critères qui font intervenir un
schéma auxiliaire variable de la forme $\operatorname{Spec}(A)$, où $A$ est
un anneau de valuation. Moyennant des hypothèses «~noethériennes~»
convenables, on peut dans ces critères se restreindre au cas où $A$ est un
anneau de valuation \emph{discrète}. Ce sera sans doute le seul cas que nous
aurons à appliquer par la suite, et nous n'avons introduit les anneaux de
valuation quelconques, dans le cas général, que pour faire le lien avec des
développements classiques.

\subsection{Rappels sur les anneaux de valuation.}
\label{subsection:II.7.1-fr}

\begin{env}[7.1.1]
\label{II.7.1.1-fr}
Parmi les diverses propriétés équivalentes qui caractérisent les anneaux de
valuation, nous utiliserons la suivante~: un anneau $A$ est dit
\emph{anneau de valuation} si c'est un anneau intègre qui n'est pas un
corps, et si, dans l'ensemble des anneaux locaux contenus dans le corps des
fractions $K$ de $A$ et distincts de $K$, $A$ est \emph{maximal} pour la
relation de domination (\hyperref[I.8.1.1-fr]{I, 8.1.1}). Rappelons qu'un
anneau de valuation est \emph{intégralement clos}. Si $A$ est un anneau de
valuation, $A_{\mathfrak{p}}$ est aussi un anneau de valuation pour tout
idéal premier $\mathfrak{p}\neq0$ de $A$.
\end{env}

\begin{env}[7.1.2]
\label{II.7.1.2-fr}
Soient $K$ un corps, $A$ un sous-anneau local de $K$ qui n'est pas un corps~;

\oldpage[II]{139}

il existe alors un anneau de valuation ayant $K$ pour corps des fractions et
qui domine $A$ (\cite[p.~1-07, lemme~2]{I-1}).

D'autre part, soient $B$ un anneau de valuation, $k$ son corps des restes,
$K$ son corps des fractions, $L$ une extension de $k$. Alors il existe un
anneau de valuation \emph{complet} $C$ qui domine $B$ et dont le corps des
restes est $L$. En effet, $L$ est extension algébrique d'une extension
transcendante pure $L'=k(T_\mu)_{\mu\in M}$~; on sait qu'on peut prolonger la
valuation de $K$ correspondant à $B$ à une valuation de
$K'=K(T_\mu)_{\mu\in M}$ de sorte que $L'$ soit le corps des restes de cette
valuation (\cite[p.~98]{II-24})~; remplaçant $B$ par le complété de l'anneau
de cette valuation prolongée, on voit qu'on peut se limiter au cas où $B$ est
complet et $L$ est une clôture algébrique de $k$. Si $\overline{K}$ est une
clôture algébrique de $K$, on peut alors prolonger à $\overline{K}$ la
valuation qui définit $B$, et le corps des restes correspondant est une
clôture algébrique de $k$, comme on le voit en relevant dans $\overline{K}$
les coefficients d'un polynôme unitaire de $k[T]$. On est donc finalement
ramené au cas où $L=k$ et il suffit alors de prendre pour $C$ le complété de
$B$ pour répondre à la question.
\end{env}

\begin{env}[7.1.3]
\label{II.7.1.3-fr}
Soient $K$ un corps, $A$ un sous-anneau de $K$~; la fermeture intégrale $A'$
de $A$ dans $K$ est l'intersection des anneaux de valuation contenant $A$ et
ayant $K$ pour corps des fractions
(\cite[p.~51, th.~2]{I-11}). La proposition
(\hyperref[II.7.1.2-fr]{7.1.2}) s'exprime sous la forme géométrique
équivalente~:
\end{env}

\begin{proposition}[7.1.4]
\label{II.7.1.4-fr}
Soient $Y$ un préschéma, $p:X\to Y$ un morphisme, $x$ un point de $X$,
$y=p(x)$, $y'\neq y$ une spécialisation
(\hyperref[0.2.1.2-fr]{0, 2.1.2}) de $y$. Il existe alors un schéma local
$Y'$, spectre d'un anneau de valuation, et un morphisme séparé $f:Y'\to Y$
tel que, si $a$ désigne l'unique point fermé de $Y'$ et $b$ le point
générique de $Y'$, on ait $f(a)=y'$ et $f(b)=y$. On peut en outre supposer
vérifiée l'une des deux propriétés additionnelles suivantes~:
\begin{enumerate}
\item[\rm{(i)}] $Y'$ est le spectre d'un anneau de valuation \emph{complet}
dont le corps des restes est algébriquement clos, et il existe un
$k(y)$-homomorphisme $k(x)\to k(b)$.
\item[\rm{(ii)}] Il existe un $k(y)$-isomorphisme
$k(x)\xrightarrow{\sim}k(b)$.
\end{enumerate}
\end{proposition}

Soit $Y_1$ le sous-préschéma fermé réduit de $Y$ ayant
$\overline{\{y\}}$ comme espace sous-jacent
(\hyperref[I.5.2.1-fr]{I, 5.2.1}), et soit $X_1$ le sous-préschéma fermé
image réciproque $p^{-1}(Y_1)$~; comme $y'\in\overline{\{y\}}$ par hypothèse
et que $k(x)$ est le même dans $X$ et dans $X_1$, on peut supposer $Y$
\emph{intègre}, de point générique $y$~; $\mathcal{O}_{y'}$ est alors un
anneau local intègre qui n'est pas un corps, et dont le corps des fractions
est $\mathcal{O}_y=k(y)$, et $k(x)$ est une extension de $k(y)$. Pour
réaliser les conditions $f(a)=y'$ et $f(b)=y$ ainsi que la condition
additionnelle (i) (resp. (ii)), on prendra
$Y'=\operatorname{Spec}(A')$, où $A'$ est un anneau de valuation dominant
$\mathcal{O}_{y'}$ et qui est complet et a un corps des restes
algébriquement clos extension de $k(x)$ (resp. un anneau de valuation $A'$
dominant $\mathcal{O}_{y'}$ et dont $k(x)$ est le corps des fractions)~;
l'existence de $A'$ est garantie par
(\hyperref[II.7.1.2-fr]{7.1.2}).

\begin{env}[7.1.5]
\label{II.7.1.5-fr}
Rappelons qu'un anneau local $A$ est dit \emph{de dimension $1$} s'il existe
un idéal premier distinct de l'idéal maximal $\mathfrak{m}$, et si tout idéal
premier de $A$ distinct de $\mathfrak{m}$ est un idéal premier
\emph{minimal}~; lorsque $A$ est \emph{intègre}, il revient au même de dire
que $\mathfrak{m}$ et $(0)$ sont les seuls idéaux premiers et
$\mathfrak{m}\neq(0)$~; autrement dit, $Y=\operatorname{Spec}(A)$ est
réduit à deux

\oldpage[II]{140}

points $a$, $b$~: $a$ est l'unique point \emph{fermé}, on a
$\mathfrak{j}_a=\mathfrak{m}$ et $k(a)=k$ est le \emph{corps résiduel}
$k=A/\mathfrak{m}$~; $b$ est le \emph{point générique} de $Y$,
$\mathfrak{j}_b=(0)$, l'ensemble $\{b\}$ étant l'unique ensemble ouvert de
$Y$ distinct de $\varnothing$ et de $Y$ (ensemble ouvert qui est donc
\emph{partout dense}), et $k(b)=K$ est le \emph{corps des fractions} de $A$.
\end{env}

\begin{env}[7.1.6]
\label{II.7.1.6-fr}
Pour un anneau local $A$, noethérien et de dimension $1$, on sait que les
conditions suivantes sont équivalentes
(\cite[p.~2-08 et 17-01]{I-1})~:
\begin{enumerate}
\item[\rm{a)}] $A$ est normal~;
\item[\rm{b)}] $A$ est régulier~;
\item[\rm{c)}] $A$ est un anneau de valuation~;
\end{enumerate}
en outre, $A$ est alors un \emph{anneau de valuation discrète}. Les prop.
(\hyperref[II.7.1.2-fr]{7.1.2}) et
(\hyperref[II.7.1.3-fr]{7.1.3}) ont alors les analogues suivants pour les
anneaux de valuation discrète~:
\end{env}

\begin{proposition}[7.1.7]
\label{II.7.1.7-fr}
Soient $A$ un anneau local intègre noethérien qui n'est pas un corps, $K$ son
corps des fractions, $L$ une extension de type fini de $K$~; il existe alors
un anneau de valuation discrète ayant $L$ pour corps des fractions et
dominant $A$.
\end{proposition}

Supposons d'abord $L=K$. Soient $\mathfrak{m}$ l'idéal maximal de $A$,
$(x_1,\ldots,x_n)$ un système de générateurs non nuls de $\mathfrak{m}$,
$B$ le sous-anneau
$A[x_2/x_1,\ldots,x_n/x_1]$ de $K$, qui est noethérien. Il est immédiat que
l'idéal $\mathfrak{m}B$ de $B$ est identique à l'idéal principal $x_1B$~;
si $\mathfrak{p}$ est un idéal premier minimal de $x_1B$, on sait que
$\mathfrak{p}$ est de rang $1$ (\cite[t.~I, p.~277]{I-13})~; autrement dit
$B_{\mathfrak{p}}$ est un anneau local noethérien \emph{de dimension $1$}~;
il est clair que
$\mathfrak{p}B_{\mathfrak{p}}\cap A$ est un idéal de $A$ contenant
$\mathfrak{m}$ et ne contenant pas $1$, donc égal à $\mathfrak{m}$, et par
suite $B_{\mathfrak{p}}$ \emph{domine} $A$
(\hyperref[I.8.1.1-fr]{I, 8.1.1}). Il résulte du th. de Krull-Akizuki
(\cite[p.~293]{II-25}) que la clôture intégrale $C$ de $B_{\mathfrak{p}}$
est un anneau noethérien (bien que $C$ ne soit pas nécessairement un
$B_{\mathfrak{p}}$-module de type fini)~; si $\mathfrak{n}$ est un idéal
maximal de $C$, $C_{\mathfrak{n}}$ est un anneau local noethérien normal de
dimension $1$ (\cite[p.~295]{II-25}), donc un anneau de valuation discrète,
qui domine $B_{\mathfrak{p}}$ et \emph{a fortiori} $A$.

Si maintenant $L$ est une extension de type fini de $K$, on peut, d'après ce
qui précède, se limiter au cas où $A$ est déjà un anneau de valuation
discrète. Soit $w$ une valuation de $K$ associée à $A$~; il existe une
valuation discrète $w'$ de $L$ qui \emph{prolonge} $w$~: on se ramène en
effet par récurrence sur le nombre de générateurs de $L$ au cas où
$L=K(\alpha)$, et alors la proposition est classique
(\cite[p.~106]{II-24}).

\begin{corollary}[7.1.8]
\label{II.7.1.8-fr}
Soient $A$ un anneau intègre noethérien, $K$ son corps des fractions, $L$ une
extension de type fini de $K$. Alors la fermeture intégrale de $A$ dans $L$
est l'intersection des anneaux de valuation discrète ayant $L$ pour corps des
fractions et contenant $A$.
\end{corollary}

En effet, un tel anneau de valuation discrète, étant normal, contient
\emph{a fortiori} tout élément de $L$ entier sur $A$. Il suffit donc de
prouver que si $x\in L$ n'est pas entier sur $A$, il existe un anneau de
valuation discrète $C$ ayant $L$ pour corps des fractions, contenant $A$ et
ne contenant pas $x$. L'hypothèse sur $x$ signifie en effet que l'on a
$x\notin B=A[1/x]$, autrement dit $1/x$ n'est pas inversible dans l'anneau
noethérien $B$. Il y a donc un idéal premier $\mathfrak{p}$ de $B$ contenant
$1/x$. L'anneau local intègre $B_{\mathfrak{p}}$ est noethérien et contenu
dans $L$, qui est une extension de type fini du corps des fractions de
$B_{\mathfrak{p}}$ (ce dernier contenant $K$). En vertu de
(\hyperref[II.7.1.7-fr]{7.1.7}), il y a donc un anneau de valuation discrète
$C$, dominant $B_{\mathfrak{p}}$ et ayant $L$ pour corps des fractions~;
comme $1/x\in\mathfrak{p}B_{\mathfrak{p}}$ appartient à l'idéal maximal de
$C$, on a $x\notin C$, ce qui conclut la démonstration.

La forme géométrique de (\hyperref[II.7.1.7-fr]{7.1.7}) est la suivante~:

\oldpage[II]{141}
\begin{proposition}[7.1.9]
\label{II.7.1.9-fr}
Soient $Y$ un préschéma localement noethérien, $p:X\to Y$ un morphisme
localement de type fini, $x$ un point de $X$, $y=p(x)$, $y'\neq y$ une
spécialisation de $y$. Il existe alors un schéma local $Y'$, spectre d'un
anneau de valuation discrète, un morphisme séparé $f:Y'\to Y$ et une
$Y$-application rationnelle $g$ de $Y'$ dans $X$, tels que, si $a$ désigne le
point fermé de $Y'$ et $b$ son point générique, on ait $f(a)=y'$,
$f(b)=y$, $g(b)=x$, et que dans le diagramme commutatif
\[
\tag{7.1.9.1}
\label{II.7.1.9.1-fr}
\begin{tikzcd}[column sep=large,row sep=large]
& k(x) \arrow[dl,"\gamma"'] \\
k(b) & k(y) \arrow[u,"\pi"] \arrow[l,"\varphi"']
\end{tikzcd}
\]
(où $\pi$, $\varphi$, $\gamma$ sont les homomorphismes correspondant
respectivement à $p$, $f$ et $g$), $\gamma$ soit une bijection.
\end{proposition}

Comme dans (\hyperref[II.7.1.4-fr]{7.1.4}), on peut se ramener au cas où $Y$
est intègre de point générique $y$ (compte tenu de
(\hyperref[I.6.3.4-fr]{I, 6.3.4, (iv)})) et comme la question est locale sur
$X$ et $Y$, on peut supposer $p$ de type fini~; on est alors dans la situation
de (\hyperref[II.7.1.4-fr]{7.1.4}), avec la propriété supplémentaire que
$k(x)$ est une extension \emph{de type fini} de $k(y)$
(\hyperref[I.6.4.11-fr]{I, 6.4.11}) et que $\mathcal{O}_{y'}$ est noethérien~;
cela permet d'appliquer (\hyperref[II.7.1.7-fr]{7.1.7}) et de prendre
$Y'=\operatorname{Spec}(A')$, où $A'$ est un anneau de valuation discrète
dominant $\mathcal{O}_{y'}$ et dont $k(x)$ est le corps des fractions. On a
donc ainsi défini un diagramme commutatif
(\hyperref[II.7.1.9.1-fr]{7.1.9.1}), où $\gamma$ est une bijection, et $\pi$
et $\varphi$ correspondent aux morphismes $p$ et $f$. En outre, comme $X$ et
$Y$ sont localement noethériens (\hyperref[I.6.6.2-fr]{I, 6.6.2}) et que $Y'$
est intègre, il y a une $Y$-application rationnelle $g$ et une seule de $Y'$
dans $X$ à laquelle correspond l'isomorphisme $\gamma$
(\hyperref[I.7.1.15-fr]{I, 7.1.15}), ce qui achève la démonstration.

\subsection{Critère valuatif de séparation}
\label{subsection:II.7.2-fr}

\begin{proposition}[7.2.1]
\label{II.7.2.1-fr}
Soient $X$, $Y$ deux préschémas, $f:X\to Y$ un morphisme quasi-compact. Les
deux conditions suivantes sont équivalentes~:
\begin{enumerate}
\item[\rm a)] Le morphisme $f$ est fermé.
\item[\rm b)] Pour tout $x\in X$ et toute spécialisation $y'$ de $y=f(x)$,
distincte de $y$, il existe une spécialisation $x'$ de $x$ telle que
$f(x')=x$.
\end{enumerate}
\end{proposition}

La condition b) exprime que $f(\overline{\{x\}})=\overline{\{y\}}$ et est donc
conséquence de a). Pour montrer que b) entraîne a), considérons une partie
fermée $X'$ de l'espace sous-jacent $X$~; soit $Y'=\overline{f(X')}$ et
montrons que $Y'=f(X')$. Considérons les sous-préschémas fermés réduits de
$X$ et $Y$ respectivement ayant pour espaces sous-jacents $X'$ et $Y'$
(\hyperref[I.5.2.1-fr]{I, 5.2.1})~; il y a alors un morphisme $f':X'\to Y'$ tel
que le diagramme
\[
\begin{tikzcd}[column sep=large,row sep=large]
X' \arrow[r,"f'"] \arrow[d] & Y' \arrow[d] \\
X \arrow[r,"f"'] & Y
\end{tikzcd}
\]
soit commutatif (\hyperref[I.5.2.2-fr]{I, 5.2.2}), et comme $f$ est
quasi-compact, il en est de même de $f'$. On est donc ramené à prouver que si
$f$ est un morphisme quasi-compact et \emph{dominant}, alors la
\oldpage[II]{142}
condition b) implique que $f(X)=Y$. Or, soit $y'$ un point de $Y$ et soit $y$
le point générique d'une composante irréductible de $Y$ contenant $y'$~; en
vertu de b), il suffit de montrer que $f^{-1}(y)$ n'est pas vide. Mais on sait
que cette propriété est conséquence du fait que $f$ est dominant et
quasi-compact (\hyperref[I.6.6.5-fr]{I, 6.6.5}).

\begin{corollary}[7.2.2]
\label{II.7.2.2-fr}
Soit $f:X\to Y$ un morphisme d'immersion quasi-compact. Pour que l'espace
sous-jacent $X$ soit fermé dans $Y$, il faut et il suffit qu'il contienne toute
spécialisation (dans $Y$) de chacun de ses points.
\end{corollary}

\begin{proposition}[7.2.3]
\label{II.7.2.3-fr}
Soient $Y$ un préschéma (resp. un préschéma localement noethérien), $f:X\to Y$
un morphisme (resp. un morphisme localement de type fini). Les conditions
suivantes sont équivalentes~:
\begin{enumerate}
\item[\rm a)] $f$ est séparé.
\item[\rm b)] Le morphisme diagonal $X\to X\times_YX$ est quasi-compact, et
pour tout $Y$-préschéma de la forme $Y'=\operatorname{Spec}(A)$, où $A$ est un
anneau de valuation (resp. un anneau de valuation discrète), deux
$Y$-morphismes de $Y'$ dans $X$ qui coïncident au point générique de $Y'$ sont
égaux.
\item[\rm c)] Le morphisme diagonal $X\to X\times_YX$ est quasi-compact, et
pour tout $Y$-préschéma de la forme $Y'=\operatorname{Spec}(A)$, où $A$ est un
anneau de valuation (resp. un anneau de valuation discrète), deux
$Y'$-sections de $X'=X_{(Y')}$ qui coïncident au point générique de $Y'$ sont
égales.
\end{enumerate}
\end{proposition}

L'équivalence de b) et de c) résulte de la correspondance biunivoque entre
$Y$-morphismes de $Y'$ dans $X$ et $Y'$-sections de $X'$
(\hyperref[I.3.3.14-fr]{I, 3.3.14}). Si $X$ est séparé sur $Y$, la condition b)
est vérifiée en vertu de (\hyperref[I.7.2.2.1-fr]{I, 7.2.2.1}), puisque $Y'$
est intègre. Il reste à prouver que b) entraîne que le morphisme diagonal
$\Delta:X\to X\times_YX$ est fermé, et il revient au même de montrer qu'il
vérifie le critère de (\hyperref[II.7.2.2-fr]{7.2.2}). Or, soit $z$ un point de
la diagonale $\Delta(X)$, $z'\neq z$ une spécialisation de $z$ dans
$X\times_YX$. Il existe alors (\hyperref[II.7.1.4-fr]{7.1.4}) un anneau de
valuation $A$ et un morphisme $f$ de $Y'=\operatorname{Spec}(A)$ dans
$X\times_YX$, qui applique le point fermé $a$ de $Y'$ sur $z'$ et le point
générique $b$ de $Y'$ sur $z$~; ce morphisme fait de $Y'$ un
$(X\times_YX)$-préschéma, et \emph{a fortiori} un $Y$-préschéma. Si on compose
avec $f$ les deux projections de $X\times_YX$, on obtient deux $Y$-morphismes
$g_1$, $g_2$ de $Y'$ dans $X$, qui par hypothèse coïncident au point $b$~; ils
sont donc égaux à un même morphisme $g$, ce qui signifie
(\hyperref[I.5.3.1-fr]{I, 5.3.1}) que $f$ se factorise en $f=\Delta\circ g$, et
par suite $z'\in\Delta(X)$. Lorsqu'on suppose $Y$ localement noethérien et $f$
localement de type fini, $X\times_YX$ est localement noethérien
(\hyperref[I.6.6.7-fr]{I, 6.6.7})~; on peut alors reprendre le raisonnement
précédent en y supposant que $A$ est un anneau de valuation discrète, en vertu
de (\hyperref[II.7.1.9-fr]{7.1.9}).

\begin{remark}[7.2.4]
\label{II.7.2.4-fr}
\begin{enumerate}
\item[\rm (i)] L'hypothèse que le morphisme $\Delta$ est quasi-compact est
toujours vérifiée lorsque $Y$ est localement noethérien et $f$ localement de
type fini, car $X\times_YX$ est alors localement noethérien
(\hyperref[I.6.6.4-fr]{I, 6.6.4, (i)}). Dans le cas général, elle signifie aussi
que pour tout recouvrement $(U_\alpha)$ de $X$ par des ouverts affines, les
ensembles $U_\alpha\cap U_\beta$ sont \emph{quasi-compacts}.
\item[\rm (ii)] Pour que $f$ soit séparé, il \emph{suffit} que la condition b)
ou la condition c) soit vérifiée pour un anneau de valuation $A$ qui est
\emph{complet} et dont le corps des restes est \emph{algébriquement clos}~;
cela résulte en effet de la démonstration de
(\hyperref[II.7.2.3-fr]{7.2.3}) et de
(\hyperref[II.7.1.4-fr]{7.1.4}).
\end{enumerate}
\end{remark}

\oldpage[II]{143}
\subsection{Critère valuatif de propreté}
\label{subsection:II.7.3-fr}

\begin{proposition}[7.3.1]
\label{II.7.3.1-fr}
Soient $A$ un anneau de valuation, $Y=\operatorname{Spec}(A)$, $b$ le point
générique de $Y$, $X$ un schéma intègre, $f:X\to Y$ un morphisme fermé tel
que $f^{-1}(b)$ se réduise à un point $x$ et que l'homomorphisme correspondant
$k(b)\to k(x)$ soit bijectif. Alors $f$ est un isomorphisme.
\end{proposition}

Comme $f$ est fermé et dominant, on a $f(X)=Y$~; il suffit
(\hyperref[I.4.2.2-fr]{I, 4.2.2}) de prouver que pour tout $y'\neq b$ dans $Y$,
il existe un \emph{seul} point $x'$ tel que $f(x')=y'$ et que l'homomorphisme
correspondant $\mathcal{O}_{y'}\to\mathcal{O}_{x'}$ soit bijectif, car $f$ sera
alors un homéomorphisme. Or, si $f(x')=y'$, $\mathcal{O}_{x'}$ est un anneau
local contenu dans $K=k(x)=k(b)$ et dominant $\mathcal{O}_{y'}$~; ce dernier
est l'anneau local $A_{y'}$, donc est un anneau de valuation
(\hyperref[II.7.1.1-fr]{7.1.1}) ayant $K$ pour corps des fractions. Par
ailleurs on a $\mathcal{O}_{x'}\neq K$ puisque $x'$ n'est pas le point
générique de $X$ (\hyperref[0.2.1.3-fr]{0, 2.1.3})~; on en conclut que
$\mathcal{O}_{x'}=\mathcal{O}_{y'}$. Comme $X$ est un schéma intègre, la
relation $\mathcal{O}_{x'}=\mathcal{O}_{x''}$ entraîne $x'=x''$
(\hyperref[I.8.2.2-fr]{I, 8.2.2}), ce qui achève la démonstration.

\begin{env}[7.3.2]
\label{II.7.3.2-fr}
Soient $A$ un anneau de valuation, $Y=\operatorname{Spec}(A)$, $b$ le point
générique de $Y$, de sorte que $\mathcal{O}_b=k(b)$ est égal à $K$, corps des
fractions de $A$~; soit $f:X\to Y$ un morphisme. On sait
(\hyperref[I.7.1.4-fr]{I, 7.1.4}) que les $Y$-sections \emph{rationnelles} de
$X$ sont en correspondance biunivoque avec les \emph{germes} de $Y$-sections
(définies dans des voisinages de $b$) au point $b$, d'où une application
canonique
\[
\tag{7.3.2.1}
\label{II.7.3.2.1-fr}
\Gamma_{\mathrm{rat}}(X/Y)\longrightarrow
\Gamma(f^{-1}(b)/\operatorname{Spec}(K))
\]
les éléments de $\Gamma(f^{-1}(b)/\operatorname{Spec}(K))$ s'identifiant
d'ailleurs, par définition (\hyperref[I.3.4.5-fr]{I, 3.4.5}) aux \emph{points
de $f^{-1}(b)=X\otimes_AK$ rationnels sur $K$}. Lorsque $f$ est \emph{séparé},
il résulte de (\hyperref[I.5.4.7-fr]{I, 5.4.7}) que l'application
(\hyperref[II.7.3.2.1-fr]{7.3.2.1}) est \emph{injective}, puisque $Y$ est un
schéma intègre.

Composant (\hyperref[II.7.3.2.1-fr]{7.3.2.1}) avec l'application canonique
$\Gamma(X/Y)\to\Gamma_{\mathrm{rat}}(X/Y)$
(\hyperref[I.7.1.2-fr]{I, 7.1.2}), on obtient une application canonique
\[
\tag{7.3.2.2}
\label{II.7.3.2.2-fr}
\Gamma(X/Y)\longrightarrow \Gamma(f^{-1}(b)/\operatorname{Spec}(K)).
\]
Lorsque $f$ est \emph{séparé}, cette application est encore \emph{injective}
(\hyperref[I.5.4.7-fr]{I, 5.4.7}).
\end{env}

\begin{proposition}[7.3.3]
\label{II.7.3.3-fr}
Soient $A$ un anneau de valuation de corps des fractions $K$,
$Y=\operatorname{Spec}(A)$, $b$ le point générique de $Y$, $f:X\to Y$ un
morphisme \emph{séparé} et \emph{fermé}. Alors l'application canonique
(\hyperref[II.7.3.2.2-fr]{7.3.2.2}) est bijective (ce qui revient à dire
qu'elle est \emph{surjective} et entraîne que les $Y$-sections rationnelles de
$X$ sont \emph{partout définies}).
\end{proposition}

Soit donc $x$ un point de $f^{-1}(b)$, \emph{rationnel sur $K$}. Comme $f$ est
séparé, il en est de même du morphisme
$f^{-1}(b)\to\operatorname{Spec}(K)$ correspondant à $f$
(\hyperref[I.5.5.1-fr]{I, 5.5.1, (iv)}), et toute section de $f^{-1}(b)$ étant
une immersion fermée (\hyperref[I.5.4.6-fr]{I, 5.4.6}), $\{x\}$ est fermé
\emph{dans $f^{-1}(b)$}. Considérons le sous-préschéma fermé réduit $X'$ de
$X$ ayant pour espace sous-jacent l'adhérence $\overline{\{x\}}$ de $\{x\}$
dans $X$. Il est clair que la restriction de $f$ à $X'$ vérifie les hypothèses
de (\hyperref[II.7.3.1-fr]{7.3.1}), donc est un \emph{isomorphisme} de $X'$ sur
$Y$, dont l'isomorphisme réciproque est la $Y$-section de $X$ cherchée.

\begin{env}[7.3.4]
\label{II.7.3.4-fr}
Pour énoncer les deux résultats suivants, nous nous servirons d'une
terminologie qui sera justifiée et discutée dans le chapitre IV~: si $F$ est
une partie
\oldpage[II]{144}
d'un préschéma $Y$, on appelle \emph{codimension} de $F$ dans $Y$ et on note
$\operatorname{codim}_YF$ la borne inférieure des entiers
$\dim(\mathcal{O}_z)$, où $z$ parcourt $F$.
\end{env}

\begin{corollary}[7.3.5]
\label{II.7.3.5-fr}
Soient $Y$ un préschéma localement noethérien réduit, $N$ l'ensemble des
points $y\in Y$ où $Y$ n'est pas régulier
(\hyperref[0.4.1.4-fr]{0, 4.1.4})~; on suppose que
$\operatorname{codim}_YN\geq2$. Soit $f:X\to Y$ un morphisme de type fini,
séparé et fermé et soit $g$ une $Y$-section rationnelle de $X$~; si $Y'$ est
l'ensemble des points de $Y$ où $g$ n'est pas définie (ensemble qui est
\emph{fermé} (\hyperref[I.7.2.1-fr]{I, 7.2.1})) on a
$\operatorname{codim}_YY'\geq2$.
\end{corollary}

Il suffit de prouver que $g$ est définie en tout point $z\in Y$ tel que
$\dim\mathcal{O}_z\leq1$. Si $\dim\mathcal{O}_z=0$, $z$ est point générique
d'une composante irréductible de $Y$ (\hyperref[I.1.1.14-fr]{I, 1.1.14}), donc
appartient à tout ouvert partout dense de $Y$, et en particulier au domaine de
définition de $g$. Supposons donc que $\dim\mathcal{O}_z=1$~;
$\mathcal{O}_z$ est alors par hypothèse un anneau local noethérien régulier,
donc (\hyperref[II.7.1.6-fr]{7.1.6}) un \emph{anneau de valuation discrète}.
Soit $Z=\operatorname{Spec}(\mathcal{O}_z)$~; comme $U=Y-Y'$ est partout
dense, $U\cap Z$ n'est pas vide (\hyperref[I.2.4.2-fr]{I, 2.4.2})~; soit $g'$
l'application rationnelle de $Z$ dans $X$ induite par $g$
(\hyperref[I.7.2.8-fr]{I, 7.2.8})~; il suffit de montrer que $g'$ est un
\emph{morphisme} (\hyperref[I.7.2.9-fr]{I, 7.2.9}). Or, $g'$ peut être
considérée comme une $Z$-section rationnelle du $Z$-préschéma
$f^{-1}(Z)=X\times_YZ$~; il est clair que le morphisme $f^{-1}(Z)\to Z$
correspondant à $f$ est fermé, et il résulte de
(\hyperref[I.5.5.1-fr]{I, 5.5.1, (i)}) qu'il est séparé~; on conclut alors de
(\hyperref[II.7.3.3-fr]{7.3.3}) que $g'$ est partout définie~; comme $Z$ est
réduit et $X$ séparé sur $Y$, $g'$ est un morphisme
(\hyperref[I.7.2.2-fr]{I, 7.2.2}).

\begin{corollary}[7.3.6]
\label{II.7.3.6-fr}
Soient $S$ un préschéma localement noethérien, $X$ et $Y$ deux
$S$-préschémas~; on suppose $Y$ réduit, et en outre que l'ensemble $N$ des
points $y\in Y$ où $Y$ n'est pas régulier soit tel que
$\operatorname{codim}_YN\geq2$~; on suppose enfin que le morphisme structural
$X\to S$ soit propre. Soit $f$ une $S$-application rationnelle de $Y$ dans
$X$ et soit $Y'$ l'ensemble des points de $Y$ où $f$ n'est pas définie~;
alors on a $\operatorname{codim}_YY'\geq2$.
\end{corollary}

On sait (\hyperref[I.7.1.2-fr]{I, 7.1.2}) qu'on peut identifier les
$S$-applications rationnelles de $Y$ dans $X$ aux $Y$-sections rationnelles
de $X\times_SY$~; comme le morphisme structural $X\times_SY\to Y$ est fermé
(\hyperref[II.5.4.1-fr]{5.4.1}), on peut appliquer
(\hyperref[II.7.3.5-fr]{7.3.5}), d'où le corollaire.

\begin{remark}[7.3.7]
\label{II.7.3.7-fr}
Les hypothèses faites sur $Y$ dans (\hyperref[II.7.3.5-fr]{7.3.5}) et
(\hyperref[II.7.3.6-fr]{7.3.6}) seront en particulier réalisées lorsque $Y$
est \emph{normal} (\hyperref[0.4.1.4-fr]{0, 4.1.4}), en vertu de
(\hyperref[II.7.1.6-fr]{7.1.6}).
\end{remark}

Nous pouvons caractériser les morphismes universellement fermés (resp.
propres) par une réciproque de (\hyperref[II.7.3.3-fr]{7.3.3})~:

\begin{theorem}[7.3.8]
\label{II.7.3.8-fr}
Soient $Y$ un préschéma (resp. un préschéma localement noethérien), $f:X\to Y$
un morphisme quasi-compact séparé (resp. de type fini). Les conditions
suivantes sont équivalentes~:
\begin{enumerate}
\item[\rm a)] $f$ est universellement fermé (resp. propre).
\item[\rm b)] Pour tout $Y$-schéma de la forme
$Y'=\operatorname{Spec}(A)$, où $A$ est un anneau de valuation (resp. un anneau
de valuation discrète) de corps des fractions $K$, l'application canonique
\[
\operatorname{Hom}_Y(Y',X)\longrightarrow
\operatorname{Hom}_Y(\operatorname{Spec}(K),X)
\]
correspondant à l'injection canonique $A\to K$, est surjective (resp.
bijective).
\oldpage[II]{145}
\item[\rm c)] Pour tout $Y$-schéma de la forme $Y'=\operatorname{Spec}(A)$,
où $A$ est un anneau de valuation (resp. un anneau de valuation discrète)
l'application canonique (\hyperref[II.7.3.2.2-fr]{7.3.2.2}) relative au
$Y'$-préschéma $X_{(Y')}$ est surjective (resp. bijective).
\end{enumerate}
\end{theorem}

L'équivalence de b) et c) résulte aussitôt de
(\hyperref[I.3.3.14-fr]{I, 3.3.14})~; a) implique c), car a) entraîne, dans
l'une ou l'autre hypothèse, que $f_{(Y')}$ est séparé
(\hyperref[I.5.5.1-fr]{I, 5.5.1, (iv)}) et fermé, et il suffit d'appliquer
(\hyperref[II.7.3.3-fr]{7.3.3}). Reste à prouver que b) entraîne a).
Plaçons-nous d'abord dans le cas où $Y$ est quelconque, $f$ séparé et
quasi-compact. Si la condition b) est vérifiée par $f$, elle l'est aussi par
$f_{(Y'')}:X_{(Y'')}\to Y''$, où $Y''$ est un $Y$-préschéma quelconque, en
vertu de l'équivalence de b) et c), et du fait que
$X_{(Y'')}\times_{Y''}Y'=X\times_YY'$ pour tout morphisme $Y'\to Y''$
(\hyperref[I.3.3.9.1-fr]{I, 3.3.9.1})~; comme en outre $f_{(Y'')}$ est séparé
et quasi-compact quand $f$ l'est
(\hyperref[I.5.5.1-fr]{I, 5.5.1, (iv)} et
\hyperref[I.6.6.4-fr]{6.6.4, (iii)}), on est ramené à prouver que b) entraîne
que $f$ est \emph{fermé}. Pour cela, il suffit de vérifier la condition b) de
(\hyperref[II.7.2.1-fr]{7.2.1}). Soient donc $x\in X$, $y'$ une spécialisation
de $y=f(x)$, distincte de $y$~; en vertu de
(\hyperref[II.7.1.4-fr]{7.1.4}), il y a un schéma $Y'$, spectre d'un anneau de
valuation, et un morphisme séparé $g:Y'\to Y$ tels que, si $a$ désigne le point
fermé et $b$ le point générique de $Y'$, on ait $g(a)=y'$, $g(b)=y$, et qu'il
existe un $k(y)$-homomorphisme $k(x)\to k(b)$. Ce dernier correspond
canoniquement à un $Y$-morphisme $\operatorname{Spec}(k(b))\to X$
(\hyperref[I.2.4.6-fr]{I, 2.4.6}), et il résulte donc de b) qu'il existe un
$Y$-morphisme $h:Y'\to X$ auquel correspond le morphisme précédent. On a
alors $h(b)=x$~; si l'on pose $h(a)=x'$, $x'$ est spécialisation de $x$, et
l'on a $f(x')=f(h(a))=g(a)=y'$.

Si maintenant $Y$ est localement noethérien et $f$ de type fini, l'hypothèse
b) entraîne d'abord que $f$ est \emph{séparé}, en vertu de
(\hyperref[II.7.2.3-fr]{7.2.3}), le morphisme diagonal $X\to X\times_YX$
étant quasi-compact (\hyperref[II.7.2.4-fr]{7.2.4}). En outre, pour vérifier
que $f$ est propre, il suffit de montrer que $f_{(Y'')}:X_{(Y'')}\to Y''$ est
\emph{fermé} pour tout $Y$-préschéma $Y''$ \emph{de type fini}, compte tenu de
(\hyperref[II.5.6.3-fr]{5.6.3}). Comme alors $Y''$ est localement noethérien,
on peut reprendre le raisonnement fait dans le premier cas en prenant pour
$Y'$ un spectre d'anneau de valuation discrète, et en appliquant
(\hyperref[II.7.1.9-fr]{7.1.9}) au lieu de
(\hyperref[II.7.1.4-fr]{7.1.4}).

\begin{remark}[7.3.9]
\label{II.7.3.9-fr}
\begin{enumerate}
\item[\rm (i)] Lorsque $Y$ est un préschéma quelconque et $f$ un morphisme
séparé, pour que $f$ soit universellement fermé, il \emph{suffit} que la
condition b) ou la condition c) soit vérifiée pour les anneaux de valuation
$A$ \emph{complets} et dont le corps des restes est \emph{algébriquement
clos}~; cela résulte en effet de la démonstration et de
(\hyperref[II.7.1.4-fr]{7.1.4}).
\item[\rm (ii)] On déduit du critère c) de
(\hyperref[II.7.3.8-fr]{7.3.8}) une nouvelle démonstration du fait qu'un
morphisme projectif $X\to Y$ est fermé (\hyperref[II.5.5.3-fr]{5.5.3}), plus
proche des méthodes classiques. On peut en effet supposer $Y$ affine, et par
suite $X$ identifié à un sous-préschéma fermé d'un fibré projectif
$\mathbf{P}_Y^n$ (\hyperref[II.5.3.3-fr]{5.3.3})~; pour prouver que $X\to Y$
est fermé, il suffit de vérifier qu'il en est ainsi du morphisme structural
$\mathbf{P}_Y^n\to Y$, et le critère c) de
(\hyperref[II.7.3.8-fr]{7.3.8}), joint à
(\hyperref[II.4.1.3.1-fr]{4.1.3.1}), prouve qu'on est ramené à démontrer le
fait suivant~: \emph{si $Y$ est le spectre d'un anneau de valuation $A$, de
corps des fractions $K$, tout point de $\mathbf{P}_Y^n$ à valeurs dans $K$
provient (par restriction au point générique de $Y$) d'un point de
$\mathbf{P}_Y^n$ à valeurs dans $A$}. Or, tout $\mathcal{O}_Y$-Module
inversible est trivial (\hyperref[I.2.4.8-fr]{I, 2.4.8})~; donc, il résulte de
(\hyperref[II.4.2.6-fr]{4.2.6}) qu'un point de $\mathbf{P}_Y^n$ à valeurs dans
$K$ s'identifie à une classe d'éléments $(\zeta c_0,\zeta c_1,\ldots,\zeta
c_n)$ de $K$, où $\zeta\neq0$ et les $c_i$ sont des éléments non tous nuls de
$K$. Or, en multipliant les $c_i$ par un élément de $A$ de
\oldpage[II]{146}
valuation convenable, on peut supposer que les $c_i$ appartiennent tous à
$A$, et que l'un au moins est inversible. Mais alors
(\hyperref[II.4.2.6-fr]{4.2.6}), le système $(c_0,\ldots,c_n)$ définit aussi
un point de $\mathbf{P}_Y^n$ à valeurs dans $A$, ce qui démontre notre
assertion.
\item[\rm (iii)] Les critères (\hyperref[II.7.2.3-fr]{7.2.3}) et
(\hyperref[II.7.3.8-fr]{7.3.8}) sont surtout commodes quand on considère la
donnée d'un $Y$-préschéma $X$ comme équivalente à la donnée du foncteur
\[
X(Y')=\operatorname{Hom}_Y(Y',X)
\]
en un $Y$-préschéma $Y'$~; ces critères nous permettront par exemple de
prouver que sous certaines conditions les «~schémas de Picard~» sont propres.
\end{enumerate}
\end{remark}

\begin{corollary}[7.3.10]
\label{II.7.3.10-fr}
Soient $Y$ un schéma intègre (resp. un schéma intègre localement noethérien),
$X$ un schéma intègre, $f:X\to Y$ un morphisme dominant.
\begin{enumerate}
\item[\rm (i)] Si $f$ est quasi-compact et universellement fermé, tout anneau
de valuation dont le corps des fractions est le corps $R(X)$ des fonctions
rationnelles sur $X$, et qui domine un anneau local de $Y$, domine aussi un
anneau local de $X$.
\item[\rm (ii)] Inversement, supposons que $f$ soit de type fini, et que la
propriété énoncée dans (i) soit vérifiée par tout anneau de valuation (resp.
tout anneau de valuation discrète) ayant $R(X)$ pour corps des fractions.
Alors $f$ est propre.
\end{enumerate}
\end{corollary}

Notons d'abord que les hypothèses impliquent dans tous les cas que $f$ est
séparé (\hyperref[I.5.5.9-fr]{I, 5.5.9}).
\begin{enumerate}
\item[\rm (i)] Soient $K=R(Y)$, $L=R(X)$, $y$ un point de $Y$, $A$ un anneau
de valuation ayant $L$ pour corps des fractions et dominant
$\mathcal{O}_y$~; l'injection $\mathcal{O}_y\to A$ définit donc un morphisme
$h$ de $Y'=\operatorname{Spec}(A)$ dans $Y$
(\hyperref[I.2.4.4-fr]{I, 2.4.4}) tel que $h(a)=y$, en désignant par $a$ le
point fermé de $Y'$~; en outre, si $\eta$ est le point générique de $Y$, qui
est aussi celui de $\operatorname{Spec}(\mathcal{O}_y)$, on a $h(b)=\eta$, en
désignant par $b$ le point générique de $Y'$ (puisque $K\subset L$ par
hypothèse). Si $\xi$ est le point générique de $X$, on a $k(\xi)=k(b)=L$ par
hypothèse, d'où un $Y$-morphisme $g:\operatorname{Spec}(L)\to X$ tel que
$g(b)=\xi$~; en vertu de (\hyperref[II.7.3.8-fr]{7.3.8}), $g$ provient d'un
$Y$-morphisme $g':Y'\to X$. Si $x=g'(a)$, il est clair que $A$ domine
$\mathcal{O}_x$.
\item[\rm (ii)] La question étant locale sur $Y$, on peut toujours supposer
que $Y$ est affine (resp. affine et noethérien). Comme $f$ est de type fini,
on peut appliquer dans les deux cas le lemme de Chow
(\hyperref[II.5.6.1-fr]{5.6.1}). Il y a donc un morphisme projectif
$p:P\to Y$, un morphisme d'immersion $j:X'\to P$ et un morphisme projectif,
surjectif et birationnel $g:X'\to X$ ($X'$ étant intègre) tels que le diagramme
\[
\begin{tikzcd}[column sep=large,row sep=large]
P \arrow[d,"p"'] & X' \arrow[l,"j"'] \arrow[d,"g"] \\
Y & X \arrow[l,"f"]
\end{tikzcd}
\]
soit commutatif. Il suffit de prouver que $j$ est une immersion \emph{fermée},
car alors $f\circ g=p\circ j$ sera un morphisme projectif, donc propre, et
comme $g$ est surjectif, $f$ sera aussi propre
(\hyperref[II.5.4.3-fr]{5.4.3}). Soit $Z$ le sous-préschéma fermé réduit de
$P$ ayant $\overline{j(X')}$ comme espace sous-jacent
(\hyperref[I.5.2.1-fr]{I, 5.2.1})~; comme $X'$ est intègre, $j$ se factorise
en $i\circ h$, où $i:Z\to P$ est l'injection canonique, $h:X'\to Z$ une
immersion ouverte dominante (\hyperref[I.5.2.3-fr]{I, 5.2.3}), et $Z$ est
intègre~;

\oldpage[II]{147}
en outre $Z$ est projectif sur $Y$, et on voit qu'on peut se borner au cas où
$P$ est \emph{intègre} et $j$ \emph{dominant et birationnel}, et tout revient à
voir que $j$ est \emph{surjectif}. Or, soit $z\in P$~; $\mathcal{O}_z$ est un
anneau local intègre (resp. intègre et noethérien) dont le corps des fractions
est
\[
L=R(P)=R(X')=R(X).
\]
On peut se borner au cas où $z$ n'est pas le point générique de $P$. Il y a par
suite (\hyperref[II.7.1.2-fr]{7.1.2} et
\hyperref[II.7.1.7-fr]{7.1.7}) un anneau de valuation (resp. un anneau de
valuation discrète) $A$, ayant $L$ pour corps des fractions et dominant
$\mathcal{O}_z$. \emph{A fortiori}, $A$ domine $\mathcal{O}_y$, en posant
$y=p(z)$, et par hypothèse il y a donc un $x\in X$ tel que $A$ domine
$\mathcal{O}_x$. Comme $g$ est propre, la première partie de la démonstration
prouve que $A$ domine aussi un $\mathcal{O}_{x'}$, pour un $x'\in X'$~; il
s'ensuit que $\mathcal{O}_z$ et $\mathcal{O}_{j(x')}=\mathcal{O}_{x'}$ sont
apparentés (\hyperref[I.8.1.4-fr]{I, 8.1.4}), et, comme $P$ est un schéma, cela
entraîne $z=j(x')$ (\hyperref[I.8.2.2-fr]{I, 8.2.2}) et achève la démonstration.
\end{enumerate}

\begin{corollary}[7.3.11]
\label{II.7.3.11-fr}
Soient $X$, $Y$ deux schémas intègres, $f:X\to Y$ un morphisme dominant,
quasi-compact et universellement fermé. Supposons en outre $Y$ affine d'anneau
(intègre) $B$. Alors $\Gamma(X,\mathcal{O}_X)$ est canoniquement isomorphe à
un sous-anneau de la fermeture intégrale de $B$ dans $R(X)$.
\end{corollary}

En effet (\hyperref[I.8.2.1.1-fr]{I, 8.2.1.1}),
$\Gamma(X,\mathcal{O}_X)$ s'identifie à l'intersection des $\mathcal{O}_x$ pour
$x\in X$~; en vertu de (\hyperref[II.7.3.10-fr]{7.3.10}), de
(\hyperref[II.7.1.2-fr]{7.1.2}) et (\hyperref[II.7.1.3-fr]{7.1.3}),
$\Gamma(X,\mathcal{O}_X)$ est par suite contenu dans l'intersection des
anneaux de valuation contenant $B$ et ayant $R(X)$ pour corps des fractions~;
la conclusion résulte alors de (\hyperref[II.7.1.3-fr]{7.1.3}).

\begin{remark}[7.3.12]
\label{II.7.3.12-fr}
Sous les hypothèses de (\hyperref[II.7.3.11-fr]{7.3.11}), et lorsqu'on suppose
que $R(X)$ est une extension de type fini de $R(Y)$, on pourra dans de nombreux
cas conclure que $\Gamma(X,\mathcal{O}_X)$ est un \emph{module de type fini}
sur l'anneau $B=\Gamma(Y,\mathcal{O}_Y)$. Ce sera le cas par exemple lorsque
$B$ est une \emph{algèbre de type fini sur un corps}, car on sait alors que la
fermeture intégrale de $B$ dans une extension de type fini de son corps des
fractions est un $B$-module de type fini ([13], t.~I, p.~267, th.~9)~; la
conclusion résulte alors de (\hyperref[II.7.3.11-fr]{7.3.11}) et du fait que
$B$ est noethérien.

En particulier, \emph{un schéma $X$ propre et affine sur un corps $K$ est
fini}. En effet, en vertu de (\hyperref[II.1.6.4-fr]{1.6.4}),
(\hyperref[II.5.4.6-fr]{5.4.6}) et
(\hyperref[I.6.4.4-fr]{I, 6.4.4, c)}), on peut se borner au cas où $X$ est
\emph{réduit}. En outre, il suffira de prouver que chacun des sous-préschémas
fermés de $X$ ayant pour espace sous-jacent une composante irréductible de $X$
(en nombre fini) est fini sur $K$, si bien que (tenant compte de
(\hyperref[II.5.4.5-fr]{5.4.5})) on est finalement ramené au cas où $X$ est
\emph{intègre}. Mais alors le résultat découle des remarques faites ci-dessus.

Au chapitre III, nous retrouverons cette dernière proposition par d'autres
méthodes et comme conséquence de résultats plus généraux, montrant que si
$f:X\to Y$ est propre et $Y$ localement noethérien, $f_*(\mathcal{F})$ est
cohérent pour tout $\mathcal{O}_X$-Module cohérent $\mathcal{F}$
(\hyperref[III.4.4.2-fr]{III, 4.4.2}).

Notons enfin que le critère (\hyperref[II.7.3.10-fr]{7.3.10}) est pris comme
\emph{définition} des morphismes propres en Géométrie algébrique classique.
Nous ne l'avons mentionné que pour cette raison, le critère
(\hyperref[II.7.3.8-fr]{7.3.8}) semblant plus maniable dans toutes les
applications à notre connaissance.
\end{remark}

\oldpage[II]{148}
\subsection{Courbes algébriques et corps de fonctions de dimension 1}
\label{subsection:II.7.4-fr}

Le but de ce numéro est de montrer comment se formule en langage des schémas
la notion classique de courbe algébrique (telle qu'elle est exposée par
exemple dans le livre de C.~Chevalley [23]). Dans tout le numéro, \emph{on
désigne par $k$ un corps, tous les schémas considérés sont des $k$-schémas de
type fini, et tous les morphismes des $k$-morphismes}.

\begin{proposition}[7.4.1]
\label{II.7.4.1-fr}
Soit $X$ un préschéma de type fini sur $k$ (donc noethérien)~; soient $x_i$
$(1\leq i\leq n)$ les points génériques des composantes irréductibles $X_i$ de
$X$, et soit $K_i=k(x_i)$ $(1\leq i\leq n)$. Les conditions suivantes sont
équivalentes~:
\begin{enumerate}
\item[\rm a)] Chacun des $K_i$ est une extension de $k$ dont le degré de
transcendance est égal à 1.
\item[\rm b)] Pour tout point fermé $x$ de $X$, l'anneau local
$\mathcal{O}_x$ est de dimension 1 (\hyperref[II.7.1.5-fr]{7.1.5}).
\item[\rm c)] Les parties fermées irréductibles de $X$ distinctes des $X_i$
sont les points fermés de $X$.
\end{enumerate}
\end{proposition}

Comme $X$ est quasi-compact, toute partie fermée irréductible $F$ de $X$
contient un point fermé (\hyperref[0.2.1.3-fr]{0, 2.1.3}). En vertu de
(\hyperref[I.2.4.2-fr]{I, 2.4.2}), il y a correspondance biunivoque entre les
idéaux premiers de $\mathcal{O}_x$ et les parties fermées irréductibles de $X$
contenant $x$ (\hyperref[I.1.1.14-fr]{I, 1.1.14})~; l'équivalence de b) et c)
en découle aussitôt. D'autre part, si $\mathfrak{p}_\alpha$
$(1\leq\alpha\leq r)$ sont les idéaux premiers minimaux de l'anneau local
noethérien $\mathcal{O}_x$, les anneaux locaux
$\mathcal{O}_x/\mathfrak{p}_\alpha$ sont intègres, et ont pour corps des
fractions ceux des $K_i$ tels que $x\in X_i$. En outre, on sait ([1],
p.~4-06, th.~2) que la dimension d'une $k$-algèbre locale intègre de type fini
est égale au degré de transcendance sur $k$ de son corps des fractions. Enfin,
la dimension de $\mathcal{O}_x$ est borne supérieure des dimensions des
$\mathcal{O}_x/\mathfrak{p}_\alpha$~; or, la condition a) implique que ces
dimensions sont égales à 1, donc a) implique b)~; inversement, si
$\mathcal{O}_x$ est de dimension 1, aucun des $\mathfrak{p}_\alpha$ ne peut
être égal à l'idéal maximal de $\mathcal{O}_x$, sans quoi $\mathcal{O}_x$
serait de dimension 0~; donc chacun des
$\mathcal{O}_x/\mathfrak{p}_\alpha$ est de dimension 1, ce qui montre que b)
entraîne a).

On notera que sous les conditions de (\hyperref[II.7.4.1-fr]{7.4.1}),
l'ensemble $X$ est \emph{vide ou infini}, comme il résulte aussitôt de
(\hyperref[I.6.4.4-fr]{I, 6.4.4}).

\begin{definition}[7.4.2]
\label{II.7.4.2-fr}
On appelle courbe algébrique sur $k$ un schéma algébrique non vide sur $k$
vérifiant les conditions de (\hyperref[II.7.4.1-fr]{7.4.1}).
\end{definition}

Dans le langage de la dimension, qui sera introduit au chapitre IV, cela
s'exprimera en disant qu'une courbe algébrique sur $k$ est un $k$-schéma
algébrique non vide dont \emph{toutes les composantes irréductibles sont de
dimension 1}.

On notera que si $X$ est une courbe algébrique sur $k$, les sous-préschémas
fermés réduits $X_i$ $(1\leq i\leq n)$ de $X$ ayant pour espaces sous-jacents
les composantes irréductibles de $X$ sont des courbes algébriques sur $k$.

\begin{corollary}[7.4.3]
\label{II.7.4.3-fr}
Soit $X$ une courbe algébrique irréductible. Le seul point non fermé de $X$ est
son point générique. Les parties fermées de $X$ distinctes de $X$ sont les
ensembles finis de points fermés~; ce sont aussi les seules parties de $X$ qui
ne soient pas partout denses.
\end{corollary}

Si un point $x\in X$ n'est pas fermé, son adhérence dans $X$ est une partie
fermée irréductible de $X$, donc nécessairement $X$ tout entier en vertu de
(\hyperref[II.7.4.1-fr]{7.4.1}), et par suite $x$ est le point générique de
$X$. Une partie fermée $F$ de $X$, distincte de $X$, ne peut contenir

\oldpage[II]{149}
le point générique de $X$, donc tous ses points sont fermés (dans $X$ et
\emph{a fortiori} dans $F$)~; en considérant le sous-préschéma fermé réduit de
$X$ ayant $F$ pour espace sous-jacent (\hyperref[I.5.2.1-fr]{I, 5.2.1}), il
résulte donc de (\hyperref[I.6.2.2-fr]{I, 6.2.2}) que $F$ est fini et discret.
L'adhérence dans $X$ d'une partie infinie de $X$ est donc nécessairement égale
à $X$.

Lorsque $X$ est une courbe algébrique quelconque, en appliquant
(\hyperref[II.7.4.3-fr]{7.4.3}) aux composantes irréductibles de $X$, on voit
que les seuls points non fermés de $X$ sont les points génériques de ces
composantes.

\begin{corollary}[7.4.4]
\label{II.7.4.4-fr}
Soient $X$ et $Y$ deux courbes algébriques irréductibles sur $k$, $f:X\to Y$
un $k$-morphisme. Pour que $f$ soit dominant, il faut et il suffit que
$f^{-1}(y)$ soit fini pour tout $y\in Y$.
\end{corollary}

En effet, si $f$ n'est pas dominant, $f(X)$ est nécessairement une partie
\emph{finie} de $Y$ en vertu de (\hyperref[II.7.4.3-fr]{7.4.3}), donc il n'est
pas possible que $f^{-1}(y)$ soit fini pour tout point de $Y$, sinon $X$ serait
fini, ce qui est absurde (\hyperref[II.7.4.1-fr]{7.4.1}). Inversement, si $f$
est dominant, pour tout $y\in Y$ distinct du point générique $\eta$ de $Y$,
$f^{-1}(y)$ est fermé dans $X$ puisque $\{y\}$ est fermé dans $Y$
(\hyperref[II.7.4.3-fr]{7.4.3})~; d'autre part, par hypothèse, $f^{-1}(y)$ ne
contient pas le point générique $\xi$ de $X$, donc est fini en vertu de
(\hyperref[II.7.4.3-fr]{7.4.3}). Enfin, pour voir que lorsque $f$ est dominant,
$f^{-1}(\eta)$ est fini, on note que la fibre $f^{-1}(\eta)$ est un schéma
irréductible de type fini sur $k(\eta)$, et de point générique $\xi$
(\hyperref[I.6.3.9-fr]{I, 6.3.9} et
\hyperref[I.6.4.11-fr]{6.4.11}). Comme $k(\xi)$ et $k(\eta)$ sont des
extensions de type fini de $k$, de même degré de transcendance 1, $k(\xi)$ est
nécessairement une extension de degré fini de $k(\eta)$, donc $\xi$ est fermé
dans $f^{-1}(\eta)$ (\hyperref[I.6.4.2-fr]{I, 6.4.2}), et $f^{-1}(\eta)$ est
par suite réduit au point $\xi$.

Nous verrons au chapitre III qu'un morphisme \emph{propre} $f:X\to Y$ de
préschémas noethériens, tel que $f^{-1}(y)$ soit fini pour tout $y\in Y$, est
nécessairement \emph{fini}~; il s'ensuivra donc de
(\hyperref[II.7.4.4-fr]{7.4.4}) qu'un morphisme dominant propre d'une courbe
algébrique irréductible dans une courbe algébrique est \emph{fini}.

\begin{corollary}[7.4.5]
\label{II.7.4.5-fr}
Soit $X$ une courbe algébrique sur $k$. Pour que $X$ soit régulière, il faut et
il suffit que $X$ soit normale, ou encore que les anneaux locaux de ses points
fermés soient des anneaux de valuation discrète.
\end{corollary}

Cela résulte aussitôt de la condition b) de
(\hyperref[II.7.4.1-fr]{7.4.1}) et de
(\hyperref[II.7.1.6-fr]{7.1.6}).

\begin{corollary}[7.4.6]
\label{II.7.4.6-fr}
Soient $X$ une courbe algébrique réduite, $\mathcal{A}$ une
$\mathcal{R}(X)$-Algèbre cohérente réduite~; alors la fermeture intégrale $X'$
de $X$ relativement à $\mathcal{A}$ (\hyperref[II.6.3.4-fr]{6.3.4}) est une
courbe algébrique normale, et le morphisme canonique $X'\to X$ est fini.
\end{corollary}

Le fait que $X'\to X$ soit fini résulte de
(\hyperref[II.6.3.10-fr]{6.3.10})~; $X'$ est donc un $k$-schéma algébrique~;
en outre, si $x_i$ $(1\leq i\leq n)$ sont les points génériques des composantes
irréductibles de $X$, $x'_j$ $(1\leq j\leq m)$ ceux des composantes
irréductibles de $X'$, chacun des $k(x'_j)$ est extension algébrique finie de
l'un des $k(x_i)$ (\hyperref[II.6.3.6-fr]{6.3.6}), donc a un degré de
transcendance 1 sur $k$. $X'$ est donc bien une courbe algébrique sur $k$, et
en outre on sait que $X'$ est somme d'un nombre fini de schémas intègres et
normaux (\hyperref[II.6.3.6-fr]{6.3.6} et
\hyperref[II.6.3.7-fr]{6.3.7}).

\begin{env}[7.4.7]
\label{II.7.4.7-fr}
On dit qu'une courbe algébrique $X$ sur $k$ est \emph{complète} si elle est
\emph{propre} sur $k$.
\end{env}

\begin{corollary}[7.4.8]
\label{II.7.4.8-fr}
Pour qu'une courbe algébrique réduite $X$ sur $k$ soit complète, il faut et il
suffit que sa normalisée $X'$ le soit.
\end{corollary}

\oldpage[II]{150}
En effet, le morphisme canonique $f:X'\to X$ est alors fini
(\hyperref[II.7.4.6-fr]{7.4.6}), donc propre
(\hyperref[II.6.1.11-fr]{6.1.11}) et surjectif
(\hyperref[II.6.3.8-fr]{6.3.8})~; si $g:X\to\operatorname{Spec}(k)$ est le
morphisme structural, $g$ et $g\circ f$ sont donc propres simultanément,
comme il résulte de (\hyperref[II.5.4.2-fr]{5.4.2, (iii)}) et
(\hyperref[II.5.4.3-fr]{5.4.3, (ii)}), $g$ étant séparé par hypothèse.

\begin{proposition}[7.4.9]
\label{II.7.4.9-fr}
Soient $X$ une courbe algébrique normale sur $k$, $Y$ un $k$-schéma algébrique
propre sur $k$. Toute $k$-application rationnelle de $X$ dans $Y$ est alors
partout définie, autrement dit est un morphisme.
\end{proposition}

En effet, il résulte de (\hyperref[II.7.3.7-fr]{7.3.7}) qu'aux points $x\in X$
où une telle application n'est pas définie, la dimension de $\mathcal{O}_x$
devrait être $\geq2$, donc l'ensemble de ces points est vide~; la dernière
assertion provient de (\hyperref[I.7.2.3-fr]{I, 7.2.3}).

\begin{corollary}[7.4.10]
\label{II.7.4.10-fr}
Une courbe algébrique normale sur $k$ est quasi-projective sur $k$.
\end{corollary}

Comme $X$ est somme d'un nombre fini de courbes algébriques intègres et
normales (\hyperref[II.6.3.8-fr]{6.3.8}), on peut se borner au cas où $X$ est
intègre (\hyperref[II.5.3.6-fr]{5.3.6}). Comme $X$ est quasi-compact, il est
recouvert par un nombre fini d'ouverts affines $U_i$ $(1\leq i\leq n)$, et
comme chacun de ces derniers est de type fini sur $k$, il existe un entier
$n_i$ et une $k$-immersion $f_i:U_i\to\mathbf{P}_k^{n_i}$
(\hyperref[II.5.3.3-fr]{5.3.3} et
\hyperref[II.5.3.4-fr]{5.3.4, (i)}). Comme $U_i$ est dense dans $X$, il résulte
de (\hyperref[II.7.4.9-fr]{7.4.9}) que $f_i$ se prolonge en un $k$-morphisme
$g_i:X\to\mathbf{P}_k^{n_i}$, d'où un $k$-morphisme
$g=(g_1,\ldots,g_n)_k$ de $X$ dans le produit $P$ des
$\mathbf{P}_k^{n_i}$ sur $k$. En outre, pour chaque indice $i$, comme la
restriction de $g_i$ à $U_i$ est une immersion, il en est de même de la
restriction de $g$ à $U_i$ (\hyperref[I.5.3.14-fr]{I, 5.3.14}). Comme les
$U_i$ recouvrent $X$ et que $g$ est séparé
(\hyperref[I.5.5.1-fr]{I, 5.5.1, (v)}), $g$ est une immersion de $X$ dans $P$
(\hyperref[I.8.2.8-fr]{I, 8.2.8}). Comme le morphisme de Segre
(\hyperref[II.4.3.3-fr]{4.3.3}) fournit une immersion de $P$ dans un
$\mathbf{P}_k^N$, cela achève de prouver que $X$ est quasi-projectif.

\begin{corollary}[7.4.11]
\label{II.7.4.11-fr}
Une courbe algébrique normale $X$ est isomorphe au schéma induit sur un ouvert
partout dense d'une courbe algébrique normale et complète $\widehat{X}$,
déterminée à un isomorphisme unique près.
\end{corollary}

Si $X_1$, $X_2$ sont deux courbes normales et complètes, il résulte aussitôt
de (\hyperref[II.7.4.9-fr]{7.4.9}) que tout isomorphisme d'un ouvert $U_1$
dense dans $X_1$, sur un ouvert $U_2$ dense dans $X_2$, se prolonge de façon
unique en un isomorphisme de $X_1$ sur $X_2$~; d'où l'assertion d'unicité. Pour
prouver l'existence de $\widehat{X}$, il suffit de remarquer que l'on peut
considérer $X$ comme un sous-schéma d'un fibré projectif $\mathbf{P}_k^n$
(\hyperref[II.7.4.10-fr]{7.4.10}). Soit $\overline{X}$ l'adhérence de $X$ dans
$\mathbf{P}_k^n$ (\hyperref[I.9.5.11-fr]{I, 9.5.11})~; comme $X$ est induit
par $\overline{X}$ sur un ouvert dense dans $\overline{X}$
(\hyperref[I.9.5.10-fr]{I, 9.5.10}), les points génériques $x_i$ des
composantes irréductibles de $X$ sont aussi ceux des composantes irréductibles
de $\overline{X}$, et les $k(x_i)$ sont les mêmes pour ces deux schémas, donc
(\hyperref[II.7.4.1-fr]{7.4.1}) $\overline{X}$ est une courbe algébrique sur
$k$, qui est réduite (\hyperref[I.9.5.9-fr]{I, 9.5.9}) et projective sur $k$
(\hyperref[II.5.5.1-fr]{5.5.1}), donc complète
(\hyperref[II.5.5.3-fr]{5.5.3}). Prenons alors pour $\widehat{X}$ la
\emph{normalisée} de $\overline{X}$, qui est encore complète
(\hyperref[II.7.4.8-fr]{7.4.8})~; en outre, si
$h:\widehat{X}\to\overline{X}$ est le morphisme canonique, la restriction de
$h$ à $h^{-1}(X)$ est un isomorphisme sur $X$ puisque $X$ est normale
(\hyperref[II.6.3.4-fr]{6.3.4}), et comme $h^{-1}(X)$ contient les points
génériques des composantes irréductibles de $\widehat{X}$
(\hyperref[II.6.3.8-fr]{6.3.8}), il est dense dans $\widehat{X}$, ce qui achève
la démonstration.

\oldpage[II]{151}

\begin{remark}[7.4.12]
\label{II.7.4.12-fr}
Nous montrerons au chapitre V que la conclusion de
(\hyperref[II.7.4.10-fr]{7.4.10}) est encore valable sans supposer la courbe
normale (ni même réduite)~; nous montrerons aussi que pour qu'une courbe
algébrique (réduite ou non) soit affine, il faut et il suffit que ses
composantes irréductibles (réduites) ne soient pas complètes.
\end{remark}

\begin{corollary}[7.4.13]
\label{II.7.4.13-fr}
Soient $X$ une courbe normale irréductible de corps $R(X)=K$, $Y$ une courbe
intègre complète de corps $L=R(Y)$. Il y a une correspondance biunivoque
canonique entre les $k$-morphismes dominants $X\to Y$ et les
$k$-monomorphismes $L\to K$.
\end{corollary}

En vertu de (\hyperref[II.7.4.9-fr]{7.4.9}), les $k$-applications rationnelles
de $X$ dans $Y$ s'identifient aux $k$-morphismes $u:X\to Y$. Les morphismes
$u$ dominants étant caractérisés par le fait que $u(x)=y$ (en désignant par
$x$ et $y$ les points génériques respectifs de $X$ et de $Y$), le corollaire
résulte de ces remarques et de (\hyperref[I.7.1.13-fr]{I, 7.1.13}).

\begin{env}[7.4.14]
\label{II.7.4.14-fr}
On peut préciser le résultat de (\hyperref[II.7.4.13-fr]{7.4.13}) lorsqu'on
prend pour $Y$ la \emph{droite projective}
$\mathbf{P}_k^1=\operatorname{Proj}(k[T_0,T])$, $T_0$ et $T$ étant deux
indéterminées. C'est bien là un schéma intègre
(\hyperref[II.2.4.4-fr]{2.4.4}), et le schéma induit sur l'ouvert $D_+(T_0)$
de $Y$ est isomorphe à $\operatorname{Spec}(k[T])$
(\hyperref[II.2.3.6-fr]{2.3.6}), donc le point générique de $Y$ est l'idéal
$(0)$ de $k[T]$ et son corps de fonctions rationnelles $k(T)$, ce qui montre
que $Y$ est une courbe algébrique complète sur $k$. En outre, le seul idéal
premier gradué de $S=k[T_0,T]$ contenant $T_0$ et distinct de $S_+$ est
l'idéal principal $(T_0)$, donc le complémentaire de $D_+(T_0)$ dans
$Y=\mathbf{P}_k^1$ est réduit à \emph{un point fermé}, dit «~point à
l'infini~» que nous noterons $\infty$ (pour une étude générale des rapports
entre fibrés vectoriels et fibrés projectifs, voir
\hyperref[subsection:II.8.4-fr]{8.4}). Avec ces notations~:
\end{env}

\begin{corollary}[7.4.15]
\label{II.7.4.15-fr}
Soit $X$ une courbe normale irréductible de corps $R(X)=K$. Il existe une
application bijective canonique de $K$ sur l'ensemble des morphismes $u$ de
$X$ dans $\mathbf{P}_k^1$, distincts du morphisme constant de valeur
$\infty$. Pour que $u$ soit dominant, il faut et il suffit que l'élément
correspondant de $K$ soit transcendant sur $k$.
\end{corollary}

Cet énoncé résulte immédiatement de (\hyperref[II.7.4.9-fr]{7.4.9}) et du

\begin{lemma}[7.4.15.1]
\label{II.7.4.15.1-fr}
Soit $X$ un préschéma intègre sur $k$, et soit $K=R(X)$ son corps des fonctions
rationnelles. Il existe une application bijective canonique de l'ensemble
$K$ sur l'ensemble des applications rationnelles $u$ de $X$ dans
$\mathbf{P}_k^1$, distinctes du morphisme constant de valeur $\infty$. Pour
qu'une telle application rationnelle soit dominante, il faut et il suffit que
l'élément correspondant de $K$ soit transcendant sur $k$.
\end{lemma}

Tout d'abord, les applications rationnelles de $X$ dans $\mathbf{P}_k^1$
correspondent biunivoquement aux points de $\mathbf{P}_k^1$ à valeurs dans
l'extension $K$ de $k$ (\hyperref[I.7.1.12-fr]{I, 7.1.12}). Si un tel point
est localisé (\hyperref[I.3.4.5-fr]{I, 3.4.5}) au point générique de
$\mathbf{P}_k^1$, l'application rationnelle correspondante est évidemment
dominante. Dans le cas contraire, comme tout point de $\mathbf{P}_k^1$
distinct du point générique est fermé (\hyperref[II.7.4.3-fr]{7.4.3}),
l'image du domaine de définition $U$ de $u$ par l'unique morphisme
$U\to\mathbf{P}_k^1$ de la classe $u$ (\hyperref[I.7.2.2-fr]{I, 7.2.2}),
est réduite à un point fermé $y$ de $\mathbf{P}_k^1$, et ce morphisme (qui
n'est pas nécessairement partout défini dans $X$) n'est donc pas dominant~;
par abus de langage, on dit alors que l'application rationnelle $u$ est
«~constante, de valeur $y$~». Il reste à mettre en correspondance biunivoque
les points de $\mathbf{P}_k^1$ à valeurs dans $K$, de localité
(\hyperref[I.3.4.5-fr]{I, 3.4.5}) distincte de $\infty$, et les éléments de
$K$, et à vérifier que la localité d'un tel point est le point générique de
$\mathbf{P}_k^1$ si et seulement s'il correspond à un

\oldpage[II]{152}
élément transcendant sur $k$. Or, cette vérification est immédiate sur
(\hyperref[II.4.2.6-fr]{4.2.6, exemple 1\textsuperscript{o}}).

\begin{corollary}[7.4.16]
\label{II.7.4.16-fr}
Soient $X$, $Y$ deux courbes algébriques sur $k$, normales, complètes et
irréductibles~; soient $K=R(X)$, $L=R(Y)$ leurs corps. Il existe une
correspondance biunivoque canonique entre l'ensemble des $k$-isomorphismes
$X\xrightarrow{\sim}Y$ et l'ensemble des $k$-isomorphismes
$L\xrightarrow{\sim}K$.
\end{corollary}

C'est une conséquence évidente de (\hyperref[II.7.4.13-fr]{7.4.13}).

\begin{env}[7.4.17]
\label{II.7.4.17-fr}
Le corollaire (\hyperref[II.7.4.16-fr]{7.4.16}) montre qu'une courbe algébrique
sur $k$, normale, complète et irréductible, est \emph{déterminée par son corps
de fonctions rationnelles $K$ à un isomorphisme unique près}~; par définition,
$K$ est une extension de type fini de $k$, de degré de transcendance $1$ (ce
qu'on appelle classiquement un \emph{corps de fonctions algébriques d'une
variable}). De plus~:
\end{env}

\begin{proposition}[7.4.18]
\label{II.7.4.18-fr}
Pour toute extension $K$ de $k$, de type fini et de degré de transcendance
$1$, il existe une courbe algébrique normale, complète et irréductible $X$
telle que $R(X)=K$ (déterminée à isomorphisme unique près). L'ensemble des
anneaux locaux de $X$ s'identifie (\hyperref[I.8.2.1-fr]{I, 8.2.1}) à
l'ensemble formé de $K$ et des anneaux de valuation contenant $k$ et ayant
$K$ comme corps des fractions.
\end{proposition}

En effet, $K$ est une extension de degré fini d'une extension transcendante
pure $k(T)$ de $k$, qui s'identifie, comme on l'a vu, au corps des fonctions
rationnelles de la droite projective $Y=\mathbf{P}_k^1$. Soit $X$ la
\emph{fermeture intégrale} de $Y$ relativement à $K$
(\hyperref[II.6.3.4-fr]{6.3.4})~; $X$ est une courbe algébrique normale de
corps $K$ (\hyperref[II.6.3.7-fr]{6.3.7}), et elle est complète puisque le
morphisme $X\to Y$ est fini (\hyperref[II.7.4.6-fr]{7.4.6}). Les anneaux
locaux $\mathcal{O}_x$ de $X$ sont~: le corps $K$ lorsque $x$ est le point
générique~; si $x$ est distinct du point générique, $\mathcal{O}_x$ est un
anneau de valuation discrète contenant $k$ et ayant $K$ comme corps des
fractions (\hyperref[II.7.4.5-fr]{7.4.5}). Inversement, soit $A$ un tel
anneau~; comme le morphisme $X\to\operatorname{Spec}(k)$ est propre et que
$A$ domine $k$, il domine un anneau local $\mathcal{O}_x$ de $X$
(\hyperref[II.7.3.10-fr]{7.3.10})~; ce dernier étant un anneau de valuation
ayant $K$ comme corps des fractions, est donc nécessairement égal à $A$.

\begin{namedtheorem}[7.4.19]{Remarques}
\label{II.7.4.19-fr}
Il résulte de (\hyperref[II.7.4.16-fr]{7.4.16}) et
(\hyperref[II.7.4.18-fr]{7.4.18}) que \emph{la donnée d'une courbe algébrique
sur $k$, normale, complète et irréductible, est essentiellement équivalente à
la donnée d'une extension $K$ de $k$, de type fini et de degré de
transcendance $1$}. On notera que si $k'$ est une extension du corps de base
$k$, $X\otimes_k k'$ sera encore une courbe algébrique complète sur $k'$
(\hyperref[II.5.4.2-fr]{5.4.2, (iii)}), mais en général elle ne sera ni réduite
ni irréductible. Elle le sera cependant si $K$ est une extension séparable de
$k$ et si $k$ est algébriquement fermé dans $K$ (ce qui s'exprime, dans une
terminologie classique que nous ne suivrons pas, en disant que $K$ est une
«~extension régulière~» de $k$). Mais même dans ce cas, il peut se faire que
$X\otimes_k k'$ ne soit pas normale. Le lecteur trouvera des détails sur ces
questions dans le chapitre IV.
\end{namedtheorem}

\section{Schémas éclatés~; cônes projetants~; fermeture projective}
\label{section:II.8-fr}

\subsection{Préschémas éclatés}
\label{subsection:II.8.1-fr}

\begin{env}[8.1.1]
\label{II.8.1.1-fr}
Soient $Y$ un préschéma, et pour tout entier $n\geq0$, soit $\mathcal{I}_n$ un
Idéal quasi-cohérent de $\mathcal{O}_Y$~; on suppose vérifiées les conditions
suivantes~:
\begin{equation}
\label{II.8.1.1.1-fr}
  \mathcal{I}_0=\mathcal{O}_Y,\qquad
  \mathcal{I}_n\subset\mathcal{I}_m\quad\text{pour }m\leq n
\tag{8.1.1.1}
\end{equation}
\begin{equation}
\label{II.8.1.1.2-fr}
  \mathcal{I}_m\mathcal{I}_n\subset\mathcal{I}_{m+n}
  \quad\text{quels que soient }m,n.
\tag{8.1.1.2}
\end{equation}

\oldpage[II]{153}
On notera que ces hypothèses entraînent
\begin{equation}
\label{II.8.1.1.3-fr}
  \mathcal{I}_1^n\subset\mathcal{I}_n.
\tag{8.1.1.3}
\end{equation}

Posons
\begin{equation}
\label{II.8.1.1.4-fr}
  \mathcal{S}=\bigoplus_{n\geq0}\mathcal{I}_n.
\tag{8.1.1.4}
\end{equation}

Il résulte de (\hyperref[II.8.1.1.1-fr]{8.1.1.1}) et
(\hyperref[II.8.1.1.2-fr]{8.1.1.2}) que $\mathcal{S}$ est une
$\mathcal{O}_Y$-Algèbre graduée quasi-cohérente, et définit donc un $Y$-schéma
$X=\operatorname{Proj}(\mathcal{S})$. Si $\mathcal{J}$ est un Idéal
\emph{inversible} de $\mathcal{O}_Y$,
$\mathcal{I}_n\otimes_{\mathcal{O}_Y}\mathcal{J}^{\otimes n}$ s'identifie
canoniquement à $\mathcal{I}_n\mathcal{J}^n$. Si on remplace donc les
$\mathcal{I}_n$ par les $\mathcal{I}_n\mathcal{J}^n$, ce qui remplace
$\mathcal{S}$ par une $\mathcal{O}_Y$-Algèbre quasi-cohérente
$\mathcal{S}_{(\mathcal{J})}$,
$X_{(\mathcal{J})}=\operatorname{Proj}(\mathcal{S}_{(\mathcal{J})})$ est
canoniquement isomorphe à $X$ (\hyperref[II.3.1.8-fr]{3.1.8}).
\end{env}

\begin{env}[8.1.2]
\label{II.8.1.2-fr}
Supposons $Y$ \emph{localement intègre}, de sorte que le faisceau
$\mathcal{R}(Y)$ des fonctions rationnelles est une $\mathcal{O}_Y$-Algèbre
quasi-cohérente (\hyperref[I.7.3.7-fr]{I, 7.3.7}). Nous dirons qu'un
sous-$\mathcal{O}_Y$-Module $\mathcal{I}$ de $\mathcal{R}(Y)$ est un
\emph{Idéal fractionnaire} de $\mathcal{R}(Y)$ s'il est de \emph{type fini}
(\hyperref[0.5.2.1-fr]{0, 5.2.1}). Supposons donné, pour tout $n\geq0$, un
Idéal fractionnaire quasi-cohérent $\mathcal{I}_n$ de $\mathcal{R}(Y)$, tel
que $\mathcal{I}_0=\mathcal{O}_Y$, et que la condition
(\hyperref[II.8.1.1.2-fr]{8.1.1.2}) soit vérifiée (mais non nécessairement la
seconde condition (\hyperref[II.8.1.1.1-fr]{8.1.1.1}))~; on peut encore alors
définir par la formule (\hyperref[II.8.1.1.4-fr]{8.1.1.4}) une
$\mathcal{O}_Y$-Algèbre graduée quasi-cohérente, et le $Y$-schéma correspondant
$X=\operatorname{Proj}(\mathcal{S})$~; on aura encore un isomorphisme
canonique de $X$ sur $X_{(\mathcal{J})}$ pour tout Idéal fractionnaire
inversible $\mathcal{J}$ de $\mathcal{R}(Y)$.
\end{env}

\begin{definition}[8.1.3]
\label{II.8.1.3-fr}
Soit $Y$ un préschéma (resp. un préschéma localement intègre), et soit
$\mathcal{I}$ un Idéal quasi-cohérent de $\mathcal{O}_Y$ (resp. un Idéal
fractionnaire quasi-cohérent de $\mathcal{R}(Y)$). On dit que le $Y$-schéma
$X=\operatorname{Proj}(\bigoplus_{n\geq0}\mathcal{I}^n)$ est obtenu en
\emph{faisant éclater l'Idéal $\mathcal{I}$}, ou est le \emph{préschéma éclaté
de $Y$ relativement à $\mathcal{I}$}. Lorsque $\mathcal{I}$ est un Idéal
quasi-cohérent de $\mathcal{O}_Y$ et $Y'$ le sous-préschéma fermé de $Y$ défini
par $\mathcal{I}$, on dit aussi que $X$ est le $Y$-schéma obtenu en faisant
éclater $Y'$.
\end{definition}

Par définition, $\mathcal{S}=\bigoplus_{n\geq0}\mathcal{I}^n$ est alors
engendrée par $\mathcal{S}_1=\mathcal{I}$~; si $\mathcal{I}$ est un
$\mathcal{O}_Y$-Module de \emph{type fini}, $X$ est donc \emph{projectif} sur
$Y$ (\hyperref[II.5.5.2-fr]{5.5.2}). Sans hypothèse sur $\mathcal{I}$, le
$\mathcal{O}_X$-Module $\mathcal{O}_X(1)$ est \emph{inversible}
(\hyperref[II.3.2.5-fr]{3.2.5}) et \emph{très ample} en vertu de
(\hyperref[II.4.4.3-fr]{4.4.3}) pour le morphisme structural $X\to Y$.

On notera que si $f:X\to Y$ est le morphisme structural, la restriction de
$f$ à $f^{-1}(Y-Y')$ est un \emph{isomorphisme} sur $Y-Y'$ lorsque
$\mathcal{I}$ est un Idéal de $\mathcal{O}_Y$ et $Y'$ le sous-préschéma fermé
qu'il définit~: en effet, la question étant locale sur $Y$, il suffit de
supposer $\mathcal{I}=\mathcal{O}_Y$, et notre assertion résulte de
(\hyperref[II.3.1.7-fr]{3.1.7}).

Lorsqu'on remplace $\mathcal{I}$ par $\mathcal{I}^d$ ($d>0$), le $Y$-schéma
éclaté $X$ est remplacé par un $Y$-schéma canoniquement isomorphe $X'$
(\hyperref[II.8.1.1-fr]{8.1.1})~; de même, pour tout Idéal (resp. Idéal
fractionnaire) inversible $\mathcal{J}$, le préschéma éclaté
$X_{(\mathcal{J})}$ relativement à l'Idéal $\mathcal{I}\mathcal{J}$ est
canoniquement isomorphe à $X$ (\hyperref[II.8.1.1-fr]{8.1.1}).

En particulier, lorsque $\mathcal{I}$ est un Idéal (resp. un Idéal
fractionnaire) \emph{inversible}, le $Y$-schéma obtenu en faisant éclater
$\mathcal{I}$ est \emph{isomorphe} à $Y$ (\hyperref[II.3.1.7-fr]{3.1.7}).

\begin{proposition}[8.1.4]
\label{II.8.1.4-fr}
Soit $Y$ un préschéma intègre.
\begin{enumerate}
\item[\rm(i)] Pour toute suite $(\mathcal{I}_n)$ d'Idéaux fractionnaires
quasi-cohérents de $\mathcal{R}(Y)$ vérifiant
(\hyperref[II.8.1.1.2-fr]{8.1.1.2})

\oldpage[II]{154}
et tels que $\mathcal{I}_0=\mathcal{O}_Y$, le $Y$-schéma
$X=\operatorname{Proj}(\bigoplus_{n\geq0}\mathcal{I}_n)$ est intègre et le
morphisme structural $f:X\to Y$ dominant.
\item[\rm(ii)] Soit $\mathcal{I}$ un Idéal fractionnaire quasi-cohérent de
$\mathcal{R}(Y)$, et soit $X$ le $Y$-schéma éclaté de $Y$ relativement à
$\mathcal{I}$. Si $\mathcal{I}\neq0$, le morphisme structural $f:X\to Y$ est
alors birationnel et surjectif.
\end{enumerate}
\end{proposition}

\begin{enumerate}
\item[\rm(i)] résulte de ce que $\mathcal{S}=\bigoplus_{n\geq0}\mathcal{I}_n$
est une $\mathcal{O}_Y$-Algèbre \emph{intègre}
(\hyperref[II.3.1.12-fr]{3.1.12} et \hyperref[II.3.1.14-fr]{3.1.14}) puisque
pour tout $y\in Y$, $\mathcal{O}_y$ est un anneau intègre
(\hyperref[I.5.1.4-fr]{I, 5.1.4}).
\item[\rm(ii)] D'après (i), $X$ est intègre~; si en outre, $x$ et $y$ sont les
points génériques de $X$ et $Y$, on a $f(x)=y$, et il faut prouver que $k(x)$
est de rang $1$ sur $k(y)$. Or $x$ est aussi le point générique de la fibre
$f^{-1}(y)$~; si $\psi$ est le morphisme canonique $Z\to Y$, où
$Z=\operatorname{Spec}(k(y))$, le préschéma $f^{-1}(y)$ s'identifie à
$\operatorname{Proj}(\mathcal{S}')$, où
$\mathcal{S}'=\psi^*(\mathcal{S})$ (\hyperref[II.3.5.3-fr]{3.5.3}). Mais il est
clair que $\mathcal{S}'=\bigoplus_{n\geq0}(\mathcal{I}_y)^n$, et comme
$\mathcal{I}$ est un Idéal fractionnaire quasi-cohérent non nul de
$\mathcal{R}(Y)$, $\mathcal{I}_y\neq0$
(\hyperref[I.7.3.6-fr]{I, 7.3.6}), d'où $\mathcal{I}_y=k(y)$~;
$\operatorname{Proj}(\mathcal{S}')$ s'identifie donc à
$\operatorname{Spec}(k(y))$ (\hyperref[II.3.1.7-fr]{3.1.7}), d'où la
conclusion.
\end{enumerate}

Nous démontrerons une \emph{réciproque} de
(\hyperref[II.8.1.4-fr]{8.1.4}) dans
(\hyperref[III.2.3.8-fr]{III, 2.3.8}).

\begin{env}[8.1.5]
\label{II.8.1.5-fr}
Revenons à la situation et aux notations de
(\hyperref[II.8.1.1-fr]{8.1.1}). Par définition, les homomorphismes
d'injection $\mathcal{I}_{n+1}\to\mathcal{I}_n$
(\hyperref[II.8.1.1.1-fr]{8.1.1.1}) définissent pour chaque $k\in\mathbf{Z}$ un
homomorphisme injectif de degré $0$ de $\mathcal{S}$-Modules gradués
\begin{equation}
\label{II.8.1.5.1-fr}
  u_k:\mathcal{S}_+(k+1)\to\mathcal{S}(k)~;
\tag{8.1.5.1}
\end{equation}
comme $\mathcal{S}_+(k+1)$ et $\mathcal{S}(k+1)$ sont canoniquement
(TN)-isomorphes, il correspond canoniquement à $u_k$ un homomorphisme injectif
de $\mathcal{O}_X$-Modules (\hyperref[II.3.4.2-fr]{3.4.2})~:
\begin{equation}
\label{II.8.1.5.2-fr}
  \widetilde{u}_k:\mathcal{O}_X(k+1)\to\mathcal{O}_X(k).
\tag{8.1.5.2}
\end{equation}

Rappelons d'autre part (\hyperref[II.3.2.6-fr]{3.2.6}) que l'on a défini des
homomorphismes canoniques
\begin{equation}
\label{II.8.1.5.3-fr}
  \lambda:\mathcal{O}_X(h)\otimes_{\mathcal{O}_X}\mathcal{O}_X(k)
  \to\mathcal{O}_X(h+k)
\tag{8.1.5.3}
\end{equation}
et comme le diagramme
\[
\xymatrix{
  \mathcal{S}(h)\otimes_{\mathcal{S}}\mathcal{S}(k)
    \otimes_{\mathcal{S}}\mathcal{S}(l)\ar[r]\ar[d]
  & \mathcal{S}(h+k)\otimes_{\mathcal{S}}\mathcal{S}(l)\ar[d]\\
  \mathcal{S}(h)\otimes_{\mathcal{S}}\mathcal{S}(k+l)\ar[r]
  & \mathcal{S}(h+k+l)
}
\]
est commutatif, il résulte de la fonctorialité des $\lambda$
(\hyperref[II.3.2.6-fr]{3.2.6}) que les homomorphismes
(\hyperref[II.8.1.5.3-fr]{8.1.5.3}) définissent sur
\begin{equation}
\label{II.8.1.5.4-fr}
  \mathcal{S}_X=\bigoplus_{n\in\mathbf{Z}}\mathcal{O}_X(n)
\tag{8.1.5.4}
\end{equation}

\oldpage[II]{155}
une structure de $\mathcal{O}_X$-Algèbre graduée quasi-cohérente. En outre, le
diagramme
\[
\xymatrix{
  \mathcal{S}(h)\otimes_{\mathcal{S}}\mathcal{S}_+(k+1)\ar[r]\ar[d]_{1\otimes u_k}
  & \mathcal{S}_+(h+k+1)\ar[d]^{u_{k+h}}\\
  \mathcal{S}(h)\otimes_{\mathcal{S}}\mathcal{S}(k)\ar[r]
  & \mathcal{S}(h+k)
}
\]
est commutatif~; la fonctorialité des $\lambda$ montre alors que l'on a un
diagramme commutatif
\begin{equation}
\label{II.8.1.5.5-fr}
\xymatrix{
  \mathcal{O}_X(h)\otimes_{\mathcal{O}_X}\mathcal{O}_X(k+1)
    \ar[r]^{\lambda}\ar[d]_{1\otimes\widetilde{u}_k}
  & \mathcal{O}_X(h+k+1)\ar[d]^{\widetilde{u}_{k+h}}\\
  \mathcal{O}_X(h)\otimes_{\mathcal{O}_X}\mathcal{O}_X(k)
    \ar[r]^{\lambda}
  & \mathcal{O}_X(h+k)
}
\tag{8.1.5.5}
\end{equation}
les flèches horizontales étant les homomorphismes canoniques. On peut donc
dire que les $\widetilde{u}_k$ définissent un \emph{homomorphisme injectif} (de
degré $0$) \emph{de $\mathcal{S}_X$-Modules gradués}
\begin{equation}
\label{II.8.1.5.6-fr}
  \widetilde{u}:\mathcal{S}_X(1)\to\mathcal{S}_X.
\tag{8.1.5.6}
\end{equation}
\end{env}

\begin{env}[8.1.6]
\label{II.8.1.6-fr}
Gardant les notations de (\hyperref[II.8.1.5-fr]{8.1.5}), remarquons maintenant
que, pour $n\geq0$, l'homomorphisme composé
$\widetilde{v}_n=\widetilde{u}_{n-1}\widetilde{u}_{n-2}\cdots\widetilde{u}_0$
est un homomorphisme \emph{injectif} $\mathcal{O}_X(n)\to\mathcal{O}_X$~;
nous désignerons par $\mathcal{I}_{n,X}$ son image, qui est donc un Idéal
quasi-cohérent de $\mathcal{O}_X$, \emph{isomorphe} à $\mathcal{O}_X(n)$. En
outre, le diagramme
\[
\xymatrix{
  \mathcal{O}_X(m)\otimes_{\mathcal{O}_X}\mathcal{O}_X(n)
    \ar[r]^{\lambda}\ar[d]_{\widetilde{v}_m\otimes\widetilde{v}_n}
  & \mathcal{O}_X(m+n)\ar[d]^{\widetilde{v}_{m+n}}\\
  \mathcal{O}_X\ar[r]_{\mathrm{id}}
  & \mathcal{O}_X
}
\]
est commutatif pour $m\geq0$, $n\geq0$. On en conclut les inclusions suivantes
\begin{equation}
\label{II.8.1.6.1-fr}
  \mathcal{I}_{0,X}=\mathcal{O}_X,\qquad
  \mathcal{I}_{n,X}\subset\mathcal{I}_{m,X}
  \quad\text{pour }0\leq m\leq n
\tag{8.1.6.1}
\end{equation}
\begin{equation}
\label{II.8.1.6.2-fr}
  \mathcal{I}_{m,X}\mathcal{I}_{n,X}\subset\mathcal{I}_{m+n,X}
  \quad\text{pour }m\geq0,\ n\geq0.
\tag{8.1.6.2}
\end{equation}
\end{env}

\oldpage[II]{156}

\begin{proposition}[8.1.7]
\label{II.8.1.7-fr}
Soient $Y$ un préschéma, $\mathcal{I}$ un Idéal quasi-cohérent de
$\mathcal{O}_Y$, $X=\operatorname{Proj}(\bigoplus_{n\geq0}\mathcal{I}^n)$ le
$Y$-schéma obtenu en faisant éclater $\mathcal{I}$. On a alors, pour tout
$n>0$, un isomorphisme canonique
\begin{equation}
\label{II.8.1.7.1-fr}
  \mathcal{O}_X(n)\xrightarrow{\sim}\mathcal{I}^n\mathcal{O}_X
  =\mathcal{I}_{n,X}
\tag{8.1.7.1}
\end{equation}
(cf. (\hyperref[0.4.3.5-fr]{0, 4.3.5})), et par suite
$\mathcal{I}^n\mathcal{O}_X$ est un $\mathcal{O}_X$-Module inversible très
ample si $n>0$.
\end{proposition}

La dernière assertion est immédiate, puisque $\mathcal{O}_X(1)$ est inversible
(\hyperref[II.3.2.5-fr]{3.2.5}) et très ample pour $Y$ par définition
(\hyperref[II.4.4.3-fr]{4.4.3} et \hyperref[II.4.4.9-fr]{4.4.9}). D'autre
part, par définition, l'image de $v_n$ n'est autre que
$\mathcal{I}^n\mathcal{S}$, et (\hyperref[II.8.1.7.1-fr]{8.1.7.1}) résulte donc
de l'exactitude du foncteur $\widetilde{\mathcal{M}}$
(\hyperref[II.3.2.4-fr]{3.2.4}) et de la formule
(\hyperref[II.3.2.4.1-fr]{3.2.4.1}).

\begin{corollary}[8.1.8]
\label{II.8.1.8-fr}
Sous les hypothèses de (\hyperref[II.8.1.7-fr]{8.1.7}), si $f:X\to Y$ est le
morphisme structural et $Y'$ le sous-préschéma fermé de $Y$ défini par
$\mathcal{I}$, le sous-préschéma fermé $X'=f^{-1}(Y')$ de $X$ est défini par
$\mathcal{I}\mathcal{O}_X$ (isomorphe canoniquement à $\mathcal{O}_X(1)$),
d'où une suite exacte canonique
\begin{equation}
\label{II.8.1.8.1-fr}
  0\to\mathcal{O}_X(1)\to\mathcal{O}_X\to\mathcal{O}_{X'}\to0.
\tag{8.1.8.1}
\end{equation}
\end{corollary}

Cela résulte de (\hyperref[II.8.1.7.1-fr]{8.1.7.1}) et de
(\hyperref[I.4.4.5-fr]{I, 4.4.5}).

\begin{env}[8.1.9]
\label{II.8.1.9-fr}
Sous les hypothèses de (\hyperref[II.8.1.7-fr]{8.1.7}), on peut préciser la
structure des $\mathcal{I}_{n,X}$. Remarquons en effet que l'homomorphisme
\[
  \widetilde{u}_{-1}:\mathcal{O}_X\to\mathcal{O}_X(-1)
\]
correspond canoniquement à une section $s$ de $\mathcal{O}_X(-1)$ au-dessus de
$X$, que nous appellerons la \emph{section canonique} (relative à
$\mathcal{I}$) (\hyperref[0.5.1.1-fr]{0, 5.1.1}). D'autre part, dans le
diagramme (\hyperref[II.8.1.5.5-fr]{8.1.5.5}), les flèches horizontales sont
des isomorphismes (\hyperref[II.3.2.7-fr]{3.2.7})~; en remplaçant dans ce
diagramme $h$ par $k$ et $k$ par $-1$, on obtient
$\widetilde{u}_k=1_k\otimes\widetilde{u}_{-1}$ ($1_h$ désignant l'identité de
$\mathcal{O}_X(h)$), autrement dit l'homomorphisme $\widetilde{u}_k$ n'est
autre que la \emph{multiplication tensorielle par la section canonique $s$}
(pour tout $k\in\mathbf{Z}$). L'homomorphisme $\widetilde{u}$
(\hyperref[II.8.1.5.6-fr]{8.1.5.6}) s'interprète donc aussi de la même manière.

Par suite, pour tout $n\geq0$, l'homomorphisme
$\widetilde{v}_n:\mathcal{O}_X(n)\to\mathcal{O}_X$ n'est autre que la
multiplication tensorielle par $s^{\otimes n}$~; on en déduit~:
\end{env}

\begin{corollary}[8.1.10]
\label{II.8.1.10-fr}
Avec les notations de (\hyperref[II.8.1.8-fr]{8.1.8}), l'espace sous-jacent à
$X'$ est l'ensemble des $x\in X$ tels que $s(x)=0$, $s$ désignant la section
canonique de $\mathcal{O}_X(-1)$.
\end{corollary}

En effet, si $c_x$ est un générateur de la fibre $(\mathcal{O}_X(1))_x$ en un
point $x$, $s_x\otimes c_x$ s'identifie canoniquement à un générateur de la
fibre de $\mathcal{I}_{1,X}$ au point $x$, et est donc inversible si et
seulement si l'on a
$s_x\notin\mathfrak{m}_x(\mathcal{O}_X(-1))_x$, autrement dit si
$s(x)\neq0$.

\begin{proposition}[8.1.11]
\label{II.8.1.11-fr}
Soient $Y$ un préschéma intègre, $\mathcal{I}$ un Idéal fractionnaire
quasi-cohérent de $\mathcal{R}(Y)$, $X$ le $Y$-schéma obtenu en faisant éclater
$\mathcal{I}$. Alors $\mathcal{I}\mathcal{O}_X$ est un
$\mathcal{O}_X$-Module inversible très ample pour $Y$.
\end{proposition}

La question étant locale sur $Y$ (\hyperref[II.4.4.5-fr]{4.4.5}), on peut se
borner au cas où $Y=\operatorname{Spec}(A)$, où $A$ est un anneau intègre de
corps des fractions $K$ et $\mathcal{I}=\widetilde{\mathfrak{I}}$,
$\mathfrak{I}$ étant un idéal fractionnaire de $K$~; il existe par suite un
élément $a\neq0$ de $A$ tel que $a\mathfrak{I}\subset A$. Posons
$S=\bigoplus_{n\geq0}\mathfrak{I}^n$~; l'application $x\mapsto ax$ est un
$A$-isomorphisme de $\mathfrak{I}^{n+1}=(S(1))_n$ sur
$a\mathfrak{I}^{n+1}=a\mathfrak{I}S_n\subset\mathfrak{I}^n=S_n$,

\oldpage[II]{157}
donc définit un (TN)-isomorphisme de degré $0$ de $S$-modules gradués
$S_+(1)\to a\mathfrak{I}S$. Par ailleurs $x\mapsto a^{-1}x$ est un
isomorphisme de degré $0$ de $S$-modules gradués
$a\mathfrak{I}S\xrightarrow{\sim}\mathfrak{I}S$. On obtient donc ainsi par
composition (\hyperref[II.3.2.4-fr]{3.2.4}) un isomorphisme de
$\mathcal{O}_X$-Modules
$\mathcal{O}_X(1)\xrightarrow{\sim}\mathcal{I}\mathcal{O}_X$, et comme $S$
est engendrée par $S_1=\mathfrak{I}$, $\mathcal{O}_X(1)$ est inversible
(\hyperref[II.3.2.5-fr]{3.2.5}) et très ample
(\hyperref[II.4.4.3-fr]{4.4.3} et \hyperref[II.4.4.9-fr]{4.4.9}), d'où notre
assertion.

\subsection{Résultats préliminaires sur la localisation dans les anneaux gradués}
\label{subsection:II.8.2-fr}

\begin{env}[8.2.1]
\label{II.8.2.1-fr}
Soit $S$ un anneau gradué, où nous ne supposons pas pour le moment les degrés
positifs. On posera
\begin{equation}
\label{II.8.2.1.1-fr}
  S^{\geq}=\bigoplus_{n\geq0}S_n,
  \qquad
  S^{\leq}=\bigoplus_{n\leq0}S_n,
\tag{8.2.1.1}
\end{equation}
qui sont des sous-anneaux gradués de $S$, à degrés respectivement tous
positifs et tous négatifs. Si $f$ est un élément homogène de degré $d$
(positif ou négatif) de $S$, l'anneau des fractions $S_f=S'$ est encore muni
d'une structure d'anneau gradué, en prenant pour $S'_n$ ($n\in\mathbf{Z}$)
l'ensemble des $x/f^k$, où $x\in S_{n+kd}$ ($k\geq0$)~; on posera encore
$S_{(f)}=S'_0$, et on écrira $S_f^{\geq}$ et $S_f^{\leq}$ pour $S'^{\geq}$ et
$S'^{\leq}$ respectivement. Si $d>0$, on a
\begin{equation}
\label{II.8.2.1.2-fr}
  (S^{\geq})_f=S_f
\tag{8.2.1.2}
\end{equation}
puisque, si $x\in S_{n+kd}$, avec $n+kd<0$, on peut écrire
$x/f^k=xf^h/f^{h+k}$ et que $n+(h+k)d>0$ pour $h$ assez grand et $>0$. On en
conclut par définition que
\begin{equation}
\label{II.8.2.1.3-fr}
  (S^{\geq})_{(f)}=(S_f^{\geq})_0=S_{(f)}.
\tag{8.2.1.3}
\end{equation}

Si $M$ est un $S$-module gradué, on posera de même
\begin{equation}
\label{II.8.2.1.4-fr}
  M^{\geq}=\bigoplus_{n\geq0}M_n,
  \qquad
  M^{\leq}=\bigoplus_{n\leq0}M_n
\tag{8.2.1.4}
\end{equation}
qui sont respectivement un $S^{\geq}$-module gradué et un $S^{\leq}$-module
gradué et ont pour intersection le $S_0$-module $M_0$. Si $f\in S_d$, on
définit encore $M_f$ comme $S_f$-module gradué en prenant comme éléments de
degré $n$ les $z/f^k$, où $z\in M_{n+kd}$ ($k\geq0$)~; on désignera par
$M_{(f)}$ l'ensemble des éléments de degré $0$ de $M_f$, qui est un
$S_{(f)}$-module, et on écrira $M_f^{\geq}$ et $M_f^{\leq}$ au lieu de
$(M_f)^{\geq}$ et $(M_f)^{\leq}$ respectivement. Si $d>0$, on voit comme
ci-dessus que
\begin{equation}
\label{II.8.2.1.5-fr}
  (M^{\geq})_f=M_f
\tag{8.2.1.5}
\end{equation}
et
\begin{equation}
\label{II.8.2.1.6-fr}
  (M^{\geq})_{(f)}=(M_f^{\geq})_0=M_{(f)}.
\tag{8.2.1.6}
\end{equation}
\end{env}

\begin{env}[8.2.2]
\label{II.8.2.2-fr}
Soit $z$ une indéterminée, que nous nommerons \emph{variable
d'homogénéisation}. Si $S$ est un anneau gradué (à degrés positifs ou
négatifs), l'algèbre de polynômes\footnote{Il ne saurait ici y avoir de
confusion avec l'usage de la notation $\widehat{S}$ pour désigner le séparé
complété d'un anneau.}
\begin{equation}
\label{II.8.2.2.1-fr}
  \widehat{S}=S[z]
\tag{8.2.2.1}
\end{equation}

\oldpage[II]{158}
est une $S$-algèbre graduée, quand on prend pour degré de $fz^n$ ($n\geq0$) où
$f$ est homogène,
\begin{equation}
\label{II.8.2.2.2-fr}
  \deg(fz^n)=n+\deg f.
\tag{8.2.2.2}
\end{equation}
\end{env}

\begin{lemma}[8.2.3]
\label{II.8.2.3-fr}
\begin{enumerate}
\item[\rm(i)] On a des isomorphismes canoniques d'anneaux (non gradués)
\begin{equation}
\label{II.8.2.3.1-fr}
  \widehat{S}_{(z)}\xrightarrow{\sim}
  \widehat{S}/(z-1)\widehat{S}\xrightarrow{\sim}S.
\tag{8.2.3.1}
\end{equation}
\item[\rm(ii)] On a un isomorphisme canonique d'anneaux (non gradués)
\begin{equation}
\label{II.8.2.3.2-fr}
  \widehat{S}_{(f)}\xrightarrow{\sim}S_f^{\leq}
\tag{8.2.3.2}
\end{equation}
pour tout $f\in S_d$, avec $d>0$.
\end{enumerate}
\end{lemma}

Le premier des isomorphismes de (\hyperref[II.8.2.3.1-fr]{8.2.3.1}) a été
défini dans (\hyperref[II.2.2.5-fr]{2.2.5}) et le second est trivial~;
l'isomorphisme $\widehat{S}_{(z)}\xrightarrow{\sim}S$ ainsi défini fait donc
correspondre à $xz^n/z^{n+k}$, où $\deg(x)=k$ ($k\geq-n$) l'élément $x$.
L'homomorphisme (\hyperref[II.8.2.3.2-fr]{8.2.3.2}) fait correspondre à
$xz^n/f^k$, où $\deg(x)=kd-n$, l'élément $x/f^k$, de degré $-n$ dans
$S_f^{\leq}$, et il est encore évident qu'on a bien ici un isomorphisme.

\begin{env}[8.2.4]
\label{II.8.2.4-fr}
Soit $M$ un $S$-module gradué. Il est clair que le $S$-module
\begin{equation}
\label{II.8.2.4.1-fr}
  \widehat{M}=M\otimes_S\widehat{S}=M\otimes_S S[z]
\tag{8.2.4.1}
\end{equation}
est somme directe des $S$-modules $M\otimes_S Sz^n$, donc des groupes abéliens
$M_k\otimes_S Sz^n$ ($k\in\mathbf{Z}$, $n\geq0$)~; on définit sur
$\widehat{M}$ une structure de $\widehat{S}$-module gradué en prenant
\begin{equation}
\label{II.8.2.4.2-fr}
  \deg(x\otimes z^n)=n+\deg x
\tag{8.2.4.2}
\end{equation}
pour tout $x$ homogène dans $M$. Nous laissons au lecteur le soin de démontrer
l'analogue de (\hyperref[II.8.2.3-fr]{8.2.3})~:
\end{env}

\begin{lemma}[8.2.5]
\label{II.8.2.5-fr}
\begin{enumerate}
\item[\rm(i)] On a un di-isomorphisme canonique de modules (non gradués)
\begin{equation}
\label{II.8.2.5.1-fr}
  \widehat{M}_{(z)}\xrightarrow{\sim}M.
\tag{8.2.5.1}
\end{equation}
\item[\rm(ii)] Pour tout $f\in S_d$ ($d>0$), on a un di-isomorphisme de
modules (non gradués)
\begin{equation}
\label{II.8.2.5.2-fr}
  \widehat{M}_{(f)}\xrightarrow{\sim}M_f^{\leq}.
\tag{8.2.5.2}
\end{equation}
\end{enumerate}
\end{lemma}

\begin{env}[8.2.6]
\label{II.8.2.6-fr}
Soit $S$ un anneau gradué à degrés \emph{positifs}, et considérons dans $S$
la suite décroissante d'idéaux gradués
\begin{equation}
\label{II.8.2.6.1-fr}
  S_{[n]}=\bigoplus_{m\geq n}S_m\qquad(n\geq0)
\tag{8.2.6.1}
\end{equation}
(en particulier on a $S_{[0]}=S$, $S_{[1]}=S_+$). Comme il est clair que
$S_{[m]}S_{[n]}\subset S_{[m+n]}$, on peut définir un \emph{anneau gradué}
$S^\natural$ en prenant
\begin{equation}
\label{II.8.2.6.2-fr}
  S^\natural=\bigoplus_{n\geq0}S_n^\natural
  \quad\text{avec}\quad S_n^\natural=S_{[n]}.
\tag{8.2.6.2}
\end{equation}
$S_0^\natural$ est donc l'anneau $S$ considéré comme anneau \emph{non gradué},
et $S^\natural$ est par suite une $S_0^\natural$-algèbre. Pour tout élément
homogène $f\in S_d$ ($d>0$), nous désignerons par $f^\natural$ l'élément
$f$ considéré comme appartenant à $S_{[d]}=S_d^\natural$. Avec ces
notations~:
\end{env}

\oldpage[II]{159}

\begin{lemma}[8.2.7]
\label{II.8.2.7-fr}
Soient $S$ un anneau gradué à degrés positifs, $f$ un élément homogène de
$S_d$ ($d>0$). On a des isomorphismes canoniques d'anneaux
\begin{equation}
\label{II.8.2.7.1-fr}
  S_f\xrightarrow{\sim}\bigoplus_{n\in\mathbf{Z}}S(n)_{(f)}
\tag{8.2.7.1}
\end{equation}
\begin{equation}
\label{II.8.2.7.2-fr}
  (S_f^{\geq})_{f/1}\xrightarrow{\sim}S_f
\tag{8.2.7.2}
\end{equation}
\begin{equation}
\label{II.8.2.7.3-fr}
  S^\natural_{(f^\natural)}\xrightarrow{\sim}S_f^{\geq}
\tag{8.2.7.3}
\end{equation}
dont les deux premiers sont des isomorphismes d'anneaux gradués.
\end{lemma}

Il est immédiat par définition que l'on a
$(S_f)_n=(S(n)_f)_0$, d'où l'isomorphisme
(\hyperref[II.8.2.7.1-fr]{8.2.7.1}), qui n'est autre que l'identité. D'autre
part, comme $f/1$ est inversible dans $S_f$, il y a un isomorphisme canonique
$S_f\xrightarrow{\sim}(S_f^{\geq})_{f/1}=(S_f)_{f/1}$ en vertu de
(\hyperref[II.8.2.1.2-fr]{8.2.1.2}) appliqué à $S_f$~; l'isomorphisme
réciproque est par définition l'isomorphisme
(\hyperref[II.8.2.7.2-fr]{8.2.7.2}). Enfin, si
$x=\sum_{m\geq n}y_m$ est un élément de $S_{[n]}$, avec $n=kd$, on fait
correspondre à l'élément $x/(f^\natural)^k$ l'élément
$\sum_m y_m/f^k$ de $S_f^{\geq}$, et on vérifie aussitôt que cela définit un
isomorphisme (\hyperref[II.8.2.7.3-fr]{8.2.7.3}).

\begin{env}[8.2.8]
\label{II.8.2.8-fr}
Si $M$ est un $S$-module gradué, on pose de même pour tout $n\in\mathbf{Z}$
\begin{equation}
\label{II.8.2.8.1-fr}
  M_{[n]}=\bigoplus_{m\geq n}M_m
\tag{8.2.8.1}
\end{equation}
et comme $S_{[m]}M_{[n]}\subset M_{[m+n]}$ ($m\geq0$), on peut définir un
$S^\natural$-module gradué $M^\natural$ en prenant
\begin{equation}
\label{II.8.2.8.2-fr}
  M^\natural=\bigoplus_{n\in\mathbf{Z}}M_n^\natural,
  \quad\text{avec}\quad M_n^\natural=M_{[n]}.
\tag{8.2.8.2}
\end{equation}
\end{env}

Nous laissons encore au lecteur la démonstration du~:

\begin{lemma}[8.2.9]
\label{II.8.2.9-fr}
Avec les notations de (\hyperref[II.8.2.7-fr]{8.2.7}) et
(\hyperref[II.8.2.8-fr]{8.2.8}), on a les di-isomorphismes canoniques de
modules
\begin{equation}
\label{II.8.2.9.1-fr}
  M_f\xrightarrow{\sim}\bigoplus_{n\in\mathbf{Z}}M(n)_{(f)}
\tag{8.2.9.1}
\end{equation}
\begin{equation}
\label{II.8.2.9.2-fr}
  (M_f^{\geq})_{f/1}\xrightarrow{\sim}M_f
\tag{8.2.9.2}
\end{equation}
\begin{equation}
\label{II.8.2.9.3-fr}
  M^\natural_{(f^\natural)}\xrightarrow{\sim}M_f^{\geq}
\tag{8.2.9.3}
\end{equation}
dont les deux premiers sont des di-isomorphismes de modules gradués.
\end{lemma}

\begin{lemma}[8.2.10]
\label{II.8.2.10-fr}
Soit $S$ un anneau gradué à degrés positifs.
\begin{enumerate}
\item[\rm(i)] Pour que $S^\natural$ soit une $S_0^\natural$-algèbre de type
fini (resp. noethérienne) il faut et il suffit que $S$ soit une $S_0$-algèbre
de type fini (resp. noethérienne).
\item[\rm(ii)] Pour que $S_{n+1}^\natural=S_1^\natural S_n^\natural$ pour
$n\geq n_0$, il faut et il suffit que $S_{n+1}=S_1S_n$ pour $n\geq n_0$.
\item[\rm(iii)] Pour que $S_n^\natural=S_1^{\natural n}$ pour $n\geq n_0$, il
faut et il suffit que $S_n=S_1^n$ pour $n\geq n_0$.
\item[\rm(iv)] Si $(f_\alpha)$ est un ensemble d'éléments homogènes de $S_+$
telle que $S_+$ soit la racine dans $S_+$ de l'idéal de $S_+$ engendré par
les $f_\alpha$, alors $S_+^\natural$ est la racine dans $S_+^\natural$ de
l'idéal de $S_+^\natural$ engendré par les $f_\alpha^\natural$.
\end{enumerate}
\end{lemma}

\begin{enumerate}
\item[\rm(i)] Si $S^\natural$ est une $S_0^\natural$-algèbre de type fini,
$S_+=S_1^\natural$ est un module de type fini sur $S=S_0^\natural$ par
(\hyperref[II.2.1.6-fr]{2.1.6}, (i)), donc $S$ est une $S_0$-algèbre de type
fini (\hyperref[II.2.1.4-fr]{2.1.4})~; si $S^\natural$ est un anneau
noethérien, il en est de même de $S_0^\natural=S$
(\hyperref[II.2.1.5-fr]{2.1.5}). Inversement, si $S$ est une $S_0$-algèbre

\oldpage[II]{160}
de type fini, on sait (\hyperref[II.2.1.6-fr]{2.1.6}, (ii)) qu'il existe
$h>0$ et $m_0>0$ tels que $S_{n+h}=S_hS_n$ pour $n\geq m_0$~; on peut
évidemment supposer $m_0\geq h$. En outre, les $S_m$ sont des $S_0$-modules
de type fini (\hyperref[II.2.1.6-fr]{2.1.6}, (i)). Cela étant, si
$n\geq m_0+h$, on a $S_n^\natural=S_hS_{n-h}^\natural
=S_h^\natural S_{n-h}^\natural$~; et si $m<m_0+h$, on a, en posant
$E=S_{m_0}+\cdots+S_{m_0+h-1}$,
\[
  S_m^\natural=S_m+\cdots+S_{m_0+h-1}+S_hE+S_h^2E+\cdots.
\]
Pour $1\leq m\leq m_0$, soit $G_m$ la réunion de systèmes finis de
générateurs des $S_0$-modules $S_i$, pour $m\leq i\leq m_0+h-1$,
considérée comme partie de $S_{[m]}$. Pour $m_0+1\leq m\leq m_0+h-1$, soit
de même $G_m$ la réunion de systèmes finis de générateurs des $S_0$-modules
$S_i$, pour $m\leq i\leq m_0+h-1$, et de $S_hE$, considérée comme partie de
$S_{[m]}$. Il est clair que l'on a $S_m^\natural=S_0^\natural G_m$ pour
$1\leq m\leq m_0+h-1$, et par suite la réunion $G$ des $G_m$ pour
$1\leq m\leq m_0+h-1$ est un système de générateurs de la
$S_0^\natural$-algèbre $S^\natural$. On en conclut que si $S=S_0^\natural$
est un anneau noethérien, il en est de même de $S^\natural$.

\item[\rm(ii)] Il est clair que si $S_{n+1}=S_1S_n$ pour $n\geq n_0$, on a
$S_{n+1}^\natural=S_1S_n^\natural$, et \emph{a fortiori}
$S_{n+1}^\natural=S_1^\natural S_n^\natural$ pour $n\geq n_0$.
Inversement, cette dernière relation s'écrit
\[
  S_{n+1}+S_{n+2}+\cdots=(S_1+S_2+\cdots)(S_n+S_{n+1}+\cdots)
\]
et la comparaison des termes de degré $n+1$ (dans $S$) aux deux membres donne
$S_{n+1}=S_1S_n$.

\item[\rm(iii)] Si $S_n=S_1^n$ pour $n\geq n_0$, on a
$S_n^\natural=S_1^n+S_1^{n+1}+\cdots$~; comme $S_1^\natural$ contient
$S_1+S_1^2+\cdots$, on a $S_n^\natural\subset S_1^{\natural n}$, et par suite
$S_n^\natural=S_1^{\natural n}$ pour $n\geq n_0$. Inversement, les seuls
termes de $S_1^{\natural n}=(S_1+S_2+\cdots)^n$ qui soient de degré $n$ dans
$S$ sont ceux de $S_1^n$~; la relation $S_n^\natural=S_1^{\natural n}$
implique donc $S_n=S_1^n$.

\item[\rm(iv)] Il suffit de prouver que si un élément $g\in S_{k+h}$ est
considéré comme un élément de $S_k^\natural$ ($k>0$, $h\geq0$), alors il
existe un entier $n>0$ tel que dans $S_{kn}^\natural$, $g^n$ soit combinaison
linéaire des $f_\alpha^\natural$ à coefficients dans $S^\natural$. Par
hypothèse, il y a un entier $m_0$ tel que pour $m\geq m_0$, on ait, \emph{dans
$S$}, $g^m=\sum_\alpha c_{\alpha m}f_\alpha$, les indices $\alpha$ figurant
dans cette formule étant \emph{indépendants de $m$}~; en outre, on peut
évidemment supposer les $c_{\alpha m}$ homogènes, avec
\[
  \deg(c_{\alpha m})=m(k+h)-\deg f_\alpha
\]
dans $S$. Prenons alors $m_0$ assez grand pour que l'on ait
$km_0>\deg f_\alpha$ pour tous les $f_\alpha$ qui figurent dans $g^{m_0}$~;
pour tout $\alpha$, soit $c'_{\alpha m}$ l'élément $c_{\alpha m}$ considéré
comme de degré $km-\deg(f_\alpha)$ dans $S^\natural$~; on a alors, dans
$S^\natural$, $g^m=\sum_\alpha c'_{\alpha m}f_\alpha^\natural$, ce qui
termine la démonstration.
\end{enumerate}

\begin{env}[8.2.11]
\label{II.8.2.11-fr}
Considérons la $S_0$-algèbre graduée
\begin{equation}
\label{II.8.2.11.1-fr}
  S^\natural\otimes_S S_0=S^\natural/S_+S^\natural
  =\bigoplus_{n\geq0}S_{[n]}/S_+S_{[n]}.
\tag{8.2.11.1}
\end{equation}
Comme $S_n$ est un $S_0$-module quotient de $S_{[n]}/S_+S_{[n]}$, on a un
homomorphisme canonique de $S_0$-algèbres graduées
\begin{equation}
\label{II.8.2.11.2-fr}
  S^\natural\otimes_S S_0\to S
\tag{8.2.11.2}
\end{equation}
qui est évidemment \emph{surjectif}, et correspond par suite
(\hyperref[II.2.9.2-fr]{2.9.2}) à une \emph{immersion fermée} canonique
\begin{equation}
\label{II.8.2.11.3-fr}
  \operatorname{Proj}(S)\to
  \operatorname{Proj}(S^\natural\otimes_S S_0).
\tag{8.2.11.3}
\end{equation}
\end{env}

\oldpage[II]{161}

\begin{proposition}[8.2.12]
\label{II.8.2.12-fr}
Le morphisme canonique (\hyperref[II.8.2.11.3-fr]{8.2.11.3}) est bijectif.
Pour que l'homomorphisme (\hyperref[II.8.2.11.2-fr]{8.2.11.2}) soit
(TN)-bijectif, il faut et il suffit qu'il existe $n_0$ tel que
$S_{n+1}=S_1S_n$ pour $n\geq n_0$. Si cette dernière condition est vérifiée,
(\hyperref[II.8.2.11.3-fr]{8.2.11.3}) est un isomorphisme~; la réciproque est
vraie lorsque $S$ est noethérien.
\end{proposition}

Pour démontrer la première assertion, il suffit
(\hyperref[II.2.8.3-fr]{2.8.3}) de prouver que le noyau $\mathfrak{I}$ de
l'homomorphisme (\hyperref[II.8.2.11.2-fr]{8.2.11.2}) est formé d'éléments
\emph{nilpotents}. Or si $f\in S_{[n]}$ est un élément dont la classe mod.
$S_+S_{[n]}$ appartient à ce noyau, cela signifie que $f\in S_{[n+1]}$~;
l'élément $f^{n+1}$, considéré comme appartenant à $S_{[n(n+1)]}$, appartient
alors à $S_+S_{[n(n+1)]}$, puisqu'il s'écrit $f\cdot f^n$~; donc la classe de
$f^{n+1}$ mod. $S_+S_{[n(n+1)]}$ est nulle, ce qui prouve notre assertion.
Comme l'hypothèse $S_{n+1}=S_1S_n$ pour $n\geq n_0$ équivaut à
$S_{n+1}^{\natural}=S_1^{\natural}S_n^{\natural}$ pour $n\geq n_0$
(\hyperref[II.8.2.10-fr]{8.2.10}, (ii)), cette hypothèse équivaut par définition
au fait que (\hyperref[II.8.2.11.2-fr]{8.2.11.2}) est (TN)-injectif, donc
(TN)-bijectif, et alors (\hyperref[II.8.2.11.3-fr]{8.2.11.3}) est un
isomorphisme en vertu de (\hyperref[II.2.9.1-fr]{2.9.1}). Inversement, si
(\hyperref[II.8.2.11.3-fr]{8.2.11.3}) est un isomorphisme, le faisceau
$\widetilde{\mathfrak{I}}$ sur
$\operatorname{Proj}(S^{\natural}\otimes_S S_0)$ est nul
(\hyperref[II.2.9.2-fr]{2.9.2}, (i))~; comme
$S^{\natural}\otimes_S S_0$ est noethérien comme quotient de $S^{\natural}$
(\hyperref[II.8.2.10-fr]{8.2.10}, (i)), on conclut de
(\hyperref[II.2.7.3-fr]{2.7.3}) que $\mathfrak{I}$ vérifie la condition (TN),
donc $S_{n+1}^{\natural}=S_1^{\natural}S_n^{\natural}$ pour $n\geq n_0$, et
cela achève la démonstration en vertu de
(\hyperref[II.8.2.10-fr]{8.2.10}, (ii)).

\begin{env}[8.2.13]
\label{II.8.2.13-fr}
Considérons maintenant les injections canoniques $(S_+)^n\to S_{[n]}$, qui
définissent un homomorphisme injectif de degré $0$ d'anneaux gradués
\begin{equation}
\label{II.8.2.13.1-fr}
  \bigoplus_{n\geq0}(S_+)^n\to S^{\natural}.
\tag{8.2.13.1}
\end{equation}
\end{env}

\begin{proposition}[8.2.14]
\label{II.8.2.14-fr}
Pour que l'homomorphisme (\hyperref[II.8.2.13.1-fr]{8.2.13.1}) soit un
(TN)-isomorphisme, il faut et il suffit qu'il existe $n_0$ tel que
$S_n=S_1^n$ pour tout $n\geq n_0$. Lorsqu'il en est ainsi, le morphisme
correspondant à (\hyperref[II.8.2.13.1-fr]{8.2.13.1}) est partout défini et
est un isomorphisme
\[
  \operatorname{Proj}(S^{\natural})\xrightarrow{\sim}
  \operatorname{Proj}\left(\bigoplus_{n\geq0}(S_+)^n\right)~;
\]
la réciproque est vraie lorsque $S$ est noethérien.
\end{proposition}

Les deux premières assertions sont évidentes, vu
(\hyperref[II.8.2.10-fr]{8.2.10}, (iii)) et
(\hyperref[II.2.9.1-fr]{2.9.1}). La troisième résultera de
(\hyperref[II.8.2.10-fr]{8.2.10}, (i) et (iii)) et du lemme suivant~:

\begin{lemma}[8.2.14.1]
\label{II.8.2.14.1-fr}
Soit $T$ un anneau gradué en degrés positifs qui est une $T_0$-algèbre de type
fini. Si le morphisme correspondant à l'homomorphisme injectif
$\bigoplus_{n\geq0}T_1^n\to T$ est partout défini et est un isomorphisme
\[
  \operatorname{Proj}(T)\to
  \operatorname{Proj}\left(\bigoplus_{n\geq0}T_1^n\right),
\]
il existe $n_0$ tel que $T_n=T_1^n$ pour $n\geq n_0$.
\end{lemma}

En effet, soient $g_i$ ($1\leq i\leq r$) des générateurs du $T_0$-module
$T_1$. L'hypothèse entraîne d'abord que les $D_+(g_i)$ recouvrent
$\operatorname{Proj}(T)$ (\hyperref[II.2.8.1-fr]{2.8.1}). Soit
$(h_j)_{1\leq j\leq s}$ un système d'éléments homogènes de $T_+$, avec
$\deg(h_j)=n_j$, qui forment avec les $g_i$ un système de générateurs de
l'idéal $T_+$, ou (ce qui revient au même
(\hyperref[II.2.1.3-fr]{2.1.3})) un système de générateurs de $T$ en tant que
$T_0$-algèbre~; si on pose $T'=\bigoplus_{n\geq0}T_1^n$, l'élément
$h_j/g_i^{n_j}$ de l'anneau $T_{(g_i)}$ doit par hypothèse appartenir au
sous-anneau $T'_{(g_i)}$, donc il existe un entier $k$ tel que
$T_1^k h_j\subset T_1^{k+n_j}$ pour tout $j$. On en conclut par récurrence
sur $r$ que $T_1^k h_j^r\subset T'$ pour tout $r\geq1$, et par définition des
$h_j$, on a donc $T_1^kT\subset T'$. Il y a d'autre part pour tout $j$ un
entier $m_j$ tel que $h_j^{m_j}$ appartienne à l'idéal de $T$ engendré par
les $g_i$ (\hyperref[II.2.3.14-fr]{2.3.14}), donc $h_j^{m_j}\in T_1T$, et

\oldpage[II]{162}
$h_j^{m_jk}\in T_1^kT\subset T'$. Il y a par suite un entier $m_0\geq k$ tel
que $h_j^m\in T_1^{mn_j}$ pour $m\geq m_0$. Cela étant, si $q$ est le plus
grand des entiers $n_j$, le nombre $n_0=qsm_0+k$ répond à la question. En
effet, un élément de $T_n$, pour $n\geq n_0$, est somme de monômes appartenant
à $T_1^\alpha u$, où $u$ est un produit de puissances des $h_j$~; si
$\alpha\geq k$, il résulte de ce qui précède que $T_1^\alpha u\subset T_1^n$~;
dans le cas contraire, un des exposants des $h_j$ est $\geq m_0$, donc
$u\in T_1^\beta v$, où $\beta\geq k$ et $v$ est encore un produit de
puissances des $h_j$~; on est alors ramené au cas précédent, et on voit
finalement que $T_1^\alpha u\subset T_1^n$ dans tous les cas.

\begin{remark}[8.2.15]
\label{II.8.2.15-fr}
La condition $S_n=S_1^n$ pour $n\geq n_0$ entraîne évidemment
$S_{n+1}=S_1S_n$ pour $n\geq n_0$, mais la réciproque est inexacte, même
quand on suppose $S$ noethérien. Par exemple, soient $K$ un corps,
$A=K[x]$, $B=K[y]/y^2K[y]$, où $x$ et $y$ sont deux indéterminées, $x$ étant
prise de degré $1$ et $y$ de degré $2$, et soit $S=A\otimes_KB$, de sorte que
$S$ est une algèbre graduée sur $K$ ayant une base formée des éléments $1$,
$x^n$ ($n\geq1$) et $x^ny$ ($n\geq0$). Il est immédiat que
$S_{n+1}=S_1S_n$ pour $n\geq2$, mais $S_1^n=Kx^n$ tandis que
$S_n=Kx^n+Kx^{n-2}y$ pour $n\geq2$.
\end{remark}

\subsection{Cônes projetants}
\label{subsection:II.8.3-fr}

\begin{env}[8.3.1]
\label{II.8.3.1-fr}
Soit $Y$ un préschéma~; dans tout ce numéro, il ne sera question que de
\emph{$Y$-préschémas} et de \emph{$Y$-morphismes}. Soit $\mathcal{S}$ une
$\mathcal{O}_Y$-Algèbre quasi-cohérente graduée à \emph{degrés positifs}~;
\emph{nous supposerons en outre que l'on a $\mathcal{S}_0=\mathcal{O}_Y$}.
Conformément aux notations introduites en
(\hyperref[II.8.2.2-fr]{8.2.2}), nous poserons
\begin{equation}
\label{II.8.3.1.1-fr}
  \widehat{\mathcal{S}}=\mathcal{S}[z]
  =\mathcal{S}\otimes_{\mathcal{O}_Y}\mathcal{O}_Y[z]
\tag{8.3.1.1}
\end{equation}
que l'on considère comme $\mathcal{O}_Y$-Algèbre graduée à degrés positifs en
définissant les degrés par la formule
(\hyperref[II.8.2.2.2-fr]{8.2.2.2}), de sorte que pour tout ouvert affine $U$
de $Y$, on a
\[
  \Gamma(U,\widehat{\mathcal{S}})=(\Gamma(U,\mathcal{S}))[z].
\]
Dans ce qui suit, nous poserons
\begin{equation}
\label{II.8.3.1.2-fr}
  X=\operatorname{Proj}(\mathcal{S}),\qquad
  C=\operatorname{Spec}(\mathcal{S}),\qquad
  \widehat{C}=\operatorname{Proj}(\widehat{\mathcal{S}})
\tag{8.3.1.2}
\end{equation}
(où dans la définition de $C$, $\mathcal{S}$ est considérée comme
$\mathcal{O}_Y$-Algèbre non graduée), et nous dirons que $C$ (resp.
$\widehat{C}$) est le \emph{cône affine} (resp. \emph{cône projectif}) défini
par $\mathcal{S}$~; on dira aussi «~cône~» au lieu de «~cône affine~». Par
abus de langage, nous dirons encore que $C$ (resp. $\widehat{C}$) est le
\emph{cône projetant affine de $X$} (resp. le \emph{cône projetant projectif
de $X$}), étant sous-entendu que le préschéma $X$ est donné sous la forme
$\operatorname{Proj}(\mathcal{S})$~; enfin, nous dirons que $\widehat{C}$ est
la \emph{fermeture projective} de $C$ (étant entendu que la donnée de
$\mathcal{S}$ figure dans la structure de $C$).
\end{env}

\begin{proposition}[8.3.2]
\label{II.8.3.2-fr}
Il existe des $Y$-morphismes canoniques
\begin{equation}
\label{II.8.3.2.1-fr}
  Y\xrightarrow{\varepsilon}C\xrightarrow{i}\widehat{C}
\tag{8.3.2.1}
\end{equation}
\begin{equation}
\label{II.8.3.2.2-fr}
  X\xrightarrow{j}\widehat{C}
\tag{8.3.2.2}
\end{equation}

\oldpage[II]{163}
tels que $\varepsilon$ et $j$ soient des immersions fermées, et $i$ un
morphisme affine, qui est une immersion ouverte dominante, pour laquelle
\begin{equation}
\label{II.8.3.2.3-fr}
  i(C)=\widehat{C}-j(X)~;
\tag{8.3.2.3}
\end{equation}
en outre $\widehat{C}$ est le plus petit sous-préschéma fermé de $\widehat{C}$
majorant $i(C)$.
\end{proposition}

Pour définir $i$, on considère l'ouvert de $\widehat{C}$,
\begin{equation}
\label{II.8.3.2.4-fr}
  \widehat{C}_z=
  \operatorname{Spec}(\widehat{\mathcal{S}}/(z-1)\widehat{\mathcal{S}})
\tag{8.3.2.4}
\end{equation}
(\hyperref[II.3.1.4-fr]{3.1.4}), $z$ s'identifiant canoniquement à une
section de $\widehat{\mathcal{S}}$ au-dessus de $Y$. L'isomorphisme
$i:C\xrightarrow{\sim}\widehat{C}_z$ correspond alors à l'isomorphisme
canonique (\hyperref[II.8.2.3.1-fr]{8.2.3.1})
\[
  \widehat{\mathcal{S}}/(z-1)\widehat{\mathcal{S}}
  \xrightarrow{\sim}\mathcal{S}.
\]
Le morphisme $\varepsilon$ correspond à l'homomorphisme d'augmentation
$\mathcal{S}\to\mathcal{S}_0=\mathcal{O}_Y$ ayant pour noyau
$\mathcal{S}_+$ (\hyperref[II.1.2.7-fr]{1.2.7}) et comme ce dernier est
surjectif, $\varepsilon$ est une immersion fermée
(\hyperref[II.1.4.10-fr]{1.4.10}). Enfin, $j$ correspond de même
(\hyperref[II.3.5.1-fr]{3.5.1}) à l'homomorphisme surjectif de degré $0$,
$\widehat{\mathcal{S}}\to\mathcal{S}$, qui se réduit à l'identité dans
$\mathcal{S}$ et est nul dans $z\widehat{\mathcal{S}}$, qui est son noyau~;
$j$ est partout défini et est une immersion fermée en vertu de
(\hyperref[II.3.6.2-fr]{3.6.2}).

Pour démontrer les autres assertions de
(\hyperref[II.8.3.2-fr]{8.3.2}), on peut évidemment se borner au cas où
$Y=\operatorname{Spec}(A)$ est affine, $\mathcal{S}=\widetilde{S}$, où $S$
est une $A$-algèbre graduée, d'où
$\widehat{\mathcal{S}}=(\widehat{S})^{\sim}$~; les éléments homogènes $f$ de
$S_+$ s'identifient alors à des sections de $\widehat{\mathcal{S}}$ au-dessus
de $Y$, et l'ouvert de $\widehat{C}$ noté $D_+(f)$ dans
(\hyperref[II.2.3.3-fr]{2.3.3}) s'écrit aussi $\widehat{C}_f$
(\hyperref[II.3.1.4-fr]{3.1.4})~; de même l'ouvert de $C$ noté $D(f)$ dans
(\hyperref[I.1.1.1-fr]{I, 1.1.1}) s'écrit aussi $C_f$
(\hyperref[0.5.5.2-fr]{0, 5.5.2}). Cela étant, il résulte de
(\hyperref[II.2.3.14-fr]{2.3.14}) et de la définition de $\widehat{S}$ que
dans le cas présent, les ouverts $\widehat{C}_z=i(C)$ et $\widehat{C}_f$
($f$ homogène dans $S_+$) constituent un \emph{recouvrement} de
$\widehat{C}$. En outre, on a, avec ces notations,
\begin{equation}
\label{II.8.3.2.5-fr}
  i^{-1}(\widehat{C}_f)=C_f~;
\tag{8.3.2.5}
\end{equation}
en effet,
$\widehat{C}_f\cap i(C)=\widehat{C}_f\cap\widehat{C}_z
=\widehat{C}_{fz}=\operatorname{Spec}(\widehat{S}_{(fz)})$. Or, si
$d=\deg(f)$, $\widehat{S}_{(fz)}$ est canoniquement isomorphe à
$(\widehat{S}_{(z)})_{f/z^d}$ (\hyperref[II.2.2.2-fr]{2.2.2}), et il résulte
de la définition de l'isomorphisme
(\hyperref[II.8.2.3.1-fr]{8.2.3.1}) que l'image de
$(\widehat{S}_{(z)})_{f/z^d}$ par l'isomorphisme correspondant des anneaux de
fractions est exactement $S_f$. Comme $C_f=\operatorname{Spec}(S_f)$, cela
prouve (\hyperref[II.8.3.2.5-fr]{8.3.2.5}) et montre en même temps que le
morphisme $i$ est affine~; en outre, la restriction de $i$ à $C_f$,
considérée comme morphisme dans $\widehat{C}_f$, correspond
(\hyperref[I.1.7.3-fr]{I, 1.7.3}) à l'homomorphisme canonique
$\widehat{S}_{(f)}\to\widehat{S}_{(fz)}$ et en vertu de ce qui précède et de
(\hyperref[II.8.2.3.2-fr]{8.2.3.2}), on peut énoncer le résultat suivant~:

\begin{env}[8.3.2.6]
\label{II.8.3.2.6-fr}
Si $Y=\operatorname{Spec}(A)$ est affine, et $\mathcal{S}=\widetilde{S}$,
alors, pour tout $f$ homogène dans $S_+$, $\widehat{C}_f$ s'identifie
canoniquement à $\operatorname{Spec}(S_f^{\leq})$, et le morphisme
$C_f\to\widehat{C}_f$ restriction de $i$ correspond alors à l'injection
canonique $S_f^{\leq}\to S_f$.
\end{env}

Notons maintenant que (pour $Y$ affine) le complémentaire de
$\widehat{C}_z$ dans $\widehat{C}=\operatorname{Proj}(\widehat{S})$

\oldpage[II]{164}
est, par définition, l'ensemble des idéaux premiers gradués de $\widehat{S}$
contenant $z$, c'est-à-dire $j(X)$ par définition de $j$, ce qui démontre
(\hyperref[II.8.3.2.3-fr]{8.3.2.3}).

Pour prouver enfin la dernière assertion de
(\hyperref[II.8.3.2-fr]{8.3.2}), on peut toujours supposer $Y$ affine. Avec
les notations précédentes, remarquons que dans l'anneau $\widehat{S}$, $z$
n'est pas diviseur de $0$~; comme $\overline{i(C)}=\widehat{C}$, il suffira
de prouver le

\begin{lemma}[8.3.2.7]
\label{II.8.3.2.7-fr}
Soient $T$ un anneau gradué à degrés positifs,
$Z=\operatorname{Proj}(T)$, $g$ un élément homogène de $T$ de degré $d>0$.
Si $g$ n'est pas diviseur de $0$ dans $T$, $Z$ est le plus petit
sous-préschéma fermé de $Z$ qui majore $Z_g=D_+(g)$.
\end{lemma}

En vertu de (\hyperref[I.4.1.9-fr]{I, 4.1.9}) la question est locale sur $Z$~;
pour tout élément homogène $h\in T_e$ ($e>0$), il suffit donc de prouver que
$Z_h$ est le plus petit sous-préschéma fermé de $Z_h$ qui majore $Z_{gh}$~;
il résulte des définitions et de (\hyperref[I.4.3.2-fr]{I, 4.3.2}) que cette
condition équivaut au fait que l'homomorphisme canonique
$T_{(h)}\to T_{(gh)}$ est \emph{injectif}. Or cet homomorphisme s'identifie à
l'homomorphisme canonique
$T_{(h)}\to(T_{(h)})_{g^e/h^d}$ (\hyperref[II.2.2.3-fr]{2.2.3}). Mais comme
$g^e$ n'est pas diviseur de $0$ dans $T$, $g^e/h^d$ ne l'est pas dans $T_h$
(ni \emph{a fortiori} dans $T_{(h)}$), car la relation
$(g^e/h^d)(t/h^m)=0$ avec $t\in T$ et $m>0$ entraîne l'existence d'un $n>0$
tel que $h^ng^et=0$, d'où $h^nt=0$, et par suite $t/h^m=0$ dans $T_h$. Cela
achève donc la démonstration (\hyperref[0.1.2.2-fr]{0, 1.2.2}).

\begin{env}[8.3.3]
\label{II.8.3.3-fr}
On identifiera souvent le cône affine $C$ au sous-préschéma induit par le cône
projectif $\widehat{C}$ sur l'ouvert $i(C)$, au moyen de l'immersion ouverte
$i$. Le sous-préschéma fermé de $C$, associé à l'immersion fermée
$\varepsilon$, est appelé le \emph{préschéma des sommets de $C$}~; on dit
aussi que $\varepsilon$, qui est une $Y$-section de $C$, est la \emph{section
sommet} ou la \emph{section nulle de $C$}~; on peut identifier $Y$ au
préschéma des sommets de $C$ au moyen de $\varepsilon$. D'ailleurs
$i\circ\varepsilon$ est une $Y$-section de $\widehat{C}$, donc aussi une
immersion fermée (\hyperref[I.5.4.6-fr]{I, 5.4.6}), correspondant à
l'homomorphisme canonique surjectif de degré $0$,
$\widehat{\mathcal{S}}=\mathcal{S}[z]\to\mathcal{O}_Y[z]$
(cf. (\hyperref[II.3.1.7-fr]{3.1.7})), de noyau
$\mathcal{S}_+[z]=\mathcal{S}_+\widehat{\mathcal{S}}$~; le sous-préschéma de
$\widehat{C}$ associé à cette immersion fermée est encore appelé le
\emph{préschéma des sommets de $\widehat{C}$}, et $i\circ\varepsilon$ la
\emph{section sommet de $\widehat{C}$}~; il peut s'identifier par
$i\circ\varepsilon$ à $Y$. Enfin, le sous-préschéma fermé de $\widehat{C}$
associé à $j$ est appelé le \emph{lieu à l'infini de $\widehat{C}$}, et peut
être identifié à $X$ au moyen de $j$.
\end{env}

\begin{env}[8.3.4]
\label{II.8.3.4-fr}
Les sous-préschémas de $C$ (resp. $\widehat{C}$) induits respectivement sur
les \emph{ouverts}
\begin{equation}
\label{II.8.3.4.1-fr}
  E=C-\varepsilon(Y),\qquad
  \widehat{E}=\widehat{C}-i(\varepsilon(Y))
\tag{8.3.4.1}
\end{equation}
sont appelés respectivement (par abus de langage) le \emph{cône affine
épointé} et le \emph{cône projectif épointé} définis par $\mathcal{S}$~; on
notera qu'en dépit de cette terminologie, $E$ \emph{n'est pas nécessairement
affine} sur $Y$, ni $\widehat{E}$ projectif sur $Y$
(cf. (\hyperref[II.8.4.3-fr]{8.4.3})). Lorsqu'on identifie $C$ à $i(C)$, on a
donc pour les espaces sous-jacents
\begin{equation}
\label{II.8.3.4.2-fr}
  C\cup\widehat{E}=\widehat{C},\qquad C\cap\widehat{E}=E
\tag{8.3.4.2}
\end{equation}
de sorte que $\widehat{C}$ peut être considéré comme obtenu par
\emph{recollement} des sous-préschémas ouverts $C$ et $\widehat{E}$~; en
outre, en vertu de (\hyperref[II.8.3.2.3-fr]{8.3.2.3}),
\begin{equation}
\label{II.8.3.4.3-fr}
  E=\widehat{E}-j(X).
\tag{8.3.4.3}
\end{equation}

\oldpage[II]{165}
Lorsque $Y=\operatorname{Spec}(A)$ est affine, on a, avec les notations de
(\hyperref[II.8.3.2-fr]{8.3.2})
\begin{equation}
\label{II.8.3.4.4-fr}
  E=\bigcup C_f,\qquad
  \widehat{E}=\bigcup\widehat{C}_f,\qquad
  C_f=C\cap\widehat{C}_f
\tag{8.3.4.4}
\end{equation}
$f$ parcourant l'ensemble des éléments homogènes de $S_+$ (ou seulement une
partie $M$ de cet ensemble engendrant un idéal de $S_+$ dont la racine dans
$S_+$ est $S_+$ lui-même, autrement dit, telle que les $X_f$ pour $f\in M$
recouvrent $X$ (\hyperref[II.2.3.14-fr]{2.3.14})). Le recollement de $C$ et de
$\widehat{C}_f$ le long de $C_f$ est alors déterminé par des morphismes
d'injection $C_f\to C$, $C_f\to\widehat{C}_f$, qui, ainsi qu'on l'a vu
(\hyperref[II.8.3.2.6-fr]{8.3.2.6}) correspondent respectivement aux
homomorphismes canoniques $S\to S_f$, $S_f^{\leq}\to S_f$.
\end{env}

\begin{proposition}[8.3.5]
\label{II.8.3.5-fr}
Avec les notations de (\hyperref[II.8.3.1-fr]{8.3.1}) et
(\hyperref[II.8.3.4-fr]{8.3.4}), le morphisme associé
(\hyperref[II.3.5.1-fr]{3.5.1}) à l'injection canonique
$\varphi:\mathcal{S}\to\widehat{\mathcal{S}}=\mathcal{S}[z]$ est un morphisme
affine surjectif (dit \emph{rétraction canonique})
\begin{equation}
\label{II.8.3.5.1-fr}
  p:\widehat{E}\to X
\tag{8.3.5.1}
\end{equation}
tel que l'on ait
\begin{equation}
\label{II.8.3.5.2-fr}
  p\circ j=1_X.
\tag{8.3.5.2}
\end{equation}
\end{proposition}

Pour démontrer la proposition, on peut se borner au cas où $Y$ est affine.
Compte tenu de l'expression (\hyperref[II.8.3.4.4-fr]{8.3.4.4}) de
$\widehat{E}$, le fait que le domaine de définition $G(\varphi)$ de $p$ est
égal à $\widehat{E}$ résultera de la première des assertions suivantes~:

\begin{env}[8.3.5.3]
\label{II.8.3.5.3-fr}
Si $Y=\operatorname{Spec}(A)$ est affine, et $\mathcal{S}=\widetilde{S}$, on
a, pour tout $f\in S_+$ homogène
\begin{equation}
\label{II.8.3.5.4-fr}
  p^{-1}(X_f)=\widehat{C}_f
\tag{8.3.5.4}
\end{equation}
et la restriction de $p$ à
$\widehat{C}_f=\operatorname{Spec}(S_f^{\leq})$, considérée comme morphisme de
$\widehat{C}_f$ dans $X_f$, correspond à l'injection canonique
$S_{(f)}\to S_f^{\leq}$. Si de plus $f\in S_1$, $\widehat{C}_f$ est isomorphe
à $X_f\otimes_{\mathbf{Z}}\mathbf{Z}[T]$ ($T$ indéterminée).
\end{env}

En effet, la formule (\hyperref[II.8.3.5.4-fr]{8.3.5.4}) n'est qu'un cas
particulier de (\hyperref[II.2.8.1.1-fr]{2.8.1.1}) et la seconde assertion
n'est autre que la définition de $\operatorname{Proj}(\varphi)$ lorsque $Y$
est affine (\hyperref[II.2.8.1-fr]{2.8.1}). Alors la formule
(\hyperref[II.8.3.5.2-fr]{8.3.5.2}) et le fait que $p$ est surjectif résultent
de ce que le composé $\mathcal{S}\to\widehat{\mathcal{S}}\to\mathcal{S}$ des
homomorphismes canoniques est l'identité dans $\mathcal{S}$. Enfin, la
dernière assertion de (\hyperref[II.8.3.5.3-fr]{8.3.5.3}) provient de ce que
$S_f^{\leq}$ est isomorphe à $S_{(f)}[T]$ lorsque $f\in S_1$
(\hyperref[II.2.2.1-fr]{2.2.1}).

\begin{corollary}[8.3.6]
\label{II.8.3.6-fr}
La restriction
\begin{equation}
\label{II.8.3.6.1-fr}
  \pi:E\to X
\tag{8.3.6.1}
\end{equation}
de $p$ à $E$ est un morphisme affine surjectif. Si $Y$ est affine et $f$
homogène dans $S_+$, on a
\begin{equation}
\label{II.8.3.6.2-fr}
  \pi^{-1}(X_f)=C_f
\tag{8.3.6.2}
\end{equation}
et la restriction de $\pi$ à $C_f$ correspond à l'injection canonique
$S_{(f)}\to S_f$. Si en outre $f\in S_1$, $C_f$ est isomorphe à
$X_f\otimes_{\mathbf{Z}}\mathbf{Z}[T,T^{-1}]$ ($T$ indéterminée).
\end{corollary}

La formule (\hyperref[II.8.3.6.2-fr]{8.3.6.2}) résulte aussitôt de
(\hyperref[II.8.3.5.3-fr]{8.3.5.3}) et
(\hyperref[II.8.3.2.5-fr]{8.3.2.5}), et démontre la surjectivité de $\pi$~; on
a vu par ailleurs que l'immersion $i$, restreinte à $C_f$, correspondait

\oldpage[II]{166}
à l'injection $S_f^{\leq}\to S_f$ (\hyperref[II.8.3.2-fr]{8.3.2}). Enfin, la
dernière assertion est conséquence de ce que pour $f\in S_1$, $S_f$ est
isomorphe à $S_{(f)}[T,T^{-1}]$ (\hyperref[II.2.2.1-fr]{2.2.1}).

\begin{remark}[8.3.7]
\label{II.8.3.7-fr}
Lorsque $Y$ est affine, les éléments de l'espace sous-jacent à $E$ sont les
idéaux premiers $\mathfrak{p}$ (non nécessairement gradués) de $S$ ne
contenant pas $S_+$, en vertu de la définition de l'immersion $\varepsilon$
(\hyperref[II.8.3.2-fr]{8.3.2}). Pour un tel idéal $\mathfrak{p}$, les
$\mathfrak{p}\cap S_n$ vérifient de façon évidente les conditions de
(\hyperref[II.2.1.9-fr]{2.1.9}), donc il existe un seul idéal premier
\emph{gradué} $\mathfrak{q}$ de $S$ tel que
$\mathfrak{q}\cap S_n=\mathfrak{p}\cap S_n$ pour tout $n$~; l'application
$\pi:E\to X$ des espaces sous-jacents s'interprète alors grâce à la relation
\begin{equation}
\label{II.8.3.7.1-fr}
  \pi(\mathfrak{p})=\mathfrak{q}.
\tag{8.3.7.1}
\end{equation}
En effet, pour vérifier cette relation, il suffit de considérer un $f$
homogène dans $S_+$ tel que $\mathfrak{p}\in D(f)$, et d'observer que
$\mathfrak{q}_{(f)}$ est l'image réciproque de $\mathfrak{p}_f$ par
l'injection $S_{(f)}\to S_f$.
\end{remark}

\begin{corollary}[8.3.8]
\label{II.8.3.8-fr}
Si $\mathcal{S}$ est engendrée par $\mathcal{S}_1$, les morphismes $p$ et
$\pi$ sont de type fini~; pour tout $x\in X$, la fibre $p^{-1}(x)$ est
isomorphe à $\operatorname{Spec}(k(x)[T])$ et la fibre $\pi^{-1}(x)$ est
isomorphe à $\operatorname{Spec}(k(x)[T,T^{-1}])$.
\end{corollary}

Cela résulte aussitôt de (\hyperref[II.8.3.5-fr]{8.3.5}) et
(\hyperref[II.8.3.6-fr]{8.3.6}) en observant que, lorsque $Y$ est affine et
$S$ engendrée par $S_1$, les $X_f$, pour $f\in S_1$, forment un recouvrement
de $X$ (\hyperref[II.2.3.14-fr]{2.3.14}).

\begin{remark}[8.3.9]
\label{II.8.3.9-fr}
Le cône épointé affine correspondant à la $\mathcal{O}_Y$-Algèbre graduée
$\mathcal{O}_Y[T]$ ($T$ indéterminée) s'identifie à
$G_m=\operatorname{Spec}(\mathcal{O}_Y[T,T^{-1}])$, puisqu'il n'est autre que
$C_T$, comme on l'a vu dans (\hyperref[II.8.3.2-fr]{8.3.2}) (voir
(\hyperref[II.8.4.4-fr]{8.4.4}) pour un résultat plus général). Ce préschéma
est canoniquement muni d'une structure de «~\emph{$Y$-schéma en groupes
commutatif}~». Il s'agit là d'une notion qui sera développée en détail plus
tard, et que l'on peut ici résumer rapidement comme suit. Un $Y$-schéma en
groupes est un $Y$-schéma $G$, muni de deux $Y$-morphismes
$p:G\times_YG\to G$ et $s:G\to G$ qui satisfont aux conditions formellement
analogues aux axiomes de la loi de composition et de l'application de
symétrie dans un groupe~: le diagramme
\[
  \xymatrix{
    G\times G\times G \ar[r]^{p\times1}\ar[d]_{1\times p}
      & G\times G \ar[d]^{p} \\
    G\times G \ar[r]_{p} & G
  }
\]
doit être commutatif («~associativité~») et l'on doit avoir une condition qui
correspond dans les groupes au fait que les applications
\[
  (x,y)\mapsto(x,x^{-1},y)\mapsto(x,x^{-1}y)\mapsto x(x^{-1}y)
\]
et
\[
  (x,y)\mapsto(x,x^{-1},y)\mapsto(x,yx^{-1})\mapsto(yx^{-1})x
\]
doivent toutes deux se réduire à $(x,y)\mapsto y$~; la suite de morphismes
correspondant par exemple à la première application composée est
\[
  G\times G\xrightarrow{(1,s)\times1}G\times G\times G
  \xrightarrow{1\times p}G\times G\xrightarrow{p}G
\]
et le lecteur écrira de même la seconde suite.

\oldpage[II]{167}
Il est immédiat (\hyperref[I.3.4.3-fr]{I, 3.4.3}) que la donnée d'une structure
de $Y$-schéma en groupes sur un $Y$-schéma $G$ équivaut à la donnée, pour
\emph{tout} $Y$-préschéma $Z$, d'une structure de \emph{groupe} sur
l'ensemble $\operatorname{Hom}_Y(Z,G)$, ces structures devant être telles que
pour tout $Y$-morphisme $Z\to Z'$, l'application correspondante
$\operatorname{Hom}_Y(Z',G)\to\operatorname{Hom}_Y(Z,G)$ soit un
homomorphisme de groupes. Dans le cas particulier de $G_m$ que nous
considérons ici, $\operatorname{Hom}_Y(Z,G)$ s'identifie à l'ensemble des
$Z$-sections de $Z\times_YG_m$ (\hyperref[I.3.3.14-fr]{I, 3.3.14}), donc à
l'ensemble des $Z$-sections de
$\operatorname{Spec}(\mathcal{O}_Z[T,T^{-1}])$~; finalement, le même
raisonnement que dans (\hyperref[I.3.3.15-fr]{I, 3.3.15}) montre que cet
ensemble s'identifie canoniquement à l'ensemble des éléments \emph{inversibles}
de l'anneau $\Gamma(Z,\mathcal{O}_Z)$, et la structure de groupe sur cet
ensemble est la structure provenant de la multiplication dans l'anneau
$\Gamma(Z,\mathcal{O}_Z)$. Le lecteur pourra vérifier que les morphismes $p$
et $s$ considérés plus haut sont obtenus de la façon suivante~: ils
correspondent d'après (\hyperref[II.1.2.7-fr]{1.2.7}) et
(\hyperref[II.1.4.6-fr]{1.4.6}) à des homomorphismes de $\mathcal{O}_Y$-Algèbres
\begin{align*}
  \pi &: \mathcal{O}_Y[T,T^{-1}]
  \to\mathcal{O}_Y[T,T^{-1},T',T'^{-1}],\\
  \sigma &: \mathcal{O}_Y[T,T^{-1}]\to\mathcal{O}_Y[T,T^{-1}]
\end{align*}
et sont entièrement définis par les données de $\pi(T)=TT'$ et
$\sigma(T)=T^{-1}$.

Cela étant, $G_m$ peut être considéré comme un «~\emph{domaine d'opérateurs
universel}~» pour tout \emph{cône affine}
$C=\operatorname{Spec}(\mathcal{S})$, où $\mathcal{S}$ est une
$\mathcal{O}_Y$-Algèbre graduée quasi-cohérente à degrés positifs. Cela
signifie qu'on peut définir canoniquement un $Y$-morphisme
$G_m\times_YC\to C$ qui a les propriétés formelles d'une loi externe d'un
ensemble muni d'un groupe d'opérateurs~; ou encore, comme ci-dessus pour les
schémas en groupes, on se donne pour tout $Y$-préschéma $Z$ une loi externe
sur $\operatorname{Hom}_Y(Z,C)$, ayant le groupe
$\operatorname{Hom}_Y(Z,G_m)$ comme ensemble d'opérateurs, avec les axiomes
usuels des ensembles munis d'un groupe d'opérateurs, et une condition de
compatibilité avec les $Y$-morphismes $Z\to Z'$. Dans le cas présent, le
morphisme $G_m\times_YC\to C$ est défini par la donnée de l'homomorphisme de
$\mathcal{O}_Y$-Algèbres
\[
  \mathcal{S}\to
  \mathcal{S}\otimes_{\mathcal{O}_Y}\mathcal{O}_Y[T,T^{-1}]
  =\mathcal{S}[T,T^{-1}]
\]
qui, à toute section $s_n\in\Gamma(U,\mathcal{S}_n)$ ($U$ ouvert de $Y$) fait
correspondre la section
$s_nT^n\in\Gamma(U,\mathcal{S}\otimes_{\mathcal{O}_Y}
\mathcal{O}_Y[T,T^{-1}])$.

Inversement, supposons donnée une $\mathcal{O}_Y$-Algèbre quasi-cohérente
$\mathcal{S}$ \emph{non graduée a priori}, et, sur
$C=\operatorname{Spec}(\mathcal{S})$, une structure de «~$Y$-schéma en
ensembles \emph{munis de groupes d'opérateurs}~» ayant pour domaine
d'opérateurs le $Y$-schéma en groupes $G_m$~; alors on en déduit
canoniquement une \emph{graduation} de $\mathcal{O}_Y$-Algèbre sur
$\mathcal{S}$. En effet, la donnée d'un $Y$-morphisme $G_m\times_YC\to C$
équivaut à celle d'un homomorphisme de $\mathcal{O}_Y$-Algèbres
$\psi:\mathcal{S}\to\mathcal{S}[T,T^{-1}]$, qui s'écrira donc
$\psi=\sum_{n\in\mathbf{Z}}\psi_nT^n$, où les
$\psi_n:\mathcal{S}\to\mathcal{S}$ sont des homomorphismes de
$\mathcal{O}_Y$-Modules (avec $\psi_n(s)=0$ sauf pour un nombre fini d'indices
pour toute section $s\in\Gamma(U,\mathcal{S})$, $U$ ouvert dans $Y$). On peut
alors vérifier que les axiomes des ensembles munis d'un groupe d'opérateurs
entraînent que les $\psi_n(\mathcal{S})=\mathcal{S}_n$ définissent une
graduation (à degrés positifs ou négatifs) de $\mathcal{O}_Y$-Algèbre sur
$\mathcal{S}$, les $\psi_n$ étant les projecteurs correspondants. On a ainsi
une notion de structure de «~\emph{cône affine}~» sur tout $Y$-schéma affine,
définie de façon «~géométrique~» sans appel à une graduation préalable.

\oldpage[II]{168}
Nous ne développerons pas davantage ce point de vue ici et laissons au
lecteur le soin de formuler de façon précise les définitions et résultats qui
correspondent aux indications précédentes.
\end{remark}

\subsection{Fermeture projective d'un fibré vectoriel}
\label{subsection:II.8.4-fr}

\begin{env}[8.4.1]
\label{II.8.4.1-fr}
Soient $Y$ un préschéma, $\mathcal{E}$ un $\mathcal{O}_Y$-Module
quasi-cohérent. Si on prend pour $\mathcal{S}$ la $\mathcal{O}_Y$-Algèbre
graduée $\mathbf{S}_{\mathcal{O}_Y}(\mathcal{E})$, la définition
(\hyperref[II.8.3.1.1-fr]{8.3.1.1}) montre que
$\widehat{\mathcal{S}}$ s'identifie à
$\mathbf{S}_{\mathcal{O}_Y}(\mathcal{E}\oplus\mathcal{O}_Y)$. Le cône affine
$\operatorname{Spec}(\mathcal{S})$ défini par $\mathcal{S}$ étant alors par
définition $\mathbf{V}(\mathcal{E})$, et $\operatorname{Proj}(\mathcal{S})$
étant par définition $\mathbf{P}(\mathcal{E})$, on voit donc que~:
\end{env}

\begin{proposition}[8.4.2]
\label{II.8.4.2-fr}
La fermeture projective d'un fibré vectoriel $\mathbf{V}(\mathcal{E})$ sur
$Y$ est canoniquement isomorphe à
$\mathbf{P}(\mathcal{E}\oplus\mathcal{O}_Y)$, et le lieu à l'infini de cette
dernière est canoniquement isomorphe à $\mathbf{P}(\mathcal{E})$.
\end{proposition}

\begin{remark}[8.4.3]
\label{II.8.4.3-fr}
Prenons par exemple $\mathcal{E}=\mathcal{O}_Y^r$ avec $r\geq2$~; alors les
cônes épointés $E$, $\widehat{E}$ définis par $\mathcal{S}$ ne sont ni affines
ni projectifs sur $Y$ si $Y\neq\varnothing$. La seconde assertion est
immédiate, car $\widehat{C}=\mathbf{P}(\mathcal{O}_Y^{r+1})$ est projectif sur
$Y$ et les espaces sous-jacents à $E$ et $\widehat{E}$ sont des ouverts non
fermés dans $\widehat{C}$, donc les immersions canoniques
$E\to\widehat{E}$ et $\widehat{E}\to\widehat{C}$ ne sont pas projectives
(\hyperref[II.5.5.3-fr]{5.5.3}), et on conclut par
(\hyperref[II.5.5.5-fr]{5.5.5, (v)}). D'autre part, supposant
$Y=\operatorname{Spec}(A)$ affine et par exemple $r=2$, on a
$C=\operatorname{Spec}(A[T_1,T_2])$ et $E$ est alors le préschéma induit par
$C$ sur l'ouvert $D(T_1)\cup D(T_2)$~; or nous avons vu que ce dernier n'est
pas affine (\hyperref[I.5.5.11-fr]{I, 5.5.11})~; 
\emph{a fortiori} $\widehat{E}$ ne peut être affine, puisque $E$ est l'ouvert
où ne s'annule pas la section $z$ au-dessus de $\widehat{E}$
(\hyperref[II.8.3.2-fr]{8.3.2}).

Toutefois~:
\end{remark}

\begin{proposition}[8.4.4]
\label{II.8.4.4-fr}
Si $\mathcal{L}$ est un $\mathcal{O}_Y$-Module inversible, on a des
isomorphismes canoniques pour les cônes épointés $E$ et $\widehat{E}$
correspondants à $C=\mathbf{V}(\mathcal{L})$,
\begin{equation}
\label{II.8.4.4.1-fr}
  \operatorname{Spec}\left(
    \bigoplus_{n\in\mathbf{Z}}\mathcal{L}^{\otimes n}
  \right)\xrightarrow{\sim}E
\tag{8.4.4.1}
\end{equation}
\begin{equation}
\label{II.8.4.4.2-fr}
  \mathbf{V}(\mathcal{L}^{-1})\xrightarrow{\sim}\widehat{E}.
\tag{8.4.4.2}
\end{equation}

En outre, il existe un isomorphisme canonique de la fermeture projective de
$\mathbf{V}(\mathcal{L})$ sur celle de $\mathbf{V}(\mathcal{L}^{-1})$,
transformant la section nulle (resp. le lieu à l'infini) de la première en le
lieu à l'infini (resp. la section nulle) de la seconde.
\end{proposition}

\begin{proof}
On a ici
$\mathcal{S}=\bigoplus_{n\geq0}\mathcal{L}^{\otimes n}$~; l'injection
canonique
\[
  \mathcal{S}\to\bigoplus_{n\in\mathbf{Z}}\mathcal{L}^{\otimes n}
\]
définit un morphisme dominant canonique
\begin{equation}
\label{II.8.4.4.3-fr}
  \operatorname{Spec}\left(
    \bigoplus_{n\in\mathbf{Z}}\mathcal{L}^{\otimes n}
  \right)
  \to\mathbf{V}(\mathcal{L})
  =\operatorname{Spec}\left(
    \bigoplus_{n\geq0}\mathcal{L}^{\otimes n}
  \right)
\tag{8.4.4.3}
\end{equation}
et il suffit de prouver que ce morphisme est un isomorphisme du schéma
$\operatorname{Spec}(\bigoplus_{n\in\mathbf{Z}}\mathcal{L}^{\otimes n})$
sur $E$. La question étant locale sur $Y$, on peut supposer que
$Y=\operatorname{Spec}(A)$ est affine

\oldpage[II]{169}
et $\mathcal{L}=\mathcal{O}_Y$, donc $\mathcal{S}=(A[T])^{\sim}$ et
$\bigoplus_{n\in\mathbf{Z}}\mathcal{L}^{\otimes n}
=(A[T,T^{-1}])^{\sim}$. Or $A[T,T^{-1}]$ est l'anneau de fractions
$A[T]_T$ de $A[T]$, donc (\hyperref[II.8.4.4.3-fr]{8.4.4.3}) identifie le
premier membre au préschéma induit par $C=\mathbf{V}(\mathcal{L})$ sur
l'ouvert $D(T)$~; le complémentaire $V(T)$ de cet ouvert dans $C$ est
l'espace sous-jacent au sous-préschéma fermé de $C$ défini par l'idéal
$TA[T]$, c'est-à-dire la section nulle de $C$, donc $E=D(T)$.

L'isomorphisme (\hyperref[II.8.4.4.2-fr]{8.4.4.2}) sera d'autre part
conséquence de la dernière assertion, $\mathbf{V}(\mathcal{L}^{-1})$ étant le
complémentaire du lieu à l'infini de sa fermeture projective et $\widehat{E}$
le complémentaire de la section nulle de la fermeture projective de
$C=\mathbf{V}(\mathcal{L})$. Or, ces fermetures projectives sont
respectivement $\mathbf{P}(\mathcal{L}^{-1}\oplus\mathcal{O}_Y)$ et
$\mathbf{P}(\mathcal{L}\oplus\mathcal{O}_Y)$~; mais on peut écrire
$\mathcal{L}\oplus\mathcal{O}_Y
=\mathcal{L}\otimes(\mathcal{L}^{-1}\oplus\mathcal{O}_Y)$. L'existence de
l'isomorphisme canonique cherché résulte alors de
(\hyperref[II.4.1.4-fr]{4.1.4}), et tout revient à voir que cet isomorphisme
échange sections nulles et lieux à l'infini. On peut pour cela se borner au
cas où $Y=\operatorname{Spec}(A)$ est affine,
$L=Ac$, $L^{-1}=Ac'$, l'isomorphisme canonique
$L\otimes L^{-1}\to A$ faisant correspondre à $c\otimes c'$ l'élément $1$
de $A$. Alors $\mathbf{S}(L\oplus A)$ est produit tensoriel de $A[z]$ et de
$\bigoplus_{n\geq0}Ac^{\otimes n}$, $\mathbf{S}(L^{-1}\oplus A)$ produit
tensoriel de $A[z]$ et de $\bigoplus_{n\geq0}Ac'^{\otimes n}$, et
l'isomorphisme défini dans (\hyperref[II.4.1.4-fr]{4.1.4}) fait correspondre à
$z^h\otimes c'^{\otimes(n-h)}$ l'élément
$z^{n-h}\otimes c^{\otimes h}$. Or, dans
$\mathbf{P}(\mathcal{L}^{-1}\oplus\mathcal{O}_Y)$, le lieu à l'infini est
l'ensemble des points où s'annule la section $z$ et la section nulle est
l'ensemble des points où s'annule la section $c'$~; comme on a des définitions
analogues pour $\mathbf{P}(\mathcal{L}\oplus\mathcal{O}_Y)$, notre conclusion
résulte aussitôt de l'explicitation qui précède.
\end{proof}

\subsection{Comportements fonctoriels}
\label{subsection:II.8.5-fr}

\begin{env}[8.5.1]
\label{II.8.5.1-fr}
Soient $Y$, $Y'$ deux préschémas, $q:Y'\to Y$ un morphisme,
$\mathcal{S}$ (resp. $\mathcal{S}'$) une $\mathcal{O}_Y$-Algèbre
(resp. une $\mathcal{O}_{Y'}$-Algèbre) graduée quasi-cohérente à degrés
\emph{positifs}. Considérons un $q$-morphisme d'Algèbres graduées
\begin{equation}
\label{II.8.5.1.1-fr}
  \varphi:\mathcal{S}\to\mathcal{S}'.
\tag{8.5.1.1}
\end{equation}
On sait (\hyperref[II.1.5.6-fr]{1.5.6}) qu'il lui correspond canoniquement un
morphisme
\[
  \Phi=\operatorname{Spec}(\varphi):
  \operatorname{Spec}(\mathcal{S}')\to\operatorname{Spec}(\mathcal{S})
\]
tel que le diagramme
\begin{equation}
\label{II.8.5.1.2-fr}
  \xymatrix{
    C'\ar[r]^{\Phi}\ar[d] & C\ar[d]\\
    Y'\ar[r]_{q} & Y
  }
\tag{8.5.1.2}
\end{equation}
où on a posé $C=\operatorname{Spec}(\mathcal{S})$,
$C'=\operatorname{Spec}(\mathcal{S}')$, soit commutatif. 
\emph{Supposons en outre que
$\mathcal{S}_0=\mathcal{O}_Y$, $\mathcal{S}'_0=\mathcal{O}_{Y'}$}~; soient
$\varepsilon:Y\to C$ et $\varepsilon':Y'\to C'$ les immersions canoniques
(\hyperref[II.8.3.2-fr]{8.3.2})~; on a alors un diagramme commutatif
\begin{equation}
\label{II.8.5.1.3-fr}
  \xymatrix{
    Y'\ar[r]^{q}\ar[d]_{\varepsilon'} & Y\ar[d]^{\varepsilon}\\
    C'\ar[r]_{\Phi} & C
  }
\tag{8.5.1.3}
\end{equation}

\oldpage[II]{170}
qui correspond au diagramme
\[
  \xymatrix{
    \mathcal{S}\ar[r]^{\varphi}\ar[d] & \mathcal{S}'\ar[d]\\
    \mathcal{O}_Y\ar[r] & \mathcal{O}_{Y'}
  }
\]
où les flèches verticales sont les homomorphismes d'augmentation, et dont la
commutativité résulte de l'hypothèse que $\varphi$ est supposé être un
homomorphisme d'Algèbres \emph{graduées}.
\end{env}

\begin{proposition}[8.5.2]
\label{II.8.5.2-fr}
Si $E$ (resp. $E'$) est le cône épointé affine défini par $\mathcal{S}$
(resp. $\mathcal{S}'$), on a $\Phi^{-1}(E)\subset E'$~; si en outre
$\operatorname{Proj}(\varphi):G(\varphi)\to\operatorname{Proj}(\mathcal{S})$
est partout défini (autrement dit si
$G(\varphi)=\operatorname{Proj}(\mathcal{S}')$), alors on a
$\Phi^{-1}(E)=E'$, et réciproquement.
\end{proposition}

\begin{proof}
La première assertion résulte de la commutativité de
(\hyperref[II.8.5.1.3-fr]{8.5.1.3}). Pour démontrer la seconde, on peut se
borner au cas où $Y=\operatorname{Spec}(A)$ et
$Y'=\operatorname{Spec}(A')$ sont affines,
$\mathcal{S}=\widetilde{S}$, $\mathcal{S}'=\widetilde{S'}$. Pour tout $f$
homogène dans $S_+$, si on pose $f'=\varphi(f)$, on a
$\Phi^{-1}(C_f)=C'_{f'}$ (\hyperref[I.2.2.4.1-fr]{I, 2.2.4.1})~; dire que
$G(\varphi)=\operatorname{Proj}(S')$ signifie que dans $S'_+$ la racine de
l'idéal engendré par les $f'=\varphi(f)$ est $S'_+$ lui-même
(\hyperref[II.2.8.1-fr]{2.8.1}) et (\hyperref[II.2.3.14-fr]{2.3.14}), et cela
équivaut encore à dire que les $C'_{f'}$ recouvrent $E'$
(\hyperref[II.8.3.4.4-fr]{8.3.4.4}).
\end{proof}

\begin{env}[8.5.3]
\label{II.8.5.3-fr}
Le $q$-morphisme $\varphi$ se prolonge canoniquement en un $q$-morphisme
d'Algèbres graduées
\begin{equation}
\label{II.8.5.3.1-fr}
  \widehat{\varphi}:\widehat{\mathcal{S}}
  \to\widehat{\mathcal{S}'}
\tag{8.5.3.1}
\end{equation}
en posant $\widehat{\varphi}(z)=z$. On en déduit par suite un morphisme
\[
  \widehat{\Phi}=\operatorname{Proj}(\widehat{\varphi}):
  G(\widehat{\varphi})\to\widehat{C}
  =\operatorname{Proj}(\widehat{\mathcal{S}})
\]
tel que le diagramme
\[
  \xymatrix{
    G(\widehat{\varphi})\ar[r]^{\widehat{\Phi}}\ar[d] &
    \widehat{C}\ar[d]\\
    Y'\ar[r]_{q} & Y
  }
\]
soit commutatif (\hyperref[II.3.5.6-fr]{3.5.6}). Il résulte aussitôt des
définitions que si on désigne par $i:C\to\widehat{C}$ et
$i':C'\to\widehat{C'}$ les immersions ouvertes canoniques
(\hyperref[II.8.3.2-fr]{8.3.2}), on a
$i'(C')\subset G(\widehat{\varphi})$ et le diagramme
\begin{equation}
\label{II.8.5.3.2-fr}
  \xymatrix{
    C'\ar[r]^{\Phi}\ar[d]_{i'} & C\ar[d]^{i}\\
    G(\widehat{\varphi})\ar[r]_{\widehat{\Phi}} & \widehat{C}
  }
\tag{8.5.3.2}
\end{equation}
est commutatif. Enfin, si on pose
$X=\operatorname{Proj}(\mathcal{S})$, $X'=\operatorname{Proj}(\mathcal{S}')$,
et si $j:X\to\widehat{C}$, $j':X'\to\widehat{C'}$ sont les immersions
fermées canoniques (\hyperref[II.8.3.2-fr]{8.3.2}), il résulte encore des
définitions de ces immersions que l'on a
$j'(G(\varphi))\subset G(\widehat{\varphi})$ et que le diagramme

\oldpage[II]{171}
\begin{equation}
\label{II.8.5.3.3-fr}
  \xymatrix{
    G(\varphi)\ar[r]^{\operatorname{Proj}(\varphi)}\ar[d]_{j'} &
    X\ar[d]^{j}\\
    G(\widehat{\varphi})\ar[r]_{\widehat{\Phi}} & \widehat{C}
  }
\tag{8.5.3.3}
\end{equation}
est commutatif.
\end{env}

\begin{proposition}[8.5.4]
\label{II.8.5.4-fr}
Si $\widehat{E}$ (resp. $\widehat{E'}$) est le cône épointé projectif défini
par $\mathcal{S}$ (resp. $\mathcal{S}'$), on a
$\widehat{\Phi}^{-1}(\widehat{E})\subset\widehat{E'}$~; en outre, si
$p:\widehat{E}\to X$ et $p':\widehat{E'}\to X'$ sont les rétractions
canoniques, on a
$p'(\widehat{\Phi}^{-1}(\widehat{E}))\subset G(\varphi)$, et le diagramme
\begin{equation}
\label{II.8.5.4.1-fr}
  \xymatrix{
    \widehat{\Phi}^{-1}(\widehat{E})
      \ar[r]^{\widehat{\Phi}}\ar[d]_{p'} &
    \widehat{E}\ar[d]^{p}\\
    G(\varphi)\ar[r]_{\operatorname{Proj}(\varphi)} & X
  }
\tag{8.5.4.1}
\end{equation}
est commutatif. Si $\operatorname{Proj}(\varphi)$ est partout défini, il en
est de même de $\widehat{\Phi}$, et on a
$\widehat{\Phi}^{-1}(\widehat{E})=\widehat{E'}$.
\end{proposition}

\begin{proof}
La première assertion résulte de la commutativité des diagrammes
(\hyperref[II.8.5.1.3-fr]{8.5.1.3}) et
(\hyperref[II.8.5.3.2-fr]{8.5.3.2}), et les deux suivantes de la définition
des rétractions canoniques (\hyperref[II.8.3.5-fr]{8.3.5}) et de la définition
de $\widehat{\varphi}$. D'autre part, pour voir que $\widehat{\Phi}$ est
partout défini lorsqu'il en est ainsi de $\operatorname{Proj}(\varphi)$, on
peut se borner à considérer le cas où $Y=\operatorname{Spec}(A)$ et
$Y'=\operatorname{Spec}(A')$ sont affines,
$\mathcal{S}=\widetilde{S}$, $\mathcal{S}'=\widetilde{S'}$~; l'hypothèse est
que lorsque $f$ parcourt l'ensemble des éléments homogènes de $S_+$, l'idéal
engendré dans $S'_+$ par les $\varphi(f)$ a pour racine dans $S'_+$ l'idéal
$S'_+$ lui-même~; on en conclut aussitôt que la racine dans $(S'[z])_+$ de
l'idéal engendré par $z$ et les $\varphi(f)$ est $(S'[z])_+$ lui-même, d'où
notre assertion~; cela prouve en même temps que $\widehat{E'}$ est réunion des
$\widehat{C'}_{\varphi(f)}$, donc égal à
$\widehat{\Phi}^{-1}(\widehat{E})$.
\end{proof}

\begin{corollary}[8.5.5]
\label{II.8.5.5-fr}
Lorsque $\operatorname{Proj}(\varphi)$ est partout défini, l'image réciproque
par $\widehat{\Phi}$ de l'espace sous-jacent au lieu à l'infini (resp. au
préschéma des sommets) de $\widehat{C}$ est l'espace sous-jacent au lieu à
l'infini (resp. au préschéma des sommets) de $\widehat{C'}$.
\end{corollary}

\begin{proof}
Cela résulte aussitôt de (\hyperref[II.8.5.4-fr]{8.5.4}) et
(\hyperref[II.8.5.2-fr]{8.5.2}), compte tenu des relations
(\hyperref[II.8.3.4.1-fr]{8.3.4.1}) et
(\hyperref[II.8.3.4.2-fr]{8.3.4.2}).
\end{proof}

\subsection{Un isomorphisme canonique pour les cônes épointés}
\label{subsection:II.8.6-fr}

\begin{env}[8.6.1]
\label{II.8.6.1-fr}
Soient $Y$ un préschéma, $\mathcal{S}$ une $\mathcal{O}_Y$-Algèbre graduée
quasi-cohérente à degrés positifs \emph{telle que
$\mathcal{S}_0=\mathcal{O}_Y$}, $X$ le $Y$-schéma
$\operatorname{Proj}(\mathcal{S})$. Nous allons appliquer les résultats de
(\hyperref[subsection:II.8.5-fr]{8.5}) au cas où on prend $Y'=X$,
$q:X\to Y$ étant le morphisme structural~; soit
\begin{equation}
\label{II.8.6.1.1-fr}
  \mathcal{S}_X=\bigoplus_{n\in\mathbf{Z}}\mathcal{O}_X(n)
\tag{8.6.1.1}
\end{equation}

\oldpage[II]{172}
qui est une $\mathcal{O}_X$-Algèbre graduée quasi-cohérente, la multiplication
y étant définie au moyen des homomorphismes canoniques
(\hyperref[II.3.2.6.1-fr]{3.2.6.1})
\[
  \mathcal{O}_X(m)\otimes_{\mathcal{O}_X}\mathcal{O}_X(n)
  \to\mathcal{O}_X(m+n)
\]
dont l'associativité est garantie par le diagramme commutatif
(\hyperref[II.2.5.11.4-fr]{2.5.11.4}). Prenons pour $\mathcal{S}'$ la
sous-$\mathcal{O}_X$-Algèbre graduée quasi-cohérente
$\mathcal{S}_X^{\geq}=\bigoplus_{n\geq0}\mathcal{O}_X(n)$ de
$\mathcal{S}_X$, à degrés positifs.

Enfin, nous considérerons le $q$-morphisme canonique
\begin{equation}
\label{II.8.6.1.2-fr}
  \alpha:\mathcal{S}\to\mathcal{S}_X^{\geq}
\tag{8.6.1.2}
\end{equation}
défini dans (\hyperref[II.3.3.2.3-fr]{3.3.2.3}) comme un homomorphisme
$\mathcal{S}\to q_*(\mathcal{S}_X)$, mais qui applique évidemment
$\mathcal{S}$ dans $q_*(\mathcal{S}_X^{\geq})$. Nous poserons
\begin{equation}
\label{II.8.6.1.3-fr}
  C_X=\operatorname{Spec}(\mathcal{S}_X^{\geq}),\qquad
  \widehat{C}_X=\operatorname{Proj}(\mathcal{S}_X^{\geq}[z]),\qquad
  X'=\operatorname{Proj}(\mathcal{S}_X^{\geq})
\tag{8.6.1.3}
\end{equation}
et nous désignerons par $E_X$ et $\widehat{E}_X$ le cône épointé affine et le
cône épointé projectif correspondants~; on notera
$\varepsilon_X:X\to C_X$, $i_X:C_X\to\widehat{C}_X$,
$j_X:X'\to\widehat{C}_X$, $p_X:\widehat{E}_X\to X'$,
$\pi_X:E_X\to X'$ les morphismes canoniques définis dans
(\hyperref[subsection:II.8.3-fr]{8.3}).
\end{env}

\begin{proposition}[8.6.2]
\label{II.8.6.2-fr}
Le morphisme structural $u:X'\to X$ est un isomorphisme, et le morphisme
$\operatorname{Proj}(\alpha)$ est partout défini et identique à $u$. Le
morphisme
$\operatorname{Proj}(\widehat{\alpha}):\widehat{C}_X\to\widehat{C}$ est
partout défini, et ses restrictions à $\widehat{E}_X$ et $E_X$ sont des
isomorphismes sur $\widehat{E}$ et $E$ respectivement. Enfin, lorsqu'on
identifie $X'$ et $X$ au moyen de $u$, les morphismes $p_X$ et $\pi_X$
s'identifient aux morphismes structuraux des $X$-préschémas $\widehat{E}_X$ et
$E_X$.
\end{proposition}

\begin{proof}
On peut évidemment se borner au cas où $Y=\operatorname{Spec}(A)$ est affine
et $\mathcal{S}=\widetilde{S}$~; alors $X$ est réunion des ouverts affines
$X_f$, où $f$ parcourt l'ensemble des éléments homogènes de $S_+$, l'anneau
de $X_f$ étant $S_{(f)}$. Il résulte en outre de
(\hyperref[II.8.2.7.1-fr]{8.2.7.1}) que l'on a
\begin{equation}
\label{II.8.6.2.1-fr}
  \Gamma(X_f,\mathcal{S}_X^{\geq})=S_f^{\geq}.
\tag{8.6.2.1}
\end{equation}

On a donc $u^{-1}(X_f)=\operatorname{Proj}(S_f^{\geq})$. Mais si
$f\in S_d$ ($d>0$), $\operatorname{Proj}(S_f^{\geq})$ est canoniquement
isomorphe à $\operatorname{Proj}((S_f^{\geq})^{(d)})$
(\hyperref[II.2.4.7-fr]{2.4.7}), et d'autre part
$(S_f^{\geq})^{(d)}=(S^{(d)})_f^{\geq}$ s'identifie à $S_{(f)}[T]$
(\hyperref[II.2.2.1-fr]{2.2.1}) par l'application $T\to f/1$~; on en conclut
(\hyperref[II.3.1.7-fr]{3.1.7}) que le morphisme structural
$u^{-1}(X_f)\to X_f$ est un isomorphisme, d'où la première assertion. Pour
démontrer la seconde, remarquons que la restriction
$u^{-1}(X_f)\cap G(\alpha)\to X=\operatorname{Proj}(S)$ de
$\operatorname{Proj}(\alpha)$ correspond à l'application canonique
$x\to x/1$ de $S$ dans $S_f^{\geq}$
(\hyperref[II.2.6.2-fr]{2.6.2})~; on en déduit tout d'abord que
$G(\alpha)=X'$, puis, en tenant compte de ce que
$u^{-1}(X_f)=(u^{-1}(X_f))_{f/1}$, il résulte de
(\hyperref[II.2.8.1.1-fr]{2.8.1.1}) que l'image de $u^{-1}(X_f)$ par
$\operatorname{Proj}(\alpha)$ est contenue dans $X_f$, et la restriction de
$\operatorname{Proj}(\alpha)$ à $u^{-1}(X_f)$, considérée comme morphisme
dans $X_f=\operatorname{Spec}(S_{(f)})$, est bien identique à celle de $u$.
Enfin, la formule (\hyperref[II.8.3.5.4-fr]{8.3.5.4}), appliquée à $p_X$ au
lieu de $p$, montre que
$p_X^{-1}(u^{-1}(X_f))=
\operatorname{Spec}((S_f^{\geq})_{f/1}^{\leq})$, et cet ouvert est, d'après
(\hyperref[II.8.5.4.1-fr]{8.5.4.1}), l'image réciproque par
$\operatorname{Proj}(\widehat{\alpha})$ de
$p^{-1}(X_f)=\operatorname{Spec}(S_f^{\leq})$
(\hyperref[II.8.3.5.3-fr]{8.3.5.3}). Compte tenu de
(\hyperref[II.8.2.3.2-fr]{8.2.3.2}), la restriction de
$\operatorname{Proj}(\widehat{\alpha})$ à
$p_X^{-1}(u^{-1}(X_f))$ correspond à l'isomorphisme réciproque de
(\hyperref[II.8.2.7.2-fr]{8.2.7.2}), restreint à $S_f^{\leq}$, d'où le
troisième point~; la dernière assertion est évidente par définition.

\oldpage[II]{173}
On notera aussi qu'il résulte du diagramme commutatif
(\hyperref[II.8.5.3.2-fr]{8.5.3.2}) que \emph{la restriction à $C_X$ de
$\operatorname{Proj}(\widehat{\alpha})$ n'est autre que le morphisme
$\operatorname{Spec}(\alpha)$}.
\end{proof}

\begin{corollary}[8.6.3]
\label{II.8.6.3-fr}
Considérés comme $X$-schémas, $\widehat{E}_X$ est canoniquement isomorphe à
$\operatorname{Spec}(\mathcal{S}_X^{\leq})$, et $E_X$ à
$\operatorname{Spec}(\mathcal{S}_X)$.
\end{corollary}

\begin{proof}
Comme on sait que les morphismes $p_X$ et $\pi_X$ sont affines
(\hyperref[II.8.3.5-fr]{8.3.5} et
\hyperref[II.8.3.6-fr]{8.3.6}), il suffit (vu
(\hyperref[II.1.3.1-fr]{1.3.1})) de vérifier le corollaire lorsque
$Y=\operatorname{Spec}(A)$ est affine et $\mathcal{S}=\widetilde{S}$. La
première assertion résulte de l'existence de l'isomorphisme canonique
(\hyperref[II.8.2.7.2-fr]{8.2.7.2})
$(S_f^{\geq})_{f/1}^{\leq}\xrightarrow{\sim}S_f^{\leq}$ et du fait que ces
isomorphismes sont compatibles avec le passage de $f$ à $fg$ ($f$ et $g$
homogènes dans $S_+$). De même, la formule
(\hyperref[II.8.3.6.2-fr]{8.3.6.2}), appliquée à $\pi_X$ au lieu de $\pi$,
montre que
$\pi_X^{-1}(u^{-1}(X_f))=
\operatorname{Spec}((S_f^{\geq})_{f/1})$ pour $f$ homogène dans $S_+$, et la
seconde assertion découle de l'existence de l'isomorphisme canonique
(\hyperref[II.8.2.7.2-fr]{8.2.7.2})
$(S_f^{\geq})_{f/1}\xrightarrow{\sim}S_f$.

On peut ainsi dire que $\widehat{C}_X$, considéré comme $X$-schéma, s'obtient
par \emph{recollement} des deux $X$-schémas affines sur $X$,
$C_X=\operatorname{Spec}(\mathcal{S}_X^{\geq})$ et
$\widehat{E}_X=\operatorname{Spec}(\mathcal{S}_X^{\leq})$, dont
l'intersection est l'ouvert $E_X=\operatorname{Spec}(\mathcal{S}_X)$.
\end{proof}

\begin{corollary}[8.6.4]
\label{II.8.6.4-fr}
Supposons que $\mathcal{O}_X(1)$ soit un $\mathcal{O}_X$-Module inversible et
que $\mathcal{S}_X$ soit isomorphe à
$\bigoplus_{n\in\mathbf{Z}}(\mathcal{O}_X(1))^{\otimes n}$ (ce qui sera en
particulier le cas si $\mathcal{S}$ est engendrée par $\mathcal{S}_1$
(\hyperref[II.3.2.5-fr]{3.2.5} et
\hyperref[II.3.2.7-fr]{3.2.7})). Alors le cône projectif épointé
$\widehat{E}$ s'identifie au fibré vectoriel de rang un
$\mathbf{V}(\mathcal{O}_X(-1))$ sur $X$, et le cône affine épointé $E$ au
sous-préschéma de ce fibré vectoriel induit sur le complémentaire de la
section nulle. Avec cette identification, la rétraction canonique
$\widehat{E}\to X$ s'identifie au morphisme structural du $X$-schéma
$\mathbf{V}(\mathcal{O}_X(-1))$. Enfin, il existe un $Y$-morphisme canonique
$\mathbf{V}(\mathcal{O}_X(1))\to C$, dont la restriction au complémentaire de
la section nulle de $\mathbf{V}(\mathcal{O}_X(1))$ est un isomorphisme de ce
complémentaire sur le cône affine épointé $E$.
\end{corollary}

\begin{proof}
En effet, si on pose $\mathcal{L}=\mathcal{O}_X(1)$,
$\mathcal{S}_X^{\geq}$ est alors identique à
$\mathbf{S}_{\mathcal{O}_X}(\mathcal{L})$, donc $\widehat{E}_X$ s'identifie
canoniquement à $\mathbf{V}(\mathcal{L}^{-1})$ en vertu de
(\hyperref[II.8.6.3-fr]{8.6.3}) et $C_X$ à $\mathbf{V}(\mathcal{L})$. Le
morphisme $\mathbf{V}(\mathcal{L})\to C$ est la restriction de
$\operatorname{Proj}(\widehat{\alpha})$, et les assertions du corollaire sont
donc des cas particuliers de (\hyperref[II.8.6.2-fr]{8.6.2}).
\end{proof}

On notera que l'image réciproque par le morphisme
$\mathbf{V}(\mathcal{O}_X(1))\to C$ de l'espace sous-jacent au préschéma des
sommets de $C$ est l'espace sous-jacent à la section nulle de
$\mathbf{V}(\mathcal{O}_X(1))$ (\hyperref[II.8.5.5-fr]{8.5.5})~; mais en
général les sous-préschémas correspondants de $C$ et de
$\mathbf{V}(\mathcal{O}_X(1))$ ne sont pas isomorphes. Cette question va être
étudiée ci-dessous.

\subsection{Éclatement de cônes projetants}
\label{subsection:II.8.7-fr}

\begin{env}[8.7.1]
\label{II.8.7.1-fr}
Sous les conditions de (\hyperref[II.8.6.1-fr]{8.6.1}), on a, en posant
$r=\operatorname{Proj}(\widehat{\alpha})$, un diagramme commutatif
\begin{equation}
\label{II.8.7.1.1-fr}
  \xymatrix{
    X\ar[r]^{i_X\circ\varepsilon_X}\ar[d]_{q} &
    \widehat{C}_X\ar[d]^{r}\\
    Y\ar[r]_{i\circ\varepsilon} & \widehat{C}
  }
\tag{8.7.1.1}
\end{equation}

\oldpage[II]{174}
en vertu de (\hyperref[II.8.5.1.3-fr]{8.5.1.3}) et
(\hyperref[II.8.5.3.2-fr]{8.5.3.2})~; en outre, la restriction de $r$ au
complémentaire $\widehat{C}_X-i_X(\varepsilon_X(X))$ de la section nulle est
un \emph{isomorphisme} sur le complémentaire
$\widehat{C}-i(\varepsilon(Y))$ de la section nulle en vertu de
(\hyperref[II.8.6.2-fr]{8.6.2}). Si on suppose pour simplifier $Y$ affine,
$\mathcal{S}$ de type fini et engendré par $\mathcal{S}_1$, $X$ est projectif
sur $Y$ et $\widehat{C}_X$ projectif sur $X$
(\hyperref[II.5.5.1-fr]{5.5.1}), donc $\widehat{C}_X$ est projectif sur $Y$
(\hyperref[II.5.5.5-fr]{5.5.5, (ii)}), et \emph{a fortiori} sur
$\widehat{C}$ (\hyperref[II.5.5.5-fr]{5.5.5, (v)}). On a ainsi un
$Y$-morphisme projectif $r:\widehat{C}_X\to\widehat{C}$ (dont la restriction
à $C_X$ est un $Y$-morphisme projectif $C_X\to C$) qui \emph{contracte $X$
dans $Y$} et induit un \emph{isomorphisme} quand on le restreint aux
\emph{complémentaires de $X$ et de $Y$}. On a donc une relation entre $C_X$
et $C$, analogue à celle qui a lieu entre un préschéma éclaté et son préschéma
de départ (\hyperref[II.8.1.3-fr]{8.1.3}). Nous allons effectivement montrer
qu'on peut identifier $C_X$ au spectre homogène d'une
$\mathcal{O}_C$-Algèbre graduée.
\end{env}

\begin{env}[8.7.2]
\label{II.8.7.2-fr}
Gardant les notations de (\hyperref[II.8.6.1-fr]{8.6.1}), considérons pour
chaque $n\geqslant 0$, l'Idéal quasi-cohérent
\begin{equation}
\label{II.8.7.2.1-fr}
  \mathcal{S}_{[n]}=\bigoplus_{m\geqslant n}\mathcal{S}_m
\tag{8.7.2.1}
\end{equation}
de la $\mathcal{O}_Y$-Algèbre graduée $\mathcal{S}$. Il est clair que l'on a
\begin{equation}
\label{II.8.7.2.2-fr}
  \mathcal{S}_{[0]}=\mathcal{S},\qquad
  \mathcal{S}_{[n]}\subset\mathcal{S}_{[m]}
  \qquad\text{pour }m\leqslant n
\tag{8.7.2.2}
\end{equation}
\begin{equation}
\label{II.8.7.2.3-fr}
  \mathcal{S}_{[n]}\mathcal{S}_{[m]}\subset\mathcal{S}_{[m+n]}.
\tag{8.7.2.3}
\end{equation}

Considérons le $\mathcal{O}_C$-Module associé à $\mathcal{S}_{[n]}$, qui est un
Idéal quasi-cohérent de $\mathcal{O}_C=\widetilde{\mathcal{S}}$
(\hyperref[II.1.4.4-fr]{1.4.4})
\begin{equation}
\label{II.8.7.2.4-fr}
  \mathcal{I}_n=(\mathcal{S}_{[n]})^{\sim}.
\tag{8.7.2.4}
\end{equation}

On déduit alors de (\hyperref[II.8.7.2.2-fr]{8.7.2.2}) et
(\hyperref[II.8.7.2.3-fr]{8.7.2.3}), utilisant
(\hyperref[II.1.4.4-fr]{1.4.4}) et (\hyperref[II.1.4.8.1-fr]{1.4.8.1}),
les formules analogues
\begin{equation}
\label{II.8.7.2.5-fr}
  \mathcal{I}_0=\mathcal{O}_C,\qquad
  \mathcal{I}_n\subset\mathcal{I}_m
  \qquad\text{pour }m\leqslant n
\tag{8.7.2.5}
\end{equation}
\begin{equation}
\label{II.8.7.2.6-fr}
  \mathcal{I}_n\mathcal{I}_m\subset\mathcal{I}_{n+m}.
\tag{8.7.2.6}
\end{equation}

On est donc dans les conditions de (\hyperref[II.8.1.1-fr]{8.1.1}), ce qui
conduit à introduire la $\mathcal{O}_C$-Algèbre graduée quasi-cohérente
\begin{equation}
\label{II.8.7.2.7-fr}
  \mathcal{S}^{\natural}
  =\bigoplus_{n\geqslant 0}\mathcal{I}_n
  =\left(\bigoplus_{n\geqslant 0}\mathcal{S}_{[n]}\right)^{\sim}.
\tag{8.7.2.7}
\end{equation}
\end{env}

\begin{proposition}[8.7.3]
\label{II.8.7.3-fr}
Il existe un $C$-isomorphisme canonique
\begin{equation}
\label{II.8.7.3.1-fr}
  h:C_X\xrightarrow{\sim}\operatorname{Proj}(\mathcal{S}^{\natural}).
\tag{8.7.3.1}
\end{equation}
\end{proposition}

\begin{proof}
Supposons d'abord $Y=\operatorname{Spec}(A)$ affine, d'où
$\mathcal{S}=\widetilde{S}$, où $S$ est une $A$-algèbre graduée à degrés
positifs et $C=\operatorname{Spec}(S)$. La définition
(\hyperref[II.8.2.7.4-fr]{8.2.7.4}) montre que l'on a alors, avec les
notations de (\hyperref[II.8.2.6-fr]{8.2.6}),
$\mathcal{S}^{\natural}=(S^{\natural})^{\sim}$. Pour définir
(\hyperref[II.8.7.3.1-fr]{8.7.3.1}), considérons un élément homogène
$f\in S_d$ ($d>0$) et l'élément correspondant $f^{\natural}\in S^{\natural}$
(\hyperref[II.8.2.6-fr]{8.2.6})~; le $S$-isomorphisme
(\hyperref[II.8.2.7.3-fr]{8.2.7.3}) définit donc un $C$-isomorphisme
\begin{equation}
\label{II.8.7.3.2-fr}
  \operatorname{Spec}(S_f^{\geqslant})
  \xrightarrow{\sim}
  \operatorname{Spec}(S_{(f^{\natural})}^{\natural}).
\tag{8.7.3.2}
\end{equation}

\oldpage[II]{175}
Mais avec les notations de (\hyperref[II.8.6.2-fr]{8.6.2}), si
$v:C_X\to X$ est le morphisme structural, il résulte de
(\hyperref[II.8.6.2.1-fr]{8.6.2.1}) que l'on a
$v^{-1}(X_f)=\operatorname{Spec}(S_f^{\geqslant})$. On a d'autre part
$\operatorname{Spec}(S_{(f^{\natural})}^{\natural})=D_+(f^{\natural})$, si
bien que (\hyperref[II.8.7.3.2-fr]{8.7.3.2}) définit un isomorphisme
$v^{-1}(X_f)\to D_+(f^{\natural})$. En outre, si $g\in S_e$ ($e>0$), le
diagramme
\[
  \xymatrix{
    v^{-1}(X_{fg})\ar[r]^{\sim}\ar[d] &
    D_+(f^{\natural}g^{\natural})\ar[d]\\
    v^{-1}(X_f)\ar[r]^{\sim} & D_+(f^{\natural})
  }
\]
est commutatif, comme il résulte aussitôt de la définition de l'isomorphisme
(\hyperref[II.8.2.7.3-fr]{8.2.7.3}). Enfin par définition $S_+$ est engendré
par les $f$ homogènes, donc il résulte de
(\hyperref[II.8.2.10-fr]{8.2.10, (iv)}) et de
(\hyperref[II.2.3.14-fr]{2.3.14}) que les $D_+(f^{\natural})$ forment un
recouvrement de $\operatorname{Proj}(S^{\natural})$ et les $v^{-1}(X_f)$ un
recouvrement de $C_X$, puisque les $X_f$ forment un recouvrement de $X$~;
cela achève donc de définir dans ce cas l'isomorphisme
(\hyperref[II.8.7.3.1-fr]{8.7.3.1}).

Pour démontrer (\hyperref[II.8.7.3-fr]{8.7.3}) dans le cas général, il suffit
de voir que si $U$, $U'$ sont deux ouverts affines de $Y$ tels que
$U'\subset U$, d'anneaux $A$ et $A'$, et si on pose
$\mathcal{S}|U=\widetilde{S}$, $\mathcal{S}|U'=\widetilde{S'}$, le diagramme
\begin{equation}
\label{II.8.7.3.3-fr}
  \xymatrix{
    C_{U'}\ar[r]\ar[d] & \operatorname{Proj}(S'^{\natural})\ar[d]\\
    C_U\ar[r] & \operatorname{Proj}(S^{\natural})
  }
\tag{8.7.3.3}
\end{equation}
est commutatif. Mais $S'$ s'identifie canoniquement à $S\otimes_A A'$, donc
$S'^{\natural}$ à
\[
  S^{\natural}\otimes_S S'=S^{\natural}\otimes_A A'~;
\]
on a donc
$\operatorname{Proj}(S'^{\natural})=\operatorname{Proj}(S^{\natural})
\times_U U'$ (\hyperref[II.2.8.10-fr]{2.8.10})~; de même, si
$X=\operatorname{Proj}(S)$, $X'=\operatorname{Proj}(S')$, on a
$X'=X\times_U U'$ et
$\mathcal{S}_{X'}=\mathcal{S}_X\otimes_{\mathcal{O}_U}\mathcal{O}_{U'}$
(\hyperref[II.3.5.4-fr]{3.5.4}), autrement dit
$\mathcal{S}_{X'}=j^*(\mathcal{S}_X)$, où $j$ est la projection $X'\to X$.
On a par suite (\hyperref[II.1.5.2-fr]{1.5.2})
$C_{U'}=C_U\times_X X'=C_U\times_U U'$, et la commutativité de
(\hyperref[II.8.7.3.3-fr]{8.7.3.3}) est alors immédiate.
\end{proof}

\begin{remark}[8.7.4]
\label{II.8.7.4-fr}
\begin{enumerate}
  \item[\rm{(i)}] La fin du raisonnement de
  (\hyperref[II.8.7.3-fr]{8.7.3}) se généralise aussitôt de la façon
  suivante. Soient $g:Y'\to Y$ un morphisme,
  $\mathcal{S}'=g^*(\mathcal{S})$, $X'=\operatorname{Proj}(\mathcal{S}')$~;
  on a alors un diagramme commutatif
  \begin{equation}
  \label{II.8.7.4.1-fr}
    \xymatrix{
      C_{X'}\ar[r]\ar[d] &
      \operatorname{Proj}(\mathcal{S}'^{\natural})\ar[d]\\
      C_X\ar[r] & \operatorname{Proj}(\mathcal{S}^{\natural})
    }
  \tag{8.7.4.1}
  \end{equation}

  D'autre part, soit $\varphi:\mathcal{S}''\to\mathcal{S}$ un homomorphisme
  de $\mathcal{O}_Y$-Algèbres graduées, tel que si on pose
  $X''=\operatorname{Proj}(\mathcal{S}'')$,
  $u=\operatorname{Proj}(\varphi):X\to X''$ soit partout défini~; on a
  d'autre part

  \oldpage[II]{176}
  un $Y$-morphisme $v:C\to C''$ (avec
  $C''=\operatorname{Spec}(\mathcal{S}'')$) tel que
  $\mathcal{A}(v)=\varphi$, et comme $\varphi$ est un homomorphisme
  d'Algèbres graduées, on déduit de $\varphi$ un $v$-morphisme d'Algèbres
  graduées
  $\psi:\mathcal{S}''^{\natural}\to\mathcal{S}^{\natural}$
  (\hyperref[II.1.4.1-fr]{1.4.1}). En outre, il résulte aisément de
  (\hyperref[II.8.2.10-fr]{8.2.10, (iv)}) et de l'hypothèse sur $\varphi$,
  que $\operatorname{Proj}(\psi)$ est partout défini. Enfin, compte tenu de
  (\hyperref[II.3.5.6.1-fr]{3.5.6.1}), on a un $u$-morphisme canonique
  $\mathcal{S}_{X''}\to\mathcal{S}_X$, d'où
  (\hyperref[II.1.5.6-fr]{1.5.6}) un morphisme $w:C_{X''}\to C_X$. Cela
  étant, le diagramme
  \begin{equation}
  \label{II.8.7.4.2-fr}
    \xymatrix{
      C_{X''}\ar[r]^{\sim}\ar[d]_{w} &
      \operatorname{Proj}(\mathcal{S}''^{\natural})
        \ar[d]^{\operatorname{Proj}(\psi)}\\
      C_X\ar[r]^{\sim} & \operatorname{Proj}(\mathcal{S}^{\natural})
    }
  \tag{8.7.4.2}
  \end{equation}
  est commutatif, comme on le vérifie aussitôt en se ramenant au cas où $Y$
  est affine.

  \item[\rm{(ii)}] Notons qu'en vertu de
  (\hyperref[II.8.7.2.5-fr]{8.7.2.5}) et
  (\hyperref[II.8.7.2.6-fr]{8.7.2.6}), on a
  $\mathcal{I}_1^m\subset\mathcal{I}_m\subset\mathcal{I}_1$ pour tout $m>0$.
  Or, par définition, $\mathcal{I}_1=(\mathcal{S}_+)^{\sim}$, donc
  $\mathcal{I}_1$ définit dans $C$ le sous-préschéma fermé
  $\varepsilon(Y)$ ((\hyperref[II.1.4.10-fr]{1.4.10}) et
  (\hyperref[II.8.3.2-fr]{8.3.2}))~; on en conclut que pour tout $m>0$,
  \emph{le support de $\mathcal{O}_C/\mathcal{I}_m$ est contenu dans l'espace
  sous-jacent au préschéma des sommets $\varepsilon(Y)$}~; dans l'image
  réciproque du cône épointé affine $E$, le morphisme structural
  $\operatorname{Proj}(\mathcal{S}^{\natural})\to C$ se réduit donc à un
  \emph{isomorphisme} (comme il résulte d'ailleurs de
  (\hyperref[II.8.7.3-fr]{8.7.3}) et
  (\hyperref[II.8.7.1-fr]{8.7.1})). En outre, identifiant canoniquement $C$ à
  un ouvert de $\widehat{C}$ (\hyperref[II.8.3.3-fr]{8.3.3}), on peut
  évidemment prolonger les Idéaux $\mathcal{I}_m$ de $\mathcal{O}_C$ en des
  Idéaux $\mathcal{J}_m$ de $\mathcal{O}_{\widehat{C}}$, par la condition de
  coïncider avec $\mathcal{O}_{\widehat{C}}$ dans l'ouvert $\widehat{E}$ de
  $\widehat{C}$. Si on pose
  $\mathcal{T}=\bigoplus_{n\geqslant 0}\mathcal{J}_m$, qui est une
  $\mathcal{O}_{\widehat{C}}$-Algèbre graduée quasi-cohérente, on peut
  prolonger l'isomorphisme (\hyperref[II.8.7.3.1-fr]{8.7.3.1}) en un
  $\widehat{C}$-isomorphisme
  \begin{equation}
  \label{II.8.7.4.3-fr}
    \widehat{C}_X\xrightarrow{\sim}\operatorname{Proj}(\mathcal{T}).
  \tag{8.7.4.3}
  \end{equation}

  En effet, au-dessus de $\widehat{E}$, il résulte de ce qui précède que
  $\operatorname{Proj}(\mathcal{T})$ s'identifie canoniquement à
  $\widehat{E}$, et on définit donc l'isomorphisme
  (\hyperref[II.8.7.4.3-fr]{8.7.4.3}) au-dessus de $\widehat{E}$ en lui
  imposant de coïncider avec l'isomorphisme canonique
  $\widehat{E}_X\to\widehat{E}$ (\hyperref[II.8.6.2-fr]{8.6.2})~; il est clair
  qu'alors cet isomorphisme et (\hyperref[II.8.7.3.1-fr]{8.7.3.1}) coïncident
  au-dessus de $E$.
\end{enumerate}
\end{remark}

\begin{corollary}[8.7.5]
\label{II.8.7.5-fr}
Supposons qu'il existe $n_0>0$ tel que
\begin{equation}
\label{II.8.7.5.1-fr}
  \mathcal{S}_{n+1}=\mathcal{S}_1\mathcal{S}_n
  \qquad\text{pour }n\geqslant n_0.
\tag{8.7.5.1}
\end{equation}
Alors le sous-préschéma des sommets de $C_X$ (isomorphe à $X$) est l'image
réciproque par le morphisme canonique $r:C_X\to C$ du sous-préschéma des
sommets de $C$ (isomorphe à $Y$). Inversement, si cette propriété a lieu et
si on suppose de plus $Y$ noethérien et $\mathcal{S}$ de type fini, il existe
$n_0>0$ tel que l'on ait (\hyperref[II.8.7.5.1-fr]{8.7.5.1}).
\end{corollary}

\begin{proof}
La première assertion étant locale sur $Y$, on peut pour l'établir supposer
$Y=\operatorname{Spec}(A)$ affine, et $\mathcal{S}=\widetilde{S}$, où $S$ est
une $A$-algèbre graduée à degrés positifs. Elle résulte alors de
(\hyperref[II.8.2.12-fr]{8.2.12}), car on a
$\operatorname{Proj}(S^{\natural}\otimes_S S_0)=
C_X\times_C\varepsilon(Y)$ (en vertu de l'identification
(\hyperref[II.8.7.3.1-fr]{8.7.3.1})), c'est-à-dire que ce préschéma est
l'image réciproque de $\varepsilon(Y)$ dans $C_X$
(\hyperref[I.4.4.1-fr]{I, 4.4.1}). La réciproque résulte aussi de
(\hyperref[II.8.2.12-fr]{8.2.12}) lorsque $Y$ est affine noethérien et $S$ de
type fini.

\oldpage[II]{177}
Si $Y$ est noethérien (non nécessairement affine) et $\mathcal{S}$ de type
fini, il y a un recouvrement fini de $Y$ par des ouverts affines noethériens
$U_i$, et on déduit alors de ce qui précède que pour tout $i$, il y a un
entier $n_i$ tel que
$\mathcal{S}_{n+1}|U_i=(\mathcal{S}_1|U_i)(\mathcal{S}_n|U_i)$ dès que
$n\geqslant n_i$~; le plus grand des $n_i$ vérifie donc
(\hyperref[II.8.7.5.1-fr]{8.7.5.1}).
\end{proof}

\begin{env}[8.7.6]
\label{II.8.7.6-fr}
Considérons maintenant le $C$-préschéma $Z$ obtenu en faisant \emph{éclater}
dans le cône affine $C$ le \emph{sous-préschéma des sommets}
$\varepsilon(Y)$~; par définition (\hyperref[II.8.1.3-fr]{8.1.3}) c'est donc
le préschéma
$\operatorname{Proj}(\bigoplus_{n\geqslant 0}\mathcal{S}_+^n)$~; l'injection
canonique
\begin{equation}
\label{II.8.7.6.1-fr}
  \iota:\bigoplus_{n\geqslant 0}\mathcal{S}_+^n
  \longrightarrow\mathcal{S}^{\natural}
\tag{8.7.6.1}
\end{equation}
définit (grâce à l'identification (\hyperref[II.8.7.3-fr]{8.7.3})) un
$C$-morphisme dominant canonique
\begin{equation}
\label{II.8.7.6.2-fr}
  G(\iota)\longrightarrow Z
\tag{8.7.6.2}
\end{equation}
où $G(\iota)$ est un ouvert de $C_X$ (\hyperref[II.3.5.1-fr]{3.5.1})~; on
notera qu'on peut avoir $G(\iota)\ne C_X$, comme le montre l'exemple où
$Y=\operatorname{Spec}(K)$, $K$ étant un corps, $\mathcal{S}=\widetilde{S}$
avec $S=K[\mathbf{y}]$, où $\mathbf{y}$ est une indéterminée \emph{de degré
2}~; si $R_n$ désigne l'ensemble $(S_+)^n$, considéré comme partie de
$S_{[n]}=S_n^{\natural}$, $S_+^{\natural}$ n'est pas la racine dans
$S^{\natural}$ de l'idéal engendré par la réunion des $R_n$ (cf.
(\hyperref[II.2.3.14-fr]{2.3.14})).
\end{env}

\begin{corollary}[8.7.7]
\label{II.8.7.7-fr}
Supposons qu'il existe $n_0>0$ tel que
\begin{equation}
\label{II.8.7.7.1-fr}
  \mathcal{S}_n=\mathcal{S}_1^n
  \qquad\text{pour }n\geqslant n_0.
\tag{8.7.7.1}
\end{equation}
Alors le morphisme canonique (\hyperref[II.8.7.6.2-fr]{8.7.6.2}) est partout
défini et est un isomorphisme $C_X\xrightarrow{\sim}Z$. Inversement, si cette
propriété a lieu et si on suppose de plus $Y$ noethérien et $\mathcal{S}$ de
type fini, il existe $n_0$ tel que l'on ait
(\hyperref[II.8.7.7.1-fr]{8.7.7.1}).
\end{corollary}

\begin{proof}
En effet, la première assertion est locale sur $Y$, et découle donc de
(\hyperref[II.8.2.14-fr]{8.2.14})~; il en est de même de la réciproque, en
raisonnant comme dans (\hyperref[II.8.7.5-fr]{8.7.5}).
\end{proof}

\begin{remark}[8.7.8]
\label{II.8.7.8-fr}
Comme la condition (\hyperref[II.8.7.7.1-fr]{8.7.7.1}) implique
(\hyperref[II.8.7.5.1-fr]{8.7.5.1}), on voit que lorsqu'elle est vérifiée,
non seulement $C_X$ s'identifie au préschéma obtenu en faisant éclater le
sommet (identifié à $Y$) du cône affine $C$, mais encore le sommet (identifié
à $X$) de $C_X$ s'identifie au sous-préschéma fermé image réciproque du
sommet $Y$ de $C$. En outre, l'hypothèse
(\hyperref[II.8.7.7.1-fr]{8.7.7.1}) implique que sur
$X=\operatorname{Proj}(\mathcal{S})$, les $\mathcal{O}_X$-Modules
$\mathcal{O}_X(n)$ sont inversibles (\hyperref[II.3.2.5-fr]{3.2.5} et
\hyperref[II.3.2.9-fr]{3.2.9}) et que l'on a
$\mathcal{O}_X(n)=\mathcal{L}^{\otimes n}$ avec
$\mathcal{L}=\mathcal{O}_X(1)$ (\hyperref[II.3.2.7-fr]{3.2.7} et
\hyperref[II.3.2.9-fr]{3.2.9})~; par définition
(\hyperref[II.8.6.1.1-fr]{8.6.1.1}), $C_X$ est donc le \emph{fibré vectoriel}
$\mathbf{V}(\mathcal{L})$ sur $X$, et son sommet la \emph{section nulle} de ce
fibré vectoriel.
\end{remark}

\subsection{Faisceaux amples et contractions}
\label{subsection:II.8.8-fr}

\begin{env}[8.8.1]
\label{II.8.8.1-fr}
Soient $Y$ un préschéma, $f:X\to Y$ un morphisme \emph{séparé et
quasi-compact}, $\mathcal{L}$ un $\mathcal{O}_X$-Module inversible \emph{ample
relativement à $f$}. Considérons la $\mathcal{O}_Y$-Algèbre graduée à degrés
positifs
\begin{equation}
\label{II.8.8.1.1-fr}
  \mathcal{S}=\mathcal{O}_Y\oplus
  \bigoplus_{n\geqslant 1}f_*(\mathcal{L}^{\otimes n})
\tag{8.8.1.1}
\end{equation}

\oldpage[II]{178}
qui est quasi-cohérente (\hyperref[I.9.2.2-fr]{I, 9.2.2, a)}). On a un
homomorphisme canonique de $\mathcal{O}_X$-Algèbres graduées
\begin{equation}
\label{II.8.8.1.2-fr}
  \tau:f^*(\mathcal{S})\longrightarrow
  \bigoplus_{n\geqslant 0}\mathcal{L}^{\otimes n}
\tag{8.8.1.2}
\end{equation}
qui, pour les degrés $\geqslant 1$, coïncide avec l'homomorphisme canonique
$\sigma:f^*(f_*(\mathcal{L}^{\otimes n}))\to\mathcal{L}^{\otimes n}$
(\hyperref[0.4.4.3-fr]{0, 4.4.3}), et pour le degré $0$ est l'identité.
L'hypothèse que $\mathcal{L}$ est $f$-ample entraîne alors
(\hyperref[II.4.6.3-fr]{4.6.3} et \hyperref[II.3.6.1-fr]{3.6.1}) que le
$Y$-morphisme correspondant
\begin{equation}
\label{II.8.8.1.3-fr}
  r=r_{\mathcal{L},\tau}:X\longrightarrow
  P=\operatorname{Proj}(\mathcal{S})
\tag{8.8.1.3}
\end{equation}
est partout défini et est une \emph{immersion ouverte dominante}, et que l'on
a
\begin{equation}
\label{II.8.8.1.4-fr}
  r^*(\mathcal{O}_P(n))=\mathcal{L}^{\otimes n}
  \qquad\text{pour tout }n\in\mathbf{Z}.
\tag{8.8.1.4}
\end{equation}
\end{env}

\begin{proposition}[8.8.2]
\label{II.8.8.2-fr}
Soit $C=\operatorname{Spec}(\mathcal{S})$ \emph{le cône affine défini par
$\mathcal{S}$}~; si $\mathcal{L}$ est $f$-ample, il existe un $Y$-morphisme
canonique
\begin{equation}
\label{II.8.8.2.1-fr}
  g:V=\mathbf{V}(\mathcal{L})\longrightarrow C
\tag{8.8.2.1}
\end{equation}
tel que le diagramme
\begin{equation}
\label{II.8.8.2.2-fr}
  \xymatrix{
    X\ar[r]^{j}\ar[d]_{f} &
    \mathbf{V}(\mathcal{L})\ar[r]^{\pi}\ar[d]^{g} & X\ar[d]^{f}\\
    Y\ar[r]_{\varepsilon} & C\ar[r]_{\psi} & Y
  }
\tag{8.8.2.2}
\end{equation}
soit commutatif, $\psi$ et $\pi$ étant les morphismes structuraux, $j$ et
$\varepsilon$ les immersions canoniques appliquant respectivement $X$ et $Y$
sur la section nulle de $\mathbf{V}(\mathcal{L})$ et le préschéma des sommets
de $C$. En outre, la restriction de $g$ à
$\mathbf{V}(\mathcal{L})-j(X)$ est une immersion ouverte
\begin{equation}
\label{II.8.8.2.3-fr}
  \mathbf{V}(\mathcal{L})-j(X)\longrightarrow
  E=C-\varepsilon(Y)
\tag{8.8.2.3}
\end{equation}
dans le cône épointé affine $E$ correspondant à $\mathcal{S}$.
\end{proposition}

\begin{proof}
Avec les notations de (\hyperref[II.8.8.1-fr]{8.8.1}), soient
$\mathcal{S}_P^{\geqslant}=\bigoplus_{n\geqslant 0}\mathcal{O}_P(n)$, et
$C_P=\operatorname{Spec}(\mathcal{S}_P^{\geqslant})$. On sait
(\hyperref[II.8.6.2-fr]{8.6.2}) qu'on a un morphisme canonique
$h=\operatorname{Spec}(\alpha):C_P\to C$ tel que le diagramme
\begin{equation}
\label{II.8.8.2.4-fr}
  \xymatrix{
    C_P\ar[r]\ar[d]_{h} & P\ar[d]^{p}\\
    C\ar[r]_{\psi} & Y
  }
\tag{8.8.2.4}
\end{equation}
soit commutatif~; en outre, si $\varepsilon_P:P\to C_P$ est l'immersion
canonique, le diagramme
\begin{equation}
\label{II.8.8.2.5-fr}
  \xymatrix{
    P\ar[r]^{\varepsilon_P}\ar[d]_{p} & C_P\ar[d]^{h}\\
    Y\ar[r]_{\varepsilon} & C
  }
\tag{8.8.2.5}
\end{equation}
est commutatif (\hyperref[II.8.7.1.1-fr]{8.7.1.1}), et enfin, la restriction
de $h$ au cône épointé affine $E_P$ est un \emph{isomorphisme}
$E_P\xrightarrow{\sim}E$ (\hyperref[II.8.6.2-fr]{8.6.2}). D'autre part, il
résulte de (\hyperref[II.8.8.1.4-fr]{8.8.1.4}) que l'on a
\[
  r^*(\mathcal{S}_P^{\geqslant})=
  \mathbf{S}_{\mathcal{O}_X}(\mathcal{L})
\]

\oldpage[II]{179}
et par suite on a un $P$-morphisme canonique
$q:\mathbf{V}(\mathcal{L})\to C_P$, le diagramme commutatif
\begin{equation}
\label{II.8.8.2.6-fr}
  \xymatrix{
    \mathbf{V}(\mathcal{L})\ar[r]^{\pi}\ar[d]_{q} & X\ar[d]^{r}\\
    C_P\ar[r] & P
  }
\tag{8.8.2.6}
\end{equation}
identifiant $\mathbf{V}(\mathcal{L})$ au produit $C_P\times_P X$
(\hyperref[II.1.5.2-fr]{1.5.2})~; comme $r$ est une immersion ouverte, il en
est donc de même de $q$ (\hyperref[I.4.3.2-fr]{I, 4.3.2}). En outre, la
restriction de $q$ à $\mathbf{V}(\mathcal{L})-j(X)$ applique ce préschéma dans
$E_P$ par (\hyperref[II.8.5.2-fr]{8.5.2}), et le diagramme
\begin{equation}
\label{II.8.8.2.7-fr}
  \xymatrix{
    X\ar[r]^{j}\ar[d]_{r} & \mathbf{V}(\mathcal{L})\ar[d]^{q}\\
    P\ar[r]_{\varepsilon_P} & C_P
  }
\tag{8.8.2.7}
\end{equation}
est commutatif (cas particulier de
(\hyperref[II.8.5.1.3-fr]{8.5.1.3})). Les assertions de
(\hyperref[II.8.8.2-fr]{8.8.2}) résultent aussitôt de ces faits, en prenant
pour $g$ le morphisme composé $h\circ q$.
\end{proof}

\begin{remark}[8.8.3]
\label{II.8.8.3-fr}
Supposons en outre que $Y$ soit un préschéma \emph{noethérien} et $f$ un
morphisme \emph{propre}. Comme $r$ est alors \emph{propre}
(\hyperref[II.5.4.4-fr]{5.4.4}), donc fermé, et que c'est une immersion ouverte
dominante, $r$ est nécessairement un \emph{isomorphisme}
$X\xrightarrow{\sim}P$. On verra en outre au chapitre III
(\hyperref[III.2.3.5.1-fr]{III, 2.3.5.1}), que $\mathcal{S}$ est alors
nécessairement une $\mathcal{O}_Y$-Algèbre \emph{de type fini}. Alors il en
résulte que $\mathcal{S}^{\natural}$ est une
$\mathcal{S}_0^{\natural}$-Algèbre \emph{de type fini}
(\hyperref[II.8.2.10-fr]{8.2.10, (i)} et
\hyperref[II.8.7.2.7-fr]{8.7.2.7})~; comme $C_P$ est $C$-isomorphe à
$\operatorname{Proj}(\mathcal{S}^{\natural})$
(\hyperref[II.8.7.3-fr]{8.7.3}), on voit que le morphisme $h:C_P\to C$ est
\emph{projectif}~; comme le morphisme $r$ est un isomorphisme, il en est de
même de $q:\mathbf{V}(\mathcal{L})\to C_P$, et on en conclut que le morphisme
$g:\mathbf{V}(\mathcal{L})\to C$ est \emph{projectif}. En outre, comme la
restriction de $h$ à $E_P$ est un isomorphisme sur $E$ et que $q$ est un
isomorphisme, la restriction (\hyperref[II.8.8.2.3-fr]{8.8.2.3}) de $g$ est un
\emph{isomorphisme}
$\mathbf{V}(\mathcal{L})-j(X)\xrightarrow{\sim}E$.

Si on suppose de plus que $\mathcal{L}$ est \emph{très ample} pour $f$, on
verra aussi au chapitre III (\hyperref[III.2.3.5.1-fr]{III, 2.3.5.1}), qu'il
existe un entier $n_0>0$ tel que
$\mathcal{S}_n=\mathcal{S}_1^n$ pour $n\geqslant n_0$. On en conclura par
(\hyperref[II.8.7.7-fr]{8.7.7}) que $\mathbf{V}(\mathcal{L})$ s'identifie au
préschéma $Z$ obtenu en \emph{faisant éclater dans le cône affine $C$ le
sous-préschéma des sommets} (identifié à $Y$), et que la \emph{section nulle}
de $\mathbf{V}(\mathcal{L})$ (identifiée à $Y$) est \emph{l'image réciproque}
du sous-préschéma des sommets $Y$ de $C$. Une partie des résultats précédents
peut en fait s'établir sans hypothèse noethérienne~:
\end{remark}

\begin{corollary}[8.8.4]
\label{II.8.8.4-fr}
Soient $Y$ un préschéma (resp. un schéma quasi-compact), $f:X\to Y$ un
morphisme propre, $\mathcal{L}$ un $\mathcal{O}_X$-Module inversible ample pour
$f$. Alors le morphisme (\hyperref[II.8.8.2.1-fr]{8.8.2.1}) est propre (resp.
projectif) et sa restriction (\hyperref[II.8.8.2.3-fr]{8.8.2.3}) est un
isomorphisme.
\end{corollary}

\begin{proof}
Pour prouver que $g$ est propre, on peut se ramener au cas où $Y$ est affine,
et il suffit donc de considérer le cas où $Y$ est un schéma quasi-compact. Le
même raisonnement que dans (\hyperref[II.8.8.3-fr]{8.8.3}) montre d'abord que
$r$ est un \emph{isomorphisme} $X\xrightarrow{\sim}P$~; par suite $q$ est aussi
un isomorphisme, et comme la restriction de $h$ à $E_P$ est un isomorphisme
$E_P\xrightarrow{\sim}E$, on voit déjà que
(\hyperref[II.8.8.2.3-fr]{8.8.2.3}) est un isomorphisme. Reste à prouver que
$g$ est \emph{projectif}.

Comme $f$ est de type fini par hypothèse, on peut appliquer
(\hyperref[II.3.8.5-fr]{3.8.5}) à l'homomorphisme

\oldpage[II]{180}
$\tau$ de (\hyperref[II.8.8.1.2-fr]{8.8.1.2})~: il y a par suite un entier
$d>0$, un sous-$\mathcal{O}_Y$-Module quasi-cohérent de type fini
$\mathcal{E}$ de $\mathcal{S}_d$ tel que si $\mathcal{S}'$ est la
sous-$\mathcal{O}_Y$-Algèbre de $\mathcal{S}$ engendrée par $\mathcal{E}$ et
$\tau'=\tau\circ f^*(\varphi)$ ($\varphi$ injection canonique
$\mathcal{S}'\to\mathcal{S}$), $r'=r_{\mathcal{L},\tau'}$ soit une immersion
\[
  X\longrightarrow P'=\operatorname{Proj}(\mathcal{S}').
\]
En outre, comme $\varphi$ est injectif, $r'$ est aussi une \emph{immersion
dominante} (\hyperref[II.3.7.6-fr]{3.7.6})~; le même raisonnement que pour $r$
montre donc que $r'$ est une \emph{immersion fermée surjective}~; comme $r'$ se
factorise en
$X\xrightarrow{r}\operatorname{Proj}(\mathcal{S})
\xrightarrow{\Phi}\operatorname{Proj}(\mathcal{S}')$, où
$\Phi=\operatorname{Proj}(\varphi)$, on en conclut que $\Phi$ est aussi une
\emph{immersion fermée surjective}. Mais cela entraîne que $\Phi$ est un
\emph{isomorphisme}~; en effet, on peut se ramener au cas où
$Y=\operatorname{Spec}(A)$ est affine, $\mathcal{S}=\widetilde{S}$,
$\mathcal{S}'=\widetilde{S'}$, où $S$ est une $A$-algèbre graduée, $S'$ une
sous-algèbre graduée de $S$. Pour tout élément homogène $t\in S'$,
$S'_{(t)}$ est alors un sous-anneau de $S_{(t)}$~; si on revient à la
définition de $\operatorname{Proj}(\varphi)$
(\hyperref[II.2.8.1-fr]{2.8.1}), on voit qu'on est ramené à prouver que si
$B'$ est un sous-anneau d'un anneau $B$, et si le morphisme
$\operatorname{Spec}(B)\to\operatorname{Spec}(B')$ correspondant à l'injection
canonique $B'\to B$ est une immersion fermée, ce morphisme est nécessairement
un \emph{isomorphisme}~; mais cela résulte de
(\hyperref[I.4.2.3-fr]{I, 4.2.3}).

On a en outre
$\Phi^*(\mathcal{O}_{P'}(n))=\mathcal{O}_P(n)$
(\hyperref[II.3.5.2-fr]{3.5.2, (ii)} et
\hyperref[II.3.5.4-fr]{3.5.4}), donc
$r'^*(\mathcal{O}_{P'}(n))$ est isomorphe à $\mathcal{L}^{\otimes n}$
(\hyperref[II.4.6.3-fr]{4.6.3}). Posons
$\mathcal{S}''=\mathcal{S}'^{(d)}$, de sorte
(\hyperref[II.3.1.8-fr]{3.1.8, (i)}) que $X$ s'identifie canoniquement à
$P''=\operatorname{Proj}(\mathcal{S}'')$ et
$\mathcal{L}''=\mathcal{L}^{\otimes d}$ à $\mathcal{O}_{P''}(1)$
(\hyperref[II.3.2.9-fr]{3.2.9, (ii)}).

Cela étant, si $C''=\operatorname{Spec}(\mathcal{S}'')$,
$\mathcal{S}_{P''}^{\geqslant}=
\bigoplus_{n\geqslant 0}\mathcal{O}_{P''}(n)$ s'identifie à
$\bigoplus_{n\geqslant 0}\mathcal{L}''^{\otimes n}$, donc
$C_{P''}=\operatorname{Spec}(\mathcal{S}_{P''}^{\geqslant})$ à
$\mathbf{V}(\mathcal{L}'')$~; on sait d'autre part
(\hyperref[II.8.7.3-fr]{8.7.3}) que $C_{P''}$ est $C''$-isomorphe à
$\operatorname{Proj}(\mathcal{S}''^{\natural})$~; par définition de
$\mathcal{S}''$, $\mathcal{S}''^{\natural}$ est engendré par
$\mathcal{S}_1''^{\natural}$ et $\mathcal{S}_1''^{\natural}$ est de type fini
sur $\mathcal{S}_0''^{\natural}=\mathcal{S}''$
(\hyperref[II.8.2.10-fr]{8.2.10, (i) et (iii)}), donc
$\operatorname{Proj}(\mathcal{S}''^{\natural})$ est \emph{projectif} sur
$C''$ (\hyperref[II.5.5.1-fr]{5.5.1}). Considérons alors le diagramme
\begin{equation}
\label{II.8.8.4.1-fr}
  \xymatrix{
    \mathbf{V}(\mathcal{L})\ar[r]^{g}\ar[d]_{u} &
    \operatorname{Spec}(\mathcal{S})=C\ar[d]^{v}\\
    \mathbf{V}(\mathcal{L}'')\ar[r]_{g''} &
    \operatorname{Spec}(\mathcal{S}'')=C''
  }
\tag{8.8.4.1}
\end{equation}
où $g$ et $g''$ correspondent par (\hyperref[II.1.5.6-fr]{1.5.6}) aux
$f$-morphismes canoniques
\[
  \mathcal{S}\longrightarrow
  \bigoplus_{n\geqslant 0}\mathcal{L}^{\otimes n}
  \quad\text{et}\quad
  \mathcal{S}''\longrightarrow
  \bigoplus_{n\geqslant 0}\mathcal{L}''^{\otimes n}
\]
(\hyperref[II.3.3.2.3-fr]{3.3.2.3}) (voir
(\hyperref[II.8.8.5-fr]{8.8.5}) ci-dessous), $v$ au morphisme d'inclusion
$\mathcal{S}''\to\mathcal{S}$, et $u$ au morphisme d'inclusion
$\bigoplus_{n\geqslant 0}\mathcal{L}^{\otimes nd}\to
\bigoplus_{n\geqslant 0}\mathcal{L}^{\otimes n}$~; il est immédiat
(\hyperref[II.3.3.2-fr]{3.3.2}) que ce diagramme est commutatif. Nous venons
de voir que $g''$ est un morphisme projectif~; d'autre part, $u$ est un
morphisme \emph{fini}. En effet, la question étant locale sur $X$, on peut
supposer que $X$ est affine d'anneau $A$ et que $\mathcal{L}=\mathcal{O}_X$~;
tout revient alors à remarquer que l'anneau $A[T]$ est un module de type fini
sur son sous-anneau $A[T^d]$ ($T$ indéterminée). Comme $Y$ est un schéma
quasi-compact, et que $C''$ est affine sur $Y$, $C''$ est aussi un schéma
quasi-compact,

\oldpage[II]{181}
et par suite $g''\circ u$ est un morphisme projectif
(\hyperref[II.5.5.5-fr]{5.5.5, (ii)})~; par la commutativité de
(\hyperref[II.8.8.4.1-fr]{8.8.4.1}), il en est de même de $v\circ g$, et comme
$v$ est affine, donc séparé, on en conclut finalement que $g$ est projectif
(\hyperref[II.5.5.5-fr]{5.5.5, (v)}).
\end{proof}

\begin{env}[8.8.5]
\label{II.8.8.5-fr}
Reprenons la situation de (\hyperref[II.8.8.1-fr]{8.8.1}). Nous allons voir
qu'on peut donner du morphisme $g:\mathbf{V}(\mathcal{L})\to C$ une autre
définition valable pour tout $\mathcal{O}_X$-Module inversible $\mathcal{L}$
(non nécessairement ample). Pour cela, considérons le $f$-morphisme
\begin{equation}
\label{II.8.8.5.1-fr}
  \tau^\flat:\mathcal{S}\longrightarrow
  \bigoplus_{n\geqslant 0}\mathcal{L}^{\otimes n}
\tag{8.8.5.1}
\end{equation}
correspondant au morphisme $\tau$ de
(\hyperref[II.8.8.1.2-fr]{8.8.1.2}). On en déduit
(\hyperref[II.1.5.6-fr]{1.5.6}) un morphisme $g':V\to C$ tel que, si
$\pi:V\to X$ et $\psi:C\to Y$ sont les morphismes structuraux, les diagrammes
\begin{equation}
\label{II.8.8.5.2-fr}
  \xymatrix{
    X\ar[d]_{f} & V\ar[l]^{\pi}\ar[d]^{g'}\\
    Y & C\ar[l]_{\psi}
  }
  \qquad
  \xymatrix{
    X\ar[r]^{j}\ar[d]_{f} & V\ar[d]^{g'}\\
    Y\ar[r]_{\varepsilon} & C
  }
\tag{8.8.5.2}
\end{equation}
soient commutatifs
(\hyperref[II.8.5.1.2-fr]{8.5.1.2} et
\hyperref[II.8.5.1.3-fr]{8.5.1.3}). Nous allons montrer que (si l'on suppose
$\mathcal{L}$ ample pour $f$) les morphismes $g$ et $g'$ sont identiques.

En effet, la question est locale sur $Y$, et on peut donc supposer que
$Y=\operatorname{Spec}(A)$ est affine, et (en raison de
(\hyperref[II.8.8.1.3-fr]{8.8.1.3})) identifier $X$ à un ouvert de
$P=\operatorname{Proj}(S)$, où
\[
  S=A\oplus\bigoplus_{n\geqslant 1}
  \Gamma(X,\mathcal{L}^{\otimes n})~;
\]
on déduit alors de (\hyperref[II.8.8.1.4-fr]{8.8.1.4}) que
$\Gamma(X,\mathcal{O}_P(n))=\Gamma(X,\mathcal{L}^{\otimes n})$ pour tout
$n\in\mathbf{Z}$. Tenant compte de la définition de
$h=\operatorname{Spec}(\alpha)$, où $\alpha$ est le $p$-morphisme canonique
$\widetilde{S}\to\mathcal{S}_P^{\geqslant}$
(\hyperref[II.8.6.1.2-fr]{8.6.1.2}), il faut vérifier que la restriction à
$X$ de
$\alpha^\sharp:p^*(\widetilde{S})\to\mathcal{S}_P^{\geqslant}$ est
identique à $\tau$. Compte tenu de
(\hyperref[0.4.4.3-fr]{0, 4.4.3}), on est ramené à voir que si on compose
l'homomorphisme canonique
$\alpha_n:S_n\to\Gamma(P,\mathcal{O}_P(n))$ avec l'homomorphisme de
restriction
\[
  \Gamma(P,\mathcal{O}_P(n))\longrightarrow
  \Gamma(X,\mathcal{O}_P(n))=
  \Gamma(X,\mathcal{L}^{\otimes n}),
\]
on obtient l'identité pour tout $n>0$~; or, cela résulte aussitôt de la
définition de l'algèbre $S$ et de celle de $\alpha_n$
(\hyperref[II.2.6.2-fr]{2.6.2}).
\end{env}

\begin{proposition}[8.8.6]
\label{II.8.8.6-fr}
Supposons (avec les notations de (\hyperref[II.8.8.5-fr]{8.8.5})) que si
l'on pose $f=(f_0,\lambda)$, l'homomorphisme
$\lambda:\mathcal{O}_Y\to f_*(\mathcal{O}_X)$ soit bijectif~; alors~:
\begin{enumerate}
  \item[\rm{(i)}] Si on pose $g=(g_0,\mu)$,
  $\mu:\mathcal{O}_C\to g_*(\mathcal{O}_V)$ est un isomorphisme.
  \item[\rm{(ii)}] Si $X$ est intègre (resp. localement intègre et normal),
  $C$ est intègre (resp. normal).
\end{enumerate}
\end{proposition}

\begin{proof}
En effet, le $f$-morphisme $\tau^\flat$ est alors un \emph{isomorphisme}
\[
  \tau^\flat:\mathcal{S}=\psi_*(\mathcal{O}_C)
  \longrightarrow f_*(\pi_*(\mathcal{O}_V))
  =\psi_*(g_*(\mathcal{O}_V))
\]
et le $Y$-morphisme $g$ peut être considéré comme celui pour lequel
l'homomorphisme $\mathcal{A}(g)$
(\hyperref[II.1.1.2-fr]{1.1.2}) est égal à $\tau^\flat$. Pour voir que
$\mu$ est un isomorphisme de $\mathcal{O}_C$-Modules, il suffit
(\hyperref[II.1.4.2-fr]{1.4.2}) de voir que
\[
  \mathcal{A}(\mu):\psi_*(\mathcal{O}_C)
  \longrightarrow\psi_*(g_*(\mathcal{O}_V))
\]
est un isomorphisme. Mais par définition
(\hyperref[II.1.1.2-fr]{1.1.2}), on a
$\mathcal{A}(\mu)=\mathcal{A}(g)$, d'où la conclusion de (i).

Pour prouver (ii), on peut se ramener au cas où $Y$ est affine, donc
$\mathcal{S}=\widetilde{S}$, avec

\oldpage[II]{182}
$S=\bigoplus_{n\geqslant 0}\Gamma(X,\mathcal{L}^{\otimes n})$~;
l'hypothèse que $X$ est intègre entraîne que l'anneau $S$ est intègre
(\hyperref[I.7.4.4-fr]{I, 7.4.4}), donc il en est de même de $C$
(\hyperref[I.5.1.4-fr]{I, 5.1.4}). Pour démontrer que $C$ est normal, nous
utiliserons le lemme suivant~:

\begin{lemma}[8.8.6.1]
\label{II.8.8.6.1-fr}
Soit $Z$ un préschéma intègre normal. Alors l'anneau
$\Gamma(Z,\mathcal{O}_Z)$ est intègre et intégralement clos.
\end{lemma}

\begin{proof}
En effet, il résulte de
(\hyperref[I.8.2.1.1-fr]{I, 8.2.1.1}) que
$\Gamma(Z,\mathcal{O}_Z)$ est l'intersection, dans le corps des fonctions
rationnelles $R(Z)$, des anneaux intégralement clos $\mathcal{O}_z$ pour
$z\in Z$.
\end{proof}

Cela étant, montrons d'abord que $V$ est \emph{localement intègre et normal}~;
on peut pour cela se borner au cas où $X=\operatorname{Spec}(A)$ est affine,
d'anneau $A$ intègre et intégralement clos
(\hyperref[II.6.3.8-fr]{6.3.8}), et où
$\mathcal{L}=\mathcal{O}_X$. Comme alors
$V=\operatorname{Spec}(A[T])$ et que $A[T]$ est intègre et intégralement clos
([24], p.~99), cela établit notre assertion. D'autre part, pour tout ouvert
affine $U$ de $C$, $g^{-1}(U)$ est quasi-compact puisque le morphisme $g$ est
quasi-compact~; puisque $V$ est localement intègre, les composantes connexes
de $g^{-1}(U)$ sont des préschémas intègres ouverts dans $g^{-1}(U)$, donc en
nombre fini, et comme $V$ est normal, ces préschémas sont normaux
(\hyperref[II.6.3.8-fr]{6.3.8}). Alors $\Gamma(U,\mathcal{O}_C)$, qui est égal
à $\Gamma(g^{-1}(U),\mathcal{O}_V)$ en vertu de (i), est composé direct d'un
nombre fini d'anneaux intègres et intégralement clos
(\hyperref[II.8.8.6.1-fr]{8.8.6.1}), ce qui montre que $C$ est normal
(\hyperref[II.6.3.4-fr]{6.3.4}).
\end{proof}

\subsection{Le critère d'amplitude de Grauert : énoncé}
\label{subsection:II.8.9-fr}

Nous nous proposons de montrer que les propriétés démontrées dans
(\hyperref[II.8.8.2-fr]{8.8.2}) \emph{caractérisent} les
$\mathcal{O}_X$-Modules amples pour $f$, et de façon précise, de démontrer le
critère suivant~:

\begin{theorem}[8.9.1]
\label{II.8.9.1-fr}
\emph{(critère de Grauert).} Soient $Y$ un préschéma,
$p:X\to Y$ un morphisme séparé et quasi-compact, $\mathcal{L}$ un
$\mathcal{O}_X$-Module inversible. Pour que $\mathcal{L}$ soit ample
relativement à $p$, il faut et il suffit qu'il existe un $Y$-préschéma $C$,
une $Y$-section $\varepsilon:Y\to C$ de $C$, et un $Y$-morphisme
$q:\mathbf{V}(\mathcal{L})\to C$, ayant les propriétés suivantes~:
\begin{enumerate}
  \item[\rm{(i)}] Le diagramme
  \begin{equation}
  \label{II.8.9.1.1-fr}
    \xymatrix{
      X\ar[r]^{j}\ar[d]_{p} & \mathbf{V}(\mathcal{L})\ar[d]^{q}\\
      Y\ar[r]_{\varepsilon} & C
    }
  \tag{8.9.1.1}
  \end{equation}
  où $j$ est la section nulle du fibré vectoriel $\mathbf{V}(\mathcal{L})$,
  est commutatif.
  \item[\rm{(ii)}] La restriction de $q$ à
  $\mathbf{V}(\mathcal{L})-j(X)$ est une immersion ouverte quasi-compacte
  \[
    \mathbf{V}(\mathcal{L})-j(X)\longrightarrow C
  \]
  dont l'image ne rencontre pas $\varepsilon(Y)$.
\end{enumerate}
\end{theorem}

On notera que si $C$ est \emph{séparé} sur $Y$, on peut, dans la condition
(ii), supprimer l'hypothèse que l'immersion ouverte est quasi-compacte~; pour
voir en effet que cette dernière propriété est conséquence des autres
conditions, on peut se borner au cas où $Y$ est affine, et l'assertion découle
alors de (\hyperref[I.5.5.1-fr]{I, 5.5.1, (i)} et
\hyperref[I.5.5.10-fr]{5.5.10}). On peut aussi supprimer

\oldpage[II]{183}
la même hypothèse lorsque l'on suppose $X$ noethérien, car alors $V$ est aussi
noethérien et l'assertion découle cette fois de
(\hyperref[I.6.3.5-fr]{I, 6.3.5}).

\begin{corollary}[8.9.2]
\label{II.8.9.2-fr}
Si le morphisme $p:X\to Y$ est propre, on peut dans l'énoncé de
(\hyperref[II.8.9.1-fr]{8.9.1}), supposer $q$ propre, et remplacer
«~immersion ouverte~» par «~isomorphisme~».
\end{corollary}

De façon imagée, on peut dire (lorsque $p:X\to Y$ est propre) que
$\mathcal{L}$ est ample relativement à $p$ si et seulement si on peut
«~contracter~» la section nulle du fibré vectoriel $\mathbf{V}(\mathcal{L})$
en le préschéma de base $Y$. Un cas particulier important est celui où $Y$
est le spectre d'un corps, et où l'opération de «~contraction~» consiste donc
à contracter la section nulle de $\mathbf{V}(\mathcal{L})$ \emph{en un seul
point}.

\begin{env}[8.9.3]
\label{II.8.9.3-fr}
La nécessité des conditions de (\hyperref[II.8.9.1-fr]{8.9.1}) et le
corollaire (\hyperref[II.8.9.2-fr]{8.9.2}) découlent aussitôt de
(\hyperref[II.8.8.2-fr]{8.8.2}) et
(\hyperref[II.8.8.4-fr]{8.8.4}).

Pour démontrer la suffisance de (\hyperref[II.8.9.1-fr]{8.9.1}), nous
considérerons une situation un peu plus générale. Pour cela, posons (avec les
notations de (\hyperref[II.8.8.2-fr]{8.8.2}))
\[
  \mathcal{S}'=\bigoplus_{n\geqslant 0}\mathcal{L}^{\otimes n}
\]
et
\[
  V=\mathbf{V}(\mathcal{L})=\operatorname{Spec}(\mathcal{S}').
\]
Le sous-préschéma fermé $j(X)$, section nulle de $\mathbf{V}(\mathcal{L})$,
est défini par le faisceau quasi-cohérent d'idéaux
$\mathcal{J}=(\mathcal{S}'_+)^{\sim}$ de $\mathcal{O}_V$
(\hyperref[II.1.4.10-fr]{1.4.10}). Ce $\mathcal{O}_V$-Module est
\emph{inversible}, car la question est locale sur $X$ et cela revient à
remarquer que l'idéal $TA[T]$ dans un anneau de polynômes $A[T]$ est un
$A[T]$-module libre monogène. En outre, il est immédiat (toujours parce que la
question est locale sur $X$) que
\[
  \mathcal{L}=j^*(\mathcal{J})
\]
et
\[
  j_*(\mathcal{L})=\mathcal{J}/\mathcal{J}^2.
\]
D'autre part, si
\[
  \pi:\mathbf{V}(\mathcal{L})\longrightarrow X
\]
est le morphisme structural, on a
$\pi_*(\mathcal{J})=\mathcal{S}'_+$ et
$\pi_*(\mathcal{J}/\mathcal{J}^2)=\mathcal{L}$~; on a donc deux
homomorphismes canoniques
$\mathcal{L}\to\pi_*(\mathcal{J})\to\mathcal{L}$, le premier étant
l'injection canonique $\mathcal{L}\to\mathcal{S}'_+$, le second la projection
canonique de $\mathcal{S}'_+$ sur $\mathcal{S}'_1=\mathcal{L}$, et leur
composé est l'identité. On peut d'ailleurs plonger canoniquement
\[
  \pi_*(\mathcal{J})=\mathcal{S}'_+
  =\bigoplus_{n\geqslant 1}\mathcal{L}^{\otimes n}
\]
dans le \emph{produit}
\[
  \prod_{n\geqslant 1}\mathcal{L}^{\otimes n}
  =\varprojlim_n\pi_*(\mathcal{J}/\mathcal{J}^{n+1})
\]
(puisque
$\pi_*(\mathcal{J}/\mathcal{J}^{n+1})=
\mathcal{L}\oplus\mathcal{L}^{\otimes2}\oplus\cdots\oplus
\mathcal{L}^{\otimes n}$), et on a donc encore deux homomorphismes canoniques
\begin{equation}
\label{II.8.9.3.1-fr}
  \mathcal{L}\longrightarrow
  \varprojlim_n\pi_*(\mathcal{J}/\mathcal{J}^{n+1})
  \longrightarrow\mathcal{L}
\tag{8.9.3.1}
\end{equation}
dont le composé est l'identité.

Cela étant, la généralisation de (\hyperref[II.8.9.1-fr]{8.9.1}) que nous
allons prouver est la suivante~:
\end{env}

\begin{proposition}[8.9.4]
\label{II.8.9.4-fr}
Soient $Y$ un préschéma, $V$ un $Y$-préschéma, $X$ un sous-préschéma fermé de
$V$, défini par un Idéal $\mathcal{J}$ de $\mathcal{O}_V$, qui est un
$\mathcal{O}_V$-Module inversible~; si $j:X\to V$ est

\oldpage[II]{184}
l'injection canonique, on pose
$\mathcal{L}=j^*(\mathcal{J})=
\mathcal{J}\otimes_{\mathcal{O}_V}\mathcal{O}_X$, de sorte que
$j_*(\mathcal{L})=\mathcal{J}/\mathcal{J}^2$. On suppose que le morphisme
structural $p:X\to Y$ est séparé et quasi-compact, et que les conditions
suivantes sont remplies~:
\begin{enumerate}
  \item[\rm{(i)}] Il existe un $Y$-morphisme de type fini $\pi:V\to X$ tel
  que $\pi\circ j=1_X$, et par suite
  $\pi_*(\mathcal{J}/\mathcal{J}^2)=\mathcal{L}$.
  \item[\rm{(ii)}] Il existe un homomorphisme de $\mathcal{O}_X$-Modules
  \[
    \varphi:\mathcal{L}\longrightarrow
    \varprojlim_n\pi_*(\mathcal{J}/\mathcal{J}^{n+1})
  \]
  tel que le composé
  \[
    \mathcal{L}\xrightarrow{\varphi}
    \varprojlim_n\pi_*(\mathcal{J}/\mathcal{J}^{n+1})
    \xrightarrow{\alpha}\pi_*(\mathcal{J}/\mathcal{J}^2)=\mathcal{L}
  \]
  (où $\alpha$ est l'homomorphisme canonique) soit l'identité.
  \item[\rm{(iii)}] Il existe un $Y$-préschéma $C$, une $Y$-section
  $\varepsilon$ de $C$ et un $Y$-morphisme $q:V\to C$ tels que le diagramme
  \begin{equation}
  \label{II.8.9.4.1-fr}
    \xymatrix{
      X\ar[r]^{j}\ar[d]_{p} & V\ar[d]^{q}\\
      Y\ar[r]_{\varepsilon} & C
    }
  \tag{8.9.4.1}
  \end{equation}
  soit commutatif.
  \item[\rm{(iv)}] La restriction de $q$ à $W=V-j(X)$ est une immersion
  ouverte quasi-compacte dans $C$, dont l'image ne rencontre pas
  $\varepsilon(Y)$.
\end{enumerate}
Dans ces conditions, $\mathcal{L}$ est ample relativement à $p$.
\end{proposition}

\subsection{Le critère d'amplitude de Grauert : démonstration}
\label{subsection:II.8.10-fr}

\begin{lemma}[8.10.1]
\label{II.8.10.1-fr}
Soient $\pi:V\to X$ un morphisme, $j:X\to V$ une $X$-section de $V$ qui soit
une immersion fermée, $\mathcal{J}$ l'Idéal quasi-cohérent de
$\mathcal{O}_V$ définissant le sous-préschéma fermé de $V$ associé à $j$.
\begin{enumerate}
  \item[\rm{(i)}] Pour tout $n\geqslant 0$,
  $\pi_*(\mathcal{O}_V/\mathcal{J}^{n+1})$ et
  $\pi_*(\mathcal{J}/\mathcal{J}^{n+1})$ sont des
  $\mathcal{O}_X$-Modules quasi-cohérents, et on a
  $\pi_*(\mathcal{O}_V/\mathcal{J})=\mathcal{O}_X$,
  $\pi_*(\mathcal{J}/\mathcal{J}^2)=j^*(\mathcal{J})$.
  \item[\rm{(ii)}] Si $X=\{\xi\}=\operatorname{Spec}(k)$, où $k$ est un
  corps, $\varprojlim_n\pi_*(\mathcal{O}_V/\mathcal{J}^{n+1})$ est isomorphe
  au séparé complété de l'anneau local $\mathcal{O}_{j(\xi)}$ pour la
  topologie $\mathfrak{m}_{j(\xi)}$-préadique.
  \item[\rm{(iii)}] Supposons que $\mathcal{J}$ soit un
  $\mathcal{O}_V$-Module inversible (ce qui entraîne que
  \[
    \mathcal{L}=j^*(\mathcal{J})=
    \pi_*(\mathcal{J}/\mathcal{J}^2)
  \]
  est un $\mathcal{O}_X$-Module inversible) et qu'il existe un homomorphisme
  \[
    \varphi:\mathcal{L}\longrightarrow
    \varprojlim_n\pi_*(\mathcal{J}/\mathcal{J}^{n+1})
  \]
  tel que le composé
  \[
    \mathcal{L}\xrightarrow{\varphi}
    \varprojlim_n\pi_*(\mathcal{J}/\mathcal{J}^{n+1})
    \xrightarrow{\alpha}\pi_*(\mathcal{J}/\mathcal{J}^2)
  \]
  ($\alpha$ homomorphisme canonique) soit l'identité. Si on pose
  $\mathcal{S}=\bigoplus_{n\geqslant0}\mathcal{L}^{\otimes n}$, on déduit
  alors canoniquement de $\varphi$ un isomorphisme de
  $\mathcal{O}_X$-Algèbres du complété $\widehat{\mathcal{S}}$ de
  $\mathcal{S}$ relatif à sa filtration canonique (complété qui est isomorphe
  au produit $\prod_{n\geqslant0}\mathcal{L}^{\otimes n}$) sur
  $\varprojlim_n\pi_*(\mathcal{O}_V/\mathcal{J}^{n+1})$.
\end{enumerate}
\end{lemma}

\begin{proof}
Notons en premier lieu que le support du $\mathcal{O}_V$-Module
$\mathcal{O}_V/\mathcal{J}^{n+1}$ est $j(X)$ et celui de
$\mathcal{J}/\mathcal{J}^{n+1}$ est contenu dans $j(X)$. Dans le cas (ii),
$j(X)$ est un point fermé $j(\xi)$ de $V$,

\oldpage[II]{185}
et par définition $\pi_*(\mathcal{O}_V/\mathcal{J}^{n+1})$ est la fibre de
$\mathcal{O}_V/\mathcal{J}^{n+1}$ au point $j(\xi)$, c'est-à-dire, en posant
$C=\mathcal{O}_{j(\xi)}$ et en désignant par $\mathfrak{m}$ l'idéal maximal
de $C$, le $C$-module $C/\mathfrak{m}^{n+1}$~; l'assertion (ii) est alors
évidente.

Pour démontrer (i), notons que la question est locale sur $X$~; on peut donc
se restreindre au cas où $X$ est affine. Soit $U$ un ouvert affine dans $V$~;
$j(X)\cap U$ est un ouvert affine dans $j(X)$, donc
$U_0=\pi(j(X)\cap U)$, qui lui est isomorphe, un ouvert affine dans $X$~;
pour tout ouvert affine $W_0\subset U_0$ dans $X$,
$W=\pi^{-1}(W_0)\cap U$ est un ouvert affine dans $V$, puisque $X$ est un
schéma (\hyperref[I.5.5.10-fr]{I, 5.5.10})~; en particulier
$U'=U\cap\pi^{-1}(U_0)$ est un ouvert affine dans $V$ et on a évidemment
$\pi(U')=U_0$ et $j(U_0)=j(X)\cap U$. Cela étant, par définition
\[
  \Gamma(W_0,\pi_*(\mathcal{O}_V/\mathcal{J}^{n+1}))
  =\Gamma(\pi^{-1}(W_0),\mathcal{O}_V/\mathcal{J}^{n+1})~;
\]
mais comme tout point de $\pi^{-1}(W_0)$ n'appartenant pas à $j(W_0)$ a un
voisinage ouvert dans $\pi^{-1}(W_0)$ ne rencontrant pas $j(X)$ et dans
lequel $\mathcal{O}_V/\mathcal{J}^{n+1}$ est donc nul, il est clair que les
sections de $\mathcal{O}_V/\mathcal{J}^{n+1}$ au-dessus de $\pi^{-1}(W_0)$ et
au-dessus de $W$ se correspondent biunivoquement. Autrement dit, si $\pi'$ est
la restriction de $\pi$ à $U'$, les $(\mathcal{O}_X|U_0)$-Modules
\[
  \pi_*(\mathcal{O}_V/\mathcal{J}^{n+1})|U_0
  \quad\hbox{et}\quad
  \pi'_*\bigl((\mathcal{O}_V/\mathcal{J}^{n+1})|U'\bigr)
\]
sont identiques. Comme $U'$ et $U_0$ sont affines et que les $U_0$ recouvrent
$X$, on en conclut (\hyperref[I.1.6.3-fr]{I, 1.6.3}) que
$\pi_*(\mathcal{O}_V/\mathcal{J}^{n+1})$ est quasi-cohérent, et la
démonstration est identique pour
$\pi_*(\mathcal{J}/\mathcal{J}^{n+1})$.

Pour démontrer enfin (iii), notons que $\mathcal{S}$ n'est autre que
$\mathbf{S}_{\mathcal{O}_X}(\mathcal{L})$~; on déduit donc canoniquement de
$\varphi$ un homomorphisme de $\mathcal{O}_X$-Algèbres
\[
  \psi:\mathcal{S}\longrightarrow
  \varprojlim_n\pi_*(\mathcal{O}_V/\mathcal{J}^{n+1})
\]
(\hyperref[II.1.7.4-fr]{1.7.4})~; en outre, cet homomorphisme applique
$\mathcal{L}^{\otimes n}$ dans
$\varprojlim_m\pi_*(\mathcal{J}^n/\mathcal{J}^{m+1})$, donc est continu pour
les topologies considérées, et se prolonge bien par suite en un homomorphisme
\[
  \widehat{\psi}:\widehat{\mathcal{S}}\longrightarrow
  \varprojlim_n\pi_*(\mathcal{O}_V/\mathcal{J}^{n+1}).
\]
Pour voir qu'il s'agit d'un isomorphisme, on peut, comme dans la démonstration
de (i), se réduire au cas où $X=\operatorname{Spec}(A)$ et
$V=\operatorname{Spec}(B)$ sont affines, avec
$\mathcal{J}=\widetilde{\mathfrak{J}}$, où $\mathfrak{J}$ est un idéal de
$B$~; il correspond à $\pi$ une injection $A\to B$, identifiant $A$ à un
sous-anneau de $B$ supplémentaire de $\mathfrak{J}$, et $\mathcal{L}$ (resp.
$\pi_*(\mathcal{O}_V/\mathcal{J}^{n+1})$) est le $\mathcal{O}_X$-Module
quasi-cohérent associé au $A$-module $L=\mathfrak{J}/\mathfrak{J}^2$ (resp.
$B/\mathfrak{J}^{n+1}$). Comme $\mathcal{J}$ est un
$\mathcal{O}_V$-Module \emph{inversible}, on peut en outre supposer que
$\mathfrak{J}=Bt$, où $t$ est non diviseur de 0 dans $B$. De la relation
$B=A\oplus Bt$ on déduit alors, pour tout $n>0$,
\[
  B=A\oplus At\oplus At^2\oplus\cdots\oplus At^n\oplus Bt^{n+1}
\]
donc on a un $A$-isomorphisme canonique de l'anneau de séries formelles
$A[[T]]$ sur $C=\varprojlim B/\mathfrak{J}^{n+1}$ faisant correspondre $t$ à
$T$. D'autre part, on a $L=A\bar t$, où $\bar t$ est la classe de $t$
mod. $Bt^2$, et l'homomorphisme $\varphi$ applique par hypothèse $\bar t$ sur
un élément $t'\in C$ congru à $t$ mod. $Ct^2$. On en déduit par récurrence
sur $n$ que
\[
  A\oplus At'\oplus\cdots\oplus At'{}^n\oplus Ct^{n+1}
  =A\oplus At\oplus\cdots\oplus At^n\oplus Ct^{n+1}
\]
ce qui prouve que l'homomorphisme $\widehat{\psi}$ correspond bien à un
isomorphisme de $\prod_{n\geqslant0}L^{\otimes n}$ sur $C$.
\end{proof}

\begin{lemma}[8.10.2]
\label{II.8.10.2-fr}
Sous les hypothèses de (\hyperref[II.8.10.1-fr]{8.10.1}), soit
$g:X'\to X$ un morphisme,

\oldpage[II]{186}
soit $V'=V\times_X X'$, et soient $\pi':V'\to X'$,
$g':V'\to V$ les projections canoniques, de sorte qu'on a le diagramme
commutatif
\[
  \xymatrix{
    V\ar[d]_{\pi} & V'\ar[l]^{g'}\ar[d]^{\pi'}\\
    X & X'\ar[l]_{g}
  }
\]
Alors $j'=j\times1_{X'}$ est une $X'$-section de $V'$ qui est une immersion
fermée, et $\mathcal{J}'=(g')^*(\mathcal{J})\mathcal{O}_{V'}$ est l'Idéal
quasi-cohérent de $\mathcal{O}_{V'}$ définissant le sous-préschéma fermé de
$V'$ associé à $j'$. En outre, on a
\[
  \pi'_*(\mathcal{O}_{V'}/\mathcal{J}'^{n+1})
  =g^*(\pi_*(\mathcal{O}_V/\mathcal{J}^{n+1})).
\]
Enfin, $\mathcal{J}'$ est un $\mathcal{O}_{V'}$-Module canoniquement isomorphe
à $(g')^*(\mathcal{J})$ et est en particulier inversible si $\mathcal{J}$ est
un $\mathcal{O}_V$-Module inversible.
\end{lemma}

\begin{proof}
Le fait que $j'$ est une immersion fermée découle de
(\hyperref[I.4.3.1-fr]{I, 4.3.1}), et c'est une $X'$-section de $V'$ par
fonctorialité de l'extension du préschéma de base. En outre, si $Z$ (resp.
$Z'$) est le sous-préschéma fermé de $V$ (resp. $V'$) associé à $j$ (resp.
$j'$), on a $Z'=g'{}^{-1}(Z)$
(\hyperref[I.4.3.1-fr]{I, 4.3.1}), et la seconde assertion résulte alors de
(\hyperref[I.4.4.5-fr]{I, 4.4.5}). Pour démontrer les autres assertions, on
voit comme dans (\hyperref[II.8.10.1-fr]{8.10.1}) qu'on peut se ramener au cas
où $X$, $V$ et $X'$ (donc aussi $V'$) sont affines~; gardons les notations de
la démonstration de (\hyperref[II.8.10.1-fr]{8.10.1}) et soit
$X'=\operatorname{Spec}(A')$. On a alors $V'=\operatorname{Spec}(B')$ où
$B'=B\otimes_A A'$, et $\mathcal{J}'=\widetilde{\mathfrak{J}'}$, avec
$\mathfrak{J}'=\operatorname{Im}(\mathfrak{J}\otimes_A A')$. On a donc
\[
  B'/\mathfrak{J}'^{n+1}=(B/\mathfrak{J}^{n+1})\otimes_A A'~;
\]
en outre, comme $\mathfrak{J}$ est facteur direct (en tant que $A$-module) de
$B$, $\mathfrak{J}\otimes_A A'$ est facteur direct (en tant que $A'$-module)
de $B'$, donc s'identifie canoniquement à $\mathfrak{J}'$.
\end{proof}

\begin{corollary}[8.10.3]
\label{II.8.10.3-fr}
Supposons vérifiées les hypothèses de (\hyperref[II.8.10.1-fr]{8.10.1}) et
supposons en outre que $\pi$ soit de type fini et que $\mathcal{J}$ soit un
$\mathcal{O}_V$-Module inversible. Alors, pour tout $x\in X$, l'anneau local
au point $j(x)$ de la fibre $\pi^{-1}(x)$ est un anneau régulier (donc
intègre) de dimension 1, dont le complété est isomorphe à l'anneau de séries
formelles $k(x)[[T]]$ ($T$ indéterminée)~; en outre il n'existe qu'une seule
composante irréductible de $\pi^{-1}(x)$ contenant $j(x)$.
\end{corollary}

\begin{proof}
Comme $\pi^{-1}(x)=V\times_X\operatorname{Spec}(k(x))$, on est ramené par
(\hyperref[II.8.10.2-fr]{8.10.2}) au cas où $X$ est le spectre d'un corps
$K$. Comme $\pi$ est de type fini
(\hyperref[I.6.3.4-fr]{I, 6.3.4, (iv)}), $\mathcal{O}_{j(x)}$ est alors un
anneau local noethérien, donc séparé pour la topologie
$\mathfrak{m}_{j(x)}$-préadique
(\hyperref[0.7.3.5-fr]{0, 7.3.5})~; il résulte de
(\hyperref[II.8.10.1-fr]{8.10.1, (ii) et (iii)}) que le complété de cet anneau
est isomorphe à $K[[T]]$, et par suite $\mathcal{O}_{j(x)}$ est régulier et de
dimension 1 ([1], p.~17-01, th.~1)~; enfin, puisque $\mathcal{O}_{j(x)}$ est
intègre, $j(x)$ n'appartient qu'à une seule des composantes irréductibles (en
nombre fini) de $V$ (\hyperref[I.5.1.4-fr]{I, 5.1.4}).
\end{proof}

\begin{corollary}[8.10.4]
\label{II.8.10.4-fr}
Supposons vérifiées les hypothèses de (\hyperref[II.8.10.1-fr]{8.10.1}) et en
outre supposons que $\mathcal{J}$ soit un $\mathcal{O}_V$-Module inversible.
Soit $W=V-j(X)$~; pour tout Idéal quasi-cohérent $\mathcal{K}$ de
$\mathcal{O}_X$, posons
$\mathcal{K}_V=\pi^*(\mathcal{K})\mathcal{O}_V$ et
$\mathcal{K}_W=\mathcal{K}_V|W$. Alors $\mathcal{K}_V$ est le plus grand
Idéal quasi-cohérent de $\mathcal{O}_V$ dont la restriction à $W$ soit
$\mathcal{K}_W$.
\end{corollary}

\begin{proof}
En effet, on voit comme dans (\hyperref[II.8.10.1-fr]{8.10.1}) que la question
est locale sur $X$ et sur $V$~; on peut donc garder les notations de la
démonstration de (\hyperref[II.8.10.1-fr]{8.10.1}), avec
$\mathfrak{J}=Bt$, où $t$ n'est pas diviseur de 0 dans $B$. En outre, on a
$W=\operatorname{Spec}(B_t)$ et $\mathcal{K}=\widetilde{\mathfrak{K}}$, où
$\mathfrak{K}$ est un idéal de $A$~; d'où
$\pi^*(\mathcal{K})\mathcal{O}_V=(\mathfrak{K}.B)^{\sim}$
(\hyperref[I.1.6.9-fr]{I, 1.6.9}),
$\mathcal{K}_W=(\mathfrak{K}.B_t)^{\sim}$, et le plus grand idéal

\oldpage[II]{187}
de $B$ dont l'image canonique dans $B_t$ soit $\mathfrak{K}.B_t$ est l'image
réciproque de $\mathfrak{K}.B_t$, c'est-à-dire l'ensemble des $s\in B$ tels
que pour un entier $n>0$, on ait $t^ns\in\mathfrak{K}.B$. Il faut montrer que
cette dernière relation entraîne $s\in\mathfrak{K}.B$, ou encore que l'image
canonique de $t$ n'est pas diviseur de 0 dans
$B/\mathfrak{K}B=(A/\mathfrak{K})\otimes_A B$, ce qui résulte de
(\hyperref[II.8.10.2-fr]{8.10.2}) appliqué à
$X'=\operatorname{Spec}(A/\mathfrak{K})$.
\end{proof}

\begin{corollary}[8.10.5]
\label{II.8.10.5-fr}
Supposons vérifiées les hypothèses de (\hyperref[II.8.10.3-fr]{8.10.3})~;
soient $W=V-j(X)$, $x$ un point de $X$, $\mathcal{K}$ un Idéal
quasi-cohérent de $\mathcal{O}_X$, $z$ le point générique de la composante
irréductible de $\pi^{-1}(x)$ contenant $j(x)$
(\hyperref[II.8.10.3-fr]{8.10.3}).
\begin{enumerate}
  \item[\rm{(i)}] Soit $g$ une section de $\mathcal{O}_V$ au-dessus de $V$
  telle que $g|W$ soit une section de $\mathcal{K}_W$ au-dessus de $W$
  (notations de (\hyperref[II.8.10.4-fr]{8.10.4})). Alors $g$ est une section
  de $\mathcal{K}_V$~; si de plus $g(z)\neq0$, et si, pour tout entier $m>0$,
  on désigne par $g_m^x$ le germe au point $x$ de l'image canonique $g_m$ de
  $g$ dans
  $\Gamma(X,\pi_*(\mathcal{O}_V/\mathcal{J}^{m+1}))$, alors il existe un
  entier $m>0$ tel que l'image de $g_m^x$ dans
  \[
    (\pi_*(\mathcal{O}_V/\mathcal{J}^{m+1}))_x
    \otimes_{\mathcal{O}_x}k(x)
  \]
  soit $\neq0$.
  \item[\rm{(ii)}] Supposons en outre que les conditions de
  (\hyperref[II.8.10.1-fr]{8.10.1, (iii)}) soient remplies. Alors, s'il existe
  une section $g$ de $\mathcal{K}_V$ au-dessus de $V$ telle que $g(z)\neq0$,
  il existe un entier $n\geqslant0$ et une section $f$ de
  $\mathcal{K}.\mathcal{L}^{\otimes n}=
  \mathcal{K}\otimes\mathcal{L}^{\otimes n}\subset\mathcal{L}^{\otimes n}$
  telle que $f(x)\neq0$. Si $g$ est une section de $\mathcal{J}$, on peut
  prendre $n>0$.
\end{enumerate}
\end{corollary}

\begin{proof}
\begin{enumerate}
  \item[\rm{(i)}] Comme l'Idéal de $\mathcal{O}_W$ engendré par $g|W$ est
  contenu dans $\mathcal{K}_W$ par hypothèse, l'Idéal de $\mathcal{O}_V$
  engendré par $g$ est contenu dans $\mathcal{K}_V$ par
  (\hyperref[II.8.10.4-fr]{8.10.4}), autrement dit $g$ est une section de
  $\mathcal{K}_V$. Pour démontrer la seconde assertion de (i), on peut encore
  supposer $X$ et $V$ affines et conserver les notations de
  (\hyperref[II.8.10.1-fr]{8.10.1})~; la fibre $\pi^{-1}(x)$ est alors affine
  d'anneau $B'=B\otimes_A k(x)$, et il existe dans $B'$ un élément $t'$ non
  diviseur de 0 tel que $B'=k(x)\oplus B't'$. Comme $j(x)$ est une
  spécialisation de $z$ et que $g(z)\neq0$, on a nécessairement
  $g_{j(x)}\neq0$. Mais $\mathcal{O}_{j(x)}$ est un anneau local séparé
  (\hyperref[II.8.10.3-fr]{8.10.3}), qui se plonge donc dans son complété, et
  l'image de $g$ dans ce complété n'est par suite pas nulle. Or, ce complété
  est isomorphe à $\varprojlim_n(B'/B't'^{n+1})$
  (\hyperref[II.8.10.3-fr]{8.10.3})~; si $g'=g\otimes1\in B'$, il existe donc
  un entier $m$ tel que $g'\notin B't'^{m+1}$, ou encore l'image $g'_m$ de
  $g'$ dans $B'/B't'^{m+1}$ n'est pas nulle. Mais comme $g'_m$ n'est autre que
  l'image de $g_m^x$, notre assertion est démontrée.
  \item[\rm{(ii)}] En vertu de
  (\hyperref[II.8.10.1-fr]{8.10.1, (iii)})
  $\pi_*(\mathcal{O}_V/\mathcal{J}^{m+1})$ est isomorphe à la somme directe
  des $\mathcal{L}^{\otimes k}$ pour $0\leqslant k\leqslant m$~; désignons
  par $f_k$ la section de $\mathcal{L}^{\otimes k}$ au-dessus de $X$,
  composante de l'élément de
  $\bigoplus_{k=0}^m\Gamma(X,\mathcal{L}^{\otimes k})$ qui correspond à $g_m$
  par cet isomorphisme. Choisissant $m$ comme dans (i), il y a donc un indice
  $k$ tel que $f_k(x)\neq0$ d'après (i). Pour voir que $f_k$ est une section
  de $\mathcal{K}\mathcal{L}^{\otimes k}$, il suffit de considérer comme
  ci-dessus le cas où $X$ et $V$ sont affines, et cela résulte aussitôt alors
  de ce que $g\in\mathfrak{K}.B$ (notations de
  (\hyperref[II.8.10.4-fr]{8.10.4})). La dernière assertion résulte de ce que
  l'hypothèse $g\in\Gamma(V,\mathcal{J})$ entraîne $f_0=0$.
\end{enumerate}
\end{proof}

\begin{env}[8.10.6]
\label{II.8.10.6-fr}
\emph{Démonstration de (\hyperref[II.8.9.4-fr]{8.9.4}).} La question est
locale sur $Y$ (\hyperref[II.4.6.4-fr]{4.6.4})~; comme $\varepsilon$ est une
$Y$-section, on peut donc remplacer $C$ par un voisinage ouvert affine $U$
d'un point de $\varepsilon(Y)$ tel que $\varepsilon(Y)\cap U$ soit fermé dans
$U$. Autrement dit, on peut supposer $C$ affine, et $Y$ sous-préschéma fermé
de $C$ (donc affine) défini par un faisceau quasi-

\oldpage[II]{188}
cohérent $\mathcal{I}$ d'idéaux de $\mathcal{O}_C$. Comme $p$ est séparé et
quasi-compact, $X$ est alors un schéma quasi-compact, et on est ramené à
prouver que $\mathcal{L}$ est ample
(\hyperref[II.4.6.4-fr]{4.6.4}). En vertu du critère
(\hyperref[II.4.5.2-fr]{4.5.2, a)}), on doit donc prouver qui ce suit~: pour
tout Idéal quasi-cohérent $\mathcal{K}$ de $\mathcal{O}_X$ et tout point
$x\in X$ n'appartenant pas au support de $\mathcal{O}_X/\mathcal{K}$, il
existe un entier $n>0$ et une section $f$ de
$\mathcal{K}\otimes\mathcal{L}^{\otimes n}$ au-dessus de $X$ telle que
$f(x)\neq0$.

Pour cela, posons
\[
  \mathcal{K}_V=\pi^*(\mathcal{K})\mathcal{O}_V
\]
et
\[
  \mathcal{K}_W=\mathcal{K}_V|W,\qquad\hbox{où}\qquad W=V-j(X)~;
\]
comme la restriction de $q$ à $W$ est une immersion quasi-compacte dans $C$,
il résulte de (\hyperref[I.9.4.2-fr]{I, 9.4.2}) que $\mathcal{K}_W$ est la
restriction à $W$ d'un Idéal quasi-cohérent $\mathcal{K}'_V$ de
$\mathcal{O}_V$, de la forme
\[
  \mathcal{K}'_V=q^*(\mathcal{K}_C)\mathcal{O}_V
\]
où $\mathcal{K}_C$ est un Idéal quasi-cohérent de $\mathcal{O}_C$. En outre,
comme par hypothèse $q^{-1}(Y)\subset j(X)$ et que $Y$ est défini par l'Idéal
$\mathcal{I}$, la restriction à $W$ de $q^*(\mathcal{I})\mathcal{O}_V$ est
identique à celle de $\mathcal{O}_V$, donc $\mathcal{K}_W$ est aussi la
restriction à $W$ de $q^*(\mathcal{I}\mathcal{K}_C)\mathcal{O}_V$, et l'on
peut donc supposer que l'on a $\mathcal{K}_C\subset\mathcal{I}$, d'où
\begin{equation}
\label{II.8.10.6.1-fr}
  \mathcal{K}'_V\subset q^*(\mathcal{I})\mathcal{O}_V\subset\mathcal{J}
\tag{8.10.6.1}
\end{equation}
compte tenu de (\hyperref[I.4.4.6-fr]{I, 4.4.6}) et de la commutativité de
(\hyperref[II.8.9.4.1-fr]{8.9.4.1}). En outre, on déduit de
(\hyperref[II.8.10.4-fr]{8.10.4}) que l'on a
\begin{equation}
\label{II.8.10.6.2-fr}
  \mathcal{K}'_V\subset\mathcal{K}_V.
\tag{8.10.6.2}
\end{equation}

Cela étant, il résulte de (\hyperref[II.8.10.3-fr]{8.10.3}) que $j(x)$
appartient à une seule composante irréductible de $\pi^{-1}(x)$~; soit $z$ le
point générique de cette composante, et posons $z'=q(z)$. En vertu de
(\hyperref[II.8.10.5-fr]{8.10.5}), la démonstration sera achevée (compte tenu
de (\hyperref[II.8.10.6.1-fr]{8.10.6.1}) et
(\hyperref[II.8.10.6.2-fr]{8.10.6.2})) si nous prouvons l'existence d'une
section $g$ de $\mathcal{K}'_V$ au-dessus de $V$ telle que $g(z)\neq0$. Or,
par hypothèse, $\mathcal{K}$ a une restriction égale à celle de
$\mathcal{O}_X$ dans un voisinage ouvert de $x$~; d'autre part, il résulte de
(\hyperref[II.8.10.3-fr]{8.10.3}) que l'on a $z\neq j(x)$, donc $z\in W$, et
par suite $(\mathcal{K}_W)_z=\mathcal{O}_{V,z}$, d'où par définition
$(\mathcal{K}_C)_{z'}=\mathcal{O}_{C,z'}$. Puisque $C$ est affine, il y a donc
une section $g'$ de $\mathcal{K}_C$ au-dessus de $C$ telle que
$g'(z')\neq0$, et en prenant pour $g$ la section de $\mathcal{K}'_V$
correspondant canoniquement à $g'$, on a bien $g(z)\neq0$, ce qui achève la
démonstration.
\end{env}

\begin{remark}[8.10.7]
\label{II.8.10.7-fr}
Nous ignorons si, dans l'énoncé de (\hyperref[II.8.9.4-fr]{8.9.4}), la
condition (ii) est ou non superflue. En tout cas, la conclusion ne subsiste
plus si l'on ne suppose pas l'existence d'un $Y$-morphisme $\pi:V\to X$ tel
que $\pi\circ j=1_X$~; indiquons brièvement comment on peut en effet former un
contre-exemple, dont les détails ne pourront être développés que plus tard.
On prend $Y=\operatorname{Spec}(k)$ où $k$ est un corps,
$C=\operatorname{Spec}(A)$, où $A=k[T_1,T_2]$, la $Y$-section $\varepsilon$
correspondant à l'homomorphisme d'augmentation $A\to k$. On désigne par $C'$
le schéma déduit de $C$ par éclatement du point fermé $a=\varepsilon(Y)$ de
$C$~; si $D$ est l'image réciproque de $a$ dans $C'$, on considère dans $D$
un point fermé $b$,

\oldpage[II]{189}
et on désigne par $V$ le schéma déduit de $C'$ par éclatement de $b$~; $X$ est
le sous-préschéma fermé de $V$, image réciproque de $a$ par le morphisme
structural $q:V\to C$. On montre que $X$ est réunion de deux composantes
irréductibles $X_1$, $X_2$, où $X_1$ est l'image réciproque de $b$ dans $V$.
Il est immédiat que l'Idéal $\mathcal{J}$ de $\mathcal{O}_V$ qui définit $X$
est encore inversible, mais on prouve que
$j^*(\mathcal{J})=\mathcal{L}$ ($j$ injection canonique $X\to V$) n'est pas
ample, par la considération du «~degré~» de l'image réciproque de
$\mathcal{L}$ sur $X_1$, qui devrait être $>0$ si $\mathcal{L}$ était ample,
et qu'on montre (par un calcul élémentaire d'intersections) être égal à 0.
\end{remark}

\subsection{Unicité des contractions}
\label{subsection:II.8.11-fr}

\begin{lemma}[8.11.1]
\label{II.8.11.1-fr}
Soient $U$, $V$ deux préschémas,
$h=(h_0,\lambda):U\to V$ un morphisme surjectif. On suppose que~:
\begin{enumerate}
  \item[$1^\circ$]
  $\lambda:\mathcal{O}_V\to h_*(\mathcal{O}_U)
  =(h_0)_*(\mathcal{O}_U)$ est un isomorphisme.
  \item[$2^\circ$] L'espace sous-jacent à $V$ s'identifie à l'espace quotient
  de l'espace sous-jacent à $U$ par la relation
  $R:h_0(x)=h_0(y)$ (condition toujours vérifiée lorsque le morphisme $h$ est
  ouvert, ou fermé, ou a fortiori lorsque $h$ est propre).
\end{enumerate}
Alors, pour tout préschéma $W$, l'application
\begin{equation}
\label{II.8.11.1.1-fr}
  \operatorname{Hom}(V,W)\longrightarrow\operatorname{Hom}(U,W)
\tag{8.11.1.1}
\end{equation}
qui, à tout morphisme $v=(v_0,\nu)$ de $V$ dans $W$, fait correspondre le
morphisme $u=v\circ h=(u_0,\mu)$, est une bijection de
$\operatorname{Hom}(V,W)$ sur l'ensemble des $u$ tels que $u_0$ soit
constante sur toute fibre $h_0^{-1}(x)$.
\end{lemma}

\begin{proof}
Il est clair que si $u=v\circ h$, donc $u_0=v_0\circ h_0$, $u_0$ est
constante sur tout ensemble $h_0^{-1}(x)$. Inversement, si $u$ a cette
propriété, montrons qu'il existe un $v\in\operatorname{Hom}(V,W)$ et un seul
tel que $u=v\circ h$. L'existence et l'unicité de l'application continue
$v_0:V\to W$ telle que $u_0=v_0\circ h_0$ résultent de l'hypothèse, puisque
$h_0$ s'identifie à l'application canonique de $U$ sur $U/R$. On peut d'autre
part, en remplaçant au besoin $V$ par un préschéma isomorphe, supposer que
$\lambda$ est l'identité~; par hypothèse, $\mu$ est alors un homomorphisme
\[
  \mu:\mathcal{O}_W\longrightarrow(u_0)_*(\mathcal{O}_U)
  =(v_0)_*((h_0)_*(\mathcal{O}_U))
\]
tel que l'homomorphisme correspondant
$\mu^\sharp:u_0^*(\mathcal{O}_W)\to\mathcal{O}_U$ soit local sur chaque
fibre. Comme
$(v_0)_*((h_0)_*(\mathcal{O}_U))=(v_0)_*(\mathcal{O}_V)$, on doit
nécessairement avoir $\nu=\mu$, et tout revient à voir que l'homomorphisme
correspondant $\nu^\sharp:v_0^*(\mathcal{O}_W)\to\mathcal{O}_V$ est local sur
chaque fibre. Or, tout $y\in V$ est de la forme $h_0(x)$ pour un $x\in U$~;
posons $z=v_0(y)=u_0(x)$. Alors
(\hyperref[0.3.5.5-fr]{0, 3.5.5}) l'homomorphisme $\mu_x^\sharp$ se factorise
en
\[
  \mu_x^\sharp:\mathcal{O}_z
  \xrightarrow{\nu_y^\sharp}\mathcal{O}_y
  \xrightarrow{\lambda_x^\sharp}\mathcal{O}_x.
\]
Par hypothèse $\lambda_x^\sharp$ et $\mu_x^\sharp$ sont des homomorphismes
locaux~; donc $\lambda_x^\sharp$ transforme tout élément inversible de
$\mathcal{O}_y$ en un élément inversible de $\mathcal{O}_x$~; si
$\nu_y^\sharp$ transformait un élément non inversible de $\mathcal{O}_z$ en
un élément inversible de $\mathcal{O}_y$, $\mu_x^\sharp$ transformerait cet
élément de $\mathcal{O}_z$ en un élément inversible de $\mathcal{O}_x$,
contrairement à l'hypothèse, d'où le lemme.
\end{proof}

\begin{corollary}[8.11.2]
\label{II.8.11.2-fr}
Soient $U$ un préschéma intègre, $V$ un préschéma normal~; alors tout morphisme
$h:U\to V$ qui est universellement fermé, birationnel et radiciel est un
isomorphisme.
\end{corollary}

\begin{proof}
Si $h=(h_0,\lambda)$, il résulte des hypothèses que $h_0$ est injectif et
fermé, et que $h_0(U)$

\oldpage[II]{190}
est dense dans $V$, donc $h_0$ est un \emph{homéomorphisme} de $U$ sur $V$.
Pour démontrer le corollaire, il suffira de voir que
$\lambda:\mathcal{O}_V\to(h_0)_*(\mathcal{O}_U)$ est un isomorphisme~: on
pourra appliquer alors (\hyperref[II.8.11.1-fr]{8.11.1}) qui prouvera que
l'application (\hyperref[II.8.11.1.1-fr]{8.11.1.1}) est bijective (les fibres
$h_0^{-1}(x)$ étant chacune réduite à un seul point)~; donc $h$ sera un
isomorphisme. La question étant évidemment locale sur $V$, on peut supposer
que $V=\operatorname{Spec}(A)$ est affine d'anneau intègre et intégralement
clos (\hyperref[II.8.8.6.1-fr]{8.8.6.1})~; $h$ correspond alors
(\hyperref[I.2.2.4-fr]{I, 2.2.4}) à un homomorphisme
$\varphi:A\to\Gamma(U,\mathcal{O}_U)$ et tout revient à voir que $\varphi$
est un isomorphisme. Or, si $K$ est le corps des fractions de $A$,
$\Gamma(U,\mathcal{O}_U)$ a par hypothèse $K$ pour corps des fractions et $A$
est un sous-anneau de $\Gamma(U,\mathcal{O}_U)$, $\varphi$ étant l'injection
canonique (\hyperref[I.8.2.7-fr]{I, 8.2.7}). Comme le morphisme $h$ vérifie
les hypothèses de (\hyperref[II.7.3.11-fr]{7.3.11}),
$\Gamma(U,\mathcal{O}_U)$ est un sous-anneau de la fermeture intégrale de $A$
dans $K$, donc est identique à $A$ par hypothèse.
\end{proof}

\begin{remark}[8.11.3]
\label{II.8.11.3-fr}
On verra au chapitre III (\hyperref[III.4.4.11-fr]{III, 4.4.11}) que lorsque
$V$ est un préschéma \emph{localement noethérien}, tout morphisme $h:U\to V$
qui est propre et quasi-fini (en particulier tout morphisme vérifiant les
hypothèses de (\hyperref[II.8.11.2-fr]{8.11.2})) est nécessairement
\emph{fini}. La conclusion de (\hyperref[II.8.11.2-fr]{8.11.2}) résulte donc
dans ce cas de (\hyperref[II.6.1.15-fr]{6.1.15}).
\end{remark}

\begin{env}[8.11.4]
\label{II.8.11.4-fr}
Nous allons maintenant voir que, dans le critère de Grauert
(\hyperref[II.8.9.1-fr]{8.9.1}), on peut souvent affirmer que le préschéma
$C$ et la «~contraction~» $q$ sont déterminés de façon essentiellement
unique.
\end{env}

\begin{lemma}[8.11.5]
\label{II.8.11.5-fr}
Soient $Y$ un préschéma, $p:X\to Y$ un morphisme propre, $\mathcal{L}$ un
$\mathcal{O}_X$-Module inversible $p$-ample, $C$ un $Y$-préschéma,
$\varepsilon:Y\to C$ une $Y$-section,
$q:V=\mathbf{V}(\mathcal{L})\to C$ un $Y$-morphisme, tels que le diagramme
(\hyperref[II.8.9.1.1-fr]{8.9.1.1}) soit commutatif. On suppose en outre que,
si $p=(p_0,\theta)$,
$\theta:\mathcal{O}_Y\to p_*(\mathcal{O}_X)$ est un isomorphisme. Soient alors
$\mathcal{S}'=\bigoplus_{n\geqslant0}p_*(\mathcal{L}^{\otimes n})$,
$C'=\operatorname{Spec}(\mathcal{S}')$, et
$q':\mathbf{V}(\mathcal{L})\to C'$ le $Y$-morphisme canonique
(\hyperref[II.8.8.5-fr]{8.8.5}). Il existe un $Y$-morphisme $u:C'\to C$ et
un seul tel que $q=u\circ q'$.
\end{lemma}

\begin{proof}
L'hypothèse sur $\theta$ entraîne en particulier que $p$ est surjectif~;
comme en vertu de (\hyperref[II.8.8.4-fr]{8.8.4}) la restriction de $q'$ à
$\mathbf{V}(\mathcal{L})-j(X)$ est un \emph{isomorphisme} sur
$C'-\varepsilon'(Y)$ ($\varepsilon'$ étant la section sommet de $C'$), il
résulte de (\hyperref[II.8.8.4-fr]{8.8.4}) que $q'$ est \emph{propre et
surjectif}~; en outre, en vertu de (\hyperref[II.8.8.6-fr]{8.8.6}), si on pose
$q'=(q'_0,\tau)$,
$\tau:\mathcal{O}_{C'}\to q'_*(\mathcal{O}_V)$ est un isomorphisme. On est
donc dans les conditions d'application du lemme
(\hyperref[II.8.11.1-fr]{8.11.1}), et le lemme sera démontré si l'on prouve
que $q$ est constant sur toute fibre $q'{}^{-1}(z')$, où $z'\in C'$. Or, la
condition est trivialement vérifiée pour $z'\notin\varepsilon'(Y)$. D'autre
part, si $z'\in\varepsilon'(Y)$, il existe un $y\in Y$ et un seul tel que
$z'=\varepsilon'(y)$~; et en vertu de la commutativité de
(\hyperref[II.8.8.5.2-fr]{8.8.5.2}) et le fait que $q'$ applique
$\mathbf{V}(\mathcal{L})-j(X)$ dans $C'-\varepsilon'(Y)$,
$q'{}^{-1}(z')=j(p^{-1}(y))$~; la commutativité du diagramme
(\hyperref[II.8.9.1.1-fr]{8.9.1.1}) établit donc notre assertion.
\end{proof}

\begin{corollary}[8.11.6]
\label{II.8.11.6-fr}
Sous les hypothèses de (\hyperref[II.8.11.5-fr]{8.11.5}), supposons outre que
$q$ soit propre et que la restriction de $q$ à
$\mathbf{V}(\mathcal{L})-j(X)$ soit un isomorphisme sur
$C-\varepsilon(Y)$. Alors le morphisme $u$ est universellement fermé,
surjectif et radiciel, et sa restriction à $C'-\varepsilon'(Y)$ est un
isomorphisme sur $C-\varepsilon(Y)$.
\end{corollary}

\begin{proof}
Comme $q'$ est un isomorphisme de $\mathbf{V}(\mathcal{L})-j(X)$ sur
$C'-\varepsilon'(Y)$ (\hyperref[II.8.8.4-fr]{8.8.4}), la dernière assertion
résulte aussitôt de $q=u\circ q'$. En outre, la commutativité des diagrammes

\oldpage[II]{191}
(\hyperref[II.8.8.5.2-fr]{8.8.5.2}) et
(\hyperref[II.8.9.1.1-fr]{8.9.1.1}) montre que la restriction de $u$ au
sous-préschéma fermé $\varepsilon'(Y)$ de $C'$ est un isomorphisme sur le
sous-préschéma fermé $\varepsilon(Y)$ de $C$, d'où résulte aussitôt que pour
tout $z'\in\varepsilon'(Y)$, si $z=u(z')$, $u$ définit un isomorphisme de
$k(z)$ sur $k(z')$. Ces remarques prouvent que $u$ est bijectif et radiciel~;
en outre, si $\psi$ et $\psi'$ sont les morphismes structuraux $C\to Y$,
$C'\to Y$, on a $\psi'=\psi\circ u$ et comme $\psi'$ est séparé
(\hyperref[II.1.2.4-fr]{1.2.4}), il en est de même de $u$
(\hyperref[I.5.5.1-fr]{I, 5.5.1, (v)}). On a vu d'autre part dans la
démonstration du lemme (\hyperref[II.8.11.5-fr]{8.11.5}) que $q'$ est
surjectif~; comme $q=u\circ q'$ est propre, on conclut finalement de
(\hyperref[II.5.4.3-fr]{5.4.3}) et (\hyperref[II.5.4.9-fr]{5.4.9}) que $u$
est universellement fermé.
\end{proof}

\begin{proposition}[8.11.7]
\label{II.8.11.7-fr}
Soient $Y$ un préschéma, $X$ un préschéma intègre, $p:X\to Y$ un morphisme
propre, $\mathcal{L}$ un $\mathcal{O}_X$-Module inversible $p$-ample, $C$ un
$Y$-préschéma normal, $\varepsilon:Y\to C$ une $Y$-section,
$q:V=\mathbf{V}(\mathcal{L})\to C$ un $Y$-morphisme, tels que le diagramme
(\hyperref[II.8.9.1.1-fr]{8.9.1.1}) soit commutatif. On suppose en outre que,
si $p=(p_0,\theta)$, $\theta:\mathcal{O}_Y\to p_*(\mathcal{O}_X)$ soit un
isomorphisme. Soient alors
$\mathcal{S}'=\bigoplus_{n\geqslant0}p_*(\mathcal{L}^{\otimes n})$,
$C'=\operatorname{Spec}(\mathcal{S}')$, et
$q':\mathbf{V}(\mathcal{L})\to C'$ le $Y$-morphisme canonique
(\hyperref[II.8.8.5-fr]{8.8.5}). Alors l'unique $Y$-morphisme $u:C'\to C$ tel
que $q=u\circ q'$ est un isomorphisme.
\end{proposition}

\begin{proof}
Il résulte de (\hyperref[II.8.8.6-fr]{8.8.6}) que $C'$ est intègre~; comme
$u$ est un homéomorphisme des espaces sous-jacents $C'\to C$ ($u$ étant
bijectif et fermé d'après (\hyperref[II.8.11.6-fr]{8.11.6})), $C$ est
irréductible, donc intègre, et puisque la restriction de $u$ à un ouvert non
vide de $C'$ est un isomorphisme sur un ouvert de $C$, $u$ est birationnel.
Puisque $C$ est supposé normal, il suffit d'appliquer
(\hyperref[II.8.11.2-fr]{8.11.2}) pour obtenir la conclusion.
\end{proof}

\begin{remark}[8.11.8]
\label{II.8.11.8-fr}
\begin{enumerate}
  \item[\rm{(i)}] Notons que l'hypothèse que $C$ est normal implique qu'il en
  est de même de $X$. En effet, $C'=\operatorname{Spec}(\mathcal{S}')$ est
  alors normal, étant isomorphe à $C$, et intègre en vertu de
  (\hyperref[II.8.8.6-fr]{8.8.6})~; on en conclut que
  $\operatorname{Proj}(\mathcal{S}')$ est \emph{normal}. En effet, la question
  est locale sur $Y$~; si $Y$ est affine, $\mathcal{S}'=\widetilde{S'}$,
  l'anneau $S'=\Gamma(C',\mathcal{O}_{C'})$ est intègre et intégralement clos
  (\hyperref[II.8.8.6.1-fr]{8.8.6.1}), donc, pour tout élément homogène
  $f\in S'_+$, l'anneau gradué $S'_f$ est intègre et intégralement clos
  ([13], t.~I, p.~257 et 261), et par suite aussi l'anneau $S'_{(f)}$ de ses
  termes de degré 0, car l'intersection de $S'_f$ et du corps des fractions de
  $S'_{(f)}$ est égale à $S'_{(f)}$~; ce qui prouve notre assertion
  (\hyperref[II.6.3.4-fr]{6.3.4}). Enfin, comme $X$ est isomorphe à un
  sous-préschéma ouvert de $\operatorname{Proj}(\mathcal{S}')$
  (\hyperref[II.8.8.1-fr]{8.8.1}), $X$ est bien normal. On peut donc exprimer
  la prop. (\hyperref[II.8.11.7-fr]{8.11.7}) sous la forme suivante~:
  \emph{Si $X$ est intègre et normal,
  $p=(p_0,\theta):X\to Y$ un morphisme propre tel que
  $\theta:\mathcal{O}_Y\to p_*(\mathcal{O}_X)$ soit un isomorphisme, alors,
  pour tout $\mathcal{O}_X$-Module $p$-ample $\mathcal{L}$, il existe une
  façon et une seule de contracter la section nulle de
  $V=\mathbf{V}(\mathcal{L})$ de façon à obtenir un $Y$-schéma normal $C$ et
  un $Y$-morphisme propre $q:V\to C$.}
  \item[\rm{(ii)}] Lorsque $p$ est propre, l'hypothèse
  $p_*(\mathcal{O}_X)=\mathcal{O}_Y$ peut être considérée comme une hypothèse
  auxiliaire, ne constituant pas une vraie restriction de la généralité du
  résultat. En effet, si elle n'est pas vérifiée, il suffit de remplacer $Y$
  par le $Y$-schéma
  $Y'=\operatorname{Spec}(p_*(\mathcal{O}_X))$, et de considérer $X$ comme un
  $Y'$-schéma. Nous reviendrons sur cette méthode générale au chapitre III,
  \S~4.
\end{enumerate}
\end{remark}
\subsection{Faisceaux quasi-cohérents sur les cônes projetants}
\label{subsection:II.8.12-fr}

\begin{env}[8.12.1]
\label{II.8.12.1-fr}
Reprenons les hypothèses et notations de
(\hyperref[II.8.3.1-fr]{8.3.1}). Soit $\mathcal{M}$ un
$\mathcal{S}$-Module gradué quasi-cohérent~; pour éviter toute confusion,
nous désignerons par $\widetilde{\mathcal{M}}$ le $\mathcal{O}_C$-Module

\oldpage[II]{192}
quasi-cohérent associé à $\mathcal{M}$
(\hyperref[II.1.4.3-fr]{1.4.3}) lorsque $\mathcal{M}$ est considéré comme
$\mathcal{S}$-Module \emph{non gradué}, et par
$\operatorname{Proj}_0(\mathcal{M})$ le $\mathcal{O}_X$-Module
quasi-cohérent associé à $\mathcal{M}$, $\mathcal{M}$ étant cette fois
considéré comme $\mathcal{S}$-Module gradué (autrement dit le
$\mathcal{O}_X$-Module noté $\widetilde{\mathcal{M}}$ dans
(\hyperref[II.3.2.2-fr]{3.2.2})). Nous poserons de plus
\begin{equation}
\label{II.8.12.1.1-fr}
  \mathcal{M}_X=\operatorname{Proj}(\mathcal{M})
  =\bigoplus_{n\in\mathbf{Z}}\operatorname{Proj}_0(\mathcal{M}(n))~;
\tag{8.12.1.1}
\end{equation}
la $\mathcal{O}_X$-Algèbre graduée quasi-cohérente $\mathcal{S}_X$ étant
définie par (\hyperref[II.8.6.1.1-fr]{8.6.1.1}),
$\operatorname{Proj}(\mathcal{M})$ est muni d'une structure de
$\mathcal{S}_X$-Module gradué (quasi-cohérent), au moyen des homomorphismes
canoniques (\hyperref[II.3.2.6.1-fr]{3.2.6.1})
\begin{equation}
\label{II.8.12.1.2-fr}
  \mathcal{O}_X(m)\otimes_{\mathcal{O}_X}
  \operatorname{Proj}_0(\mathcal{M}(n))
  \to \operatorname{Proj}_0(\mathcal{S}(m)\otimes_{\mathcal{S}}\mathcal{M}(n))
  \to \operatorname{Proj}_0(\mathcal{M}(m+n))
\tag{8.12.1.2}
\end{equation}
la vérification des axiomes des Modules se faisant à l'aide du diagramme
commutatif (\hyperref[II.2.5.11.4-fr]{2.5.11.4}).

Si $Y=\operatorname{Spec}(A)$ est affine, $\mathcal{S}=\widetilde{S}$ et
$\mathcal{M}=\widetilde{M}$, où $S$ est une $A$-algèbre graduée et $M$ un
$S$-module gradué, alors, pour tout élément homogène $f\in S_+$, on a
\begin{equation}
\label{II.8.12.1.3-fr}
  \Gamma(X_f,\operatorname{Proj}(\widetilde{M}))=M_f
\tag{8.12.1.3}
\end{equation}
en vertu des définitions et de (\hyperref[II.8.2.9.1-fr]{8.2.9.1}).

Considérons maintenant le $\widehat{\mathcal{S}}$-Module gradué
quasi-cohérent
\begin{equation}
\label{II.8.12.1.4-fr}
  \widehat{\mathcal{M}}
  =\mathcal{M}\otimes_{\mathcal{S}}\widehat{\mathcal{S}}
\tag{8.12.1.4}
\end{equation}
($\widehat{\mathcal{S}}$ étant défini par
(\hyperref[II.8.3.1.1-fr]{8.3.1.1}))~; on en déduit un
$\mathcal{O}_{\widehat{C}}$-Module gradué quasi-cohérent
$\operatorname{Proj}_0(\widehat{\mathcal{M}})$, que nous noterons aussi
\begin{equation}
\label{II.8.12.1.5-fr}
  \mathcal{M}^{\square}=\operatorname{Proj}_0(\widehat{\mathcal{M}}).
\tag{8.12.1.5}
\end{equation}
Il est clair (\hyperref[II.3.2.4-fr]{3.2.4}) que
$\mathcal{M}^{\square}$ est un foncteur additif exact en $\mathcal{M}$,
commutant aux sommes directes et aux limites inductives.
\end{env}

\begin{proposition}[8.12.2]
\label{II.8.12.2-fr}
Avec les notations de (\hyperref[II.8.3.2-fr]{8.3.2}), on a des
isomorphismes canoniques fonctoriels
\begin{equation}
\label{II.8.12.2.1-fr}
  i^*(\mathcal{M}^{\square})\xrightarrow{\sim}\widetilde{\mathcal{M}},
  \qquad
  j^*(\mathcal{M}^{\square})\xrightarrow{\sim}
  \operatorname{Proj}_0(\mathcal{M}).
\tag{8.12.2.1}
\end{equation}
\end{proposition}

\begin{proof}
En effet, $i^*(\mathcal{M}^{\square})$ s'identifie canoniquement à
$\left(\widehat{\mathcal{M}}/(z-1)\widehat{\mathcal{M}}\right)^{\sim}$
sur
$\operatorname{Spec}(\widehat{\mathcal{S}}/(z-1)\widehat{\mathcal{S}})$
en vertu de (\hyperref[II.3.2.3-fr]{3.2.3})~; le premier des isomorphismes
canoniques (\hyperref[II.8.12.2.1-fr]{8.12.2.1}) se déduit alors aussitôt
(\hyperref[II.1.4.1-fr]{1.4.1}) de l'isomorphisme canonique
$\widehat{\mathcal{M}}/(z-1)\widehat{\mathcal{M}}
\xrightarrow{\sim}\mathcal{M}$. D'autre part, l'immersion canonique
$j:X\to\widehat{C}$ correspond à l'homomorphisme canonique
$\widehat{\mathcal{S}}\to\mathcal{S}$ de noyau
$z\widehat{\mathcal{S}}$ (\hyperref[II.8.3.2-fr]{8.3.2})~; le second
homomorphisme (\hyperref[II.8.12.2.1-fr]{8.12.2.1}) est le cas particulier
de l'homomorphisme canonique (\hyperref[II.3.5.2-fr]{3.5.2, (ii)}), du fait
que l'on a ici
$\widehat{\mathcal{M}}\otimes_{\widehat{\mathcal{S}}}\mathcal{S}
=\mathcal{M}$~; pour vérifier que c'est un isomorphisme, on peut se borner au
cas où $Y=\operatorname{Spec}(A)$ est affine,
$\mathcal{S}=\widetilde{S}$, $\mathcal{M}=\widetilde{M}$~; en se reportant à
(\hyperref[II.2.8.8-fr]{2.8.8}), la vérification que pour tout $f$ homogène
dans $S_+$, l'homomorphisme précédent, restreint à $X_f$, se réduit à un
isomorphisme, est alors immédiate.
\end{proof}

\oldpage[II]{193}
Par abus de langage, on dira encore, en raison de l'existence du premier
isomorphisme (\hyperref[II.8.12.2.1-fr]{8.12.2.1}), que
$\mathcal{M}^{\square}$ est la \emph{fermeture projective} du
$\mathcal{O}_C$-Module $\widetilde{\mathcal{M}}$ (étant sous-entendu que dans
la donnée du $\mathcal{O}_C$-Module $\widetilde{\mathcal{M}}$ figure la
graduation du $\mathcal{S}$-Module $\mathcal{M}$).

\begin{env}[8.12.3]
\label{II.8.12.3-fr}
Avec les notations de (\hyperref[II.8.3.5-fr]{8.3.5}), on a un
homomorphisme canonique fonctoriel
\begin{equation}
\label{II.8.12.3.1-fr}
  p^*(\operatorname{Proj}_0(\mathcal{M}))
  \to\mathcal{M}^{\square}|\widehat{E}.
\tag{8.12.3.1}
\end{equation}
En effet, c'est un cas particulier de l'homomorphisme $u^{\sharp}$ défini de
façon générale dans (\hyperref[II.3.5.6-fr]{3.5.6}). Si
$Y=\operatorname{Spec}(A)$ est affine, $\mathcal{S}=\widetilde{S}$,
$\mathcal{M}=\widetilde{M}$, on voit, en se reportant à
(\hyperref[II.2.8.8-fr]{2.8.8}), que la restriction de
(\hyperref[II.8.12.3.1-fr]{8.12.3.1}) à
$p^{-1}(X_f)=\widehat{C}_f$ (pour un $f$ homogène dans $S_+$) correspond à
l'homomorphisme canonique
\begin{equation}
\label{II.8.12.3.2-fr}
  M_{(f)}\otimes_{S_{(f)}}S_f^{\leq}\to M_f^{\leq}
\tag{8.12.3.2}
\end{equation}
compte tenu de (\hyperref[II.8.2.3.2-fr]{8.2.3.2}) et
(\hyperref[II.8.2.5.2-fr]{8.2.5.2}).
\end{env}

\begin{env}[8.12.4]
\label{II.8.12.4-fr}
Plaçons-nous dans les hypothèses de (\hyperref[II.8.5.1-fr]{8.5.1}), dont
nous conservons les notations. Il résulte de
(\hyperref[II.1.5.6-fr]{1.5.6}) que pour tout $\mathcal{S}$-Module gradué
quasi-cohérent $\mathcal{M}$, on a d'une part un isomorphisme canonique
\begin{equation}
\label{II.8.12.4.1-fr}
  \Phi^*(\widetilde{\mathcal{M}})
  \xrightarrow{\sim}
  \left(q^*(\mathcal{M})\otimes_{q^*(\mathcal{S})}\mathcal{S}'\right)^{\sim}
\tag{8.12.4.1}
\end{equation}
de $\mathcal{O}_{C'}$-Modules~; et d'autre part
(\hyperref[II.3.5.6-fr]{3.5.6}) entraîne l'existence d'un
$\operatorname{Proj}(\varphi)$-morphisme canonique
\begin{equation}
\label{II.8.12.4.2-fr}
  \operatorname{Proj}_0(\mathcal{M})\to
  \left(\operatorname{Proj}_0
  (q^*(\mathcal{M})\otimes_{q^*(\mathcal{S})}\mathcal{S}')\right)
  |G(\varphi)
\tag{8.12.4.2}
\end{equation}
et aussi d'un $\widehat{\Phi}$-morphisme canonique
\begin{equation}
\label{II.8.12.4.3-fr}
  \operatorname{Proj}_0(\widehat{\mathcal{M}})\to
  \left(\operatorname{Proj}_0
  (q^*(\widehat{\mathcal{M}})
  \otimes_{q^*(\widehat{\mathcal{S}})}\widehat{\mathcal{S}'})\right)
  |G(\widehat{\varphi}).
\tag{8.12.4.3}
\end{equation}
\end{env}

\begin{env}[8.12.5]
\label{II.8.12.5-fr}
Considérons maintenant la situation de (\hyperref[II.8.6.1-fr]{8.6.1}),
avec les mêmes notations~; on prend donc dans ce qui précède $Y'=X$, le
morphisme $q:X\to Y$ étant le morphisme structural, et $\varphi$ est le
$q$-morphisme canonique (\hyperref[II.8.6.1.2-fr]{8.6.1.2}). On a alors un
isomorphisme canonique
\begin{equation}
\label{II.8.12.5.1-fr}
  q^*(\mathcal{M})\otimes_{q^*(\mathcal{S})}\mathcal{S}_X^{\geq}
  \xrightarrow{\sim}\mathcal{M}_X^{\geq}
\tag{8.12.5.1}
\end{equation}
en posant
$\mathcal{M}_X^{\geq}=\bigoplus_{n\geq0}\operatorname{Proj}_0(\mathcal{M}(n))$.
On peut en effet se borner au cas où $Y=\operatorname{Spec}(A)$ est affine,
$\mathcal{S}=\widetilde{S}$ et $\mathcal{M}=\widetilde{M}$, et définir
l'isomorphisme (\hyperref[II.8.12.5.1-fr]{8.12.5.1}) dans chacun des ouverts
affines $X_f$ ($f$ homogène dans $S_+$), en vérifiant la compatibilité avec le
passage à un multiple homogène de $f$. Or, la restriction à $X_f$ du premier
membre de (\hyperref[II.8.12.5.1-fr]{8.12.5.1}) est
$\widetilde{M}'=
\left((M\otimes_A S_{(f)})
\otimes_{S\otimes_A S_{(f)}}S_f^{\geq}\right)^{\sim}$ par
(\hyperref[II.8.6.2.1-fr]{8.6.2.1})~; comme on a un isomorphisme canonique de
$M\otimes_A S_{(f)}$ sur $M\otimes_S(S\otimes_A S_{(f)})$, on en déduit un
isomorphisme de $\widetilde{M}'$ sur
$\left(M\otimes_S S_f^{\geq}\right)^{\sim}$, et ce dernier est
canoniquement isomorphe par (\hyperref[II.8.2.9.1-fr]{8.2.9.1}) à la
restriction à $X_f$ du second membre de
(\hyperref[II.8.12.5.1-fr]{8.12.5.1}), et satisfait aux conditions de
compatibilité requises.

\oldpage[II]{194}
Remplaçant $\mathcal{M}$ par $\widehat{\mathcal{M}}$, $\mathcal{S}$ par
$\widehat{\mathcal{S}}$ et $\mathcal{S}_X$ par
$(\mathcal{S}_X^{\geq})^{\wedge}$ dans le raisonnement précédent, on a de
même un isomorphisme canonique
\begin{equation}
\label{II.8.12.5.2-fr}
  q^*(\widehat{\mathcal{M}})
  \otimes_{q^*(\widehat{\mathcal{S}})}(\mathcal{S}_X^{\geq})^{\wedge}
  \xrightarrow{\sim}(\mathcal{M}_X^{\geq})^{\wedge}.
\tag{8.12.5.2}
\end{equation}
Si l'on se souvient (\hyperref[II.8.6.2-fr]{8.6.2}) que le morphisme
structural $u:\operatorname{Proj}(\mathcal{S}_X^{\geq})\to X$ est un
isomorphisme, on déduit d'abord de ce qui précède que l'on a un
$u$-isomorphisme canonique
\begin{equation}
\label{II.8.12.5.3-fr}
  \operatorname{Proj}_0\mathcal{M}
  \xrightarrow{\sim}\operatorname{Proj}_0(\mathcal{M}_X^{\geq})
\tag{8.12.5.3}
\end{equation}
comme cas particulier de (\hyperref[II.8.12.4.2-fr]{8.12.4.2}). On constate
en effet, avec les notations de la démonstration de
(\hyperref[II.8.6.2-fr]{8.6.2}), que cela revient à voir que
l'homomorphisme canonique
$M_{(f)}\otimes_{S_{(f)}}(S_f^{\geq})^{(d)}\to(M_f^{\geq})^{(d)}$ est un
isomorphisme lorsque $f\in S_d$, ce qui est immédiat.

En second lieu, l'isomorphisme (\hyperref[II.8.12.5.2-fr]{8.12.5.2}) permet,
en appliquant cette fois (\hyperref[II.8.12.4.3-fr]{8.12.4.3}) au morphisme
canonique $r=\operatorname{Proj}(\widehat{\alpha}):\widehat{C}_X\to
\widehat{C}$, d'obtenir un $r$-morphisme canonique
\begin{equation}
\label{II.8.12.5.4-fr}
  \mathcal{M}^{\square}\to(\mathcal{M}_X^{\geq})^{\square}.
\tag{8.12.5.4}
\end{equation}
Rappelons maintenant (\hyperref[II.8.6.2-fr]{8.6.2}) que les restrictions de
$r$ aux cônes épointés $\widehat{E}_X$ et $E_X$ sont des \emph{isomorphismes}
sur $\widehat{E}$ et $E$ respectivement. En outre~:
\end{env}

\begin{proposition}[8.12.6]
\label{II.8.12.6-fr}
Les restrictions à $\widehat{E}_X$ et à $E_X$ du $r$-morphisme canonique
(\hyperref[II.8.12.5.4-fr]{8.12.5.4}) sont des isomorphismes
\begin{equation}
\label{II.8.12.6.1-fr}
  \mathcal{M}^{\square}|\widehat{E}
  \xrightarrow{\sim}(\mathcal{M}_X^{\geq})^{\square}|\widehat{E}_X
\tag{8.12.6.1}
\end{equation}
\begin{equation}
\label{II.8.12.6.2-fr}
  \widetilde{\mathcal{M}}|E
  \xrightarrow{\sim}(\mathcal{M}_X^{\geq})^{\sim}|E_X.
\tag{8.12.6.2}
\end{equation}
\end{proposition}

\begin{proof}
On se ramène au cas où $Y$ est affine comme dans la démonstration de
(\hyperref[II.8.6.2-fr]{8.6.2})~; avec les notations de cette dernière, et
en se ramenant aux définitions (\hyperref[II.2.8.8-fr]{2.8.8}), on doit
montrer que l'homomorphisme canonique
\[
  \widehat{M}_{(f)}\otimes_{\widehat{S}_{(f)}}
  (S_f^{\geq})_{(f/1)}^{\wedge}
  \to(M\otimes_S S_f^{\geq})_{(f/1)}^{\wedge}
\]
est un isomorphisme~; mais en vertu de
(\hyperref[II.8.2.3.2-fr]{8.2.3.2}) et
(\hyperref[II.8.2.5.2-fr]{8.2.5.2}), le premier membre s'identifie
canoniquement à
$M_f^{\leq}\otimes_{S_f^{\leq}}(S_f^{\geq})_{f/1}^{\leq}$, donc à
$M_f^{\leq}$ en vertu de (\hyperref[II.8.2.7.2-fr]{8.2.7.2}), et le second à
$(M_f^{\geq})_{f/1}^{\leq}$, donc aussi à $M_f^{\leq}$ en vertu de
(\hyperref[II.8.2.9.2-fr]{8.2.9.2}), d'où la conclusion en ce qui concerne
(\hyperref[II.8.12.6.1-fr]{8.12.6.1})~; la formule
(\hyperref[II.8.12.6.2-fr]{8.12.6.2}) résulte alors de
(\hyperref[II.8.12.6.1-fr]{8.12.6.1}) et de
(\hyperref[II.8.12.2.1-fr]{8.12.2.1}).
\end{proof}

\begin{corollary}[8.12.7]
\label{II.8.12.7-fr}
Avec les identifications de (\hyperref[II.8.6.3-fr]{8.6.3}), la restriction
de $(\mathcal{M}_X^{\geq})^{\square}$ à $\widehat{E}_X$ s'identifie à
$(\mathcal{M}_X^{\leq})^{\sim}$ et la restriction de
$(\mathcal{M}_X^{\geq})^{\square}$ à $E_X$ s'identifie à
$\widetilde{\mathcal{M}}_X$.
\end{corollary}

\begin{proof}
On est en effet ramené au cas affine, et cela résulte de l'identification de
$(M_f^{\geq})_{f/1}^{\leq}$ à $M_f^{\leq}$ et de
$(M_f^{\geq})_{f/1}$ à $M_f$ (\hyperref[II.8.2.9.2-fr]{8.2.9.2}).
\end{proof}

\begin{proposition}[8.12.8]
\label{II.8.12.8-fr}
Sous les hypothèses de (\hyperref[II.8.6.4-fr]{8.6.4}), l'homomorphisme
canonique (\hyperref[II.8.12.3.1-fr]{8.12.3.1}) est un isomorphisme.
\end{proposition}

\begin{proof}
Compte tenu du fait que $\operatorname{Proj}(\mathcal{S}_X^{\geq})\to X$
est un isomorphisme (\hyperref[II.8.6.2-fr]{8.6.2}) et des

\oldpage[II]{195}
isomorphismes (\hyperref[II.8.12.5.4-fr]{8.12.5.4}) et
(\hyperref[II.8.12.6.1-fr]{8.12.6.1}), on est ramené à démontrer la
proposition correspondante pour l'homomorphisme canonique
$p_X^*(\operatorname{Proj}_0(\mathcal{M}_X^{\geq}))
\to(\mathcal{M}_X^{\geq})^{\square}|\widehat{E}_X$, autrement dit, on est
ramené au cas où $\mathcal{S}_1$ est un $\mathcal{O}_Y$-Module inversible et
où $\mathcal{S}$ est engendrée par $\mathcal{S}_1$. Avec les notations de
(\hyperref[II.8.12.3-fr]{8.12.3}), on a alors, pour un $f\in S_1$,
$S_f^{\leq}=S_{(f)}[1/f]$ et l'homomorphisme canonique
$M_{(f)}\otimes_{S_{(f)}}S_f^{\leq}\to M_f^{\leq}$ est un isomorphisme par
définition de $M_f^{\leq}$.
\end{proof}

\begin{env}[8.12.9]
\label{II.8.12.9-fr}
Considérons maintenant les $\mathcal{S}$-Modules quasi-cohérents
\[
  \mathcal{M}_{[n]}=\bigoplus_{m\geqslant n}\mathcal{M}_m
\]
et (avec les notations de (\hyperref[II.8.7.2-fr]{8.7.2})) le
$\mathcal{S}^{\natural}$-Module gradué quasi-cohérent
\begin{equation}
\label{II.8.12.9.1-fr}
  \mathcal{M}^{\natural}
  =\left(\bigoplus_{n\geqslant0}\mathcal{M}_{[n]}\right)^{\sim}.
\tag{8.12.9.1}
\end{equation}
On a vu (\hyperref[II.8.7.3-fr]{8.7.3}) qu'il existe un $C$-isomorphisme
canonique $h:C_X\xrightarrow{\sim}\operatorname{Proj}(\mathcal{S}^{\natural})$.
En outre~:
\end{env}

\begin{proposition}[8.12.10]
\label{II.8.12.10-fr}
Il existe un $h$-isomorphisme canonique
\begin{equation}
\label{II.8.12.10.1-fr}
  \operatorname{Proj}_0(\mathcal{M}^{\natural})
  \xrightarrow{\sim}\widetilde{\mathcal{M}}_X.
\tag{8.12.10.1}
\end{equation}
\end{proposition}

\begin{proof}
On raisonne comme dans (\hyperref[II.8.7.3-fr]{8.7.3}), en utilisant cette
fois l'existence du di-isomorphisme (\hyperref[II.8.2.9.3-fr]{8.2.9.3}) au
lieu de (\hyperref[II.8.2.7.3-fr]{8.2.7.3}). Nous laissons les détails au
lecteur.
\end{proof}

\subsection{Fermetures projectives de sous-faisceaux et de sous-schémas
fermés}
\label{subsection:II.8.13-fr}

\begin{env}[8.13.1]
\label{II.8.13.1-fr}
Les hypothèses et notations étant celles de
(\hyperref[II.8.12.1-fr]{8.12.1}), considérons un sous-$\mathcal{S}$-Module
quasi-cohérent $\mathcal{N}$ de $\mathcal{M}$, \emph{non nécessairement
gradué}. On peut donc considérer le $\mathcal{O}_C$-Module quasi-cohérent
$\widetilde{\mathcal{N}}$ associé à $\mathcal{N}$, qui est un
sous-$\mathcal{O}_C$-Module de $\widetilde{\mathcal{M}}$. On a vu d'autre
part (\hyperref[II.8.12.2.1-fr]{8.12.2.1}) que
$\widetilde{\mathcal{M}}$ s'identifie à la restriction de
$\mathcal{M}^{\square}$ à $C$. Comme l'injection canonique
$i:C\to\widehat{C}$ est un morphisme affine
(\hyperref[II.8.3.2-fr]{8.3.2}), et \emph{a fortiori} quasi-compact, le
\emph{prolongement canonique} $(\widetilde{\mathcal{N}})^{-}$, plus grand
sous-$\mathcal{O}_{\widehat{C}}$-Module contenu dans
$\mathcal{M}^{\square}$ et induisant $\widetilde{\mathcal{N}}$ sur $C$, est
un $\mathcal{O}_{\widehat{C}}$-Module quasi-cohérent
(\hyperref[I.9.4.2-fr]{I, 9.4.2}). Nous allons en donner une description
explicite à l'aide d'un $\widehat{\mathcal{S}}$-Module gradué.
\end{env}

\begin{env}[8.13.2]
\label{II.8.13.2-fr}
Pour cela, considérons, pour tout entier $n\geqslant0$, l'homomorphisme
$\bigoplus_{i\leqslant n}\mathcal{M}_i\to\mathcal{M}$ qui, pour tout ouvert
$U$ de $Y$, fait correspondre à la famille
\[
  (s_i)\in\bigoplus_{i\leqslant n}\Gamma(U,\mathcal{M}_i)
\]
la section $\sum_i s_i\in\Gamma(U,\mathcal{M})$. Désignons par
$\mathcal{N}'_n$ l'image réciproque de $\mathcal{N}$ par cet homomorphisme,
qui est un sous-$\mathcal{O}_Y$-Module quasi-cohérent de
$\bigoplus_{i\leqslant n}\mathcal{M}_i$. Considérons maintenant
l'homomorphisme
$\bigoplus_{i\leqslant n}\mathcal{M}_i\to
\widehat{\mathcal{M}}=\mathcal{M}[z]$ qui fait correspondre à $(s_i)$ la
section
$\sum_{i\leqslant n}s_i z^{n-i}\in\Gamma(U,\widehat{\mathcal{M}}_n)$, et
soit $\mathcal{N}_n$ l'image de $\mathcal{N}'_n$ par cet homomorphisme~; on
vérifie aussitôt que
$\overline{\mathcal{N}}=\bigoplus_{n\geqslant0}\mathcal{N}_n$ est un
sous-$\widehat{\mathcal{S}}$-Module gradué (quasi-cohérent) de
$\widehat{\mathcal{M}}$~; on dit que $\overline{\mathcal{N}}$ est déduit de
$\mathcal{N}$ par \emph{homogénéisation}, à l'aide de la «~variable
homogénéisante~» $z$. On notera

\oldpage[II]{196}
que si $\mathcal{N}$ est déjà un sous-$\mathcal{S}$-Module \emph{gradué} de
$\mathcal{M}$, alors $\overline{\mathcal{N}}$ s'identifie à la somme directe
des composantes $\widehat{\mathcal{N}}_n$ de degré $n\geqslant0$ dans
$\widehat{\mathcal{N}}=\mathcal{N}[z]$.
\end{env}

\begin{proposition}[8.13.3]
\label{II.8.13.3-fr}
Le $\mathcal{O}_{\widehat{C}}$-Module
$\operatorname{Proj}_0(\overline{\mathcal{N}})$ est le prolongement canonique
$(\widetilde{\mathcal{N}})^{-}$ de $\widetilde{\mathcal{N}}$ à
$\widehat{C}$.
\end{proposition}

\begin{proof}
La question est locale sur $Y$ et $\widehat{C}$ en vertu de la définition du
prolongement canonique (\hyperref[I.9.4.1-fr]{I, 9.4.1}). On peut donc déjà
supposer que $Y=\operatorname{Spec}(A)$ est affine, avec
$\mathcal{S}=\widetilde{S}$, $\mathcal{M}=\widetilde{M}$ et
$\mathcal{N}=\widetilde{N}$, où $N$ est un sous-$S$-module non
nécessairement gradué de $M$. En outre
(\hyperref[II.8.3.2.6-fr]{8.3.2.6}) $\widehat{C}$ est réunion des ouverts
affines $\widehat{C}_z=C$ et
$\widehat{C}_f=\operatorname{Spec}(S_f^{\leq})$ ($f$ homogène dans $S_+$).
Il suffira donc de montrer que~: $1^\circ$ la restriction de
$\operatorname{Proj}_0(\overline{\mathcal{N}})$ à $C$ est
$\widetilde{\mathcal{N}}$~; $2^\circ$ la restriction de
$\operatorname{Proj}_0(\overline{\mathcal{N}})$ à chaque $\widehat{C}_f$ est
le prolongement canonique de la restriction de $\mathcal{N}$ à
$C\cap\widehat{C}_f=\operatorname{Spec}(S_f)$
(\hyperref[II.8.3.2.6-fr]{8.3.2.6}). Pour le premier point, notons que
$\operatorname{Proj}_0(\overline{\mathcal{N}})|C$ s'identifie à
$(\overline{N}_{(z)})^{\sim}$
(\hyperref[II.8.3.2.4-fr]{8.3.2.4})~; or $\overline{N}_{(z)}$ s'identifie
canoniquement (\hyperref[II.2.2.5-fr]{2.2.5}) à l'image de
$\overline{N}$ dans $\widehat{M}/(z-1)\widehat{M}$, et par l'isomorphisme
canonique de ce dernier sur $M$ (\hyperref[II.8.2.5-fr]{8.2.5}), cette image
s'identifie à $N$, en vertu de la définition de $\overline{N}$ donnée dans
(\hyperref[II.8.13.2-fr]{8.13.2}).

Pour prouver le second point, notons que l'injection
$i:C\cap\widehat{C}_f\to\widehat{C}_f$ correspond à l'injection canonique
$S_f^{\leq}\to S_f$ (\hyperref[II.8.3.2.6-fr]{8.3.2.6})~; d'autre part, on
a $\Gamma(\widehat{C}_f,\mathcal{M}^{\square})=M_f^{\leq}$,
$\Gamma(\widehat{C}_f,i_*(\widetilde{\mathcal{N}}))=N_f$ et (par
(\hyperref[II.8.12.2.1-fr]{8.12.2.1}))
$\Gamma(\widehat{C}_f,i_*(i^*(\mathcal{M}^{\square})))=M_f$. Compte tenu de
(\hyperref[I.9.4.2-fr]{I, 9.4.2}), on est donc ramené à montrer que
$\overline{N}_{(f)}\subset\widehat{M}_{(f)}=M_f^{\leq}$ s'identifie
canoniquement à l'image réciproque de $N_f$, par l'injection canonique
$M_f^{\leq}\to M_f$. En effet, soit $d=\deg(f)>0$, et supposons qu'un
élément $(\sum_{k\leqslant md}x_k)/f^m$ de $M_f$ (avec $x_k\in M_k$) soit de
la forme $y/f^m$ avec $y\in N$. En multipliant $y$ et les $x_k$ par un même
$f^h$ convenable, on peut déjà supposer que $\sum_{k\leqslant md}x_k=y$.
Mais dans l'identification (\hyperref[II.8.2.5.2-fr]{8.2.5.2}),
$(\sum_{k\leqslant md}x_k)/f^m$ correspond à
$\sum_{k\leqslant md}x_kz^{md-k}/f^m$, et cet élément appartient bien à
$\overline{N}_{(f)}$, puisque $\sum_{k\leqslant md}x_k\in N$~; la réciproque
est évidente.
\end{proof}

\begin{remark}[8.13.4]
\label{II.8.13.4-fr}
\begin{enumerate}
  \item[\rm{(i)}] Le cas d'application le plus important de
  (\hyperref[II.8.13.3-fr]{8.13.3}) est celui où
  $\mathcal{M}=\mathcal{S}$, $\widetilde{\mathcal{N}}$ étant donc un faisceau
  quasi-cohérent \emph{arbitraire} d'idéaux $\mathcal{J}$ de
  $\mathcal{O}_C$ (\hyperref[II.1.4.3-fr]{1.4.3}), correspondant
  biunivoquement à un \emph{sous-préschéma fermé} $Z$ de $C$. Alors le
  prolongement canonique $\overline{\mathcal{J}}$ de $\mathcal{J}$ est le
  faisceau quasi-cohérent d'idéaux de $\mathcal{O}_{\widehat{C}}$ qui définit
  l'\emph{adhérence} $\overline{Z}$ de $Z$ dans $\widehat{C}$
  (\hyperref[I.9.5.10-fr]{I, 9.5.10})~; la prop.
  (\hyperref[II.8.13.3-fr]{8.13.3}) donne un moyen canonique de définir
  $\overline{Z}$ à l'aide d'un Idéal gradué dans
  $\widehat{\mathcal{S}}=\mathcal{S}[z]$.
  \item[\rm{(ii)}] Supposons pour simplifier que $Y$ soit affine, et gardons
  les notations de la démonstration de
  (\hyperref[II.8.13.3-fr]{8.13.3}). Pour tout $x\in N$ non nul, soit $d(x)$
  le plus grand des degrés des composantes homogènes $x_i$ de $x$ dans $M$~;
  $\overline{N}$ est par définition le sous-module de $\widehat{M}$ formé de
  $0$ et des éléments de la forme
  $h(x,k)=z^k\sum_{i\leqslant d(x)}x_i z^{d(x)-i}$ ($k$ entier
  $\geqslant0$)~; il est donc engendré, en tant que module sur
  $\widehat{S}=S[z]$, par les éléments de la forme
  \[
    h(x,0)=\sum_{i\leqslant d(x)}x_i z^{d(x)-i}.
  \]

  \oldpage[II]{197}
  On dit que $h(x,0)$ est déduit de $x$ par \emph{homogénéisation} à l'aide
  de la «~variable homogénéisante~» $z$. Mais comme $h(x,0)$ ne dépend pas
  additivement de $x$ (ni \emph{a fortiori} $S$-linéairement), \emph{on se
  gardera de croire} (même lorsque $M=S$) que les $h(x,0)$ parcourent un
  \emph{système de générateurs} du $\widehat{S}$-module gradué
  $\overline{N}$ lorsqu'on fait parcourir à $x$ un \emph{système de
  générateurs} du $S$-module $N$. Il en est cependant bien ainsi dans le cas
  (seul considéré dans la géométrie algébrique élémentaire) où $N$ est un
  $S$-module \emph{libre monogène}, car si $t$ est une base de $N$, $h(t,0)$
  engendre le $\widehat{S}$-module $\overline{N}$.
\end{enumerate}
\end{remark}

\subsection{Compléments sur les faisceaux associés aux
$\mathcal{S}$-Modules gradués}
\label{subsection:II.8.14-fr}

\begin{env}[8.14.1]
\label{II.8.14.1-fr}
Soient $Y$ un préschéma, $\mathcal{S}$ une $\mathcal{O}_Y$-Algèbre graduée
quasi-cohérente à \emph{degrés positifs},
$X=\operatorname{Proj}(\mathcal{S})$, $q:X\to Y$ le morphisme structural
(séparé en vertu de (\hyperref[II.3.1.3-fr]{3.1.3})). Conservons les
notations de (\hyperref[II.8.12.1-fr]{8.12.1}), de sorte que nous avons
défini un foncteur
$\mathcal{M}_X=\operatorname{Proj}(\mathcal{M})$ en $\mathcal{M}$, de la
catégorie des $\mathcal{S}$-Modules gradués quasi-cohérents dans celle des
$\mathcal{S}_X$-Modules gradués quasi-cohérents~; il est clair en outre
(\hyperref[II.3.2.4-fr]{3.2.4}) que c'est un foncteur \emph{additif exact},
commutant aux limites inductives.

Notons en outre qu'il résulte aussitôt de la définition
(\hyperref[II.8.12.1.1-fr]{8.12.1.1}) que l'on a
\begin{equation}
\label{II.8.14.1.1-fr}
  \operatorname{Proj}(\mathcal{M}(n))
  =(\operatorname{Proj}(\mathcal{M}))(n)
  \qquad\text{pour tout }n\in\mathbf{Z}.
\tag{8.14.1.1}
\end{equation}
\end{env}

\begin{env}[8.14.2]
\label{II.8.14.2-fr}
Nous allons d'abord étendre aux $\mathcal{S}_X$-Modules de la forme
$\operatorname{Proj}(\mathcal{M})$ les homomorphismes canoniques $\lambda$
et $\mu$ définis dans (\hyperref[II.3.2.6-fr]{3.2.6}). Pour cela, notons
que, quels que soient $m\in\mathbf{Z}$ et $n\in\mathbf{Z}$, on a, en vertu
de (\hyperref[II.2.1.2.1-fr]{2.1.2.1}), un homomorphisme canonique de
$\mathcal{O}_X$-Modules
\begin{equation}
\label{II.8.14.2.1-fr}
  \lambda_{mn}:\operatorname{Proj}_0(\mathcal{M}(m))
  \otimes_{\mathcal{O}_X}\operatorname{Proj}_0(\mathcal{N}(n))
  \to\operatorname{Proj}_0
  ((\mathcal{M}\otimes_{\mathcal{S}}\mathcal{N})(m+n))
\tag{8.14.2.1}
\end{equation}
et de même, en vertu de (\hyperref[II.2.1.2.2-fr]{2.1.2.2}), un
homomorphisme canonique de $\mathcal{O}_X$-Modules
\begin{equation}
\label{II.8.14.2.2-fr}
  \mu_{mn}:\operatorname{Proj}_0
  ((\mathcal{H}om_{\mathcal{S}}(\mathcal{M},\mathcal{N}))(n-m))
  \to\mathcal{H}om_{\mathcal{O}_X}
  (\operatorname{Proj}_0(\mathcal{M}(m)),
   \operatorname{Proj}_0(\mathcal{N}(n)))
\tag{8.14.2.2}
\end{equation}
quels que soient les $\mathcal{S}$-Modules gradués quasi-cohérents
$\mathcal{M}$, $\mathcal{N}$. On en déduit un homomorphisme
\[
  \mu_k:\operatorname{Proj}_0
  ((\mathcal{H}om_{\mathcal{S}}(\mathcal{M},\mathcal{N}))(k))
  \to
  (\mathcal{H}om_{\mathcal{S}_X}
  (\operatorname{Proj}(\mathcal{M}),\operatorname{Proj}(\mathcal{N})))_k
\]
obtenu en faisant correspondre à tout
$u\in\Gamma(U,\operatorname{Proj}_0
((\mathcal{H}om_{\mathcal{S}}(\mathcal{M},\mathcal{N}))(k)))$
l'homomorphisme $\mu_k(u)$, de degré $k$, de $\mathbf{Z}$-modules gradués
$\Gamma(U,\operatorname{Proj}(\mathcal{M}))\to
\Gamma(U,\operatorname{Proj}(\mathcal{N}))$ ($U$ ouvert de $X$) qui dans
chaque $\Gamma(U,\operatorname{Proj}_0(\mathcal{M}(m)))$ coïncide avec
$\mu_{m,m+k}(u)$~; en outre, en revenant à la définition des $\mu_{mn}$
(\hyperref[II.2.5.12.1-fr]{2.5.12.1}), on voit aussitôt que $\mu_k(u)$ est
en fait un homomorphisme de degré $k$ de
$\Gamma(U,\mathcal{S}_X)$-modules gradués, et en outre que les $\mu_k$
définissent un homomorphisme de $\mathcal{S}_X$-Modules gradués
\begin{equation}
\label{II.8.14.2.3-fr}
  \operatorname{Proj}
  (\mathcal{H}om_{\mathcal{S}}(\mathcal{M},\mathcal{N}))
  \to\mathcal{H}om_{\mathcal{S}_X}
  (\operatorname{Proj}(\mathcal{M}),\operatorname{Proj}(\mathcal{N})).
\tag{8.14.2.3}
\end{equation}

De même, compte tenu du diagramme d'associativité
(\hyperref[II.2.5.11.4-fr]{2.5.11.4}), les homomorphismes
(\hyperref[II.8.14.2.1-fr]{8.14.2.1}) donnent un homomorphisme de
$\mathcal{S}_X$-Modules gradués
\begin{equation}
\label{II.8.14.2.4-fr}
  \lambda:\operatorname{Proj}(\mathcal{M})
  \otimes_{\mathcal{S}_X}\operatorname{Proj}(\mathcal{N})
  \to\operatorname{Proj}(\mathcal{M}\otimes_{\mathcal{S}}\mathcal{N}).
\tag{8.14.2.4}
\end{equation}
\end{env}

\oldpage[II]{198}
\begin{proposition}[8.14.3]
\label{II.8.14.3-fr}
L'homomorphisme (\hyperref[II.8.14.2.4-fr]{8.14.2.4}) est bijectif~; il en
est de même de (\hyperref[II.8.14.2.3-fr]{8.14.2.3}) lorsque le
$\mathcal{S}$-Module gradué $\mathcal{M}$ admet une présentation finie
(\hyperref[II.3.1.1-fr]{3.1.1}).
\end{proposition}

\begin{proof}
La question est évidemment locale sur $X$ et sur $Y$~; on peut donc supposer
$Y=\operatorname{Spec}(A)$ affine, $\mathcal{S}=\widetilde{S}$,
$\mathcal{M}=\widetilde{M}$, $\mathcal{N}=\widetilde{N}$, où $S$ est une
$A$-algèbre graduée à degrés positifs, $M$ et $N$ deux $S$-modules gradués.
Si $f$ est un élément homogène de $S_+$, les homomorphismes
(\hyperref[II.8.14.2.1-fr]{8.14.2.1}) et
(\hyperref[II.8.14.2.2-fr]{8.14.2.2}), restreints à l'ouvert affine
$D_+(f)$, correspondent aux homomorphismes canoniques
(\hyperref[II.2.5.11.1-fr]{2.5.11.1}) et
(\hyperref[II.2.5.12.1-fr]{2.5.12.1})~:
\begin{align*}
  M(m)_{(f)}\otimes_{S_{(f)}}N(n)_{(f)}
  &\to(M\otimes_S N)(m+n)_{(f)},\\
  (\operatorname{Hom}_S(M,N))(n-m)_{(f)}
  &\to\operatorname{Hom}_{S_{(f)}}(M(m)_{(f)},N(n)_{(f)}).
\end{align*}
Si on se reporte aux définitions de ces homomorphismes, on voit donc (compte
tenu de (\hyperref[II.8.2.9.1-fr]{8.2.9.1})) que la restriction de
(\hyperref[II.8.14.2.4-fr]{8.14.2.4}) à $D_+(f)$ correspond à
l'homomorphisme canonique
\[
  M_f\otimes_{S_f}N_f\to(M\otimes_S N)_f
\]
défini dans (\hyperref[0.1.3.4-fr]{0, 1.3.4}), et on sait que ce dernier est
un isomorphisme. De même, la restriction de
(\hyperref[II.8.14.2.3-fr]{8.14.2.3}) à $D_+(f)$ correspond à
l'homomorphisme canonique (\hyperref[0.1.3.5-fr]{0, 1.3.5})
\[
  (\operatorname{Hom}_S(M,N))_f\to
  \operatorname{Hom}_{S_f}(M_f,N_f)
\]
compte tenu de ce que, $M$ étant de type fini, le module
$\operatorname{Hom}_S(M,N)$, somme directe des sous-groupes formés des
homomorphismes \emph{homogènes} de $S$-modules
(\hyperref[II.2.1.2-fr]{2.1.2}), coïncide avec l'ensemble de \emph{tous} les
homomorphismes $M\to N$ de $S$-modules. L'hypothèse que $M$ admet une
présentation finie entraîne alors (\hyperref[0.1.3.5-fr]{0, 1.3.5}) que
l'homomorphisme canonique considéré est bien un isomorphisme.
\end{proof}

\begin{proposition}[8.14.4]
\label{II.8.14.4-fr}
Si $U$ est un ouvert quasi-compact de $X$, il existe un entier $d$ tel que
pour tout entier $n$ multiple de $d$, $\mathcal{O}_X(n)|U$ soit inversible,
son inverse étant $\mathcal{O}_X(-n)|U$.
\end{proposition}

\begin{proof}
Comme $q(U)$ est quasi-compact, il est recouvert par un nombre fini d'ouverts
affines $V_i$, donc tout $x\in U$ est contenu dans un ouvert affine de la
forme $D_+(f)$, où $f$ est un élément homogène de degré $>0$ de l'un des
anneaux $\Gamma(V_i,\mathcal{S})$. Comme $U$ est quasi-compact, on peut le
recouvrir par un nombre fini de tels ouverts $D_+(f_j)$~; soit $d$ un
multiple commun des degrés des $f_j$. Le nombre $d$ répond à la question en
vertu de (\hyperref[II.2.5.17-fr]{2.5.17}).
\end{proof}

\begin{env}[8.14.5]
\label{II.8.14.5-fr}
Avec les hypothèses et notations de (\hyperref[II.8.14.1-fr]{8.14.1}), on a
défini dans (\hyperref[II.3.3.2-fr]{3.3.2}) des homomorphismes canoniques de
$\mathcal{O}_Y$-Modules
\begin{equation}
\label{II.8.14.5.1-fr}
  \alpha_n:\mathcal{M}_n\to
  q_*(\operatorname{Proj}_0(\mathcal{M}(n)))
  \qquad(n\in\mathbf{Z})
\tag{8.14.5.1}
\end{equation}
Généralisant les notations de (\hyperref[II.3.3.1-fr]{3.3.1}), nous
poserons, pour tout $\mathcal{S}_X$-Module gradué $\mathcal{F}$,
\begin{equation}
\label{II.8.14.5.2-fr}
  \Gamma_*(\mathcal{F})=
  \bigoplus_{n\in\mathbf{Z}}q_*(\mathcal{F}_n).
\tag{8.14.5.2}
\end{equation}
En particulier,
$\Gamma_*(\mathcal{S}_X)=\bigoplus_{n\in\mathbf{Z}}q_*(\mathcal{O}_X(n))$
est la $\mathcal{O}_Y$-Algèbre graduée notée
$\Gamma_*(\mathcal{O}_X)$ dans
(\hyperref[II.3.3.1.2-fr]{3.3.1.2})~; il est clair que
$\Gamma_*(\mathcal{F})$ est un $\Gamma_*(\mathcal{S}_X)$-Module gradué
(\hyperref[0.4.2.2-fr]{0, 4.2.2}). Lorsque

\oldpage[II]{199}
dans les homomorphismes (\hyperref[II.8.14.5.1-fr]{8.14.5.1}), on prend
$\mathcal{M}=\mathcal{S}$, on obtient l'homomorphisme de
$\mathcal{O}_Y$-Algèbres graduées
\begin{equation}
\label{II.8.14.5.3-fr}
  \alpha:\mathcal{S}\to\Gamma_*(\mathcal{S}_X)
\tag{8.14.5.3}
\end{equation}
déjà défini dans (\hyperref[II.3.3.2-fr]{3.3.2}), et qui fait donc de
$\Gamma_*(\mathcal{F})$ un $\mathcal{S}$-Module gradué~; les homomorphismes
(\hyperref[II.8.14.5.1-fr]{8.14.5.1}) définissent alors un homomorphisme (de
degré $0$) de $\mathcal{S}$-Modules gradués
\begin{equation}
\label{II.8.14.5.4-fr}
  \alpha:\mathcal{M}\to\Gamma_*(\operatorname{Proj}(\mathcal{M})).
\tag{8.14.5.4}
\end{equation}
\end{env}

\begin{env}[8.14.6]
\label{II.8.14.6-fr}
En général, pour un $\mathcal{S}_X$-Module gradué quasi-cohérent
$\mathcal{F}$, il n'est pas certain que le $\mathcal{S}$-Module gradué
$\Gamma_*(\mathcal{F})$ soit nécessairement quasi-cohérent. Considérons un
ouvert $X'$ de $X$ tel que la restriction $q':X'\to Y$ de $q$ à $X'$ soit
un morphisme \emph{quasi-compact}. Comme $q'$ est en outre séparé,
$q'_*(\mathcal{F}')$ est alors un $\mathcal{O}_Y$-Module quasi-cohérent pour
tout $\mathcal{O}_{X'}$-Module quasi-cohérent $\mathcal{F}'$
(\hyperref[I.9.2.2-fr]{I, 9.2.2, b)}). Nous poserons
\begin{equation}
\label{II.8.14.6.1-fr}
  \mathcal{S}_{X'}=\mathcal{S}_X|X'
  =\bigoplus_{n\in\mathbf{Z}}\mathcal{O}_X(n)|X'
\tag{8.14.6.1}
\end{equation}
et, pour tout $\mathcal{S}_{X'}$-Module gradué $\mathcal{F}'$,
\begin{equation}
\label{II.8.14.6.2-fr}
  \Gamma'_*(\mathcal{F}')=
  \bigoplus_{n\in\mathbf{Z}}q'_*(\mathcal{F}'_n).
\tag{8.14.6.2}
\end{equation}
La remarque précédente montre alors que si $\mathcal{F}'$ est un
$\mathcal{S}_{X'}$-Module quasi-cohérent, $\Gamma'_*(\mathcal{F}')$ est un
$\mathcal{S}$-Module gradué quasi-cohérent
(\hyperref[I.9.6.1-fr]{I, 9.6.1}).

On notera d'ailleurs que l'injection canonique $j:X'\to X$ est
\emph{quasi-compacte} puisque $q'=q\circ j$ l'est et que $q$ est séparé
(\hyperref[I.6.6.4-fr]{I, 6.6.4, (v)}). Par suite
$\mathcal{F}=j_*(\mathcal{F}')$ est un $\mathcal{S}_X$-Module gradué
quasi-cohérent pour tout $\mathcal{S}_{X'}$-Module gradué quasi-cohérent
$\mathcal{F}'$, et il résulte des définitions précédentes que l'on a
\begin{equation}
\label{II.8.14.6.3-fr}
  \Gamma'_*(\mathcal{F}')=\Gamma_*(\mathcal{F}).
\tag{8.14.6.3}
\end{equation}
Avec les mêmes hypothèses sur $X'$, pour tout $\mathcal{S}$-Module gradué
quasi-cohérent $\mathcal{M}$, nous poserons
\begin{equation}
\label{II.8.14.6.4-fr}
  \operatorname{Proj}'(\mathcal{M})
  =\operatorname{Proj}(\mathcal{M})|X'
\tag{8.14.6.4}
\end{equation}
qui est un $\mathcal{S}_{X'}$-Module gradué quasi-cohérent. L'homomorphisme
canonique
\[
  \operatorname{Proj}(\mathcal{M})\to
  j_*(\operatorname{Proj}'(\mathcal{M}))
\]
(\hyperref[0.4.4.3-fr]{0, 4.4.3}) donne donc un homomorphisme canonique
$\Gamma_*(\operatorname{Proj}(\mathcal{M}))\to
\Gamma'_*(\operatorname{Proj}'(\mathcal{M}))$ de $\mathcal{S}$-Modules
gradués, et par suite, en le composant avec
(\hyperref[II.8.14.5.4-fr]{8.14.5.4}), on obtient un homomorphisme canonique
fonctoriel (de degré $0$) de $\mathcal{S}$-Modules gradués quasi-cohérents
\begin{equation}
\label{II.8.14.6.5-fr}
  \alpha':\mathcal{M}\to
  \Gamma'_*(\operatorname{Proj}'(\mathcal{M})).
\tag{8.14.6.5}
\end{equation}
\end{env}

\begin{env}[8.14.7]
\label{II.8.14.7-fr}
Conservons sur $X'$ les hypothèses de
(\hyperref[II.8.14.6-fr]{8.14.6}), et soit $\mathcal{F}'$ un
$\mathcal{S}_{X'}$-Module gradué quasi-cohérent, de sorte que
$\operatorname{Proj}'(\Gamma'_*(\mathcal{F}'))$ est aussi un
$\mathcal{S}_{X'}$-Module gradué quasi-cohérent. Nous allons définir un
homomorphisme canonique fonctoriel (de degré $0$) de
$\mathcal{S}_{X'}$-Modules gradués
\begin{equation}
\label{II.8.14.7.1-fr}
  \beta':\operatorname{Proj}'(\Gamma'_*(\mathcal{F}'))\to\mathcal{F}'.
\tag{8.14.7.1}
\end{equation}

\oldpage[II]{200}
Supposons d'abord $Y=\operatorname{Spec}(A)$ affine,
$\mathcal{S}=\widetilde{S}$, où $S$ est une $A$-algèbre graduée à degrés
positifs~; alors $\Gamma'_*(\mathcal{F}')=\widetilde{M}$, où
$M=\bigoplus_{n\in\mathbf{Z}}\Gamma(X',\mathcal{F}'_n)$ est un $S$-module
gradué. Soit $f\in S_d$ tel que $D_+(f)\subset X'$~; par définition
(\hyperref[II.2.6.2-fr]{2.6.2}), $\alpha_d(f)$ restreinte à $D_+(f)$ est la
section de $\mathcal{O}_X(d)$ au-dessus de $D_+(f)$ correspondant à
l'élément $f/1$ de $(S(d))_{(f)}$ et est par suite inversible~; il en est donc
de même de $\alpha_d(f^n)$ pour tout $n>0$. On en conclut aussitôt que
l'on définit un $S_f$-homomorphisme (de degré $0$) de modules gradués
$\beta_f:M_f\to\Gamma(D_+(f),\mathcal{F}')$ en faisant correspondre à tout
élément $z/f^n\in M_f$ (où $z\in M$), la section
$(z|D_+(f))(\alpha_d(f^n)|D_+(f))^{-1}$ de $\mathcal{F}'$ au-dessus de
$D_+(f)$. En outre, on a un diagramme commutatif correspondant à
(\hyperref[II.2.6.4.1-fr]{2.6.4.1}), d'où la définition de $\beta'$ dans ce
cas. Pour passer au cas général, il faut considérer une $A$-algèbre $A'$,
la $A'$-algèbre graduée $S'=S\otimes_A A'$, et utiliser un diagramme
commutatif analogue à (\hyperref[II.2.8.13.2-fr]{2.8.13.2})~; nous laissons
les détails aux lecteurs.
\end{env}

\begin{proposition}[8.14.8]
\label{II.8.14.8-fr}
Si $X'$ est un ouvert de $X=\operatorname{Proj}(\mathcal{S})$ tel que
$q':X'\to Y$ soit quasi-compact, l'homomorphisme $\beta'$ défini dans
(\hyperref[II.8.14.7-fr]{8.14.7}) est bijectif.
\end{proposition}

\begin{proof}
On peut évidemment se restreindre au cas où $Y$ est affine et tout revient à
prouver (avec les notations de (\hyperref[II.8.14.7-fr]{8.14.7})) que
l'homomorphisme $\beta_f:M_f\to\Gamma(D_+(f),\mathcal{F}')$ est un
isomorphisme. Or, lorsqu'on remplace $f$ par une de ses puissances, on ne
change pas $D_+(f)$ ni $\beta_f$~; comme $X'$ est \emph{quasi-compact} par
hypothèse, on peut donc toujours supposer, en vertu de
(\hyperref[II.8.14.4-fr]{8.14.4}), que le faisceau $\mathcal{O}_X(d)$ est
\emph{inversible}. Comme $X'$ est un schéma (puisque $q'$ est séparé), la
proposition n'est autre alors que (\hyperref[I.9.3.1-fr]{I, 9.3.1}).
\end{proof}

\begin{corollary}[8.14.9]
\label{II.8.14.9-fr}
Sous les hypothèses de (\hyperref[II.8.14.8-fr]{8.14.8}), tout
$\mathcal{S}_{X'}$-Module gradué quasi-cohérent $\mathcal{F}'$ est isomorphe
à un $\mathcal{S}_{X'}$-Module gradué de la forme
$\operatorname{Proj}'(\mathcal{M})$, où $\mathcal{M}$ est un
$\mathcal{S}$-Module gradué quasi-cohérent. Si en outre $\mathcal{F}'$ est de
type fini, et si on suppose que $Y$ est un schéma quasi-compact ou un
préschéma dont l'espace sous-jacent est noethérien, on peut supposer
$\mathcal{M}$ de type fini.
\end{corollary}

\begin{proof}
La démonstration à partir de (\hyperref[II.8.14.8-fr]{8.14.8}) suit
exactement la marche de celle de (\hyperref[II.3.4.5-fr]{3.4.5}) à partir de
(\hyperref[II.3.4.4-fr]{3.4.4}), et nous en laissons les détails au lecteur.
\end{proof}

\begin{proposition}[8.14.10]
\label{II.8.14.10-fr}
Sous les hypothèses de (\hyperref[II.8.14.7-fr]{8.14.7}), soient
$\mathcal{M}$ un $\mathcal{S}$-Module gradué quasi-cohérent,
$\mathcal{F}'$ un $\mathcal{S}_{X'}$-Module gradué quasi-cohérent~; les
homomorphismes composés
\begin{equation}
\label{II.8.14.10.1-fr}
  \operatorname{Proj}'(\mathcal{M})
  \xrightarrow{\operatorname{Proj}'(\alpha')}
  \operatorname{Proj}'(\Gamma'_*(\operatorname{Proj}'(\mathcal{M})))
  \xrightarrow{\beta'}\operatorname{Proj}'(\mathcal{M})
\tag{8.14.10.1}
\end{equation}
\begin{equation}
\label{II.8.14.10.2-fr}
  \Gamma'_*(\mathcal{F}')
  \xrightarrow{\alpha'}
  \Gamma'_*(\operatorname{Proj}'(\Gamma'_*(\mathcal{F}')))
  \xrightarrow{\Gamma'_*(\beta')}\Gamma'_*(\mathcal{F}')
\tag{8.14.10.2}
\end{equation}
sont les isomorphismes identiques.
\end{proposition}

\begin{proof}
La question est locale sur $Y$ et la vérification se fait comme dans
(\hyperref[II.2.6.5-fr]{2.6.5})~; nous en laissons les détails au lecteur.
\end{proof}

\begin{remark}[8.14.11]
\label{II.8.14.11-fr}
Au chapitre III (\hyperref[III.2.3.1-fr]{III, 2.3.1}), nous verrons que
lorsque $Y$ est \emph{localement noethérien} et $\mathcal{S}$ une
$\mathcal{O}_Y$-Algèbre graduée quasi-cohérente \emph{de type fini} (auquel
cas

\oldpage[II]{201}
on peut prendre $X'=X$), alors l'homomorphisme $\alpha$
(\hyperref[II.8.14.5.4-fr]{8.14.5.4}) est \textbf{(TN)}-\emph{bijectif} pour
tout $\mathcal{S}$-Module gradué quasi-cohérent $\mathcal{M}$ vérifiant la
condition \textbf{(TF)}.
\end{remark}

\begin{remark}[8.14.12]
\label{II.8.14.12-fr}
La situation décrite dans (\hyperref[II.8.14.4-fr]{8.14.4}) est un cas
particulier de la suivante. Soient $X$ un espace annelé, $\mathcal{S}$ une
$\mathcal{O}_X$-Algèbre graduée (à degrés positifs et négatifs)~; supposons
qu'il existe un entier $d>0$ tel que $\mathcal{S}_d$ et
$\mathcal{S}_{-d}$ soient \emph{inversibles}, l'homomorphisme canonique
\begin{equation}
\label{II.8.14.12.1-fr}
  \mathcal{S}_d\otimes_{\mathcal{O}_X}\mathcal{S}_{-d}
  \to\mathcal{O}_X
\tag{8.14.12.1}
\end{equation}
étant un \emph{isomorphisme} (de sorte que $\mathcal{S}_{-d}$ s'identifie à
$\mathcal{S}_d^{-1}$). On dit alors que la $\mathcal{O}_X$-Algèbre graduée
$\mathcal{S}$ est \emph{périodique de période $d$}. Cette dénomination
provient de la propriété suivante~: \emph{sous les hypothèses précédentes,
pour tout $\mathcal{S}$-Module gradué $\mathcal{F}$, l'homomorphisme
canonique}
\begin{equation}
\label{II.8.14.12.2-fr}
  \mathcal{S}_d\otimes\mathcal{F}_n\to\mathcal{F}_{n+d}
\tag{8.14.12.2}
\end{equation}
\emph{est un isomorphisme pour tout $n\in\mathbf{Z}$.} En effet, la question
est locale sur $X$ et on peut donc supposer que $\mathcal{S}_d$ possède une
section \emph{inversible} $s$ au-dessus de $X$, son inverse $s'$ étant une
section de $\mathcal{S}_{-d}$. L'homomorphisme
$\mathcal{F}_{n+d}\to\mathcal{S}_d\otimes\mathcal{F}_n$ qui, à toute section
$z\in\Gamma(U,\mathcal{F}_{n+d})$ fait correspondre la section
$(s|U)\otimes(s'|U)z$ de $\mathcal{S}_d\otimes\mathcal{F}_n$ sur $U$, est
alors réciproque de (\hyperref[II.8.14.12.2-fr]{8.14.12.2}), d'où notre
assertion. On en déduit pour tout $k\in\mathbf{Z}$ un isomorphisme canonique
\[
  (\mathcal{S}_d)^{\otimes k}\otimes\mathcal{F}_n
  \xrightarrow{\sim}\mathcal{F}_{n+kd}.
\]
Par suite, \emph{un $\mathcal{S}$-Module gradué $\mathcal{F}$ est connu
lorsqu'on connaît les $\mathcal{S}_0$-Modules $\mathcal{F}_i$
$(0\leqslant i\leqslant d-1)$ et les homomorphismes canoniques}
\[
  \mathcal{S}_i\otimes\mathcal{F}_j\to\mathcal{F}_{i+j}
  \qquad\text{pour }0\leqslant i,j\leqslant d-1
\]
(en posant
$\mathcal{F}_{i+j}=\mathcal{S}_d\otimes_{\mathcal{S}_0}
\mathcal{F}_{i+j-d}$ lorsque $i+j\geqslant d$). Bien entendu, pour que ces
homomorphismes définissent bien une structure de $\mathcal{S}$-Module sur la
somme directe des $(\mathcal{S}_d)^{\otimes k}\otimes\mathcal{F}_i$
$(k\in\mathbf{Z},\ 0\leqslant i\leqslant d-1)$, ils doivent satisfaire à des
conditions d'associativité que nous n'expliciterons pas.

Dans le cas où $d=1$ (qui est celui considéré dans
(\hyperref[subsection:II.3.3-fr]{3.3})) on peut donc dire que la catégorie des
$\mathcal{S}$-Modules gradués (resp. quasi-cohérents si $X$ est un préschéma
et $\mathcal{S}$ est quasi-cohérente) est \emph{équivalente} à la catégorie
des $\mathcal{S}_0$-Modules quelconques (resp. quasi-cohérents)~; c'est dans
ce sens qu'on peut considérer les développements de ce numéro comme
généralisant ceux du \S~\hyperref[section:II.3-fr]{3}. En outre, on voit que,
sous des conditions de finitude convenables, les résultats de ce numéro
(joints à (\hyperref[II.8.14.11-fr]{8.14.11})) ramènent dans une certaine
mesure l'étude des Algèbres graduées quasi-cohérentes sur un préschéma, et des
Modules gradués «~mod. \textbf{(TN)}~» sur de telles Algèbres, à l'étude du
cas particulier où les Algèbres considérées sont \emph{périodiques} (et où
la condition \textbf{(TN)} pour $\mathcal{M}$
(\hyperref[II.3.4.2-fr]{3.4.2}) implique donc $\mathcal{M}=0$).
\end{remark}

\begin{remark}[8.14.13]
\label{II.8.14.13-fr}
Sous les hypothèses de (\hyperref[II.8.14.1-fr]{8.14.1}), soit $d$ un entier
$>0$~; on a défini un $Y$-isomorphisme canonique $h$ de $X$ sur
$X^{(d)}=\operatorname{Proj}(\mathcal{S}^{(d)})$
(\hyperref[II.3.1.8-fr]{3.1.8}). Pour tout

\oldpage[II]{202}
$\mathcal{S}$-Module gradué quasi-cohérent $\mathcal{M}$ et tout entier $k$
tel que $0\leqslant k\leqslant d-1$, on a aussi un $h$-isomorphisme canonique
(avec les notations de (\hyperref[II.3.1.1-fr]{3.1.1}))
\begin{equation}
\label{II.8.14.13.1-fr}
  (\operatorname{Proj}(\mathcal{M}))^{(d,k)}
  \xleftarrow{\sim}\operatorname{Proj}(\mathcal{M}^{(d,k)}).
\tag{8.14.13.1}
\end{equation}

Supposons d'abord $Y=\operatorname{Spec}(A)$ affine,
$\mathcal{S}=\widetilde{S}$, $\mathcal{M}=\widetilde{M}$, où $S$ est une
$A$-algèbre graduée à degrés positifs et $M$ un $S$-module gradué. On sait
que pour tout $f\in S_e$ ($e>0$) $h$ applique $D_+(f)$ sur $D_+(f^d)$ et
correspond à l'isomorphisme canonique
$S_{(f^d)}\xrightarrow{\sim}S_{(f)}$
(\hyperref[II.2.2.2-fr]{2.2.2}). La restriction de
(\hyperref[II.8.14.13.1-fr]{8.14.13.1}) à $D_+(f^d)$ correspond alors au
di-isomorphisme canonique $M_{f^d}\xrightarrow{\sim}M_f$, restreint aux
éléments de $M_{f^d}$ de degré congru à $k$ (mod. $d$). Nous laissons au
lecteur le soin de vérifier que ces isomorphismes sont compatibles avec le
passage de $f$ à un de ses multiples homogènes $fg$, puis qu'on a une
compatibilité analogue avec le passage de $S$ à une $A'$-algèbre graduée
$S'=S\otimes_A A'$, lorsque $A'$ est une $A$-algèbre.

En particulier, on a ainsi un $h$-isomorphisme
\begin{equation}
\label{II.8.14.13.2-fr}
  (\mathcal{S}^{(d)})_{X^{(d)}}
  \xrightarrow{\sim}(\mathcal{S}_X)^{(d)}
\tag{8.14.13.2}
\end{equation}
qui respecte les structures multiplicatives des deux membres, et grâce auquel
(\hyperref[II.8.14.13.1-fr]{8.14.13.1}) devient un $h$-di-isomorphisme d'un
$(\mathcal{S}^{(d)})_{X^{(d)}}$-Module gradué sur un
$(\mathcal{S}_X)^{(d)}$-Module gradué. De même, on a un $h$-isomorphisme
\begin{equation}
\label{II.8.14.13.3-fr}
  \operatorname{Proj}_0(\mathcal{S}^{(d,k)}(n))
  \xrightarrow{\sim}\mathcal{O}_X(nd+k)
\tag{8.14.13.3}
\end{equation}
qui complète le résultat de (\hyperref[II.3.2.9-fr]{3.2.9, (ii)}).

On déduit immédiatement de l'isomorphisme
(\hyperref[II.8.14.13.1-fr]{8.14.13.1}) un isomorphisme de
$\mathcal{S}^{(d)}$-Modules gradués
\begin{equation}
\label{II.8.14.13.4-fr}
  \Gamma_*^{(d)}(\operatorname{Proj}(\mathcal{M}^{(d,k)}))
  \xrightarrow{\sim}
  \Gamma_*((\operatorname{Proj}(\mathcal{M}))^{(d,k)})
\tag{8.14.13.4}
\end{equation}
où $\Gamma_*^{(d)}$ correspond au morphisme structural
$q^{(d)}:X^{(d)}\to Y$~; il est immédiat de vérifier que l'homomorphisme
canonique $\alpha$ (\hyperref[II.8.14.5.4-fr]{8.14.5.4}) et l'homomorphisme
analogue $\alpha^{(d)}$ pour $X^{(d)}$ rendent commutatif le diagramme
\begin{equation}
\label{II.8.14.13.5-fr}
  \xymatrix{
    & \mathcal{M}^{(d,k)}\ar[dl]_{\alpha^{(d)}}\ar[dr]^{\alpha} & \\
    \Gamma_*^{(d)}(\operatorname{Proj}(\mathcal{M}^{(d,k)}))
      \ar[rr]^{\sim}
      && \Gamma_*((\operatorname{Proj}(\mathcal{M}))^{(d,k)})
  }
\tag{8.14.13.5}
\end{equation}
la vérification se faisant en supposant $Y$ affine, et calculant les
restrictions des images par $\alpha^{(d)}$ et $\alpha$ d'un même élément de
$M^{(d,k)}$ aux ouverts $D_+(f^d)$ et $D_+(f)$ (avec les mêmes notations que
ci-dessus).
\end{remark}

\begin{proposition}[8.14.14]
\label{II.8.14.14-fr}
Soient $Y$ un préschéma quasi-compact, $\mathcal{S}$ une
$\mathcal{O}_Y$-Algèbre graduée quasi-cohérente de type fini, $\mathcal{M}$
un $\mathcal{S}$-Module gradué quasi-cohérent vérifiant la condition
\textbf{(TF)}~; soit $X=\operatorname{Proj}(\mathcal{S})$. Alors
$\mathcal{S}_X$ est une $\mathcal{O}_X$-Algèbre graduée périodique
(\hyperref[II.8.14.12-fr]{8.14.12}), et il existe une

\oldpage[II]{203}
période $d$ de $\mathcal{S}_X$ telle que les
$(\operatorname{Proj}(\mathcal{M}))^{(d,k)}$
$(0\leqslant k\leqslant d-1)$ soient des
$(\mathcal{S}_X)^{(d)}$-Modules de type fini.
\end{proposition}

\begin{proof}
En effet, (\hyperref[II.3.1.10-fr]{3.1.10}) prouve qu'il existe $d$ tel que
$\mathcal{S}^{(d)}$ soit engendrée par
$\mathcal{S}_d=(\mathcal{S}^{(d)})_1$, ce dernier étant un
$\mathcal{S}_0$-Module de type fini. Pour démontrer la première assertion, on
peut donc, en vertu de (\hyperref[II.8.14.13.2-fr]{8.14.13.2}), se borner au
cas où $d=1$, et alors la proposition résulte de
(\hyperref[II.3.2.7-fr]{3.2.7}). En outre, compte tenu de
(\hyperref[II.8.14.13.1-fr]{8.14.13.1}), la seconde assertion est conséquence
de (\hyperref[II.2.1.6-fr]{2.1.6, (iii)}) et de
(\hyperref[II.3.4.3-fr]{3.4.3}).
\end{proof}
